% =====================================================================
% main.tex – Soongsil Univ. PhD Thesis (DRTO v2)
% !TeX program = xelatex
% 빌드: latexmk -xelatex main.tex
% =====================================================================
\documentclass[11pt,a4paper]{report}

% ---------------- 기본 패키지 ----------------
\usepackage[a4paper,top=30mm,bottom=30mm,left=25mm,right=25mm]{geometry}
\setlength{\headheight}{13.6pt}
\addtolength{\topmargin}{-1.6pt}

% ---------------- 페이지 구분선 (상·하단 가로줄) ----------------
\usepackage{fancyhdr}
\fancypagestyle{borderstyle}{%
  \fancyhf{}%
  \fancyhead[C]{\rule{\textwidth}{0.4pt}}%
  \fancyfoot[C]{\rule{\textwidth}{0.4pt}\\\vspace{2pt}\thepage}%
  \renewcommand{\headrulewidth}{0pt}%
  \renewcommand{\footrulewidth}{0pt}%
}
\pagestyle{borderstyle}
% chapter 첫 페이지(plain)에도 동일 적용
\fancypagestyle{plain}{%
  \fancyhf{}%
  \fancyhead[C]{\rule{\textwidth}{0.4pt}}%
  \fancyfoot[C]{\rule{\textwidth}{0.4pt}\\\vspace{2pt}\thepage}%
  \renewcommand{\headrulewidth}{0pt}%
  \renewcommand{\footrulewidth}{0pt}%
}

% ---------------- 폰트 ----------------
% === Engine / Font ===
\usepackage{fontspec}
\usepackage{xeCJK}
\xeCJKsetup{CJKspace=true}

% === 한글 폰트 ===
% 본문: 명조
\setCJKmainfont[
  UprightFont = *,
  BoldFont    = * Bold
]{Noto Serif CJK KR}

% 중고딕: 고딕
\setCJKsansfont[
  UprightFont = *,
  BoldFont    = * Bold
]{Noto Sans CJK KR}

% === 줄간격 (가이드라인) ===
\usepackage{setspace}

% ---------------- 기타 패키지 ----------------
\usepackage{graphicx}
\usepackage{amsmath,amssymb}
\usepackage{caption,subcaption}
\usepackage{booktabs,multirow,tabularx,array,longtable}
\usepackage{threeparttable}
\usepackage{float}
\usepackage{placeins}

% --- 플로팅 환경(표/그림) 배치 최적화: 지면 절약 ---
\renewcommand{\topfraction}{0.9}
\renewcommand{\bottomfraction}{0.9}
\renewcommand{\textfraction}{0.1}
\renewcommand{\floatpagefraction}{0.7}
\setcounter{topnumber}{4}
\setcounter{bottomnumber}{4}
\setcounter{totalnumber}{8}
\setcounter{dbltopnumber}{4}

% 기본 플로팅 배치 옵션: htbp
\makeatletter
\renewcommand{\fps@figure}{htbp}
\renewcommand{\fps@table}{htbp}
\makeatother
\usepackage{titlesec}
\usepackage{tocloft}
\usepackage{indentfirst}
\usepackage[inline]{enumitem}
\usepackage{ragged2e}
\usepackage{microtype}
\usepackage{pifont}
\usepackage{makecell}
\usepackage{algorithm}
\usepackage{algpseudocode}
\usepackage{tcolorbox}
\tcbuselibrary{breakable}
\usepackage{colortbl}
\usepackage{amsthm}
\usepackage{listings}
\usepackage[normalem]{ulem}            % \sout 취소선 (4단계: CVR 삭제 표시용)

% ── 참고문헌 리뷰 모드: \refchg 만 빨강, \vrev / \color{red} 는 검정 ──
\definecolor{v2red}{rgb}{0,0,0}       % 1단계 완료: v2red 비활성
\definecolor{p2red}{rgb}{0,0,0}       % 2·3단계 완료: p2red 비활성 (검정)
\colorlet{red}{black}                  % 기존 \color{red} → 검은색으로 무효화
% 4단계 (CVR 삭제 + 마크다운 별표 정리) 변경 표시: 검정 환원 + 삭제 완전 제거 ─
\definecolor{p4red}{rgb}{0,0,0}       % 4단계 완료: p4red 비활성 (검정)
\newcommand{\pdel}[1]{}               % 4단계 완료: 삭제 내용 완전 제거
\newcommand{\padd}[1]{#1}             % 4단계 완료: 추가 내용 평문화
% (본문/부록 \pdel{...}은 모두 사라지고, \padd{...}는 평문으로 출력됨)
% 5단계 5-1단계: 출처 마커 (5-2 진입으로 비활성화 — 모두 사라짐)
\definecolor{srcred}{rgb}{0,0,0}      % srcred 비활성 (검정)
\newcommand{\src}[1]{}                % 비활성화: 출처 마커 제거
% 5단계 완료: 검정 환원 + 삭제 완전 제거 ──────────────────────────────
\definecolor{p5red}{rgb}{0,0,0}                         % p5red 비활성 (검정)
\newcommand{\fdel}[1]{}                                  % 삭제: 완전 제거
\newcommand{\fadd}[1]{#1}                                % 추가: 평문 출력 (검정)
\newcommand{\fmod}[2]{#2}                                % 수정: 새 내용만 평문
\newenvironment{p5del}[1]{}{}                            % no-op
\newenvironment{p5add}{}{}                               % no-op (평문 출력)
% 검토 마크업 완료: 검정 환원 + 삭제 완전 제거 ────────────────────────
\definecolor{revred}{rgb}{0,0,0}                         % revred 비활성 (검정)
\newcommand{\rev}[1]{}                                   % 비활성
\newcommand{\revdel}[1]{}                                % 삭제 완전 제거
\newcommand{\revadd}[1]{#1}                              % 추가 평문화
\newenvironment{revnote}{}{}                             % no-op
\newenvironment{revdelblock}[1]{}{}                      % no-op
% 이전 임시 표시 (검정 환원 완료) ────────────────────────────────
\definecolor{tmpred}{rgb}{0,0,0}                         % tmpred 비활성 (검정)
\newcommand{\tmp}[1]{#1}                                 % 평문화
% 이번 세션 변경분 표시 (미참조 표 참조 추가 완료: 검정 환원) ─────
\definecolor{nref}{rgb}{0,0,0}                           % nref 비활성 (검정)
\newcommand{\nref}[1]{#1}                                % 평문화
% 김영철 교수 대응 (v3.1) 변경분 표시 — 검정 환원 ──────────────────
\definecolor{kimred}{rgb}{0,0,0}                         % kimred 비활성 (검정)
\newcommand{\kim}[1]{#1}                                 % 평문화
\newcommand{\kcitep}[1]{\citep{#1}}                      % 평문
\newcommand{\kcitet}[1]{\citet{#1}}                      % 평문
% 한글본 v3 정합 작업 마킹 — 검정 환원 (한글본 body와 일치 확인 완료) ──
\definecolor{syncr}{rgb}{0,0,0}                          % syncr 비활성 (검정)
\newcommand{\syncr}[1]{#1}                               % 평문화
% J.4 수정 — 3차: 검정 환원 (기존 빨간 표시 제거) ──
\definecolor{jrev}{rgb}{0,0,0}                           % 3차: jrev 비활성 (검정)
\newcommand{\jrev}[1]{#1}                                % 3차: 평문화
\newcommand{\syncp}[1]{\citep{#1}}                       % 평문
\newcommand{\synct}[1]{\citet{#1}}                       % 평문
\newcommand{\vrev}[1]{#1}  % 3차: 검정으로 변환  % 1단계 리뷰 빨강 (toc/lot 안전)
% 3차 한글 작업 완료분 — 검정으로 환원 (사용자 한글본 작업 완료)
\definecolor{nrevred}{rgb}{0,0,0}                        % 3차: 검정으로 환원
\newcommand{\nrev}[1]{#1}                                % 3차: 평문화
% 3차 신규 추가 변경 (가설 보완안 + C.2.2 상세화) — 빨강
\definecolor{mrevred}{rgb}{0,0,0}                        % 3차: 검정 환원
\newcommand{\mrev}[1]{#1}                                % 3차: 평문화
% C.2.2 전용 빨강 마커 (3차 심사 완료 후 검정 환원: 2026-05-09)
\definecolor{prevred}{rgb}{0,0,0}                         % 3차 심사 후 검정 환원
\newcommand{\prev}[1]{#1}                                 % 평문화
% ── 3차 심사 결과 반영분 (검정 환원 완료: 2026-06-01, 4차 작업 진입) ──
\definecolor{trevred}{rgb}{0,0,0}                         % 4차 진입: trevred 비활성 (검정)
\newcommand{\tnew}[1]{#1}                                 % 4차: 평문화
\newcommand{\tdel}[1]{}                                   % 삭제 내용 (제거)
\newcommand{\tmod}[2]{#2}                                 % 4차: 새 내용만 평문
\newcommand{\tnewp}[1]{#1}                                % 4차: 평문화
\newcommand{\trefcol}[1]{\ref{#1}}                        % 4차: 평문 \ref
% ── ICC 정정 라운드 (검정 환원 완료: 2026-06-01, 4차 작업 진입) ──
\definecolor{iccred}{rgb}{0,0,0}                          % 4차 진입: iccred 비활성 (검정)
\newcommand{\inew}[1]{#1}                                 % 4차: 평문화
\newcommand{\idel}[1]{}                                   % 삭제 (제거)
\newcommand{\imod}[2]{#2}                                 % 4차: 새 내용만 평문
% ── 4차(최종) 심사 대비 수정분 (2026-06-01 신설) — 이번 라운드만 빨간색 ──
\definecolor{frevred}{rgb}{0,0,0}                         % 4차 빨강 비활성 (전체 검정 환원)
\newcommand{\frev}[1]{{\color{frevred}#1}}                % 단어 수준 수정 (빨강)
\newcommand{\fpara}[1]{#1}                                % 신규 추가 단락 (검정, 식별만)
\newcommand{\fdelf}[1]{}                                  % 삭제 (제거)
\newcommand{\fmodf}[2]{{\color{frevred}#2}}              % 수정 (새 내용만 빨강)
\newcommand{\frevp}[1]{\begingroup\color{frevred}#1\endgroup}  % 단락용
\newcommand{\frefcol}[1]{\begingroup\hypersetup{linkcolor=frevred}\color{frevred}\ref{#1}\endgroup}
% ── 4차(최종) 심사 반영: 단계 완료 — 검정 환원 (2026-06-11, 최종제출본 정합 라운드 진입) ──
\definecolor{stepred}{rgb}{0,0,0}                % 검정 환원
\newcommand{\chg}[1]{#1}                          % 평문화
\newcommand{\del}[1]{}                            % 삭제 완전 제거
% ── 최종제출본7 정합 라운드: 잔존 오류 표시 활성 (2026-06-12) ──
\definecolor{subred}{rgb}{0,0,0}                  % 교수수정본(0622) 정합: 전체 검정 환원
\newcommand{\pdferr}[1]{{\color{subred}#1}}      % 제출본 PDF 오류 정정 (빨강)
% 참조형 오류 정정 (내부 ref 매크로의 p2red(검정)에 덮이지 않도록 전용 정의)
\newcommand{\pdferrsubsub}[1]{\begingroup\hypersetup{linkcolor=subred}\color{subred}\ref{#1}목\endgroup}
\newcommand{\pdferrsec}[1]{\begingroup\hypersetup{linkcolor=subred}\color{subred}\ref{#1}절\endgroup}
\newcommand{\pdferrappch}[1]{\begingroup\hypersetup{linkcolor=subred}\color{subred}부록~\ref{#1}장\endgroup}
\newcommand{\pdferrR}[1]{\begingroup\hypersetup{linkcolor=subred}\color{subred}#1\endgroup} % HWP 참조오류 빨강(제출본12)
% ── 14버전(0622): 지도교수 의도를 제출본13에 적용 — 수정 대상 빨강 마킹 ──
\definecolor{m14}{rgb}{0,0,0}                  % 공개본: \m 검정 환원
\newcommand{\m}[1]{{\color{m14}#1}}               % 교수 의도 반영(추가/교체) — 빨강
% ── 정오(errata) 토글: #1=공개본 오류형, #2=정정형 ──
% main_doc.tex(오류본)은 #1을, main_doc_revised.tex(수정본)은 #2를 출력
\ifdefined\ERRATAFIX\newcommand{\erf}[2]{#2}\else\newcommand{\erf}[2]{#1}\fi
\newcommand{\mdel}[1]{{\color{m14}\sout{#1}}}     % 삭제 대상 — 빨강 취소선
% 4차 진입: trevred 기반 helper도 검정 환원 (frevred 4차 신설은 위 참조)
\newcommand{\tsecref}[1]{\ref{#1}절}
\newcommand{\tsubsecref}[1]{\ref{#1}항}
\newcommand{\tchref}[1]{제\ref{#1}장}
\newcommand{\ttabref}[1]{[표~\ref{#1}]}
\newcommand{\tfigref}[1]{[그림~\ref{#1}]}
% ref/autoref/eqref 전용: hyperref linkcolor까지 빨강 전환 (본문에서만 사용)
\newcommand{\vref}[1]{\begingroup\hypersetup{linkcolor=v2red}\color{v2red}\ref{#1}\endgroup}
\newcommand{\vautoref}[1]{\begingroup\hypersetup{linkcolor=v2red}\color{v2red}\autoref{#1}\endgroup}
\newcommand{\veqref}[1]{\begingroup\hypersetup{linkcolor=v2red}\color{v2red}\eqref{#1}\endgroup}
% 2단계 변경 표시 매크로
\newcommand{\p}[1]{{\color{p2red}#1}}
\newcommand{\pref}[1]{\begingroup\hypersetup{linkcolor=p2red}\color{p2red}\ref{#1}\endgroup}
\newcommand{\pautoref}[1]{\begingroup\hypersetup{linkcolor=p2red}\color{p2red}\autoref{#1}\endgroup}
% 3단계 통합 참조 매크로 (장·절·항·목·식 표기 규칙)
\newcommand{\chref}[1]{\begingroup\hypersetup{linkcolor=p2red}\color{p2red}제\ref{#1}장\endgroup}
\newcommand{\secref}[1]{\begingroup\hypersetup{linkcolor=p2red}\color{p2red}\ref{#1}절\endgroup}
\newcommand{\subsecref}[1]{\begingroup\hypersetup{linkcolor=p2red}\color{p2red}\ref{#1}항\endgroup}
\newcommand{\subsubsecref}[1]{\begingroup\hypersetup{linkcolor=p2red}\color{p2red}\ref{#1}목\endgroup}
\newcommand{\appref}[1]{\begingroup\hypersetup{linkcolor=p2red}\color{p2red}\ref{#1}\endgroup}
\newcommand{\appchref}[1]{\begingroup\hypersetup{linkcolor=p2red}\color{p2red}부록~\ref{#1}장\endgroup}
\newcommand{\appsecref}[1]{\begingroup\hypersetup{linkcolor=p2red}\color{p2red}부록~\ref{#1}절\endgroup}
\newcommand{\appsubsecref}[1]{\begingroup\hypersetup{linkcolor=p2red}\color{p2red}부록~\ref{#1}항\endgroup}
% 표/그림 참조 (대괄호 + 빨강)
\newcommand{\tabref}[1]{\begingroup\hypersetup{linkcolor=p2red}\color{p2red}[표~\ref{#1}]\endgroup}
\newcommand{\figref}[1]{\begingroup\hypersetup{linkcolor=p2red}\color{p2red}[그림~\ref{#1}]\endgroup}
% 수식 참조: \eqref{X} → (식 X.Y) 형식 (hyperref linkcolor도 빨강)
\renewcommand{\eqref}[1]{\begingroup\hypersetup{linkcolor=p2red}\color{p2red}(식~\ref{#1})\endgroup}
% 한국어 조사 자동 처리: \refpo{label}{받침있음조사}{받침없음조사}
% xstring으로 참조 번호 마지막 문자를 추출해 조사 분기
% 받침 있음: 끝자리 0,1,3,6,7,8 → 은/이/을
% 받침 없음: 끝자리 2,4,5,9   → 는/가/를
\usepackage{xstring}
\usepackage{refcount}
\makeatletter
\newcommand{\refpo}[3]{%
  \begingroup
  \edef\refpo@tmp{\getrefnumber{#1}}%
  \expandafter\StrRight\expandafter{\refpo@tmp}{1}[\refpo@last]%
  \IfEqCase{\refpo@last}{%
    {0}{#2}{1}{#2}{2}{#3}{3}{#2}{4}{#3}%
    {5}{#3}{6}{#2}{7}{#2}{8}{#2}{9}{#3}%
  }[#3]%
  \endgroup
}
\makeatother
\newcommand{\reFneun}[1]{\refpo{#1}{은}{는}}
\newcommand{\reFyi}[1]{\refpo{#1}{이}{가}}
\newcommand{\reFeul}[1]{\refpo{#1}{을}{를}}
\newcommand{\refchg}[1]{#1}  % 참고문헌 표시 비활성 (1단계에서는 구조 변경만 강조)

% Theorem-like environments
\newtheorem{theorem}{Theorem}[chapter]
\newtheorem{proposition}[theorem]{\p{명제}}
\newtheorem{lemma}[theorem]{Lemma}
\newtheorem{corollary}[theorem]{Corollary}
\theoremstyle{definition}
\newtheorem{definition}{Definition}[section]
\newtheorem{hypothesis}{\p{가설 H}}  % 2단계 #11: "Hypothesis X.Y.Z" → "가설 HN" (H1..H11 연속 번호)
\newtheorem{testspec}{Test Specification}[section]

% --- Figures (TikZ) ---
\usepackage{tikz}
\usetikzlibrary{arrows.meta, positioning, fit, calc, shapes.geometric, backgrounds, shadows}

% --- Plots (pgfplots) ---
\usepackage{pgfplots}
\pgfplotsset{compat=1.18}

% 4단계: TikZ 그림 내부 기본 텍스트 색 = 검정
\tikzset{every picture/.append style={execute at begin picture={\color{black}}}}
\tikzset{
  box/.style={draw, rounded corners=2pt, align=center, inner sep=6pt, minimum width=28mm, minimum height=8mm},
  sbox/.style={draw, rounded corners=2pt, align=center, inner sep=5pt, minimum width=24mm, minimum height=7mm},
  gbox/.style={draw, rounded corners=2pt, align=left, inner sep=7pt, text width=0.82\linewidth},
  arrow/.style={-{Stealth[length=2.2mm]}, thick},
  dashedarrow/.style={-{Stealth[length=2.2mm]}, thick, dashed},
  group/.style={draw, rounded corners=3pt, inner sep=6pt},
  note/.style={align=left, inner sep=2pt}
}

\setmainfont{TeX Gyre Termes}     % 영문
\setCJKmonofont[Script=Hangul]{Noto Sans CJK KR}

\newfontfamily\ChapterFont{Noto Serif CJK KR}[BoldFont=Noto Serif CJK KR Bold]
\newfontfamily\JungGothic{Noto Sans CJK KR}[BoldFont=Noto Sans CJK KR Bold]
\newfontfamily\TOCMincho{Noto Serif CJK KR}
\newcommand{\mathkr}[1]{\text{\CJKfontspec{Noto Serif CJK KR}{#1}}}

\usepackage{unicode-math}
\setmathfont{TeX Gyre Termes Math}
\newcommand{\cmark}{{\JungGothic ✓}}
\newcommand{\circnum}[1]{{\JungGothic #1}}
% 원문자(①②③…)가 라틴 폰트에 없어 깨지는 문제 → CJK 폰트(Noto)로 자동 라우팅
\usepackage{newunicodechar}
\newunicodechar{①}{{\JungGothic\char"2460}}
\newunicodechar{②}{{\JungGothic\char"2461}}
\newunicodechar{③}{{\JungGothic\char"2462}}
\newunicodechar{④}{{\JungGothic\char"2463}}
\newunicodechar{⑤}{{\JungGothic\char"2464}}
\newunicodechar{⑥}{{\JungGothic\char"2465}}
\newunicodechar{⑦}{{\JungGothic\char"2466}}
\newunicodechar{⑧}{{\JungGothic\char"2467}}
\newunicodechar{⑨}{{\JungGothic\char"2468}}
\newunicodechar{⑩}{{\JungGothic\char"2469}}
\newunicodechar{⑪}{{\JungGothic\char"246A}}
\newunicodechar{⑫}{{\JungGothic\char"246B}}
\newunicodechar{⑬}{{\JungGothic\char"246C}}
\newunicodechar{⑭}{{\JungGothic\char"246D}}
\newunicodechar{⑮}{{\JungGothic\char"246E}}
% DRTO 표기 통일 매크로
\newcommand{\DRTO}{\mathrm{DRTO}}

\newcommand{\IACS}{산업자동화제어시스템(IACS: Industrial Automation and Control Systems)}
\newenvironment{tablecellblock}{%
  \begin{minipage}[t]{\linewidth}
  \setlength{\parindent}{0pt}%
  \setlength{\parskip}{2pt}}%
  {\end{minipage}}
\newenvironment{tablecellitemize}{%
  \begin{itemize}[leftmargin=1.2em,itemsep=2pt,topsep=0pt,partopsep=0pt,parsep=0pt]}%
  {\end{itemize}}

% ---------------- 한글 띄어쓰기/줄바꿈 ----------------
\spaceskip=1.20\fontdimen2\font plus 0.25\fontdimen3\font minus 0.25\fontdimen4\font
\AtBeginDocument{%
  \spaceskip=1.6\fontdimen2\font plus 0.35\fontdimen3\font minus 0.35\fontdimen4\font
}
\setlength{\xspaceskip}{0.8em plus 0.25em minus 0.1em}
\XeTeXlinebreaklocale "ko"
\XeTeXlinebreakskip = 0.5em plus 0.2em minus 0.1em

% === 줄간격 (가이드라인) ===
\setstretch{2.0}
\setlength{\parskip}{0pt}
\setlength{\parindent}{2em}
\hyphenpenalty=5000
\exhyphenpenalty=5000
\sloppy
\hfuzz=200pt
\vfuzz=50pt
\hbadness=10000
\vbadness=10000
\setlength{\emergencystretch}{3em}
\newcommand{\CurrentAppendixLetter}{}
\newcounter{appsection}
\newcounter{appsubsection}[appsection]
\newcounter{appsubsubsection}[appsubsection]
% 부록 다중 장(A·B·C)에서 하이퍼앵커 중복 방지
\makeatletter
\newcommand{\theHappsection}{\CurrentAppendixLetter.\arabic{appsection}}
\newcommand{\theHappsubsection}{\CurrentAppendixLetter.\arabic{appsection}.\arabic{appsubsection}}
\newcommand{\theHappsubsubsection}{\CurrentAppendixLetter.\arabic{appsection}.\arabic{appsubsection}.\arabic{appsubsubsection}}
\makeatother
\newcommand{\SetAppendixLetter}[1]{%
  \renewcommand{\CurrentAppendixLetter}{#1}%
  \setcounter{appsection}{0}%
  \setcounter{appsubsection}{0}%
  \setcounter{appsubsubsection}{0}}
\makeatletter
\newcommand{\AppSection}[1]{%
  \setcounter{appsubsection}{0}%
  \setcounter{appsubsubsection}{0}%
  \refstepcounter{appsection}%
  \def\@currentlabel{\CurrentAppendixLetter.\arabic{appsection}}%
  \phantomsection
  % [제출본13 정합] 부록 하위는 목차 미등재
  {\par\noindent\hspace{2em}\JungGothic\bfseries\fontsize{12}{16}\selectfont
  \CurrentAppendixLetter.\arabic{appsection}\ #1\par}\vspace{6pt}}
\newcommand{\AppSubsection}[1]{%
  \setcounter{appsubsubsection}{0}%
  \refstepcounter{appsubsection}%
  \def\@currentlabel{\CurrentAppendixLetter.\arabic{appsection}.\arabic{appsubsection}}%
  \phantomsection
  % [제출본13 정합] 부록 하위는 목차 미등재
  {\par\noindent\hspace{3em}\JungGothic\bfseries\fontsize{11}{14}\selectfont
  \CurrentAppendixLetter.\arabic{appsection}.\arabic{appsubsection}\ #1\par}\vspace{4pt}}
\newcommand{\AppSubsubsection}[1]{%
  \refstepcounter{appsubsubsection}%
  \def\@currentlabel{\CurrentAppendixLetter.\arabic{appsection}.\arabic{appsubsection}.\arabic{appsubsubsection}}%
  \phantomsection
  {\par\noindent\hspace{4em}\JungGothic\fontsize{11}{14}\selectfont
  \CurrentAppendixLetter.\arabic{appsection}.\arabic{appsubsection}.\arabic{appsubsubsection}\ #1\par}\vspace{3pt}}
\makeatother
\DeclareRobustCommand{\PlainTOCtext}[1]{%
  {\TOCMincho\fontsize{11}{14}\selectfont #1}%
}


% =====================================================================
%  목차 / 표목차 / 그림목차 설정
% =====================================================================
\renewcommand{\contentsname}{}
\renewcommand{\listfigurename}{}
\renewcommand{\listtablename}{}

\renewcommand{\cftaftertoctitle}{\par\vskip 0.35cm}
\renewcommand{\cftafterloftitle}{\par\vskip 0.35cm}
\renewcommand{\cftafterlottitle}{\par\vskip 0.35cm}

\newcommand{\FrontListTitle}[1]{%
  \begin{center}
    {\JungGothic\bfseries\fontsize{16}{22}\selectfont #1}
  \end{center}
  \vspace{0.35cm}
}

% 목차 항목 폰트/간격 설정
\renewcommand{\cftchapfont}{\TOCMincho\bfseries\fontsize{13}{18}\selectfont}
\renewcommand{\cftchappagefont}{\TOCMincho\bfseries\fontsize{13}{18}\selectfont}
\renewcommand{\cftsecfont}{\JungGothic\fontsize{11}{15}\selectfont}
\renewcommand{\cftsecpagefont}{\JungGothic\fontsize{11}{15}\selectfont}
\renewcommand{\cftsubsecfont}{\JungGothic\fontsize{11}{15}\selectfont}
\renewcommand{\cftsubsecpagefont}{\JungGothic\fontsize{11}{15}\selectfont}
\renewcommand{\cftsubsubsecfont}{\JungGothic\fontsize{11}{15}\selectfont}
\renewcommand{\cftsubsubsecpagefont}{\JungGothic\fontsize{11}{15}\selectfont}

\setlength{\cftbeforechapskip}{0.8cm}
\setlength{\cftbeforesecskip}{0pt}
\setlength{\cftbeforesubsecskip}{0pt}

\renewcommand{\cftchappresnum}{제}
\renewcommand{\cftchapaftersnum}{장\ }

\renewcommand{\cftchapfont}{\ChapterFont\bfseries\fontsize{13}{18}\selectfont}
\renewcommand{\cftchappagefont}{\cftchapfont}
\setlength{\cftchapnumwidth}{5em}

% 표목차: [표 2-1] 형식
\renewcommand{\cfttabpresnum}{[표\ }
\renewcommand{\cfttabaftersnum}{]\ }
\setlength{\cfttabnumwidth}{5em}
\setlength{\cftbeforetabskip}{0pt}
\setlength{\cfttabindent}{0pt}

% 그림목차: [그림 3-1] 형식
\renewcommand{\cftfigpresnum}{[그림\ }
\renewcommand{\cftfigaftersnum}{]\ }
\setlength{\cftfignumwidth}{5em}
\setlength{\cftbeforefigskip}{0pt}
\setlength{\cftfigindent}{0pt}

\renewcommand{\cftsecfont}{\JungGothic\fontsize{11}{14}\selectfont}
\renewcommand{\cftsecpagefont}{\cftsecfont}
\renewcommand{\cftsubsecfont}{\JungGothic\fontsize{11}{14}\selectfont}
\renewcommand{\cftsubsecpagefont}{\cftsubsecfont}
\renewcommand{\cftchapleader}{\cftdotfill{\cftdotsep}}
\renewcommand{\cftsecleader}{\cftdotfill{\cftdotsep}}
\renewcommand{\cftsubsecleader}{\cftdotfill{\cftdotsep}}

\setcounter{tocdepth}{3}
\setcounter{secnumdepth}{3}

% =====================================================================
%  장/절/소제목 스타일 ("제 1 장 서론", 1.6.1 형식)
% =====================================================================
\titleformat{\chapter}[hang]
  {\centering\ChapterFont\bfseries\fontsize{16}{22}\selectfont\ifnum\value{chapter}>1\color{v2red}\fi}
  {제~\thechapter~장}{0.8em}{}
\titlespacing*{\chapter}{0pt}{0pt}{18pt}

\titleformat{\section}
  {\bfseries\fontsize{13}{18}\selectfont\ifnum\value{chapter}>2\color{v2red}\fi}
  {\thesection}{0.8em}{}
\titlespacing*{\section}{0pt}{24pt}{12pt}

\makeatletter
\@beginparpenalty=0
\@secpenalty=-300
\makeatother
\widowpenalty=10000
\clubpenalty=10000

\titleformat{\subsection}
  {\JungGothic\bfseries\fontsize{11}{16}\selectfont\ifnum\value{chapter}>2\color{v2red}\fi}
  {\thesubsection}{0.8em}{}
\titlespacing*{\subsection}{0pt}{18pt}{9pt}

\titleformat{\subsubsection}
  {\JungGothic\fontsize{11}{15}\selectfont\ifnum\value{chapter}>2\color{v2red}\fi}
  {\thesubsubsection}{0.8em}{}
\titlespacing*{\subsubsection}{0pt}{14pt}{7pt}


% =====================================================================
%  표/그림 번호 형식 – "표 2-1", "그림 3-1"
% =====================================================================
\renewcommand{\figurename}{그림}
\renewcommand{\tablename}{표}
\numberwithin{figure}{chapter}
\numberwithin{table}{chapter}
\renewcommand{\thefigure}{\thechapter-\arabic{figure}}
\renewcommand{\thetable}{\thechapter-\arabic{table}}
\captionsetup[figure]{labelsep=space}
\captionsetup[table]{labelsep=space}

% 1단계 완료: 캡션 라벨 (번호 변경) — 비활성화, 2단계에서 대괄호로 대체
% 2단계 #8: 그림/표 캡션 라벨을 [그림 X-Y]/[표 X-Y] 대괄호 형식 + 빨강 표시
\DeclareCaptionLabelFormat{p2bracket}{{\color{p2red}[#1 #2]}}
\captionsetup[figure]{labelformat=p2bracket}
\captionsetup[table]{labelformat=p2bracket}

% =====================================================================
%  hyperref
% =====================================================================
\usepackage{hyperref}
\pdfstringdefDisableCommands{%
  \let\(\relax
  \let\)\relax
  \let\[\relax
  \let\]\relax
  \let\$ \relax
  \def\alpha{alpha}%
  \def\beta{beta}%
  \def\gamma{gamma}%
  \def\delta{delta}%
  \def\epsilon{epsilon}%
  \def\sigma{sigma}%
  \def\mitsigma{sigma}%
  \def\mu{mu}%
  \def\lambda{lambda}%
  \def\eta{eta}%
  \def\miteta{eta}%
  \def\kappa{kappa}%
  \def\mitkappa{kappa}%
  \def\cdot{.}%
  \def\times{x}%
  \def\leq{<=}%
  \def\geq{>=}%
  \def\mathrm#1{#1}%
  \def\vrev#1{#1}%
  \def\refchg#1{#1}%
  \def\p#1{#1}%
  \def\pref#1{#1}%
  \def\pautoref#1{#1}%
  \def\chref#1{제#1장}%
  \def\secref#1{#1절}%
  \def\subsecref#1{#1항}%
  \def\subsubsecref#1{#1목}%
  \def\appref#1{#1}%
  \def\appchref#1{부록 #1장}%
  \def\appsecref#1{부록 #1절}%
  \def\appsubsecref#1{부록 #1항}%
  \def\tabref#1{[표 #1]}%
  \def\figref#1{[그림 #1]}%
  \def\refpo#1#2#3{}%
  \def\tnew#1{#1}%
  \def\tdel#1{}%
  \def\tmod#1#2{#2}%
  \def\tnewp#1{#1}%
  \def\trefcol#1{#1}%
  \def\inew#1{#1}%
  \def\idel#1{}%
  \def\imod#1#2{#2}%
  \def\tsecref#1{#1절}%
  \def\tsubsecref#1{#1항}%
  \def\tchref#1{제#1장}%
  \def\ttabref#1{[표 #1]}%
  \def\tfigref#1{[그림 #1]}%
  \def\nrev#1{#1}%
  \def\jrev#1{#1}%
  \def\mrev#1{#1}%
  \def\prev#1{#1}%
  \let\textcolor\@gobbletwo
  \def\color#1{}%
  \def\_{}%
  \def\TOCMincho{}%
}
\hypersetup{
  colorlinks=true,
  linkcolor=black,
  citecolor=black,
  urlcolor=black,
  bookmarks=true,
  bookmarksnumbered=true,
  bookmarksopen=true,
  bookmarksopenlevel=2,
  pdfencoding=auto
}
\usepackage{bookmark}

% -----------------------------
% 참고문헌(biblatex + biber)
% 스타일: 숭실대 재난안전관리학과 관례 (송은정 박사 2022 참조)
% 형식: 저자. (연도). 제목. 출처, 권(호), p.페이지.
% [국내 문헌] / [국외 문헌] 분리
% -----------------------------
\usepackage[
  backend=biber,
  style=authoryear,
  citestyle=authoryear,
  maxcitenames=2,
  maxbibnames=99,
  sorting=nyt,
  sortcites=true,
  uniquename=false,
  uniquelist=false
]{biblatex}

\AtBeginDocument{%
  \DeclareDelimFormat[cite]{nameyeardelim}{\addcomma\space}%
  \DeclareDelimFormat[parencite]{nameyeardelim}{\addcomma\space}%
  % 한국 학회지 양식: 본문 인용 "저자(연도)" — 괄호 앞 공백 제거 (2026-06-12 권고 반영)
  \DeclareDelimFormat[textcite]{nameyeardelim}{}%
}

\DeclareSourcemap{
  \maps[datatype=bibtex]{
    \map{\step[typesource=article_foreign, typetarget=article]}
    \map{\step[typesource=book_domestic, typetarget=book]}
  }
}
\addbibresource{references.bib}

\newcommand{\citep}{\parencite}
% \citet "저자(연도)" 무공백 양식 (한국 학회지) — 제출본10 반영 확인, 검정 환원
\DeclareRobustCommand{\citet}[1]{\textcite{#1}}

\DeclareFieldFormat{titlecase}{#1}
\DeclareFieldFormat{booktitlecase}{#1}
\DeclareFieldFormat{journaltitlecase}{#1}

% --- 송은정 박사 스타일 적용: ISBN, DOI, URL 제거 ---
\AtEveryBibitem{%
  \clearfield{isbn}%
  \clearfield{issn}%
  \clearfield{doi}%
  \clearfield{url}%
  \clearfield{urldate}%
  \clearfield{eprint}%
  \clearfield{eprinttype}%
  \clearfield{note}%
}
% 저널명·서적명 이탤릭 제거 (송은정 스타일: 이탤릭 없음)
\DeclareFieldFormat{journaltitle}{#1}
\DeclareFieldFormat{booktitle}{#1}
\DeclareFieldFormat[book]{title}{#1}
\DeclareFieldFormat[inbook]{title}{#1}
% 논문 제목 인용부호 제거
\DeclareFieldFormat[article]{title}{#1}
\DeclareFieldFormat[inproceedings]{title}{#1}
\DeclareFieldFormat[thesis]{title}{#1}
% pp. → p. 변환
\DeclareFieldFormat{pages}{p.#1}
% edition 형식
\DeclareFieldFormat{edition}{#1}

% Manual numbering macros
% =====================================================================
% preamble_manual_numbering.tex – 한국어 학술문서 번호체계
% =====================================================================
% 문서 구조: 제1장 → 1.1 → 1.1.1 → 1.1.1.1
% 본문 넘버링 (enumerate): 1) → 가) → (1) → (가) → ① → ⓐ
% 불릿 포인트 (itemize): • (가운데점) → - (대쉬)
% 들여쓰기: 각 레벨은 상위에서 2em 들여쓰기
% =====================================================================

% --- 한글 가나다 카운터 정의 ---
\makeatletter
\newcommand*{\korean@char}[1]{%
  \ifcase#1\or 가\or 나\or 다\or 라\or 마\or 바\or 사\or 아\or 자\or 차%
  \or 카\or 타\or 파\or 하\fi}
\AddEnumerateCounter*{\gana}{\korean@char}{가}
\makeatother

% --- Enumerate 설정 (6단계) ---
% Level 1: 1), 2), 3)... (기본 들여쓰기 2em)
% Level 2: 가), 나), 다)... (상위에서 +2em)
% Level 3: (1), (2), (3)... (상위에서 +2em)
% Level 4: (가), (나), (다)... (상위에서 +2em)
% Level 5: ①, ②, ③... (상위에서 +2em)
% Level 6: ⓐ, ⓑ, ⓒ... (상위에서 +2em)

\setlist[enumerate,1]{%
  label=\arabic*),%
  leftmargin=2em,%
  labelsep=0.5em,%
  itemsep=0.3em,%
  parsep=0pt,%
  topsep=0.5em%
}

\setlist[enumerate,2]{%
  label=\gana*),%
  leftmargin=2em,%
  labelsep=0.5em,%
  itemsep=0.2em,%
  parsep=0pt,%
  topsep=0.3em%
}

\setlist[enumerate,3]{%
  label=(\arabic*),%
  leftmargin=2em,%
  labelsep=0.5em,%
  itemsep=0.2em,%
  parsep=0pt,%
  topsep=0.3em%
}

\setlist[enumerate,4]{%
  label=(\gana*),%
  leftmargin=2em,%
  labelsep=0.5em,%
  itemsep=0.2em,%
  parsep=0pt,%
  topsep=0.3em%
}

% Level 5: ①, ②, ③... (원문자)
\setlist[enumerate,5]{%
  label=\ding{\numexpr171+\value{enumv}\relax},%
  leftmargin=2em,%
  labelsep=0.5em,%
  itemsep=0.2em,%
  parsep=0pt,%
  topsep=0.3em%
}

% Level 6: ⓐ, ⓑ, ⓒ... (원소문자 - Unicode 직접 사용)
\newcommand*{\circledletter}[1]{%
  \ifcase#1\or ⓐ\or ⓑ\or ⓒ\or ⓓ\or ⓔ\or ⓕ\or ⓖ\or ⓗ\or ⓘ\or ⓙ%
  \or ⓚ\or ⓛ\or ⓜ\or ⓝ\or ⓞ\or ⓟ\or ⓠ\or ⓡ\or ⓢ\or ⓣ%
  \or ⓤ\or ⓥ\or ⓦ\or ⓧ\or ⓨ\or ⓩ\fi}

\setlist[enumerate,6]{%
  label=\circledletter{\value{enumvi}},%
  leftmargin=2em,%
  labelsep=0.5em,%
  itemsep=0.2em,%
  parsep=0pt,%
  topsep=0.3em%
}

% --- Itemize 설정 (2단계) ---
% Level 1: • (가운데점) - 자료 분류용
% Level 2: - (대쉬) - 하위 내용용
% 들여쓰기: 상위에서 2em

\setlist[itemize,1]{%
  label=\textbullet,%
  leftmargin=2em,%
  labelsep=0.5em,%
  itemsep=0.3em,%
  parsep=0pt,%
  topsep=0.5em%
}

\setlist[itemize,2]{%
  label=--,%
  leftmargin=2em,%
  labelsep=0.5em,%
  itemsep=0.2em,%
  parsep=0pt,%
  topsep=0.3em%
}

\setlist[itemize,3]{%
  label=\textbullet,%
  leftmargin=2em,%
  labelsep=0.5em,%
  itemsep=0.2em,%
  parsep=0pt,%
  topsep=0.3em%
}

\setlist[itemize,4]{%
  label=--,%
  leftmargin=2em,%
  labelsep=0.5em,%
  itemsep=0.2em,%
  parsep=0pt,%
  topsep=0.3em%
}

% --- enumerate 최대 깊이 확장 ---
\setlistdepth{6}

% =====================================================================
% 사용 예시:
% =====================================================================
% \begin{enumerate}
%   \item 첫 번째 항목 (1))
%   \begin{enumerate}
%     \item 하위 항목 (가))
%     \begin{enumerate}
%       \item 세부 항목 ((1))
%     \end{enumerate}
%   \end{enumerate}
% \end{enumerate}
%
% \begin{itemize}
%   \item 분류 항목 (•)
%   \begin{itemize}
%     \item 하위 내용 (-)
%   \end{itemize}
% \end{itemize}
%
% enumerate 내에서 itemize 시작 시 동일한 2em 들여쓰기 적용
% =====================================================================


\begin{document}
\justifying
\raggedbottom

% ===============================
% Front matter (Roman numerals)
% ===============================
\pagenumbering{roman}

% front/01_cover.tex
% 모든 페이지와 동일 geometry(4/4/5cm) 사용

\thispagestyle{empty}

\begingroup
\setlength{\parskip}{0pt}
\setlength{\parindent}{0pt}

% --------------------------------
% 0) 박사 학위 논문  – 상단에서 0.5cm
% (top=4cm 이 이미 적용되어 있으니, 텍스트 블록 위에서 0.5cm만 더 내려감)
% --------------------------------
\vspace*{0.5cm}
{\fontsize{14}{20}\selectfont 박사 학위 논문}

% --------------------------------
% 1) 국문 제목 – 상단에서 3.5cm
% → 이미 0.5cm 내려왔으니, 3.5 - 0.5 = 3.0cm 더 내려주면 됨
% --------------------------------
\vspace*{3.0cm}

\begin{center}
  \begin{minipage}{13cm}  % 한글/영문 제목 폭 = 13cm
    \centering
    {\fontsize{18}{18}\selectfont\bfseries
    관측성과 불확실성 기반 동적\par
    복구목표시간(DRTO) 프레임워크 연구\par}
  \end{minipage}
\end{center}

% --------------------------------
% 2) 영문 제목 – 한글 제목 블록과 1cm 간격
% --------------------------------
\vspace*{1.0cm}

\begin{center}
  \begin{minipage}{13cm}
    \centering
    {\fontsize{18}{18}\selectfont\bfseries
      A Dynamic RTO Framework\par
      Based on Observability and Uncertainty\par}
  \end{minipage}
\end{center}

% --------------------------------
% 3) 아래 정보 (날짜 ~ 이름)
%    간격: 날짜-숭실대 3.5cm, 숭실대-학과 1.5cm, 학과-이름 1.5cm
%    "박스 하단 → 이름 1.5cm"는 보고 눈으로 미세조정하는 쪽이 현실적이라
%    아래 마지막 \vspace* 값을 조금씩 조절해 쓰시면 됩니다.
% --------------------------------

% 영문 제목 아래에서 날짜까지 여백 (눈으로 맞춰볼 값, 대략 3.0~3.5cm 정도에서 시작)
\vspace*{3.0cm}

\begin{center}
  {\fontsize{14}{20}\selectfont 2026년 6월}\par

  \vspace{2.0cm} % 날짜 → 숭실대 대학원 (3.5cm)
  {\fontsize{16}{22}\selectfont\bfseries 숭실대학교 대학원}\par

  \vspace{1.0cm} % 숭실대 대학원 → 학과 (1.5cm)
  {\fontsize{14}{20}\selectfont 재난안전관리학과}\par

  \vspace{1.0cm} % 학과 → 이름 (1.5cm)
  {\fontsize{14}{20}\selectfont 이\ \ \ 창\ \ \ 호}\par

  % 박스 하단에서 이름까지 1.5cm 목표
  % 부족하거나 넘치면 여기 값을 ±0.5cm 정도씩 조절해 보세요.
  \vspace*{1.5cm}
\end{center}

\endgroup

\clearpage

% front/02_blank_page.tex
\thispagestyle{empty}
\null
\clearpage

\input{front/03_titlepage}
% front/04_inner_title.tex
% 모든 페이지와 동일 geometry(4/4/5cm) 사용

\thispagestyle{empty}

\begingroup
\setlength{\parskip}{0pt}
\setlength{\parindent}{0pt}

% --------------------------------
% 0) 박사 학위 논문  – 상단에서 0.5cm
% (top=4cm 이 이미 적용되어 있으니, 텍스트 블록 위에서 0.5cm만 더 내려감)
% --------------------------------
\vspace*{0.5cm}
{\fontsize{14}{20}\selectfont 박사 학위 논문}

% --------------------------------
% 1) 국문 제목 – 상단에서 3.5cm
% → 이미 0.5cm 내려왔으니, 3.5 - 0.5 = 3.0cm 더 내려주면 됨
% --------------------------------
\vspace*{1.5cm}

\begin{center}
  \begin{minipage}{13cm}  % 한글/영문 제목 폭 = 13cm
    \centering
    {\fontsize{18}{18}\selectfont\bfseries
    관측성과 불확실성 기반 동적\par
    복구목표시간(DRTO) 프레임워크 연구\par}
  \end{minipage}
\end{center}

% --------------------------------
% 2) 지도교수 – 한글 제목 블록과 1cm 간격
% --------------------------------
\vspace*{1.0cm}

\begin{center}
  \begin{minipage}{13cm}
    \centering
    {\fontsize{14}{26}\selectfont\bfseries
      지도교수 정  종  수
    }\par
  \end{minipage}
\end{center}

% --------------------------------
% 3) 논문 제출 – 지도교수 제목 블록과 1cm 간격
% --------------------------------
\vspace*{0.5cm}

\begin{center}
  \begin{minipage}{13cm}
    \centering
    {\fontsize{14}{26}\selectfont\bfseries
      이 논문은 박사학위 논문으로 제출함
    }\par
  \end{minipage}
\end{center}

% --------------------------------
% 3) 아래 정보 (날짜 ~ 이름)
%    간격: 날짜-숭실대 3.5cm, 숭실대-학과 1.5cm, 학과-이름 1.5cm
%    "박스 하단 → 이름 1.5cm"는 보고 눈으로 미세조정하는 쪽이 현실적이라
%    아래 마지막 \vspace* 값을 조금씩 조절해 쓰시면 됩니다.
% --------------------------------

% 영문 제목 아래에서 날짜까지 여백 (눈으로 맞춰볼 값, 대략 3.0~3.5cm 정도에서 시작)
\vspace*{3.0cm}

\begin{center}
  {\fontsize{14}{20}\selectfont 2026년 6월}\par

  \vspace{2.0cm} % 날짜 → 숭실대 대학원 (3.5cm)
  {\fontsize{18}{22}\selectfont\bfseries 숭실대학교 대학원}\par

  \vspace{1.0cm} % 숭실대 대학원 → 학과 (1.5cm)
  {\fontsize{14}{20}\selectfont\bfseries 재난안전관리학과}\par

  \vspace{1.0cm} % 학과 → 이름 (1.5cm)
  {\fontsize{14}{20}\selectfont\bfseries 이\ \ \ 창\ \ \ 호}\par

  % 박스 하단에서 이름까지 1.5cm 목표
  % 부족하거나 넘치면 여기 값을 ±0.5cm 정도씩 조절해 보세요.
  % \vspace*{1.5cm}
\end{center}

\endgroup

\clearpage

% thesis/front/05_approval.tex
\begin{titlepage}
\thispagestyle{empty}

\vspace*{15mm}

{\fontsize{18}{26}\selectfont\bfseries
\noindent\hspace*{20mm}
\makebox[0.78\textwidth][s]{%
  이\quad 창\quad 호의\quad 박사학위\quad 논문을\quad 인준함.
}}
\par

\vspace{15mm}
\raggedright
\vspace{10mm}\hspace*{64mm}{\large 심 사 위 원 장\; \rule{45mm}{0.4pt}\; (인)}\\[6mm]
\vspace{10mm}\hspace*{70mm}{\large 심  사  위  원\; \rule{45mm}{0.4pt}\; (인)}\\[6mm]
\vspace{10mm}\hspace*{70mm}{\large 심  사  위  원\; \rule{45mm}{0.4pt}\; (인)}\\[6mm]
\vspace{10mm}\hspace*{70mm}{\large 심  사  위  원\; \rule{45mm}{0.4pt}\; (인)}\\[6mm]
\vspace{10mm}\hspace*{70mm}{\large 심  사  위  원\; \rule{45mm}{0.4pt}\; (인)}\\[8mm]

\centering
\vspace*{5mm}
{\fontsize{14}{26}\selectfont 2026년 \hspace{8mm} 8월\par}

\vspace{5mm}
{\fontsize{18}{26}\selectfont\bfseries 숭실대학교 대학원}\par

\end{titlepage}

% thesis/front/06_acknowledgements.tex
% 감사의 글 (목차에 포함되지 않음)

\thispagestyle{empty}  % 페이지 번호 숨김

\begin{center}
  {\fontsize{16}{22}\selectfont\bfseries 감사의 글}
\end{center}

\vspace{1.5cm}

% — 아래에 내용 작성 —
{\fontsize{11}{22}\selectfont\setstretch{2.0}\setlength{\parindent}{2em}\setlength{\parskip}{6pt}\justifying
\spaceskip=1.6\fontdimen2\font plus 0.35\fontdimen3\font minus 0.35\fontdimen4\font
본 논문이 완성되기까지 한결같은 지도와 아낌없는 격려를 보내주신
지도교수 정종수 교수님께 진심 어린 감사의 뜻을 표합니다.
연구 주제 선정과 방향 설정에서부터 핵심 이론의 정립, 시뮬레이션 설계,
데이터 분석, 그리고 논문의 세부 구성에 이르기까지
세심하고도 깊이 있는 지도를 해주셨기에
본 연구를 현재의 수준으로 마무리할 수 있었습니다.
특히 DRTO 프레임워크의 수학적 모델링과 시그모이드 바운딩 메커니즘의
이론적 기반을 정립하는 과정에서 보여주신 학문적 통찰력은
연구의 방향을 올바르게 이끄는 나침반이 되었습니다.
교수님께서 보여주신 학문적 열정과 인내 어린 가르침은
저에게 큰 울림과 배움의 기회가 되었으며, 앞으로의 연구 인생에서도
잊지 못할 귀중한 자산이 될 것입니다.

아울러 바쁜 일정 중에도 논문 심사를 맡아
귀중한 조언과 격려를 아끼지 않으신 심사위원 교수님들께도 감사드립니다.
교수님들께서 주신 건설적인 비판과 소중한 의견은
본 연구의 완성도를 높이고 학문적 깊이를 더하는 데 큰 도움이 되었습니다.
심사 과정에서 제기해 주신 날카로운 질문들은
연구의 논리적 허점을 보완하고 학술적 엄밀성을 강화하는 계기가 되었으며,
이를 통해 논문의 전반적인 품질을 한층 끌어올릴 수 있었습니다.

특별히, DRTO 프레임워크의 실무적 타당성 검증을 위해
귀한 시간을 내어 전문가 평가에 참여해 주신 다섯 분야의 전문가 여러분께
깊은 감사의 말씀을 드립니다.
업무연속성관리(BCM) 분야에서는 BCP 수립 및 운영 관점에서
프레임워크의 BCM 프로세스 통합 가능성을 평가해 주셨고,
재해복구(DR) 분야에서는 IT 인프라 복구 계획의 기술적 타당성과
운영 적용성에 대한 심도 있는 검토를 제공해 주셨습니다.
리스크 관리 분야에서는 파라미터 적정성과 리스크 모델링의 현실성을,
IT 운영 관리 분야에서는 실무 환경에서의 적용 가능성과 운영 통합 방안을,
그리고 정보보안 분야에서는 사이버 위협 시나리오와 복원력 관점에서
프레임워크의 보안 측면을 면밀히 검토해 주셨습니다.
각 분야에서 오랜 경력과 전문성을 쌓아 오신 전문가분들께서
보내주신 통찰력 있는 조언과 세심한 검토는
본 연구의 학술적 가치와 현업 적용 가능성을 한층 높이는 데
결정적인 기여를 하였습니다.
현장의 생생한 경험에서 우러나온 귀중한 피드백은
이론과 실무의 간극을 메우는 소중한 가교가 되었으며,
향후 DRTO 프레임워크가 실제 BCM/DR 환경에 적용될 때
실무자들이 직면할 수 있는 과제들을 미리 예측하고 대비할 수 있게 해주었습니다.

연구 과정 전반에서 상담, 논의, 협업 등 다양한 방식으로 도움을 주신
동료들과 여러 교수님들께도 깊이 감사드립니다.
학문적 토론과 아이디어 공유를 통해 새로운 시각을 열어주셨고,
연구의 어려운 고비마다 함께 고민하고 격려해 주셨습니다.
여러분과의 소중한 인연과 협력은 본 연구를 이어가는 데 큰 힘이 되었습니다.

마지막으로, 언제나 믿고 응원해 주며 연구 여정의 든든한 버팀목이 되어 준
가족들에게 이 논문의 영광을 돌리고자 합니다.
학업과 연구에 전념할 수 있도록 묵묵히 희생하고 기다려 준 가족들,
지칠 때마다 다시 일어설 수 있는 용기를 준 그 따뜻한 사랑에 깊이 감사드립니다.
가족들의 헌신과 변함없는 격려가 있었기에 어려움 속에서도 흔들림 없이
끝까지 연구를 완수할 수 있었습니다.
다시 한 번 이 논문이 나오기까지 도와주신 모든 분들께 진심으로 감사드립니다.
\par}

\vspace{6mm}

\begin{center}
  2026년 6월
\end{center}

\vspace*{4mm}

\begin{flushright}
  \hspace*{-2cm}
  이\ \ 창\ \ 호
\end{flushright}

% front/07_kor_abstract.tex

\clearpage
\addcontentsline{toc}{chapter}{\PlainTOCtext{국문초록}}

% ① "국문초록" 제목 (왼쪽, 16pt Bold)
{\fontsize{16}{20}\selectfont\bfseries 국문초록}\par
\vspace{15mm}  % 1.5cm

% ② 한글 논문 제목 (16pt Bold)
\begin{center}
  \begin{minipage}{0.9\textwidth}
    \centering
    {\fontsize{16}{20}\selectfont\bfseries
    관측성과 불확실성 기반 동적\par
    복구목표시간(DRTO) 프레임워크 연구
    }\par
  \end{minipage}
\end{center}
\vspace{15mm}  % 1.5cm

% ③ 이름/학과/대학원 – 오른쪽 정렬 + 2cm 안쪽
\begin{flushright}
  \hspace*{-2cm}
  {\fontsize{14}{22}\selectfont
    \setlength{\baselineskip}{22pt}
    이\ \ 창\ \ 호\par
    재난안전관리학과\par
    숭실대학교 대학원\par
  }
\end{flushright}

\vspace{15mm}  % 1.5cm

% ④ 본문 (11pt, 줄간 200%)
{\fontsize{11}{22}\selectfont
\setstretch{2.0}\justifying
\setlength{\parindent}{2em}

업무연속성관리(BCM)에서 복구목표시간(RTO)은 업무영향분석(BIA) 시점에 고정값으로 설정되어, 이후 운영 환경의 변화와 무관하게 유지된다. 이러한 정적 RTO 접근법은 시스템 관측성 수준의 동적 변동, 복구 과정에 내재된 불확실성의 확률적 특성, 사이버 공격이나 복합 재난 등 극단적 사건의 파레토 위험을 구조적으로 반영하지 못하는 한계를 지닌다.

본 연구는 이러한 한계를 극복하기 위해, \chg{관측성과 불확실성을 개별 시그모이드 바운딩 함수로 정식화한 수식 모델, 정적 기준선에서 파레토 불확실성까지의 4단계 운영 조건(C0--C3), 조건당 10,000회 몬테카를로 시뮬레이션, 효과 크기와 CDF 교차점 기반 분할 회귀 및 CVaR 꼬리 위험 정량화를 포함하는 다차원 비교 분석, 다중 시드($10^8$회) 강건성 검증, 기초 검증과 10개 연구가설 검증, 전문가 평가를 결합한 6단계 다층적 타당성 검증 체계(L0--L5), 6개 산업 시나리오 범용성 검증}으로 구성되는 DRTO 프레임워크를 제안한다. 핵심 수식은 $\DRTO = R_0 + E_{O,\mathrm{bnd}} + E_{U,\mathrm{bnd}}$로 정의되며, 관측성 효과 $E_{O,\mathrm{bnd}} = -R_0 \cdot S(z_O)$는 복구 시간을 기준선 $R_0$로부터 하향 조정하고, 불확실성 효과 $E_{U,\mathrm{bnd}} = (R_{\max} - R_0) \cdot S(z_U)$는 상향 조정한다. 수정 시그모이드 $S(z) = 2/(1 + e^{-z}) - 1$에 의해 $\DRTO$는 물리적 경계 $(0,\; R_{\max})$ 내에서 점근적으로 보장되며, $R_0 = 6.0$시간, $R_{\max} = 24.0$시간($= 4R_0$), $\kappa = 1.2$를 기본 설정으로 채택하였다. \fpara{$\kappa = 1.2$는 C1 조건의 격자 탐색에서 도출된 결과이며, 본 연구는 이를 \emph{환경별 최적값}이 아닌 BCM 일반 운영의 \emph{\chg{잠정적이고 보수적인} 실무 기본값}으로 채택한다. C2/C3 조건에서의 적합성은 가설 H2--H10 PASS의 통합 해석으로 \emph{운영 적합성을 시사}하며, 환경 무관 $\kappa$의 \emph{존재 여부 탐색}은 후속 연구의 과제로 남긴다.}

모델은 효율비 $\eta = \DRTO / R_0$를 기준으로 네 가지 운영 조건을 정의한다. \chg{정적 기준선인 C0는 $\eta = 1.000$이고, 관측성만 활성화한 C1은 $\eta \leq 1.000$이며, 가우시안 불확실성을 결합한 C2는 $\eta \geq \eta_{\mathrm{C1}}$, 파레토 불확실성의 C3는 꼬리 영역에서 $\eta \geq \eta_{\mathrm{C2}}$를 보인다.} 두 분포의 기대값을 $3.0$으로 통일하여 분포 형태의 차이만을 비교하도록 설계하였다.

조건당 $n = 10{,}000$ 몬테카를로 시뮬레이션 결과, 공통 경로에서 C0($1.000$) $\to$ C1($0.853$)으로 관측성에 의해 평균 $14.7\%$ 단축되었으며(Cohen's $d = -1.180$), C1에서 분기하여 C2($1.682$)와 C3($1.514$) 모두 불확실성 프리미엄에 의해 상향되었다. C2와 C3는 동일한 평균값($3.0$)을 갖도록 설계되었음에도 불구하고, 극단 백분위 영역에서는 현저한 차이를 보였다. P99 기준에서 C3의 $\DRTO$는 약 $13.2$시간으로, C2의 약 $8.4$시간보다 약 $57\%$ 높은 값을 나타냈으며, 이는 파레토 분포의 $U_{\mathrm{raw}}$ 99백분위가 가우시안 대비 약 $4.1$배($21.85$ vs. $5.33$) 큰 꼬리 특성에 기인한다.

두 불확실성 분포의 CDF가 교차하는 P85.8($\eta \approx 2.42$)을 기준으로, 핵심 영역(Core, $\leq$ P85.8)과 꼬리 영역(Tail, $>$ P85.8)으로 분할 회귀를 수행한 결과, 관측성 가중치 $k_O$의 회귀 계수가 Core($-0.798$)에서 Tail($-0.415$)로 $0.52$배 감쇠되는 이중채널 감쇠 효과를 발견하였다. 이는 극단적 불확실성 하에서 관측성의 복구 단축 효과가 구조적으로 약화됨을 의미한다. CVaR$_{99}$ 기준 C3가 C2 대비 $+24.0\%$ 높아, 파레토 위험이 극단 영역의 복구 부담을 불균형적으로 증폭시킴을 정량적으로 입증하였다.

본 프레임워크는 다층적 검증 체계를 통해 타당성을 확보하였다. 10개 \m{연구 가설}(H1--H10)이 모두 채택되었으며, BCM/DR 분야 전문가 5명의 정량적 설문은 전체 평균 $4.46/5.0$점($p = 0.011$, Cohen's $d = 1.65$)과 높은 내적 일관성(Cronbach's $\alpha = 0.89$, S-CVI/Ave$\,= 0.985$)을 보였고, 심층 인터뷰에서는 5대 주제 15개 하위 주제 모두에서 $80\%$ 이상 동의($100\%$ 충족)를 확인하였다. 2차 회귀 분석은 C1에서 $R^2 = 1.000$, C2에서 $R^2 = 0.998$의 사실상 완전한 설명력을 달성하였고, C3는 전체 $R^2 = 0.858$이나 P85.8 분할 시 Core $R^2 = 0.999$, Tail $R^2 = 0.710$으로서 분할 회귀의 유효성을 입증하였다. 또한 6개 산업 시나리오(\chg{$R_0 = 0.5$--$8.0$시간})에서 P99 비율이 평균 $1.57$배로 일관되어 모델의 범용 적용 가능성을 확인하였다.

본 연구는 관측성의 이점과 불확실성의 패널티를 개별 바운딩으로 통합하고, 특히 복구 계획에 대한 Heavy-Tail 분포의 불균형적 영향과 이중채널 감쇠 메커니즘을 규명함으로써, \chg{통계적 근거에 기반한 DRTO 조정 방법론과 더불어, 회귀식 기반 $\DRTO$ 산출 도구와, 이중채널 감쇠에 기반한 자원 배분 원칙 및 스트레스 테스트 설계 기준을} 업무연속성관리 분야에 기여한다. \fpara{본 연구의 학술적 한계로서, $\kappa = 1.2$의 환경 무관 보편 최적성은 본 연구에서 입증되지 않으며 이는 후속 연구의 탐색 과제로 남긴다.}

\vspace{10mm}
\noindent\textbf{주제어:} 복구목표시간(RTO), 동적복구목표시간(DRTO), 업무연속성관리(BCM), 관측성, 불확실성정량화, 시그모이드바운딩, 중꼬리분포

}

% front/08_eng_abstract.tex

\clearpage
\newcommand{\EnglishAbstractTitle}{\mbox{영문초록}}
\newcommand{\EnglishAbstractTOC}{\PlainTOCtext{영문초록}}
\let\EngAbstractHeading\EnglishAbstractTitle
\phantomsection
\addcontentsline{toc}{chapter}{\EnglishAbstractTOC}

% ① "Abstract" title (left aligned, 16pt bold)
{\fontsize{16}{20}\selectfont\bfseries\EnglishAbstractTitle}\par
\vspace{15mm}  % 1.5cm

\begin{center}
  \begin{minipage}{0.9\textwidth}
    \centering
    {\fontsize{16}{20}\selectfont\bfseries
      A Dynamic RTO Framework\par
      Based on Observability and Uncertainty\par}
  \end{minipage}
\end{center}

% ③ Name/department block – right aligned, inset by 2cm
\vspace{15mm}  % 1.5cm

\begin{flushright}
  \hspace*{-2cm}
  {\fontsize{14}{22}\selectfont
    \setlength{\baselineskip}{22pt}
    LEE,\ CHANGHO\par
    Department of Disaster and Safety Management\par
    Graduate School of Soongsil University\par
  }
\end{flushright}

\vspace{15mm}  % 1.5cm

% ④ Main text (11pt, line spacing 200%)
{\fontsize{11}{22}\selectfont
\setstretch{2.0}\justifying
\setlength{\parindent}{2em}

In business continuity management (BCM), the Recovery Time Objective (RTO) is determined as a fixed value during business impact analysis (BIA) and remains unchanged regardless of subsequent changes in the operational environment. This static RTO approach has structural limitations in that it fails to account for dynamic fluctuations in system observability, the probabilistic nature of uncertainty inherent in recovery processes, and heavy-tail risks from extreme events such as cyberattacks and compound disasters.

To address these limitations, this study proposes a Dynamic RTO (DRTO) framework comprising seven components: (i) a mathematical model that formulates observability and uncertainty via individual sigmoid bounding functions, (ii) four operational conditions (C0--C3) spanning from a static baseline to heavy-tail uncertainty, (iii) Monte Carlo simulation with $n = 10{,}000$ per condition, (iv) multi-dimensional comparative analysis including effect sizes, CDF-crossover-based split regression, and CVaR tail-risk quantification, (v) robustness verification through multi-seed testing ($10^8$ runs), (vi) a six-stage multi-layer validation framework (L0--L5) combining Golden Compare foundation verification, ten hypothesis tests, and expert evaluation, and (vii) generalizability verification across six industry scenarios. The core equation is defined as $\DRTO = R_0 + E_{O,\mathrm{bnd}} + E_{U,\mathrm{bnd}}$, where the observability effect $E_{O,\mathrm{bnd}} = -R_0 \cdot S(z_O)$ adjusts recovery time downward from the baseline $R_0$, and the uncertainty effect $E_{U,\mathrm{bnd}} = (R_{\max} - R_0) \cdot S(z_U)$ adjusts it upward. The modified sigmoid $S(z) = 2/(1 + e^{-z}) - 1$ asymptotically guarantees that $\DRTO$ remains within the physical bounds $(0,\; R_{\max})$. The default parameters are $R_0 = 6.0$ hours, $R_{\max} = 24.0$ hours ($= 4R_0$), and $\kappa = 1.2$. \fpara{The value $\kappa = 1.2$ was derived from grid search in environment C1 (observability-only, Gaussian baseline) and is adopted not as an \emph{environment-specific optimum} but as a \emph{provisional and conservative operational default} for general BCM operations. Its applicability in environments C2/C3 is interpreted as \emph{suggesting operational adequacy} through the integrated PASS of hypotheses H2--H10; the \emph{investigation of the existence} of an environment-invariant $\kappa$ is deferred to future research.}

The model defines four operational conditions in terms of the efficiency ratio $\eta = \DRTO / R_0$. \chg{The static baseline C0 has $\eta = 1.000$, the observability-only C1 has $\eta \leq 1.000$, the Gaussian-uncertainty C2 has $\eta \geq \eta_{\mathrm{C1}}$, and the Pareto-uncertainty C3 has $\eta \geq \eta_{\mathrm{C2}}$ in the tail region.} Both distributions share a unified expected value of $3.0$, enabling fair comparison of distributional shape effects.

Monte Carlo simulation with $n = 10{,}000$ per condition revealed a two-stage branching structure in the efficiency ratio $\eta = \DRTO / R_0$: the common path C0 ($1.000$) $\to$ C1 ($0.853$) showed a mean reduction of $14.7\%$ (Cohen's $d = -1.180$), confirming that observability significantly decreases recovery time. From C1, two parallel branches emerged: C2 ($1.682$, Gaussian) and C3 ($1.514$, Pareto), both elevated by uncertainty premiums. Despite the common mean design, C2 and C3 exhibited marked differences at extreme percentiles: at the 99th percentile, C3 $\DRTO$ reached approximately $13.2$ hours compared to C2 at approximately $8.4$ hours---a $57\%$ gap attributable to the Pareto distribution's $U_{\mathrm{raw}}$ 99th percentile being approximately $4.1$ times larger than the Gaussian ($21.85$ vs.\ $5.33$).

Using the CDF crossover point at P85.8 ($\eta \approx 2.42$), split regression between Core ($\leq$ P85.8) and Tail ($>$ P85.8) regions revealed a dual-channel attenuation effect: the regression coefficient of observability weight $k_O$ attenuated from $-0.798$ (Core) to $-0.415$ (Tail), a $0.52\times$ reduction indicating that observability's recovery-shortening effect is structurally weakened under extreme uncertainty. CVaR$_{99}$ for C3 exceeded that of C2 by $24.0\%$, quantitatively demonstrating that heavy-tail risk disproportionately amplifies recovery burden in extreme regions.

The framework was validated through a six-stage multi-layer verification system (L0--L5). Foundation verification (L0, Golden Compare) confirmed implementation correctness via 47 deterministic checkpoints (100\% PASS). All ten hypotheses (H1--H10) were accepted. Expert evaluation by five BCM/DR specialists (mean $15.2$ years of experience) yielded a quantitative survey mean of $4.46/5.0$ ($p = 0.011$, Cohen's $d = 1.65$) with high internal consistency (Cronbach's $\alpha = 0.89$, S-CVI/Ave$\,= 0.985$), and in-depth interviews confirmed $\geq 80\%$ agreement on all 15 sub-themes across 5 major themes ($100\%$ met). Full quadratic regression (including interaction terms) achieved $R^2 = 1.000$ for C1, $R^2 = 0.998$ for C2, and $R^2 = 0.858$ for C3 overall, with the P85.8 split yielding Core $R^2 = 0.999$ and Tail $R^2 = 0.710$, validating the effectiveness of split regression. Additionally, six industry scenarios ($R_0$: $0.5$--$8.0$ hours) showed a consistent P99 ratio averaging $1.57\times$, confirming the model's general applicability.

This research contributes to business continuity management by integrating observability benefits and uncertainty penalties through individual bounding, and by elucidating the disproportionate impact of heavy-tail distributions and the dual-channel attenuation mechanism. It thereby provides a statistically grounded methodology for dynamic RTO adjustment, a regression-based DRTO computation tool, and---grounded in the dual-channel attenuation finding---resource-allocation principles and stress-test design criteria. \fpara{As an acknowledged academic limitation, the environment-invariant universal optimality of $\kappa = 1.2$ is not proven in this study and is deferred to future research.}

\vspace{10mm}
\noindent\textbf{Keywords:} Recovery Time Objective(RTO), Dynamic RTO(DRTO), Business Continuity Management(BCM), Observability, Uncertainty Quantification, Sigmoid Bounding, Heavy-tailed Distribution

}

\clearpage

\FrontListTitle{목\ \ 차}
\tableofcontents
\clearpage
\FrontListTitle{그\ \ 림\ \ 목\ \ 차}
\listoffigures
\clearpage
\FrontListTitle{표\ \ 목\ \ 차}
\listoftables

\clearpage
\pagenumbering{arabic}
\setcounter{page}{1}

% ─── 본문 (5장 체계, 1차 심사 재편) ───
% 제1장 서론
% 참조: doctoral/research_model/research_model.tex, doctoral/ontology/ontology_analysis.tex
% v9.3 기준값: qa/v9.3_band_seed/reference_values.json
% κ = 1.2, R₀ = 6.0h, R_max = 24.0h (= 4R₀)

\chapter{서론}
\label{ch:introduction}

\section{연구 배경}
\label{sec:background}

현대 사회는 자연재해, 사이버공격, 공급망 중단, 감염병 대유행 등 다양한 재난 및
업무 중단 위험에 직면하고 있다. 이러한 위험은 단일 사건으로 발생하기보다는
복합적으로 상호작용하며, 조직의 핵심 업무 프로세스에 심각한 영향을 미친다.
BCMS(Business Continuity Management System, 업무연속성관리체계)는 이러한 재난
상황에서 조직의 핵심 기능을 유지하고 신속하게 복구하기 위한 체계적 관리
프레임워크로서, \nrev{ISO 22301:2019 국제표준}을 중심으로
전 세계적으로 확산되고 있다.

BCMS의 수립 과정에서 BIA(Business Impact Analysis, 업무영향분석)은 핵심
구성 요소로 기능하며, 이를 통해 MTPD(Maximum Tolerable
Period of Disruption, 최대허용중단기간)과 RTO(Recovery Time Objective, 복구목표시간)이 결정된다
\citep{hiles2010bcm, snedaker2013bcdr}.
RTO는 재난 발생 후 업무 프로세스가 수용 가능한 수준으로 복구되어야 하는
목표 시간을 의미하며, $\text{RTO} \leq \text{MTPD}$의 관계가 성립한다
\nrev{(ISO 22313, 2020)}.
현행 \nrev{ISO 22301:2019 국제표준} 및 대부분의 BCM 실무에서 RTO는 BIA 단계에서 고정값으로
설정되며, BIA 주기 사이의 평상 운영 중에는 실시간으로 변경되지
않는다~\citep{herbane2010disaster}.

\chg{그러나 복구가 이루어지는 실제 운영 환경은 BIA 시점에 고정한
가정과 달리 끊임없이 변동하므로, 한 번 설정한 RTO는 실제 복구 상황과 어긋날
수 있다. 이러한 한계를 다루기 위해 \citet{lee2026drto}은 조직의
장애를 탐지하고 상황을 파악하는 능력인 관측성(observability)을 반영하여
RTO를 동적으로 조정하는 DRTO 모델을 제안하였다.

이 모델은 평상시 기준 복구시간 $R_0$에서 출발하여, 관측성의 효과를 시그모이드
함수로 변환해 기준 시간에서 차감하는 형태이다. 관측성이 높을수록 목표 시간을
더 많이 단축하되, 시그모이드 바운딩에 의해 복구시간이 $0$ 아래로 내려가지
않으며, 곡률은 매개변수 $\kappa$로 조절한다. 해당 연구에서는 정적 기준(C0)에서
관측성을 활성화한 조건(C1)으로의 전환을 설정하고, 대규모 몬테카를로
시뮬레이션과 비중첩 대역 시드 체계로 그 효과를 검증함으로써, 고정된 정적 RTO를
상황에 반응하는 동적 모델로 전환하는 진전을 이루었다.

그러나 이 모델은 관측성이라는 단일 차원에 머물렀다. 복구에 걸리는 시간은
장비 가용성과 인력 동원, 외부 공급망 등 여러 요인의 불확실성(uncertainty)
아래 놓인다. \chg{이 불확실성은 일상적 변동을 나타내는 가우시안(정규) 분포뿐
아니라, 그 범위를 크게 벗어나는 극단 사건을 나타내는 파레토(중꼬리) 분포까지를
포함한다.} 관측성에 국한된 모델은 이러한 불확실성을 반영하지 못한다. \chg{관측성이
복구시간을 단축하는 방향으로 작용한 것과 달리, 불확실성은 복구시간을 늘려
MTPD에 근접시키는 방향으로 작용하며, 특히 극단을 나타내는
파레토 영역에서 그 증가 폭이 커진다.} 본 연구는 이 관측성 모델에 불확실성
차원을 더하여 모델을 완성하는 데에서 출발한다. 다음 절에서는 관측성과 불확실성의
두 측면에서 현행 접근의 한계를 선행연구에 비추어 구체화하되, 불확실성은 다시
가우시안과 파레토의 두 종류로 나누어 살펴보고, 이를 통합적으로 해소할 필요성을
논증한다.}


\section{연구 필요성}
\label{sec:problem_statement}

\chg{연구 배경에서 제기한 관측성과 불확실성의 두 측면은 업무연속성 연구에서
개별적으로는 다루어져 왔으나, 그동안의 연구 및 실무에서는 RTO를 결정하는 하나의 틀 안에서
통합적으로 반영되지는 못하였다. 그 결과 정적 RTO는 두 측면 모두에서 현실과
어긋나며, 이는 복구 자원의 과부족과 목표 시간의 신뢰성 저하로 이어진다.
따라서 관측성의 한계를 먼저 검토한 뒤, 불확실성을 일상적 변동의 가우시안과
극단 사건의 파레토라는 두 종류로 나누어 선행연구에서 없었던 방법으로 구체화하고, 이들이
결합될 때 개별적 보정만으로는 해결되지 않는다는 점에서 RTO 산정에 대한
통합적 접근의 필요성을 다음과 같이 세 가지로 제시한다.}

\chg{첫째, 일상적 변동을 나타내는 가우시안과 극단 사건을 나타내는 파레토의
두 불확실성을 관측성 모델에 통합할 필요가 있다. 관측성은
\citet{lee2026drto}의 동적 모델 연구에서 이미 반영한 차원이다. 이 모델은 관측성이 높을수록 이상 징후를 조기에 탐지하여
복구를 가속할 수 있다는 점을 시그모이드 바운딩으로 포착함으로써, 정적 RTO가
야기하던 관측성 높은 조건의 과다 예비(over-provisioning)와 낮은 조건의 과소
예비(under-provisioning)를 완화하였다~\citep{majors2022observability}. 그러나
이 모델은 관측성이라는 단일 차원에 국한되어, 복구 과정의 불확실성을 통합하지
못하였다. 이 불확실성은 일상적 변동을 나타내는 가우시안과 극단성을 나타내는
파레토의 두 가지로 나타나므로, 이 둘을 하나의 모델로 통합하여 반영하는 것이
본 연구가 해소하고자 하는 첫째 필요성이 된다.

둘째, 불확실성 정량화 부재 문제이다. 복구 과정에는 장비 가용성과 인력 동원
속도, 외부 공급망 의존도, 규제 대응 절차 등 다양한 요인이 작용하며, \chg{이들이
복구 시간에 미치는 영향은 재난의 종류에 따라 크게 달라진다. 예컨대 사이버공격은
데이터 복원과 시스템 재구성에, 자연재해는 물리적 설비와 인력 접근성에, 감염병
대유행은 인력 가용성과 외부 의존성에 각기 다른 부담을 지운다. 이처럼 같은
조직이라도 재난 유형과 상황에 따라 복구 소요가 달라지므로, 이 요인들은
본질적으로 확률적으로 변동한다.} 그럼에도 기존 RTO 프레임워크는 이러한
불확실성을 ``최악의 경우(worst case)''나 ``평균적 시나리오''라는 하나의
시나리오로 단순화할 뿐, 확률 분포에 기반한 정량화를 수행하지
않는다~\citep{cerullo2004bcp, gibb2006bcm}. 그 결과 BIA에서 설정한 RTO와 실제
복구 시간 사이에 구조적 괴리가 생기고, 조직은 그 괴리가 어느 방향으로 얼마나
벌어질지를 사전에 가늠할 수 없다. 다수 요인의 합으로 근사되는 일상적 복구
변동은 가우시안(정규) 분포로 표현될 수 있으므로, 불확실성을 확률 분포로
명시하면 RTO를 단일 점 추정이 아니라 분포의 함수로 다룰 수 있다. 이 차원은
관측성 기반 모델에서는 다루어지지 않았으므로, 일상적 복구 변동의 불확실성을
확률 분포로 정량화하여 복구목표시간 모형에 반영하는 일이 요구된다.

셋째, 극단 사건 미모형화 문제이다. 이는 앞의 일상적 변동과 더불어 불확실성이
갖는 또 다른 형태로, 불확실성을 확률 분포로 다루더라도 대부분의 접근이
정규분포만을 전제하는 데에서 비롯된다~\citep{coles2001extremes, embrechts1997extremal, taleb2007blackswan}.
그러나 사이버공격에 의한 데이터 센터 전면 장애, 복합 재난에 따른 공급망 동시
중단, 핵심 인력의 대규모 이동 제한과 같은 극단 사건은 정규분포의 가정을 크게
벗어난다. 이러한 사건의 복구 시간은 평균에서 멀어질수록 발생 확률이 급격히
줄지 않는 중꼬리(heavy-tail) 특성을 보이며, 극단값이론(Extreme Value Theory, EVT)에 따르면 파레토
분포가 이를 적절히 기술한다~\citep{coles2001extremes, embrechts1997extremal}. 중꼬리 분포에서는 고분위 영역의 복구 시간이
정규분포에 비해 현저히 커지므로, 정규분포를 전제한 RTO 산정은 극단 상황에서
실제 소요 시간을 크게 과소평가한다. 정적 RTO와 관측성 단일 모델은 모두 이 꼬리
위험(tail risk)을 구조적으로 포착하지 못하여, 업무연속성이 가장 절실한 극단
상황에서 보호에 실패할 수 있다. 따라서 정규분포의 범위를 넘어, 극단값이론에
근거하여 극단 사건을 중꼬리 분포로 별도로 모형화하고 그 꼬리 위험을
RTO에 반영하는 일이 요구된다.}

\nrev{\chg{나아가 관측성과 불확실성은 서로 독립적인(직교하는) 두 차원으로서
복구 시간에 서로 반대 방향으로 작용한다. 관측성이 높을수록 이상 징후를 일찍
탐지하여 복구 시간을 단축하는 반면, 불확실성이 클수록 복구 시간의 변동과 지연을
키운다. 특히 불확실성의 분포가 정규분포를 벗어나 중꼬리(파레토) 특성을 가지면
극단 영역에서 복구 시간이 급증하므로, 정규분포를 전제한 통상적 위험 평가는 그
타당성을 잃는다. 두 차원은 서로 독립적이더라도 복구 시간이라는 하나의 결과
위에서 함께 작용하므로, 각 차원을 따로 다루어서는 하나의 동적복구목표시간에 온전히
반영하기 어렵다.

이때 각 요인을 서로 분리하여 독립적·선형 보정량으로 추정한 뒤 경계 없이 단순
합산하는 접근(decoupled approach)은 다음과 같은 구조적 한계를 지니며, 바로
여기에서 본 연구의 필요성이 비롯된다. 첫째, 보정량을 경계 없이 더하는 방식은
복구 시간이 음수가 되거나 최대허용중단기간을 넘어서는 등 물리적으로 타당한
범위를 벗어날 수 있다. 둘째, 실무에서는 동적복구목표시간을 하나의 시간값으로
결정해야 하므로, 여러 보정 계수를 사후에 합산하는 구조는 의사결정 절차상
비현실적이며 책임 소재가 분산된다. 셋째, 각 차원을 서로 다른 척도와 절차로
분리해 보정하면, 관측성의 단축 효과와 불확실성의 연장 효과를 하나의 복구 시간
위에서 일관되게 비교하고 결합하기 어렵다. 요컨대 독립적·선형·무경계 보정의 단순 합산만으로는 두 차원을 물리적
경계 안에서 하나의 동적복구목표시간으로 통합할 수 없으며, 여기에서 통합적 접근의
필요성이 도출된다.}

\chg{다음 절에서는 이러한 필요성으로부터 도출되는 본 연구의 목적을
구체화한다.}}


\section{연구 목적}
\label{sec:research_objectives}

\chg{본 연구는 \citet{lee2026drto}이 제안한 관측성 기반
동적 복구목표시간($\DRTO$) 모델을 불확실성까지 반영하는 모델로 확장하고 그
타당성을 검증하는 것을 목적으로 한다. 기존 모델은 기준 복구시간에 관측성의
효과를 더하는 단일 채널 구조로, 관측성을 시그모이드 바운딩으로 반영하여
복구시간을 물리적으로 타당한 범위 안에 유지한다. 본 연구는 동일한 바운딩
원리를 불확실성에 적용하여, 관측성과 불확실성의 이중 채널 모델로 격상한다.
확장된 모델의 구체적 수리 구조와 도출 근거의 이론적 배경은
\chref{ch:theory}에서 설명한다. 이러한 세 가지 필요성(한계)에
대응하는 구체적 연구 목적은 다음과 같이 다섯 가지로 설정하였다.

첫째, 관측성 단일 채널 모델을 관측성과 불확실성의 이중 채널 $\DRTO$ 모델로
확장한다. 관측성을 복구 능력의 척도로 반영하고(관측성 미반영 문제 대응),
불확실성을 가우시안 분포로 정량화하며(정량화 부재 문제 대응), 일상적 변동을
넘는 극단을 파레토 분포로 포착한다(극단 미모형화 문제 대응). 두 채널은 동일한
시그모이드 바운딩 원리를 공유하므로, 불확실성 채널을 더한 뒤에도 복구시간이
물리적으로 타당한 범위를 벗어나지 않는 일관된 구조를 유지한다. 이로써 세 한계를
하나의 모형 안에서 동시에 해소한다.

둘째, 정적 기준에서 관측성과 불확실성을 단계적으로 도입하는 네 개 운영 조건
체계(C0--C3)를 설계한다. 이 조건 체계는 i) 정적 기준(C0), ii) 관측성(C1)
기준, iii) 관측성(C1)과 가우시안 불확실성(C2)을 함께 적용하는 기준, 그리고
iv) 관측성(C1)과 파레토 불확실성(C3)을 함께 적용하는 기준을 차례로 하여 각
요인이 추가될 때의 한계 기여를 분리하여 관측함으로써, 요인별 효과와 결합 효과를
명확히 귀속하고 비교할 수 있는 분석 틀이 된다.

셋째, 확장된 모델의 타당성을 검증할 연구 가설을 수립하고 체계적으로 검증한다.
가설은 모델의 핵심 효과와 결과의 강건성, 회귀 설명력, 이중채널 비선형 효과,
전문가 타당성의 다섯 측면을 포괄하며, 이를 대규모 몬테카를로 시뮬레이션과
전문가 평가를 결합한 다층 검증으로 확인한다.

넷째, 가우시안 불확실성(C2)과 파레토 불확실성(C3)을 각각 독립적으로 적용한
결과를 분포 전 구간에 걸쳐 비교하여, 극단 영역에서 나타나는 이중채널 비선형
감쇠 메커니즘을 규명하고 정량화한다. 두 조건은 평균 수준에서는 거의 구분되지
않으나 특정 백분위(약 P85.8)를 경계로 분포가 역전되어, 극단 영역에서는 파레토
불확실성이 복구 시간을 급격히 연장하는 한편 관측성 채널의 한계 효과는 그
영역에서 감쇠한다. 이러한 핵심 영역과 꼬리 영역 사이의 비선형 거동은 두 조건을
분포 전 구간에서 함께 비교·분석할 때에만 드러나며, 이는 확장 모델이 갖는 고유한
설명력에 해당한다.

다섯째, 검증된 모델을 바탕으로 제조업, 금융, 의료, 정보기술 서비스, 물류,
소매 등 주요 산업 분야에 적용할 수 있는 실무 가이드라인을 제시한다. 산업마다
허용 가능한 복구시간과 위험 노출이 다르므로, 가이드라인은 산업별 특성에 맞추어
모델의 매개변수를 조정하는 방식으로 제공한다.

\chg{요컨대 본 연구는 관측성에 국한된 기존 동적 복구목표시간 모델을 관측성과
불확실성의 이중 채널 모델로 확장하고, 이를 검증할 네 개 운영 조건 체계와 연구
가설을 수립하여 다층 검증으로 그 타당성과 극단 영역의 이중채널 비선형 감쇠를
확인하며, 검증된 모델을 산업별 실무 가이드라인으로 구체화하는 것을 목적으로
한다.}

이상의 목적에서 확장 모델의 수리적 구조와 매개변수, 조건 체계, 연구 가설의
도출은 \chref{ch:theory} 이론적 배경에서, 검증 절차와 다층 검증 체계는
\chref{ch:method} 연구 방법에서 선행 이론에 근거하여 상세히 전개한다.}


\section{연구 범위 및 방법}
\label{sec:research_scope}

\chg{이러한 목적을 달성하기 위해 본 연구의 범위와 방법을 다음과 같이 정한다.
모델의 범위는 정적 기준(C0)에서 관측성(C1), 가우시안 불확실성(C2), 파레토
불확실성(C3)에 이르는 네 개 운영 조건을 대상으로 한다. 이 네 조건은 앞서
설정한 이중 채널 모델의 구성 요소가 하나씩 활성화되는 단계에 대응하며,
복잡성이 점증하는 순서로 배열된다. 확장된 $\DRTO$ 모델은 관측성과 불확실성을
개별 시그모이드 바운딩으로 결합하는 구조이며, 그 수리적 정의와 매개변수 설정
근거는 \chref{ch:theory} 이론적 배경에서 전개한다.

검증 방법으로는 \citet{lee2026drto}에서 정립한 몬테카를로 시뮬레이션과
비중첩 대역 시드(band seeding) 체계를 따른다. 모델이 확률 분포를 입력으로
받으므로, 그 결과인 $\DRTO$의 분포를 하나의 공식으로 직접 구하기는
어렵다. 따라서 난수 표본을 많이 생성하여 그 결과를 통계적으로 추정하는
몬테카를로 방법이 적합하다. 본 연구는 이 방법론을 확장한 조건 체계에
적용하여, $n=10{,}000$ 규모의 시뮬레이션으로 조건별로 모델이 산출하는 결과를
통계적으로 확인하고, 여기에 업무연속성관리 및 재해복구
분야 전문가 평가를 결합하여 그 실무 타당성을 보완적으로 검증한다. 특히
시뮬레이션과 회귀, 전문가 평가 등 서로 다른 방법의 결과가 동일한 결론으로
수렴하는지를 확인하는 삼각검증(triangulation)을 통해 결론의 타당성을 높인다.
\chg{이러한 검증은 여섯 층위로 이루어진 다층 검증 체계(L0--L5)로 구조화된다.
먼저 기반 검증(L0)에서 모델 구현의 정합성과 물리적 경계를 결정론적으로
확인하고, 시뮬레이션 검증(L1)에서 관측성 효과와 경계 보장 등 모델의 기본
작동을 확인한다. 이어 회귀 검증(L2)으로 변수 간 설명력과 교호작용을, 강건성
검증(L3)으로 결과의 재현성을, 이중채널 검증(L4)으로 관측성과 극단 불확실성의
비선형 상호작용을 살핀 뒤, 전문가 평가(L5)로 실무 타당성을 확인한다. 각 층위의
구체적 정의와 절차는 \chref{ch:method} 연구 방법에서 제시한다.} 이를 통해
모델의 핵심 효과와 강건성, 실무 타당성을 포괄하는 연구 가설을 검증하며, 그
가설의 도출 근거는 \chref{ch:theory} 이론적 배경에서 상세히 기술한다.

한편 본 연구는 이론적 모델 수립과 시뮬레이션 검증, 전문가 타당성 평가에
초점을 맞추며, 다음 세 가지는 범위에 포함하지 않는다. 즉 i) 실제 조직
환경에서의 현장 실험(field test), ii) 시간에 따라 RTO가 연속적으로 변화하는
시계열 DRTO 모델, iii) 실제 단일 조직에 대한 심층 사례 연구(in-depth
single-case study)이다. 이들은 모델의 이론적 타당성과 여러 상황에서의 일반적
작동을 먼저 확립한다는 본 연구의 목적에 비추어 후속 과제로 두며, 검증된 모델은 이러한
현장 적용 연구의 토대가 된다.}



\section{논문 구성}
\label{sec:thesis_structure}

본 논문은 총 5개 장과 부록으로 구성되며, 각 장의 내용은 다음과 같다.

\vrev{%
\textbf{\chref{ch:introduction} 서론}에서는 연구 배경, 연구 필요성, 연구 목적, 연구 범위 및
방법, 논문 구성을 서술한다.

\textbf{\chref{ch:theory} 이론적 배경}에서는 업무연속성관리체계(BCMS)와
복구목표시간(RTO), 시스템 관측성 이론, 불확실성 정량화, 극단값
이론(EVT), 시그모이드 바운딩 기법에 관한 기존 연구를 검토하고
연구 격차를 도출한다. 이어서 개별 시그모이드 바운딩 모델의
수리적 정의, 4단계 조건 체계(C0--C3), 매개변수 설정($R_0 = 6.0$h,
$R_{\max} = 24.0$h, $\kappa = 1.2$), 시뮬레이션 설계 요소(비중첩
대역 시드 체계), 그리고 \jrev{10개} 가설(H1--\jrev{H10})을 수립한다.

\textbf{\chref{ch:method} 연구 방법}에서는 4단계 조건 체계(C0--C3)에 대응하는
연구 설계와 $n = 10{,}000$ 몬테카를로 시뮬레이션의 실행 절차,
분포 매개변수, 시드 배정 규칙, 통계량 산출 방법을 기술한다.

\textbf{\chref{ch:results} 연구 결과}에서는 C0--C3 조건별 시뮬레이션 결과, 회귀분석,
수렴 진단, 가설 검증을 순차적으로 보고하고, 조건 간 효과 크기 분석,
회귀 모델의 설명력 구조, 이중채널 비선형 메커니즘, 전문가 검증,
선행연구와의 교차 정합성을 제시한다.

\textbf{\chref{ch:discussion} 논의 및 결론}에서는 연구 결과에 대한 종합 논의와
\chg{이론적, 방법론적, 실무적} 시사점,
6개 주요 산업 가이드라인(기준선 제외)을 제시하고,
연구 한계와 주요 성과를 정리하며 향후 연구 방향을 제시한다.%
}

\textbf{부록}은 세 개 장으로 구성하였다. 첫째는 기호·용어·수식 목록(부록 A), 둘째는 조직 성숙도 자가 진단 체크리스트(부록 B), 셋째는 전문가 검증 상세 설문 구성(부록 C)이다. 전체 시뮬레이션 코드와 대규모 재현 자료는 공개
재현 패키지에 수록하여 GitHub에 공개한다.

\tabref{tab:chapter_overview}는 \chg{이상에서 기술한 장과 부록의 구성, 그리고 각 장과 가설(H1--H10) 및 분석 조건(C0--C3)의 대응 관계를 요약한 것이다}.

\begin{table}[htbp]
  \centering
  \caption{논문 구성과 가설/조건 대응 관계}
  \label{tab:chapter_overview}
  \small
  \begin{tabularx}{\textwidth}{@{}c >{\raggedright\arraybackslash}p{2.4cm} X >{\centering\arraybackslash}p{1.5cm}@{}}
    \toprule
    장 & 제목 & 주요 내용 & 관련 가설/조건 \\
    \midrule
    1 & 서론 & 배경, 연구 필요성, 연구 목적, 범위 및 방법 & --- \\
    2 & 이론적 배경 & BCMS, 관측성, 불확실성, EVT, 바운딩 기법, 모델 설계, C0--C3, H1--\jrev{H10} & H1--\jrev{H10}, C0--C3 \\
    3 & 연구 방법 & 연구 설계, MC(몬테카를로 시뮬레이션) $n{=}10{,}000$ 실행 절차 & C0--C3 \\
    4 & 연구 결과 & 조건별 결과, 회귀, 수렴, 가설 검증, 교차분석, 전문가 검증 & H1--\jrev{H10}, C0--C3 \\
    5 & 논의 및 결론 & 종합 논의, 이론적/방법론적/실무적 시사점, 6개 주요 산업 가이드라인(기준선 제외), 한계, 성과 요약, 향후 연구 & 전체 \\
    \addlinespace
    부록 & 기호·수식·용어 일람, 전문가 평가 설문 문항, BCMS 성숙도 자가 진단 체크리스트 & & --- \\
    \bottomrule
  \end{tabularx}
\end{table}
  % 제1장 서론
% 제2장 이론적 배경
% 1차 심사 반영: "이론적 배경 및 연구 프레임워크" → "이론적 배경" (축약)

\chapter{\vrev{이론적 배경}}
\label{ch:theory}

\vrev{%
\chg{1.2(연구의 필요성)에서 제기한 정적 RTO의 세 한계와 통합적 접근의 필요성에 따라서,
제2장(이론적 배경)은 동적 복구목표시간($\DRTO$) 모델과 그 조건 체계, 연구 가설을 선행
이론에 근거하여 처음으로 도출하고 확립한다.}
본 장에서는 BCMS와 RTO의 기존 한계,
관측성 이론, 불확실성 모델링, 시그모이드 함수와 바운딩 기법에 대한
선행연구를 검토하고, 본 연구의 이론적 토대를 정립한다.
이어서 관측성 \m{기반 $\DRTO$ 모델} 구조와 수학적 성질을 제시하고,
\jrev{10개 \m{연구 가설}}을 체계적으로 수립한다.
}

\section{BCM(업무연속성관리)}
\label{sec:lit-bcm}
%% ===================================================================

\chg{본 연구가 다루는 RTO가 자리하는 제도적 맥락을 분명히 하기 위해, BCM의 개념과 역사적 발전, 국제 표준 체계, 그리고 RTO가 도출되는 BIA 절차를 차례로 검토한다.}

\chg{본 연구에서 사용하는 BCM과 BCMS의 의미를 구분하여 밝혀 둔다. \m{BCM은} 조직이 재해와 장애, 위협으로부터 핵심 업무를
지속하고 복구하기 위한 관리 활동과 개념 전반을 가리킨다. 이에 비해
\m{BCMS는} 이러한 BCM 활동을
정책과 절차, 자원, 문서, 평가의 형태로 체계화한 경영 시스템을 가리키며,
ISO 22301:2019 인증의 직접 대상이 된다. 따라서 본 연구는 일반적 활동이나
개념을 가리킬 때는 BCM을, 정형화된 시스템이나 표준, 인증을 가리킬 때는 BCMS를
사용하여 일관되게 구분한다.}
\citet{herbane2010disaster}\m{은} BCM의 역사적 발전 과정을 세 단계로 구분하였다.
첫째, 1970--1980년대의 IT 재해복구(Disaster Recovery) 중심 단계에서는 데이터 백업과
시스템 이중화 등 기술적 복구에 초점을 두었다.
둘째, 1990년대에는 \m{BIA의 도입}과 함께
IT를 넘어 전사적 업무 프로세스로 관리 범위가 확대되었다.
셋째, 2000년대 이후에는 \citet{bhamra2011resilience}이 정리한
\textbf{조직 복원력}(organizational resilience) 관점에서
BCM을 전략적 경영 기능으로 통합하는 방향으로 진화하였다.

이러한 발전의 제도적 결실로서 영국 BSI는 2006년에 BS~25999를 발표하여
BCM의 국가 표준 단계의 핵심 프레임워크를 제시하였다\citep{bs25999}.
이후 국제표준화기구(ISO)는 2012년에 \nrev{ISO 22301:2012}를 제정하였으며,
2019년에 개정판을 공표하여 현재의 국제 표준으로 자리매김하였다.
\nrev{ISO 22301:2019}는 BCMS의 수립·구현·유지·지속적 개선을 위한 요구사항을 \m{규정하며}, Plan--Do--Check--Act(PDCA) 주기에
기반한 관리체계를 \m{요구하고 있다}.
\nrev{ISO 22313:2020}은 이에 대한 실행 지침서로서 BIA 수행 방법과
복구 전략 수립 절차를 구체적으로 \m{안내한다}.

\figref{fig:bcm-pdca}\refpo{fig:bcm-pdca}{은}{는} \nrev{ISO 22301:2019}의 PDCA 주기에서
RTO가 결정되는 위치를 도식화한 것이다. 정적 RTO는 Do(운영) 단계의 BIA에서 한 번
결정되어 다음 BIA까지 고정되는 반면, 본 연구의 $\DRTO$는 운영과 모니터링이
이루어지는 Do--Check 구간에서 실시간으로 파악된다.

\begin{figure}[htbp]
  \centering
  \begin{tikzpicture}[
    phase/.style={draw, semithick, rounded corners=8pt, minimum width=34mm, minimum height=16mm,
                  align=center, font=\small, inner sep=4pt},
    cyc/.style={-{Stealth[length=3.2mm]}, very thick, gray!55!black},
    annot/.style={draw=red!55, fill=red!7, semithick, rounded corners=4pt,
                  align=center, font=\footnotesize, inner sep=5pt},
    darr/.style={-{Stealth[length=2.8mm]}, very thick, dashed, red!70!black}
  ]
    % --- PDCA 순환 (시계방향: 상→우→하→좌) ---
    \node[phase, fill=blue!10]   (plan)  at (0,  2.6) {\textbf{Plan}\\[1pt]\footnotesize BCMS 수립};
    \node[phase, fill=green!13]  (do)    at (4.2, 0) {\textbf{Do}\\[1pt]\footnotesize BIA 수행\\\footnotesize \textbf{정적 RTO 결정}};
    \node[phase, fill=yellow!22] (check) at (0, -2.6) {\textbf{Check}\\[1pt]\footnotesize 모니터링, 측정};
    \node[phase, fill=orange!14] (act)   at (-4.2, 0) {\textbf{Act}\\[1pt]\footnotesize 지속적 개선};
    \node[font=\small\bfseries, align=center] at (0,0) {ISO 22301:2019\\PDCA 주기};
    \draw[cyc] (plan.east)  to[bend left=22] (do.north);
    \draw[cyc] (do.south)   to[bend left=22] coordinate[midway] (dc) (check.east);
    \draw[cyc] (check.west) to[bend left=22] (act.south);
    \draw[cyc] (act.north)  to[bend left=22] (plan.west);
    % --- 정적 RTO vs 본 연구의 DRTO (우측 분리 배치, 직각 연결) ---
    \node[annot] (static) at (9.3, 2.6)
      {\textbf{정적 RTO}\\BIA에서 결정};
    \draw[darr] (static.west) -| (do.north);
    \node[annot, fill=red!11] (drto) at (9.3, -2.6)
      {\textbf{본 연구: $\DRTO$}\\운영--모니터링 중\\실시간 동적 파악};
    \draw[darr] (drto.west) -| (dc);
  \end{tikzpicture}
  \caption{\nrev{ISO 22301:2019} PDCA 주기 기반 RTO 결정 위치}
  \label{fig:bcm-pdca}
\end{figure}


따라서 BIA는 BCMS의 핵심 구성 요소로서 중단 사태가 조직에 미치는 영향을 분석하고,
핵심 업무 기능의 복구 우선순위와 복구 목표 시간을 결정하는 기초를 제공한다
\citep{cerullo2004bcp, snedaker2013bcdr}.
\citet{gibb2006bcm}은 BIA 결과가 조직의 복구 전략과 자원 배분에 직접적으로 연결되어야
함을 \m{강조하며}, 정보관리 관점에서의 BCM 프레임워크를 제안하였다.

RTO는 중단 사태 발생 이후 핵심 업무 기능을
사전에 정의된 수준까지 복구하는 데 소요되는 목표 시간을 의미한다.
\nrev{ISO 22301:2019}에서는 RTO를 ``제품이나 서비스의 납품 또는 수행, 활동의 재개,
또는 자원의 복구가 이루어져야 하는 시점까지의 기간''으로 정의한다.
미국 NIST SP~800-34~Rev.~1에서는 RTO를 ``시스템 중단 이후 시스템과 관련 자원이
복구되어야 하는 최대 허용 시간''으로 \m{규정하고 있다}\citep{nist800-34r1}.

MTPD는 조직이 중단 사태를
견딜 수 있는 최대 허용 기간을 의미하며, RTO는 반드시 MTPD 이내에 설정되어야 한다.
즉, $\text{RTO} \leq \text{MTPD}$의 관계가 성립한다\nrev{(ISO 22313, 2020; \citeauthor{snedaker2013bcdr}, \citeyear{snedaker2013bcdr})}.
\citet{hiles2010bcm}은 BIA에서 도출된 MTPD가 RTO 설정의 상한 기준임을
강조하였다.

국내 BCM 연구는 \chg{BIA, 위험평가, 평가지표} 측면에서 점진적 진전을 이루었으나, RTO의 동적 설정 문제는 다루지 않았다.
\citet{cheung2019quickhit}는 BIA 및 위험평가를
효율적으로 수행하기 위한 Quick-Hit 프레임워크를 제안하여 BCMS 도입 초기의
실무적 장벽을 낮추고자 하였다. \citet{jang2021bcms}은 BCMS의 위험평가가
경영성과에 미치는 영향을 실증 분석하여 BCM과 조직 성과 간의
정(+)의 관계를 확인하였으며, \citet{jung2023bcms}은 BCMS의 체계적 평가를 위한
평가지표를 개발하여 국내 조직의 BCM 성숙도 진단에 기여하였다.

\chg{이처럼 BCM 연구는 체계 도입과 성과 평가로 꾸준히 확장되어 왔으나, 그 핵심 통제 변수인 RTO를 어떻게 설정할 것인가의 문제는 여전히 정적 관점에 머물러 있다. 2.2(RTO의 정적 한계)에서는 이 정적 RTO 설정이 지닌 구조적 한계를 살펴본다.}
  % 2.1 업무연속성관리(BCM)


%% ===================================================================
\section{RTO의 정적 한계}
\label{sec:lit-rto-static}
%% ===================================================================

\chg{앞 절에서 살펴본 BCM의 핵심 산출물 가운데 하나가 RTO이다. 그러나 현행 RTO는 복구 환경의 변화와 무관하게 단일 정적 값으로 고정된다는 구조적 제약을 안고 있다. 이 절에서는 그 한계를 정적 설정의 구조적 문제와 선행연구의 접근을 중심으로 검토한다.}

%% -------------------------------------------------------------------
\subsection{정적 RTO 설정의 구조적 문제}
\label{subsec:static-rto-problem}
%% -------------------------------------------------------------------

전통적 RTO 설정 방식에는 세 가지 근본적 한계가 존재하며, 이는 본 논문
\secref{sec:problem_statement}(연구 필요성)에서 제시한 세 가지 연구 필요성과 직접 대응된다.

\chg{첫째는 관측성 미반영이다. 정적 RTO는 BIA에서 도출된 고정값으로서 조직의 운영
가시성, 곧 관측성(observability)을 반영하지 못한다. 시스템 복잡도와 인력 역량, 기술 환경 등
운영 환경이 시간과 상황에 따라 변화하는 현실에서 관측성이 높은 시스템은 장애 원인을 신속히
파악하여 복구 시간을 단축할 수 있으나, 전통적 BIA는 이러한 운영 역량 차이를 RTO에
반영하는 메커니즘을 갖추지 못한다. 그 결과 관측성이 높은 조건에서는 복구 자원의 과다
예비(over-provisioning)가, 낮은 조건에서는 과소 예비(under-provisioning)가 동시에 초래된다.
둘째는 일상적 불확실성의 정량화 부재이다. 복구 과정에서 발생하는 장비 가용성과 인력 동원 속도,
외부 공급망 의존도 같은 일상적 불확실성이 정량적으로 모델링되지 않는다. 실제 복구 시간은 다양한
예측 불가능한 요인에 의해 목표 시간을 초과할 수 있으나, 기존 접근법은 이러한 정규분포(Gaussian)적
변동을 체계적으로 반영하지 못한다. 셋째는 극단 사건의 미모형화이다. 사이버공격에 의한 데이터
센터 전면 장애나 복합 재난으로 인한 공급망 동시 중단처럼 일상적 변동을 넘는 극단적 사건에 의한
복구 지연은, 정규분포 가정으로는 포착할 수 없는 파레토(Pareto, 중꼬리, heavy-tail) 분포 특성을
보인다~\citep{coles2001extremes, embrechts1997extremal, taleb2007blackswan}.}


%% -------------------------------------------------------------------
\subsection{선행 연구에서의 RTO 접근}
\label{subsec:prior-rto-approaches}
%% -------------------------------------------------------------------

기존의 RTO 관련 연구는 대체로 BIA 방법론의 개선, RTO 달성을 위한 기술적 수단
(이중화, 백업 전략), 또는 특정 산업 분야에서의 사례 분석에 집중하고 있다.
\citet{snedaker2013bcdr}는 IT 재해복구와 업무연속성 계획 수립의 체계적 방법론을
제시하였으나, RTO 자체를 동적으로 조정하는 수리 모델은 다루지 않았다.
\citet{cerullo2004bcp}는 금융 서비스 산업에서의 BCP 수립 사례를 분석하여
RTO의 실무적 중요성을 강조하였으나, RTO의 동적 특성에 대한 논의는 부재하였다.

\chg{최근에는 업무연속성과 재해복구를 통합하여 조직 복원력 관점에서 정량적으로 접근하려는 시도가 이어지고 있다. \citet{sahebjamnia2015integrated}는 핵심 업무의 재개와 복구를 동시에 고려하는 다목적 혼합정수계획 모형으로 복구 자원의 최적 배분 문제를 정식화하였고, \citet{torabi2016enhanced}는 최선-최악 방법(Best--Worst Method)에 기반한 향상된 위험평가 체계를 제안하여 BCM의 위험 정량화 수준을 끌어올렸다. 그러나 이들 정량적 프레임워크에서도 RTO 자체는 위험평가와 자원배분의 고정된 입력 매개변수로 취급되며, 운영 환경의 관측성이나 복구 시간의 확률적 변동에 따라 RTO를 동적으로 조정하는 메커니즘은 여전히 다루어지지 않는다.}

\citet{cheung2019quickhit}, \citet{jang2021bcms} 그리고 \citet{jung2023bcms}을 포함한 국내 BCM
연구 또한 표준 적용과 조직 성과 평가에 집중하고 있어, RTO의 동적 조정이나
운영 조건에 따른 적응적 복구 모델에 대한 연구가 상대적으로 부재함을 드러낸다.
\citet{lee2026drto}이 관측성 채널(C0$\to$C1)을 도입하여
초기 접근을 수행하였으나, 불확실성 채널(C2/C3)의 체계적 통합은 향후 과제로
남아있었다.
이러한 세 가지 한계와 향후과제가 본 연구에서 $\DRTO$ 프레임워크를 통해 해소하고자 하는
핵심 문제의식의 출발점이다.

\chg{이러한 적응적 복구를 실현하려면 복구 환경의 상태를 실시간으로 관측할 수 있어야 하며, 다음 절에서는 그 이론적 토대인 관측성 개념을 검토한다.}
  % 2.2 복구목표시간(RTO)의 정적 한계


%% ===================================================================
\section{관측성(Observability) 이론}
\label{sec:lit-observability}
%% ===================================================================

\chg{정적 RTO의 한계를 극복하려면 복구 환경의 상태 변화를 실시간으로 파악할 수 있어야 한다. 이러한 상태 파악 능력을 이론적으로 뒷받침하는 개념이 관측성(observability)이며, 이 절에서는 제어이론에서 출발한 관측성 개념과 그 현대적 확장을 검토한다.}

%% -------------------------------------------------------------------
\subsection{제어 이론에서의 관측성 (Kalman, 1960)}
\label{subsec:kalman-observability}
%% -------------------------------------------------------------------

관측성(observability)의 개념은 \citet{kalman1960control}이 제어 이론에서 체계적으로 형식화하였다.
Kalman은 선형 시불변(Linear Time-Invariant, LTI) 시스템
\begin{equation}
  \dot{x}(t) = Ax(t) + Bu(t), \quad y(t) = Cx(t)
  \label{eq:lti-system}
\end{equation}
에서 시스템의 출력 $y(t)$를 관찰함으로써 초기 상태 $x(0)$를 유일하게 결정할 수 있는지를
판별하는 기준으로 관측성 행렬(observability matrix)의 계수(rank) 조건을 제시하였다.
관측성 행렬
\begin{equation}
  \mathcal{O} = \begin{bmatrix} C \\ CA \\ CA^2 \\ \vdots \\ CA^{n-1} \end{bmatrix}
  \label{eq:obs-matrix}
\end{equation}
이 완전 계수(full rank)일 때 시스템이 완전 관측 가능(completely observable)하다.

이 개념은 시스템의 내부 상태를 외부 출력만으로 추론할 수 있는 역량을 수학적으로
정의한 것으로, 단순한 모니터링(monitoring)과는 본질적으로 구별된다~\citep{beyer2016sre}.
모니터링이 사전에 정의된 지표의 임계값 초과 여부를 감시하는 반응적(reactive) 방식인 반면,
관측성은 시스템 출력의 충분성(sufficiency)에 관한 구조적 속성이다.

\chg{Kalman 이후 제어 이론에서 관측성은 상태 추정(state estimation)의 토대로 발전하였다. \citet{luenberger1966observers}는 관측성 조건이 충족되는 선형 시스템에서 직접 측정되지 않는 내부 상태를, 가용한 출력과 입력만으로 구동되는 보조 동역학계인 상태 관측기(state observer)를 통해 재구성할 수 있음을 보였다. 이로써 관측성은 가부 판별을 넘어 관측 가능한 시스템에서는 내부 상태를 실제로 복원할 수 있다는 구성적 의미를 획득하였고, 이후 상태 추정과 피드백 제어의 핵심 원리로 자리 잡았다. 본 연구가 관측성을 복구 환경의 내부 상태에 대한 가시성으로 해석하는 것은, 바로 이 출력으로부터 내부 상태를 복원하는 역량이라는 제어 이론적 함의를 BCM 영역으로 확장한 것이다.}


%% -------------------------------------------------------------------
\subsection{BCM 영역으로의 확장}
\label{subsec:observability-bcm}
%% -------------------------------------------------------------------

%% --- BCM 맥락의 정규화 ---
\subsubsection{BCM 맥락에서의 관측성 정규화}
\label{subsubsec:obs-normalization}

본 연구에서는 Kalman의 제어 이론적 관측성 개념을 BCM 영역으로 확장하여 적용한다.
\chg{Kalman의 원래 정식화는 관측성 행렬의 계수 조건을 통해 시스템이 ``완전 관측
가능한가''를 가능과 불가능의 이분법으로 판별한다. 그러나 \chg{BCM과 재해복구} 실무에서 조직의 관측성은 이분법으로 주어지지 않는다.
\chg{현대 IT 운영에서 관측성은 로그(log), 메트릭(metric), 트레이스(trace)라는 세 가지
텔레메트리(telemetry)를 통해 확보되며~\citep{majors2022observability}, 조직마다 센서와
모니터링 도구의 구비 정도, 그리고 이 세 텔레메트리의 수집과 보존 수준이 달라 그 확보
정도가 연속적으로 변한다.} 그 결과 조직의 관측성은 0과 1 사이에서 연속적으로 분포한다. 본 연구는 바로 이 부분을 확장하여 적용하는데, Kalman의
이진 판별 대신 관측성의 \emph{정도}를 $O_{\text{raw}} \in [0, 1]$의 연속 점수로 정량화한다.
이렇게 연속화하는 이유는, 복구 시간의 단축 효과가 관측성의 유무가 아니라 그
수준에 연속적으로 의존한다고 보기 때문이며, 연속 점수라야 시그모이드 바운딩을 통한 복구 시간
단축량의 정량적 모형화가 가능해지기 때문이다.}
이때 $O_{\text{raw}} = 0$은 시스템 상태를 전혀 파악할 수 없는 상태를,
$O_{\text{raw}} = 1$은 완전히 관측 가능한 상태를 의미한다.


%% --- 관측성-복구 능력 구조적 연결 ---
\subsubsection{관측성과 복구 능력의 구조적 연결}
\label{subsubsec:obs-recovery-link}

\citet{avizienis2004basic}와 \citet{hollnagel2006resilience}이 각각
의존성과 복원공학 문헌에서 논의한 바와 같이, 관측성과 복구 능력 사이에는
구조적 연결이 존재한다.
시스템 상태를 보다 정확히 파악할 수 있을수록 장애의 근본 원인(root cause)을
신속히 진단할 수 있고, 이에 따라 복구에 소요되는 시간이 단축된다.
\citet{avizienis2004basic}는 가용성(availability), 신뢰성(reliability),
안전성(safety), 무결성(integrity)을 포괄하는 의존성(dependability)의 분류 체계에서
장애 진단 능력이 시스템의 가용성에 직접적으로 영향을 미침을 논의하였다.
\citet{hollnagel2006resilience}는 복원공학(Resilience Engineering) 관점에서 시스템의
예측, 감시, 대응, 학습 능력을 강조하며, 이 중 감시(monitoring) 능력이 본질적으로
관측성에 해당한다고 볼 수 있다.
\chg{본 연구는 이 두 이론을 관측성과 복구 시간 사이의 단조 감소 관계로 반영하여 적용한다.
\citet{avizienis2004basic}에서는 장애 진단 능력이 가용성에 기여한다는 인과를,
\citet{hollnagel2006resilience}에서는 감시 능력이 복원력의 핵심이라는 명제를 각각
가져온다. 본 모델은 이 ``관측성에서 진단으로, 진단에서 복구 시간 단축으로''의 인과를
관측성 채널의 하향 조정량 $E_{O,\text{bnd}} \leq 0$으로 구현한다. 즉 관측성이 높을수록
$\DRTO$를 기준 RTO $R_0$ 아래로 끌어내리는 음의 효과로 정량화하는 것이다. 두 이론을
이렇게 적용한 이유는, 이론이 제시한 인과 관계를 모델에서 관측성과 복구 시간 사이의
단조 감소 관계로 직접 반영함으로써, 관측성 투자에 따른 복구 시간 단축을 정량적으로
예측할 수 있게 하기 위함이다.}


%% --- SRE와 세 가지 핵심 요소 ---
\subsubsection{\chg{현대 IT 운영의 관측성 SRE와 세 가지 핵심 요소(three pillars)}}
\label{subsubsec:obs-sre-pillars}

현대 IT 운영에서의 관측성은 Kalman의 제어 이론적 개념을 분산 시스템
환경으로 재해석한 것으로 볼 수 있으며, \citet{beyer2016sre}와
\citet{majors2022observability}이 이를 SRE 실무 체계로 정착시켰다.
\citet{beyer2016sre}는 Google의 SRE(Site Reliability Engineering)
실무를 통해 모니터링과 관측성의 차이를 구분하였고,
\citet{majors2022observability}는 관측성의 세 가지 핵심 구성 요소(three pillars), \chg{곧
메트릭(metrics)과 로그(logs), 트레이스(traces)를} 체계화하였다.
이 세 구성 요소의 통합적 활용을 통해 조직은 예상치 못한 장애 상황에서도
시스템의 내부 상태를 효과적으로 파악할 수 있다.
\p{[}그림~\ref{fig:obs-three-pillars}\p{]}는 이러한 세 가지 핵심 요소(three pillars)의 구성 관계를 도식화한 것이다.


%% --- 운영화 매핑 ---
\subsubsection{\chg{데이터 유형에서} \chg{도구와 목적} \chg{차원으로의 운영화 매핑}}
\label{subsubsec:obs-operationalization}

이러한 세 가지 핵심 요소(three pillars) 분류는 관측성을 구성하는 \textbf{데이터 유형(data primitives)} 차원의
정의이다. 본 연구는 \chg{BCM과 재해복구} 맥락의 운영화 단계에서 관측성을 \chg{도구와 목적}별 통합 시스템
차원, \chg{곧 APM(Application Performance Monitoring, 서비스 신뢰성 영역)과 SIEM(Security
Information and Event Management, 보안 이벤트 대응 영역)의 두 축으로} 재구성하여 정량화한다.
APM과 SIEM은 각각 \chg{메트릭, 로그, 트레이스}를 통합 활용하되 측정 목적과 주된 데이터 출처가
상이하므로 분리 산출 후 가중 합성($O_{\text{raw}} = \alpha \cdot O_{\text{APM}}
+ (1-\alpha) \cdot O_{\text{SIEM}}$)으로 통합하며, 이론적 분류와 운영화 모델의 매핑 관계는
본문 \secref{app:operationalization}(제5장 실무적 기여)의 조작화 예시에서 참고용으로 다룬다.

\begin{figure}[htbp]
\centering
\begin{tikzpicture}[
  every node/.style={text=black},
  pillar/.style={draw, thick, rounded corners=4pt, fill=blue!10, minimum width=40mm,
    minimum height=22mm, font=\Large\bfseries, align=center, text=black},
  obs/.style={draw, very thick, rounded corners=4pt, fill=green!15, minimum width=64mm,
    minimum height=18mm, font=\Large\bfseries, align=center, text=black},
  arr/.style={->, >=Stealth, very thick, black}
]
  \node[pillar] (m) at (0, 0)   {Metrics\\\normalsize\mdseries(메트릭)};
  \node[pillar] (l) at (5.5, 0) {Logs\\\normalsize\mdseries(로그)};
  \node[pillar] (t) at (11, 0)  {Traces\\\normalsize\mdseries(트레이스)};
  \node[obs]    (o) at (5.5, -3.4)
    {Observability\\\normalsize\mdseries(관측성)};
  \draw[arr] (m.south) -- (o.north);
  \draw[arr] (l.south) -- (o.north);
  \draw[arr] (t.south) -- (o.north);
\end{tikzpicture}
\caption{관측성의 세 가지 핵심 구성 요소(three pillars)}
\label{fig:obs-three-pillars}
\end{figure}

관측성 수준이 높은 조직은 장애 발생 시 원인 파악에 소요되는 시간(Mean Time to Detect,
MTTD)과 복구에 소요되는 시간(Mean Time to Recovery, MTTR)을 모두 단축할 수 있다~\citep{forsgren2018accelerate, kim2016devops}.
이는 관측성이 단순한 기술적 역량을 넘어 조직의 복구 목표 달성에 직접적으로
기여하는 운영 변수임을 시사한다.
본 연구에서는 이러한 이론적 근거를 바탕으로 관측성 점수 $O_{\text{raw}} \in [0,1]$을
$\DRTO$ 모델의 핵심 입력 변수로 도입한다.


%% -------------------------------------------------------------------
\subsection{모니터링과 관측성의 개념적 구분}
\label{subsec:monitoring-vs-observability}
%% -------------------------------------------------------------------

실무에서 모니터링(monitoring)과 관측성(observability)은 혼용되는 경우가 많으나,
BCM 맥락에서 양자는 분석 범위와 목적이 구별된다. 모니터링은 사전에 정의된
지표를 수집하여 임계값 초과 여부를 감시하는 활동인 반면, 관측성은 시스템 외부
출력만으로 내부 상태를 추정할 수 있는 시스템 고유의 속성이다. 즉 모니터링이
관측성을 구현하기 위한 수단(도구와 프로세스)이라면, 관측성은 그 수단이 확보하는
결과(시스템의 파악 가능 정도)에 해당한다.
\p{[}표~\ref{tab:monitoring-vs-observability}\p{]}은 두 개념의 핵심 차이를 비교한 것이다.

\begin{table}[htbp]
\centering
\caption{모니터링과 관측성의 비교}
\label{tab:monitoring-vs-observability}
\small
\begin{tabularx}{\textwidth}{l X X}
\toprule
\textbf{비교 기준} & \textbf{모니터링 (Monitoring)} & \textbf{관측성 (Observability)} \\
\midrule
정의 &
  사전에 정의된 지표(메트릭, 로그, 알림 규칙)를 수집·추적하는 활동 &
  시스템 외부 출력만으로 내부 상태를 추정할 수 있는 시스템 고유 속성 \\
\addlinespace
관점 &
  ``무엇을 측정할 것인가'' (도구·프로세스 중심) &
  ``시스템이 얼마나 파악 가능한가'' (시스템 능력 중심) \\
\addlinespace
범위 &
  사전 정의된 이상 징후 탐지 (known-unknowns) &
  예상하지 못한 장애의 근본 원인 추적 포함 (unknown-unknowns) \\
\addlinespace
BCM 역할 &
  장애 감지 $\to$ 알림 발생 $\to$ 대응 절차 개시 &
  복구 과정에서 시스템 상태 가시성 확보 $\to$ 복구 시간 단축 \\
\addlinespace
$\DRTO$ 대응 &
  모니터링 도구의 존재 자체는 $O_{\mathrm{raw}}$의 필요조건 &
  $O_{\mathrm{raw}} \in [0,1]$로 정량화, 관측성 채널($E_{O,\mathrm{bnd}}$)에 직접 반영 \\
\bottomrule
\end{tabularx}
\end{table}

$\DRTO$ 모델에서는 모니터링 인프라의 수준이 $O_{\mathrm{raw}}$ 값으로 변환되어
관측성 채널에 반영되므로, 모니터링 투자가 곧 $\DRTO$ 감소로 이어지는 정량적
경로를 제공한다. 이러한 관측성의 차이는 실무 환경에서도 뚜렷하게 나타난다.
관측성이 높은($O_{\mathrm{raw}} \approx 1$) 환경의 전형적 예는
실시간 서버 모니터링 시스템(Prometheus, Grafana)이 구축된 IT 인프라로,
장애 발생 시 원인을 수 분 내에 특정하여 복구 시간을 단축하며,
금융권의 실시간 트랜잭션 모니터링이나 SRE 조직의 가시성 대시보드가 여기에 해당한다.
반대로 관측성이 낮은($O_{\mathrm{raw}} \approx 0$) 환경에서는
센서가 미비한 노후 제조 설비나 수동 점검에 의존하는 중소기업 IT 환경처럼
장비 고장의 근본 원인 파악에 수 시간이 소요되어 복구가 지연된다.

\chg{이처럼 관측성은 복구 환경의 가시성을 좌우하여 복구 시간에 직접 영향을 미치지만, 복구 시간에 내재한 변동성 자체를 설명하지는 못한다. 다음 절에서는 이 변동성을 정량화하는 불확실성 모델링을 검토한다.}
  % 2.3 관측성(Observability) 이론


%% ===================================================================
\section{불확실성(Uncertainty) 모델링}
\label{sec:lit-uncertainty}
%% ===================================================================

\chg{관측성이 복구 환경의 가시성을 높여 주더라도, 복구 소요 시간 자체에 내재한 확률적 변동까지 제거하지는 못한다. 이 절에서는 이러한 변동을 설명하는 불확실성 모델링을 일상적 변동을 다루는 가우시안 분포와 극단 사건을 다루는 파레토 분포를 중심으로 검토한다.}

%% -------------------------------------------------------------------
\subsection{가우시안 불확실성 분포}
\label{subsec:gaussian-uncertainty}
%% -------------------------------------------------------------------

불확실성 정량화의 가장 기본적인 도구는 가우시안(정규) 분포이다.
\citet{jaynes2003probability}는 확률 이론의 논리적 기초에 관한 체계적 논의를 통해
중심극한정리(Central Limit Theorem, CLT)의 정보이론적 정당성을 제시하였다.
다수의 독립적인 확률 변수의 합이 가우시안 분포에 수렴한다는 CLT는,
실제 복구 과정에서 발생하는 다양한 소규모 불확실성 요인, \chg{곧 인력 가용성과
장비 상태, 네트워크 지연, 절차적 오류 등의} 복합적 효과가
가우시안 형태로 모델링될 수 있음을 이론적으로 뒷받침한다.

본 연구에서는 일상적 운영 불확실성을 가우시안 분포 \chg{$U^{(G)} \sim N(\mu, \sigma^2)$}로
모델링한다. 이 분포는 대칭적 특성을 가지므로, 복구 시간의 단축과 연장 가능성을
동등하게 반영할 수 있다. 구체적으로 $\DRTO$ 프레임워크에서는
$U_{\text{raw}} \sim N(3, 1)$의 원시 불확실성 값을 생성한 후
$z_U = \kappa \cdot k_U \cdot R_0 \cdot U_{\text{raw}} / (R_{\max} - R_0)$으로
시그모이드 입력값을 산출한다.
\chg{가우시안을 C2 조건에 채택한 근거는 앞서 언급한 중심극한정리에 있다. 복구 과정에서
작용하는 다수의 독립적 소규모 불확실성 요인이 누적될 때 그 합이 가우시안에 수렴하므로,
어느 하나의 지배적 요인 없이 여러 작은 변동이 합산되는 일상적 운영 상황을 가우시안으로
표현하는 것이 이론적으로 타당하다.}
그러나 가우시안 분포는 꼬리 확률(tail probability)이 지수적으로 감소하므로,
대규모 중단이나 극단적 지연과 같은 드문 사건(rare event)의 발생 가능성을
구조적으로 과소평가하는 한계를 지닌다.


%% -------------------------------------------------------------------
\subsection{파레토 분포}
\label{subsec:heavy-tail}
%% -------------------------------------------------------------------

현실 세계의 재해 및 중단 사태는 가우시안 분포가 예측하는 것보다 훨씬 빈번하게
극단적 양상을 보인다. \citet{taleb2007blackswan}은 ``블랙 스완''의 개념을 통해
극도로 낮은 확률이지만 발생 시 막대한 영향을 미치는 사건들이 금융, 기술, 자연재해
영역에서 체계적으로 과소평가되고 있음을 경고하였다.

멱법칙(power-law) 분포는 이러한 극단적 사건을 수학적으로 표현하는 핵심 도구이다.
\citet{newman2005powerlaws}은 자연과학, 사회과학, 기술 시스템에서 발견되는
멱법칙 분포의 보편성을 체계적으로 검토하고, 그 확률밀도함수가
\begin{equation}
  f(x) = \frac{\alpha \, x_m^{\alpha}}{x^{\alpha+1}}, \quad x \geq x_m
  \label{eq:pareto-pdf}
\end{equation}
의 형태를 가지는 파레토(Pareto) 분포로 대표됨을 보였다.
여기서 $\alpha$는 형상 모수(shape parameter) 또는 꼬리 지수(tail index)로서,
분포의 꼬리 두께와 모멘트 존재 조건을 결정하는 핵심 파라미터이다.
$\alpha$가 작을수록 꼬리가 두꺼워져 극단값의 출현 빈도가 높아진다.

%% --- 형상 모수 α의 역할 ---
\subsubsection{형상 모수 $\alpha$의 역할}
\label{subsubsec:pareto-alpha}

파레토 분포의 통계적 성질은 형상 모수 $\alpha$의 값에 따라 질적으로 달라진다.
$n$차 모멘트 $E[X^n]$이 유한하기 위해서는 $\alpha > n$이어야 하며,
\chg{이로부터 다음과 같은 임계 구간이 도출된다. $\alpha$가 1 이하이면 평균조차 발산하여
모든 모멘트가 무한하고, 1과 2 사이이면 평균은 정의되나 분산이 발산한다. 2와 3 사이이면
평균과 분산은 유한하나 왜도가 발산하여 분포 형태가 극도로 비대칭이 되고, 3과 4 사이이면
왜도까지 유한하나 첨도가 발산하여 정규분포 대비 극단값의 빈도가 현저히 높아진다.
$\alpha$가 4를 넘으면 4차 이하의 모든 모멘트가 유한해져 꼬리가 상대적으로 얇아지면서
가우시안에 근사한다.}
\p{[}표~\ref{tab:pareto-alpha-properties}\p{]}는 대표적 $\alpha$ 값에 따른
파레토 분포($x_m = 1$)의 통계적 성질을 정리한 것이다.

\begin{table}[htbp]
  \centering
  \caption{형상 모수 $\alpha$에 따른 파레토 분포의 통계적 성질 ($x_m = 1$)}
  \label{tab:pareto-alpha-properties}
  \small
  \begin{tabular}{ccccc}
    \toprule
    $\alpha$ & 평균 $E[X]$ & 분산 $\text{Var}[X]$ & 왜도 & 첨도 \\
    \midrule
    1 & $\infty$ & $\infty$ & --- & --- \\
    2 & 2.000 & $\infty$ & $\infty$ & --- \\
    \textbf{3} & \textbf{1.500} & \textbf{0.750} & $\boldsymbol{\infty}$ & $\boldsymbol{\infty}$ \\
    5 & 1.250 & 0.104 & 4.649 & 70.800 \\
    10 & 1.111 & 0.015 & 2.811 & 14.828 \\
    \bottomrule
  \end{tabular}
\end{table}

$\alpha$가 증가할수록 평균은 $x_m$에 수렴하고 분산은 급격히 감소하여
분포가 $x_m$ 부근에 집중된다.
반대로 $\alpha$가 감소하면 꼬리가 두꺼워져
극단적으로 큰 값의 출현 확률이 높아진다.
이러한 $\alpha$의 역할은 생존함수(survival function)
$\bar{F}(x) = (x_m/x)^{\alpha}$를 통해 직관적으로 확인된다\chg{. 예컨대 $\alpha = 2$일 때
$\bar{F}(10) = 0.01$인 반면, $\alpha = 5$에서는 $\bar{F}(10) = 10^{-5}$로 극단값의 출현
확률이 1/1{,}000 이하로 감소한다.}

사이버 위험(cyber risk) 분야에서 멱법칙 분포의 적합성은 다수의 실증 연구에 의해
확인되어 왔다. 데이터 유출 규모, 사이버공격에 의한 재무 손실, 시스템 장애 지속 시간 등은
정규분포보다 멱법칙 분포에 의해 더 잘 설명되며,
이는 BCM 맥락에서 극단적 복구 지연 시나리오의 모델링에
파레토 분포가 적합함을 시사한다.
이에 본 연구에서는 멱법칙 분포의 대표적 형태인 파레토 분포를 채택하여
C3 조건의 불확실성을 모형화하며,
이후 본문에서 ``파레토 분포''로 통일 표기한다.

극단값 이론(Extreme Value Theory, EVT)은 이러한 파레토 분포의 극단적 사건을
체계적으로 분석하는 통계적 프레임워크를 제공한다~\citep{coles2001extremes}.
\citet{embrechts1997extremal}은 보험 및 금융 분야에서의 극단값 모델링을 위한
수학적 기초를 수립하였으며, \citet{mcneil2015risk}는 정량적 위험관리에서
EVT의 활용 방안을 구체적으로 논의하였다.
\chg{본 연구는 극단값 이론을 직접적인 추정 기법으로 사용하지는 않으나, 그 핵심 결과는
C3 조건에서 파레토 분포를 채택한 이론적 정당성을 제공한다. 극단값 이론에 따르면 광범위한
분포에서 충분히 높은 임계값을 초과하는 부분이 일반화 파레토 분포로 수렴하므로, 복구 시간의
극단적 꼬리를 파레토 분포로 모형화하는 것은 임의의 선택이 아니라 극단값의 점근적 거동에
부합하는 선택이다. 본 연구가 멱법칙과 극단값 이론에서 가져온 것은 ``극단적 사건의 꼬리는
지수적으로가 아니라 멱법칙으로 감소한다''는 명제이며, 이를 C3 채널의 분포 형태로 직접
반영한 것이다.}

%% --- α = 3 선택 근거 ---
\subsubsection{$\alpha = 3$ 선택 근거}
\label{subsubsec:alpha3-rationale}

본 연구에서는 $\DRTO$ 프레임워크의 조건 \textbf{C3}에서 파레토 분포
\chg{$U^{(P)} \sim 6 \cdot \text{Pareto}(\alpha = 3)$}에 기반한 불확실성 채널을 도입하여,
가우시안 분포(\textbf{C2})가 포착하지 못하는 극단적 복구 지연 시나리오를 모델링한다.
\chg{$\alpha = 3$의 선택은 네 가지 근거에 기반한다. 첫째, $\alpha > 2$이면 평균과 분산이
모두 유한하여 표본 통계량의 안정적 추정과 회귀분석이 가능하다. 둘째, $\alpha \leq 4$이면
첨도가 발산하여 가우시안 분포 대비 극단값이 현저히 빈번하게 출현하는 현실적 시나리오를
표현한다. 셋째, $\alpha = 3$은 유한 분산과 무한 왜도 및 첨도가 공존하는 경계에 위치하여
통계적 분석 가능성과 중꼬리(heavy-tail) 특성 사이의 최적 균형점을 제공한다. 넷째,
사이버 보안 사고의 규모 분포에서 보고되는 $\alpha$ 추정값이 2와 4 사이의 중앙 영역에
해당한다는 실증적 정합성을 갖는다.}


%% --- C2·C3 정량적 비교 및 첨도-꼬리 관계 ---
\subsubsection{가우시안과 파레토 분포의 정량적 비교}
\label{subsubsec:c2c3-comparison}

$\DRTO$ 프레임워크가 채택한 두 불확실성 채널, 즉 가우시안 분포
\chg{$U^{(G)} \sim N(3,\,1)$}(C2)과 파레토 분포 \chg{$U^{(P)} \sim 6 \cdot \text{Pareto}(\alpha=3)$}(C3)은
동일한 기대값($3.0$)을 갖도록 척도를 조정하였으나, 꼬리 영역의 거동에서 질적으로
구별된다. \p{[}표~\ref{tab:c2c3-comparison}\p{]}은 두 분포의 꼬리 성질과 주요 분위수를
비교한 것이다.

\begin{table}[htbp]
\centering
\caption{가우시안(C2)과 파레토(C3) 불확실성 분포의 비교}
\label{tab:c2c3-comparison}
\small
\begin{tabular}{l c c l}
\toprule
\textbf{구분} & \textbf{가우시안 (C2)} & \textbf{파레토 (C3)} & \textbf{의미} \\
\midrule
꼬리 분류 & 경꼬리 (light-tail) & 중꼬리 (heavy-tail) & 감소 속도 기반 성질 분류 \\
꼬리 감소 속도 & $\sim e^{-x^2/2}$ (지수적) & $\sim x^{-\alpha}$ (멱법칙) & 중꼬리가 훨씬 느리게 감소 \\
첨도 & $3$ (기준값) & $\infty$ ($\alpha = 3$) & 꼬리 두께의 수치 지표 \\
분산 & $1.0$ & $27.0$ (약 $27$배) & 변동 폭의 차이 \\
왜도 & $0$ (대칭) & 양수 (우편향) & 극단 지연의 비대칭성 \\
P95 & $4.645$ & $10.29$ (약 $2.2$배) & 상위 $5\%$ 복구 시간 \\
P99 & $5.326$ & $21.85$ (약 $4.1$배) & 상위 $1\%$ 복구 시간 \\
중앙값 & $\approx 3.0$ & $1.56$ & 대부분은 빠르되 소수가 급격히 지연 \\
\bottomrule
\end{tabular}
\end{table}

여기서 첨도와 꼬리 분류는 관련은 있으나 서로 다른 개념임에 유의해야 한다.
첨도는 분포의 꼬리 두께를 하나의 수치로 측정한 지표인 반면, 경꼬리와 중꼬리는
꼬리의 감소 속도에 따른 성질 분류이다. 중꼬리 분포는 첨도가 높은 경향이 있으나,
첨도가 높다고 반드시 중꼬리인 것은 아니다. 예컨대 유한 지지(bounded support)를
갖는 분포도 극단적 이봉 구조를 통해 높은 첨도를 가질 수 있으나 경꼬리에 속한다.
파레토($\alpha=3$)는 첨도가 무한대이므로 중꼬리의 전형적 사례에 해당한다.
파레토 분포의 중앙값($1.56$)이 평균($3.0$)보다 현저히 낮다는 사실은, 대부분의 사건이
빠르게 복구되지만 소수의 극단 사례에서 복구가 급격히 지연되어 평균을 끌어올리는
중꼬리의 구조적 특성을 직접 보여준다. 두 분포의 시각적 비교와 P99 수치는
\figref{fig:pareto_concept}(본문 \secref{sec:result_c3})를 함께 참조하라.

이러한 분포 특성의 차이는 실무 맥락에서도 뚜렷이 구별된다. 가우시안 불확실성은
일상적 운영 변동을 대표하는데, 일반적 라우터 장애의 복구 시간이나 정기적 소프트웨어
업데이트의 롤백, 표준화된 서버 교체 절차처럼 복구 시간이 평균 근처에 집중되고
극단적 편차의 확률이 지수적으로 감소하여 대부분 예측 가능한 범위에서 마무리되는
경우가 여기에 해당한다. 경꼬리 특성상 이러한 사건은 P99에서도 평균 대비
$2.3\sigma$ 이내로 한정된다. 반면 파레토 불확실성은 극단적 사건을 대표하는데,
전사 시스템이 암호화되는 랜섬웨어 공격에서 대부분은 백업 복원으로 수 시간 내에
해결되나 소수의 극단 사례에서는 데이터 복호화와 포렌식 분석으로 복구가 수일까지
연장되며(P99에서 약 $13.2$시간 이상), 지진 후 화재와 정전이 동시에 발생하는 복합
재난이나 핵심 부품 공급사의 공장 화재로 인한 공급망 연쇄 장애 또한 개별 사건
복구 시간의 합을 초과하는 비선형적 지연을 유발하여 꼬리 영역의 복구 부담을 급증시킨다.


%% -------------------------------------------------------------------
\subsection{불확실성 분류 체계의 통합적 이해}
\label{subsec:uncertainty-classification}
%% -------------------------------------------------------------------

본 절에서 다룬 가우시안 분포(C2)와 파레토 분포(C3)는 불확실성을 \textbf{분포 형태}
(즉 \textbf{강도 분포}의 모양) 차원에서 분류한 것이다. 그러나 BCM 운영 환경에서 불확실성 전반을 체계적으로
이해하기 위해서는 \citet{walker2003uncertainty}이 제시한 \textbf{불확실성 단계 분류}
프레임워크와 함께 살펴볼 필요가 있다(\p{[}그림~\ref{fig:uncertainty-levels}\p{]}).
이 프레임워크는 우리가 알고 있는 정도(level of knowledge)에 따라 불확실성을 네 단계로
구분한다. 이는 리스크 평가의 표준 분해(리스크 = 빈도 × 강도)와 대응되며, 분포는 강도
(Severity)에, 시나리오·확률은 빈도(Frequency)에 해당한다. \chg{네 단계는 i) 분포와
확률을 모두 추정할 수 있는 통계적 불확실성, ii) 대안 시나리오는 열거할 수 있으나 확률이
미정인 시나리오 불확실성, iii) 분포 형태 자체를 특정할 수 없는 인지된 무지, iv) 모델과
시나리오와 확률이 모두 불확실한 전적 무지로 구분된다.}

본 연구의 $\DRTO$ 모델은 이 분류 체계 중 양 극단을 정량적으로 다룬다.
통계적 불확실성은 가우시안 분포(C2, \chg{$U^{(G)} \sim N(3,1)$})에 대응하며, 인지된 무지는
파레토 분포(C3, \chg{$U^{(P)} \sim 6 \cdot \text{Pareto}(\alpha\!=\!3)$})로 보수적으로 표현된다.
시나리오 불확실성은 관측성 채널($O_{\text{raw}}$, $k_O$)의 정보 확보를 통해 부분적으로
흡수되며, 전적 무지(블랙 스완 사건)는 본 모델의 적용 한계로 4.3절에서 별도 논의한다.
\chg{한편 운영 단계에서 $U_{\text{raw}}$를 추정하기 위한 데이터 출처(과거 장애 이력 DB,
외부 위협 지수)와 합성 방식은 본문 \subsecref{app:operationalization} 조작화 예시에서 참고용으로 다룬다.}

\begin{figure}[htbp]
\centering
\begin{tikzpicture}[
  every node/.style={text=black},
  level/.style={draw, thick, rounded corners=3pt, text width=125mm,
    minimum height=24mm, font=\large\bfseries, align=left, inner sep=6pt, text=black},
  arr/.style={-{Stealth[length=5mm]}, ultra thick, black}
]
  \node[level, fill=green!15]                   (l1) {Statistical uncertainty (통계적 불확실성)\\
    \normalsize\mdseries 분포(강도)·확률(빈도) 추정 가능. 반복 관측으로 정량화. 예: 장비 고장률.
    \;\textit{$\to$ DRTO C2 가우시안}};
  \node[level, fill=blue!12, below=2mm of l1]   (l2) {Scenario uncertainty (시나리오 불확실성)\\
    \normalsize\mdseries 대안 시나리오 열거 가능, 확률(빈도) 미정. 예: 미경험 재난 유형.
    \;\textit{$\to$ 관측성 채널}};
  \node[level, fill=orange!18, below=2mm of l2] (l3) {Recognized ignorance (인지된 무지)\\
    \normalsize\mdseries 분포(강도) 형태 자체 특정 불가. 예: 신종 사이버공격, 블랙 스완.
    \;\textit{$\to$ DRTO C3 파레토}};
  \node[level, fill=red!12, below=2mm of l3]    (l4) {Total ignorance (전적 무지)\\
    \normalsize\mdseries 모델·시나리오·확률(빈도·강도) 모두 불확실. 예: 기후 복합재난.
    \;\textit{$\to$ 4.3절 한계}};

  \coordinate (atop) at ($(l1.north west) + (-12mm, 0)$);
  \coordinate (abot) at ($(l4.south -| atop)$);
  \draw[arr] (atop) -- (abot);
  \node[font=\normalsize\bfseries, rotate=90, text=black, anchor=south]
    at ($(atop)!0.5!(abot) + (-3mm, 0)$) {불확실성 정도 증가};
\end{tikzpicture}
\caption{불확실성의 단계별 분류 체계와 $\DRTO$ 모델 대응}
\label{fig:uncertainty-levels}
\end{figure}


%% -------------------------------------------------------------------
\subsection{불확실성 하위개념과 업무연속성 리스크와의 관계}
\label{subsec:uncertainty-risk}
%% -------------------------------------------------------------------

\chg{본 연구에서 불확실성은 복구시간의 확률적 변동을 지칭하며, BCM에서
다루는 리스크(재난)와는 분석 수준이 다르다. 리스크는 무엇이 발생하는가라는
사건(event)에 초점을 두는 반면, $\DRTO$에서의 불확실성은 복구시간이 얼마나 변동하는가라는
결과(outcome)에 초점을 둔다. 즉 리스크 사건이 발생하면 복구시간에 불확실성이
유입되고, $\DRTO$가 이를 확률분포로 정량화한다.

경제학자 \citet{knight1921risk}는 리스크와 불확실성을 다음과 같이 구분하였다. 리스크는 결과의
확률분포가 알려져 정량적 측정이 가능한 상황을, 불확실성은 확률분포 자체를 특정할 수
없는 상황을 의미한다. 이 구분은 본 연구의 분류 체계와 대응한다. C2의 가우시안
불확실성은 분포 형태와 매개변수가 사전에 특정되므로 나이트적 의미의 리스크에 가깝고,
C3의 파레토 불확실성은 분포 형태를 확정할 수 없는 나이트적 불확실성을 보수적으로
근사한 것이다. 불확실성을 학술적으로 더 세분한 분류와 본 모델의 대응은
\p{[}표~\ref{tab:app-uncertainty-taxonomy}\p{]}과 같다.}

\begin{table}[htbp]
\centering
\caption{불확실성 하위개념 분류와 $\DRTO$ 모델 대응}
\label{tab:app-uncertainty-taxonomy}
\small
\begin{tabular}{p{1.8cm} p{2.0cm} p{3.8cm} p{2.8cm} p{2.8cm}}
\toprule
\textbf{분류} & \textbf{영어 표기} & \textbf{정의} & \textbf{BCM 예시} & \textbf{DRTO 대응} \\
\midrule
인식론적 불확실성 & Epistemic & 지식 부족에 기인, 정보 확보로 축소 가능 & 미경험 재난 유형, 불완전한 BIA & 관측성 채널 ($O_{\text{raw}}$, $k_O$) \\
\midrule
우연적 불확실성 & Aleatory & 본질적 무작위성, 데이터를 추가해도 축소 불가 & 장비 고장률, 네트워크 지연 변동 & C2 가우시안 ($U^{(G)}$) \\
\midrule
나이트적 불확실성 & Knightian & 확률분포 자체를 특정할 수 없는 불확실성 & 블랙 스완 사건, 신종 사이버공격 & C3 파레토 ($U^{(P)}$) \\
\midrule
심층 불확실성 & Deep Uncertainty & 모델·시나리오·확률 모두 불확실 & 기후변화 복합재난, 팬데믹 & 한계 사항 (\secref{sec:limitations}) \\
\bottomrule
\end{tabular}
\end{table}

\chg{표에서 우연적 불확실성은 반복 관측으로 분포를 추정할 수 있으나 변동 자체를 제거할
수는 없는 유형이고, 인식론적 불확실성은 관측성($O_{\text{raw}}$)을 높임으로써 축소할 수
있어 $\DRTO$ 모델의 관측성 채널이 직접 반영하는 메커니즘에 해당한다. 나이트적 불확실성은
분포 형태를 사전에 특정할 수 없으므로 본 연구에서는 파레토 분포로 보수적으로 모델링하며,
심층 불확실성은 현 모델의 범위를 벗어나 \secref{sec:limitations}에서 한계로 명시하였다.}

\chg{한편 BCM에서 다루는 세 가지 재난 유형과 본 모델의 불확실성은
\p{[}표~\ref{tab:app-disaster-mapping}\p{]}과 같이 대응한다. 동일한 재난 유형이라도
발생 규모에 따라 불확실성 특성이 달라져, 일상적 서버 장애는 가우시안(C2)으로,
대규모 사이버공격은 파레토(C3)로 모형화된다.}

\begin{table}[htbp]
\centering
\caption{BCMS 재난 유형과 $\DRTO$ 불확실성 매핑}
\label{tab:app-disaster-mapping}
\small
\begin{tabular}{p{2cm} p{3.2cm} p{2.8cm} p{2.5cm} p{2.5cm}}
\toprule
\textbf{재난 유형} & \textbf{세부 예시} & \textbf{빈도-영향 특성} & \textbf{불확실성 유형} & \textbf{DRTO 조건} \\
\midrule
자연재난 & 지진, 태풍, 홍수, 산사태 & 빈도 낮음, 영향 극대 & 나이트적 & C3 (파레토) \\
\midrule
사회재난 & 대규모 정전, 통신장애, 공급망 단절, 테러 & 빈도 중간, 영향 중$\sim$대 & 우연적 + 나이트적 & C2 $\sim$C3 \\
\midrule
기술적 재난 & 서버 장애, 사이버공격, 소프트웨어 결함, 데이터 유출 & 빈도 높음, 영향 소$\sim$극대 & 우연적(일상) / 나이트적(극단) & C2(일상) / C3(극단) \\
\bottomrule
\end{tabular}
\end{table}

\chg{사회재난의 경우에도 공급망 단절의 파급 범위에 따라 가우시안(C2)에서 파레토(C3)로
전이될 수 있다.}

\chg{끝으로 BCM에서의 리스크와 $\DRTO$에서의 불확실성의 유사점과 차이점을
\p{[}표~\ref{tab:app-risk-vs-uncertainty}\p{]}에 비교하였다. 리스크가 사건의 발생
가능성과 영향을 다루는 데 비해 불확실성은 그 결과인 복구시간의 확률적 변동을 다루며,
양자는 모두 복구시간에 영향을 미치므로 $\DRTO$가 이를 통합 정량화한다.}

\begin{table}[htbp]
\centering
\caption{BCMS 리스크(재난)와 $\DRTO$ 불확실성 비교}
\label{tab:app-risk-vs-uncertainty}
\small
\begin{tabular}{p{2.5cm} p{5.5cm} p{5.5cm}}
\toprule
\textbf{비교 기준} & \textbf{BCMS 리스크 (재난)} & \textbf{DRTO 불확실성} \\
\midrule
정의 & 조직 목표에 부정적 영향을 미치는 사건의 발생 가능성과 결과~\nrev{(ISO 31000, 2018)} & 복구시간($\DRTO$)의 확률적 변동 범위 \\
\midrule
분석 단위 & 사건(event): ``무엇이 발생하는가'' & 결과(outcome): ``복구시간이 얼마나 변동하는가'' \\
\midrule
방향성 & 주로 부정적 (위협 중심) & 양방향 (단축도 연장도 가능) \\
\midrule
정량화 방법 & 발생가능성 $\times$ 영향도 매트릭스 & 확률분포 (가우시안 / 파레토) \\
\midrule
시간적 성격 & BIA 시점 고정 평가 (정적) & 실시간 동적 변동 \\
\midrule
축소 방법 & 예방·완화 통제 수단 & 관측성 향상 ($O_{\text{raw}}$ $\uparrow$) \\
\midrule
\multicolumn{3}{l}{공통점: 둘 다 복구시간에 영향 $\to$ $\DRTO$가 양자를 통합 정량화} \\
\bottomrule
\end{tabular}
\end{table}

\chg{이상에서 살펴본 가우시안과 파레토 불확실성은 모두 복구 시간에 확률적 변동을 유발하므로, 이를 $\DRTO$ 모델에 반영하려면 그 영향을 일정한 경계 안에서 변환하는 장치가 필요하다. 다음 절에서는 이러한 변환을 담당하는 시그모이드 함수와 바운딩 기법을 검토한다.}
  % 2.4 불확실성(Uncertainty) 모델링


%% ===================================================================
\section{시그모이드 함수와 바운딩}
\label{sec:lit-sigmoid}
%% ===================================================================

%% -------------------------------------------------------------------
\subsection{시그모이드 함수의 수학적 성질}
\label{subsec:sigmoid-properties}
%% -------------------------------------------------------------------

$\DRTO$ 모델에서는 관측성과 불확실성의 복합 효과를 물리적으로 타당한 범위 내로
제한하는 바운딩(bounding) 기법이 필수적이다. 복구 시간은 물리적으로
음수가 될 수 없으며, 조직이 허용하는 최대 복구 시간($R_{\max}$)을 초과할 수도 없기 때문이다.
기준 RTO($R_0$)와 최대 복구 시간 MTPD($R_{\max}$)의 상세 정의는 2.7절에서 제시한다.

\chg{본 연구의 바운딩 함수는 \citet{verhulst1838notice}가 인구 성장 모형으로 처음
도입한 표준 로지스틱 함수 $\sigma(z) = 1/(1+e^{-z})$에 이론적 뿌리를 둔다. 로지스틱
함수는 매끄러운 S자 형태로 출력을 유한 구간 안에 제한하는 성질 덕분에, 생물학적 성장
모형에서 출발하여 통계학과 여러 응용 분야의 표준적 변환으로 정착해 왔다~\citep{cramer2002origins}.}
본 연구에서 채택한 수정 시그모이드(modified sigmoid) 함수는 다음과 같이 정의된다.
\begin{equation}
  S(z) = \frac{2}{1 + e^{-z}} - 1, \quad S: \mathbb{R} \to (-1,\; 1)
  \label{eq:sigmoid-def}
\end{equation}

이 함수는 표준 로지스틱 함수를 $(-1, +1)$ 구간으로 사상한 것이다.
\chg{표준 로지스틱 함수의 치역이 $(0,1)$인 것과 달리 이를 $(-1,+1)$로 사상하여 채택한
이유는, $\DRTO$ 모델에서 관측성은 복구 시간을 기준선 \emph{아래}로, 불확실성은 기준선
\emph{위}로 작용하므로 $S(0)=0$을 중심으로 음과 양의 양방향으로 대칭인 출력이 필요하기
때문이다. 즉 로지스틱 함수의 ``매끄러운 유한 바운딩''이라는 성질은 그대로 가져오되,
두 채널의 상하 대칭 작용을 표현하기 위해 치역만 원점 대칭으로 \chg{평행이동하고 확대한} 것이다.}
%% [G19/일원화] 구 성질 열거(S(0)=0·극한·도함수·홀함수·C∞)는 2.9절 subsec:basic-properties(2계 도함수·변곡점 포함 엄밀 정리)와 중복 → 한 문장으로 축약·전방참조.
\chg{이 함수는 $S(0) = 0$의 중립성, $(-1, +1)$로의 매끄러운 유한 바운딩, 전 구간 단조증가,
$S(-z) = -S(z)$의 원점 대칭(홀함수), $C^{\infty}$ 미분 가능성 등의 핵심 성질을 가지며, 이들
성질의 엄밀한 정리와 2계 도함수와 변곡점 분석은 모델의 수학적 성질을 다루는 2.9절에서 종합한다.}


%% -------------------------------------------------------------------
\subsection{\chg{선형과 클리핑, 시그모이드 바운딩 기법의 비교}}
\label{subsec:bounding-comparison}
%% -------------------------------------------------------------------

\chg{$\DRTO$ 모델의 바운딩 함수는 다음 다섯 가지 설계 요구조건을 충족해야 한다. 첫째,
출력이 언제나 물리적으로 타당한 범위 $0 < \DRTO < R_{\max}$ 안에 머물러야 한다(경계 보장).
둘째, 입력이 작은 평균 영역에서는 근사적으로 선형으로 반응하여 해석이 직관적이어야 한다
(평균 영역 선형성). 셋째, 입력이 극단적으로 커지는 영역에서는 출력이 포화하여 경계를
넘지 않아야 한다(극단 영역 포화). 넷째, 관측성과 불확실성 두 채널의 기여가 각각의 물리적
범위 안에서 분리되어 작용해야 한다(채널 분리). 다섯째, 입력 변화에 대해 출력이 단조적으로
반응하여 입력의 순서가 보존되어야 한다(단조성). 이 다섯 요구조건을 기준으로 세 가지 후보
바운딩 기법을 비교한다.}

\textbf{선형 모델}은 $\DRTO = R_0 + \alpha \cdot x$ 형태로서, 구현이 가장 단순하고
해석이 용이하다. 그러나 입력값이 극단적인 경우 $\DRTO$가 음수가 되거나
$R_{\max}$를 초과하는 물리적 비현실성이 발생할 수 있어 \chg{첫 번째 요구조건인 경계
보장을 충족하지 못한다}.

\textbf{클리핑(hard clipping)}은 $f(z) = \min(\max(z, z_{\min}), z_{\max})$로서
입력값이 경계 $[z_{\min}, z_{\max}]$ 내이면 그대로 통과시키고 경계를 초과하면 강제로
절단하는 경성 바운딩 함수이다. 클리핑은 경계 조건을 강제로 보장하나, 경계 지점에서
도함수가 불연속이 되고 \chg{경계 외의 다수 입력값이 단일 경계값으로 매핑되면서} 경계
근방의 미세한 입력 변화가 출력에 전혀 반영되지 않는 정보 손실이 발생한다.
\chg{즉 클리핑은 경계 보장 요구조건은 충족하지만, 경계 지점에서 매끄럽지 않아 평균 영역의
선형성과 극단 영역의 점진적 포화를 함께 만족하지 못하며, 경계 외에서 입력 변화가 출력에
반영되지 않아 단조성도 부분적으로 훼손된다.}

\textbf{시그모이드 바운딩}은 매끄러운 S자 곡선을 통해 출력을 유한 범위 내로
자연스럽게 제한하며, 전 구간에서 미분 가능하여 입력 변수의 변화에 대한
출력의 민감도를 연속적으로 분석할 수 있다.
\chg{시그모이드 바운딩은 매끄러운 유한 제한으로 경계를 보장하고, 원점 부근에서 근사적으로
선형이며, 입력이 커질수록 출력이 점진적으로 포화하고, 두 채널을 각각의 물리적 범위에서
독립적으로 바운딩할 수 있으며, 전 구간에서 단조증가한다. 따라서 앞서 제시한 다섯 가지
설계 요구조건을 모두 충족하는 것은 세 후보 중 시그모이드 바운딩뿐이며, 본 연구가 수정
시그모이드를 채택한 근거가 여기에 있다.}

\p{[}표~\ref{tab:bounding-comparison}\p{]}에 세 기법의 특성을 종합적으로 비교하였다.

\begin{table}[htbp]
\centering
\caption{바운딩 기법의 비교}
\label{tab:bounding-comparison}
\small
\begin{tabularx}{\textwidth}{@{}>{\raggedright\arraybackslash}p{2.4cm}
  >{\raggedright\arraybackslash}X
  >{\raggedright\arraybackslash}X
  >{\raggedright\arraybackslash}X@{}}
\toprule
\textbf{특성} & \textbf{선형 모델} & \textbf{클리핑(경성)} & \textbf{시그모이드 바운딩} \\
\midrule
수식 &
$R_0 + \alpha x$ &
$\min(\max(\cdot, 0), R_{\max})$ &
$R_0 + E_{O,\text{bnd}} + E_{U,\text{bnd}}$ \\
\addlinespace
경계 보장 &
미보장 (음수/초과 가능) &
보장 (강제 절단) &
보장 (점근적 접근) \\
\addlinespace
미분 가능성 &
전 구간 ($C^{\infty}$) &
경계점 불연속 &
전 구간 ($C^{\infty}$) \\
\addlinespace
경계 정보 보존 &
해당 없음 &
경계에서 기울기 $= 0$ (정보 손실) &
경계에서도 기울기 $> 0$ (정보 보존) \\
\addlinespace
포화 효과 &
없음 (무한 선형 증가) &
경성 포화 (즉시 절단) &
연성 포화 (점진적 둔화) \\
\addlinespace
물리적 해석 &
비현실적 극단값 가능 &
현실적이나 불연속 &
현실적이고 연속적 \\
\bottomrule
\end{tabularx}
\end{table}

\figref{fig:bounding-comparison}\refpo{fig:bounding-comparison}{은}{는} 세 바운딩 기법의 출력 특성을 비교한 것이다.
\chg{선형 모델은 경계를 초과할 수 있고 클리핑은 경계점에서 불연속인 반면, 시그모이드는 $(-1, +1)$ 범위 안에서 경계에 매끄럽게 점근하면서도 전 구간에서 기울기가 양으로 유지되는 점근적 바운딩을 제공한다.}
시그모이드 바운딩은 경계 영역에서의 정보 보존성, 미분 가능성,
운영 해석 가능성 측면에서 $\DRTO$ 프레임워크에 가장 적합하다.

\begin{figure}[htbp]
  \centering
  \begin{tikzpicture}
    \begin{axis}[
      width=0.88\textwidth,
      height=0.52\textwidth,
      xlabel={입력 $z$},
      ylabel={출력},
      xmin=-4.5, xmax=4.5,
      ymin=-1.6, ymax=1.9,
      grid=major,
      grid style={black},
      legend pos=south east,
      legend style={font=\small, rounded corners=2pt, fill=white, fill opacity=0.9},
      axis lines=middle,
      every axis x label/.style={at={(axis description cs:1.02,0.5)}, anchor=west, font=\small},
      every axis y label/.style={at={(axis description cs:0.5,1.02)}, anchor=south, font=\small},
      clip=false
    ]
      % 선형 모델 (제한 없음)
      \addplot[thick, blue!70, domain=-4.5:4.5, samples=100]
        {0.4*x};
      \addlegendentry{선형 ($0.4z$)}

      % 클리핑
      \addplot[thick, red!70, domain=-4.5:4.5, samples=200]
        {min(max(0.4*x, -1), 1)};
      \addlegendentry{클리핑 (hard)}

      % 시그모이드 S(z) = 2/(1+exp(-z)) - 1
      \addplot[thick, black, domain=-4.5:4.5, samples=200]
        {2/(1+exp(-x)) - 1};
      \addlegendentry{시그모이드 $S(z)$}

      % 경계 초과 영역 채우기 (위쪽: x>2.5에서 0.4x > 1)
      \fill[blue!12, opacity=0.5]
        (axis cs:2.5,1) -- (axis cs:4.5,1.8) -- (axis cs:4.5,1) -- cycle;

      % 경계 초과 영역 채우기 (아래쪽: x<-2.5에서 0.4x < -1)
      \fill[blue!12, opacity=0.5]
        (axis cs:-2.5,-1) -- (axis cs:-4.5,-1.8) -- (axis cs:-4.5,-1) -- cycle;

      % 경계선 표시
      \addplot[thin, dashed, black] coordinates {(-4.5, 1) (4.5, 1)};
      \addplot[thin, dashed, black] coordinates {(-4.5, -1) (4.5, -1)};

      % 경계 라벨
      \node[font=\scriptsize, text=black, anchor=west] at (axis cs:4.6, 1) {$+1$};
      \node[font=\scriptsize, text=black, anchor=west] at (axis cs:4.6, -1) {$-1$};

      % 선형 경계 초과 영역 라벨 (위쪽)
      \node[font=\scriptsize, text=black, fill=white, inner sep=1.5pt,
        anchor=south west] at (axis cs:2.8, 1.45)
        {선형: 경계 초과};
      \draw[-{Stealth[length=1.5mm]}, blue!70, thin]
        (axis cs:3.2, 1.42) -- (axis cs:3.5, 1.15);

      % 선형 경계 초과 영역 라벨 (아래쪽)
      \node[font=\scriptsize, text=black, fill=white, inner sep=1.5pt,
        anchor=north east] at (axis cs:-2.8, -1.45)
        {선형: 경계 초과};
      \draw[-{Stealth[length=1.5mm]}, blue!70, thin]
        (axis cs:-3.2, -1.42) -- (axis cs:-3.5, -1.15);

    \end{axis}
  \end{tikzpicture}
  \caption{바운딩 기법 비교}
  \label{fig:bounding-comparison}
\end{figure}


%% -------------------------------------------------------------------
\subsection{개별 시그모이드 바운딩의 적용}
\label{subsec:individual-sigmoid}
%% -------------------------------------------------------------------

본 연구에서 채택한 개별 시그모이드 바운딩(individual sigmoid bounding) 모델은
관측성과 불확실성의 효과를 각각 고유한 물리적 범위(span)에서
독립적으로 바운딩하는 구조이다. 대안인 결합(joint) 시그모이드 $S(z_O + z_U)$가
두 효과를 단일 함수로 압축하는 반면, 개별 시그모이드는 각 효과의 물리적 범위를
독립적으로 보장한다.
%% [G19/일원화] 구 핵심수식(eq:drto-individual-sigmoid, eq:lit-E_O_bnd, eq:lit-E_U_bnd, eq:lit-z_O, eq:lit-z_U)은
%% 2.7절 모델 구조(eq:drto, eq:E_O_bnd, eq:E_U_bnd, eq:z_O, eq:z_U; canonical)와 완전 중복 → 수식 삭제·전방참조(라벨 참조 0건 확인).
\chg{구체적으로 최종 $\DRTO$는 기준선 $R_0$에 관측성 바운딩 조정량
$E_{O,\text{bnd}} = -R_0 \cdot S(z_O)$(물리적 범위 $R_0$)와 불확실성 바운딩 조정량
$E_{U,\text{bnd}} = (R_{\max} - R_0) \cdot S(z_U)$(물리적 범위 $R_{\max} - R_0$)를 가산하여
산출한다. 음수 부호 $-R_0$는 관측성이 높을수록 복구 시간 단축 효과로 작용함을, $z_U$의
분모 $R_{\max} - R_0$는 불확실성 효과를 가용 상향 여유 범위로 정규화함을 뜻한다. 시그모이드
입력 인자 $z_O$, $z_U$의 전체 정식과 변수 및 매개변수의 형식적 정의, 경계 보장의 엄밀한
증명, 곡률 $\kappa = 1.2$의 선정 근거는 모델 구조와 최적화를 다루는 2.7절과 2.10절에서
종합하므로, 본 절은 그 형태만 바운딩 기법의 귀결로 제시한다.}

\chg{이 구조는 세 가지 핵심적 장점을 갖는다. 첫째, 경계 보장의 보편성으로서, 관측성 효과의
바운딩 범위($R_0$)와 불확실성 효과의 바운딩 범위($R_{\max} - R_0$)를 각각의 물리적 의미에
맞게 독립적으로 설정하므로 $\DRTO \in [0,\; R_{\max}]$가 $\kappa$ 값에 무관하게 수학적으로
보장된다. 둘째, 경사 최적화와 민감도 분석의 용이성으로서, 각 시그모이드 함수가 전체
정의역에서 $C^{\infty}$이므로 경사 기반 최적화와 민감도 분석이 용이하다. 셋째, C1 조건
분리 가능성으로서, C1 조건($E_{U,\text{bnd}} = 0$)에서 관측성만의 순수 효과를 불확실성
상한과 독립적으로 분석할 수 있어 모델의 단계적 확장이 자연스럽다.}


%% -------------------------------------------------------------------
\subsection{이중채널 감쇠 메커니즘}
\label{subsec:dual-channel-attenuation}
%% -------------------------------------------------------------------

개별 시그모이드 바운딩 구조는 관측성 채널($E_{O,\text{bnd}} = -R_0 \cdot S(z_O)$,
복구 시간 단축)과 불확실성 채널($E_{U,\text{bnd}} = (R_{\max}-R_0) \cdot S(z_U)$,
복구 시간 연장)이 동시에 작용하는 이중채널 구조를 형성한다. 두 채널이 각자의 물리적
범위에서 독립적으로 바운딩되므로, 불확실성이 극단적으로 커지는 꼬리(tail) 영역에서는
불확실성 채널이 지배적으로 작용하면서 관측성 채널의 단축 효과가 구조적으로 약화되는
현상이 나타난다. 이를 이중채널 감쇠(dual-channel attenuation)라 한다.

이 메커니즘은 실무적으로 중요한 함의를 갖는다. 고도의 모니터링 체계를 갖춘
조직(높은 관측성)이라 하더라도 대규모 사이버 공격과 같은 극단적 불확실성 상황에서는
관측성 투자만으로 복구 시간을 충분히 단축하기 어려우며, 불확실성 자체에 대한 관리,
이를테면 업무연속성계획(BCP) 시나리오의 다양화나 사이버 보험 등이 병행되어야 함을
시사한다. 이중채널 감쇠가 분포 수준에서 발현되는 전환점과 그 정량적 크기(꼬리
영역에서의 관측성 효과 감쇠 비율, 조건부 위험가치 격차 등)는 본 연구의 실증 결과에서
분할 회귀를 통해 확인되며, 이는 \secref{sec:ch5_dual_channel}에서 상세히 분석한다.
  % 2.5 시그모이드 함수와 바운딩


%% ===================================================================
\section{선행연구 검토 및 연구의 차별점}
\label{sec:lit-gap}
%% ===================================================================

앞에서 검토한 BCM, 정적 RTO의 한계, 관측성, 불확실성, 시그모이드 바운딩의 이론적 논의를 종합하고 선행연구의 흐름을 정리하여 본 연구가 수행하고자 하는 차별점을 [표 2-8]에 주제별로 정리하였다.

%% -------------------------------------------------------------------
%% -------------------------------------------------------------------

\begin{table}[htbp]
\centering
\caption{주제별 선행연구 요약}
\label{tab:prior-research}
\small
\begin{tabularx}{\textwidth}{@{}>{\raggedright\arraybackslash}p{1.8cm}
  >{\raggedright\arraybackslash}p{3.0cm}
  >{\raggedright\arraybackslash}X
  >{\raggedright\arraybackslash}p{4.0cm}@{}}
\toprule
\textbf{주제} & \textbf{주요 문헌} & \textbf{핵심 내용} & \textbf{한계} \\
\midrule
BCM/BCMS 표준 &
\citet{herbane2010disaster}, ISO 22301(2019) &
PDCA 기반 BCMS 프레임워크, BCM 발전 3단계 분류 &
RTO를 정적 고정값으로 설정, 동적 조정 메커니즘 부재 \\
\addlinespace
BIA 및 RTO &
\citet{snedaker2013bcdr}, \citet{hiles2010bcm} &
BIA 방법론 체계화, MTPD-RTO 관계 정립 &
RTO 동적 특성 미논의, 관측성 미반영 \\
\addlinespace
\m{관측성(제어 이론)} &
\citet{kalman1960control}, \citet{avizienis2004basic} &
LTI 시스템 관측성 형식화, 의존성 분류 체계 &
BCM 영역 적용 부재, RTO와의 수리적 연결 미시도 \\
\addlinespace
\m{관측성(IT 운영)} &
\citet{beyer2016sre}, \citet{majors2022observability} &
SRE 관측성 실무, 세 가지 핵심 요소(three pillars) 체계 &
정량적 모델화 미수행, MTTD/MTTR 관계만 정성적 \\
\addlinespace
가우시안 불확실성 &
\citet{jaynes2003probability}, \citet{gelman2013bayesian} &
CLT 기반 불확실성 정량화, 베이지안 추론 체계 &
극단값 과소평가, 파레토 현상 미포착 \\
\addlinespace
파레토 분포/EVT &
\citet{taleb2007blackswan}, \citet{embrechts1997extremal}, \citet{newman2005powerlaws} &
블랙 스완 개념, EVT 수학적 기초, 멱법칙 보편성 &
금융/보험 중심 적용, BCM/RTO 영역 미적용 \\
\addlinespace
국내 BCM &
\citet{cheung2019quickhit}, \citet{jang2021bcms}, \citet{jung2023bcms} &
BIA Quick-Hit 프레임워크, BCMS 성과 실증, BCM 평가지표 개발 &
정적 RTO 전제, 동적 조정 모델 미제시 \\
\addlinespace
DRTO, 시그모이드 바운딩 &
\citet{lee2026drto} &
관측성 기반 DRTO에 시그모이드 바운딩 도입, 곡률 매개변수 $\kappa^*$ 도출 &
관측성 단일 채널 한정, 불확실성 채널 미포함 \\
\bottomrule
\end{tabularx}
\end{table}

[표 2-8]에서 확인되는 선행연구의 공통적 한계는 \chg{다음 세 가지로 요약되며 2.7(DRTO 모델 구조)에서 상세히 다룬다. 첫째는 관측성을 통합한 RTO 모델의 부재로서 부분적인 단일 채널
모델만 존재하고~\citep{lee2026drto}, 둘째는 불확실성이 정량화된 동적 RTO 모델의 부재이며,
셋째는 극단 시나리오에 대한 파레토 분석의 부재이다.} 본 연구의 $\DRTO$ 프레임워크는 이
\chg{세 가지 한계}를 통합적으로 해소하여 선행연구와 차별화하는 것을 목표로 한다.


%% -------------------------------------------------------------------
%% -------------------------------------------------------------------

\p{[}표~\ref{tab:prior-research}\p{]}의 선행연구를 종합하면, 기존 BCM/RTO 연구와
본 연구의 $\DRTO$ 프레임워크 사이에는 다음 세 가지 핵심적 차별점이 존재한다.

\chg{첫째, 관측성을 통합한 RTO 모델의 부재이다.} 기존 BCM 연구에서 RTO는 BIA를 통해 고정적으로 결정되며, 시스템의 관측성 수준이 복구 시간에 미치는 영향은 정량적으로 모델링되지 않았다. 제어 이론\citep{kalman1960control}과 SRE\citep{beyer2016sre, majors2022observability}에서 관측성의 중요성이 \m{인정되며}, \citet{lee2026drto}은 관측성 단일 채널을 시그모이드 바운딩으로 RTO에 연결하는 수리적 프레임워크를 처음 제시하였으나, 불확실성 채널과 결합한 다채널 통합 모델은 여전히 부재하였다.

\chg{둘째, 불확실성이 정량화된 동적 RTO($\DRTO$) 모델의 부재이다.} 기존 RTO는 결정론적(deterministic) 단일값으로 설정되어 복구 과정에서 발생하는 불확실성을 반영하지 못하였다. 확률론적 사고방식\citep{gelman2013bayesian, jaynes2003probability}이 위험관리 분야에 일부 \m{적용되나}, RTO 자체를 확률 분포로 모델링하여 불확실성에 따른 동적 조정을 수행하는 연구는 부재하였다.

\chg{셋째, 극단적 시나리오에 대한 파레토 분석의 부재이다.} \m{극단값이론(EVT)}과 파레토 분포는 금융 위험관리\citep{embrechts1997extremal, mcneil2015risk}에서 활발히 \m{활용되나}, BCM 영역에서 극단적 복구 지연 시나리오를 파레토 분포로 모델링한 연구는 찾아보기 어렵다. \citet{taleb2007blackswan}이 지적한 극단 사건의 과소평가 문제가 RTO 설정에서도 동일하게 \m{발생함에도} 불구하고, 이에 대한 체계적 분석 및 대응 방안은 선행 연구에서 다루어지지 않았다.

[표 2-9]는 이상의 논의를 종합하여 기존 연구의 접근 방식과 본 연구의 DRTO 프레임워크가 제시하는 해결 방안을 대비하여 비교한 것이다.

\begin{table}[htbp]
\centering
\caption{선행 연구의 한계와 본 연구의 접근 비교}
\label{tab:research-gap}
\small
\begin{tabularx}{\textwidth}{@{}>{\raggedright\arraybackslash}p{2.6cm}
  >{\raggedright\arraybackslash}X
  >{\raggedright\arraybackslash}X
  >{\raggedright\arraybackslash}p{2.3cm}@{}}
\toprule
\textbf{연구 차별점} & \textbf{기존 연구} & \textbf{본 \m{연구($\DRTO$)}} & \textbf{관련 조건} \\
\midrule
차별점 1:\newline 관측성 미반영 &
단일 채널 관측성 모델~\citep{lee2026drto}만 존재;\newline 다채널 관측성-불확실성\newline 통합 부재 &
$O_{\text{raw}} \in [0,1]$ 관측성 원시값을\newline 개별 시그모이드 바운딩의\newline 핵심 입력으로 통합 &
\p{C1 (관측성)} \\
\addlinespace
차별점 2:\newline 불확실성 미정량 &
결정론적 단일값 RTO;\newline 확률적 변동 미고려 &
가우시안 $N(\mu,\sigma^2)$ 및\newline 파레토 $\text{Pareto}(\alpha)$\newline 이중채널 모델링 &
\p{C2 (가우시안),\newline C3 (파레토)} \\
\addlinespace
차별점 3:\newline 극단 시나리오\newline 미분석 &
BCM 영역에서 EVT/\newline 파레토 미적용;\newline 꼬리 위험 미포착 &
EVT 기반 파레토 분석;\newline 분할 회귀(Core/Tail);\newline 이중채널 비선형 효과 정량화 &
\p{C3 tail (파레토\newline P85.8 이상)} \\
\addlinespace
바운딩 기법 &
클리핑(hard boundary)\newline 또는 미적용 &
개별 시그모이드 바운딩\newline $S(z) = 2/(1+e^{-z})-1$\newline $\kappa = 1.2$, 전 구간 미분 가능 &
전 조건 공통 \\
\bottomrule
\end{tabularx}
\end{table}

\chg{[표 2-9]의 관련 조건 열에 나타나는 조건 표기는 관측성과 불확실성의 활성화 조합에 따라 정의한 네 가지 실험 조건을 가리킨다. C0는 관측성과 불확실성을 모두 비활성화한 기준선이고, C1은 관측성만 활성화한 조건이며, C2는 C1에 가우시안 불확실성을 더한 조건, C3는 C1에 파레토 불확실성을 더한 조건이다.}

본 연구의 $\DRTO$ 프레임워크는 위의 세 가지 차별점을 통합적으로 구현하는 것을 목표로 한다.
구체적으로, 관측성을 핵심 입력 변수로 도입하고,
가우시안 및 파레토 분포를 통해 불확실성을 이중 채널로 정량화하며,
파레토 분포를 BCM 영역에 도입하여 극단 시나리오 분석의 공백을 채운다.
다음으로 이러한 이론적 기반 위에 구축되는 DRTO 프레임워크의 구체적 수리 모델과 가설 체계를 상세히 기술한다.
  % 2.6 선행연구 검토 및 연구의 차별점
          % 2.1--2.6 이론적 배경
%% ===================================================================
\section{DRTO 모델 구조}
\label{sec:drto-model}
%% ===================================================================

앞서 도출한 이론적 배경과 선행 연구의 한계를 토대로, 본 절부터는 동적
복구목표시간($\DRTO$) 프레임워크의 수학적 구조와 설계 원리를 체계적으로 기술한다.
이후 계속하여 구체적으로 (i) 시그모이드 정식화(2.7), (ii) 4단계 조건 체계 C0--C3(2.8),
(iii) 수학적 성질(2.9), (iv) 곡률 매개변수 $\kappa^{*}=1.2$ 선정 근거(2.10),
(v) \jrev{10개 연구가설} H1--\jrev{H10}(2.11)을 차례로 다룬다.

%% -------------------------------------------------------------------
\subsection{\chg{개별 시그모이드 바운딩 핵심 수식}}
\label{subsec:core-equation}
%% -------------------------------------------------------------------

$\DRTO$ 모델의 핵심 과제는 관측성과 불확실성이라는 두 입력 차원의 효과를
물리적으로 타당한 범위 내에서 결합하는 것이다. 복구목표시간은 음수가 될 수 없으며,
조직이 허용할 수 있는 최대허용중단기간(MTPD)을 초과해서도 안 된다.
이러한 물리적 제약 조건을 만족하면서도 입력 변수의 변화에 연속적이고
미분 가능한 응답을 제공하는 모델을 구축하는 것이 설계의 핵심 요구사항이다.

본 연구에서는 \textbf{개별 시그모이드 바운딩(individual sigmoid bounding)} 모델을 제안한다.
선행연구\citep{lee2026drto}에서 수립한 관측성 기반 모델(C0$\to$C1)과
최적 곡률 $\kappa^{*} = 1.2$를 기반으로,
본 연구에서는 불확실성 채널(C2, C3)을 추가하여 다채널 구조로 확장한다.
관측성과 불확실성의 효과를 각각 고유한 물리적 범위(span)에서
독립적으로 바운딩하여, 물리적 경계 보장, 연속 미분 가능성,
이론적 유도 가능성이라는 세 요구사항을 동시에 충족한다.

\chg{모델의 기저 함수는} 수정 시그모이드(modified sigmoid)로서,
표준 로지스틱 함수를 $(-1, 1)$ 구간으로 사상한다.

\begin{equation}
  S(z) = \frac{2}{1 + e^{-z}} - 1, \quad S: \mathbb{R} \to (-1,\; 1)
  \label{eq:sigmoid}
\end{equation}

\chg{관측성의 효과는} 시그모이드 함수를 통해 $(-R_0, 0]$ 범위로 바운딩된다\chg{. 이때 $0$은
달성 가능한 값이지만 $-R_0$에는 한없이 접근할 뿐 도달하지는 못한다}.
관측성은 $\DRTO$를 $R_0$에서 최대 $0$까지 감소시킬 수 있으므로,
물리적 하향 범위(span)는 $R_0$이다.
\chg{이와 대칭적으로 불확실성의 효과는} 시그모이드 함수를 통해 $[0, R_{\max} - R_0)$ 범위로 바운딩된다\chg{. $0$은
달성 가능하나 $R_{\max} - R_0$에는 도달하지 못하며, $R_0 = 6.0$h와 $R_{\max} = 24.0$h를
가정하면 그 범위는 $[0, 18.0)$h이다}.
불확실성은 $\DRTO$를 $R_0$에서 최대 $R_{\max}$(MTPD)까지 증가시킬 수 있으므로,
물리적 상향 범위(span)는 $R_{\max} - R_0$이다.

\medskip
두 채널의 바운딩 수식은 다음과 같이 대칭적 구조를 갖는다.
\begin{align}
  E_{O,\text{bnd}} &= -R_0 \cdot S(z_O),
    &\quad \text{span} &= R_0
    \label{eq:E_O_bnd} \\[4pt]
  E_{U,\text{bnd}} &= \;(R_{\max} - R_0) \cdot S(z_U),
    &\quad \text{span} &= R_{\max} - R_0
    \label{eq:E_U_bnd}
\end{align}
시그모이드 입력 인자 $z_O$, $z_U$는 각각
\begin{align}
  z_O &= \kappa \cdot k_O \cdot O_{\text{raw}}
    \label{eq:z_O} \\[4pt]
  z_U &= \frac{\kappa \cdot k_U \cdot R_0 \cdot U_{\text{raw}}}{R_{\max} - R_0},
    \quad k_U = 1 - k_O
    \label{eq:z_U}
\end{align}
이다. $z_O = \kappa \cdot k_O \cdot O_{\text{raw}} \geq 0$은
입력 변수의 정의역에 의해 항상 성립하며,
$U_{\text{raw}} \geq 0$일 때 $z_U \geq 0$이므로
$E_{O,\text{bnd}} \in (-R_0, 0]$,\;
$E_{U,\text{bnd}} \in [0, R_{\max} - R_0)$이다.
$N(3,1)$에서의 음수 $U_{\text{raw}}$ 처리에 대해서는
\ref{subsec:c2}항에서 후술한다.

\chg{두 입력 인자 $z_O$, $z_U$는 세 가지 원칙에 따라 설계된다. 첫째는 무차원
정규화이다. 관측성 점수와 불확실성 원시값은 단위와 척도가 서로 다르므로, 각각
가중치를 곱한 뒤 불확실성은 물리적 상향 범위 $R_{\max} - R_0$로 나누어 두 입력을
모두 무차원량으로 변환한다. 이로써 시그모이드가 비교 가능한 척도 위에서
작동한다. 둘째는 물리적 범위 분리이다. 관측성은 복구 시간을 기준선 아래로
($\mathrm{span} = R_0$), 불확실성은 기준선 위로($\mathrm{span} = R_{\max} - R_0$)
작용하도록 두 채널의 방향과 범위를 분리한다. 셋째는 상보적 가중이다. 관측성
가중치와 불확실성 가중치가 $k_O + k_U = 1$을 이루어, 한 채널의 비중이 커지면
다른 채널의 비중이 줄어드는 상보 관계를 갖는다.}

두 입력 인자는 표면적 형태는 다르나 동일한 구조 $z = \kappa \cdot k \cdot (R_0/\mathrm{span})
\cdot X_{\text{raw}}$를 공유한다. 관측성 채널은 $\mathrm{span} = R_0$이므로 정규화 인자가
$R_0/R_0 = 1$이 되어 수식에서 생략되고, 불확실성 채널은 $\mathrm{span} = R_{\max} - R_0$이므로
정규화 인자가 $R_0/(R_{\max}-R_0) = 1/3$로 나타난다. 두 원시 입력은 모두 무차원이며,
시간 단위(h)는 바운딩 수식의 span에 의해 부여된다.
\p{[}표~\ref{tab:zo-zu-symmetry}\p{]}은 두 채널의 항목별 대칭 구조를 정리한 것이다.

\begin{table}[htbp]
\centering
\caption{관측성 채널($z_O$)과 불확실성 채널($z_U$)의 대칭 구조 비교}
\label{tab:zo-zu-symmetry}
\small
\begin{tabular}{l c c}
\toprule
\textbf{구분} & \textbf{관측성 ($z_O$)} & \textbf{불확실성 ($z_U$)} \\
\midrule
원시 입력     & $O_{\text{raw}} \in [0,1]$ (무차원) & $U_{\text{raw}} \geq 0$ (무차원) \\
정규화 인자   & $R_0 / R_0 = 1$              & $R_0 / (R_{\max} - R_0) = 1/3$ \\
인자 수식     & $\kappa \cdot k_O \cdot (R_0/R_0) \cdot O_{\text{raw}}$
              & $\kappa \cdot k_U \cdot (R_0/(R_{\max}\!-\!R_0)) \cdot U_{\text{raw}}$ \\
가중치        & $k_O$                               & $k_U = 1 - k_O$ \\
span (바깥)   & $R_0$                               & $R_{\max} - R_0$ \\
효과 방향     & 하향 ($\DRTO$ 감소)                 & 상향 ($\DRTO$ 증가) \\
\bottomrule
\end{tabular}
\end{table}

\chg{최종 $\DRTO$는} 기준선 $R_0$에 개별 바운딩된 두 효과를 가산하여 산출한다.

\begin{equation}
  \DRTO = R_0 + E_{O,\text{bnd}} + E_{U,\text{bnd}}
  \label{eq:drto}
\end{equation}

두 바운딩 효과의 범위가 물리적으로 분리되어 있으므로,
$\DRTO$는 항상 다음의 경계 내에 위치한다.

\begin{equation}
  E_{O,\text{bnd}} \in (-R_0, 0], \quad
  E_{U,\text{bnd}} \in [0, R_{\max} - R_0)
  \quad \therefore\;\; \DRTO \in (0,\; R_{\max})
  \label{eq:bounds}
\end{equation}

이 경계는 $\kappa$ 값에 무관하게 항상 성립하며,
$R_{\max} = 4R_0 = 24.0$h로 설정할 경우 $\DRTO \in (0,\; 24.0)$h가 보장된다.
특히 C1 조건($E_{U,\text{bnd}} = 0$)에서는 $\DRTO \in (0, R_0]$이 되어
$R_{\max}$가 불필요하다. 이는 관측성만의 효과를 불확실성 상한과
독립적으로 분석할 수 있게 하는 개별 바운딩의 핵심 장점이다.


\chg{조건 간 비교를 위해} 무차원 정규화 비율 $\eta$를 정의한다.

\begin{equation}
  \eta = \frac{\DRTO}{R_0}
  \label{eq:eta}
\end{equation}

$\eta = 1$은 정적 기준선과 동일함을, $\eta < 1$은 관측성에 의한 복구 시간 단축을,
$\eta > 1$은 불확실성에 의한 복구 시간 연장을 의미한다.
$\DRTO \in (0, R_{\max})$이므로 $\eta$의 이론적 범위는
$(0,\; R_{\max}/R_0)$이며, $R_{\max} = 4R_0$ 가정 하 $(0,\; 4)$이다.


%% -------------------------------------------------------------------
\subsection{모델 구조 도식}
\label{subsec:model-diagram}
%% -------------------------------------------------------------------

\p{[}그림~\ref{fig:model-structure}\p{]}는 $\DRTO$ 개별 시그모이드 바운딩 모델의
전체 구조를 도식화한 것이다. 두 입력 차원(관측성, 불확실성)이
각각 독립적인 시그모이드 바운딩 경로를 통과한 후 기준선 $R_0$에
가산되어 최종 $\DRTO$를 산출하는 구조를 보여준다.
\chg{그림에서 관측성 경로는 상단에 파란색으로, 불확실성 경로는 하단에 주황색으로 표시되며, 두 경로는 각각 독립적인 span에서 시그모이드 바운딩을 거쳐 가산 노드에서 기준선 $R_0$와 합산된다. 곡률 매개변수 $\kappa = 1.2$는 두 경로에 공통으로 적용된다.}

\begin{figure}[htbp]
\centering
\resizebox{0.95\textwidth}{!}{%
\begin{tikzpicture}[
  inputbox/.style={draw, rounded corners=5pt, minimum width=26mm,
    minimum height=13mm, align=center, font=\small, line width=0.6pt},
  procbox/.style={draw, rounded corners=4pt, minimum width=36mm,
    minimum height=12mm, align=center, font=\small, line width=0.6pt},
  outbox/.style={draw, rounded corners=5pt, minimum width=40mm,
    minimum height=16mm, align=center, font=\small\bfseries,
    line width=0.8pt},
  sumbox/.style={draw, circle, minimum size=10mm, align=center,
    font=\normalsize\bfseries, line width=0.8pt},
  arr/.style={-{Stealth[length=2.5mm]}, thick},
  lbl/.style={font=\scriptsize, fill=white, inner sep=1.5pt}
]

% ── 관측성 경로 (상단, y=3.0) ──
\node[inputbox, fill=blue!10] (ko) at (0, 4.0) {\textbf{가중치} $k_O$\\$\in [0,1]$};
\node[inputbox, fill=blue!10] (oraw) at (0, 2.0) {\textbf{원시 관측성} $O_{\text{raw}}$\\$\in [0,1]$};
\node[procbox, fill=blue!5] (zo) at (5.0, 3.0) {$z_O = \kappa \cdot k_O \cdot O_{\text{raw}}$};
\node[procbox, fill=blue!12] (sobnd) at (10.5, 3.0)
  {$E_{O,\text{bnd}} = -R_0 \cdot S(z_O)$\\{\scriptsize span $= R_0$}};

% ── 기준선 (중앙, y=0) ──
\node[procbox, fill=gray!15] (r0) at (5.0, 0) {기준선 $R_0 = 6.0$h};

% ── 불확실성 경로 (하단, y=-3.0) ──
\node[inputbox, fill=orange!10] (ku) at (0, -1.5) {\textbf{가중치} $k_U$\\$= 1 - k_O$};
\node[inputbox, fill=orange!10] (unorm) at (0, -4.0) {\textbf{원시 불확실성} $U_{\text{raw}}$\\$\geq 0$};
\node[procbox, fill=orange!5] (zu) at (5.0, -3.0)
  {$z_U = \kappa \cdot k_U \cdot \frac{R_0 \cdot U_{\text{raw}}}{R_{\max} - R_0}$};
\node[procbox, fill=orange!12] (subnd) at (10.5, -3.0)
  {$E_{U,\text{bnd}} = (R_{\max}\!-\!R_0) \cdot S(z_U)$\\{\scriptsize span $= R_{\max} - R_0$}};

% ── 가산 노드 ──
\node[sumbox, fill=green!8] (sum) at (14.5, 0) {$+$};

% ── 출력 ──
\node[outbox, fill=green!8] (drto) at (19.0, 0)
  {$\DRTO = R_0 + E_{O} + E_{U}$\\$\eta = \DRTO / R_0 \in (0,\; 4)$};

% ── 화살표: 관측성 경로 ──
\draw[arr, blue!60] (ko.east) -- (zo.150);
\draw[arr, blue!60] (oraw.east) -- (zo.210);
\draw[arr, blue!70] (zo.east) -- (sobnd.west);

% ── 화살표: 불확실성 경로 ──
\draw[arr, orange!60] (ku.east) -- (zu.150);
\draw[arr, orange!60] (unorm.east) -- (zu.210);
\draw[arr, orange!70] (zu.east) -- (subnd.west);

% ── 화살표: 3경로 → 가산 노드 (진입점 분리, 겹침 방지) ──
\draw[arr, blue!70]   (sobnd.east) -- ++(0.8,0) |- (sum.130);
\draw[arr, black]  (r0.east) -- ++(6.5,0) -- (sum.west);
\draw[arr, orange!70] (subnd.east) -- ++(0.8,0) |- (sum.230);

% ── 화살표: 가산 노드 → 출력 ──
\draw[arr, black] (sum.east) -- (drto.west);

% ── 매개변수 κ (중앙, 두 경로 사이) ──
\node[draw, rounded corners=3pt, fill=red!5, inner sep=3pt,
  font=\scriptsize, text=black] at (5.0, 0.9)
  {$\kappa = 1.2$};

% ── κ → 두 z 노드 연결 (점선) ──
\draw[densely dashed, red!40, thin, -{Stealth[length=1.5mm]}]
  (5.0, 0.7) -- (zo.south);
\draw[densely dashed, red!40, thin, -{Stealth[length=1.5mm]}]
  (5.0, -0.7) -- (5.0, -2.2);

% ── 범위 라벨 ──
\node[font=\scriptsize, text=black] at (10.5, 4.5)
  {$E_{O,\text{bnd}} \in (-R_0,\; 0]$};
\node[font=\scriptsize, text=black] at (10.5, -4.5)
  {$E_{U,\text{bnd}} \in [0,\; R_{\max}-R_0)$};

% ── 물리적 경계 ──
\node[draw, dashed, rounded corners=3pt, inner sep=5pt,
  font=\scriptsize, text=black] at (19.0, -2.5)
  {물리적 경계: $0 < \DRTO < R_{\max} = 24.0$h};

\end{tikzpicture}%
}% end resizebox
\caption{$\DRTO$ 개별 시그모이드 바운딩 모델 구조}
\label{fig:model-structure}
\end{figure}


%% -------------------------------------------------------------------
\subsection{변수 및 매개변수 정의}
\label{subsec:parameters}
%% -------------------------------------------------------------------

\p{[}표~\ref{tab:parameters}\p{]}은 $\DRTO$ 모델의 전체 변수와 매개변수를 정리한 것이다.
각 매개변수의 설정 근거는 연구자의 논문 \citet{lee2026drto}, 국제 표준, 또는 수학적 제약 조건에
기반한다.

\begin{table}[htbp]
\centering
\caption{DRTO 모델 변수 및 매개변수 정의}
\label{tab:parameters}
\small
\begin{tabularx}{\textwidth}{l l c c X}
\toprule
\textbf{기호} & \textbf{명칭} & \textbf{값/범위} & \textbf{단위} & \textbf{설정 근거} \\
\midrule
$R_0$ & 기준선 RTO & 6.0 & h & \nrev{ISO 22301:2019} BIA 기반 표준 설정값 \\
$R_{\max}$ & DRTO 상한 (MTPD) & 24.0 & h & $R_{\max} = 4 \cdot R_0$ (MTPD 기반 상한) \\
$\kappa$ & 곡률 매개변수 & 1.2 & --- & 다기준 최적화($\kappa^{*}$)~\citep{lee2026drto} \\
$S(z)$ & 수정 시그모이드 & $(-1, 1)$ & --- & $S(z) = 2/(1+e^{-z})-1$ \\
$k_O$ & 관측성 가중치 & $[0, 1]$ & --- & $\mathrm{Uniform}(0, 1)$으로 시뮬레이션 \\
$O_{\text{raw}}$ & 관측성 원시값 & $[0, 1]$ & --- & $\mathrm{Uniform}(0, 1)$으로 시뮬레이션 \\
$k_U$ & 불확실성 가중치 & $[0, 1]$ & --- & $k_U = 1 - k_O$ (상보적) \\
$z_O$ & 관측성 인자 & $\geq 0$ & --- & $\kappa \cdot k_O \cdot O_{\text{raw}}$ \\
$z_U$ & 불확실성 인자 & $\geq 0$ & --- & $\kappa \cdot k_U \cdot R_0 \cdot U_{\text{raw}} / (R_{\max} - R_0)$ \\
$U_{\text{raw}}^{(G)}$ & Gaussian 불확실성 & $N(3, 1)$ & --- & 평균 $\mu=3$, 표준편차 $\sigma=1$ \\
$U_{\text{raw}}^{(P)}$ & Pareto 불확실성 & $6 \cdot \mathrm{Pareto}(3)$ & --- & $\alpha=3$, $x_m=1$, 스케일링($\times 6$) \\
$E_{O,\text{bnd}}$ & 관측성 바운딩 조정량 & $(-R_0, 0]$ & h & span $= R_0 = 6.0$h \\
$E_{U,\text{bnd}}$ & 불확실성 바운딩 조정량 & $[0, R_{\max}\!-\!R_0)$ & h & span $= R_{\max} - R_0 = 18.0$h \\
$\eta$ & 효율비 & $(0, 4)$ & --- & $\DRTO / R_0$ \\
\bottomrule
\end{tabularx}
\end{table}


%% -------------------------------------------------------------------
\subsection{물리적 경계 보장 증명}
\label{subsec:boundary-proof}
%% -------------------------------------------------------------------

\phantomsection\label{prop:boundary}%
임의의 $\kappa > 0$, $k_O \in [0, 1]$, $O_{\text{raw}} \in [0, 1]$,
$z_U \geq 0$에 대해 $0 < \DRTO < R_{\max}$가 성립한다.(명제 2.1)

이 명제의 증명은 다음과 같다.
$z_O = \kappa \cdot k_O \cdot O_{\text{raw}} \geq 0$이므로
$S(z_O) \in [0, 1)$이다. 따라서
$E_{O,\text{bnd}} = -R_0 \cdot S(z_O) \in (-R_0, 0]$이다.
$z_U \geq 0$이므로
$S(z_U) \in [0, 1)$이다. 따라서
$E_{U,\text{bnd}} = (R_{\max} - R_0) \cdot S(z_U) \in [0, R_{\max} - R_0)$이다.
이 두 부등식을 합산하면 $\DRTO = R_0 + E_{O,\text{bnd}} + E_{U,\text{bnd}} \in (0,\; R_{\max})$이므로
$0 < \DRTO < R_{\max}$가 $\kappa$에 무관하게 성립한다. (명제 증명 완료)

위 증명은 $z_U \geq 0$ (즉, $U_{\text{raw}} \geq 0$) 전제 하에 성립한다.
C2 조건에서 $N(3,1)$의 음수 표본($0.135\%$)은 $z_U < 0$을 유발하여
$E_{U,\text{bnd}} < 0$이 되나,
$|E_{O,\text{bnd}}| < R_0$이고 $|E_{U,\text{bnd}}| < R_{\max} - R_0$이므로
$\DRTO > R_0 - R_0 - (R_{\max} - R_0) = 2R_0 - R_{\max}$이다.
$R_{\max} = 4R_0$일 때 이론적 하한은 $-2R_0$이나,
실제 시뮬레이션에서 $\DRTO$ 최솟값은 $2.861$h로 $0$을 충분히 상회한다.
즉 이론 하한이 음수임에도 음수 표본 출현 확률이 $0.135\%$로 작아 실질적 영향은 무시 가능하다.

\medskip
\chg{\noindent 이 경계 보장은 세 가지 실무적 의의를 갖는다. 첫째, 양수
보장($\DRTO > 0$)으로 $\DRTO$가 음수가 되는 물리적으로 무의미한 상황을
원천적으로 방지한다. 둘째, 상한 보장($\DRTO < R_{\max}$)으로
\nrev{ISO 22301:2019}의 MTPD 제약을 수학적으로 충족한다.
셋째, 클리핑(hard bounding)과 달리 경계 근처에서 미분 가능하여 매개변수 변화에
연속적으로 응답한다.}


%% ===================================================================
\section{4개 실험 조건 (C0--C3)}
\label{sec:conditions}
%% ===================================================================

$\DRTO$ 모델의 검증을 위해 관측성과 불확실성의 활성화 조합에 따라
네 개의 실험 조건(C0--C3)을 체계적으로 정의한다.
이 중 C0(기준선)과 C1(관측성 활성화)은 \citet{lee2026drto}에서
\chg{수립되고 검증된} 조건이며, C2(가우시안 불확실성)와 C3(파레토 불확실성)는
본 연구에서 새롭게 도입한 조건이다.
조건 설계의 핵심 원리는 \textbf{점진적 활성화(progressive activation)}로서,
C0에서 C3까지 한 번에 하나의 요소를 추가하여 각 요소의 개별 효과와
결합 효과를 분리 분석할 수 있도록 한다.

\p{[}표~\ref{tab:conditions}\p{]}는 각 조건의 구성과 특성을 비교한 것이다.

\begin{table}[htbp]
\centering
\caption{4개 실험 조건(C0--C3) 정의 및 시뮬레이션 결과 요약}
\label{tab:conditions}
\small
\begin{tabularx}{\textwidth}{c l c c l c X}
\toprule
\textbf{조건} & \textbf{명칭} & \textbf{관측성} & \textbf{불확실성}
  & \textbf{$U_{\text{raw}}$ 분포} & \textbf{평균 $\eta$} & \textbf{수식} \\
\midrule
\textbf{C0} & 정적 기준선 & OFF & OFF & --- & 1.000 &
  $\DRTO = R_0$ \\
\textbf{C1} & 관측성 전용 & ON & OFF & --- & 0.853 &
  $\DRTO = R_0 + E_{O,\text{bnd}}$ \\
\textbf{C2} & Gaussian 불확실성 & ON & ON & $N(3, 1)$ & 1.682 &
  $\DRTO = R_0 + E_{O,\text{bnd}} + E_{U,\text{bnd}}^{(G)}$ \\
\textbf{C3} & Pareto 불확실성 & ON & ON & $6 \cdot \mathrm{Pa}(3)$ & 1.514 &
  $\DRTO = R_0 + E_{O,\text{bnd}} + E_{U,\text{bnd}}^{(P)}$ \\
\bottomrule
\end{tabularx}
\end{table}


\p{[}그림~\ref{fig:conditions-progression}\p{]}은 C0에서 C3까지의 조건 진행과
각 단계에서 활성화되는 요소를 도식화한 것이다.
\chg{정적 조건 C0에서 관측성을 도입한 C1로 이어지는 공통 경로를 거친 뒤, 가우시안 불확실성의 C2와 파레토 불확실성의 C3가 C1에서 병렬로 분기하는 구조이다.}

\begin{figure}[htbp]
\centering
\resizebox{0.95\textwidth}{!}{%
\begin{tikzpicture}[
  cond/.style={draw, rounded corners=6pt, minimum width=30mm,
    minimum height=20mm, align=center, font=\small, line width=0.7pt},
  arr/.style={-{Stealth[length=3mm]}, thick},
  arrlbl/.style={font=\scriptsize, fill=white, inner sep=2pt,
    rounded corners=1pt},
  comp/.style={font=\scriptsize, text=black, align=center},
  efflbl/.style={font=\scriptsize, align=center, text=black}
]

% 조건 노드
\node[cond, fill=gray!15] (c0) at (0, 0)
  {\textbf{C0}\\정적 기준선\\$\eta = 1.000$\\$\DRTO = 6.0$h};
\node[cond, fill=blue!12] (c1) at (5.5, 0)
  {\textbf{C1}\\관측 전용\\$\bar{\eta} = 0.853$\\$\overline{\DRTO} = 5.12$h};
\node[cond, fill=green!12] (c2) at (11, 0)
  {\textbf{C2}\\Gaussian 불확실성\\$\bar{\eta} = 1.682$\\$\overline{\DRTO} = 10.09$h};
\node[cond, fill=red!12] (c3) at (16.5, 0)
  {\textbf{C3}\\Pareto 불확실성\\$\bar{\eta} = 1.514$\\$\overline{\DRTO} = 9.09$h};

% 화살표 및 라벨 (상단: 전이 조건, 중앙: Cohen's d)
\draw[arr] (c0) -- (c1);
\node[arrlbl, above=3pt] at ($(c0)!0.5!(c1)$) {$+O_{\text{raw}},\;k_O$};
\node[efflbl] at ($(c0)!0.5!(c1) + (0, 1.3)$)
  {Cohen's $d = -1.180$\\(C1 vs C0)};

\draw[arr] (c1) -- (c2);
\node[arrlbl, above=3pt] at ($(c1)!0.5!(c2)$) {$+U \sim N(3,1)$};
\node[efflbl] at ($(c1)!0.5!(c2) + (0, 1.3)$)
  {Cohen's $d = +1.871$\\(C2 vs C1)};

%% 분기 B: C1 → C3 (ㄷ자 직각 경로: 아래→오른쪽→위)
\draw[arr] (c1.south) -- ++(0,-0.6)
          -- ($(c3.south) + (0,-0.6)$) -- (c3.south);
\node[arrlbl, above=2pt] at ($(c1.south)!0.5!(c3.south) + (0,-0.6)$)
  {$+U \sim 6\!\cdot\!\mathrm{Pa}(3)$};
\node[efflbl] at ($(c2)!0.5!(c3) + (0, 1.3)$)
  {CVaR$_{99}$ 증가\\(C3 vs C2)};

% 하단 구성 요소
\node[comp, below=5mm of c0] {관측성: OFF\\불확실성: OFF};
\node[comp, below=5mm of c1] {관측성: ON\\불확실성: OFF};
\node[comp, below=5mm of c2] {관측성: ON\\불확실성: 가우시안(경꼬리)};
\node[comp, below=5mm of c3] {관측성: ON\\불확실성: 파레토(중꼬리)};

\end{tikzpicture}%
}% end resizebox
\caption{4개 실험 조건의 점진적 활성화}
\label{fig:conditions-progression}
\end{figure}


%% -------------------------------------------------------------------
\subsection{\chg{정적 기준선 C0}}
\label{subsec:c0}
%% -------------------------------------------------------------------

C0 조건은 관측성과 불확실성이 모두 비활성화된 상태로서,
기존 \nrev{ISO 22301:2019} 기반 정적 RTO 체계를 재현한다.
$O_{\text{raw}} = 0$, $U_{\text{raw}} = 0$이므로 $z_O = z_U = 0$이 되어
$S(0) = 0$이 성립하며, 따라서 $\DRTO = R_0 = 6.0$h, $\eta = 1.000$이다.
C0는 후속 조건들의 비교 기준선(baseline)으로 기능하며,
모든 동적 효과의 출발점을 제공한다.


%% -------------------------------------------------------------------
\subsection{\chg{C1의 관측성 기반 동적 조정}}
\label{subsec:c1}
%% -------------------------------------------------------------------

C1 조건은 관측성만 활성화하고 불확실성은 비활성화한 상태이다.
$U_{\text{raw}} = 0$이므로 $E_{U,\text{bnd}} = 0$이 되어
$\DRTO = R_0 - R_0 \cdot S(\kappa \cdot k_O \cdot O_{\text{raw}})$이다.
$S(\cdot) \geq 0$이므로 $\DRTO \leq R_0$가 보장되며,
이는 관측성 활성화가 복구 시간을 단축하는 방향으로만 작용함을 의미한다.
개별 바운딩의 핵심 장점으로, C1에서는 $R_{\max}$가 수식에 등장하지 않아
관측성만의 효과를 불확실성 상한과 독립적으로 분석할 수 있다.

$k_O$와 $O_{\text{raw}}$가 모두 최대값(1.0)일 때
$z_O = 1.2$가 되어
$S(1.2) \approx 0.537$이므로
$\DRTO \approx 2.78$h가 된다.
이는 완전 관측 상태에서 약 $53.7\%$의 복구 시간 단축이 가능함을 나타낸다.

\chg{이처럼 C1에서는 관측성 활성화가 복구 시간을 단축하는 방향으로만 작용하며, 그 평균적 단축 효과와 효과 크기가 기준선(C0) 대비 통계적으로 유의한가는 본 장이 제시한 이론적 예측에 해당한다. 표본 기반 실증값(평균 효율비, 단축률, Cohen's $d$)은 제4장에서 정량적으로 검증한다(\secref{sec:result_c0c1} 참조).}


%% -------------------------------------------------------------------
\subsection{\chg{C2의 가우시안 불확실성 결합}}
\label{subsec:c2}
%% -------------------------------------------------------------------

C2 조건은 관측성에 더하여 가우시안(Gaussian) 불확실성을 결합한 상태이다.
불확실성 원시값은 $U_{\text{raw}}^{(G)} \sim N(3, 1)$으로 생성되며,
$z_U = \kappa \cdot k_U \cdot R_0 \cdot U_{\text{raw}}^{(G)} / (R_{\max} - R_0)$으로 변환한 후
수식~(\ref{eq:E_U_bnd})에 따라 불확실성 효과 $E_{U,\text{bnd}}$를 산출한다.

\chg{한편 가우시안 분포는 음수 표본을 낼 수 있어 그 처리를 짚어 둔다.} 가우시안 분포의 지지(support)는 $(-\infty, +\infty)$이므로
$U_{\text{raw}}^{(G)} < 0$인 표본이 이론적으로 발생할 수 있다.
그러나 $N(3,1)$에서 $\Pr(U < 0) = \Phi(-3) \approx 0.135\%$로서
$n = 10{,}000$ 표본 중 13건에 불과하다.
해당 표본에서는 $z_U < 0$이 되어 $E_{U,\text{bnd}} < 0$이 산출되나,
실측 범위는 $[-1.13,\; -0.01]$h에 그치며
$\DRTO$ 최솟값이 $2.861$h로 물리적 하한($0$)을 충분히 상회한다.
따라서 절단정규분포(truncated normal)나 클리핑(clipping) 등의
인위적 분포 변형을 적용하지 않는다.
\chg{이는 세 가지 근거에 의한다. 첫째, 빈도 측면에서 음수 표본은 $0.13\%$(실측
13/10{,}000, 이론 $0.135\%$)에 불과하여 통계적 결론에 무시할 수 있는 영향만 미친다.
둘째, 경계 측면에서 해당 표본에서도 $\DRTO = 2.861$h로 0보다 충분히 크다. 셋째, 분포
보존 측면에서 절단을 적용하면 $N(3,1)$의 확률 구조가 왜곡되어 C2와 C3 사이의 공정한
비교(동일 기대값 $3.0$ 설계)가 훼손된다.}

가우시안 분포는 대칭적이고 경꼬리(light-tail) 특성을 가지므로,
C2 조건은 일상적 운영 변동, \chg{곧 인력 가용성의 일시적 변화와
장비 성능의 통상적 편차, 절차 수행 시간의 자연적 변이를} 반영한다.

\chg{C2에서는 이러한 가우시안 불확실성의 결합이 복구 시간을 연장하는 방향으로 작용하리라 예측되며, 그 평균적 연장 크기와 C1 대비 효과 크기는 제4장에서 정량적으로 검증한다(\secref{sec:result_c2} 참조).}


%% -------------------------------------------------------------------
\subsection{\chg{C3의 파레토 Heavy-tail 결합}}
\label{subsec:c3}
%% -------------------------------------------------------------------

C3 조건은 C1에서 분기하여 가우시안(C2) 대신 파레토(중꼬리, heavy-tail) 분포로
불확실성을 모델링한 변종이다. 불확실성 원시값은
$U_{\text{raw}}^{(P)} \sim 6 \cdot \mathrm{Pareto}(\alpha = 3)$으로
생성되며, $z_U = \kappa \cdot k_U \cdot R_0 \cdot U_{\text{raw}}^{(P)} / (R_{\max} - R_0)$으로 변환된다.
$\alpha = 3$은 평균은 유한하나 분산이 가우시안($\sigma^2=1$) 대비 약 27배인 분포를 제공한다.

C2와 C3의 불확실성 분포는 모집단 평균($\mu = 3.0$)이 동일하도록 설계되어
공정한 비교가 가능하다. 두 분포의 차이는 분산과 극단 백분위에서 발현된다.
가우시안의 99백분위가 약 5.33인 반면, 파레토의 99백분위는 약 21.85로서
약 4.1배의 차이를 보인다.

\chg{파레토 분포는 중꼬리 특성을 가지므로, C3는 평균 수준에서는 C2와 유사하더라도 극단 백분위로 갈수록 $\DRTO$가 C2를 크게 상회하리라 예측된다. 두 분포의 누적분포함수가 교차하는 지점이 존재하며, 본 연구에서는 이를 ``CDF 교차점(crossover)''으로 명명한다. 교차점을 경계로 핵심 영역(Core)과 꼬리 영역(Tail)을 분할하면, 꼬리 영역에서는 시그모이드 포화로 인해 관측성 가중치 $k_O$의 회귀 계수가 감쇠하는 이중채널 비선형 상호작용(dual-channel nonlinear interaction)이 나타나리라 예측된다. 교차점의 구체적 위치와 조건별 평균 효율비, 그리고 감쇠의 실증값은 \jrev{4.10절}(이중채널)에서 상세히 분석한다.}
           % 2.7--2.8 DRTO 모델·실험 조건


%% ===================================================================
\section{시그모이드 바운딩의 수학적 성질}
\label{sec:sigmoid-properties}
%% ===================================================================

본 절에서는 $\DRTO$ 모델의 기저 함수인 수정 시그모이드 $S(z)$의
수학적 성질을 체계적으로 분석한다.

%% -------------------------------------------------------------------
\subsection{기본 해석학적 성질}
\label{subsec:basic-properties}
%% -------------------------------------------------------------------

수정 시그모이드 함수 $S(z) = 2/(1+e^{-z}) - 1$은 다음의 성질을 갖는다.
\chg{앞서 2.5(시그모이드 함수와 바운딩)에서 바운딩 기법의 관점에서 개관한 기본 성질을, 2.9(시그모이드 바운딩의 수학적 성질)에서는 DRTO 바운딩 함수로서 2계 도함수와 변곡점, 곡률 거동까지 포함하여 엄밀하게 정리한다.}

\chg{첫째, 전역 미분 가능성($C^{\infty}$) 측면에서 $S(z)$는 지수 함수의 합성으로 구성되므로
전체 정의역에서 무한 번 미분 가능하며, 1차 도함수가 $S'(z) = 2e^{-z}/(1+e^{-z})^2 > 0$
($\forall z \in \mathbb{R}$)이므로 전 구간에서 단조증가한다. 둘째, 단조증가성에 의해 입력의
순서가 출력에 보존되어 $z_1 < z_2$이면 $S(z_1) < S(z_2)$가 성립하며, 이는 관측성이 높을수록
$\DRTO$가 더 감소하고 불확실성이 클수록 더 증가하는 물리적 직관과 일치한다. 셋째, 경계
보장 측면에서 $z \to +\infty$일 때 $S(z) \to +1$, $z \to -\infty$일 때 $S(z) \to -1$이므로
$S: \mathbb{R} \to (-1, 1)$이고, 점근선에 도달하되 등호가 성립하지 않으므로 $\DRTO$가 정확히
$0$ 또는 $R_{\max}$에 도달하는 것은 불가능하다. 넷째, 원점 대칭성(홀함수)로서
$S(-z) = -S(z)$이고 $S(0) = 0$이어서 효과가 없을 때 기준선을 유지하는 중립 조건을 보장하며,
$z_O = z_U = 0$일 때 $E_{O,\text{bnd}} = E_{U,\text{bnd}} = 0$이 되어 $\DRTO = R_0$로 기준선이
유지된다. 다섯째, 2계 도함수가 $S''(z) = -2 e^{-z} (1 - e^{-z})/(1 + e^{-z})^3$로서 $z = 0$에서
부호가 바뀌는 변곡점을 가지며, 이 변곡점이 ``효과 없음'' 지점($S(0)=0$)과 일치하므로 작은
입력에서는 함수 거동이 선형에 가깝다가 $|z|$가 커질수록 비선형성이 강화되는 구조를 갖는다.}


%% -------------------------------------------------------------------
\subsection{작은 $\kappa \cdot z$ 값 영역에서의 선형 근사}
\label{subsec:linear-approx}
%% -------------------------------------------------------------------

$|z|$가 충분히 작을 때, $S(z)$는 \m{식~(\ref{eq:linear-approx})에 나타낸 바와 같이} 1차 테일러 전개로 근사된다.

\begin{equation}
  S(z) \approx S'(0) \cdot z = \frac{1}{2} z, \quad |z| \ll 1
  \label{eq:linear-approx}
\end{equation}

이 근사는 $|z| \leq 0.5$에서 약 $2\%$ 이하의 오차로 성립하며,
$|z| \leq 0.8$에서는 약 $5\%$ 이하의 오차를 갖는다.

이 근사를 C1 수식에 적용하면
$E_{O,\text{bnd}} \approx -R_0 \cdot \frac{1}{2} \cdot \kappa \cdot k_O \cdot O_{\text{raw}}$
$= -\frac{\kappa}{2} \cdot R_0 \cdot k_O \cdot O_{\text{raw}}$가 되어
선형 모델과의 대응 관계가 명확해진다.

$\kappa = 1.2$일 때 $z_O = \kappa \cdot k_O \cdot O_{\text{raw}} \in [0, 1.2]$이다.
$z_O$의 상한($k_O = O_{\text{raw}} = 1$일 때 $z_O = 1.2$)에서 $S(1.2) = 0.537$이고
선형 근사는 $\frac{1}{2} \times 1.2 = 0.600$으로서, 근사 오차는 약 $10.5\%$이다.
따라서 $\kappa = 1.2$는 선형 영역에서 비선형 전이 영역으로 진입하는 경계에 위치하며,
시그모이드의 비선형 압축 효과(선형 근사 대비 출력값 감소)가 의미 있게 발현되면서도
포화에는 도달하지 않는 균형 영역에 해당한다.

\p{[}그림~\ref{fig:sigmoid-curve}\p{]}은 $\kappa = 1.2$에서의 시그모이드 함수
$S(z)$와 선형 근사 $z/2$를 비교한 것이다.
\chg{그림에서 관측성 인수 $z_O$의 범위 $[0, 1.2]$가 시그모이드의 선형 영역에서 비선형 전이 영역으로 확장됨을 볼 수 있으며, $S(1.2) = 0.537$로서 선형 근사값 $0.600$ 대비 약 $10.5\%$의 압축이 발생한다.}

\begin{figure}[htbp]
\centering
\begin{tikzpicture}
\begin{axis}[
  width=0.82\textwidth,
  height=0.50\textwidth,
  xlabel={$z$},
  ylabel={$S(z)$},
  xmin=-4, xmax=4,
  ymin=-1.15, ymax=1.15,
  grid=major,
  grid style={black},
  legend pos=north west,
  legend style={font=\small, rounded corners=2pt},
  axis lines=middle,
  every axis x label/.style={at={(axis description cs:1.02,0.5)},
    anchor=west, font=\small},
  every axis y label/.style={at={(axis description cs:0.5,1.04)},
    anchor=south, font=\small},
  clip=false
]
  % 시그모이드 곡선
  \addplot[thick, black, domain=-4:4, samples=200]
    {2/(1+exp(-x)) - 1};
  \addlegendentry{$S(z) = 2/(1+e^{-z})-1$}

  % 선형 근사
  \addplot[thick, dashed, blue!70, domain=-2.2:2.2, samples=100]
    {0.5*x};
  \addlegendentry{선형 근사 $z/2$}

  % 경계선
  \addplot[thin, dashed, black] coordinates {(-4, 1) (4, 1)};
  \addplot[thin, dashed, black] coordinates {(-4, -1) (4, -1)};

  % κ·k_O·O_raw 범위 표시
  \draw[thick, red!70, {Stealth[length=2mm]}-{Stealth[length=2mm]}]
    (axis cs:0, -1.08) -- (axis cs:1.2, -1.08);
  \node[font=\scriptsize, text=black, anchor=north] at (axis cs:0.6, -1.08)
    {$\kappa k_O O_{\text{raw}} \in [0, 1.2]$};

  % S(1.2) 표시
  \draw[thin, dotted, red!50] (axis cs:1.2, 0) -- (axis cs:1.2, 0.537);
  \node[font=\tiny, text=black, anchor=west] at (axis cs:1.25, 0.537)
    {$S(1.2)=0.537$};
  \filldraw[red!70] (axis cs:1.2, 0.537) circle (1.5pt);

  % 경계 라벨
  \node[font=\scriptsize, text=black, anchor=west] at (axis cs:4.1, 1) {$+1$};
  \node[font=\scriptsize, text=black, anchor=west] at (axis cs:4.1, -1) {$-1$};

\end{axis}
\end{tikzpicture}
\caption{수정 시그모이드 함수 $S(z)$와 선형 근사 $z/2$의 비교}
\label{fig:sigmoid-curve}
\end{figure}


%% ===================================================================
\section{최적 곡률 파라미터 $\kappa^{*} = 1.2$}
\label{sec:kappa-optimization}
%% ===================================================================

곡률 매개변수 $\kappa$는 시그모이드 바운딩 모델에서 입력 변수의 변화에 대한
$\DRTO$의 반응 강도를 결정하는 핵심 설계 변수이다.
$\kappa$가 너무 작으면 시그모이드가 선형 영역에 머물러
입력 변화에 대한 $\DRTO$의 반응이 미약하고(과소 반응),
$\kappa$가 너무 크면 적은 입력에서도 $|S|$가 1에 근접하여 다수 표본이
포화 영역에 진입함으로써 극단값 간의 구분력이 상실된다(조기 포화).
본 연구는 \citet{lee2026drto}의 다기준 최적화(multi-criteria
optimization) 방법을 적용하여 $\kappa^{*} = 1.2$를 도출하였다.


%% -------------------------------------------------------------------
\subsection{4개 최적화 기준}
\label{subsec:four-criteria}
%% -------------------------------------------------------------------

$\kappa$의 최적값 선정을 위해 다음의 4개 기준을 동시에 고려한다.

\chg{첫째는 효과 크기(Effect Size)로서 C1 대 C0의 Cohen's $d$ \m{절댓값}
$|d_{\text{C1vsC0}}|$이며, $\kappa$가 클수록 관측성의 효과가 증가하여 $|d|$가 커지므로
$|d| \geq 0.8$(대효과)을 기준으로 한다. 둘째는 포화율(Saturation Rate)로서 관측성
채널에서 $|S(z_O)| \geq \erf{0.99}{0.95}$인 표본의 비율이며, $\kappa$가 너무 크면 다수의 표본이 포화
영역에 진입하여 입력 변화에 대한 구분력이 상실되므로 포화율 $\leq 5\%$를 기준으로 한다.
셋째는 회귀 설명력($R^2$)으로서 C1에서의 2차 다항 회귀 결정계수이며, 시그모이드의
비선형성이 과도하면 2차 다항 근사의 정확도가 저하되므로 $R^2 \geq 0.95$를 기준으로 한다.
넷째는 RTO 단축률(Reduction Rate)로서 C1 평균 효율성 비율에 의한 단축률
$(1 - \bar{\eta}_{\text{C1}}) \times 100\%$이며, 실무적으로 의미 있는 단축 효과를 제공하는지를
평가하기 위해 단축률 $\geq 10\%$를 기준으로 한다.}


%% -------------------------------------------------------------------
\subsection{종합점수 $F(\kappa)$와 격자 탐색}
\label{subsec:composite-score}
%% -------------------------------------------------------------------

4개 기준을 종합하는 종합 점수(composite score) $F(\kappa)$를
\m{식~(\ref{eq:composite})에 나타낸 바와 같이} 곱함수(product function)로 정의한다.

\begin{equation}
  F(\kappa) = f_1(\kappa) \cdot f_2(\kappa) \cdot f_3(\kappa) \cdot f_4(\kappa)
  \label{eq:composite}
\end{equation}

여기서 $f_i(\kappa) \in [0, 1]$은 각 기준의 부분 점수이다.
각 부분 점수는 해당 기준에서 가능한 최대값으로 정규화되어 $f_i \in [0,1]$이며,
따라서 $F \in [0,1]$이고 $F = 1$은 모든 기준이 동시에 최적치를 달성한 상태를
의미한다.
곱 구조를 채택한 근거는 다음과 같다.
하나의 기준이라도 0이면 $F = 0$이 되어 모든 기준의 동시 충족을 강제하고,
기준 간 대체(trade-off)를 허용하지 않으며,
기준별 가중치 결정이 불필요하다.
$F(\kappa)$가 최대인 $\kappa$를 $\kappa^{*}$로 선정한다.

각 부분 점수 $f_i(\kappa)$의 구체적 함수형은 다음과 같다.

\chg{첫째, 실무 유의성을 나타내는 $f_1$은} \m{식~(\ref{eq:f1})에 나타낸 바와 같이} 평균 감소율 $\bar{r}(\kappa) = (1 - \bar{\eta}) \times 100$이
    목표 범위(5--30\%)에 있을 때 최대이며 15\%에서 정점을 갖는다.
    \begin{equation}
      f_1(\kappa) = \exp\!\left(-\frac{(\bar{r}(\kappa) - 15)^2}{2 \cdot 10^2}\right)
      \label{eq:f1}
    \end{equation}
    중심값 15\%는 RTO 감소율이 너무 작으면(5\% 미만) 도입 효과가 미미하고 너무 크면
    (30\% 초과) 비현실적이므로 목표 범위 5--30\%의 중앙값을 채택한 것이다. 분모의 $10^2$은
    가우시안 커널의 표준편차 $\sigma = 10$에 해당하며, 데이터에서 추정한 값이 아니라 목표 범위의
    구조에 맞추어 설정한 설계 파라미터이다. 목표 범위의 중심이 15\%이고 하한이 5\%이므로 중심에서
    하한까지의 거리는 $15 - 5 = 10$\%p이며, 이 거리를 $1\sigma$로 놓으면($\sigma = 10$) 가우시안
    커널의 확률적 해석과 자연스럽게 대응된다. $\sigma = 10$일 때 중심으로부터의 거리별 $f_1$ 값은
    \p{[}표~\ref{tab:f1-sigma10}\p{]}과 같다.
    \begin{table}[htbp]
    \centering
    \caption{$\sigma = 10$, 중심($\bar{r}=15\%$)에서 거리별 포화도 가중치 $f_1$}
    \label{tab:f1-sigma10}
    \small
    \begin{tabular}{c c c l}
    \toprule
    $\bar{r}$ & 중심 거리 & $f_1$ & 해석 \\
    \midrule
    15\% & $0\sigma$ & $e^{0} = 1.000$ & \m{최적점(정점)} \\
    5\% 또는 25\% & $1\sigma$ & $e^{-0.5} \approx 0.607$ & 목표 범위 내, 충분히 허용 \\
    35\% & $2\sigma$ & $e^{-2} \approx 0.135$ & 목표 초과, 강하게 억제 \\
    45\% & $3\sigma$ & $e^{-4.5} \approx 0.011$ & 사실상 배제 \\
    \bottomrule
    \end{tabular}
    \end{table}

\chg{둘째, 비포화를 나타내는 $f_2$는} \m{식~(\ref{eq:f2})에 나타낸 바와 같이} 포화율(시그모이드 출력 $> 0.95$인 비율)이 5\% 미만이면 1, 아니면 0이다.
    \begin{equation}
      f_2(\kappa) = \begin{cases}
        1.0 & \text{if sat}(\kappa) < 5\% \\
        0.0 & \text{otherwise}
      \end{cases}
      \label{eq:f2}
    \end{equation}
    시그모이드 출력이 0.95를 초과하면 입력 변화에 대한 출력 변화가 거의 없는 포화(saturation)
    영역에 진입하여 $\DRTO$가 입력 변수의 차이를 변별하지 못하므로, 포화 비율을 5\% 미만으로
    억제하는 이진 게이트로 설계하여 이 조건을 위반하는 $\kappa$를 후보에서 즉시 배제한다.

\chg{셋째, 효과 크기를 나타내는 $f_3$는} \m{식~(\ref{eq:f3})에 나타낸 바와 같이} C0 대비 Cohen's $d$의 \m{절댓값}이 $\geq 0.8$이면 1이다.
    \begin{equation}
      f_3(\kappa) = \begin{cases}
        1.0 & \text{if } |d(\kappa)| \geq 0.8 \\
        |d(\kappa)| / 0.8 & \text{otherwise}
      \end{cases}
      \label{eq:f3}
    \end{equation}
    식~(\ref{eq:f3})에서 C0 조건은 $\eta \equiv 1.000$의 결정론적 상수이므로 합동 표준편차 공식의 분모가
    인위적으로 축소되는 것을 피하기 위해 C1 자체의 표준편차를 분모로 사용하는 일표본 효과
    크기를 적용한다. $\bar{\eta}_{\text{C1}} = 0.853$, $s_{\text{C1}} = 0.125$이므로
    $d = (0.853 - 1.000)/0.125 = -1.176 \approx -1.180$이고, $|d| \geq 0.8$이므로 $f_3 = 1.0$이다. Cohen's $d$의 정의와 해석 기준은 4.8(조건간 효과 크기 분석)을 참조한다.

\chg{넷째, 설명력을 나타내는 $f_4$는} \m{식~(\ref{eq:f4})에 나타낸 바와 같이} $\eta_{C1}$에 대한 회귀 $R^2$이 $\geq 0.95$이면 1이다.
    \begin{equation}
      f_4(\kappa) = \begin{cases}
        1.0 & \text{if } R^2(\kappa) \geq 0.95 \\
        R^2(\kappa) / 0.95 & \text{otherwise}
      \end{cases}
      \label{eq:f4}
    \end{equation}

\p{[}표~\ref{tab:kappa-grid}\p{]}는 $\kappa \in \{0.2, 0.5, 1.0, 1.2, 1.5, 2.0, 3.0, 5.0\}$에
대한 격자 탐색(grid search) 결과를 정리한 것이다.
\chg{탐색은 표본 크기 $n = 10{,}000$, seed $= 42$에서 수행하였으며, 표의 $R^2_{(2\mathrm{q})}$는 2차 다항 회귀의 결정계수이다.}

\p{[}표~\ref{tab:kappa-grid}\p{]}에서 확인되듯이, $\kappa = 1.2$에서 종합 점수
$F(\kappa) = 1.000$으로 최대값을 달성한다.
이 지점에서 Cohen's $|d| = 1.180$(대효과), 포화율 $0.0\%$,
$R^2 = 0.999$, 단축률 $14.7\%$가 동시에 충족된다.

\begin{table}[htbp]
\centering
\caption{$\kappa$ 격자 탐색 \m{결과($n = 10{,}000$, seed $= 42$)}}
\label{tab:kappa-grid}
\small
\begin{tabularx}{\textwidth}{c c c c c c c}
\toprule
$\kappa$ & 평균 $\bar{\eta}_{\text{C1}}$ & 단축률 (\%) & $|d|_{\text{C1vsC0}}$
  & $R^2_{(2\mathrm{q})}$ & 포화율 (\%) & $F(\kappa)$ \\
\midrule
0.2  & 0.975 &  2.5 & 1.150 & 1.000 &  0.0 & \erf{0.813}{0.459} \\
0.5  & 0.937 &  6.3 & 1.150 & 1.000 &  0.0 & \erf{0.855}{0.684} \\
1.0  & 0.876 & 12.4 & 1.170 & 0.999 &  0.0 & \erf{0.957}{0.967} \\
\rowcolor{green!10}
1.2  & 0.853 & 14.7 & 1.180 & 0.999 &  0.0 & \textbf{1.000} \\
1.5  & 0.818 & 18.2 & 1.200 & 0.997 &  0.0 & \erf{0.947}{0.951} \\
2.0  & 0.765 & 23.5 & 1.230 & 0.993 &  0.0 & \erf{0.859}{0.697} \\
3.0  & 0.673 & 32.7 & 1.310 & 0.975 &  0.0 & \erf{0.858}{0.209} \\
5.0  & 0.538 & 46.2 & 1.490 & 0.908 &  3.9 & \erf{0.566}{0.007} \\
\bottomrule
\end{tabularx}
\end{table}


$\kappa < 1.2$인 영역에서는 효과 크기와 단축률이 부족하고,
$\kappa > 1.2$인 영역에서는 포화가 시작되어 $R^2$가 하락한다.
$\kappa = 1.2$는 이 두 영역의 균형점(trade-off point)에 정확히 위치한다.
\p{[}표~\ref{tab:kappa-optimal}\p{]}는 $\kappa^{*} = 1.2$에서 4개 부분 점수의 개별
달성값과 종합점수 $F^{*}$를 분해한 것이다.

\begin{table}[htbp]
\centering
\caption{$\kappa^{*} = 1.2$에서의 4가지 기준 달성값}
\label{tab:kappa-optimal}
\begin{tabular}{llrrl}
\toprule
\textbf{기준} & \textbf{지표} & \textbf{달성값} & \textbf{$f_i$ 점수} & \textbf{판정} \\
\midrule
$f_1$ (실무 유의성)  & 평균 감소율      & 14.74 & 0.9997 & 목표 15\%에 근접 \\
$f_2$ (비포화)       & 포화율           & 0.0   & 1.000 & 포화 완전 회피 \\
$f_3$ (효과 크기)    & Cohen's $d$     & $-1.18$ & 1.000 & 대효과 달성 \\
$f_4$ (설명력)       & $R^2$           & 0.999   & 1.000 & 고설명력 달성 \\
\midrule
\multicolumn{2}{l}{\textbf{종합점수 $F^*$}} & & \textbf{0.9997} & \textbf{유일 극대} \\
\bottomrule
\end{tabular}
\end{table}

\p{[}그림~\ref{fig:kappa-composite}\p{]}은 종합점수 $F(\kappa)$의 곡선을
도시한 것으로, $\kappa^{*} = 1.2$에서 최대값을 달성하는 것을 보여준다.
\chg{이 지점에서 효과 크기와 비포화, $R^2$, 단축률의 네 기준이 균형을 이루어 $F = 1.000$을 달성하며, $\kappa < 1.0$은 과소 반응 영역에, $\kappa > 2.0$은 조기 포화 영역에 해당한다.}

\begin{figure}[htbp]
\centering
\begin{tikzpicture}
\ifdefined\ERRATAFIX\def\FKymin{0}\def\FKcoords{(0.2, 0.459) (0.5, 0.684) (1.0, 0.967) (1.2, 1.000) (1.5, 0.951) (2.0, 0.697) (3.0, 0.209) (5.0, 0.007)}\else\def\FKymin{0.4}\def\FKcoords{(0.2, 0.813) (0.5, 0.855) (1.0, 0.957) (1.2, 1.000) (1.5, 0.947) (2.0, 0.859) (3.0, 0.858) (5.0, 0.566)}\fi
\begin{axis}[
  width=0.80\textwidth,
  height=0.48\textwidth,
  xlabel={$\kappa$},
  ylabel={$F(\kappa)$},
  xmin=0, xmax=5.5,
  ymin=\FKymin, ymax=1.05,
  grid=major,
  grid style={black},
  axis lines=left,
  every axis x label/.style={at={(axis description cs:1.02,0.0)},
    anchor=west, font=\small},
  every axis y label/.style={at={(axis description cs:-0.02,1.02)},
    anchor=south, font=\small},
  clip=false,
  xtick={0, 0.5, 1.0, 1.2, 1.5, 2.0, 3.0, 5.0},
  xticklabels={0, 0.5, 1.0, \textbf{1.2}, 1.5, 2.0, 3.0, 5.0}
]
  % 데이터 점 연결 곡선 (smooth)
  \addplot[thick, black, smooth, mark=*, mark size=2pt]
    coordinates {
      \FKcoords
    };

  % 최적점 강조
  \filldraw[red!80] (axis cs:1.2, 0.996) circle (3pt);
  \node[font=\small, text=black, anchor=south west] at (axis cs:1.3, 0.996)
    {$\kappa^{*} = 1.2$, $F = 1.000$};

  % 과소 반응 영역
  \draw[thick, blue!50, {Stealth[length=2mm]}-] (axis cs:0.1, 0.45) -- (axis cs:0.8, 0.45);
  \node[font=\scriptsize, text=black, anchor=north] at (axis cs:0.45, 0.45)
    {과소 반응};

  % 조기 포화 영역
  \draw[thick, orange!60, -{Stealth[length=2mm]}] (axis cs:2.5, 0.45) -- (axis cs:5.3, 0.45);
  \node[font=\scriptsize, text=black, anchor=north] at (axis cs:3.9, 0.45)
    {조기 포화};

\end{axis}
\end{tikzpicture}
\caption{$\kappa$ 격자 탐색에 의한 종합점수 $F(\kappa)$ 곡선}
\label{fig:kappa-composite}
\end{figure}
     % 2.9--2.10 수학적 성질·κ 최적화


%% ===================================================================
\section{연구가설}
\label{sec:hypotheses}
%% ===================================================================

$\DRTO$ 프레임워크의 이론적 타당성과 실증적 유효성을 검증하기 위해
\jrev{10개의 연구가설을 수립한다.} 가설은 시뮬레이션 검증\jrev{(H1--H3, 3건)},
회귀 검증(H4--H5, 2건), 강건성 검증(H6--H7, 2건),
이중채널 비선형 효과(H8, 1건), 전문가 평가(H9--H10, 2건)의 5개 범주로 구성된다.

\chg{각 가설은 서로 다른 이론적 근거에서 도출된다. H1에서 H3까지와 H8은 본 장에서 확립한 모델의 해석적 구조, 즉 시그모이드 바운딩의 부호 조건과 경계 보장 및 포화 성질에서 연역되고, H4와 H5는 입력 변수와 응답 사이의 비선형 관계를 완전 2차 다항식으로 근사하는 반응표면법~\citep{box1951experimental}에, H6과 H7은 몬테카를로 시뮬레이션의 독립 시드 설계와 계수 안정성 원리에, H9와 H10은 정량과 정성 방법을 결합하는 혼합연구방법의 삼각검증 원리~\citep{creswell2018designing}에 각각 근거를 둔다.}

\chg{각 가설에 대해서는 도출 배경과 핵심 명제, 검증 방향의 세 요소를 일관되게
제시한다. 도출 배경은 왜 이 가설이 필요한가를, 핵심 명제는 그 수학적이고 통계적인
정의를, 검증 방향은 어느 장에서 어떻게 검증되는가를 밝힌다. 시뮬레이션 검증에
해당하는 H1에서 H3까지는 모형의 기본 동작을, 회귀 검증인 H4와 H5는 변수 간 통계적
관계를, 강건성 검증인 H6과 H7은 결과의 재현성을, 이중채널의 H8은 본 연구의 핵심 신규
발견을, 전문가 평가인 H9와 H10은 외부 타당성을 각각 담당한다.
\secref{sec:ch4_hypothesis_summary}에서 10개 가설이 모두 채택(PASS)으로 종합되며, 그
정량 근거는 가설별 정의와 시뮬레이션 및 회귀 결과로 뒷받침된다.}


%% -------------------------------------------------------------------
%% -------------------------------------------------------------------

\chg{가설 H1(관측성 효과)은 관측성이 활성화된 조건(C1)에서 $\DRTO$가 정적 기준선(C0)보다 작거나 같다는 것이다.}
\[
  H_1:\;\; \DRTO_{\textbf{C1}} \leq \DRTO_{\textbf{C0}} = R_0
  \quad (\text{준수율 } 100\%)
\]

H1은 관측성의 효과 방향성을 검증한다.
수학적으로, C1 조건에서 $z_O \geq 0$이므로 $S(z_O) \geq 0$이 되어
$E_{O,\text{bnd}} = -R_0 \cdot S(z_O) \leq 0$이 항상 성립한다.
따라서 $\DRTO_{\text{C1}} = R_0 + E_{O,\text{bnd}} \leq R_0$가
모델 구조에 의해 해석적으로 보장된다.


\chg{가설 H2(시그모이드 경계 보장)는 모든 조건(C0에서 C3까지)에서 $\DRTO$가 물리적 경계 내에 위치한다는 것이다.}
\[
  H_2:\;\; 0 < \DRTO < R_{\max} = 4R_0 = 24.0\text{h}
  \quad (\text{준수율 } 100\%)
\]

H2는 개별 시그모이드 바운딩 모델의 핵심 설계 목표인
물리적 경계 보장을 검증한다. (명제 2.1)에 의해
이 경계는 $\kappa$에 무관하게 항상 성립한다.



\chg{가설 H3(조건 순서 보존)는 네 조건이 다음의 세 가지의 순서 관계를 유지한다는 것이다. $H_{3\text{-}a}$는
$\DRTO_{\textbf{C1}} \leq \DRTO_{\textbf{C0}}$(준수율 $100\%$), $H_{3\text{-}b}$는
$\DRTO_{\textbf{C1}} \leq \DRTO_{\textbf{C2}}$(준수율 $\geq 99\%$), $H_{3\text{-}c}$는
$\DRTO_{\textbf{C2}} \leq \DRTO_{\textbf{C3}}$(조건부, 백분위 $\geq$ P85.8)을 의미한다.}

H3는 조건 간 $\DRTO$의 크기 순서가 이론적 예측과 일치하는지를 검증한다.
H3-a는 H1과 동치이며 모델 구조에 의해 보장된다.
H3-b는 불확실성의 도입이 관측성의 하향 효과를 상쇄하고 초과함을 의미하며,
$N(3,1)$의 음수 표본(약 $0.135\%$)에서 $E_{U,\text{bnd}} < 0$이 되는 경우를
허용하기 위해 준수율 기준을 100\%가 아닌 99\% 이상으로 설정한다.
H3-c는 파레토 분포 특성이 가우시안 대비 높은 $\DRTO$를 산출함을
검증하나, 이 효과는 P85.8 CDF 교차점 이상의 극단 영역에서만 발현되므로
조건부 가설로 설정한다.


%% -------------------------------------------------------------------
%% -------------------------------------------------------------------

\chg{가설 H4(회귀 설명력)는 교호 작용을 포함한 2차 회귀 모델의 결정계수가 모든 조건에서 $R^2_{(2\mathrm{q})} \geq 0.95$를 달성한다는 것이다.}
\[
  \jrev{H_4}:\;\; R^2_{(2\mathrm{q})} \geq 0.95
  \quad \forall\; \textbf{C1}, \textbf{C2}, \textbf{C3}
\]

\erf{H5}{H4}는 $\DRTO$ 모델의 비선형 응답 곡면이 2차 다항식으로 충분히
근사될 수 있는지를 검증한다. 2차 회귀는 1차항($k_O$, $O$, $U$),
2차항($k_O^2$, $O^2$, $U^2$), 교호작용항($k_O \cdot O$, $k_O \cdot U$,
$O \cdot U$)을 포함하는 완전 2차 다항 모델이다.
\chg{이러한 모델 선택의 이론적 근거는 \citet{box1951experimental}이 정립한 반응표면법에 있다. 반응표면법은 입력 변수와 응답 사이의 미지의 비선형 관계를 국소적으로 곡률까지 포착하는 완전 2차 다항식으로 근사하는 접근이 타당함을 보였으며, 본 연구의 응답 곡면 또한 시그모이드 바운딩에서 비롯한 매끄러운 곡률 구조를 가지므로 이 근사가 적합하다.}


\chg{가설 H5(교호 작용 유의성)는 교호작용항(interaction terms)이 회귀 모델의 설명력에 유의하게 기여한다는 것이다.}
\[
  \jrev{H_5}:\;\; \Delta R^2_{(\text{교호작용항 추가})} > 0,\quad F\text{-검정에서 } p < 0.05
\]

H5는 관측성과 불확실성이 독립적으로 작용하는 것이 아니라,
두 요소의 곱(교호 작용)이 $\DRTO$에 추가적인 설명력을 제공함을 검증한다.
특히 $k_O \cdot O_{\text{raw}}$ 교호작용항은 C1에서 $R^2$를 결정적으로 향상시킨다.


%% -------------------------------------------------------------------
%% -------------------------------------------------------------------

\chg{가설 H6(비중첩 대역 시드 독립성)은 비중첩 대역(non-overlapping band) 시드에서 파생된 입력 변수 $k_O$, $O_{\text{raw}}$, $U_{\text{raw}}$ 간 상관계수의 \m{절댓값}이 0.02 미만이라는 것이다.}
\[
  \jrev{H_6}:\;\; |r(X_i, X_j)| < 0.02 \quad \forall\; i \neq j
\]

H6은 시뮬레이션의 입력 변수들이 통계적으로 독립임을 검증한다.
마스터 시드 42로부터 비중첩 대역 규칙에 따라 파생된 시드 간
상관계수가 무시 가능한 수준임을 확인한다.


\chg{가설 H7(다중 시드 강건성)은 다중 시드($s = 0, 1, \ldots, 9{,}999$) 시뮬레이션에서 회귀 계수가 안정적이라는 것이다.}
\[
  \jrev{H_7}:\;\; \mathrm{CV}(\beta_{k_O}) < 0.05
\]

H7은 seed = 42에서 도출된 결과가 특정 시드에 의한 우연적 산물이 아님을
검증한다. 10{,}000개의 독립 시드 각각에서 $n = 10{,}000$ 시뮬레이션을
수행하여 총 $10^8$회의 $\DRTO$ 계산을 실시하고,
회귀계수의 분포를 산출한다.


%% -------------------------------------------------------------------
%% -------------------------------------------------------------------

\chg{가설 H8(이중채널 비선형 감쇠)은 C3 꼬리 영역($\geq \text{P85.8}$)에서 관측성 가중치 $k_O$의 회귀 계수 크기가 핵심 영역($< \text{P85.8}$) 대비 감소한다(감쇠비 $< 1$)는 것이다.}
\[
  \jrev{H_8}:\;\; \frac{|\beta_{k_O}^{\text{Tail}}|}{|\beta_{k_O}^{\text{Core}}|} < 1
\]

H8은 본 연구에서 새롭게 발견한 이중 채널 비선형 상호작용(dual-channel nonlinear
interaction) 메커니즘을 형식화한 가설이다.
``이중 채널''이란 관측성 채널($z_O$)과 불확실성 채널($z_U$)의 두 경로를 의미하며,
두 채널이 시그모이드 함수의 비선형성을 통해 결합될 때
극단 영역에서 시그모이드 포화(saturation)로 인해 관측성의 한계 효과가
오히려 감소하는 현상을 지칭한다.
감쇠비의 실증값과 P85.8 분할 회귀 분석 상세 내용은 4.10(이중채널 감쇠 메커니즘)에서 다루며,
이는 극단적 불확실성 하에서 관측성 투자의 한계 효과가 체감함을 의미한다.
이러한 비선형 감쇠 메커니즘은 조직의 위험 관리 전략에 대한
중요한 정책적 시사점을 제공하며 5.3(실무적 기여)에서 상세히 논의한다.


%% -------------------------------------------------------------------
%% -------------------------------------------------------------------

\chg{아홉째 가설 H9(전문가 정량적 검증)는 BCM/DR 전문가 설문조사에서 전체 영역 평균이 5점 리커트 척도 기준 4.0 이상이라는 것이다.}
\[
  \jrev{H_{9}}:\;\; \bar{X}_{\text{expert}} \geq 4.0 \;/\; 5.0
\]


\chg{열째 가설 H10(전문가 정성적 합의)은 전문가 심층 인터뷰에서 도출된 핵심 주제에 대한 동의율이 $80\%$ 이상이라는 것이다. 각 주제 $j$에 대한 동의 여부를 이진 변수 $a_j \in \{0, 1\}$로 정의하면 $a_j = 1$은 동의를, $a_j = 0$은 비동의를 가리킨다.}
\[
  \jrev{H_{10}}:\;\; A_{\text{agree}} \geq 0.80, \quad
  A_{\text{agree}} = \frac{1}{N_{\text{topic}}}\sum_{j=1}^{N_{\text{topic}}} a_j
\]
\chg{여기서 $A_{\text{agree}}$는 전체 $N_{\text{topic}}$개 주제 중 전문가가 동의한 주제의 비율을 의미한다.}

H9와 H10은 $\DRTO$ 모델의 실무적 타당성을 현장 전문가의 판단을 통해
삼각검증(triangulation)하기 위한 가설로서, \chg{정량적 방법과 정성적 방법을 결합하여 단일 방법의 한계를 보완하는 혼합연구방법의 삼각검증 원리~\citep{creswell2018designing}에 이론적 근거를 둔다.}
정량적 설문(H9)은 4개 평가 영역, \chg{곧 모델 타당성과 파라미터 적정성,
실무 적용성, 산업적용 및 $\eta$ 활용에 대한} 27개 문항으로 구성되며,
정성적 인터뷰(H10)는 반구조화 방식으로 수행한다.
\chg{H9와 H10의 검증 결과 상세내용은 4.11(전문가 패널 검증)에서 다룬다.}


%% -------------------------------------------------------------------
%% -------------------------------------------------------------------

\p{[}표~\ref{tab:hypotheses-summary}\p{]}에 \jrev{10개} 가설의 정의, 판정 기준, 관련 조건을 종합하였다.
\chg{[표 2-16]에서 H7의 $10^8$은 다중 시드 시뮬레이션의 총 횟수를 가리킨다.}

\begin{table}[htbp]
\centering
\caption{\jrev{10개 연구가설 종합}}
\label{tab:hypotheses-summary}
\small
\begin{tabularx}{\textwidth}{c c X l c}
\toprule
\textbf{가설} & \textbf{범주} & \textbf{핵심 내용} & \textbf{판정 기준} & \textbf{조건} \\
\midrule
H1 & 시뮬레이션 & $\DRTO_{\text{C1}} \leq R_0$ (관측성 효과) & 준수율 100\% & C0, C1 \\
H2 & 시뮬레이션 & $0 < \DRTO < R_{\max} = 24.0$h (경계 보장) & 준수율 100\% & C0--C3 \\
H3 & 시뮬레이션 & C0--C3 순서 보존 (a/b/c) & 준수율 기준 & C0--C3 \\
\midrule
H4 & 회귀 & $R^2_{(2\mathrm{q})} \geq 0.95$ & C1/C2/C3 전체 & C1--C3 \\
H5 & 회귀 & 교호작용항 \m{유의성($\Delta R^2$)} & 기여도 검정 & C1--C3 \\
\midrule
H6 & 강건성 & 시드 간 $|r| < 0.02$ (독립성) & 상관계수 검정 & 전체 \\
H7 & 강건성 & 다중 시드 계수 \m{안정성($10^8$회)} & 계수 CV 기준 & 전체 \\
\midrule
H8 & 이중채널 & Tail/Core $k_O$ 감쇠비 $< 1$ & P85.8 분할 회귀 & C3 \\
\midrule
H9 & 전문가 & 설문 평균 $\geq 4.0/5.0$ & $t$-검정 & 전체 \\
H10 & 전문가 & $A_{\text{agree}} \geq 0.80$ (15개 하위 주제) & 합의 기준 & 전체 \\
\bottomrule
\end{tabularx}
\end{table}
          % 2.11 \jrev{10개 연구가설}
        % 제2장 이론적 배경
% 제3장 연구 방법
% 1차 심사 반영: 구 2.12 연구 설계 + 구 3.1 시뮬레이션 설계 를 통합

\chapter{\vrev{연구 방법}}
\label{ch:method}

\vrev{%
본 연구의 연구 설계와 시뮬레이션 실험 방법을 기술하기 위하여 먼저
관측성과 불확실성의 조합으로 정의되는 4단계 조건 체계(C0--C3)와
이에 대응하는 \jrev{10개 가설}의 검증 설계를 제시한다.
이어서 $n = 10{,}000$ 몬테카를로 시뮬레이션의 실행 절차와
분포 매개변수, 시드 배정 규칙, 통계량 산출 방법을 제시한다.
}

\chg{본 연구는 모델 정립에서 산업 범용성 검증에 이르기까지 일곱 단계의 절차를
거치며, \figref{fig:comprehensive_framework}\refpo{fig:comprehensive_framework}{은}{는}
이 전체 절차를 한눈에 보여준다. 본 연구는 먼저 관측성과 불확실성을 시그모이드
바운딩으로 결합하는 동적 복구목표시간(DRTO) 수식 모델 2.7(DRTO 모델 구조)를 정립하는 데에서
출발한다. 다음으로 정적 기준에서 관측성과 불확실성을 단계적으로 도입하는 네 개
운영 조건(C0--C3)을 설계하여, 각 요인의 기여와 결합 효과를 분리하여 관찰할 수
있는 실험 틀을 마련한다. 본격적인 가설 검증에 앞서 기반 검증(L0)에 해당하는
Golden Compare를 수행하여, 수식의 코드 구현이 단일 기준 정의(SSOT)와 정확히
일치함을 결정론적으로 확인한다. 이 토대 위에서 검증은 네 층위로 이어진다.
$n=10{,}000$ 몬테카를로 시뮬레이션으로 관측성 효과와 경계 보장 등 모델의 기본
작동을 확률적으로 확인하고(L1), 2차 다항 회귀와 분할 회귀로 변수 간 설명력과
이중채널 비선형 상호작용을 규명하며(L2와 L4), 다중 시드 분석으로 결과가 특정
난수에 의존하지 않고 재현됨을 보인다(L3). 이어 BCM 및 재해복구 분야
전문가 평가로 모델의 이론적 타당성과 실무 적용 가능성을 삼각검증하고(L5),
마지막으로 여섯 개 산업 시나리오에 적용하여 모델이 특정 환경에 국한되지 않는
범용성을 갖는지 확인한다. 이처럼 각 단계는 앞 단계의 결과를 전제로 다음 단계로
이어지며, 하위 검증의 통과가 상위 검증의 신뢰성을 뒷받침하는 계층적 구조를
이룬다.}

% 제3장 도입 — 연구 전체 수행 절차 개관 그림 (1장에서 이동, 결과수치 제거됨)
\begin{figure}[htbp]
  \centering
  \resizebox{\textwidth}{!}{%
  \begin{tikzpicture}[
    >=Stealth,
    stepnum/.style={circle, draw=black, fill=white, font=\scriptsize\bfseries,
                    minimum size=5.5mm, inner sep=0pt, line width=0.7pt},
    mainbox/.style={draw=black, rounded corners=5pt, line width=0.7pt,
                    align=center, font=\scriptsize, minimum height=8mm,
                    inner sep=3pt},
    subbox/.style={draw=black, rounded corners=2pt, line width=0.4pt,
                   align=center, font=\tiny, minimum height=7mm,
                   minimum width=22mm, inner sep=2pt},
    condbox/.style={draw=black, rounded corners=2pt, line width=0.5pt,
                    align=center, font=\tiny, minimum width=16mm,
                    minimum height=8mm, inner sep=2pt},
    mainarr/.style={->, line width=1.0pt, black},
    subarr/.style={->, line width=0.5pt, black},
  ]

  %% 행 좌표 (상→하)
  \def\rowA{14.0}   % ① 수식 모델
  \def\rowB{11.5}   % ② 운영 조건       (①-②: 2.5)
  \def\rowC{7.5}    % ③ 시뮬레이션      (②-③: 4.0, ②-L0 화살표 공간 확보)
  \def\rowD{3.8}    % ④ 회귀·교차분석   (③-④: 3.7, H3 하위 2단 + 박스 간섭 회피)
  \def\rowE{1.6}    % ⑤ 강건성 검증     (④-⑤: 2.2)
  \def\rowF{-0.6}   % ⑥ 전문가 평가     (⑤-⑥: 2.2)
  \def\rowG{-2.8}   % ⑦ 산업 범용성     (⑥-⑦: 2.2)

  %% ====================================================================
  %% ① 수식 모델
  %% ====================================================================
  \node[stepnum] (n1) at (-1.3, \rowA) {1};
  \node[mainbox, fill=blue!8, minimum width=115mm] (m1) at (5.0, \rowA)
    {\textbf{수식 모델 정립}\;
    {\tiny $\DRTO = R_0 + E_{O,\mathrm{bnd}} + E_{U,\mathrm{bnd}}$,\;
    $S(z) = 2/(1\!+\!e^{-z})\!-\!1$,\; $\kappa = 1.2$}};

  \node[subbox, fill=blue!5] (m1a) at (1.3, \rowA-1.1)
    {관측성 채널\\$E_{O,\mathrm{bnd}} = -R_0 \!\cdot\! S(z_O)$};
  \node[subbox, fill=blue!5] (m1b) at (5.0, \rowA-1.1)
    {불확실성 채널\\$E_{U,\mathrm{bnd}} = (R_{\max}\!-\!R_0)\!\cdot\! S(z_U)$};
  \node[subbox, fill=blue!5] (m1c) at (8.7, \rowA-1.1)
    {경계 보장\\$0 < \DRTO < R_{\max}$};

  \draw[subarr] (m1.south -| m1a) -- (m1a.north);
  \draw[subarr] (m1.south -| m1b) -- (m1b.north);
  \draw[subarr] (m1.south -| m1c) -- (m1c.north);

  %% ====================================================================
  %% ② 운영 조건
  %% ====================================================================
  \node[stepnum] (n2) at (-1.3, \rowB+0.5) {2};
  \node[mainbox, fill=green!8, minimum width=115mm] (m2) at (5.0, \rowB+0.5)
    {\textbf{4단계 운영 조건 설계 (C0--C3)}};

  \node[condbox, fill=gray!15] (c0) at (0.8, \rowB-0.7)
    {\textbf{C0}\\정적 기준};
  \node[condbox, fill=blue!10] (c1) at (3.6, \rowB-0.7)
    {\textbf{C1}\\관측성};
  \node[condbox, fill=green!10] (c2) at (6.4, \rowB-0.7)
    {\textbf{C2}\\Gaussian};
  \node[condbox, fill=red!10] (c3) at (9.2, \rowB-0.7)
    {\textbf{C3}\\Pareto};

  \draw[subarr] (m2.south -| c0) -- (c0.north);
  \draw[subarr] (m2.south -| c1) -- (c1.north);
  \draw[subarr] (m2.south -| c2) -- (c2.north);
  \draw[subarr] (m2.south -| c3) -- (c3.north);

  \draw[subarr] (c0) -- (c1) node[midway, above, font=\tiny] {$+O$};
  \draw[subarr] (c1) -- (c2) node[midway, above, font=\tiny] {$+U^{(G)}$};
  \draw[subarr] (c1.south) -- ++(0,-0.7) -| (c3.south)
    node[pos=0.25, below, font=\tiny] {$+U^{(P)}$};

  %% ====================================================================
  %% L0: Golden Compare (② → ③ 사이)
  %% ====================================================================
  \node[stepnum, fill=black!8] (n0) at (-1.3, {(\rowB-1.5+\rowC)/2})
    {\tiny L0};
  \node[mainbox, fill=black!6, minimum width=115mm] (m0)
    at (5.0, {(\rowB-1.5+\rowC)/2})
    {\textbf{L0: 기반 검증 (Golden Compare)}\;
    {\tiny 47개 포인트 --- 순서 제약(H3), 물리적 경계(H2),
    $z_O$ 동일성, P99 비율 검증}};

  %% ====================================================================
  %% ③ 시뮬레이션 [L1: H1-H3] — H1~H3 각각 독립 박스, H3 하위 a/b/c
  %% ====================================================================
  \node[stepnum] (n3) at (-1.3, \rowC) {3};
  \node[mainbox, fill=yellow!10, minimum width=115mm] (m3) at (5.0, \rowC)
    {\textbf{몬테카를로 시뮬레이션}\; ($n\!=\!10{,}000 \times 4$)\;
    {\tiny [L1: H1--H3]}};

  \node[subbox, fill=yellow!6, minimum width=22mm] (m3h1) at (1.5, \rowC-1.1)
    {H1 관측성 효과};
  \node[subbox, fill=yellow!6, minimum width=22mm] (m3h2) at (5.0, \rowC-1.1)
    {H2 경계 보장};
  \node[subbox, fill=yellow!6, minimum width=28mm] (m3h4) at (8.7, \rowC-1.1)
    {H3 조건 순서};

  \draw[subarr] (m3.south -| m3h1) -- (m3h1.north);
  \draw[subarr] (m3.south -| m3h2) -- (m3h2.north);
  \draw[subarr] (m3.south -| m3h4) -- (m3h4.north);

  %% H3 하위 박스 (a, b, c)
  \node[subbox, fill=yellow!4, minimum width=12mm, minimum height=6mm]
    (h4a) at (7.4, \rowC-2.0)
    {a: $C1\!\leq\!C0$};
  \node[subbox, fill=yellow!4, minimum width=12mm, minimum height=6mm]
    (h4b) at (8.7, \rowC-2.0)
    {b: $C1\!\leq\!C2$};
  \node[subbox, fill=yellow!4, minimum width=12mm, minimum height=6mm]
    (h4c) at (10.0, \rowC-2.0)
    {c: $C2\!\leq\!C3$};

  \draw[subarr] (m3h4.south -| h4a) -- (h4a.north);
  \draw[subarr] (m3h4.south) -- (h4b.north);
  \draw[subarr] (m3h4.south -| h4c) -- (h4c.north);

  %% ====================================================================
  %% ④ 회귀·교차분석 [L2: H4-H5, L4: H8]
  %% ====================================================================
  \node[stepnum] (n4) at (-1.3, \rowD+0.5) {4};
  \node[mainbox, fill=orange!10, minimum width=115mm] (m4) at (5.0, \rowD+0.5)
    {\textbf{회귀분석 및 교차분석}\;
    {\tiny [L2: H4--H5, L4: H8]}};

  \node[subbox, fill=orange!6] (m4a) at (1.3, \rowD-0.6)
    {H4 회귀 설명력};
  \node[subbox, fill=orange!6] (m4b) at (5.0, \rowD-0.6)
    {H5 교호작용};
  \node[subbox, fill=orange!6] (m4c) at (8.7, \rowD-0.6)
    {H8 이중채널 감쇠};

  \draw[subarr] (m4.south -| m4a) -- (m4a.north);
  \draw[subarr] (m4.south -| m4b) -- (m4b.north);
  \draw[subarr] (m4.south -| m4c) -- (m4c.north);

  %% ====================================================================
  %% ⑤ 강건성 검증 [L3: H6-H7]
  %% ====================================================================
  \node[stepnum] (n5) at (-1.3, \rowE+0.5) {5};
  \node[mainbox, fill=cyan!8, minimum width=115mm] (m5) at (5.0, \rowE+0.5)
    {\textbf{강건성 검증}\;
    {\tiny [L3: H6--H7]}};

  \node[subbox, fill=cyan!5] (m5a) at (1.3, \rowE-0.6)
    {H6 시드 독립성};
  \node[subbox, fill=cyan!5] (m5b) at (5.0, \rowE-0.6)
    {H7 다중 시드 강건성};
  \node[subbox, fill=cyan!5] (m5c) at (8.7, \rowE-0.6)
    {Bonferroni 보정};

  \draw[subarr] (m5.south -| m5a) -- (m5a.north);
  \draw[subarr] (m5.south -| m5b) -- (m5b.north);
  \draw[subarr] (m5.south -| m5c) -- (m5c.north);

  %% ====================================================================
  %% ⑥ 전문가 평가 + 가설 종합 [L5: H9-H10]
  %% ====================================================================
  \node[stepnum] (n6) at (-1.3, \rowF+0.5) {6};
  \node[mainbox, fill=violet!8, minimum width=115mm] (m6) at (5.0, \rowF+0.5)
    {\textbf{다층적 타당성 검증}\;
    {\tiny [L5: H9--H10]}};

  \node[subbox, fill=violet!5] (m6a) at (1.3, \rowF-0.6)
    {H9 전문가 정량 평가};
  \node[subbox, fill=violet!5] (m6b) at (5.0, \rowF-0.6)
    {H10 전문가 정성 평가};
  \node[subbox, fill=violet!5] (m6c) at (8.7, \rowF-0.6)
    {가설 종합 판정};

  \draw[subarr] (m6.south -| m6a) -- (m6a.north);
  \draw[subarr] (m6.south -| m6b) -- (m6b.north);
  \draw[subarr] (m6.south -| m6c) -- (m6c.north);

  %% ====================================================================
  %% ⑦ 산업 범용성 검증
  %% ====================================================================
  \node[stepnum] (n7) at (-1.3, \rowG+0.5) {7};
  \node[mainbox, fill=teal!8, minimum width=115mm] (m7) at (5.0, \rowG+0.5)
    {\textbf{6개 산업 시나리오 범용성 검증}};

  \node[subbox, fill=teal!4, minimum width=16mm] (m7a) at (0.0, \rowG-0.7)
    {제조업};
  \node[subbox, fill=teal!4, minimum width=16mm] (m7b) at (2.0, \rowG-0.7)
    {금융};
  \node[subbox, fill=teal!4, minimum width=16mm] (m7c) at (4.0, \rowG-0.7)
    {의료};
  \node[subbox, fill=teal!4, minimum width=16mm] (m7d) at (6.0, \rowG-0.7)
    {IT서비스};
  \node[subbox, fill=teal!4, minimum width=16mm] (m7e) at (8.0, \rowG-0.7)
    {물류};
  \node[subbox, fill=teal!4, minimum width=16mm] (m7f) at (10.0, \rowG-0.7)
    {소매};

  \draw[subarr] (m7.south -| m7a) -- (m7a.north);
  \draw[subarr] (m7.south -| m7b) -- (m7b.north);
  \draw[subarr] (m7.south -| m7c) -- (m7c.north);
  \draw[subarr] (m7.south -| m7d) -- (m7d.north);
  \draw[subarr] (m7.south -| m7e) -- (m7e.north);
  \draw[subarr] (m7.south -| m7f) -- (m7f.north);

  %% ====================================================================
  %% 왼쪽: 단계 진행 화살표 (흐름 명시)
  %% ====================================================================
  \foreach \a/\b in {n1/n2, n2/n0, n0/n3, n3/n4, n4/n5, n5/n6, n6/n7} {
    \draw[->, black, line width=0.7pt] (\a.south) -- (\b.north);
  }

  %% ====================================================================
  %% 오른쪽: 중괄호 (왼쪽 열림 = 본문 쪽) + 가설 매핑
  %% ====================================================================
  \draw[decorate, decoration={brace, amplitude=4pt},
        line width=0.5pt, black]
    (11.0, \rowA+0.4) -- (11.0, \rowB-1.0)
    node[midway, right=5pt, font=\tiny, text=black, align=left]
    {이론적\\기반};

  \draw[decorate, decoration={brace, amplitude=4pt},
        line width=0.5pt, black]
    (11.0, \rowC+0.35) -- (11.0, \rowC-2.35)
    node[midway, right=5pt, font=\tiny, text=black, align=left]
    {L1\\시뮬레이션};

  \draw[decorate, decoration={brace, amplitude=4pt},
        line width=0.5pt, black]
    (11.0, \rowD+0.9) -- (11.0, \rowD-1.0)
    node[midway, right=5pt, font=\tiny, text=black, align=left]
    {L2--L4\\회귀$\cdot$이중채널};

  \draw[decorate, decoration={brace, amplitude=4pt},
        line width=0.5pt, black]
    (11.0, \rowE+0.9) -- (11.0, \rowE-1.0)
    node[midway, right=5pt, font=\tiny, text=black, align=left]
    {L3\\강건성};

  \draw[decorate, decoration={brace, amplitude=4pt},
        line width=0.5pt, black]
    (11.0, \rowF+0.9) -- (11.0, \rowF-1.0)
    node[midway, right=5pt, font=\tiny, text=black, align=left]
    {L5\\전문가};

  \draw[decorate, decoration={brace, amplitude=4pt},
        line width=0.5pt, black]
    (11.0, \rowG+0.9) -- (11.0, \rowG-1.0)
    node[midway, right=5pt, font=\tiny, text=black, align=left]
    {적용};

  %% 바운딩 박스 확장 (우측 라벨 잘림 방지, 하단 확장)
  \path (13.5, \rowG-1.5);

  \end{tikzpicture}%
  }% end resizebox
  \caption{\jrev{본 연구의 전체 수행 절차 개관}}
  \label{fig:comprehensive_framework}
\end{figure}
    % 연구 전체 수행 절차 개관 그림 (1장에서 이동)



%% ===================================================================
\section{\vrev{연구 설계}}
\label{sec:research-design}
%% ===================================================================

%% -------------------------------------------------------------------
\subsection{\vrev{6단계 다층 검증 프레임워크}}
\label{subsec:validation-framework}
%% -------------------------------------------------------------------

\jrev{10개 가설}을 체계적으로 검증하기 위해 6단계 다층 검증(six-stage multi-layer
validation) 프레임워크를 설계한다. 각 단계는 서로 다른 검증 관점을 제공하며,
하위 단계의 결과가 상위 단계의 전제 조건으로 작용하는
계층적 의존 구조를 갖는다~\citep{oberkampf2010verification}.

\chg{프레임워크는 결정론적 기반 검증에서 출발하여 통계적 검증과 전문가 평가로 단계적으로 심화한다. 먼저 기반 검증(L0, Foundation Verification)은 Golden Compare로서, 7개 산업 시나리오의 대표 매개변수를 $\DRTO$ 수식에 직접 대입하여 해석적 기대값과 코드 출력을 비교하는 결정론적 검증이다. 순서 제약, 물리적 경계, $z_O$ 동일성, P99 비율 일관성의 네 범주에 걸친 총 47개 검증 포인트로 구성되며, 경계 보장(H2)과 순서 보존(\jrev{H3})을 결정론적으로 사전 검증한다. 47개 포인트 전체에서 PASS를 달성해야만 이후 L1에서 L5까지의 가설 검증이 의존하는 수식 구현이 단일 기준 정의(SSOT)와 정합함이 보장된다(결과는 \subsecref{subsec:golden_compare} 참조, 7개 시나리오의 전체 계산과 47개 점별 결과는 공개 재현 패키지에 수록).}

\chg{이 기반 위에서 시뮬레이션 검증(L1)은 $n = 10{,}000$ 몬테카를로 시뮬레이션으로 관측성 효과(H1), 경계 보장(H2), 조건 순서(\jrev{H3})를 확인하여 모델의 기본 행동을 검증한다. 이때 H2와 \jrev{H3}는 L0에서 결정론적으로 사전 확인된 뒤 본 단계의 확률적 표본으로 재검증되며, 곡률 매개변수 $\kappa = 1.2$의 최적성은 \citet{lee2026drto}에서 확정된 전제로서 본 연구의 가설 검증 대상이 아니다. 이어 회귀 검증(L2)은 2차 다항 회귀 모델의 설명력(\jrev{H4})과 교호작용 유의성(\jrev{H5})을 통해 응답 곡면의 구조적 특성을 규명하고, 강건성 검증(L3)은 시드 독립성(\jrev{H6})과 다중 시드 강건성(\jrev{H7})을 통해 결과의 재현성과 대표성을 확인하는데, 표본 $n = 10{,}000$과 시드 $10{,}000$개를 곱한 총 $10^8$회 계산으로 통계적 안정성을 보장한다. 다음으로 이중채널 검증(L4)은 P85.8 기준 분할 회귀로 이중채널 비선형 감쇠 효과의 존재와 크기를 검증하며, L1에서 L3까지의 결과를 종합하는 핵심 발견으로서 실무적 정책 시사점의 근거를 제공한다. 마지막으로 전문가 평가(L5)는 BCM 및 재해복구 전문가 5명의 정량적 설문(\jrev{H9})과 정성적 인터뷰(\jrev{H10})를 통해 모델의 이론적 타당성과 실무 적용 가능성을 삼각검증한다.}

\chg{특히 최상위 단계인 L5의 전문가 평가는 본 연구가 삼각검증(triangulation)을 채택한 지점이다. 시뮬레이션과 회귀(L1에서 L4까지)가 제공하는 정량적 증거만으로는 모델이 실제 업무연속성 현장에서 타당한지를 온전히 확증하기 어렵다. 단일 방법에 내재한 한계를 보완하기 위해, 정량적 분석 결과를 전문가의 정성적 판단과 교차 대조하여 결론이 서로 수렴하는지를 확인하는 삼각검증~\citep{creswell2018designing}을 결과 단계가 아니라 설계 단계에서부터 검증 체계의 최상위 층위로 포함한다. 이로써 정량적으로 도출된 모델의 효과가 현장 전문가의 실무적 판단과 일치할 때 비로소 타당성이 확증되는 구조를 갖춘다.}

\p{[}그림~\vref{fig:validation-layers}\p{]}는 6단계 검증 프레임워크의
구조와 계층 간 의존 관계를 도식화한 것이다.

\begin{figure}[htbp]
\centering
\resizebox{0.95\textwidth}{!}{%
\begin{tikzpicture}[
  pyr/.style={
    draw, rounded corners=5pt, minimum height=1.3cm,
    align=center, font=\small, line width=0.6pt,
    drop shadow={shadow xshift=0.8mm, shadow yshift=-0.8mm, opacity=0.15}
  },
  arr/.style={-{Stealth[length=2.5mm]}, thick, black},
  lnum/.style={circle, draw, fill=white, font=\small\bfseries, inner sep=2pt,
    minimum size=7mm, line width=0.6pt},
  rlbl/.style={font=\scriptsize, text=black, align=left, anchor=west},
]

% 피라미드 계층 (아래→위, 넓은→좁은)
\node[pyr, fill=black!8, minimum width=15.5cm] (L0) at (0, -1.8) {
  \textbf{L0: 기반 검증 (Golden Compare)} \quad
  47개 포인트 \;\textbar\;
  순서 제약 \;\textbar\;
  물리적 경계 \;\textbar\;
  $z_O$ 동일성 \;\textbar\;
  P99 비율
};
\node[pyr, fill=green!15, minimum width=14cm] (L1) at (0, 0) {
  \textbf{L1: 시뮬레이션 검증} \quad
  H1 (관측성 효과) \;\textbar\;
  H2 (경계 보장) \;\textbar\;
  H3 (순서 보존)
};
\node[pyr, fill=violet!10, minimum width=12cm] (L2) at (0, 1.8) {
  \textbf{L2: 회귀 검증} \quad
  H4 ($R^2_{(2\mathrm{q})} \geq 0.95$) \;\textbar\;
  H5 (교호작용)
};
\node[pyr, fill=blue!10, minimum width=10cm] (L3) at (0, 3.6) {
  \textbf{L3: 강건성 검증} \quad
  H6 (시드 독립성, $|r| < 0.02$) \;\textbar\;
  H7 (다중 시드, $10^8$회)
};
\node[pyr, fill=orange!15, minimum width=8cm] (L4) at (0, 5.4) {
  \textbf{L4: 이중채널}\\
  H8 (감쇠비 $< 1$)
};
\node[pyr, fill=yellow!10, minimum width=6cm] (L5) at (0, 7.2) {
  \textbf{L5: 전문가 평가}\\
  H9 ($\bar{X} \geq 4.0$) \;\textbar\;
  H10 (동의율 $\geq 80\%$)
};

% 계층 번호 (왼쪽 원형 배지)
\node[lnum, text=black] at ([xshift=-9.1cm]L0.center) {L0};
\node[lnum, text=black] at ([xshift=-8.4cm]L1.center) {L1};
\node[lnum, text=black] at ([xshift=-6.8cm]L2.center) {L2};
\node[lnum, text=black] at ([xshift=-6.4cm]L3.center) {L3};
\node[lnum, text=black] at ([xshift=-4.8cm]L4.center) {L4};
\node[lnum, text=black] at ([xshift=-3.8cm]L5.center) {L5};

% 우측 핵심 산출물
\node[rlbl] at ([xshift=8.5cm]L0.center) {SSOT 구현 정합성};
\node[rlbl] at ([xshift=7.8cm]L1.center) {$n=10{,}000$ 시뮬레이션};
\node[rlbl] at ([xshift=6.8cm]L2.center) {응답 곡면 규명};
\node[rlbl] at ([xshift=5.8cm]L3.center) {재현성 $10^8$회};
\node[rlbl] at ([xshift=4.8cm]L4.center) {이중채널 효과 검증};
\node[rlbl] at ([xshift=3.8cm]L5.center) {전문가 삼각검증};

% 계층 간 화살표
\draw[arr] (L0.north) -- (L1.south);
\draw[arr] (L1.north) -- (L2.south);
\draw[arr] (L2.north) -- (L3.south);
\draw[arr] (L3.north) -- (L4.south);
\draw[arr] (L4.north) -- (L5.south);

% 좌측 방향 표시
\node[font=\footnotesize, text=black, rotate=90] at (-9.8, 2.7) {검증 심화 방향};
\draw[-{Stealth[length=2mm]}, thick, black] (-9.4, -2.3) -- (-9.4, 7.7);

\end{tikzpicture}%
}% end resizebox
\caption{6단계 다층 검증 프레임워크}
\label{fig:validation-layers}
\end{figure}

\p{[}표~\ref{tab:six-stage}\p{]}은 6단계 검증 프레임워크의 각 단계와 대응 가설,
검증 목적을 요약한 것이다. 여기서 ``L''은 단계(Level)의 약자로, L0의 결정론적
구현 정합성 검증에서 출발하여 L5의 전문가 삼각검증에 이르는 누적적 위계를 구성하며,
각 단계는 이전 단계의 결과를 전제로 한다. 이 순서는 본 연구의 실행 흐름과 일치한다.

\begin{table}[H]
\centering
\caption{6단계 다층 검증 프레임워크}
\label{tab:six-stage}
\small
\begin{tabular}{c l c l}
\toprule
\textbf{단계} & \textbf{명칭} & \textbf{가설} & \textbf{목적} \\
\midrule
L0 & 기반 검증       & (H2, \jrev{H3} 사전) & Golden Compare 47개 포인트: SSOT 구현 정합성 \\
L1 & 시뮬레이션 검증 & H1--\jrev{H3}   & 모델 기본 행동 확인 \\
L2 & 회귀 검증       & \jrev{H4}--\jrev{H5}   & 설명력과 교호작용 검증 \\
L3 & 강건성 검증     & \jrev{H6}--\jrev{H7}   & 재현성과 안정성 보장 \\
L4 & 이중채널 검증   & \jrev{H8}       & 핵심 발견 규명 \\
L5 & 전문가 평가     & \jrev{H9}--\jrev{H10} & 실무 타당성 삼각검증 \\
\bottomrule
\end{tabular}
\end{table}

이 위계에서 L0의 기반 검증과 L1 이후의 통계적 검증은 성격이 근본적으로 다르다.
L0의 결정론적 검증(deterministic verification)은 고정된 입력 매개변수에서 모델 수식의
해석적 출력과 코드 실행 결과를 직접 비교하여 정확 일치 여부를 PASS/FAIL로 판정하는
방식으로, 확률적 표본 추출이나 통계적 추론을 사용하지 않으며 동일 입력에서 항상 동일
출력이 보장된다. 반면 L1--L4의 통계적 검증은 확률 분포에서 $n = 10{,}000$ 표본을
추출하여 평균, 분산, 백분위, 회귀 계수 등 모델 거동의 통계적 성질을 분석한다.
두 방식의 차이를 \p{[}표~\ref{tab:verification-types}\p{]}에 정리하였다.

\begin{table}[H]
\centering
\caption{결정론적 검증과 통계적 검증의 비교}
\label{tab:verification-types}
\small
\begin{tabular}{l p{5.5cm} p{5.5cm}}
\toprule
\textbf{항목} & \textbf{결정론적 검증 (L0)} & \textbf{통계적 검증 (L1--L4)} \\
\midrule
입력 & 고정 매개변수 (시나리오) & 확률 분포 ($n = 10{,}000$ 표본) \\
출력 비교 & 해석적 기대값 = 코드 출력 (정확 일치) & 분포·평균·검정 통계량 \\
결론 형식 & PASS / FAIL (이진) & 가설 채택/기각, 효과크기 \\
변동성 & 없음 (재실행 시 동일) & 시드·표본에 따른 변동 \\
검증 대상 & 수식 구현 정합성 (코드와 이론) & 모델 거동의 통계적 성질 \\
대표 사례 & Golden Compare 47개 포인트 & Cohen's $d$, $R^2$, CVaR \\
\bottomrule
\end{tabular}
\end{table}

두 방식은 위계적 사전 조건 관계를 갖는다. 결정론적 검증(L0)이 PASS되어야 수식이
코드로 정확히 구현됨이 보장되며, 비로소 후속 통계적 검증(L1--L4)의 결과가 의미를
갖는다. 만약 L0에서 단 하나의 검증 포인트라도 FAIL이면 코드 구현 자체에 오류가
있다는 의미이므로 이후의 모든 분석 결과를 신뢰할 수 없게 된다~\citep{oberkampf2010verification}.
이러한 ``결정론적 사전 검증에서 통계적 본 검증으로''의 2단계 구조는 계산과학 분야의
\chg{검증과 확인}(Verification and Validation) 표준에 부합한다.


%% -------------------------------------------------------------------
\chg{이상의 검증 프레임워크를 실제로 수행하기 위한 시뮬레이션의 구체적 구성, 곧 표본 크기와 분포 매개변수, 난수 생성 방법, 변수별 시드 배정 규칙과 그 독립성 검증은 다음 \secref{sec:exp_design}에서 통합적으로 기술한다.}
%% [C3.6/일원화] 구 3.1.2 시드 설계(itemize·그림 fig:seed-bands)는 3.2.3 시드 체계(표 tab:seed_system)·3.2.7 독립성 검증과 중복 → 삭제, 3.2로 일원화(사용자 지시: 분량 무관 중복 삭제).


%% -------------------------------------------------------------------
%% [C3.6/일원화] 구 3.1.3 시뮬레이션 매개변수(itemize·μ=3.0/Lomax 도출)는 3.2.1 tab:sim_parameters(설정근거 포함)·
%% 3.2.5 난수표 각주·3.2.6 분포 설정과 중복 → 삭제, 3.2로 일원화(사용자 지시: 분량 무관 중복 삭제).


%% [C3.3] 구 3.1.4 ``장 요약'' 절 삭제 — 심사 지적(한윤영12·2-B): 절 단위 요약 불필요,
%% 동시에 본 절에 선등장하던 4장 결과수치(C1~C3 평균 효율비·Cohen's d·단축률)도 함께 제거(C4.5 정합).
   % 3.1 연구 설계 (구 2.12)
\section{\vrev{시뮬레이션 설계}}
\label{sec:exp_design}
%% ===================================================================

%% -------------------------------------------------------------------
\subsection{\vrev{몬테카를로 시뮬레이션 구성}}
\label{subsec:mc_config}
%% -------------------------------------------------------------------

본 연구의 시뮬레이션은 조건당 $n = 10{,}000$개의 독립 표본을 생성하는
몬테카를로(Monte Carlo) 방법~\citep{rubinstein2016simulation}에 기반한다.
시뮬레이션의 핵심 매개변수는 다음과 같이
기본 복구시간 $R_0 = 6.0$h, 최대 복구시간 $R_{\max} = 24.0$h($= 4R_0$, MTPD 기반 상한),
시그모이드 곡률 $\kappa = 1.2$이다.
각 표본에 대해 입력 변수를 독립적으로 생성하고,
식~(2.10)에 따라 $\DRTO$ 값과
효율성 비율 $\eta = \DRTO / R_0$를 산출한다.

\tabref{tab:sim_parameters}\refpo{tab:sim_parameters}{은}{는} 시뮬레이션의 전체 매개변수를 요약한 것이며, 각 모델 변수의 전체 정의는 2.7.3(변수 및 매개변수 정의)의 [표 2-11]에 정리되어 있다.

\begin{table}[htbp]
  \centering
  \caption{시뮬레이션 매개변수 요약}
  \label{tab:sim_parameters}
  \small
  \begin{tabularx}{\textwidth}{l c c X}
    \toprule
    \textbf{매개변수} & \textbf{기호} & \textbf{값/범위} & \textbf{설정 근거} \\
    \midrule
    기준선 RTO & $R_0$ & 6.0\,h & \nrev{ISO 22301:2019} BIA 기반 표준 설정값 \\
    DRTO 상한 & $R_{\max}$ & 24.0\,h & MTPD 기반 상한 ($= 4R_0$) \\
    곡률 매개변수 & $\kappa$ & 1.2 & 다기준 복합 점수 최대화 기준 최적값 \\
    표본 크기 & $n$ & 10{,}000 & 조건당, CVaR$_{99}$ 안정 추정 기준 \\
    관측성 원시 점수 & $O_{\text{raw}}$ & $\mathrm{Uniform}(0,1)$ & 연속 균등 분포 \\
    관측성 가중치 & $k_O$ & $\mathrm{Uniform}(0,1)$ & 연속 균등 분포 \\
    불확실성 (C2) & $U_{\text{raw}}^{(G)}$ & $N(3,\;1)$ & 경꼬리, $E[U] = 3.0$ \\
    불확실성 (C3) & $U_{\text{raw}}^{(P)}$ & $6 \cdot \mathrm{Pareto}(\alpha\!=\!3)$ & 파레토(중꼬리), $E[U] = 3.0$ \\
    \bottomrule
  \end{tabularx}
\end{table}


\chg{$n = 10{,}000$의 표본 크기는} 다음의 세 가지 기준에 의해 결정되었다.
첫째, 평균 $\bar{\eta}$의 표준오차가 $\mathrm{SE} = \mathrm{SD}/\sqrt{n}$으로서,
C3의 $\eta$ 표준편차 $\mathrm{SD}(\eta_{\text{C3}}) = 0.791$에서도 $\mathrm{SE} = 0.008$에 불과하여
소수점 둘째 자리까지의 정밀도가 보장된다.
둘째, CVaR$_{99}$의 산출에 상위 $1\%$에 해당하는 100개 이상의 표본이
필요하며, $n = 10{,}000$은 이 요구를 충족한다.
셋째, 부분 표본 분석에서 $n = 5{,}000$ 이상일 때 핵심 통계량이
소수점 셋째 자리까지 수렴함을 사전 확인하였으며 4.6(수렴 진단 및 재현성)에서 그 상세 결과를 알 수 있다.


%% -------------------------------------------------------------------
\subsection{\vrev{난수 생성 방법}}
\label{subsec:rng_method}
%% -------------------------------------------------------------------

본 시뮬레이션의 난수 생성은 [표 3-3]의 시뮬레이션 매개변수 요약 기준에 따라 설계되었으며, 난수 생성기의 상세 규격은 [표 3-4]에 정리하였다.

\begin{table}[htbp]
  \centering
  \caption{난수 생성 방법 규격}
  \label{tab:rng_spec}
  \small
  \begin{tabularx}{\textwidth}{l X}
    \toprule
    \textbf{항목} & \textbf{규격 및 근거} \\
    \midrule
    난수 생성기 & PCG64 (Permuted Congruential Generator, 64-bit). NumPy 1.17+의 \texttt{default\_rng()} 사용.
      기존 Mersenne Twister(MT19937) 대비 통계적 품질·속도·시드 공간 우수. \\
    \m{주기(Period)} & $2^{128}$ (MT19937의 $2^{19937}$보다 작으나,
      $n = 10{,}000$ 수준의 시뮬레이션에서 주기 소진 위험 없음) \\
    시드 지정 & 변수별 개별 시드를 명시적으로 지정하여 \texttt{default\_rng(seed)} 호출.
      암묵적 시스템 시각 기반 시드를 사용하지 않음. (명시 미지정 시 시스템 시각 기반 자동 할당으로 재현 불가) \\
    재현성 보장 & 동일 시드$\cdot$동일 알고리즘$\cdot$동일 호출 순서에서 비트 수준 동일 결과 보장.
      NumPy 버전 간 호환성을 위해 \texttt{Generator} API \m{사용(legacy \texttt{RandomState} 미사용)} \\
    분포 생성 & 균등분포: \texttt{rng.uniform(0, 1)},
      정규분포: \texttt{rng.normal(3, 1)},
      파레토분포: \tnew{\texttt{6 * rng.pareto(3)}} ($E[U] = 3.0$ 보장) \\
    \bottomrule
  \end{tabularx}
\end{table}

\chg{여기서 세 가지 설계 근거는 다음과 같다. 시드 지정과 관련하여, 시드를 명시하지 않으면 많은 난수 생성기가 현재 시스템 시각(system clock)을 시드로 자동 할당하여 실행 시점마다 결과가 달라지므로, 본 연구는 모든 확률 스트림에 고정 시드를 명시적으로 지정하였다. 재현성과 관련하여, NumPy의 구 API인 \texttt{RandomState}는 Mersenne Twister 엔진을 사용하여 버전 갱신 시 동일 시드에서도 난수열이 달라질 수 있는 반면, 신 API인 \texttt{Generator}(NumPy 1.17 이상)는 버전 간 비트 수준 재현성을 보장하므로 이를 채택하였다. 분포 생성과 관련하여, NumPy의 \texttt{Generator.pareto(a)}는 Lomax(파레토 II형) 분포를 반환하여 $E[\mathrm{Lomax}(\alpha=3)] = 1/(\alpha-1) = 0.5$이므로 6을 곱한 \texttt{6 * rng.pareto(3)}의 기댓값이 $3.0$이 된다.}


%% -------------------------------------------------------------------
\subsection{\vrev{비중첩 대역 시드 체계}}
\label{subsec:seed_system}
%% -------------------------------------------------------------------

시뮬레이션의 완전한 재현성을 보장하면서 변수 간 통계적 독립성을
유지하기 위해, 마스터 시드 $s = 42$에서 변수별 시드를 비중첩 대역(non-overlapping band)으로
파생하는 체계를 설계하였다. \tabref{tab:seed_system}\refpo{tab:seed_system}{은}{는}
각 변수의 시드 값과 대역 범위를 정리한 것이다.

\begin{table}[htbp]
  \centering
  \caption{비중첩 대역 시드 \m{체계(마스터 시드 $s = 42$)}}
  \label{tab:seed_system}
  \begin{tabular}{llcc}
    \toprule
    변수 & 기호 & 시드 값 & 대역 범위 \\
    \midrule
    관측성 원시 점수 & $O_{\text{raw}}$ & 42 & $[0\text{--}9999]$ \\
    관측성 가중치 & $k_O$ & 20{,}042 & $[20000\text{--}29999]$ \\
    \m{불확실성(C2, Gaussian)} & $U^{(G)}$ & 30{,}042 & $[30000\text{--}39999]$ \\
    \m{불확실성(C3, Pareto)} & $U^{(P)}$ & 40{,}042 & $[40000\text{--}49999]$ \\
    \bottomrule
  \end{tabular}
\end{table}

각 대역은 $10{,}000$개의 고유 시드를 수용하며, 대역 간 간격이 충분하여
난수 생성기(NumPy PCG64)의 상태 간 교차 오염을 구조적으로
방지한다. 마스터 시드를 변경하면 모든 변수의 시드가 동일한 오프셋 규칙으로
갱신되므로, 다중 시드 검증($s = 0, 1, \ldots, 9{,}999$) 시에도
일관된 독립성이 유지된다.
본 시드 체계의 독립성은 \pdferrR{표~\ref{tab:independence}}에서 실증적으로 검증되며 H6을 충족한다.


%% -------------------------------------------------------------------
\subsection{\padd{재현성 자료 공개}}
\label{subsec:reproducibility-release}
%% -------------------------------------------------------------------

\padd{본 연구의 후속 \chg{검증과 확장} 가능성을 확보하기 위해, 시뮬레이션 코드와
데이터를 GitHub 공개 저장소에 공개한다. \chg{저장소의 공개 주소는 \texttt{https://github.com/WindoorsLEE/drto-dynamic-rto}이며, 코드는 MIT 라이선스로, 데이터는 CC BY 4.0 라이선스로 배포한다.} \chg{공개 자료는 네 종류로 구성된다. 첫째, 시뮬레이션 코드(Python 3.12와 NumPy 2.4 기반)는 C0에서 C3까지의 조건별 몬테카를로 생성과 분할 회귀 및 CVaR 산출, 그리고 $10^8$회 다중 시드 강건성 검증 스크립트를 포함한다. 둘째, 원시 난수열(\texttt{raw\_data/})은 비중첩 대역 시드 체계로 생성된 $O_{\text{raw}}$, $k_O$, $U^{(G)}$, $U^{(P)}$ 4종을 표본 크기 $n = 10{,}000$의 CSV로 담는다. 셋째, 계산 결과(\texttt{calc\_data/})는 조건별 $\DRTO$ 값과 평균, 표준편차, 백분위 등 기준값을 정리한 \texttt{statistics.json}으로 이루어진다. 넷째, LaTeX 소스는 본 논문 본문과 부록을 빌드 가능한 형태로 제공한다.}
재현 절차(\texttt{README.md} 참조)는 저장소 클론 후 \texttt{make all} 실행 시
본 연구의 모든 수치 결과(조건별 \chg{평균, CDF, 회귀계수} 등)가 비트 수준에서
동일하게 재현된다.}


%% -------------------------------------------------------------------
\subsection{\vrev{변수 간 독립성 검증}}
\label{subsec:independence}
%% -------------------------------------------------------------------

비중첩 대역 시드 체계의 유효성을 확인하기 위해, 생성된 $n = 10{,}000$개
표본에서 변수 쌍 간 피어슨 상관계수를 산출하였다.
\tabref{tab:independence}에 제시된 바와 같이, 모든 변수 쌍의
상관계수 절댓값이 $|r| < 0.02$로 나타나 실질적 독립성이 확인되었다. \chg{여기서 기준값 $0.02$는 표본 크기 $n = 10{,}000$에서 귀무가설 $H_0\!: \rho = 0$ 하의 상관계수 표준오차가 $\mathrm{SE}(r) \approx 1/\sqrt{n} = 0.01$이고 95\% 신뢰구간이 $\pm 1.96 \times 0.01 \approx \pm 0.020$이라는 데 근거하며, 따라서 $|r| < 0.02$는 관측된 상관이 유의수준 $\alpha = 0.05$에서 영과 통계적으로 구별되지 않음을 의미한다.}

\begin{table}[htbp]
  \centering
  \caption{입력 변수 간 피어슨 상관계수 ($n = 10{,}000$)}
  \label{tab:independence}
  \begin{tabular}{lcc}
    \toprule
    변수 쌍 & 상관계수 $r$ & 독립성 판정 \\
    \midrule
    $k_O$ -- $O_{\text{raw}}$ & $0.016$ & $|r| < 0.02$ \\
    $O_{\text{raw}}$ -- $U^{(G)}$ & $0.009$ & $|r| < 0.02$ \\
    $O_{\text{raw}}$ -- $U^{(P)}$ & $-0.001$ & $|r| < 0.02$ \\
    \bottomrule
  \end{tabular}
\end{table}


%% -------------------------------------------------------------------
\subsection{\vrev{분포 설정 및 \imod{공정 비교}{공정한 비교} 설계}}
\label{subsec:dist_settings}
%% -------------------------------------------------------------------

C2 조건과 C3 조건의 불확실성 분포는 평균이 동일하되 꼬리 특성이
상이하도록 설계하였다. 이를 통해 분포 형태(shape)가 $\DRTO$에 미치는
고유한 효과를 분리하여 관측할 수 있다.

\chg{C2 조건의 가우시안 불확실성은} $U_{\text{raw}} \sim N(\mu=3,\;\sigma=1)$에서 추출한다.
시그모이드 입력 $z_U = \kappa \cdot k_U \cdot R_0 \cdot U_{\text{raw}} / (R_{\max} - R_0)$로
변환되며, $E[U_{\text{raw}}] = 3.0$의 기대값을 가진다. 시뮬레이션에서
관측된 통계량은 $\bar{U}_{\text{C2}} = 3.003$, $\mathrm{SD} = 0.996$,
$P_{99} = 5.333$이다. 가우시안 분포의 특성상 극단값이 평균에서
크게 이탈하지 않으며, $P_{99}/\mu$ 비율은 $1.78$에 불과하다. \chg{여기서 $P_{99}/\mu$ 비율은 99번째 백분위수와 모평균의 비율로서 분포의 우측 꼬리가 중심으로부터 얼마나 이탈하는지를 나타내며, 본 시뮬레이션에서 $\mu$는 추정값이 아니라 설계 파라미터 $E[U_{\text{raw}}] = 3.0$이므로 정확한 모집단 평균을 사용한다.}

\chg{C3 조건의 파레토 불확실성은} $U_{\text{raw}} \sim 6 \cdot \mathrm{Pareto}(\alpha=3)$에서 추출한다.
형상모수 $\alpha = 3$은 2차 모멘트가 존재하되($\alpha > 2$),
파레토(Heavy-Tail) 특성을 보유하는 분포를 생성한다.
시뮬레이션에서 관측된 통계량은 $\bar{U}_{\text{C3}} = 3.019$,
$\mathrm{SD} = 5.303$, $P_{99} = 23.402$이다.
$P_{99}/\mu$ 비율은 $23.402 / 3.0 = 7.80$으로서,
\textbf{가우시안의 $1.78$ 대비 $4.39$배}에 달하며,
극단 사건의 발생 빈도와 규모가 현저히 크다는 것을 확인할 수 있다.

두 분포의 평균 $\bar{U}_{\text{raw}} \approx 3.0$이 실질적으로
동일하므로, C2와 C3 간 $\eta$ 차이는 불확실성의 크기가 아닌
분포 형태, \chg{특히 우측 꼬리의 중량에} 기인하는 것으로 해석할 수 있다.

\chg{이상에서 제시한 연구 설계와 시뮬레이션 방법에 따라 수행한 실증 분석의 결과는 다음 장에서 조건별 기술통계, 회귀 및 분할 회귀, 강건성 검증의 순으로 제시한다.}
        % 3.2 시뮬레이션 설계 (구 3.1)
        % 제3장 연구 방법
% 제4장 연구 결과
% 1차 심사 반영: 구 제3장 결과 부분(3.2--3.8) + 구 제4장 교차분석(4.1--4.6) 통합
% 3차 심사 반영: 파일명을 실제 절 번호와 일치시키도록 정리 (2026-05-09)

\chapter{\vrev{연구 결과}}
\label{ch:results}

\vrev{%
본 장에서는 4단계 조건 체계(C0--C3)별 시뮬레이션 결과를 순차적으로 보고하고,
\chg{회귀분석, 강건성, 수렴 진단}을 통해 모델의 통계적 타당성을 검증한다.
이어서 가설 검증 종합과 조건 간 교차분석(효과 크기, 회귀 설명력,
이중채널 감쇠), 전문가 검증 결과, 선행연구와의 교차 정합성을 제시한다.
}
% 3차 심사 (2026-05-09): 파일명을 실제 절 번호와 일치시키기 위해 ch4 7개 파일 일괄 rename.
% sec4.7_hypothesis → sec4.13_hypothesis (실제 4.13절)
% sec4.8_c2_c3_synthesis → sec4.7_c2_c3_synthesis (실제 4.7절)
% sec4.9_effect_size → sec4.8_effect_size, sec4.10_r2_progression → sec4.9_r2_progression
% sec4.11_dual_channel → sec4.10_dual_channel, sec4.12_expert_verify → sec4.11_expert_verify
% sec4.13_cross_consistency → sec4.12_cross_consistency

% 조건별 시뮬레이션 결과


%% ===================================================================
\section{\vrev{C0와 C1의 관측성 효과}}
\label{sec:result_c0c1}
%% ===================================================================

%% -------------------------------------------------------------------
\subsection{\vrev{정적 기준선 C0}}
\label{subsec:result_c0}
%% -------------------------------------------------------------------

\textbf{C0} 조건은 관측성과 불확실성이 모두 비활성화된 정적 기준선으로서,
$O_{\text{raw}}$, $k_O$, $U_{\text{raw}}$가 모델에 입력되지 않아
$z_O = 0$, $z_U = 0$이 고정된다. 이 경우 개별 시그모이드 바운딩 함수의
인자가 $S(0) = 0$이 되어, $\DRTO = R_0 = 6.0$시간이
결정론적으로 산출된다.

따라서 효율성 비율은 $\eta_{\text{C0}} = \DRTO / R_0 = 1.000$이며,
표준편차 $\mathrm{SD} = 0.000$, 모든 표본에서 동일한 값이 관측된다.
C0는 확률적 변동이 존재하지 않는 기준점으로서, C1 이후 조건에서
나타나는 변동의 크기와 방향을 평가하기 위한 앵커(anchor) 역할을 수행한다.


%% -------------------------------------------------------------------
\subsection{\vrev{C1의 관측성 기반 동적 조정}}
\label{subsec:result_c1}
%% -------------------------------------------------------------------

\textbf{C1} 조건은 관측성($O$, $k_O$)만 활성화하고 불확실성을 비활성화한
상태($z_U = 0$)로, 관측성이 $\DRTO$에 미치는 순수한 효과를 분리하여
관측하는 조건이다.
$\kappa = 1.2$에서 관측성 바운딩 효과 $E_{O,\text{bnd}} = -R_0 \cdot S(\kappa \cdot k_O \cdot O_{\text{raw}})$만
활성화되어, $\DRTO = R_0 + E_{O,\text{bnd}} \leq R_0$가 보장된다.

\chg{기술통계 측면에서,} C1 조건의 효율성 비율 평균은 $\bar{\eta}_{\text{C1}} = 0.853$,
$\mathrm{SD} = 0.125$로 나타났다.
C0의 결정론적 $\eta = 1.000$ 대비 $14.7\%$가 감소한 것으로,
$\DRTO$ 평균은 $5.115$시간이다.
관측성이 도입되면 $\DRTO$가 기본 복구시간 $R_0$보다 체계적으로 단축됨을 의미하며,
이러한 단축 효과는 관측성이 높을수록 이상 징후의 조기 탐지를 통해
복구 시간을 절감할 수 있다는 실무적 논리와 일치한다.

\chg{효과 크기 측면에서,} C0 대비 C1의 Cohen's $d = -1.180$로서, 이는 \citet{cohen1988power}의
기준에 따라 \textbf{대효과(large effect)}에 해당한다(대효과 기준값 $|d|=0.80$의 약 $1.48$배).
방향이 음(-)인 것은 관측성 도입이 $\eta$를 감소(즉, 복구 시간 단축)
시킨다는 것을 나타낸다.
$n = 10{,}000$개 표본 전체에서 $\eta_{\text{C1}} \leq \eta_{\text{C0}}$인
비율이 $100\%$로(\jrev{H3}-a), 관측성의 복구 단축 효과가 예외 없이
관찰되었다. 이는 H1(관측성 효과 가설)의 충족을 실증적으로 확인한다.


%% -------------------------------------------------------------------
\subsection{\vrev{C0/C1 분포 시각화}}
\label{subsec:c0c1_visual}
%% -------------------------------------------------------------------

\figref{fig:eta_c0c1_hist}\refpo{fig:eta_c0c1_hist}{은}{는} C0와 C1 조건의 $\eta$ 분포를 비교한 것이다.
C0는 $\eta = 1.000$에서의 단일 수직선으로 표현되며,
C1은 $[0.468,\;1.000]$ 구간에서 좌측으로 치우친 분포를 형성한다.
\chg{그림에서 적색 점선은 C0 정적 기준선을, 청색 곡선은 C1 관측성 분포의 정규 근사를 나타낸다.}

\begin{figure}[htbp]
  \centering
  \begin{tikzpicture}
    \begin{axis}[
      width=0.9\textwidth,
      height=7cm,
      xlabel={효율성 비율 $\eta$},
      ylabel={밀도 (density)},
      xmin=0.3, xmax=2.2,
      ymin=0,
      ymajorgrids=true,
      grid style={dashed, black},
      legend style={at={(0.98,0.95)}, anchor=north east, font=\small},
      legend cell align=left,
      ]
      % C0: vertical line at eta=1.0
      \addplot[red, ultra thick, dashed, domain=0:8, samples=2]
        coordinates {(1.0,0) (1.0,7.5)};
      \addlegendentry{C0 ($\eta = 1.000$)}

      % C1: approximate density using beta-like shape (mean=0.853, SD=0.125)
      % Skewed left distribution
      \addplot[blue, thick, smooth, domain=0.4:1.0, samples=100]
        {3.2*exp(-0.5*((x-0.853)/0.125)^2)};
      \addlegendentry{C1 ($\bar{\eta} = 0.853$)}

      % 우측 하단 메모: C1 범위 및 산출 근거
      \node[anchor=south east, font=\footnotesize, align=left,
            text=black] at (rel axis cs:0.97,0.03)
        {C1: $\eta \in [0.468,\;1.000]$\\
         $\eta_{\min} = \DRTO_{\min}/R_0 = 2.810/6.0 = 0.468$\\
         $z_{O,\max} = \kappa \cdot k_O \cdot O_{\text{raw}} = 1.2 \times 1 \times 1$};
    \end{axis}
  \end{tikzpicture}
  \caption{C0/C1 조건의 효율성 비율 $\eta$ 분포 비교}
  \label{fig:eta_c0c1_hist}
\end{figure}

\chg{끝으로 H1의 검증 결론을 검증한다.} \mrev{H1은 ``$\DRTO_{\text{C1}} \leq R_0 = 6.0$h가 100\% 표본에서 성립''을
판정 기준으로 한다. $n = 10{,}000$ 표본 전체에서 $\eta_{\text{C1}} \in
[0.468, 1.000]$로 1.0을 초과하는 사례가 0건이며, 순서 준수율이 100.0\%로
H3-a와 결합하여 관측성의 복구 단축 효과가 결정론적으로 확인된다. 효과
크기 Cohen's $d = -1.180$ (매우 큰 효과; \secref{sec:ch5_effect_size}
참조)이 동시에 충족되어, \textbf{H1은 채택된다}.}

\chg{이어 H3-a의 검증 결론은 다음과 같다.} \mrev{H3-a는 ``$\eta_{\text{C1}} \leq \eta_{\text{C0}} = 1.0$인 비율
$\geq 99\%$''를 기준으로 한다. $n = 10{,}000$ 전체에서 100.0\%가 충족되어,
\textbf{H3-a는 채택된다}.}
        % 4.1 C0/C1 결과


%% ===================================================================
\section{\vrev{C2 가우시안 불확실성 결과}}
\label{sec:result_c2}
%% ===================================================================

\textbf{C2} 조건은 C1의 관측성에 가우시안 불확실성($U_{\text{raw}}^{(G)} \sim N(3,1)$)을
추가한 것으로, 경꼬리(light-tail, 가우시안) 불확실성이 $\DRTO$에 미치는 효과를
관측하는 조건이다.
개별 시그모이드 바운딩 모델에서 $E_{O,\text{bnd}}$와 $E_{U,\text{bnd}}$가
각각 독립적으로 바운딩되므로,
$\DRTO = R_0 + E_{O,\text{bnd}} + E_{U,\text{bnd}}$의 구조가
관측성의 하향 효과와 불확실성의 상향 효과를 명확히 분리한다.

\figref{fig:gaussian_concept}\refpo{fig:gaussian_concept}{은}{는} 가우시안 불확실성의 분포 특성과
시그모이드 바운딩을 통한 $\DRTO$ 상향 조정 메커니즘을 개념적으로 도식화한 것이다.
\chg{패널 (a)는 평균 $\mu = 3.0$의 대칭적 확률밀도함수 $N(3,1)$을, 패널 (b)는 불확실성 수준에 따른
$\DRTO$의 상향 조정을 보이며, 시그모이드 바운딩에 의해 $R_0 < \DRTO < R_{\max}$가 보장되어 높은
$U$ 영역에서 자연스럽게 포화한다.}

\begin{figure}[htbp]
  \centering
  \begin{tikzpicture}
    % ── Panel (a): Gaussian N(3,1) PDF ──
    \begin{axis}[
      name=panelA,
      width=0.47\textwidth,
      height=6.5cm,
      xlabel={$U_{\text{raw}}$ (불확실성 수준, 시간)},
      ylabel={확률밀도},
      title={\textbf{(a)} 가우시안 불확실성 $U \sim N(3, 1)$},
      title style={font=\small\bfseries},
      xmin=-1, xmax=7,
      ymin=0, ymax=0.46,
      ymajorgrids=true,
      grid style={dashed, black},
      legend style={at={(0.02,0.98)}, anchor=north west, font=\scriptsize},
      ]
      % Gaussian PDF (approximated with smooth curve)
      \addplot[blue!70!black, thick, smooth, domain=-1:7, samples=100]
        {1/(sqrt(2*pi)*1)*exp(-0.5*((x-3)/1)^2)};
      \addlegendentry{$N(3, 1)$}

      % Mean line
      \draw[red!70!black, dashed, thick] (axis cs:3, 0) -- (axis cs:3, 0.40);
      \node[font=\scriptsize, red!70!black, fill=white, inner sep=1.5pt,
            rounded corners=1pt] at (axis cs:3, 0.44) {$\mu = 3.0$};

      % ±1σ lines
      \draw[black, dotted, thin] (axis cs:2, 0) -- (axis cs:2, 0.245);
      \draw[black, dotted, thin] (axis cs:4, 0) -- (axis cs:4, 0.245);
      \node[font=\tiny, black] at (axis cs:5.5, 0.26) {$\pm 1\sigma$: 68.3\%};
      \draw[-{Stealth}, black, thin] (axis cs:5.2, 0.25) -- (axis cs:4.1, 0.24);

      % Symmetry annotation
      \node[font=\tiny, black, align=center] at (axis cs:-0.2, 0.30) {대칭 분포\\(Skewness $= 0$)};
    \end{axis}

    % ── Panel (b): DRTO shift via sigmoid bounding ──
    \begin{axis}[
      at={(panelA.east)},
      anchor=west,
      xshift=18mm,
      width=0.47\textwidth,
      height=6.5cm,
      xlabel={$U_{\text{raw}}$ (불확실성 수준, 시간)},
      ylabel={$\DRTO$ (시간)},
      title={\textbf{(b)} 시그모이드 바운딩에 의한 불확실성 프리미엄},
      title style={font=\small\bfseries},
      xmin=0, xmax=8,
      ymin=5.0, ymax=25,
      ymajorgrids=true,
      grid style={dashed, black},
      legend style={at={(0.98,0.98)}, anchor=north east, font=\scriptsize},
      ]
      % DRTO = R0 + E_U,bnd = R0 + (Rmax-R0)·S(κ·U/6)
      % S(z) = 2/(1+exp(-z)) - 1
      \addplot[red!70!black, thick, smooth, domain=0:8, samples=100]
        {6.0 + 18.0*(2/(1+exp(-1.2*x/6)) - 1)};
      \addlegendentry{$R_0 + E_{U,\text{bnd}}$}

      % R0 baseline
      \draw[black, dashed, thick] (axis cs:0, 6.0) -- (axis cs:8, 6.0);
      \node[font=\tiny, black] at (axis cs:1.0, 6.5) {$R_0 = 6.0$h};

      % Rmax upper bound
      \draw[black, dotted, thin] (axis cs:0, 24.0) -- (axis cs:8, 24.0);
      \node[font=\tiny, black] at (axis cs:1.0, 24.5) {$R_{\max} = 24.0$h};

      % Mean point
      \addplot[only marks, mark=*, mark size=2.5pt, red!70!black]
        coordinates {(3.0, {6.0 + 18.0*(2/(1+exp(-1.2*3.0/6)) - 1)})};
      \node[font=\tiny, red!70!black, fill=white, inner sep=1pt, rounded corners=1pt,
            anchor=north west] at (axis cs:3.3, 9.0)
        {$U_{\text{raw}} = \mu\;(3.0)$};

      % Saturation annotation
      \node[font=\tiny, blue!70!black, align=center, fill=white, inner sep=1pt,
            rounded corners=1pt] at (axis cs:6.5, 19)
        {시그모이드 포화\\($< R_{\max}$)};
      \draw[-{Stealth}, blue!70!black, thin] (axis cs:6.5, 18.2) -- (axis cs:6.5, 16.5);
    \end{axis}
  \end{tikzpicture}
  \caption{가우시안 불확실성 개념도}
  \label{fig:gaussian_concept}
\end{figure}


%% -------------------------------------------------------------------
\subsection{\vrev{기술통계 및 효과 크기}}
\label{subsec:c2_descriptive}
%% -------------------------------------------------------------------

\chg{기술통계 측면에서,} C2 조건의 효율성 비율 평균은 $\bar{\eta}_{\text{C2}} = 1.682$,
$\mathrm{SD} = 0.615$로 나타났다.
$\DRTO$ 평균은 $10.094$시간으로서, 기준선 $R_0 = 6.0$시간 대비
$68.2\%$가 증가하였다.
C1의 $\bar{\eta} = 0.853$ 대비 $\erf{0.829}{0.830}$가 상승한 것으로,
불확실성의 도입이 $\eta$를 $1.0$ 이상으로
끌어올리는 \textbf{불확실성 프리미엄(uncertainty premium)} 효과를
확인하였다. 즉, 불확실성이 존재하는 환경에서 $\DRTO$는 기본 복구 시간
$R_0$보다 길어지며, 이는 복구 과정의 불확실성을 사전에 반영한
안전 마진(safety margin)으로 해석된다.

\chg{불확실성 프리미엄 비율을 보면,} 전체 $n = 10{,}000$개 표본 중 $\eta_{\text{C2}} > 1.0$인 비율,
즉 $\DRTO$가 기준선 $R_0$를 초과하는 비율은 $84.9\%$이다.
이는 가우시안 불확실성이 존재하는 환경에서 약 $85\%$의 시나리오가
기준선보다 긴 복구 시간을 필요로 함을 의미하며,
불확실성 프리미엄이 예외적 현상이 아닌 \textit{일반적 현상}임을 보여준다.
나머지 $15.1\%$는 관측성이 매우 높은($k_O \approx 1$, $O_{\text{raw}} \approx 1$) 조합에서
발생하는 것으로, 높은 관측성이 불확실성 프리미엄을 상쇄할 수 있는
조건을 나타낸다.

\chg{효과 크기 측면에서,} C1 대비 C2의 Cohen's $d = +1.871$으로서, \textbf{대효과(large effect)}에
해당한다(대효과 기준값 $|d|=0.80$의 약 $2.34$배).
방향이 양(+)인 것은 불확실성 추가가 $\eta$를 증가(복구 시간
연장)시킨다는 것을 나타내며, C0$\to$C1의 감소 방향($d = -1.180$)과
대조를 이룬다. 이러한 방향 반전은 관측성의 단축 효과와 불확실성의
연장 효과가 구조적으로 상반됨을 보여주며,
개별 시그모이드 바운딩 모델의 이론적 설계가 실험적으로 타당함을 입증한다.


%% -------------------------------------------------------------------
\subsection{\vrev{순서 준수 및 가설 검증}}
\label{subsec:c2_order}
%% -------------------------------------------------------------------

$n = 10{,}000$개 표본 중 $\eta_{\text{C1}} \leq \eta_{\text{C2}}$인
비율이 $99.87\%$로(\jrev{H3}-b), C2의 불확실성 프리미엄은 거의 모든
표본에서 관찰되었다. $0.13\%$의 예외는 관측성이 매우 높고($O_{\text{raw}} \approx 1$,
$k_O \approx 1$) 불확실성이 매우 낮은($U_{\text{raw}}^{(G)} \approx 0$) 극단적 조합에서
발생하는 것으로, 모델의 구조적 특성과 일치한다.
\jrev{H3}-b의 기준($\geq 99\%$)을 $99.87\%$로 충족한다.

\chg{C2의 순서 준수에 대한 H3-b의 검증 결론은 다음과 같다.} \mrev{H3-b는 ``$\eta_{\text{C1}} \leq \eta_{\text{C2}}$ 비율 $\geq 99\%$''를
기준으로 한다. 99.87\%가 충족되어 ($0.13\%$ 예외는 $O_{\text{raw}} \approx 1$
극단 영역), \textbf{H3-b는 채택된다}.}
          % 4.2 C2 결과


%% ===================================================================
\section{\vrev{파레토 불확실성 환경의 C3 결과}}
\label{sec:result_c3}
%% ===================================================================

\textbf{C3} 조건은 C2의 가우시안 불확실성을 파레토 분포
($U_{\text{raw}}^{(P)} \sim 6 \cdot \mathrm{Pareto}(\alpha=3)$)로 대체한 것으로,
극단 사건(extreme events)이 $\DRTO$에 미치는 구조적 영향을
관측하는 핵심 조건이다.

\figref{fig:pareto_concept}\refpo{fig:pareto_concept}{은}{는} 가우시안과 파레토 분포의 확률밀도함수 비교 및
생존함수(survival function)를 통한 꼬리 행동의 질적 차이를 도식화한 것이다.
\chg{여기서 생존함수 $\bar{F}(u) = P(U > u) = 1 - F(u)$는 확률변수 $U$가 특정 값 $u$를 초과할 확률을
나타내며, 예컨대 $\bar{F}(10) = 0.01$이면 불확실성이 10시간을 초과할 확률이 1\%임을 의미한다. 생존함수
값이 높을수록 해당 수준 이상의 극단 사건이 빈번함을 뜻하며, 분포의 꼬리가 얼마나 느리게 감소하는지를
직접 비교할 수 있다.}
두 분포는 동일한 기대값 $E[U] = 3.0$을 공유하지만,
파레토 꼬리는 가우시안의 지수적 감쇠에 비해
극단 영역에서 현저히 높은 초과 확률을 나타낸다.
\chg{그림에서 패널 (a)는 동일 평균 $E[U] = 3.0$ 하의 확률밀도함수 비교를, 패널 (b)는 로그 스케일
생존함수 비교를 보인다. 가우시안의 지수적 감쇠($\sim e^{-x^2/2}$)에 비해 파레토 감쇠는 멱법칙($\sim x^{-3}$)을
따르며, 그 결과 P99에서 $U_{\text{raw}}$ 원시값은 가우시안 $5.326$, 파레토 $21.85$로 약 $4.1$배 차이를 보인다.}

\begin{figure}[htbp]
  \centering
  \begin{tikzpicture}
    % ── Panel (a): Gaussian vs Pareto PDF ──
    \begin{axis}[
      name=panelA,
      width=0.47\textwidth,
      height=6.5cm,
      xlabel={$U_{\text{raw}}$ (불확실성 수준, 시간)},
      ylabel={확률밀도},
      title={\textbf{(a)} 가우시안 vs.\ 파레토: 확률밀도함수 비교},
      title style={font=\small\bfseries},
      xmin=0, xmax=15,
      ymin=0, ymax=0.46,
      ymajorgrids=true,
      grid style={dashed, black},
      legend style={at={(0.98,0.88)}, anchor=north east, font=\scriptsize},
      ]
      % Gaussian N(3,1) PDF
      \addplot[blue!70!black, thick, smooth, domain=0:15, samples=120]
        {1/(sqrt(2*pi)*1)*exp(-0.5*((x-3)/1)^2)};
      \addlegendentry{Gaussian $N(3, 1)$}

      % Pareto PDF: f_Y(y) = α/(y/6+1)^(α+1) / 6, α=3
      \addplot[red!70!black, thick, smooth, domain=0.01:15, samples=120]
        {3.0 / ((x/6+1)^4) / 6.0};
      \addlegendentry{Pareto $6 \cdot \mathrm{Pareto}(3)$}

      % Mean line
      \draw[black, dashed, thick] (axis cs:3, 0) -- (axis cs:3, 0.42);
      \node[font=\scriptsize, black, fill=white, inner sep=1.5pt,
            rounded corners=1pt] at (axis cs:3, 0.44)
        {$E[U] = 3.0$h (동일 평균)};

      % Heavy tail annotation
      \node[font=\tiny, red!70!black, align=center, fill=white, inner sep=1pt,
            rounded corners=1pt] at (axis cs:11.5, 0.08)
        {중꼬리(파레토)\\멱법칙 감쇠};

      % Light tail annotation
      \node[font=\tiny, blue!70!black, align=center, fill=white, inner sep=1pt,
            rounded corners=1pt] at (axis cs:7.0, 0.16)
        {경꼬리(가우시안)\\지수적 감쇠};
    \end{axis}

    % ── Panel (b): Survival function (log scale) ──
    \begin{axis}[
      at={(panelA.east)},
      anchor=west,
      xshift=18mm,
      width=0.47\textwidth,
      height=6.5cm,
      xlabel={$U_{\text{raw}}$ (불확실성 수준, 시간)},
      ylabel={$P(U > u)$ (생존함수)},
      title={\textbf{(b)} 꼬리 행동: 생존함수 비교 (로그 스케일)},
      title style={font=\small\bfseries},
      xmin=0, xmax=20,
      ymode=log,
      ymin=1e-6, ymax=1.1,
      ymajorgrids=true,
      grid style={dashed, black},
      legend style={at={(0.98,0.98)}, anchor=north east, font=\scriptsize},
      ]
      % Gaussian survival: 1 - Phi((x-3)/1)
      % Logistic approximation: 1-Phi(z) ≈ 1/(1+exp(1.7*(x-3)))
      \addplot[blue!70!black, thick, smooth, domain=0:15, samples=150]
        {1/(1+exp(1.7*(x-3)))};
      \addlegendentry{Gaussian $N(3, 1)$}

      % Pareto survival: P(Y>y) = (1/(y/6+1))^3
      \addplot[red!70!black, thick, smooth, domain=0.01:20, samples=150]
        {(1/(x/6+1))^3};
      \addlegendentry{Pareto $6 \cdot \mathrm{Pareto}(3)$}

      % P99 line
      \draw[black, dotted, thin] (axis cs:0, 0.01) -- (axis cs:20, 0.01);
      \node[font=\tiny, black, fill=white, inner sep=1pt, anchor=east] at (axis cs:20, 0.003) {P99};

      % Fill excess tail risk
      % (visual approximation — the area where Pareto > Gaussian)
      \node[font=\tiny, red!70!black, fill=orange!10, inner sep=2pt,
            rounded corners=1pt, align=center] at (axis cs:12, 0.005)
        {초과 꼬리 위험\\(Pareto $>$ Gaussian)};
    \end{axis}
  \end{tikzpicture}
  \caption{파레토 불확실성 개념도}
  \label{fig:pareto_concept}
\end{figure}


%% -------------------------------------------------------------------
\subsection{\vrev{기술통계}}
\label{subsec:c3_descriptive}
%% -------------------------------------------------------------------

C3 조건의 효율성 비율 평균은 $\bar{\eta}_{\text{C3}} = 1.514$,
$\mathrm{SD} = 0.791$로 나타났다.
$\DRTO$ 평균은 $9.085$시간이다.
평균은 C2의 $1.682$보다 $0.168$ 낮으나,
표준편차는 C2의 $0.615$ 대비 $1.29$배로 증가하였다.
이는 파레토 분포의 특성이
평균보다는 분산과 극단값에서 현저한 차이를 만들어냄을 보여준다.

\chg{불확실성 프리미엄 비율을 보면,} 전체 표본 중 $\eta_{\text{C3}} > 1.0$인 비율은 $70.8\%$로서,
C2의 $84.9\%$보다 $14.1$\,pp 낮다.
이는 파레토 분포의 높은 분산 때문에 관측성의 하향 효과가
불확실성의 상향 효과를 상쇄하는 빈도가 C2보다 높기 때문이다.
그러나 $\eta > 1.0$인 표본들의 \textit{크기}는 C2보다 훨씬 크며,
이는 CVaR 분석에서 확인된다(\subsecref{subsec:cvar}).

\chg{효과 크기 측면에서,} C2 대비 C3의 Cohen's $d = -0.238$로서 사실상 \textbf{소효과(small effect)}에 해당한다.
이는 두 조건의 평균이 유사하기 때문이며, 파레토 효과가
평균적 수준이 아닌 극단 영역에서 작동한다는 것을 의미한다.
한편, C1 대비 C3의 Cohen's $d = +1.169$는 \textbf{대효과(large effect)}에
해당하여, 순수 관측성 조건(C1)에서 파레토 불확실성 조건(C3)으로의
전환이 $\eta$에 실질적이고 큰 변화를 야기함을 보여준다.
다만, C1 대비 C2의 $d = +1.871$과 비교하면 $0.702$만큼 낮아,
가우시안 불확실성이 평균 수준에서는 파레토보다 더 큰 효과를 산출함을 알 수 있다.
\tabref{tab:c3_effect_size}\refpo{tab:c3_effect_size}{은}{는} C3의 인접 조건 대비 효과 크기를 정리한 것이다.

\begin{table}[htbp]
  \centering
  \caption{C3 조건의 인접 조건 대비 효과 크기 (Cohen's $d$)}
  \label{tab:c3_effect_size}
  \begin{tabular}{lccl}
    \toprule
    비교 쌍 & $\Delta\bar{\eta}$ & Cohen's $d$ & 해석 \\
    \midrule
    C3 vs.\ C0 & $+0.514$ & $+0.650$ & 중효과(medium effect) \\
    C3 vs.\ C1 & $+0.661$ & $+1.169$ & 대효과(large effect) \\
    C3 vs.\ C2 & $-0.168$ & $-0.238$ & 소효과(small effect) \\
    \bottomrule
  \end{tabular}
\end{table}

C3 vs.\ C1의 대효과($d = +1.169$)와 C3 vs.\ C2의 소효과($d = -0.238$)의
대비는, 불확실성의 도입 자체가 핵심 변화 요인이며 분포 형태(가우시안 vs.\ 파레토)의
차이는 평균 수준에서는 미미하되 극단 영역에서 발현됨을 시사한다.

\chg{C1$\leq$C3 순서 준수율(\jrev{H3}-c)을 보면,} 동일 표본($i = 1, \ldots, 10{,}000$)에서 $\eta_{\text{C1},i} \leq \eta_{\text{C3},i}$인
비율\chg{,} 즉 관측성만 적용한 C1의 복구시간이 파레토 불확실성까지 반영한 C3 이하인 표본의 비율은 $100.0\%$로서,
파레토 불확실성 조건에서는 C1 대비 C3의 순서가 \textit{예외 없이} 준수된다.
이는 파레토 분포의 높은 기대값과 양의 스케일 특성에 의해,
불확실성 프리미엄이 관측성 효과를 모든 표본에서 압도하기 때문이다.
C1$\leq$C2 순서 준수율($99.87\%$)과 비교하면, 파레토의 더 큰 불확실성 규모가
예외 없는 순서를 보장하는 데 기여함을 알 수 있다.


%% -------------------------------------------------------------------
\subsection{\vrev{P85.8 CDF 교차점과 분할 회귀}}
\label{subsec:crossover}
%% -------------------------------------------------------------------

C3 조건의 핵심적 구조적 특성은 약 P85.8 지점($\eta \approx 2.42$)에서
C2의 누적분포함수(CDF)와 교차하는 것이다.
이 교차점을 기준으로 분포 수준에서
P85.8 이하 구간은 $\eta_{\text{C3}} < \eta_{\text{C2}}$인 표본이 우세하고,
P85.8 이상 구간은 $\eta_{\text{C3}} > \eta_{\text{C2}}$인 표본이 우세하여
파레토 heavy-tail 효과가 발현한다.
이는 분포 수준의 경향성이며, 개별 관측값에 항상 성립하는 것은 아니다.

\figref{fig:c2c3_cdf}\refpo{fig:c2c3_cdf}{은}{는} C2와 C3의 $\eta$ CDF를 비교하여
P85.8 CDF 교차점을 시각화한 것이다.
분할 회귀의 구조적 분리점 역시 P85.8으로, CDF 교차점과 일치한다.

\begin{figure}[htbp]
  \centering
  \begin{tikzpicture}
    \begin{axis}[
      width=0.9\textwidth,
      height=7.5cm,
      xlabel={효율성 비율 $\eta$},
      ylabel={누적 확률 $F(\eta)$},
      xmin=0.3, xmax=4.2,
      ymin=0, ymax=1.05,
      ymajorgrids=true,
      xmajorgrids=true,
      grid style={dashed, black},
      legend style={at={(0.98,0.02)}, anchor=south east, font=\small},
      legend cell align=left,
      ]
      % Core/Tail 배경 영역
      \fill[green!8] (axis cs:0.3,0) rectangle (axis cs:2.42,1.05);
      \fill[red!8] (axis cs:2.42,0) rectangle (axis cs:4.2,1.05);

      % C2: 실제 경험적 CDF (시뮬레이션 데이터)
      \addplot[blue!70!black, thick, no markers]
        table[x=eta_C2, y=F_C2] {figures/data/cdf_c2_c3.dat};
      \addlegendentry{C2 (Gaussian, $\bar{\eta}=1.682$)}

      % C3: 실제 경험적 CDF (시뮬레이션 데이터)
      \addplot[red!70!black, thick, no markers]
        table[x=eta_C3, y=F_C3] {figures/data/cdf_c2_c3.dat};
      \addlegendentry{C3 (Pareto, $\bar{\eta}=1.514$)}

      % P85.8 교차점 수평선
      \draw[black, dashed, thick] (axis cs:0.3,0.858) -- (axis cs:4.2,0.858);

      % η≈2.42 수직선 (교차점 위치)
      \draw[black, dashed, thick] (axis cs:2.42, 0) -- (axis cs:2.42, 0.858);

      % 교차점 마커
      \addplot[only marks, mark=*, mark size=3pt, black, fill=gray!15]
        coordinates {(2.42, 0.858)};

      % 교차점 통합 레이블
      \node[font=\scriptsize, align=center, fill=white, draw=black, inner sep=2pt,
            rounded corners=1pt, text=black]
        at (axis cs:1.4, 0.95) {P85.8 CDF 교차점\\$(\eta,\; F) = (2.42,\; 0.858)$};
      \draw[-{Stealth}, black, thin]
        (axis cs:1.85, 0.92) -- (axis cs:2.38, 0.862);

      % Core 영역 레이블
      \node[font=\scriptsize, text=black]
        at (axis cs:2.05, 0.45) {$\longleftarrow$ Core 영역};

      % Tail 영역 레이블
      \node[font=\scriptsize, text=black]
        at (axis cs:2.80, 0.45) {Tail 영역 $\longrightarrow$};
    \end{axis}
  \end{tikzpicture}
  \caption{C2 vs.\ C3 누적분포함수(CDF) 비교 및 P85.8 CDF 교차점}
  \label{fig:c2c3_cdf}
\end{figure}


%% -------------------------------------------------------------------
\subsection{\vrev{CVaR 분석을 통한 극단 위험 정량화}}
\label{subsec:cvar}
%% -------------------------------------------------------------------

조건부 꼬리 기대값(Conditional Value-at-Risk, CVaR)~\citep{rockafellar2000optimization}은 특정 백분위 이상의
$\eta$ 평균을 나타내는 꼬리 위험 척도로서, 평균이나 표준편차로는
포착할 수 없는 극단 영역의 위험을 정량화한다.
\chg{구체적으로 $\mathrm{CVaR}_q = E[\eta \mid \eta \geq \eta_q]$로 정의되며, $q$번째 백분위 이상에 해당하는
표본들의 조건부 평균이다. 예컨대 $\mathrm{CVaR}_{99}$는 상위 1\%($n = 10{,}000$ 중 100개) 극단값의
평균으로서, 최악의 1\% 시나리오에서 평균적으로 복구시간이 얼마나 확대되는가를 정량화한다.}
\tabref{tab:cvar_comparison}\refpo{tab:cvar_comparison}{은}{는} C2와 C3 조건에 대해 6개 수준의 CVaR을 비교한 것이며,
\figref{fig:cvar_bar}\refpo{fig:cvar_bar}{은}{는} 이를 막대 차트로 시각화한 것이다.

\begin{table}[htbp]
  \centering
  \caption{C2 vs.\ C3 조건부 꼬리 기대값(CVaR) 비교}
  \label{tab:cvar_comparison}
  \begin{tabular}{lcccc}
    \toprule
    CVaR 수준 & C2 (Gaussian) & C3 (Pareto) & $\Delta$ & 증가율 \\
    \midrule
    CVaR$_{80}$ & 2.587 & 2.866 & $+0.279$ & $+10.8\%$ \\
    CVaR$_{85}$ & 2.674 & 3.082 & $+0.408$ & $+15.3\%$ \\
    CVaR$_{90}$ & 2.778 & 3.346 & $+0.568$ & $+20.4\%$ \\
    CVaR$_{95}$ & 2.925 & 3.665 & $+0.740$ & $+25.3\%$ \\
    CVaR$_{97}$ & 3.017 & 3.808 & $+0.791$ & $+26.2\%$ \\
    CVaR$_{99}$ & 3.180 & 3.944 & $+0.764$ & $+24.0\%$ \\
    \bottomrule
  \end{tabular}
\end{table}

\begin{figure}[htbp]
  \centering
  \begin{tikzpicture}
    \begin{axis}[
      ybar,
      width=0.92\textwidth,
      height=7cm,
      bar width=9pt,
      ylabel={CVaR ($\eta$)},
      symbolic x coords={CVaR$_{80}$, CVaR$_{85}$, CVaR$_{90}$, CVaR$_{95}$, CVaR$_{97}$, CVaR$_{99}$},
      xtick=data,
      xticklabel style={font=\small},
      ymin=2.0, ymax=4.2,
      ymajorgrids=true,
      grid style={dashed, black},
      nodes near coords,
      nodes near coords style={font=\tiny, above=1pt},
      every node near coord/.append style={/pgf/number format/.cd, fixed, precision=3},
      enlarge x limits=0.12,
      legend style={at={(0.02,0.98)}, anchor=north west, font=\small},
      ]
      \addplot[fill=blue!30, draw=blue!60] coordinates {
        (CVaR$_{80}$, 2.587)
        (CVaR$_{85}$, 2.674)
        (CVaR$_{90}$, 2.778)
        (CVaR$_{95}$, 2.925)
        (CVaR$_{97}$, 3.017)
        (CVaR$_{99}$, 3.180)
      };
      \addlegendentry{C2 (Gaussian)}
      \addplot[fill=red!30, draw=red!60] coordinates {
        (CVaR$_{80}$, 2.866)
        (CVaR$_{85}$, 3.082)
        (CVaR$_{90}$, 3.346)
        (CVaR$_{95}$, 3.665)
        (CVaR$_{97}$, 3.808)
        (CVaR$_{99}$, 3.944)
      };
      \addlegendentry{C3 (Pareto)}
    \end{axis}
  \end{tikzpicture}
  \caption{C2 vs.\ C3 CVaR 비교 (CVaR$_{80}$--CVaR$_{99}$)}
  \label{fig:cvar_bar}
\end{figure}

CVaR$_{80}$에서 CVaR$_{99}$로 갈수록 C3와 C2의 격차가 확대된다.
CVaR$_{80}$ 수준에서는 C3가 C2 대비 $+10.8\%$ 높으며,
CVaR$_{95}$ 수준에서는 $+25.3\%$로 약 $2.3$배 이상 확대된다.
이러한 비선형적 격차 확대는 파레토 분포의 꼬리 특성에 기인하며,
극단 사건의 빈도가 높은 산업(금융, 의료, 통신)에서 C3 조건 적용의
필요성을 직접 뒷받침한다.

\chg{CVaR 에스컬레이션 패턴을 보면,} C2의 CVaR은 $80 \to 99$ 수준에서 $2.587 \to 3.180$으로
상승하는 반면($\Delta_{80 \to 99} = 0.593$),
C3의 CVaR은 $2.866 \to 3.944$로 상위 백분위에서
더 가파르게 상승한다($\Delta_{80 \to 99} = 1.078$).
이는 C3의 극단 영역에서 $\eta$의 에스컬레이션 속도가 C2의 약 $1.8$배에
달함을 의미하며, 꼬리 위험의 구조적 비대칭성을 보여준다.

\chg{P85.8 분할 기준의 실무적 의미는 다음과 같다.} CVaR 분석 결과는 P85.8 분할 회귀 분석과 일관된다.
전체 표본의 $85.8\%$에 해당하는 Core 영역에서는 C2와 C3의 $\eta$가
거의 구분되지 않으나,
상위 $14.2\%$의 Tail 영역에서 파레토 효과가 급격히 발현한다.
이는 극단 위험 관리가 필요한 조직\chg{(}금융기관의 시장 리스크,
의료기관의 사이버 공격 대응, 통신사의 대규모 네트워크 장애 등\chg{)}에서
C3 조건의 적용이 사실상 필수적임을 시사한다.


%% -------------------------------------------------------------------
\subsection{\vrev{C2와 C3 핵심 지표 비교}}
\label{subsec:c2_c3_comparison}
%% -------------------------------------------------------------------

\tabref{tab:c2_c3_key_metrics}\refpo{tab:c2_c3_key_metrics}{은}{는} 동일 시뮬레이션 설계($n = 10{,}000$,
$\kappa = 1.2$, $R_{\max} = 24.0$h, $E[U] = 3.0$) 하에서
가우시안(C2)과 파레토(C3) 불확실성의 핵심 지표를 직접 대비한 것이다.

\begin{table}[htbp]
  \centering
  \caption{C2(가우시안) vs C3(파레토) 분포 통계 비교}
  \label{tab:c2_c3_key_metrics}
  \small
  \begin{tabular}{l rr l}
    \toprule
    \textbf{지표} & \textbf{C2 (가우시안)} & \textbf{C3 (파레토)} & \textbf{해석} \\
    \midrule
    $\bar{\eta}$          & 1.682 & 1.514 & C2가 평균적으로 높음 \\
    $\text{SD}(\eta)$     & 0.615 & 0.791 & C3가 분산 더 큼 \\
    $\eta_{\max}$         & 3.443 & 4.000 & C3가 $R_{\max}$ 포화 접근 \\
    프리미엄 비율         & 84.9\% & 70.8\% & C2가 더 빈번하게 초과 \\
    $d$ (vs C1)           & $+1.871$ & $+1.169$ & C2의 효과 크기 더 큼 \\
    순서 준수율 (C1$\leq$C$x$) & 99.87\% & 100.0\% & 양쪽 모두 높음 \\
    \bottomrule
  \end{tabular}
\end{table}
회귀분석 측면의 비교(1차/2차 $R^2$, 전진 선택, LASSO)는 4.4절에서
상세히 다룬다.

이 비교로부터 두 가지 분포 구조의 차이가 도출된다.
\textbf{첫째}, 가우시안(C2)의 대칭적 종형 분포는 평균 중심의 예측 가능한
변동을 생성하여 ``일상적 운영 변동''의 모델링에 적합하다.
\textbf{둘째}, 비대칭 파레토(C3) 분포는 극단 꼬리에서 $\eta \to 4.0$
($= R_{\max}/R_0$)에 접근하는 소수 사례를 생성하여, 꼬리 위험 탐지에서
본질적 우위를 보인다.
이러한 ``이중축'' 구조\chg{(}일상적 변동은 경꼬리(가우시안, C2) 기반으로,
극단 위험은 중꼬리(파레토, C3) 기반으로 관리\chg{)}는 \chref{ch:discussion}의
실무 적용 가이드라인에서 구체화된다.

\chg{끝으로 H3-c의 검증 결론은 다음과 같다.} \mrev{H3-c는 ``$\eta_{\text{C2}} \leq \eta_{\text{C3}}$''를 \emph{조건부}
가설로 하며, 전체 표본에서는 31.8\%에 그치나 P85.8 CDF 교차점 이상의
극단 영역에서 94.9\%가 역전된다. 조건부 명제로서 \textbf{H3-c는 채택된다}.}
          % 4.3 C3 결과

% 통계 분석


%% ===================================================================
\section{\vrev{회귀분석 결과}}
\label{sec:regression_results}
%% ===================================================================

본 절에서는 C1, C2, C3 각 조건에 대해 3단계 회귀분석---1차 단순,
1차 다중, 2차 회귀---의 설명력 진행($R^2$ progression)을
체계적으로 보고한다. 2차 회귀는 1차항, 2차항, 교호작용항을 포함하는
완전 2차 다항 모형(complete quadratic polynomial model)이다.


%% -------------------------------------------------------------------
\subsection{\vrev{C1 회귀분석(관측성 구조)}}
\label{subsec:reg_c1}
%% -------------------------------------------------------------------

\chg{C1 회귀분석은 설명력이 점증하는 세 단계로 진행한다. 첫째, 1차 단순 회귀는}
$\hat{\eta}_{\text{C1}}^{(1s)} = 0.997 - 0.285\,k_O$ ($R^2 = 0.428$, $p < 0.001$)\chg{로서,}
관측성 가중치 $k_O$만으로 분산의 $42.8\%$를 설명하며, $k_O$가
증가할수록 관측성의 복구 단축 효과가 강화되는 선형 관계가 확인되었다.
\chg{둘째, 1차 다중 회귀는} $\hat{\eta}_{\text{C1}}^{(1m)} = 1.07 - 0.15\,k_O - 0.15\,O$ ($R^2 = 0.869$)\chg{로서,}
관측성 원시 점수 $O$를 추가하면 설명력이 $86.9\%$로 두 배 이상
증가한다. $k_O$와 $O$의 회귀 계수 크기가 $-0.15$로 거의 동일하여,
두 관측성 변수가 $\eta$ 감소에 동등하게 기여함을 보여준다.
\chg{셋째, 2차교호작용 회귀는 다음과 같다.}
\begin{equation}
  \hat{\eta}_{\text{C1}}^{(2q)} = \hat{\beta}_0
    + \hat{\beta}_1 k_O + \hat{\beta}_2 O
    + \hat{\beta}_3 k_O^2 + \hat{\beta}_4 O^2
    + \hat{\beta}_5 (k_O \cdot O)
  \quad (R^2 = 1.000)
\end{equation}
\begin{equation}
  = 1.011
    - 0.036\,k_O - 0.036\,O
    + 0.024\,k_O^2 + 0.024\,O^2
    - 0.546\,(k_O \cdot O)
\end{equation}
정밀 표기 시 $R^2 = 0.9995$로 거의 1.000에 근접하며,
핵심 교호작용항 $k_O \cdot O$(계수 $= -0.546$)가
$R^2$를 $0.869 \to 1.000$으로 향상시키는 결정적 기여를 한다.
1차항($k_O$, $O$)과 2차항($k_O^2$, $O^2$)의 계수는 모두
$|\hat{\beta}| < 0.04$로 미미한 반면,
교호작용항의 계수 절댓값이 $0.546$으로 압도적이다.
이는 C1의 $\eta$가 개별 시그모이드 바운딩 함수의 해석적 구조에 의해
$k_O$와 $O$의 2차 다항식으로 정확히 근사됨을 의미하며,
잔차가 수치 정밀도 수준에 불과하다.


%% -------------------------------------------------------------------
\subsection{\vrev{C2 회귀분석(불확실성 통합 구조)}}
\label{subsec:reg_c2}
%% -------------------------------------------------------------------

\chg{C2 회귀분석도 동일한 세 단계로 진행한다. 첫째, 1차 단순 회귀는}
$\hat{\eta}_{\text{C2}}^{(1s)} = 2.633 - 1.878\,k_O$ ($R^2 = 0.768$)\chg{로서,}
$k_O$의 회귀 계수가 C1의 $-0.285$에서 $-1.878$로 $6.6$배 증가하였다.
이는 불확실성이 존재하는 환경에서 관측성 가중치의 영향력이
증폭된다는 것을 의미하며, $R^2$도 $0.428 \to 0.768$으로
대폭 향상되었다.
\chg{둘째, 1차 다중 회귀는} $\hat{\eta}_{\text{C2}}^{(1m)} = 2.024 - 1.869\,k_O - 0.281\,O + 0.248\,U$ ($R^2 = 0.948$)\chg{로서,}
불확실성 변수 $U$가 양의 계수($+0.248$)로 유의미하게 포함되어,
불확실성이 클수록 $\eta$가 증가하는 관계가 확인되었다.
\chg{셋째, 2차교호작용 회귀는 다음과 같다.}
\begin{equation}
  \hat{\eta}_{\text{C2}}^{(2q)} = \hat{\beta}_0
    + \sum_{j} \hat{\beta}_j x_j
    + \sum_{j} \hat{\beta}_{jj} x_j^2
    + \sum_{j<k} \hat{\beta}_{jk} (x_j \cdot x_k)
  \quad (R^2 = 0.998)
\end{equation}
\begin{equation}
  \begin{aligned}
    = 1.011
    &- 0.036\,k_O - 0.036\,O + 0.598\,U \\
    &+ 0.022\,k_O^2 + 0.024\,O^2 - 0.003\,U^2 \\
    &- 0.546\,(k_O \cdot O) - 0.593\,(k_O \cdot U) + 0.001\,(O \cdot U)
  \end{aligned}
\end{equation}
정밀 표기 시 $R^2 = 0.9982$이며,
9개 변수($k_O$, $O$, $U$, $k_O^2$, $O^2$, $U^2$,
$k_O \cdot O$, $k_O \cdot U$, $O \cdot U$)를 포함한다.
$k_O \cdot O$, $k_O \cdot U$ 등의 교호작용이 핵심 구조를
이루며, C2 조건의 분산을 사실상 완전히 설명한다.
$k_O \cdot U$ 교호작용의 존재는 관측성 가중치와 불확실성이
독립적이 아니라 상호작용하여 $\eta$에 영향을 미침을 시사한다.


%% -------------------------------------------------------------------
\subsection{\vrev{C3 회귀분석(파레토 구조와 분할 회귀)}}
\label{subsec:reg_c3}
%% -------------------------------------------------------------------

\chg{C3 회귀분석은 전체 표본에 대해 세 단계로 진행한 뒤 분할 회귀로 이어진다. 첫째, 1차 단순 회귀는}
$\hat{\eta}_{\text{C3}}^{(1s)} = 2.220 - 1.394\,k_O$ ($R^2 = 0.256$)\chg{로서,}
전체 표본에 대한 설명력은 $25.6\%$에 불과하여, C3의 변동을
$k_O$ 단독으로는 충분히 설명하지 못한다.
\chg{둘째, 1차 다중 회귀는} $\hat{\eta}_{\text{C3}}^{(1m)} = 2.018 - 1.408\,k_O - 0.249\,O + 0.112\,U$ ($R^2 = 0.695$)\chg{로서,}
1차 다중 회귀에서도 C3의 설명력은 $69.5\%$로 C2의 $94.8\%$에
미치지 못하며, 파레토 분포의 극단값이 선형 모형의 한계를
드러내는 것으로 해석된다.
\chg{셋째, 2차+교호작용 회귀는 다음과 같다.}
\begin{equation}
  \hat{\eta}_{\text{C3}}^{(2q)} = \hat{\beta}_0
    + \sum_{j} \hat{\beta}_j x_j
    + \sum_{j} \hat{\beta}_{jj} x_j^2
    + \sum_{j<k} \hat{\beta}_{jk} (x_j \cdot x_k)
  \quad (R^2 = 0.858)
\end{equation}
\begin{equation}
  \begin{aligned}
    = 1.430
    &- 0.216\,k_O + 0.085\,O + 0.256\,U \\
    &- 0.513\,k_O^2 + 0.022\,O^2 - 0.002\,U^2 \\
    &- 0.627\,(k_O \cdot O) - 0.125\,(k_O \cdot U) - 0.025\,(O \cdot U)
  \end{aligned}
\end{equation}
C2와 동일한 9개 변수 구조이나, 전체 $R^2 = 0.858$로
C1($1.000$) 및 C2($0.998$)와 달리 높은 설명에는 도달하지 못한다.
이는 파레토 분포의 비선형성이 2차 회귀 다항식의 근사 능력을 초과함을 시사한다.

\chg{P85.8을 기준으로 분할 회귀를 수행하면 각 영역에서의 설명력이 현저히 개선된다.
Core 영역($\leq$ P85.8)에서는 2차교호작용 $R^2 = 0.999$로 C2와 유사한 완전 설명이 달성되며,
Tail 영역($>$ P85.8)에서도 분할 회귀 적용 시 2차교호작용 $R^2 = 0.710$으로 설명력이 개선된다.}

\p{[}표~\vref{tab:c3_split_regression}\p{]}는 C3 분할 회귀의 주요 결과를 요약한다.

\begin{table}[htbp]
  \centering
  \caption{C3 분할 회귀분석 결과 (P85.8 기준)}
  \label{tab:c3_split_regression}
  \begin{tabular}{llcccc}
    \toprule
    영역 & 모형 & $n$ & $k_O$ 계수 & $R^2$ & 비고 \\
    \midrule
    \multirow{3}{*}{전체}
      & 1차 단순 & 10{,}000 & $-1.394$ & 0.256 & --- \\
      & 1차 다중 & 10{,}000 & $-1.408$ & 0.695 & $+O,\;U$ \\
      & 2차 회귀 & 10{,}000 & --- & 0.858 & --- \\
    \midrule
    \multirow{3}{*}{Core ($\leq$ P85.8)}
      & 1차 단순 & 8{,}580 & $-0.798$ & 0.263 & --- \\
      & 1차 다중 & 8{,}580 & $-1.034$ & 0.630 & $+O,\;U$ \\
      & 2차 회귀 & 8{,}580 & --- & 0.999 & 완전 설명 \\
    \midrule
    \multirow{3}{*}{Tail ($>$ P85.8)}
      & 1차 단순 & 1{,}420 & $-0.415$ & 0.032 & $0.52$배 감쇠 \\
      & 1차 다중 & 1{,}420 & $-1.664$ & 0.860 & $+O,\;\ln U$ \\
      & 2차 회귀 & 1{,}420 & --- & 0.710 & --- \\
    \bottomrule
  \end{tabular}
\end{table}

\chg{표의 Tail 영역 1차 다중 회귀에서 불확실성 $U$를 원래 스케일로 쓰면 $R^2$가 0.457에 불과하나,
$\ln U$로 변환하면 0.860으로 크게 개선된다. 이는 파레토 분포의 멱법칙 구조($P(U > u) \propto u^{-\alpha}$)를
로그 변환이 선형화하기 때문이며, Tail의 $\DRTO$가 불확실성의 절대 크기가 아니라 로그 스케일에 비례하여
증가함을 의미한다. Core와 Tail 분할 회귀의 전체 계수는 공개 재현 패키지에 수록한다.}

Tail 영역의 1차 단순 회귀에서 $k_O$ 계수는 $-0.415$로,
Core 영역($-0.798$) 대비 $0.52$배로 감쇠된다(전체 계수는 공개 재현 패키지 수록).
이 $0.52$배 감쇠가 \textbf{이중채널 감쇠(dual-channel attenuation)}의
핵심 지표이다. 관측성 가중치 $k_O$가 동일하게 1단위 변화할 때,
극단 불확실성 영역에서는 관측성의 복구 단축 효과가 핵심 영역의 $0.52$배로 약화된다.
이는 관측성과 파레토 불확실성이 결합될 때 비선형적으로 상호작용하여,
극단 상황에서 관측성의 역할이 질적으로 변화함을 나타낸다.


%% -------------------------------------------------------------------
\subsection{\vrev{$R^2$ 진행 종합 비교}}
\label{subsec:r2_progression}
%% -------------------------------------------------------------------

\p{[}표~\vref{tab:regression_progression}\p{]}는 조건별 회귀분석 설명력($R^2$)의
3단계 진행을 비교한 것이며, \p{[}그림~\vref{fig:r2_progression}\p{]}은
이를 그룹 막대 차트로 시각화한 것이다.
\chg{표에서 (1s)는 1차 단순, (1m)은 1차 다중, (2q)는 2차교호작용 회귀를 가리킨다.}

\begin{table}[htbp]
  \centering
  \caption{조건별 회귀분석 설명력($R^2$) 진행 비교}
  \label{tab:regression_progression}
  \begin{tabular}{lccc}
    \toprule
    조건 & $R^2_{(1\mathrm{s})}$ & $R^2_{(1\mathrm{m})}$ & $R^2_{(2\mathrm{q})}$ \\
    \midrule
    C1 (관측성) & 0.428 & 0.869 & 1.000 \\
    C2 (Gaussian) & 0.768 & 0.948 & 0.998 \\
    C3 전체 (Pareto) & 0.256 & 0.695 & 0.858 \\
    \quad C3 Core ($\leq$ P85.8) & 0.263 & 0.630 & 0.999 \\
    \quad C3 Tail ($>$ P85.8) & 0.032 & 0.860 & 0.710 \\
    \bottomrule
  \end{tabular}
\end{table}

\begin{figure}[htbp]
  \centering
  \begin{tikzpicture}
    \begin{axis}[
      ybar,
      width=0.92\textwidth,
      height=7.5cm,
      bar width=8pt,
      ylabel={$R^2$},
      symbolic x coords={C1, C2, C3 전체, C3 Core, C3 Tail},
      xtick=data,
      xticklabel style={font=\small},
      ymin=0, ymax=1.15,
      ymajorgrids=true,
      grid style={dashed, black},
      nodes near coords,
      nodes near coords style={font=\tiny, anchor=south, yshift=1pt},
      every node near coord/.append style={/pgf/number format/.cd, fixed, precision=3},
      enlarge x limits=0.15,
      legend style={at={(0.5,-0.18)}, anchor=north, font=\small,
        legend columns=3, draw=black, fill=white},
      ]
      % Simple regression
      \addplot[fill=blue!20, draw=blue!50] coordinates {
        (C1, 0.428)
        (C2, 0.768)
        (C3 전체, 0.256)
        (C3 Core, 0.263)
        (C3 Tail, 0.032)
      };
      \addlegendentry{1차 단순}

      % Multiple regression
      \addplot[fill=green!25, draw=green!60] coordinates {
        (C1, 0.869)
        (C2, 0.948)
        (C3 전체, 0.695)
        (C3 Core, 0.630)
        (C3 Tail, 0.860)
      };
      \addlegendentry{1차 다중}

      % Quadratic regression
      \addplot[fill=red!25, draw=red!60] coordinates {
        (C1, 1.000)
        (C2, 0.998)
        (C3 전체, 0.858)
        (C3 Core, 0.999)
        (C3 Tail, 0.710)
      };
      \addlegendentry{2차교호작용}

      % H4 threshold line (전체 수평선)
      \draw[dashed, black, thick] (rel axis cs:0,0.826) -- (rel axis cs:1,0.826);
      \node[font=\scriptsize, fill=white, inner sep=1.5pt,
        anchor=south east] at (rel axis cs:0.99,0.826)
        {H4 판정 기준 ($R^2 \geq 0.95$)};
    \end{axis}
  \end{tikzpicture}
  \caption{조건별 $R^2$ 진행 비교 (1차 단순 $\to$ 1차 다중 $\to$ 2차 회귀)}
  \label{fig:r2_progression}
\end{figure}

C1은 2차교호작용 모형에서 $R^2 = 1.000$, C2는 $R^2 = 0.998$에 도달하여 개별 시그모이드 바운딩 함수가
2차교호작용 다항식으로 높은 정밀도로 근사됨을 보여준다. C3는 전체 표본에서 $R^2 = 0.858$로
완전 설명에 미달하며, P85.8 기준 분할 시 Core 영역에서 $R^2 = 0.999$,
Tail 영역에서 $R^2 = 0.710$을 달성하여, 분할 회귀가 C3의
이질적 구조를 효과적으로 포착함을 입증한다.


%% -------------------------------------------------------------------
\subsection{\vrev{회귀 계수 비교 시각화}}
\label{subsec:coeff_comparison}
%% -------------------------------------------------------------------

\p{[}그림~\vref{fig:regression_coefficients}\p{]}은 C1, C2, C3(전체) 1차 다중 회귀모형의
주요 계수를 비교한 것이다.
1차 단순 회귀는 $k_O$ 단일 계수만 존재하여 조건 간 비교 차트의 의미가 제한적이고,
2차 회귀는 최대 9개 항의 계수가 포함되어 막대 차트보다 표 형태가
적합하므로(\p{[}표~\vref{tab:regression_progression}\p{]} 참조),
3개 공통 변수($k_O$, $O$, $U$)를 갖는 1차 다중 회귀만 시각화한다.

\begin{figure}[htbp]
  \centering
  \begin{tikzpicture}
    \begin{axis}[
      ybar,
      width=0.88\textwidth,
      height=7.5cm,
      bar width=12pt,
      ylabel={회귀 계수 $\hat{\beta}$},
      symbolic x coords={절편, $k_O$, $O$, $U$},
      xtick=data,
      xticklabel style={font=\small},
      ymin=-2.2, ymax=2.4,
      ymajorgrids=true,
      grid style={dashed, black},
      nodes near coords,
      nodes near coords style={font=\tiny},
      every node near coord/.append style={
        /pgf/number format/.cd, fixed, fixed zerofill, precision=2,
        /tikz/anchor=south,
        /tikz/yshift=1pt,
      },
      enlarge x limits=0.2,
      legend style={at={(0.98,0.98)}, anchor=north east, font=\small},
      ]
      \addplot[fill=blue!25, draw=blue!60] coordinates {
        (절편, 1.07) ($k_O$, -0.15) ($O$, -0.15) ($U$, 0)
      };
      \addlegendentry{C1}

      \addplot[fill=green!25, draw=green!60] coordinates {
        (절편, 2.02) ($k_O$, -1.87) ($O$, -0.28) ($U$, 0.25)
      };
      \addlegendentry{C2}

      \addplot[fill=red!25, draw=red!60] coordinates {
        (절편, 2.02) ($k_O$, -1.41) ($O$, -0.25) ($U$, 0.11)
      };
      \addlegendentry{C3}

      % Zero line
      \draw[dashed, black, thin] (axis cs:절편, 0) -- (axis cs:$U$, 0);
    \end{axis}
  \end{tikzpicture}
  \caption{C1/C2/C3 1차 다중 회귀 계수 비교}
  \label{fig:regression_coefficients}
\end{figure}


%% -------------------------------------------------------------------
\subsection{\vrev{전진 선택법과 LASSO 정규화를 통한 변수 선택}}
\label{subsec:variable_selection}
%% -------------------------------------------------------------------

2차 완전 모형은 9개 독립변수($k_O$, $O$, $U$, $k_O^2$, $O^2$, $U^2$,
$k_O \cdot O$, $k_O \cdot U$, $O \cdot U$)를 포함하여 높은 설명력을 보이나,
과적합(overfitting) 위험과 해석 가능성(interpretability) 저하가 우려된다.
이에 \textbf{전진 선택법(Forward Selection)}과 \textbf{LASSO 정규화}를
적용하여 핵심 변수를 체계적으로 식별하였다.

\p{[}표~\vref{tab:forward_selection}\p{]}은 C2와 C3 각 조건에서 전진 선택법에 의한
변수 진입 순서와 누적 $R^2$를 비교한 것이다.

\begin{table}[htbp]
  \centering
  \caption{전진 선택법 변수 진입 순서 및 누적 $R^2$ (C2 vs.\ C3)}
  \label{tab:forward_selection}
  \small
  \begin{tabular}{c l c c l c}
    \toprule
    & \multicolumn{2}{c}{\textbf{C2 (Gaussian)}} & & \multicolumn{2}{c}{\textbf{C3 (Pareto)}} \\
    \cmidrule{2-3} \cmidrule{5-6}
    순서 & 변수 & 누적 $R^2$ & & 변수 & 누적 $R^2$ \\
    \midrule
    1 & $k_O$ & 0.768 & & $U$ & 0.424 \\
    2 & $U$ & 0.931 & & $k_O$ & 0.686 \\
    3 & $k_O \cdot U$ & 0.973 & & $U^2$ & 0.799 \\
    4 & $k_O \cdot O$ & 0.996 & & $k_O \cdot U$ & 0.839 \\
    5 & $k_O^2$ & 0.997 & & $k_O \cdot O$ & 0.853 \\
    6 & $U^2$ & 0.998 & & $k_O^2$ & 0.856 \\
    7 & $O$ & 0.998 & & $O \cdot U$ & 0.857 \\
    8 & $O^2$ & 0.998 & & $O$ & 0.858 \\
    9 & $O \cdot U$ & 0.998 & & $O^2$ & 0.858 \\
    \bottomrule
  \end{tabular}
\end{table}

변수 진입 순서의 차이는 두 조건의 구조적 특성을 반영한다.
C2에서는 관측성 가중치 $k_O$가 최초로 진입하여($R^2 = 0.768$)
관측성이 1차적 설명 변수임을 보여준다.
반면, C3에서는 불확실성 $U$가 최초로 진입하여($R^2 = 0.424$)
파레토 불확실성이 지배적 설명 변수로 전환됨을 나타낸다.
이러한 \textbf{지배 변수의 전환}은 분포의 꼬리 특성이
모형 구조에 질적 변화를 야기한다는 것을 의미한다.

C2에서는 2개 변수 투입($k_O$, $U$)만으로 $R^2 = 0.931$에 도달하여
핵심 구조가 상대적으로 단순한 반면,
C3에서는 4개 변수 투입 후에도 $R^2 = 0.839$에 머물러
파레토 분포의 비선형적 복잡성을 보여준다.

\p{[}그림~\vref{fig:forward_selection_r2}\p{]}은 변수 투입 수에 따른 누적 $R^2$의
진행을 비교한 것이다.

\begin{figure}[htbp]
  \centering
  \begin{tikzpicture}
    \begin{axis}[
      width=0.88\textwidth,
      height=7cm,
      xlabel={투입 변수 수},
      ylabel={누적 $R^2$},
      xmin=0.5, xmax=9.5,
      ymin=0.3, ymax=1.05,
      xtick={1,2,3,4,5,6,7,8,9},
      ymajorgrids=true,
      grid style={dashed, black},
      legend style={at={(0.98,0.15)}, anchor=south east, font=\small,
        fill=white, fill opacity=0.9, draw=black},
      ]
      % C2 (Gaussian)
      \addplot[blue, thick, mark=o, mark size=3pt] coordinates {
        (1, 0.768) (2, 0.931) (3, 0.973) (4, 0.996) (5, 0.997)
        (6, 0.998) (7, 0.998) (8, 0.998) (9, 0.998)
      };
      \addlegendentry{C2 (Gaussian)}

      % C3 (Pareto)
      \addplot[red, thick, mark=square*, mark size=3pt] coordinates {
        (1, 0.424) (2, 0.686) (3, 0.799) (4, 0.839) (5, 0.853)
        (6, 0.856) (7, 0.857) (8, 0.858) (9, 0.858)
      };
      \addlegendentry{C3 (Pareto)}

      % H4 threshold (전체 수평선)
      \draw[dashed, black, thin] (rel axis cs:0,0.867) -- (rel axis cs:1,0.867);
      \node[font=\scriptsize, fill=white, inner sep=1.5pt,
        anchor=north east] at (rel axis cs:0.99,0.867)
        {H4 판정 기준 ($R^2 \geq 0.95$)};
    \end{axis}
  \end{tikzpicture}
  \caption{전진 선택법 누적 $R^2$ 진행 비교}
  \label{fig:forward_selection_r2}
\end{figure}


\chg{다음으로 LASSO 정규화에 의한 변수 선택을 살펴본다.} 전진 선택법의 결과를 독립적으로 검증하기 위해 LASSO(Least Absolute
Shrinkage and Selection Operator) 정규화를 적용하였다.
LASSO는 $L_1$ 패널티($\lambda$)를 부과하여 기여도가 낮은 변수의
회귀계수를 정확히 0으로 수축시킴으로써,
전진 선택법과 \textbf{상이한 원리}로 핵심 변수를 식별한다.
\p{[}표~\vref{tab:lasso_results}\p{]}는 정규화 계수 $\lambda$에 따른 생존 변수 수와
$R^2$를 요약한 것이다.

\begin{table}[htbp]
  \centering
  \caption{$\lambda$에 따른 LASSO 생존 변수와 $R^2$}
  \label{tab:lasso_results}
  \small
  \begin{tabular}{c cc cc}
    \toprule
    & \multicolumn{2}{c}{\textbf{C2 (Gaussian)}} & \multicolumn{2}{c}{\textbf{C3 (Pareto)}} \\
    \cmidrule(lr){2-3} \cmidrule(lr){4-5}
    $\lambda$ & 생존 변수 & $R^2$ & 생존 변수 & $R^2$ \\
    \midrule
    0.00015 & 7 & 0.998 & 9 & 0.858 \\
    0.0005 & 6 & 0.998 & 9 & 0.858 \\
    0.001 & 7 & 0.998 & 9 & 0.857 \\
    0.005 & 5 & 0.997 & 7 & 0.856 \\
    0.01 & 5 & 0.995 & 6 & 0.852 \\
    \bottomrule
  \end{tabular}
\end{table}

LASSO 결과는 전진 선택법과 일관된 패턴을 보여준다.
C2에서는 $\lambda = 0.01$에서도 5개 변수($k_O$, $U$, $k_O^2$, $k_O \cdot O$,
$k_O \cdot U$)만으로 $R^2 = 0.995$를 유지하여,
관측성-불확실성 교호작용 중심의 핵심 구조가 확인된다.
C3에서는 동일한 $\lambda = 0.01$에서 6개 변수가 생존하며
$R^2 = 0.852$에 그쳐, 더 많은 변수가 필요하되
설명력은 상대적으로 제한적이다.
이는 C3의 파레토 구조가 2차 다항식 근사의 한계를
벗어나는 비선형성을 포함함을 재확인한다.

\p{[}그림~\vref{fig:lasso_path}\p{]}는 $\lambda$에 따른 $R^2$ 경로를 비교한 것이다.

\begin{figure}[htbp]
  \centering
  \begin{tikzpicture}
    \begin{axis}[
      width=0.88\textwidth,
      height=6.5cm,
      xlabel={정규화 계수 $\lambda$},
      ylabel={$R^2$},
      xmode=log,
      xmin=0.0001, xmax=0.02,
      ymin=0.84, ymax=1.005,
      ymajorgrids=true,
      grid style={dashed, black},
      legend style={at={(0.98,0.5)}, anchor=east, font=\small,
        fill=white, fill opacity=0.9, draw=black},
      ]
      % C2 (Gaussian)
      \addplot[blue, thick, mark=o, mark size=3pt] coordinates {
        (0.00015, 0.998) (0.0005, 0.998) (0.001, 0.998)
        (0.005, 0.997) (0.01, 0.995)
      };
      \addlegendentry{C2 (Gaussian)}

      % C3 (Pareto)
      \addplot[red, thick, mark=square*, mark size=3pt] coordinates {
        (0.00015, 0.858) (0.0005, 0.858) (0.001, 0.857)
        (0.005, 0.856) (0.01, 0.852)
      };
      \addlegendentry{C3 (Pareto)}

      % H4 threshold (전체 수평선)
      \draw[dashed, black, thin] (rel axis cs:0,0.667) -- (rel axis cs:1,0.667);
      \node[font=\scriptsize, fill=white, inner sep=1.5pt,
        anchor=south east] at (rel axis cs:0.99,0.667)
        {H4 판정 기준 ($R^2 \geq 0.95$)};
    \end{axis}
  \end{tikzpicture}
  \caption{LASSO $R^2$ 경로 (C2 vs.\ C3)}
  \label{fig:lasso_path}
\end{figure}


\chg{끝으로 두 방법의 교차검증 합의로 간명 모형을 도출한다.} 전진 선택법과 LASSO는 각각 단계별 최적(stepwise) 진입 순서와
$L_1$ 정규화라는 상이한 원리로 변수를 선택한다.
단일 방법론의 편향을 최소화하기 위해, 두 방법의
\textbf{교집합(consensus)}을 합의 변수로 채택하였다.
구체적으로, 전진 선택법 상위 4개 변수 집합과
LASSO 정규화($\lambda = 0.005$)에서의 생존 변수 집합의
교집합을 선정하였다.

\chg{C2의 합의 변수는 전진 선택법 상위 4개($k_O$, $U$, $k_O \cdot U$, $k_O \cdot O$)와 LASSO 생존
5개($k_O$, $U$, $k_O^2$, $k_O \cdot O$, $k_O \cdot U$)의 교집합인 \{$k_O$, $U$, $k_O \cdot U$, $k_O \cdot O$\}
4개로서, 전진 선택에 없는 $k_O^2$가 LASSO에서만 생존하여 탈락한다. C3의 합의 변수는 전진 선택법
상위 4개($U$, $k_O$, $U^2$, $k_O \cdot U$)와 LASSO 생존 7개($k_O$, $U$, $k_O^2$, $U^2$, $k_O \cdot O$,
$k_O \cdot U$, $O \cdot U$)의 교집합인 \{$U$, $k_O$, $U^2$, $k_O \cdot U$\} 4개로서, $k_O^2$, $k_O \cdot O$,
$O \cdot U$는 LASSO에서만 생존하여 탈락한다.}

합의 변수를 사용한 4변수 간명 모형의 회귀식은 다음과 같다\chg{.}

\begin{align}
  \hat{\eta}_{\text{C2}} &= 1.225
    - 0.291\,k_O + 0.467\,U
    - 0.433\,(k_O \cdot U) - 0.562\,(k_O \cdot O)
    & (R^2 = 0.996) \label{eq:c2_parsimonious} \\
  \hat{\eta}_{\text{C3}} &= 1.574
    + 0.240\,U - 1.053\,k_O
    - 0.002\,U^2 - 0.122\,(k_O \cdot U)
    & (R^2 = 0.839) \label{eq:c3_parsimonious}
\end{align}

간명 모형은 9변수 완전 모형(C2: $0.998$, C3: $0.858$) 대비
4개 변수만으로 C2의 $99.6\%$, C3의 $83.9\%$를 설명하여,
설명력 손실을 최소화하면서 해석 가능성을 확보한다.
C2에서는 관측성 효과가 $k_O$($-0.291$), $k_O \cdot O$($-0.562$),
$k_O \cdot U$($-0.433$)의 3개 항에 분산되어 있어,
관측성과 불확실성의 교호작용 구조가 세분화된 형태로 포착된다.
C3에서는 $k_O \cdot O$ 항이 선택되지 않아 $k_O$($-1.053$) 단일 항에
관측성 효과가 집중되며, $U^2$ 항(계수 $-0.002$)이 포함되어
파레토 분포의 비선형 구조가 반영된다.
C3에서 $k_O$ 계수의 절대값이 C2보다 큰 것은 관측성 영향력의 증대가
아니라, C2에서 교호작용 항에 분산된 효과가 C3에서 단일 항에
집중된 결과이다.
이는 파레토 극단 영역에서 관측성의 한계적 효과가 감쇠하는
이중채널 감쇠 현상(\p{[}표~\vref{tab:c3_split_regression}\p{]}, $0.52$배 감쇠)과 일관된다.

\subsection{\mrev{회귀분석 가설 검증 결과}}
\label{subsec:ch4_h4_h5_conclusion}



%% [분량보존 복원] 9변수 완전 모형 전체 계수 및 LASSO 경로 상세 (구 부록 G)
\subsubsection{9변수 완전 모형의 전체 계수와 LASSO 경로}
\label{subsec:full-coeff-detail}
4.4.6(전진 선택법과 LASSO 정규화를 통한 변수 선택)에서 요약한 9변수 완전 2차 모형과 LASSO 변수 선택의 전체 계수를 제시하는데, \m{(i)} \p{[}표~\ref{tab:app-c2-full-quad}\p{]}과 \p{[}표~\ref{tab:app-c3-full-quad}\p{]}는 C2와 C3 각 조건의 \chg{9변수 계수와 $t$-통계량, 전진 선택 진입 순서}를, \m{(ii)} \p{[}표~\ref{tab:app-c2-lasso-path}\p{]}과 \p{[}표~\ref{tab:app-c3-lasso-path}\p{]}은 정규화 강도 $\lambda$별 \chg{생존 변수와 제거 변수}를 보인다.

\chg{먼저 C2 조건의} 9변수 완전 2차 모형($R^2 = 0.998$)에서 각 변수의 회귀계수,
$t$-통계량, 전진 선택법 진입 순서를 \p{[}표~\ref{tab:app-c2-full-quad}\p{]}에 제시한다.

C2의 9변수 완전 모형에서 가장 큰 $|t|$를 보이는 변수는
$k_O \cdot U$ ($|t| = 476.1$)로, 관측성 가중치와 불확실성의
교호 작용이 C2의 $\eta$ 변동에서 가장 결정적인 역할을 한다.
$U$의 주효과($|t| = 415.4$)가 그 다음으로 크며,
$k_O \cdot O$ ($|t| = 170.8$)가 뒤를 잇는다.
한편 $k_O$는 전진 선택에서 1순위로 진입하지만($k_O$ 단독 시 관측성 효과를 포착),
완전 모형에서는 교호 작용항($k_O \cdot U$, $k_O \cdot O$)에 주효과가 흡수되어
$|t| = 2.4$로 비유의해진다.
$O \cdot U$ ($|t| = 1.4$) 역시 완전 모형에서 비유의하다.

\begin{table}[htbp]
\centering
\caption{C2 전체 9변수 2차 회귀 계수 ($R^2 = 0.998$)}
\label{tab:app-c2-full-quad}
\small
\begin{tabular}{lrrl}
\toprule
\textbf{변수} & \textbf{계수 $\hat{\beta}$} & \textbf{$t$-통계량} & \textbf{진입 순서} \\
\midrule
절편 & $1.082$ & $379.4$ & --- \\
$k_O$ & $0.012$ & $2.4$ & 1 \\
$O$ & $-0.034$ & $-7.0$ & 7 \\
$U$ & $0.542$ & $415.4$ & 2 \\
$k_O^2$ & $-0.307$ & $-86.5$ & 5 \\
$O^2$ & $0.019$ & $5.3$ & 8 \\
$U^2$ & $-0.013$ & $-67.3$ & 6 \\
$k_O \cdot O$ & $-0.545$ & $-170.8$ & 4 \\
$k_O \cdot U$ & $-0.433$ & $-476.1$ & 3 \\
$O \cdot U$ & $0.001$ & $1.4$ & 9 \\
\bottomrule
\end{tabular}
\end{table}



\chg{이어 C3 조건의} 9변수 완전 2차 모형($R^2 = 0.858$)에서 각 변수의 회귀계수,
$t$-통계량, 전진 선택법 진입 순서를 \p{[}표~\ref{tab:app-c3-full-quad}\p{]}에 제시한다.

\begin{table}[htbp]
\centering
\caption{C3 전체 9변수 2차 회귀 계수 ($R^2 = 0.858$)}
\label{tab:app-c3-full-quad}
\small
\begin{tabular}{lrrl}
\toprule
\textbf{변수} & \textbf{계수 $\hat{\beta}$} & \textbf{$t$-통계량} & \textbf{진입 순서} \\
\midrule
절편 & $1.430$ & $88.6$ & --- \\
$k_O$ & $-0.216$ & $-4.7$ & 2 \\
$O$ & $0.085$ & $1.9$ & 8 \\
$U$ & $0.256$ & $126.6$ & 1 \\
$k_O^2$ & $-0.513$ & $-12.7$ & 6 \\
$O^2$ & $0.022$ & $0.5$ & 9 \\
$U^2$ & $-0.002$ & $-85.9$ & 3 \\
$k_O \cdot O$ & $-0.627$ & $-17.4$ & 5 \\
$k_O \cdot U$ & $-0.126$ & $-53.7$ & 4 \\
$O \cdot U$ & $-0.025$ & $-11.4$ & 7 \\
\bottomrule
\end{tabular}
\end{table}

C3에서는 $U$가 최대 $|t| = 126.6$으로 1순위 진입하여,
C2와 달리 불확실성의 주효과가 지배적임을 보여준다.
$U^2$ ($|t| = 85.9$)이 3순위로 진입하여 파레토 분포의
비선형적 $U$ 효과가 2차 항으로 포착됨을 확인한다.
$O^2$ ($|t| = 0.5$)와 $O$ ($|t| = 1.9$)는 사실상 유의하지 않으며,
C3에서 관측성의 개별 효과보다 $k_O \cdot O$ 교호 작용이
훨씬 중요함을 시사한다.


\chg{다음으로 LASSO 경로를 상세히 본다.} 정규화 강도 $\lambda$를 증가시키면서 제거되는 변수와 잔존 $R^2$의 변화를
C2(\p{[}표~\ref{tab:app-c2-lasso-path}\p{]})와 C3(\p{[}표~\ref{tab:app-c3-lasso-path}\p{]}) 조건 각각에 대해 제시한다.

\begin{table}[htbp]
\centering
\caption{C2 LASSO 경로의 $\lambda$별 생존 변수와 제거 변수}
\label{tab:app-c2-lasso-path}
\small
\begin{tabular}{c c l l}
\toprule
$\lambda$ & $R^2$ & 생존 변수 & 제거 변수 \\
\midrule
0.00015 & 0.998 & $O$, $U$, $k_O^2$, $O^2$, $U^2$, $k_O \cdot O$, $k_O \cdot U$ (7) & $k_O$, $O \cdot U$ \\
0.0005 & 0.998 & $O$, $U$, $k_O^2$, $U^2$, $k_O \cdot O$, $k_O \cdot U$ (6) & $k_O$, $O^2$, $O \cdot U$ \\
0.001 & 0.998 & $k_O$, $O$, $U$, $k_O^2$, $U^2$, $k_O \cdot O$, $k_O \cdot U$ (7) & $O^2$, $O \cdot U$ \\
0.005 & 0.997 & $k_O$, $U$, $k_O^2$, $k_O \cdot O$, $k_O \cdot U$ (5) & $O$, $O^2$, $U^2$, $O \cdot U$ \\
0.01 & 0.995 & $k_O$, $U$, $k_O^2$, $k_O \cdot O$, $k_O \cdot U$ (5) & $O$, $O^2$, $U^2$, $O \cdot U$ \\
\bottomrule
\end{tabular}
\end{table}

\begin{table}[htbp]
\centering
\caption{C3 LASSO 경로의 $\lambda$별 생존 변수와 제거 변수}
\label{tab:app-c3-lasso-path}
\small
\begin{tabular}{c c l l}
\toprule
$\lambda$ & $R^2$ & 생존 변수 & 제거 변수 \\
\midrule
0.00015 & 0.858 & 전체 9개 & --- \\
0.0005 & 0.858 & 전체 9개 & --- \\
0.001 & 0.857 & 전체 9개 & --- \\
0.005 & 0.856 & $k_O$, $U$, $k_O^2$, $U^2$, $k_O \cdot O$, $k_O \cdot U$, $O \cdot U$ (7) & $O$, $O^2$ \\
0.01 & 0.852 & $k_O$, $U$, $k_O^2$, $U^2$, $k_O \cdot O$, $k_O \cdot U$ (6) & $O$, $O^2$, $O \cdot U$ \\
\bottomrule
\end{tabular}
\end{table}

C2에서는 $\lambda = 0.005$ 이상에서 $O \cdot U$, $O^2$, $O$, $U^2$가 순차적으로
제거되어도 $R^2$가 $0.997$ 이상을 유지하는 반면,
C3에서는 $\lambda = 0.005$에서야 처음으로 변수 제거($O$, $O^2$)가 시작되며,
$R^2$ 감소 폭도 미미하다($0.858 \to 0.856$).
이는 C3의 9변수 모형이 C2보다 덜 과적합(over-fitted)되어 있으며,
파레토 분포의 복잡한 비선형성을 포착하기 위해 더 많은 변수가 구조적으로
필요함을 보여준다.


%% [코드릴리스→본문 복원] C1 완전계수표 + C3 Tail 전체계수표 (구 부록 F 유일분)
\chg{C1 조건의 2차 회귀 전체 계수는 다음과 같다.} C1 조건은 불확실성이 없는 순수 관측성 조건으로서, $U$ 관련 항($\beta_3$, $\beta_6$, $\beta_8$, $\beta_9$)이
제외된 5개 계수의 2차 회귀 결과를 \p{[}표~\ref{tab:app-c1-quad}\p{]}에 제시한다.
시그모이드 입력 $z_O = \kappa \cdot k_O \cdot O_{\text{raw}}$가 $\kappa = 1.2$ 영역에서
사실상 선형으로 작동하므로, $k_O \cdot O$ 교호작용 단일 항이 $\eta_{C1}$의 변동을 결정론적으로 포착한다.

\begin{table}[htbp]
\centering
\caption{C1 조건 2차 회귀 전체 계수 ($n = 10{,}000$)}
\label{tab:app-c1-quad}
\begin{tabular}{lrrrl}
\toprule
\textbf{항} & \textbf{계수} & \textbf{SE} & \textbf{$t$} & \textbf{비고} \\
\midrule
\m{절편($\beta_0$)}      & 1.000  & $<$0.001 & ---     & 이론치 1.000 \\
$k_O$ ($\beta_1$)     & 0.000  & $<$0.001 & ---     & 주 효과 흡수됨 \\
$O$ ($\beta_2$)       & 0.000  & $<$0.001 & ---     & 주 효과 흡수됨 \\
$k_O^2$ ($\beta_4$)   & 0.000  & $<$0.001 & ---     & 비유의 \\
$O^2$ ($\beta_5$)     & 0.000  & $<$0.001 & ---     & 비유의 \\
$k_O \cdot O$ ($\beta_7$) & \erf{$-0.30$}{$-0.546$} & $<$0.001 & ---  & \textbf{핵심 교호작용} \\
\midrule
$R^2$                 & \multicolumn{3}{c}{\textbf{1.000}} & 결정론적 구조 완전 포착 \\
\bottomrule
\end{tabular}
\end{table}

\chg{이 결과는 C1 조건이 $\eta_{C1} \approx 1 - \erf{0.30}{0.546} \cdot k_O \cdot O_{\text{raw}}$로 근사됨을 보여주며,
이는 시그모이드 입력 $z_O = \kappa \cdot k_O \cdot O_{\text{raw}}$의 선형 근사가 $\kappa = 1.2$ 영역에서
사실상 정확함(비선형 오차 $\approx 0$)을 반영한다.}

\chg{끝으로 Tail 2차 회귀 전체 계수를 제시한다.} P85.8 이상의 Tail 영역($n \approx 1{,}420$)에서 $\ln(U)$ 변환을 적용한
2차 회귀의 전체 계수를 \p{[}표~\ref{tab:app-c3-tail-quad}\p{]}에 제시한다.
$k_O \cdot \ln U$ 교호작용이 꼬리 영역에서의
관측성 효과 변화를 설명하는 핵심 항이다.

\begin{table}[htbp]
\centering
\caption{C3 Tail 영역 2차 회귀 전체 계수 ($>$ P85.8)}
\label{tab:app-c3-tail-quad}
\begin{tabular}{lrl}
\toprule
\textbf{항} & \textbf{계수} & \textbf{비고} \\
\midrule
절편          & 0.50   & \\
$k_O$        & 2.42   & 양의 주 \m{효과(교호작용과 결합)} \\
$\ln U$      & 0.10   & 로그 변환된 불확실성 \\
$k_O \cdot \ln U$ & $-1.25$ & \textbf{핵심 \m{교호작용(증폭 원인)}} \\
\midrule
$R^2$        & 0.710  & \\
\bottomrule
\end{tabular}
\end{table}

\chg{\p{[}표~\ref{tab:app-c3-tail-quad}\p{]}에서 $k_O \cdot \ln U$ 교호작용($-1.25$)이 꼬리 영역에서 효과의 수학적 원인이며,
$\ln U$가 커질수록(극단 불확실성) $k_O$의 효과가 비선형적으로 변화한다.}


\chg{회귀 설명력에 관한 H4의 검증 결론을 정리한다.} \mrev{H4는 $R^2_{(2\mathrm{q})} \geq 0.95$를 기준으로 한다.
\chg{조건별 $R^2$는 C1 1.000(충족), C2 0.998(충족), C3 전체 0.858(단일 회귀 시 미달)이다.}
C3는 P85.8 분할 시 Core $R^2 = 0.999$로 충족하여, 분할 회귀 적용
조건부로 \textbf{H4는 채택된다}.}

\chg{이어 교호작용 유의성에 관한 H5의 검증 결론은 다음과 같다.} \mrev{H5는 $\Delta R^2$ 기여도 유의를 기준으로 한다. C1에서 $k_O \cdot O$
교호작용항이 $R^2$를 $0.869 \to 1.000$로 향상시켜 $\Delta R^2 = 0.131$
($p < 0.001$). C2 전진 선택법에서 $k_O \cdot U$와 $k_O \cdot O$가 3$\sim$4
순위로 진입하며, LASSO ($\lambda = 0.01$)에서도 두 항 모두 생존한다.
서로 다른 두 변수 선택 방법이 일관되게 교호작용항을 핵심으로 식별하여,
\textbf{H5는 채택된다}.}
          % 4.4 회귀분석


%% ===================================================================
\section{\vrev{기술통계 종합}}
\label{sec:descriptive_summary}
%% ===================================================================

\p{[}표~\vref{tab:descriptive_summary}\p{]}는 C0--C3 4개 조건의 효율성 비율 $\eta$에
대한 기술통계량을 종합한 것이다.
\chg{여기서 C0는 표준편차가 0이어서 왜도와 초과 첨도가 정의되지 않으므로 ``---''로 표기하였다.}

\begin{table}[htbp]
  \centering
  \caption{효율성 비율 $\eta$의 조건별 기술통계 종합}
  \label{tab:descriptive_summary}
  \small
  \begin{tabular}{lccccccc}
    \toprule
    조건 & $\bar{\eta}$ & SD & $\DRTO$ (h) & 최솟값 & 최댓값 & 왜도 & 초과 첨도 \\
    \midrule
    C0 (정적) & 1.000 & 0.000 & 6.000 & 1.000 & 1.000 & --- & --- \\
    C1 (관측성) & 0.853 & 0.125 & 5.115 & 0.468 & 1.000 & $-0.84$ & $-0.20$ \\
    C2 (Gaussian) & 1.682 & 0.615 & 10.094 & 0.477 & 3.443 & $0.26$ & $-0.77$ \\
    C3 (Pareto) & 1.514 & 0.791 & 9.085 & 0.472 & 4.000 & $1.32$ & $1.03$ \\
    \bottomrule
  \end{tabular}
\end{table}

네 조건의 기술통계에서 다음의 구조적 패턴이 도출된다.

\chg{첫째, 평균의 2단계 분기 구조를 보면,} C2와 C3는 모두 C1에서 분기하는 병렬 경로이므로,
$\eta$ 평균의 변화를 공통 경로와 분기 경로로 구분하여 해석한다.
\chg{공통 경로인 C0($1.000$)에서 C1($0.853$)로는 관측성 도입에 의해 평균이 $14.7\%$ 감소하여 복구 시간이
단축된다. 분기 A인 C1($0.853$)에서 C2($1.682$)로는 가우시안 불확실성 추가 시 불확실성 프리미엄에 의해
$\erf{97.2}{97.3}\%$ 증가하며, 분기 B인 C1($0.853$)에서 C3($1.514$)로는 파레토 불확실성 추가 시 $77.5\%$ 증가하되
C2보다 낮은 평균을 산출한다.}

C2와 C3의 평균 차이($1.682$ vs.\ $1.514$)는 분포 특성의 차이를 반영한다.
가우시안 불확실성은 대칭적으로 $\eta$를 상향시키는 반면,
파레토 불확실성은 소수 극단값에 질량이 집중되어
평균은 C2보다 낮으나 분산과 극단값 범위는 더 크다.

\chg{둘째, 표준편차의 단조 증가를 보면,} SD는 공통 경로에서 C0($0$) $\to$ C1($0.125$)로 증가하며,
분기 후 C2($0.615$) 및 C3($0.791$) 모두 추가로 증가한다.
C3의 SD가 C2의 $1.29$배인 것은 파레토 분포의
극단값이 분산을 확대시킨다는 것을 정량적으로 보여준다.

\chg{셋째, 분포 형태의 질적 전환을 보면,} 왜도(skewness)의 변화가 주목할 만하다. C1은 음의 왜도($-0.84$)를
보여 $\eta$가 좌측으로 치우치며, C2는 왜도가 거의 $0$에 가까워
($0.26$) 약한 양의 비대칭을 보인다. 반면 C3는 강한 양의 왜도($1.32$)를
보여 우측 꼬리가 길게 늘어진 비대칭 분포를 형성한다.
첨도(kurtosis)도 C3에서 $1.03$으로 가장 높아, 극단값의 발생 빈도가
정규분포 대비 높음을 나타낸다.

\p{[}그림~\vref{fig:eta_distribution_all}\p{]}은 4개 조건의 $\eta$ 분포를 밀도 곡선으로
중첩하여 비교한 것이다.
\chg{그림에서 C1은 관측성 단축에 의한 좌측 편향을, C2는 평균 중심 분포(왜도 $\approx 0$)의 상향(불확실성
프리미엄)을, C3는 우측 파레토 꼬리에 의한 극단 위험 확대를 보인다. 여기서 왜도(skewness)는 분포의 비대칭
정도를 나타내어 양의 왜도는 오른쪽 꼬리가, 음의 왜도는 왼쪽 꼬리가 긴 분포를 의미하며, 초과 첨도(excess
kurtosis)는 정규분포 대비 꼬리의 무거움을 나타내어 양수는 뾰족하고 극단값이 빈번한 분포(leptokurtic),
음수는 납작하고 극단값이 드문 분포(platykurtic)를 가리킨다. C2의 입력은 가우시안이나 시그모이드 바운딩이
극단값을 $(-1,\;+1)$ 범위로 압축하여 출력 분포의 꼬리가 정규분포보다 가벼워지므로 첨도가 음수($-0.77$)로
나타나는 반면, C3는 파레토 입력의 극단값이 시그모이드 압축을 거쳐도 꼬리 특성이 잔존하여 첨도가
양수($+1.03$)를 유지한다.}

\begin{figure}[htbp]
  \centering
  \begin{tikzpicture}
    \begin{axis}[
      width=0.92\textwidth,
      height=8cm,
      xlabel={효율성 비율 $\eta$},
      ylabel={밀도 (density)},
      xmin=0.2, xmax=4.2,
      ymin=0,
      ymajorgrids=true,
      grid style={dashed, black},
      legend style={at={(0.98,0.95)}, anchor=north east, font=\small},
      legend cell align=left,
      ]
      % C0: vertical line (Dirac delta, η=1.000 고정)
      \addplot[black, ultra thick, dashed]
        coordinates {(1.0,0) (1.0,7.0)};
      \addlegendentry{C0 ($\eta = 1.000$)}

      % C1: 실데이터 KDE — 경계 보정(reflection at η=1.0)
      \addplot[blue, thick, smooth] coordinates {
        (0.425,0.005) (0.445,0.022) (0.466,0.062) (0.486,0.130) (0.506,0.219)
        (0.526,0.320) (0.546,0.425) (0.566,0.534) (0.587,0.639) (0.607,0.734)
        (0.627,0.831) (0.647,0.936) (0.667,1.049) (0.687,1.179) (0.708,1.313)
        (0.728,1.433) (0.748,1.550) (0.768,1.687) (0.788,1.857) (0.808,2.045)
        (0.829,2.221) (0.849,2.388) (0.869,2.580) (0.889,2.841) (0.909,3.196)
        (0.929,3.678) (0.950,4.377) (0.960,4.820) (0.970,5.292) (0.980,5.728)
        (0.990,6.040) (1.000,6.154)
      };
      \addlegendentry{C1 ($\bar{\eta} = 0.853$)}

      % C2: 실데이터 KDE (n=10,000, Scott bandwidth)
      \addplot[green!60!black, thick, smooth] coordinates {
        (0.320,0.004) (0.401,0.017) (0.481,0.053) (0.561,0.114) (0.641,0.194)
        (0.722,0.277) (0.802,0.356) (0.882,0.428) (0.963,0.486) (1.043,0.523)
        (1.123,0.544) (1.203,0.555) (1.284,0.553) (1.364,0.545) (1.444,0.535)
        (1.524,0.531) (1.605,0.531) (1.685,0.534) (1.765,0.526) (1.845,0.506)
        (1.926,0.486) (2.006,0.474) (2.086,0.459) (2.167,0.433) (2.247,0.406)
        (2.327,0.383) (2.407,0.359) (2.488,0.328) (2.568,0.292) (2.648,0.253)
        (2.728,0.211) (2.809,0.167) (2.889,0.127) (2.969,0.093) (3.049,0.066)
        (3.130,0.046) (3.210,0.031) (3.290,0.019) (3.371,0.011) (3.451,0.005)
      };
      \addlegendentry{C2 ($\bar{\eta} = 1.682$)}

      % C3: 실데이터 KDE (n=10,000, Scott bandwidth)
      \addplot[red, thick, smooth] coordinates {
        (0.240,0.004) (0.347,0.024) (0.454,0.089) (0.561,0.219) (0.668,0.402)
        (0.775,0.616) (0.882,0.823) (0.936,0.896) (0.989,0.932) (1.043,0.927)
        (1.096,0.885) (1.150,0.818) (1.203,0.742) (1.257,0.667) (1.310,0.598)
        (1.364,0.537) (1.417,0.484) (1.471,0.439) (1.524,0.400) (1.578,0.368)
        (1.631,0.342) (1.685,0.319) (1.738,0.298) (1.792,0.277) (1.845,0.256)
        (1.899,0.238) (1.953,0.221) (2.006,0.208) (2.060,0.197) (2.113,0.188)
        (2.167,0.179) (2.274,0.164) (2.381,0.151) (2.488,0.137) (2.595,0.121)
        (2.702,0.110) (2.809,0.104) (2.916,0.098) (3.023,0.092) (3.130,0.086)
        (3.237,0.081) (3.344,0.077) (3.451,0.074) (3.558,0.071) (3.665,0.069)
        (3.772,0.068) (3.879,0.063) (3.986,0.044) (4.093,0.019) (4.200,0.005)
      };
      \addlegendentry{C3 ($\bar{\eta} = 1.514$)}

      % eta=1 reference line
      \draw[black, thin, dashed] (axis cs:1.0,0) -- (axis cs:1.0,7);

      % 왜도·첨도 주석
      \node[blue, font=\scriptsize, align=left, anchor=west]
        at (axis cs:0.35,3.5) {왜도 $= -0.84$\\첨도 $= -0.20$};
      \draw[blue, thin, ->] (axis cs:0.52,3.3) -- (axis cs:0.62,1.0);

      \node[green!60!black, font=\scriptsize, align=left, anchor=west]
        at (axis cs:2.5,1.2) {왜도 $= 0.26$\\첨도 $= -0.77$};
      \draw[green!60!black, thin, ->] (axis cs:2.6,1.05) -- (axis cs:2.2,0.45);

      \node[red, font=\scriptsize, align=left, anchor=west]
        at (axis cs:1.6,2.0) {왜도 $= 1.32$\\첨도 $= 1.03$};
      \draw[red, thin, ->] (axis cs:1.8,1.85) -- (axis cs:1.7,0.30);
    \end{axis}
  \end{tikzpicture}
  \caption{C0/C1/C2/C3 효율성 비율 $\eta$ 분포 비교}
  \label{fig:eta_distribution_all}
\end{figure}
 % 4.5 기술통계 종합


%% ===================================================================
\section{\vrev{수렴 진단 및 재현성}}
\label{sec:convergence}
%% ===================================================================

%% -------------------------------------------------------------------
\subsection{\vrev{표본 크기 적정성}}
\label{subsec:sample_size}
%% -------------------------------------------------------------------

$n = 10{,}000$의 표본 크기가 결과의 안정성을 보장하는지 확인하기 위해
부분 표본(subsample) 분석을 수행하였다. $n = 1{,}000$, $2{,}000$, $5{,}000$,
$10{,}000$에서의 핵심 통계량 수렴 양상을 \p{[}표~\vref{tab:convergence}\p{]}에 제시한다.

\begin{table}[htbp]
  \centering
  \caption{표본 크기별 핵심 통계량 \m{수렴($s = 42$, $\kappa = 1.2$)}}
  \label{tab:convergence}
  \small
  \begin{tabular}{lcccc}
    \toprule
    통계량 & $n = 1{,}000$ & $n = 2{,}000$ & $n = 5{,}000$ & $n = 10{,}000$ \\
    \midrule
    $\bar{\eta}_{\text{C1}}$ & 0.851 & 0.848 & 0.853 & 0.853 \\
    $\bar{\eta}_{\text{C2}}$ & 1.685 & 1.669 & 1.683 & 1.682 \\
    $\bar{\eta}_{\text{C3}}$ & 1.526 & 1.496 & 1.523 & 1.514 \\
    Cohen's $d$ (C1 vs C0) & $-1.164$ & $-1.173$ & $-1.171$ & $-1.180$ \\
    C2 2차 회귀 $R^2$ & 0.998 & 0.998 & 0.998 & 0.998 \\
    C3 2차 회귀 $R^2$ (전체) & 0.902 & 0.890 & 0.875 & 0.858 \\
    \midrule
    $\mathrm{CVaR}_{99}$ C2 & 3.196 & 3.185 & 3.149 & 3.180 \\
    $\mathrm{CVaR}_{99}$ C3 & 3.957 & 3.946 & 3.941 & 3.944 \\
    $\mathrm{CVaR}_{99}$ 격차 & $+23.8\%$ & $+23.9\%$ & $+25.2\%$ & $+24.0\%$ \\
    \bottomrule
  \end{tabular}
\end{table}

$n = 2{,}000$ 이상에서 \chg{C1, C2} 평균과 효과 크기가 소수점 셋째 자리까지 안정되며,
C2 2차 회귀 $R^2$는 전 구간에서 $0.998$로 일정하다.
C3 2차 회귀 $R^2$(전체)는 표본 크기 증가에 따라
$0.902 \to 0.858$로 점진 감소 후 안정화된다.
작은 표본에서는 극단값이 충분히 표집되지 않아 $R^2$가 과대추정된다.
표본 크기가 커지면 파레토 꼬리의 비선형성이 회귀에 반영되면서
$R^2$가 점진적으로 감소하여 본래의 설명력 수준으로 수렴한다.
$\mathrm{CVaR}_{99}$는 상위 $1\%$ 표본의 조건부 평균으로, $n = 1{,}000$에서도
C2($3.196$)와 C3($3.957$)의 격차가 $+23.8\%$로 $n = 10{,}000$의
$+24.0\%$와 유사하나, 상위 표본 수가 10개에 불과하여 추정 변동이 크다.
$n = 10{,}000$에서는 상위 100개 표본에 기반하여 $\mathrm{CVaR}_{99}$의 안정적
추정이 가능하므로, 본 연구의 표본 크기가 적절함을 확인하였다.


%% -------------------------------------------------------------------
\subsection{\vrev{다중 시드 재현성}}
\label{subsec:multi_seed}
%% -------------------------------------------------------------------

마스터 시드 $s = 42$의 결과가 특정 시드에 종속되지 않는지 확인하기 위해,
$s = 0, 1, \ldots, 9{,}999$의 $10{,}000$개 마스터 시드에 대해
동일한 시뮬레이션을 반복 수행하였다(총 $10^8$개 표본).
핵심 통계량의 시드 간 변동은 \p{[}표~\vref{tab:multi_seed_stability}\p{]}에 요약한다.

\begin{table}[htbp]
  \centering
  \caption{다중 시드 시뮬레이션 결과 안정성}
  \label{tab:multi_seed_stability}
  \small
  \begin{tabular}{lcccc}
    \toprule
    통계량 & 시드 간 평균 & 시드 간 SD & 기준 ($s=42$) & 판정 \\
    \midrule
    $\bar{\eta}_{\text{C1}}$ & \erf{0.853}{0.854} & 0.001 & 0.853 & \erf{평균 이내}{2 SD 이내} \\
    $\bar{\eta}_{\text{C2}}$ & \erf{1.682}{1.692} & 0.006 & 1.682 & \erf{평균 이내}{2 SD 이내} \\
    $\bar{\eta}_{\text{C3}}$ & \erf{1.514}{1.521} & 0.008 & 1.514 & \erf{평균 이내}{2 SD 이내} \\
    CDF 교차점 백분위 (P85.8) & 85.8 & 0.8 & 85.8 & \erf{평균 이내}{2 SD 이내} \\
    $\mathrm{CVaR}_{99}$ 격차 & $\erf{+24.0\%}{+24.7\%}$ & $\erf{1.4\%}{0.6\%}$ & $+24.0\%$ & \erf{평균 이내}{2 SD 이내} \\
    \bottomrule
  \end{tabular}
\end{table}

모든 핵심 통계량에서 기준 시드($s = 42$)의 결과가 $10{,}000$ 시드 모집단의
평균에서 \erf{1}{2}\,SD 이내에 위치하며, Bonferroni 보정 $z$-검정에서
45개 회귀계수 중 기각이 0개로(부록 참조), $s = 42$가
모집단의 대표값으로서 타당함이 확인되었다.


%% -------------------------------------------------------------------
\subsection{\vrev{소프트웨어 환경}}
\label{subsec:software}
%% -------------------------------------------------------------------

시뮬레이션, 통계 분석, 시각화에 사용된 소프트웨어 환경을
\p{[}표~\ref{tab:software-env}\p{]}에 정리하였다. 모든 시뮬레이션은 동일한
환경에서 수행되었으며, 난수 생성기는 NumPy의 기본 PCG64
(\texttt{numpy.random.default\_rng()})를 사용하여 재현성을 보장하였고,
비중첩 대역 시드 체계가 생성기의 내부 상태 독립성을 보장한다.
전체 코드와 데이터는 GitHub 공개 저장소에
공개되어 완전한 재현이 가능하며, 3.2.4(재현성 자료 공개)의 상세 공개내용을 참조한다.

\begin{table}[htbp]
\centering
\caption{시뮬레이션 및 분석 소프트웨어 환경}
\label{tab:software-env}
\begin{tabular}{lll}
\toprule
\textbf{구분} & \textbf{소프트웨어} & \textbf{버전} \\
\midrule
프로그래밍 언어 & Python          & 3.12.3 \\
수치 계산       & NumPy           & 2.4.0 \\
통계 분석       & SciPy           & 1.16.3 \\
회귀 분석       & statsmodels     & 0.14.6 \\
기계학습        & scikit-learn    & 1.8.0 \\
데이터 처리     & pandas          & 2.3.3 \\
시각화          & Matplotlib      & 3.10.8 \\
문서 조판       & \TeX\ Live (Xe\LaTeX) & 2023 \\
\bottomrule
\end{tabular}
\end{table}


%% -------------------------------------------------------------------
\subsection{\vrev{Golden Compare를 통한 L0 기반 검증}}
\label{subsec:golden_compare}
%% -------------------------------------------------------------------

6단계 다층 검증 프레임워크의 \textbf{L0(기반 검증)} 단계로서,
시뮬레이션 코드의 구현 정확성과 SSOT(Single Source of Truth) 일관성을
결정론적으로 검증하였다.
\textbf{Golden Compare}는 확률적 시뮬레이션이 아닌,
대표 매개변수를 $\DRTO$ 수식에 직접 대입하여
해석적 기댓값과 코드 출력을 비교하는 결정론적 검증 메커니즘이다.
카테고리 1(순서 제약)과 카테고리 2(물리적 경계)는 각각
\jrev{H3}(순서 보존)와 H2(경계 보장)을 결정론적으로 사전 검증하는 역할을 겸한다.

7개 산업 시나리오(①기준선, ②제조업, ③금융, ④의료, ⑤IT서비스, ⑥물류, ⑦소매)에 대해
4개 카테고리, 총 47개 검증 포인트를 \p{\tabref{tab:golden_compare_summary}에} 설정하였다.

\begin{table}[htbp]
  \centering
  \caption{Golden Compare 47개 검증 포인트 결과 요약}
  \label{tab:golden_compare_summary}
  \small
  \begin{tabular}{clccl}
    \toprule
    \textbf{카테고리} & \textbf{검증 내용} & \textbf{적용 시나리오} & \textbf{포인트 수} & \textbf{결과} \\
    \midrule
    1. 순서 제약   & $\DRTO_{C1} \leq \DRTO_{C2} \leq \DRTO_{C3}$ (P99 기준) & ①--⑦ (전체) & 7  & 100\% PASS \\
    2. 물리적 경계 & $0 \leq \DRTO \leq R_{\max}$                              & ①--⑦ (전체) & 28 & 100\% PASS \\
    3. $z_O$ 동일성 & $z_O^{(C1)} = z_O^{(C2)} = z_O^{(C3)}$ (채널 독립성)    & ②--⑦ & 6  & 100\% PASS \\
    4. P99 비율 일관성 & $\DRTO_{C3,\text{P99}} / \DRTO_{C2,\text{P99}} \approx \erf{1.57}{1.66}$ & ②--⑦ & 6  & 100\% PASS \\
    \midrule
    \textbf{합계} & & & \textbf{47} & \textbf{100\% PASS} \\
    \bottomrule
  \end{tabular}
\end{table}

\chg{[표 4-18]에서 카테고리 3과 4는 시나리오 1(기준선)을 별도로 상세 검증하고 시나리오 2--7의 6개를 대상으로
하며, 각 카테고리의 점별 검증 결과와 7개 시나리오 계산 상세는 공개 재현 패키지에 수록한다.}

첫째, \textbf{카테고리 1(순서 제약)}은 P99 기준에서
$\DRTO_{C1} \leq \DRTO_{C2} \leq \DRTO_{C3}$가 7개 시나리오 모두에서 성립함을 확인하였다.
둘째, \textbf{카테고리 2(물리적 경계)}는 7개 시나리오 $\times$ 4개 조건(C0--C3) = 28개 $\DRTO$ 값
모두에서 $0 \leq \DRTO \leq R_{\max}$를 만족하였다.
셋째, \textbf{카테고리 3($z_O$ 동일성)}은 동일 시나리오 내에서 C1, C2, C3의
관측성 입력 $z_O = \kappa \cdot k_O \cdot O_{\text{raw}}$가 조건에 무관하게
동일함을 확인하여, 관측성 채널과 불확실성 채널의 독립성을 검증하였다.
넷째, \textbf{카테고리 4(P99 비율 일관성)}는 $\DRTO_{C3,\text{P99}} / \DRTO_{C2,\text{P99}}$
비율이 \erf{시나리오에 무관하게 $\approx 1.57$로 안정됨}{시나리오 전반에서 $1.58$--$1.75$(평균 $\approx 1.66$)로 일관되게 1을 상회함}을 확인하였다.

47개 검증 포인트 전체에서 PASS를 달성하여,
$\DRTO$ 개별 시그모이드 바운딩 모델의 수식 구현이
SSOT와 정합함이 결정론적으로 확인되었다.
L0 검증이 완료됨에 따라, 이후 L1--L5 가설 검증이
올바른 구현에 기반함이 보장된다.
7개 시나리오의 $\DRTO$ 계산 상세와 4개 카테고리 47개 검증 포인트의
점별 결과는 4.6.5(7개 시나리오 DRTO 계산 상세)에 제시한다.



%% =================================================================
%% [분량보존 복원] 부록 E Golden Compare 7시나리오·47점 점별 결과
%% =================================================================
\subsection{7개 시나리오 $\DRTO$ 계산 상세}
\label{app:golden-scenarios}

\subsubsection{계산 프레임워크}
\label{app:golden-framework}

각 시나리오에 대해 고정된 입력 매개변수로 이론적 $\DRTO$를 산출한다.
$\DRTO$ 개별 시그모이드 바운딩 모델은 \m{식~(\ref{eq:app-golden-model})에 나타낸 바와 같다.}

\begin{equation}
  \DRTO = R_0
  + \underbrace{-R_0 \cdot S\!\bigl(\kappa \cdot k_O \cdot O_{\text{raw}}\bigr)}_{E_{O,bnd}}
  + \underbrace{(R_{\max} - R_0) \cdot S\!\bigl(\kappa \cdot k_U \cdot \tfrac{R_0}{R_{\max}-R_0} \cdot U_{\text{raw}}\bigr)}_{E_{U,bnd}}
  \label{eq:app-golden-model}
\end{equation}

이고, 여기서 $S(z) = 2/(1+e^{-z})-1$, $\kappa = 1.2$, $k_U = 1 - k_O$이다.

\subsubsection{시나리오 정의}
\label{app:golden-scenario-def}

7개 산업 시나리오의 매개변수를 \p{[}표~\ref{tab:app-golden-scenarios}\p{]}에 정의한다.

\begin{table}[htbp]
\centering
\caption{Golden Compare 7개 검증 시나리오 매개변수}
\label{tab:app-golden-scenarios}
\small
\begin{tabular}{clccccl}
\toprule
\textbf{\#} & \textbf{시나리오} & \textbf{$R_0$ (h)} & \textbf{$R_{\max}$ (h)}
  & \textbf{$O_{\text{raw}}$} & \textbf{$k_O$} & \textbf{특성} \\
\midrule
1 & \m{기준선(Baseline)}  & 6.0  & 24.0 & 0.50 & 0.50 & 기본 파라미터 \\
2 & 제조업             & 4.0  & 16.0 & 0.60 & 0.55 & 중간 $R_0$, 균형적 관측성 \\
3 & 금융               & 1.0  & 4.0  & 0.77 & 0.45 & 짧은 $R_0$, 높은 관측성 \\
4 & 의료               & 2.0  & 8.0  & 0.68 & 0.50 & 짧은 $R_0$, 높은 관측 필요 \\
5 & IT서비스            & 0.5  & 2.0  & 0.82 & 0.55 & 가장 짧은 $R_0$, 최고 관측성 \\
6 & 물류               & 6.0  & 24.0 & 0.53 & 0.48 & 긴 $R_0$, 중간 관측성 \\
7 & 소매               & 3.0  & 12.0 & 0.57 & 0.52 & 중간 규모, 균형적 \\
\bottomrule
\end{tabular}
\end{table}

\chg{모든 시나리오에서 $R_{\max} = 4R_0$(MTPD 기반 상한)를 적용하며, 시나리오별 $k_O$와
$O_{\text{raw}}$는 산업 특성을 반영한 대표값이다.}

\subsubsection{시나리오 1의 기준선(Baseline) 계산 상세}
\label{app:golden-baseline}

\chg{시나리오 1(기준선)의 매개변수는 $R_0 = 6.0$h, $R_{\max} = 24.0$h, $O_{\text{raw}} = 0.50$,
$k_O = 0.50$이며, 조건별 계산은 다음과 같다.}

\chg{먼저 C0(정적 기준선)에서는 $\DRTO$가 다음과 같이 기준선과 일치한다.}
\begin{equation*}
  \DRTO_{C0} = R_0 = 6.000\;\text{h}, \quad \eta_{C0} = 1.000
\end{equation*}

\chg{C1(관측성만)의 계산은 다음과 같다.}
\begin{align*}
  z_O &= \kappa \cdot k_O \cdot O_{\text{raw}} = 1.2 \times 0.50 \times 0.50 = 0.300 \\
  S(z_O) &= \frac{2}{1+e^{-0.300}} - 1 = 0.1489 \\
  E_{O,bnd} &= -6.0 \times 0.1489 = -0.894\;\text{h} \\
  \DRTO_{C1} &= 6.0 + (-0.894) + 0 = 5.106\;\text{h}, \quad \eta_{C1} = 0.851
\end{align*}

\chg{C2(가우시안 불확실성 추가, $U_{\text{raw}} = 3.0$ 대표값, $k_U = 0.5$)의 계산은 다음과 같다.}
\begin{align*}
  z_U &= \kappa \cdot k_U \cdot \frac{R_0 \cdot U_{\text{raw}}}{R_{\max} - R_0}
       = 1.2 \times 0.5 \times \frac{6.0 \times 3.0}{18.0} = 0.600 \\
  S(z_U) &= \frac{2}{1+e^{-0.600}} - 1 = 0.2913 \\
  E_{U,bnd} &= (24.0 - 6.0) \times 0.2913 = 5.244\;\text{h} \\
  \DRTO_{C2} &= 6.0 + (-0.894) + 5.244 = 10.350\;\text{h}, \quad \eta_{C2} = 1.725
\end{align*}

\chg{C3(파레토 불확실성 추가, $U_{\text{raw}} = 3.0$ 대표값, $k_U = 0.5$)에서는} 파레토의 대표값도
기댓값 $E[U_{\text{raw}}] = 3.0$을 사용하므로 $z_U$ 산출이 C2와 동일하다.
\begin{align*}
  z_U &= \kappa \cdot k_U \cdot \frac{R_0 \cdot U_{\text{raw}}}{R_{\max} - R_0}
       = 1.2 \times 0.5 \times \frac{6.0 \times 3.0}{18.0} = 0.600 \\
  S(z_U) &= \frac{2}{1+e^{-0.600}} - 1 = 0.2913 \\
  E_{U,bnd} &= (24.0 - 6.0) \times 0.2913 = 5.244\;\text{h} \\
  \DRTO_{C3} &= 6.0 + (-0.894) + 5.244 = 10.350\;\text{h}, \quad \eta_{C3} = 1.725
\end{align*}

대표값(평균)에서는 C2와 C3의 $\DRTO$가 일치한다.
두 조건의 차이는 평균이 아닌 \textbf{꼬리 영역}에서 발현된다.
파레토의 P99($U_{\text{raw}} = 21.85$)를 대입하면
$z_U = 1.2 \times 0.5 \times 21.85/3 = 4.370$,
$S(4.370) = \erf{0.9875}{0.975}$,
$E_{U,bnd} = 18.0 \times \erf{0.9875}{0.975} = \erf{17.775}{17.55}$h,
$\DRTO = 6.0 - 0.894 + \erf{17.775}{17.55} = \erf{22.881}{22.657}$h ($\eta = \erf{3.814}{3.776}$)로
C2의 P99($\DRTO = 18.35$h)보다 현저히 높다.

\chg{검증 결과} $0 < \text{C1}(5.106) < \text{C2}(10.350) < \text{C3}_{\text{P99}}(\erf{22.881}{22.657}) < 24.0$\chg{으로 물리적 경계를 충족한다.}

\subsubsection{산업별 C0, C1, C2, C3 계산 요약}
\label{app:golden-industry-c1}

7개 시나리오에 대해 C0(정적 기준선), C1(관측성), C2/C3(불확실성 추가) 조건의
$\DRTO$를 \p{[}표~\ref{tab:app-golden-c1-summary}\p{]}에 정리하였다.
C2와 C3는 등평균 설계($E[U_{\text{raw}}] = 3.0$)이므로
대표값 계산에서는 동일한 결과를 보인다.

\begin{table}[htbp]
\centering
\caption{산업별 $\DRTO$ 계산 \m{결과(대표값 기준)}}
\label{tab:app-golden-c1-summary}
\small
\begin{tabular}{clrrrrrrrr}
\toprule
\textbf{\#} & \textbf{시나리오}
  & \textbf{$\DRTO_{C0}$} & \textbf{$\DRTO_{C1}$} & \textbf{$\eta_{C1}$}
  & \textbf{감소율}
  & \textbf{$\DRTO_{\text{C2=C3}}$} & \textbf{$\eta_{\text{C2=C3}}$}
  & \textbf{$\DRTO_{C3,\text{P99}}$} & \textbf{$\eta_{C3,\text{P99}}$} \\
  & & (h) & (h) & & & (h) & & (h) & \\
\midrule
1 & 기준선    & 6.000 & 5.107 & 0.851 & 14.9\% & 10.350 & 1.725 & 22.657 & 3.776 \\
2 & 제조업    & 4.000 & 3.218 & 0.805 & 19.5\% &  6.382 & 1.595 & 14.757 & 3.689 \\
3 & 금융      & 1.000 & 0.795 & 0.795 & 20.5\% &  1.751 & 1.751 &  3.746 & 3.746 \\
4 & 의료      & 2.000 & 1.598 & 0.799 & 20.1\% &  3.345 & 1.673 &  7.448 & 3.724 \\
5 & IT서비스   & 0.500 & 0.368 & 0.736 & 26.4\% &  0.763 & 1.527 &  1.810 & 3.620 \\
6 & 물류      & 6.000 & 5.091 & 0.849 & 15.1\% & 10.532 & 1.755 & 22.713 & 3.786 \\
7 & 소매      & 3.000 & 2.472 & 0.824 & 17.6\% &  4.995 & 1.665 & 11.205 & 3.735 \\
\bottomrule
\end{tabular}
\end{table}

\chg{표의 C1은 $E_{O,bnd} = -R_0 \cdot S(\kappa \cdot k_O \cdot O_{\text{raw}})$($\kappa = 1.2$)로 계산하며,
C2와 C3는 $U_{\text{raw}} = 3.0$(대표값, 두 분포의 공통 평균)에서 계산하므로 등평균 설계상 대표값에서는
두 조건이 일치한다. $\DRTO_{C3,\text{P99}}$는 파레토 P99($U_{\text{raw}} = 21.85$) 대입값으로, 꼬리 영역에서
C3가 C2보다 현저히 높은 $\DRTO$를 유발함을 보여준다.}


\subsection{47개 검증 포인트 전체 결과}
\label{app:golden-47points}

Golden Compare의 검증 포인트는 4개 카테고리, 총 47개로 구성된다.
\chg{47개의 산출 근거는 다음과 같다. 카테고리 1(순서 제약)은 7개 시나리오에 P99 기준
$C1 \leq C2 \leq C3$ 1건씩을 적용하여 7개, 카테고리 2(물리적 경계)는 7개 시나리오에 4개
조건(C0, C1, C2, C3)을 적용하여 28개, 카테고리 3($z_O$ 동일성)은 6개 시나리오에 C1=C2=C3 확인
1건씩을 적용하여 6개, 카테고리 4(P99 비율 일관성)는 6개 시나리오에 C3/C2 비율 1건씩을 적용하여
6개로 구성된다. 카테고리 3과 4의 시나리오 수가 6개인 것은 시나리오 1(기준선)이 상세 계산에서 별도로
검증되어 시나리오 2--7만을 대상으로 하기 때문이다. 따라서 합계는 47개이며, 모든
포인트에서 PASS를 달성하였다.}

\subsubsection{카테고리 1의 순서 제약(7 포인트)}
\label{app:golden-order}

\p{[}표~\ref{tab:app-golden-order}\p{]}은 P99 기준에서 $\DRTO_{C1} \leq \DRTO_{C2} \leq \DRTO_{C3}$ 순서가 7개 시나리오 모두에서 성립하는지 검증한 결과를 보여준다.

\begin{table}[htbp]
\centering
\caption{순서 제약 검증 \m{결과(P99 기준)}}
\label{tab:app-golden-order}
\small
\begin{tabular}{clcccc}
\toprule
\textbf{\#} & \textbf{시나리오}
  & \textbf{$\DRTO_{C1,\text{P99}}$ (h)} & \textbf{$\DRTO_{C2,\text{P99}}$ (h)}
  & \textbf{$\DRTO_{C3,\text{P99}}$ (h)} & \textbf{판정} \\
\midrule
1 & 기준선    & 5.998 & 8.443  & 13.224 & PASS \\
2 & 제조업    & 3.997 & 5.670  & 8.872  & PASS \\
3 & 금융      & 0.999 & 1.398  & 2.218  & PASS \\
4 & 의료      & 1.998 & 2.830  & 4.436  & PASS \\
5 & IT서비스   & 0.500 & 0.704  & 1.109  & PASS \\
6 & 물류      & 5.998 & 8.443  & 13.224 & PASS \\
7 & 소매      & 2.998 & 4.221  & 6.612  & PASS \\
\midrule
\multicolumn{5}{l}{\textbf{합계: 7/7 PASS}} & \textbf{100\%} \\
\bottomrule
\end{tabular}
\end{table}

\subsubsection{카테고리 2의 물리적 경계(28 포인트)}
\label{app:golden-boundary}

\p{[}표~\ref{tab:app-golden-boundary}\p{]}는 7개 시나리오 $\times$ 4개 조건 $= 28$ 개 $\DRTO$ 값 모두에서 $0 \leq \DRTO \leq R_{\max}$가 성립하는지 검증한 결과를 보여준다.

\begin{table}[htbp]
\centering
\caption{물리적 경계 검증 \m{결과(전체 28 포인트)}}
\label{tab:app-golden-boundary}
\small
\begin{tabular}{clrrrrl}
\toprule
\textbf{\#} & \textbf{시나리오} & \textbf{$R_{\max}$ (h)}
  & \textbf{C0} & \textbf{C1} & \textbf{C2/C3 최대} & \textbf{경계 이내} \\
\midrule
1 & 기준선    & 24.0 & 6.000 & 5.106 & 13.224 & PASS (4건) \\
2 & 제조업    & 16.0 & 4.000 & 3.221 & 8.872  & PASS (4건) \\
3 & 금융      & 4.0  & 1.000 & 0.796 & 2.218  & PASS (4건) \\
4 & 의료      & 8.0  & 2.000 & 1.599 & 4.436  & PASS (4건) \\
5 & IT서비스   & 2.0  & 0.500 & 0.369 & 1.109  & PASS (4건) \\
6 & 물류      & 24.0 & 6.000 & 5.095 & 13.224 & PASS (4건) \\
7 & 소매      & 12.0 & 3.000 & 2.474 & 6.612  & PASS (4건) \\
\midrule
\multicolumn{5}{l}{\textbf{합계: 28/28 PASS}} & & \textbf{100\%} \\
\bottomrule
\end{tabular}
\end{table}

\subsubsection{카테고리 3의 $z_O$ 동일성(6 포인트)}
\label{app:golden-zo}

동일 시나리오 내에서 C1, C2, C3의 관측성 입력 $z_O = \kappa \cdot k_O \cdot O_{\text{raw}}$가
조건에 무관하게 동일한 값을 갖는지 검증한다.
$\DRTO$ 모델에서 관측성 채널($E_{O,bnd}$)과 불확실성 채널($E_{U,bnd}$)은
독립적으로 작동하도록 설계되었으므로,
C2/C3에서 불확실성 분포를 다르게 적용하더라도
관측성 쪽 계산($z_O$)은 C1과 완전히 동일해야 한다.
이 조건이 성립해야 ``불확실성 조건을 추가해도 관측성 채널이
오염되지 않는다''는 \textbf{채널 독립성}이 보장된다. \p{[}표~\ref{tab:app-golden-zo}\p{]}는 $z_O$ 동일성 검증 결과를 보여준다.

\begin{table}[htbp]
\centering
\caption{$z_O$ 동일성 검증 결과}
\label{tab:app-golden-zo}
\small
\begin{tabular}{clcccc}
\toprule
\textbf{\#} & \textbf{시나리오}
  & \textbf{$z_O^{(C1)}$} & \textbf{$z_O^{(C2)}$} & \textbf{$z_O^{(C3)}$} & \textbf{판정} \\
\midrule
2 & 제조업    & 0.396 & 0.396 & 0.396 & PASS \\
3 & 금융      & 0.416 & 0.416 & 0.416 & PASS \\
4 & 의료      & 0.408 & 0.408 & 0.408 & PASS \\
5 & IT서비스   & 0.541 & 0.541 & 0.541 & PASS \\
6 & 물류      & 0.306 & 0.306 & 0.306 & PASS \\
7 & 소매      & 0.356 & 0.356 & 0.356 & PASS \\
\midrule
\multicolumn{5}{l}{\textbf{합계: 6/6 PASS}} & \textbf{100\%} \\
\bottomrule
\end{tabular}
\end{table}

\chg{시나리오 1(기준선)은 대표 시나리오로서 별도로 검증하며, $z_O = \kappa \cdot k_O \cdot O_{\text{raw}}$는
조건(C1, C2, C3)에 무관하게 동일하다.}

\subsubsection{카테고리 4의 P99 비율 일관성(6 포인트)}
\label{app:golden-ratio}

$\DRTO_{C3,\text{P99}} / \DRTO_{C2,\text{P99}}$ 비율이
시나리오에 무관하게 안정적인 값($\approx 1.57$)을 유지하는지 검증한다.
원시 입력 $U_{\text{raw}}$의 P99 비율은 파레토/가우시안 $\approx 3.95$이나,
시그모이드 바운딩을 거치면서 비선형 압축이 적용되어
$\DRTO$ 수준에서의 비율은 $\approx 1.57$로 수렴한다(\p{[}표~\ref{tab:app-golden-ratio}\p{]}).

\begin{table}[htbp]
\centering
\caption{P99 \m{비율($\DRTO_{C3,\text{P99}} / \DRTO_{C2,\text{P99}}$)} 일관성 검증}
\label{tab:app-golden-ratio}
\small
\begin{tabular}{clcccl}
\toprule
\textbf{\#} & \textbf{시나리오}
  & \textbf{$\DRTO_{C2,\text{P99}}$} & \textbf{$\DRTO_{C3,\text{P99}}$}
  & \textbf{비율} & \textbf{판정} \\
\midrule
2 & 제조업    & 5.670  & 8.872  & 1.56 & PASS \\
3 & 금융      & 1.398  & 2.218  & 1.59 & PASS \\
4 & 의료      & 2.830  & 4.436  & 1.57 & PASS \\
5 & IT서비스   & 0.704  & 1.109  & 1.58 & PASS \\
6 & 물류      & 8.443  & 13.224 & 1.57 & PASS \\
7 & 소매      & 4.221  & 6.612  & 1.57 & PASS \\
\midrule
\multicolumn{4}{l}{\textbf{합계: 6/6 PASS}} & \textbf{평균 1.57} & \textbf{100\%} \\
\bottomrule
\end{tabular}
\end{table}

\chg{P99 비율의 평균은 $\approx 1.57$로 시나리오 간 높은 일관성을 보이며, 이 비율은 파레토 분포의 P99와
가우시안 P99의 구조적 차이를 반영하여 $R_0$와 $R_{\max}$의 \m{절댓값}에 무관하게 안정적이다.}

\subsubsection{검증 포인트 종합}
\label{app:golden-summary}

4개 카테고리 47개 검증 포인트의 종합 결과를 \p{[}표~\ref{tab:app-golden-total}\p{]}에 요약한다.

\begin{table}[htbp]
\centering
\caption{Golden Compare 47개 검증 포인트 종합 결과}
\label{tab:app-golden-total}
\begin{tabular}{clcl}
\toprule
\textbf{카테고리} & \textbf{검증 내용} & \textbf{포인트 수} & \textbf{결과} \\
\midrule
순서 제약       & $\DRTO_{C1} \leq \DRTO_{C2} \leq \DRTO_{C3}$ (P99 기준)      & 7  & 100\% PASS \\
물리적 경계     & $0 \leq \DRTO \leq R_{\max}$                                  & 28 & 100\% PASS \\
$z_O$ 동일성   & $z_O^{(C1)} = z_O^{(C2)} = z_O^{(C3)}$ (조건 독립)            & 6  & 100\% PASS \\
P99 비율 일관성 & $\DRTO_{C3,\text{P99}} / \DRTO_{C2,\text{P99}} \approx 1.57$           & 6  & 100\% PASS \\
\midrule
\textbf{합계}  &                                                               & \textbf{47} & \textbf{100\% PASS} \\
\bottomrule
\end{tabular}
\end{table}

\chg{47개 검증 포인트 모두에서 PASS를 달성하여, $\DRTO$ 개별 시그모이드 바운딩 모델의 구현 정확성과
SSOT(Single Source of Truth) 일관성이 검증되었다.}

\chg{Golden Compare 검증은 네 가지 결론을 뒷받침한다. 첫째, 구현 정확성으로서 7개 시나리오와 4개
조건을 곱한 28개 $\DRTO$ 값 모두에서 이론적 계산치와 시뮬레이션 출력이 일치한다. 둘째, 물리적
타당성으로서 모든 $\DRTO$가 $(0,\; R_{\max})$ 범위 내에 있어 명제 2.1(물리적 경계
보장)을 실증적으로 확인한다. 셋째, 구조적 일관성으로서 $z_O$ 동일성과 P99 비율 일관성을 통해 관측성과
불확실성 채널의 독립 작동이 검증된다. 넷째, 산업 범용성으로서 기준선과 6개 산업에서 동일한 구조적
패턴이 재현되어 $\DRTO$ 프레임워크의 산업 범용 적용성을 확인한다.}



\chg{이상을 종합한 H2의 검증 결론을 정리한다.} \mrev{H2는 ``모든 표본에서 $0 \leq \DRTO \leq R_{\max} = 24$h''를 판정
기준으로 한다. L0 Golden Compare의 28개 물리적 경계 검증 포인트 (시나리오
①$\sim$⑦ $\times$ 4개 변수)와 본 절 확률 시뮬레이션 $n = 10{,}000$ 표본
전체에서 경계 위반 사례가 0건이다. 시그모이드 바운딩 함수의 수학적
구조가 이를 결정론적으로 보장하며, \textbf{H2는 채택된다}.}

\chg{시드 독립성에 관한 H6의 검증 결론은 다음과 같다.} \mrev{H6은 ``서로 다른 시드 간 변수 상관 $|r(X_i, X_j)| < 0.02$''를
기준으로 한다. 모든 페어에서 $|r|_{\max} = 0.016 < 0.02$이므로,
$\alpha = 0.05$ 유의수준에서 영(0)과 통계적으로 구별되지 않으며,
\textbf{H6은 채택된다}.}

\chg{이어 다중 시드 강건성에 관한 H7의 검증 결론은 다음과 같다.} \mrev{H7은 ``$10^8$회 다중 시드 시뮬레이션에서 회귀 계수 안정''을 기준으로
한다. 시드 $s = 42, 43, 44$의 45개 계수 중 신뢰구간 기각이 0건이며,
계수 평균 $\pm 0.5\%$ 이내 안정. 따라서 \textbf{H7은 채택된다}.}
         % 4.6 수렴 진단

% 교차 분석 (3차 심사: 파일명을 실제 절 번호와 일치)


%% ===================================================================
\section{\vrev{가우시안 vs.\ 파레토 불확실성 종합 비교 분석}}
\label{sec:ch5_c2_c3_synthesis}
%% ===================================================================

\jrev{4.2절(C2)·4.3절(C3)}에서 C2(가우시안)와 C3(파레토) 조건의 개별 실험 결과를 각각 검토하였으며,
두 조건의 핵심 지표를 \jrev{\tabref{tab:c2_c3_key_metrics}}에서 직접 대비하였다.
본 절에서는 이를 바탕으로 두 불확실성 유형의 구조적 차이를 종합적으로 분석하고,
이후 절(효과 크기, 회귀 설명력, 이중채널 감쇠)의 교차분석으로 연결하는
통합적 시각을 제공한다.


%% -------------------------------------------------------------------
\subsection{\vrev{분포 특성과 $\eta$ 분포의 구조적 차이}}
\label{subsec:c2c3_distribution}
%% -------------------------------------------------------------------

가우시안과 파레토 불확실성은 동일한 기댓값 $E[U_{\text{raw}}] = 3.0$을
공유하지만, 분포의 형태적 특성에서 근본적으로 상이하다.
\jrev{\tabref{tab:c2_c3_distribution_properties}\refpo{tab:c2_c3_distribution_properties}{은}{는}} 두 분포의 이론적 특성과
이로부터 유도되는 $\eta$ 분포의 구조적 차이를 정리한 것이다.

\begin{table}[htbp]
  \centering
  \caption{\jrev{C2 vs.\ C3 불확실성 분포의 구조적 특성 비교}}
  \label{tab:c2_c3_distribution_properties}
  \small
  \begin{tabular}{l cc l}
    \toprule
    \textbf{특성} & \textbf{C2 (가우시안)} & \textbf{C3 (파레토)} & \textbf{함의} \\
    \midrule
    분포 & $N(3, 1)$ & \jrev{$6 \cdot \mathrm{Pareto}(3)$} & 동일 평균, 상이 형태 \\
    기댓값 $E[U]$ & 3.0 & 3.0 & 동일 \\
    분산 $\mathrm{Var}[U]$ & 1.0 & \jrev{27.0} & \jrev{C3가 가우시안의 27배} \\
    왜도(Skewness) & 0 & 양 & C3 우측 비대칭 \\
    첨도(Kurtosis) & 3 (정규) & $\infty$ & C3 극단값 빈도 높음 \\
    꼬리 감쇠 & 지수적 $e^{-x^2/2}$ & 멱법칙 $x^{-\alpha}$ & C3의 꼬리 질적 차이 \\
    \midrule
    $\bar{\eta}$ & 1.682 & 1.514 & C2가 평균적으로 높음 \\
    $\mathrm{SD}(\eta)$ & 0.615 & 0.791 & C3가 분산 $1.29$배 \\
    $\eta_{\max}$ & 3.443 & 4.000 & C3가 $R_{\max}$ 포화 접근 \\
    \bottomrule
  \end{tabular}
\end{table}

이 비교에서 핵심적인 관찰은 다음과 같다.
\textbf{첫째}, 가우시안의 유한 분산($\mathrm{Var} = 1.0$)은
$\eta$의 예측 가능한 변동을 생성하는 반면,
\jrev{파레토의 큰 분산($\mathrm{Var} = 27$, 가우시안의 27배)}은 표본 수준에서 높은 $\mathrm{SD}(\eta) = 0.791$과
$\eta_{\max} = 4.000$ (= $R_{\max}/R_0$)에 근접하는 극단값으로 발현한다.
\textbf{둘째}, C2의 평균이 C3보다 높은($1.682 > 1.514$) 이유는
가우시안의 대칭적 종형 분포가 평균 근방에 밀도를 집중시키는 반면,
파레토의 우측 비대칭이 평균을 좌측으로 편향시키기 때문이다.
이 ``평균에서의 역전''은 분포의 첨도와 왜도 차이에서 비롯되며,
극단 영역에서의 $\eta$ 행동은 정반대의 양상을 보인다.


%% -------------------------------------------------------------------
\subsection{\vrev{백분위별 $\eta$ 비교(분포 교차 구조)}}
\label{subsec:c2c3_percentile}
%% -------------------------------------------------------------------

\jrev{\tabref{tab:c2_c3_percentile}\refpo{tab:c2_c3_percentile}{은}{는}} C2와 C3의 $\eta$ 백분위 값을
주요 분위수에서 비교한 것이다. 이 표는 두 분포의 교차 구조를
체계적으로 보여준다.

\begin{table}[htbp]
  \centering
  \caption{\jrev{C2 vs.\ C3 효율성 비율 $\eta$의 백분위별 비교 ($n = 10{,}000$)}}
  \label{tab:c2_c3_percentile}
  \small
  \begin{tabular}{l rr r l}
    \toprule
    \textbf{백분위} & \textbf{C2 $\eta$} & \textbf{C3 $\eta$} &
    \textbf{$\Delta$(C3$-$C2)} & \textbf{우위 조건} \\
    \midrule
    P10 & 0.897 & 0.786 & $-0.112$ & C2 $>$ C3 \\
    P25 & 1.183 & 0.964 & $-0.219$ & C2 $>$ C3 \\
    P50 (중앙값) & 1.645 & 1.230 & $-0.415$ & C2 $>$ C3 \\
    P75 & 2.141 & 1.842 & $-0.299$ & C2 $>$ C3 \\
    P85 & 2.395 & 2.371 & $-0.024$ & C2 $\approx$ C3 \\
    \textbf{P85.8} & \textbf{2.42} & \textbf{2.42} & $\approx$\textbf{0.000} & \textbf{교차점} \\
    P90 & 2.541 & 2.772 & $+0.231$ & C3 $>$ C2 \\
    P95 & 2.731 & 3.319 & $+0.588$ & C3 $>$ C2 \\
    P99 & 3.059 & 3.880 & $+0.821$ & C3 $>$ C2 \\
    \bottomrule
  \end{tabular}
\end{table}

이 백분위 분석에서 세 가지 구조적 특성이 관찰된다.

\chg{첫째, P85.8 이하의 가우시안 우위 영역에서는} 전체 표본의 85.8\%에 해당하는 일상적 운영 범위에서
C2의 $\eta$가 C3보다 일관되게 높다. 이 영역에서 가우시안의 대칭적 분포가 평균 근방에 밀도를 집중시켜,
불확실성 프리미엄이 파레토보다 체계적으로 큰 값을 산출하며, 중앙값 기준 C2와 C3의 격차는 $0.415$에
달한다. \chg{둘째, 약 P85.8의 교차점에서는} 두 분포의 $\eta$ 백분위가 교차한다. 이 교차점은
``일상적 변동''과 ``극단적 변동''을 분리하는 자연적 경계이며, 이후 분할 회귀 분석(\jrev{\subsecref{subsec:ch5_split_regression}})의
분리 기준으로 사용된다. \chg{셋째, P85.8 이상에서는 파레토 우위가 급격히 확대되어} 교차점 이상에서
C3의 $\eta$가 C2를 초과하며, 격차는 상위 백분위로 갈수록 비선형적으로 확대된다. $\Delta$(C3$-$C2) 값이
P90에서 $+0.231$, P99에서 $+0.821$로 증가하여, 극단 영역에서 파레토 꼬리가 $\DRTO$를 $R_{\max}$에
근접시키는 효과를 보여준다.


%% -------------------------------------------------------------------
\subsection{\vrev{회귀 계수의 직접 비교(분포별 변수 영향력 차이)}}
\label{subsec:c2c3_regression_coefficients}
%% -------------------------------------------------------------------

\jrev{\tabref{tab:c2_c3_regression_comparison}\refpo{tab:c2_c3_regression_comparison}{은}{는}} C2와 C3의 1차 다중 회귀 계수를
직접 비교한 것이다.

\begin{table}[htbp]
  \centering
  \caption{\jrev{C2 vs.\ C3 1차 다중 회귀 계수 직접 비교}}
  \label{tab:c2_c3_regression_comparison}
  \small
  \begin{tabular}{l rr r l}
    \toprule
    \textbf{변수} & \textbf{C2 계수} & \textbf{C3 계수} &
    \textbf{비율 (C3/C2)} & \textbf{해석} \\
    \midrule
    절편 $\hat{\beta}_0$ & $+2.024$ & $+2.018$ & 1.00 & 기준수준 동일 \\
    $k_O$ & $-1.869$ & $-1.408$ & 0.75 & C3에서 관측성 효과 감쇠 \\
    $O$ & $-0.281$ & $-0.249$ & 0.89 & C3에서 관측성 기여 약간 감쇠 \\
    $U$ & $+0.248$ & $+0.112$ & 0.45 & C3에서 선형 기여 감소 \\
    \midrule
    $R^2$ & 0.948 & 0.695 & 0.73 & C3의 비선형성으로 설명력 감소 \\
    \bottomrule
  \end{tabular}
\end{table}

회귀 계수 비교에서 도출되는 핵심 패턴은 다음과 같다.
\chg{첫째, 관측성 변수($k_O$, $O$)의 계수가 감쇠한다. C3에서 $k_O$의 계수가 C2 대비 $0.75$배, $O$의
계수가 $0.89$배로 감쇠되는데, 이는 파레토 불확실성이 존재하는 환경에서 관측성의 복구 단축 효과가
체계적으로 약화됨을 의미한다. 나아가 C3 내부에서 P85.8 분할 회귀를 수행하면 Core($k_O = -0.798$) 대비
Tail($k_O = -0.415$)에서 $0.52$배의 추가 감쇠가 관측되어(\jrev{\subsecref{subsec:ch5_split_regression}}),
극단 영역으로 갈수록 관측성 효과가 단계적으로 약화되는 구조적 현상임을 확인한다. 둘째, 불확실성
변수($U$)의 선형 계수가 감소하되 비선형 항이 지배한다. C3에서 $U$의 1차 회귀 계수가 C2 대비 $0.45$배로
오히려 감소하는데, 이는 파레토 분포의 극단값이 시그모이드 포화 영역에서 비선형적으로 작용하여 선형 회귀
계수가 $U$의 실제 영향력을 과소 추정하기 때문이다. 그러나 전진 선택법에서 C3가 $U$를 1순위로 진입시키는
것은(\jrev{\subsecref{subsec:ch5_dominant_variable}}) $U$의 기여가 $U^2$, $k_O \cdot U$ 등 비선형 항을 통해
작용함을 보여주며, 실제로 C3에서 $U$ 관련 변수($U$, $U^2$, $k_O \cdot U$)가 전진 선택 상위 4개 중 3개를
차지한다. 셋째, 설명력 격차가 존재한다. 1차 다중 회귀 $R^2$가 C2의 $0.948$에서 C3의 $0.695$로 $0.253$만큼
하락하는데, 이 격차는 파레토 분포의 비선형 꼬리 구조가 선형 모형으로 포착되지 않는 변동을 생성하기
때문이며, 2차교호작용 모형에서 C2($R^2 = 0.998$)와 C3($R^2 = 0.858$)의 격차가 $0.140$으로 축소되지만
여전히 잔존한다.}


%% -------------------------------------------------------------------
\subsection{\vrev{이중축(Dual-Axis) 불확실성 관리 체계의 종합 해석}}
\label{subsec:c2c3_dual_axis}
%% -------------------------------------------------------------------

C2와 C3의 종합 비교에서 도출되는 핵심 결론은
\textbf{두 불확실성 유형이 상호 배타적이 아닌 상호 보완적}이라는 것이다.

\chg{이는 세 측면에서 구체화된다. 첫째, 일상적 운영 축(C2 기반)으로서, 전체 표본의 85.8\%에 해당하는
일상적 변동 범위에서는 가우시안 분포(C2)가 높은 예측력($R^2 = 0.948$)과 관측성 중심의 간명한 관리 구조를
제공하며, 이 영역에서 BCM 전략의 핵심은 관측성 투자에 의한 복구 시간 단축이다. 둘째, 극단 위험 축(C3
기반)으로서, 상위 14.2\%의 극단 영역에서는 파레토 분포(C3)가 가우시안보다 현실적인 위험 평가를 제공하는데,
$\mathrm{CVaR}_{99}$ 기준 C3가 C2 대비 $+24.0\%$ 높은 값을 보이며 이 영역에서 BCM 전략의 핵심은 불확실성
자체의 정량화와 꼬리 위험 관리이다. 셋째, P85.8 교차점의 전략적 의미로서, P85.8은 두 축을 분리하는 자연적
경계이며 조직은 자신의 불확실성 프로파일이 P85.8 이하(일상적)인지 이상(극단적)인지를 진단하여 적절한
분포를 선택할 수 있다. 이 이중축 체계는 \jrev{\secref{sec:practical-implications}}에서 2단계 운영
프레임워크와 산업별 가이드라인으로 구체화된다.}

\jrev{\figref{fig:c2_c3_synthesis}\refpo{fig:c2_c3_synthesis}{은}{는}} C2와 C3의 $\eta$ 분포 및
P85.8 교차점을 종합적으로 시각화한 것이다.

\begin{figure}[htbp]
  \centering
  \begin{tikzpicture}
    \begin{axis}[
      width=0.92\textwidth,
      height=8cm,
      xlabel={효율성 비율 $\eta$},
      ylabel={확률밀도},
      xmin=0.3, xmax=4.2,
      ymin=0, ymax=1.35,
      ymajorgrids=true,
      xmajorgrids=true,
      grid style={dashed, black},
      legend style={at={(0.5,-0.18)}, anchor=north, font=\small,
                    legend columns=2, /tikz/every even column/.append style={column sep=8pt}},
      legend cell align=left,
      ]
      % Core/Tail 배경 영역 색상 (옅게 조정 — 영역 1/2/3 표시와 양립)
      \fill[green!10] (axis cs:0.3,0) rectangle (axis cs:2.42,1.35);
      \fill[red!8] (axis cs:2.42,0) rectangle (axis cs:4.2,1.35);

      % 영역 1/2/3 경계 수직 점선 (η=1.33, η=2.96, PDF 이중 교차점)
      \draw[black, dashdotted, thick] (axis cs:1.33, 0) -- (axis cs:1.33, 1.05);
      \draw[black, dashdotted, thick] (axis cs:2.96, 0) -- (axis cs:2.96, 1.05);
      \node[font=\tiny, black, fill=white, inner sep=1pt]
        at (axis cs:1.33, 1.08) {$\eta\!=\!1.33$};
      \node[font=\tiny, black, fill=white, inner sep=1pt]
        at (axis cs:2.96, 1.08) {$\eta\!=\!2.96$};

      % C2 (Gaussian): 스무딩된 히스토그램 데이터 (200 bins, σ=2)
      \addplot[blue!70!black, thick, no markers]
        table[x=eta, y=pdf_C2] {figures/data/pdf_c2_c3.dat};
      \addlegendentry{C2 (Gaussian, solid)}

      % C3 (Pareto): 스무딩된 히스토그램 데이터 (200 bins, σ=2)
      \addplot[red!70!black, thick, dashed, no markers]
        table[x=eta, y=pdf_C3] {figures/data/pdf_c2_c3.dat};
      \addlegendentry{C3 (Pareto, dashed)}

      % 영역 1/2/3 라벨 (영역 중앙 상단, 굵은 글씨)
      \node[font=\small\bfseries, black, align=center,
            fill=yellow!20, inner sep=3pt, rounded corners=2pt,
            draw=black, thin, anchor=west]
        at (axis cs:0.32, 0.95)
        {영역 1\\$(\eta\!<\!1.33)$\\\scriptsize C3$>$C2};
      \node[font=\small\bfseries, black, align=center,
            fill=yellow!20, inner sep=3pt, rounded corners=2pt,
            draw=black, thin]
        at (axis cs:2.145, 0.95)
        {영역 2\\$(1.33\!\leq\!\eta\!\leq\!2.96)$\\\scriptsize C2$>$C3};
      \node[font=\small\bfseries, black, align=center,
            fill=yellow!20, inner sep=3pt, rounded corners=2pt,
            draw=black, thin]
        at (axis cs:3.55, 0.65)
        {영역 3\\$(\eta\!>\!2.96)$\\\scriptsize C3$>$C2 (heavy-tail)};

      % P85.8 수직선 (η ≈ 2.42, CDF 교차점)
      \draw[black, dashed, thick] (axis cs:2.42, 0) -- (axis cs:2.42, 1.25);
      \node[font=\footnotesize, fill=white, inner sep=2pt, rounded corners=1pt]
        at (axis cs:2.42, 1.30) {\textbf{P85.8} (CDF 교차)};

      % R0 baseline
      \draw[black, dotted, thick] (axis cs:1.0, 0) -- (axis cs:1.0, 1.25);
      \node[font=\scriptsize, black, fill=white, inner sep=1pt]
        at (axis cs:1.0, 1.30) {$\eta = 1.0$ ($\DRTO = R_0$)};

      % Core 영역 주석 (좌측 하단, 영역 1/2 라벨과 충돌 방지)
      \node[font=\tiny, green!50!black, align=center,
            fill=green!5, inner sep=2pt, rounded corners=1pt,
            draw=green!40!black, thin]
        at (axis cs:0.50, 0.20) {Core\\($\leq$ P85.8)};

      % Tail 영역 주석 (우측 하단)
      \node[font=\tiny, red!60!black, align=center,
            fill=red!5, inner sep=2pt, rounded corners=1pt,
            draw=red!40!black, thin]
        at (axis cs:3.80, 0.15) {Tail\\($>$ P85.8)};

      % 통계 주석 (우상단, 위치 조정)
      \node[font=\scriptsize, align=left, fill=white, inner sep=3pt,
            rounded corners=2pt, draw=black, thin]
        at (axis cs:0.85, 0.35)
        {$\bar{\eta}_{\text{C2}} = 1.682$, SD $= 0.615$\\
         $\bar{\eta}_{\text{C3}} = 1.514$, SD $= 0.791$};

      % 꼬리 확장 화살표
      \draw[-{Stealth}, red!60, thick]
        (axis cs:3.7, 0.06) -- (axis cs:4.1, 0.01);
      \node[font=\tiny, red!70!black, anchor=south] at (axis cs:3.9, 0.07)
        {$\eta \to 4.0$};
    \end{axis}
  \end{tikzpicture}
  \caption{C2 vs.\ C3 $\eta$ 분포 종합 비교}
  \label{fig:c2_c3_synthesis}
\end{figure}

\noindent
\chg{그림은 시뮬레이션 $n = 10{,}000$의 200구간 히스토그램을 $\sigma = 2$로 스무딩한 것으로, PDF 이중
교차점 $\eta = 1.33$과 $\eta = 2.96$(점쇄선)이 분포를 영역 1($\eta < 1.33$, C3 우세), 영역
2($1.33 \leq \eta \leq 2.96$, C2 우세), 영역 3($\eta > 2.96$, C3 heavy-tail)로 분할한다. 이는 CDF
교차점 P85.8($\eta \approx 2.42$, 파선)을 기준으로 한 Core/Tail 분류(옅은 녹색과 적색 배경)와는 별개의
분류이며, C3의 우측 꼬리가 $\eta \to 4.0$($= R_{\max}/R_0$)까지 확장되는 것이 파레토 불확실성의 구조적
특성이다.}
\jrev{\figref{fig:c2_c3_synthesis}}에서 C2와 C3의 확률밀도 곡선은
$\eta \approx 1.33$과 $\eta \approx 2.96$에서 두 번 교차하며,
이 \textbf{이중 교차(dual crossing)} 구조는
파레토 분포의 ``양극단 집중'' 특성을 반영한다.
\jrev{\tabref{tab:pdf_dual_crossing}\refpo{tab:pdf_dual_crossing}{은}{는}} 두 교차점이 분리하는 세 영역의
확률 비중과 그 해석을 정리한 것이다.

\begin{table}[htbp]
  \centering
  \caption{\jrev{C2 vs.\ C3 확률밀도 이중 교차 구조}}
  \label{tab:pdf_dual_crossing}
  \small
  \begin{tabular}{cl c c >{\raggedright\arraybackslash}p{3.8cm}}
    \toprule
    & \textbf{범위} & \textbf{우세} &
    \textbf{확률 비중} & \textbf{해석} \\
    \midrule
    영역 1 & $\eta < 1.33$ & C3 $>$ C2 &
    C3: 55.6\% vs.\ C2: 33.0\% &
    파레토의 좌측 편향 (낮은 $\eta$에 밀집) \\
    영역 2 & $1.33 \leq \eta \leq 2.96$ & C2 $>$ C3 &
    C2: 64.1\% vs.\ C3: 35.1\% &
    가우시안의 종형 대칭 (중간값에 밀집) \\
    영역 3 & $\eta > 2.96$ & C3 $>$ C2 &
    C3: 8.1\% vs.\ C2: 1.7\% &
    파레토의 heavy-tail (극단값) \\
    \bottomrule
  \end{tabular}
\end{table}

\chg{표에서 확률은 각 영역에서의 확률밀도 적분값(면적 비중)이며, CDF 교차점(P85.8, $\eta \approx 2.42$)은
영역 2 내부에서 발생하여 영역 2에서 C2의 누적 초과분이 영역 1의 C3 초과분을 상쇄하는 지점에 해당한다.}

\noindent
이 이중 교차 구조에서 세 가지 함의가 도출된다.
\textbf{첫째}, 영역 1에서 C3의 확률이 C2보다 $22.6\%$p 높다는 것은
파레토 불확실성 환경에서 $\eta < 1.33$ (즉 $\DRTO < 1.33 R_0$)의
비교적 낮은 복구 시간이 출현할 확률이 가우시안보다 체계적으로 높음을 의미한다.
\textbf{둘째}, 영역 2에서 C2가 $29.0\%$p 우세한 것은
가우시안의 종형 분포가 중간 $\eta$ 범위에 밀도를 집중시켜,
``전형적인'' 복구 시간의 예측 가능성이 높음을 보여준다.
\textbf{셋째}, 영역 3에서 C3의 확률이 C2의 $4.8$배($8.1\%$ vs.\ $1.7\%$)인 것은
극단적으로 긴 복구 시간($\eta > 2.96$, 즉 $\DRTO > 17.8$h)의 출현 빈도가
분포 유형에 따라 질적으로 상이함을 확인한다.
이 영역 3의 비대칭이 $\mathrm{CVaR}_{99}$ 격차(C3가 C2 대비 $+24.0\%$)의
직접적 원인이며, \jrev{\secref{sec:ch5_dual_channel}}의 이중채널 감쇠 메커니즘으로
연결된다.
     % 4.7 C2 vs C3 종합 비교
\section{\vrev{조건 간 효과 크기 분석}}
\label{sec:ch5_effect_size}
%% ===================================================================

4개 조건의 효율성 비율 $\eta = \DRTO / R_0$를 \m{가우시안 경로(C0$\to$C1$\to$C2)와 파레토 경로(C0$\to$C1$\to$C3)} 두 갈래로 구분하여
비교하면, 관측성, 불확실성, 분포 형태라는 세 요인의 효과가
구조적으로 분리된다.
\jrev{\tabref{tab:ch5_effect_size_summary}\refpo{tab:ch5_effect_size_summary}{은}{는}} 각 경로의 단계별 효과 크기와
분포 형태 간 교차 비교를 종합한 것이다.

\begin{table}[htbp]
  \centering
  \caption{\jrev{효율성 비율 $\eta$의 경로별 효과 크기 종합}}
  \label{tab:ch5_effect_size_summary}
  \small
  \begin{tabular}{ll rr cc l}
    \toprule
    & \textbf{비교 쌍} & $\Delta\bar{\eta}$ & Cohen's $d$ &
    $|d|$ & 해석 & \textbf{실무적 의미} \\
    \midrule
    \multicolumn{7}{l}{\textit{1단: 가우시안 \m{경로(C0 $\to$ C1 $\to$ C2)}}} \\
    & C0 $\to$ C1 & $-0.147$ & $-1.180$ & $1.180$ & \jrev{매우 큰} & 관측성의 복구 단축 \\
    & C1 $\to$ C2 & $+0.830$ & $+1.871$ & $1.871$ & \jrev{매우 큰} & 가우시안 불확실성 프리미엄 \\
    & C0 $\to$ C2 (순효과) & $+0.682$ & $+1.110$ & $1.110$ & \jrev{매우 큰} & 관측성$+$가우시안 결합 \\
    \midrule
    \multicolumn{7}{l}{\textit{2단: 파레토 \m{경로(C0 $\to$ C1 $\to$ C3)}}} \\
    & C0 $\to$ C1 & $-0.147$ & $-1.180$ & $1.180$ & \jrev{매우 큰} & 관측성의 복구 \m{단축(동일)} \\
    & C1 $\to$ C3 & $+0.662$ & $+1.169$ & $1.169$ & \jrev{매우 큰} & 파레토 불확실성 프리미엄 \\
    & C0 $\to$ C3 (순효과) & $+0.514$ & $+0.650$ & $0.650$ & \jrev{중간} & 관측성$+$파레토 결합 \\
    \midrule
    \multicolumn{7}{l}{\textit{3단: 분포 형태 비교}} \\
    & C2 $\leftrightarrow$ C3 & $-0.168$ & $-0.238$ & $0.238$ & \jrev{작은} & 평균 수준 비차별 \\
    \bottomrule
  \end{tabular}
\end{table}

\noindent
경로별 효과 크기에서 다음의 구조적 패턴이 도출된다.

%% --- 1단: 공통 단계 ---
\chg{먼저 C0에서 C1으로의 관측성 복구 단축 효과는 양 경로에 공통으로 작용한다.}
Cohen's $d = -1.180$으로 \jrev{대효과(large effect)}의 음의 효과가 관찰되었다.
\citet{cohen1988power}의 기준에 따르면 $|d| \geq 0.80$이 \jrev{대효과}에
해당하므로, $|d| = 1.180$은 그 기준을 크게 초과하는 극대 효과이다.
음의 방향은 관측성 도입이 $\eta$를 체계적으로 감소시킨다는 것,
즉 $\DRTO$가 기본 복구 시간 $R_0$보다 단축됨을 의미한다.
개별 시그모이드 바운딩 모델에서 C1의 $E_{U,\text{bnd}} = 0$이므로,
$\DRTO_{\text{C1}} = R_0 + E_{O,\text{bnd}} \leq R_0$가 해석적으로 보장되며,
이러한 구조적 필연은 $n = 10{,}000$ 전체 표본에서 예외 없이(100\%) 확인되었다.
이 C0$\to$C1 단계는 가우시안 경로와 파레토 경로 양쪽 모두에 공통으로 작용하는
기저 효과(baseline effect)이다.

%% --- 1단: 가우시안 경로 ---
\chg{1단 가우시안 경로에서 C1에서 C2로의 불확실성 프리미엄을 보면,}
C1$\to$C2 전환의 Cohen's $d = +1.871$은 전체 비교 중 가장 큰 효과 크기이며,
방향이 양(+)으로 반전되었다.
가우시안 불확실성의 도입은 $\bar{\eta}$를 $0.853$에서 $1.682$로
$97.2\%$ 증가시켜, $\DRTO$가 $R_0$를 초과하게 한다.
이 초과분이 \textbf{불확실성 프리미엄(uncertainty premium)}으로서,
복구 과정의 불확실성을 사전에 반영한 안전 마진(safety margin)에 해당한다.
C0$\to$C1의 음의 효과($d = -1.180$)와 C1$\to$C2의 양의 효과($d = +1.871$)가
\jrev{대조적으로 상반되어}, 관측성과 불확실성이 $\DRTO$에 구조적으로 반대 방향의
힘을 가한다는 개별 시그모이드 바운딩 모델의 설계 의도가 실험적으로 확인되었다.
가우시안 경로의 순효과(C0$\to$C2)는 $d = +1.110$으로, 불확실성 프리미엄이
관측성 단축을 압도하여 $\bar{\eta}$가 기준선 $1.0$을 $0.682$만큼 초과한다.

%% --- 2단: 파레토 경로 ---
\chg{2단 파레토 경로에서 C1에서 C3로의 불확실성 프리미엄을 보면,}
C1$\to$C3 전환의 Cohen's $d = +1.169$ 역시 \jrev{대효과}의 양의 효과이다.
파레토 불확실성 프리미엄($\Delta\bar{\eta} = +0.662$)은 가우시안($+0.830$)보다
$20.2\%$ 작으며, 이는 파레토 분포의 우측 비대칭이 평균을 좌측으로 편향시키기
때문이다. 파레토 경로의 순효과(C0$\to$C3)는 $d = +0.650$으로
\jrev{중효과(medium effect)}에 해당하여, 가우시안 경로의 순효과($d = +1.110$)에 비해
관측성 단축의 상대적 기여가 더 크다. 이 차이는 파레토 환경에서
관측성 투자의 실효성이 가우시안 환경보다 높음을 시사한다.

%% --- 3단: 분포 형태 비교 ---
\chg{3단 분포 형태 비교로 C2와 C3의 평균 수준 비차별성을 보면,}
Cohen's $d = -0.238$로 사실상 \jrev{소효과(small effect)}에 해당한다.
C2와 C3의 불확실성 분포는 $E[U_{\text{raw}}] = 3.0$으로 동일한 평균을
가지므로, 평균 수준에서의 차이가 미미한 것은 모델의 이론적 예측과 부합한다.
그러나 이 결과는 파레토 효과가 ``존재하지 않는다''는 것을 의미하지 않는다.
C3의 표준편차($0.791$)가 C2($0.615$)의 $1.29$배로서, 파레토 효과가
평균이 아닌 분산과 극단 영역에 집중되어 있음을 시사한다.
이 점은 \jrev{4.10(이중채널 감쇠 메커니즘)}에서 P85.8 교차점과
이중 채널 감쇠 메커니즘을 통해 상세히 규명된다.

\chg{끝으로 실무적 유의성(practical significance) 측면에서,}
경로별 효과 크기 비교에서 실무적으로 주목할 점은 세 가지이다.
\textbf{첫째}, C0$\to$C1의 $|d| = 1.180$은 양 경로에 공통으로 작용하므로,
불확실성의 유형과 무관하게 관측성 투자가 즉각적이고 체계적인
복구 시간 단축을 제공한다는 정량적 근거이다.
\textbf{둘째}, 불확실성 프리미엄의 효과가 가우시안($d = +1.871$)과
파레토($d = +1.169$) 양쪽 모두에서 매우 크므로, 불확실성을 무시하는 것이
$\DRTO$ 추정에 가장 큰 왜곡을 초래한다.
\textbf{셋째}, 순효과의 격차---가우시안 $d = +1.110$ vs.\ 파레토 $d = +0.650$---는
동일한 관측성 투자가 불확실성 유형에 따라 상이한 순효과를 산출함을 보여주며,
이는 조직의 불확실성 프로파일 진단이 BCM 전략 수립의 전제 조건임을 함의한다.

\jrev{\figref{fig:ch5_cohens_d_bar}\refpo{fig:ch5_cohens_d_bar}{은}{는}} 경로별 Cohen's $d$를
효과 크기 기준선과 함께 시각화한 것이다.
\chg{막대 색상은 C0$\to$C1(녹색, 관측성), C1$\to$C2(파란색, 가우시안 불확실성), C0$\to$C2(연파란, 순효과),
C1$\to$C3(빨간색, 파레토 불확실성), C0$\to$C3(연빨간, 순효과), C2$\leftrightarrow$C3(주황색, 분포 비교)를
나타내며, 수직 점선은 큰 효과 기준 $|d| = 0.80$을, 좌측($d < 0$)은 $\bar{\eta}$ 감소를, 우측($d > 0$)은
$\bar{\eta}$ 증가를 가리킨다.}

\begin{figure}[htbp]
  \centering
  \begin{tikzpicture}[
    lbl/.style={font=\small, anchor=east},
    val/.style={font=\scriptsize\bfseries, anchor=west, xshift=3pt},
    sechead/.style={font=\small\bfseries, anchor=west},
  ]
    %% 좌표 설정: x = Cohen's d, y = 행 (위→아래)
    \def\xscale{2.6}   % cm per unit of d
    \def\xorigin{5.8}  % cm position of d=0
    \def\barht{0.28}    % bar half-height in cm

    %% --- 배경: d=0 기준 좌/우 동일 (3차 수정: 검정→흰색) ---
    \fill[white] (-0.3, 1.20) rectangle (13.0, -8.80);

    %% --- x축 그리드 + 라벨 ---
    \foreach \d in {-2.0,-1.5,-1.0,-0.5,0,0.5,1.0,1.5,2.0} {
      \draw[black, thin] ({\xorigin+\d*\xscale}, 1.20)
        -- ({\xorigin+\d*\xscale}, -8.80);
      \node[font=\tiny, black, below] at ({\xorigin+\d*\xscale}, -8.80) {\d};
    }
    %% --- Cohen's d 축 라벨 (중앙) ---
    \node[font=\small, below] at ({\xorigin}, -9.30) {Cohen's $d$};

    %% --- 좌/우 방향 의미 (축 라벨 아래 별도 행) ---
    \node[font=\scriptsize, black, anchor=east]
      at ({\xorigin-0.5}, -9.30)
      {$\longleftarrow$ $\bar{\eta}$ 감소};
    \node[font=\scriptsize, black, anchor=west]
      at ({\xorigin+0.5}, -9.30)
      {$\bar{\eta}$ 증가 $\longrightarrow$};

    %% --- 효과 크기 기준선 |d|=0.80 ---
    \draw[dashed, black, thick]
      ({\xorigin+0.80*\xscale}, 1.20) -- ({\xorigin+0.80*\xscale}, -8.80);
    \draw[dashed, black, thick]
      ({\xorigin-0.80*\xscale}, 1.20) -- ({\xorigin-0.80*\xscale}, -8.80);
    \node[font=\tiny, black, fill=white, inner sep=1pt]
      at ({\xorigin+0.80*\xscale}, 1.35) {$|d|=0.80$};
    \node[font=\tiny, black, fill=white, inner sep=1pt]
      at ({\xorigin-0.80*\xscale}, 1.35) {$|d|=0.80$};

    %% --- 영 기준선 ---
    \draw[ black, thin] (\xorigin, 1.20) -- (\xorigin, -8.80);

    %% --- 구간 구분선 ---
    \draw[black, densely dotted] (-0.3, -0.05) -- (13.0, -0.05);
    \draw[black, densely dotted] (-0.3, -3.30) -- (13.0, -3.30);
    \draw[black, densely dotted] (-0.3, -6.10) -- (13.0, -6.10);

    %% ============================================================
    %% 공통: C0→C1 (양 경로 공통)
    %% ============================================================
    \node[sechead, green!50!black] at (-0.1, 0.85)
      {공통: 관측성 (C0$\to$C1)};

    \def\yR{0.15}
    \node[lbl] at ({\xorigin-0.25}, \yR) {C0 $\to$ C1};
    \fill[fill=green!45, draw=green!70!black, thick]
      ({\xorigin+(-1.180)*\xscale}, {\yR-\barht})
      rectangle (\xorigin, {\yR+\barht});
    \node[val, anchor=east, xshift=-3pt] at ({\xorigin+(-1.180)*\xscale}, \yR)
      {$-1.180$};

    %% ============================================================
    %% 1단: 가우시안 경로 (C1→C2)
    %% ============================================================
    \node[sechead, blue!60!black] at (-0.1, -0.55)
      {1단: 가우시안 경로};

    %% Row: C1→C2  d = +1.871
    \def\yR{-1.30}
    \node[lbl] at ({\xorigin-0.25}, \yR) {C1 $\to$ C2};
    \fill[blue!60, draw=blue!80, thick]
      (\xorigin, {\yR-\barht})
      rectangle ({\xorigin+(1.871)*\xscale}, {\yR+\barht});
    \node[val] at ({\xorigin+(1.871)*\xscale}, \yR) {$+1.871$};

    %% Row: C0→C2 순효과  d = +1.110
    \def\yR{-2.15}
    \node[lbl] at ({\xorigin-0.25}, \yR) {C0 $\to$ C2};
    \fill[blue!30, draw=blue!50, thick]
      (\xorigin, {\yR-\barht})
      rectangle ({\xorigin+(1.110)*\xscale}, {\yR+\barht});
    \node[val] at ({\xorigin+(1.110)*\xscale}, \yR) {$+1.110$};
    \node[font=\tiny, black, anchor=west] at ({\xorigin+(1.110)*\xscale+1.5}, \yR)
      {(순효과)};

    %% ============================================================
    %% 2단: 파레토 경로 (C1→C3)
    %% ============================================================
    \node[sechead, red!60!black] at (-0.1, -3.80)
      {2단: 파레토 경로};

    %% Row: C1→C3  d = +1.169
    \def\yR{-4.55}
    \node[lbl] at ({\xorigin-0.25}, \yR) {C1 $\to$ C3};
    \fill[red!60, draw=red!80, thick]
      (\xorigin, {\yR-\barht})
      rectangle ({\xorigin+(1.169)*\xscale}, {\yR+\barht});
    \node[val] at ({\xorigin+(1.169)*\xscale}, \yR) {$+1.169$};

    %% Row: C0→C3 순효과  d = +0.650
    \def\yR{-5.40}
    \node[lbl] at ({\xorigin-0.25}, \yR) {C0 $\to$ C3};
    \fill[red!30, draw=red!50, thick]
      (\xorigin, {\yR-\barht})
      rectangle ({\xorigin+(0.650)*\xscale}, {\yR+\barht});
    \node[val] at ({\xorigin+(0.650)*\xscale}, \yR) {$+0.650$};
    \node[font=\tiny, black, anchor=west] at ({\xorigin+(0.650)*\xscale+1.5}, \yR)
      {(순효과)};

    %% ============================================================
    %% 3단: 분포 비교 (C2↔C3)
    %% ============================================================
    \node[sechead, black] at (-0.1, -6.60)
      {3단: 분포 비교};

    %% Row: C2↔C3  d = -0.238
    \def\yR{-7.40}
    \node[lbl] at ({\xorigin-0.85}, \yR) {C2 $\leftrightarrow$ C3};
    \fill[fill=orange!50, draw=orange!70, thick]
      ({\xorigin+(-0.238)*\xscale}, {\yR-\barht})
      rectangle (\xorigin, {\yR+\barht});
    \node[val, anchor=west] at ({\xorigin+0.15}, \yR)
      {$-0.238$};
    \node[font=\tiny, black, anchor=west] at ({\xorigin+1.65}, \yR)
      {(소효과)};

  \end{tikzpicture}
  \caption{\jrev{경로별 Cohen's $d$ 수평 포레스트 플롯}}
  \label{fig:ch5_cohens_d_bar}
\end{figure}
         % 4.8 조건 간 효과 크기


%% ===================================================================
\section{\vrev{회귀 모델의 설명력 구조}}
\label{sec:ch5_r2_progression}
%% ===================================================================

각 조건에서 1차 단순$\to$1차 다중$\to$2차 회귀 분석의 설명력($R^2$) 진행은
모델의 구조적 복잡성이 조건에 따라 어떻게 변화하는지를 보여준다.
이 절에서는 $R^2$ 진행 패턴을 조건별로 비교 분석하고,
2차 회귀 모형이 필요한 구조적 이유를 규명한다.

\jrev{\tabref{tab:ch5_r2_progression}\refpo{tab:ch5_r2_progression}{은}{는}} 모든 조건의 $R^2$ 진행을 종합한 것이다.

\begin{table}[htbp]
  \centering
  \caption{\jrev{조건별 회귀분석 설명력($R^2$) 진행 \m{종합($\kappa = 1.2$)}}}
  \label{tab:ch5_r2_progression}
  \begin{tabular}{lccccc}
    \toprule
    조건 & $n$ & $R^2_{(1\mathrm{s})}$ & $R^2_{(1\mathrm{m})}$ & $R^2_{(2\mathrm{q})}$ & $\Delta R^2$
    ($R^2_{(1\mathrm{m})} \to R^2_{(2\mathrm{q})}$) \\
    \midrule
    C1 (관측성) & 10{,}000 & 0.428 & 0.869 & 1.000 & $+0.131$ \\
    C2 (Gaussian) & 10{,}000 & 0.768 & 0.948 & 0.998 & $+0.050$ \\
    C3 \m{전체(Pareto)} & 10{,}000 & 0.256 & 0.695 & 0.858 & $+0.163$ \\
    \midrule
    C3 Core ($\leq$ P85.8) & 8{,}580 & 0.263 & 0.630 & 0.999 & $+0.369$ \\
    C3 Tail ($>$ P85.8) & 1{,}420 & 0.032 & 0.860 & 0.710 & $-0.150$ \\
    \bottomrule
  \end{tabular}
\end{table}


%% -------------------------------------------------------------------
\subsection{\vrev{1차 단순 회귀에서의 조건별 차이}}
\label{subsec:ch5_simple_r2}
%% -------------------------------------------------------------------

\chg{1차 단순 회귀($\eta$를 $k_O$에 대해 회귀)에서 $R^2$는
C2($0.768$) $>$ C1($0.428$) $>$ C3($0.256$)의 순으로 나타난다. C2에서 $k_O$의
설명력이 가장 높은 것은 불확실성 도입이 관측성 가중치의 효과를
증폭시키기 때문이다.} \jrev{4.4(회귀분석)}\chg{에서 보고하였듯이 $k_O$ 계수가
C1의 $-0.285$에서 C2의 $-1.878$로, \m{절댓값} 기준 약 $6.6$배 커졌다.
반면 C3에서 $R^2 = 0.256$에 그치는 것은 파레토 분포의 극단값이
$k_O$와 $\eta$ 사이의 선형 관계에서 벗어나는 체계적
잔차(systematic residual)를 생성하기 때문이다.}


%% -------------------------------------------------------------------
\subsection{\vrev{1차 다중 회귀에서의 설명력 향상}}
\label{subsec:ch5_multi_r2}
%% -------------------------------------------------------------------

관측성 원시 점수 \jrev{$O_{\text{raw}}$}와 불확실성 변수 \jrev{$U_{\text{raw}}$}를 추가한 1차 다중 회귀에서
설명력은 C1($0.869$), C2($0.948$), C3($0.695$)로 크게 향상된다.
$O$와 $U$의 추가 기여는 조건마다 질적으로 상이하다.

\chg{먼저 C1에서는 $k_O$와 $O$의 계수가 모두 $-0.15$로 동일하여 두 관측성 변수가 동등하게(equally)
기여하며, $O$를 추가함으로써 $R^2$가 $0.428 \to 0.869$로 $+0.441$ 향상되었다. C2에서는 $k_O$의
계수($-1.869$)가 지배적이며 $O$($-0.281$)와 $U$($+0.248$)가 보조적으로 기여하여 관측성 감도가 여전히
핵심 설명 변수임을 확인하며, 1차 다중 회귀의 $R^2 = 0.948$로서 3개 입력 변수의 선형 조합만으로 분산의
$94.8\%$를 설명한다. C3에서는 $U$의 계수($+0.112$)가 C2($+0.248$)의 약 절반(45\%) 수준으로 감소하여,
파레토 불확실성이 선형 모형의 틀에서는 효과적으로 포착되지 않음을 보여준다.}


%% -------------------------------------------------------------------
\subsection{\vrev{2차 회귀 모형의 필요성과 교호 작용 효과}}
\label{subsec:ch5_quad_necessity}
%% -------------------------------------------------------------------

1차 다중 회귀에서 2차 회귀 모형으로의 전환이 필요한 구조적 이유는
개별 시그모이드 바운딩 함수의 비선형성에 있다.
시그모이드 $S(z) = 2/(1+e^{-z})-1$의 테일러 전개에서
$S(z) \approx z - z^3/3 + \cdots$로서, 2차 이상의 비선형 항이 존재한다.
$E_{O,\text{bnd}} = -R_0 \cdot S(\kappa \cdot k_O \cdot O)$에서
$k_O \cdot O$ 곱이 시그모이드의 인자가 되므로,
$\eta$는 $k_O$와 $O$의 교호작용항($k_O \cdot O$)에 의존하는 것이
해석적으로 보장된다.

\chg{첫째, C1과 C2의 완전 설명을 보면,}
C1은 2차 회귀 모형에서 $R^2 = 1.000$, C2는 $R^2 = 0.998$에 도달하여, 개별 시그모이드 바운딩 함수가
이들 조건에서 2차 회귀 다항식으로 사실상 높은 정밀도로 근사됨을 보여준다.
이는 $\kappa = 1.2$에서의 입력 범위($|z| \leq 1.2$)가 시그모이드의
선형--약비선형 영역에 해당하여, 3차 이상의 고차항의 기여가
수치 정밀도 수준에 불과하기 때문이다.
C1에서 핵심 교호작용항 $k_O \cdot O$(계수 $\approx \erf{-0.30}{-0.546}$)를 포함하면
$R^2$가 $0.869 \to 1.000$으로 완전히 향상되며,
이 교호 작용이 관측성 채널의 비선형 구조를 완벽히 포착한다.
C2에서는 $k_O \cdot O$에 더하여 $k_O \cdot U$(계수 $\approx \erf{-0.10}{-0.43}$)
교호작용이 추가로 기여하며, 이 두 교호 작용의 합이 $R^2$를
$0.948 \to 0.998$으로 향상시킨다.

\chg{둘째, C3의 잔여 비설명 분산을 보면,}
C3 전체 표본에서 2차 회귀 모형의 $R^2 = 0.858$로, $14.2\%$의 잔여 분산이
존재한다. 이 잔여 분산은 파레토 분포의 꼬리가
2차 회귀 다항식의 근사 범위를 초과하는 극단값에서 발생한다.
그러나 P85.8 기준 분할 시 Core 영역에서 $R^2 = 0.999$으로
완전 설명이 달성되고, Tail 영역에서도 $R^2 = 0.710$으로
전체 표본($0.858$) 대비 일정 수준 개선된다.
이는 P85.8이 C3의 회귀 구조를 Core와 Tail의 두 가지 질적으로
상이한 체제(regime)로 분리하는 자연적 전환점임을 입증한다.

\jrev{\figref{fig:ch5_r2_heatmap}\refpo{fig:ch5_r2_heatmap}{은}{는}} 조건별$\times$모형 단계별 $R^2$ 값을
히트맵 형태로 시각화한 것이다.
\chg{히트맵에서 녹색은 $R^2 \geq 0.95$를, 주황색은 $0.50 \leq R^2 < 0.95$를, 적색은 $R^2 < 0.50$을 나타낸다.}

\begin{figure}[htbp]
  \centering
  \begin{tikzpicture}[
      cw/.store in=\CW, cw=2.2cm,       % cell width
      ch/.store in=\CH, ch=1.0cm,        % cell height
      hw/.store in=\HW, hw=0.95cm,       % half-width for fill rectangle
    ]
    % Helper macro: \hcell{col}{row}{color}{value}{bold?}
    % col = 1..3, row = 0..4 (0=top), bold = b or empty
    \newcommand{\hcell}[5]{%
      \fill[#3] ({#1*\CW-\HW}, {-#2*\CH-0.4cm})
        rectangle ({#1*\CW+\HW}, {-#2*\CH+0.4cm});
      \ifx b#5%
        \node[font=\small\bfseries] at ({#1*\CW}, {-#2*\CH}) {#4};%
      \else
        \node[font=\small] at ({#1*\CW}, {-#2*\CH}) {#4};%
      \fi
    }

    % Column headers
    \node[font=\small\bfseries] at (1*\CW, 0.8cm) {단순};
    \node[font=\small\bfseries] at (2*\CW, 0.8cm) {다중};
    \node[font=\small\bfseries] at (3*\CW, 0.8cm) {2차};

    % Row labels
    \node[font=\small, anchor=east] at (0.3cm, 0)       {C1};
    \node[font=\small, anchor=east] at (0.3cm, -1*\CH)  {C2};
    \node[font=\small, anchor=east] at (0.3cm, -2*\CH)  {C3 전체};
    \node[font=\small, anchor=east] at (0.3cm, -3*\CH)  {C3 Core};
    \node[font=\small, anchor=east] at (0.3cm, -4*\CH)  {C3 Tail};

    % Row 0: C1
    \hcell{1}{0}{red!25}{0.428}{}
    \hcell{2}{0}{orange!40}{0.869}{}
    \hcell{3}{0}{green!50}{1.000}{b}

    % Row 1: C2
    \hcell{1}{1}{orange!30}{0.768}{}
    \hcell{2}{1}{yellow!50!green!30}{0.948}{}
    \hcell{3}{1}{green!50}{0.998}{b}

    % Row 2: C3 전체
    \hcell{1}{2}{red!40}{0.256}{}
    \hcell{2}{2}{orange!30}{0.695}{}
    \hcell{3}{2}{orange!30}{0.858}{b}

    % Row 3: C3 Core
    \hcell{1}{3}{red!30}{0.263}{}
    \hcell{2}{3}{orange!25}{0.630}{}
    \hcell{3}{3}{green!50}{0.999}{b}

    % Row 4: C3 Tail
    \hcell{1}{4}{red!50}{0.032}{}
    \hcell{2}{4}{orange!40}{0.860}{}
    \hcell{3}{4}{orange!30}{0.710}{b}

    % H4 threshold line
    \draw[dashed, red!60, thick]
      (0.5cm, {-5*\CH+0.2cm}) -- ({3*\CW+1.0cm}, {-5*\CH+0.2cm});
    \node[font=\scriptsize, red!60, anchor=west]
      at ({3*\CW+1.1cm}, {-5*\CH+0.2cm}) {H4 기준 ($R^2_{(2\mathrm{q})} \geq 0.95$)};

    % Color legend
    \fill[red!35]    (0.5cm, {-5.8*\CH}) rectangle (1.1cm, {-5.8*\CH+0.35cm});
    \node[font=\scriptsize, anchor=west] at (1.2cm, {-5.8*\CH+0.18cm}) {$R^2 < 0.50$};
    \fill[orange!35] (3.5cm, {-5.8*\CH}) rectangle (4.1cm, {-5.8*\CH+0.35cm});
    \node[font=\scriptsize, anchor=west] at (4.2cm, {-5.8*\CH+0.18cm}) {$0.50 \leq R^2 < 0.95$};
    \fill[green!50]  (7.5cm, {-5.8*\CH}) rectangle (8.1cm, {-5.8*\CH+0.35cm});
    \node[font=\scriptsize, anchor=west] at (8.2cm, {-5.8*\CH+0.18cm}) {$R^2 \geq 0.95$};
  \end{tikzpicture}
  \caption{\jrev{조건별$\times$모형 단계별 $R^2$ 히트맵}}
  \label{fig:ch5_r2_heatmap}
\end{figure}

\noindent
\jrev{\figref{fig:ch5_r2_heatmap}}에서 다음의 구조적 패턴이 관찰된다.
\chg{첫째,} 모든 조건에서 1차 단순$\to$1차 다중$\to$2차 회귀로 갈수록 $R^2$가 단조 증가한다.
\chg{둘째,} C1은 \pdferr{2차 회귀} 모형에서 녹색($R^2 = 1.000$), C2는 $R^2 = 0.998$에 도달하여
개별 시그모이드 바운딩 함수의 높은 정밀도 근사가 확인된다.
\chg{셋째,} C3 전체의 2차 회귀 $R^2 = 0.858$는 녹색 기준($0.95$)에 미달하나,
분할 시 Core가 녹색에 도달하고 Tail은 $0.710$으로 잔여 비선형성이 존재한다.
이러한 패턴은 P85.8 분할 회귀의 이론적 정당성을 시각적으로 입증한다.
      % 4.9 회귀 설명력 구조


%% ===================================================================
\section{\vrev{이중채널 감쇠 메커니즘}}
\label{sec:ch5_dual_channel}
%% ===================================================================

C2--C3 비교에서 Cohen's $d = -0.238$라는 평균 수준의 비차별성과,
C3의 표준편차가 C2의 $1.29$배라는 분산 수준의 차이 사이에는
\jrev{구조적 패러독스(structural paradox)}가 존재한다. 이 절에서는 이 \jrev{패러독스}의 원인을 P85.8 교차점에서의
분포 역전과 이중채널 감쇠 메커니즘으로 해명한다.


%% -------------------------------------------------------------------
\subsection{\vrev{P85.8 교차점(분포 역전의 전환점)}}
\label{subsec:ch5_crossover}
%% -------------------------------------------------------------------

C2와 C3의 $\eta$ \jrev{백분위 함수가 교차하는 지점\chg{(}$\text{VaR}_p^{\text{C2}} = \text{VaR}_p^{\text{C3}}$를}
만족하는 백분위수 $p$\chg{)}은 약 \textbf{P85.8}에서 관측된다.
이 교차점 이하에서는 \jrev{$\text{VaR}_p^{\text{C2}} > \text{VaR}_p^{\text{C3}}$}, 즉 C2가 C3보다
높은 $\eta$를 보이며, 이는 가우시안 분포의 중심 집중이 파레토보다
높은 중앙값을 산출하기 때문이다.
교차점 이상에서는 \jrev{$\text{VaR}_p^{\text{C3}} > \text{VaR}_p^{\text{C2}}$}로 역전되어,
파레토 효과가 급격히 발현한다.

이 구조적 역전의 실무적 함의는 중대하다.
전체 표본의 $85.8\%$에 해당하는 일상적 운영 상황에서는 C2와 C3의 차이가 미미하여
조직은 두 불확실성 유형을 구분할 실무적 필요가 없다.
그러나 상위 $14.2\%$의 극단 상황에서는 파레토 불확실성이 복구 시간을
급격히 연장한다. $\mathrm{CVaR}_{99}$ 기준 C3가 C2 대비 $+24.0\%$ 높은 것은
이 극단 영역의 격차를 정량화한 것이다.

이러한 ``평균에서의 무차별$+$극단에서의 극적 분기'' 패턴은
BCM 실무에서 극단 사건에 대한 별도의 비상 계획\chg{(}CVaR 기반의
꼬리 리스크 관리\chg{)}이 필요하다는 이론적 근거를 제공한다.


%% -------------------------------------------------------------------
\subsection{\vrev{분할 회귀와 $k_O$ 계수 감쇠}}
\label{subsec:ch5_split_regression}
%% -------------------------------------------------------------------

P85.8 교차점을 기준으로 C3의 표본을 Core($\leq$ P85.8, $n = 8{,}580$)와
Tail($>$ P85.8, $n = 1{,}420$)로 분할하여 회귀 분석을 수행하면,
관측성 가중치 $k_O$의 효과가 두 영역에서 질적으로 상이한 구조를 보인다.

Core 영역의 1차 단순 회귀에서 $k_O$ 계수는 $-0.798$이다.
$k_O$가 1단위 증가하면 $\eta$가 $0.798$만큼 감소하며,
이는 $\DRTO$가 $0.798 \times R_0 = 4.79$시간 단축됨을 의미한다.
Tail 영역의 1차 단순 회귀에서 $k_O$ 계수는 $-0.415$이다.
동일한 $k_O$ 1단위 증가가 Tail에서는 $\eta$를 $0.415$만큼 감소시켜,
$\DRTO$가 $0.415 \times R_0 = 2.49$시간 단축된다.

Tail 대비 Core의 계수 비율은 $|-0.415|/|-0.798| \approx 0.52$배이다.
이 $0.52$배 감쇠가 \textbf{이중채널 감쇠(dual-channel attenuation)}의
핵심 지표이며, 관측성 가중치와 파레토 불확실성이 결합될 때
극단 영역에서 관측성 효과가 상대적으로 약화되는 구조적 증거이다.


%% -------------------------------------------------------------------
\subsection{\vrev{이중채널 감쇠의 구조적 원인}}
\label{subsec:ch5_mechanism}
%% -------------------------------------------------------------------

이중채널 감쇠의 구조적 원인은 개별 시그모이드 바운딩 함수의 비선형성에 있다.

\chg{먼저, Core 채널은 선형 영역에서 작동한다.}
Core 영역($\leq$ P85.8)에서는 불확실성 $U_{\text{raw}}$의 값이 상대적으로 작아
$z_U = \kappa \cdot k_U \cdot R_0 \cdot U_{\text{raw}} / (R_{\max} - R_0)$이 시그모이드의 선형 근사 영역($|z| \ll 1$)에
위치한다. 이 영역에서 $S(z) \approx z$이므로,
$E_{U,\text{bnd}} \approx (R_{\max} - R_0) \cdot z_U$로 불확실성의 효과가
입력에 거의 비례한다. $k_O$의 효과는 $E_{O,\text{bnd}}$를 통해
단조적으로 전달되며, 관측성의 복구 단축 효과가 충분히 발현된다.

\chg{반면에, Tail 채널은 비선형 구간에서 포화 효과를 겪는다.}
Tail 영역($>$ P85.8)에서는 큰 $U_{\text{raw}}$ 값(파레토 분포의 우측 꼬리)에 의해
$z_U$가 시그모이드의 곡률(curvature)이 큰 영역에 진입한다.
이 상태에서 $E_{U,\text{bnd}}$가 시그모이드의 포화에 근접하여
불확실성 채널이 $\DRTO$를 지배하게 되며,
$k_O$의 변화가 $E_{O,\text{bnd}}$를 통해 전달되는 관측성의
복구 단축 효과가 불확실성의 상향 효과에 의해 상대적으로 약화된다.
구체적으로, $\DRTO = R_0 + E_{O,\text{bnd}} + E_{U,\text{bnd}}$에서
$E_{U,\text{bnd}}$가 지배적일수록 $E_{O,\text{bnd}}$의 상대적 기여가 감소한다.

정리하면, 파레토 불확실성은 시그모이드의 비선형 구간을 활성화하여
불확실성 채널을 지배적으로 만들며, 이로 인해 관측성의 상대적 역할이 축소된다.
이는 극단 불확실성 환경에서 관측성 투자만으로는 복구 시간 단축의
효과가 제한적임을 의미하며, 불확실성 자체의 관리가 우선되어야 한다는
실무적 함의로 직결된다.


%% -------------------------------------------------------------------
\subsection{\vrev{분포별 회귀 구조 차이(지배 변수의 전환)}}
\label{subsec:ch5_dominant_variable}
%% -------------------------------------------------------------------

4.4(회귀분석 결과)의 분석에서 C2와 C3의 변수 진입 순서가 질적으로
다르다는 발견은 교차 분석 관점에서 중요한 의미를 갖는다.
\jrev{\tabref{tab:ch5_dominant_var}\refpo{tab:ch5_dominant_var}{은}{는}} 두 조건의 변수 진입 1순위 및
2순위까지의 누적 $R^2$를 비교한 것이다.

\begin{table}[htbp]
  \centering
  \caption{\jrev{C2 vs.\ C3 전진 선택법 변수 진입 구조 비교}}
  \label{tab:ch5_dominant_var}
  \begin{tabular}{l cc cc}
    \toprule
    & \multicolumn{2}{c}{\textbf{C2 (Gaussian)}} & \multicolumn{2}{c}{\textbf{C3 (Pareto)}} \\
    \cmidrule(lr){2-3} \cmidrule(lr){4-5}
    진입 순서 & 변수 & 누적 $R^2$ & 변수 & 누적 $R^2$ \\
    \midrule
    1순위 & $k_O$ & 0.768 & $U$ & 0.424 \\
    2순위 & $U$ & 0.931 & $k_O$ & 0.686 \\
    \bottomrule
  \end{tabular}
\end{table}

\chg{지배 변수 전환의 의미를 보면,}
C2에서 $k_O$가 1순위($R^2 = 0.768$)인 것은 가우시안 불확실성의
분산이 상대적으로 작아($\mathrm{SD} = 1.0$) $\eta$ 변동의 주된 원천이
관측성임을 반영한다. 반면, C3에서 $U$가 1순위($R^2 = 0.424$)인 것은
파레토 분포의 높은 분산(\jrev{$U$ 값의 이론적 분산이 가우시안의 $27$배})이
$\eta$ 변동에서 불확실성의 역할을 증대시키기 때문이다.

\chg{2순위 투입 후 설명력 격차를 보면,}
C2에서 2순위($U$) 투입 후 $R^2 = 0.931$에 도달하는 반면,
C3에서 2순위($k_O$) 투입 후 $R^2 = 0.686$에 머문다.
이 $0.245$의 격차는 C3의 $\eta$ 변동 중 상당 부분이
비선형 교호 작용에 의해 설명됨을 시사하며,
\pdferr{2차 회귀} 모형의 필요성을 교차 분석 관점에서 뒷받침한다.

\chg{실무적 함의로서,}
이러한 지배 변수 전환은 BCM 실무에서 중요한 의사결정 원칙을 시사한다.
가우시안 불확실성(일상적 변동) 환경에서는 \textbf{관측성 투자가
최우선 전략}인 반면, 파레토 불확실성(극단 사건 가능성) 환경에서는
\textbf{불확실성 정량화가 최우선 전략}이다.
두 전략의 우선순위는 조직이 직면하는 불확실성의 분포 형태에 의해 결정되며,
이는 산업 특성에 따라 차별화되어야 한다.


\jrev{\figref{fig:ch5_dual_channel_mechanism}\refpo{fig:ch5_dual_channel_mechanism}{은}{는}} 이중 채널 감쇠의
메커니즘을 도식으로 나타낸 것이다.
\chg{그림은 P85.8 기준으로 분할한 후 Tail 채널에서 $k_O$ 회귀 계수가 Core 대비 $0.52$배로 감쇠함을 보인다.}

\begin{figure}[htbp]
  \centering
  \resizebox{\textwidth}{!}{%
  \begin{tikzpicture}[
    channel/.style={draw, rounded corners=5pt, minimum width=52mm, minimum height=14mm,
      align=center, font=\small, line width=0.6pt},
    input/.style={draw, rounded corners=4pt, minimum width=38mm, minimum height=11mm,
      align=center, font=\small, line width=0.6pt},
    result/.style={draw, rounded corners=5pt, minimum width=55mm, minimum height=14mm,
      align=center, font=\small\bfseries, line width=0.6pt},
    arr/.style={-{Stealth[length=2.5mm]}, thick},
    splitlbl/.style={font=\scriptsize, fill=white, inner sep=2pt,
      rounded corners=1pt},
    ]

    % Input variables
    \node[input, fill=blue!15] (ko) at (0, 4)   {관측성 가중치\\$k_O \in [0,1]$};
    \node[input, fill=green!15] (O) at (0, 2.2)  {관측성 원시 점수\\$O_{\text{raw}} \in [0,1]$};
    \node[input, fill=orange!15] (U) at (0, 0.4) {불확실성\\$U_{\text{raw}}^{(P)} \sim 6 \cdot \mathrm{Pareto}(\alpha\!=\!3)$};

    % Sigmoid bounding
    \node[channel, fill=gray!15] (sig) at (7, 2.2) {개별 시그모이드 바운딩\\$\DRTO = R_0 + E_{O,\text{bnd}} + E_{U,\text{bnd}}$};

    % Fork point
    \coordinate (fork) at (10.0, 2.2);

    % Split channels
    \node[channel, fill=blue!10] (core) at (14.5, 4.2)
      {\textbf{Core 채널} ($\leq$ P85.8)\\$k_O$ 계수 $= -0.798$,\; $R^2 = 0.999$};
    \node[channel, fill=red!10] (tail) at (14.5, 0.2)
      {\textbf{Tail 채널} ($>$ P85.8)\\$k_O$ 계수 $= -0.415$,\; $R^2 = 0.710$};

    % Attenuation result
    \node[result, fill=yellow!15] (amp) at (14.5, -2.2)
      {이중채널 감쇠 비율\\$|-0.415| / |-0.798| \approx 0.52$배};

    % Input arrows — 3개 모두 바운딩 박스 좌측면(west)으로 수렴
    \draw[arr, blue!60, line width=1.2pt]
      (ko.east) -- ++(1.8, 0) |- (sig.170);
    \draw[arr, green!50!black, line width=1.2pt]
      (O.east) -- (sig.west);
    \draw[arr, orange!60, line width=1.2pt]
      (U.east) -- ++(1.8, 0) |- (sig.190);
    % Arrow labels
    \node[font=\tiny, fill=white, inner sep=1pt, text=black] at (3.5, 3.3) {관측성 ($k_O$)};
    \node[font=\tiny, fill=white, inner sep=1pt, text=black] at (3.5, 2.2) {관측성 ($O$)};
    \node[font=\tiny, fill=white, inner sep=1pt, text=black] at (3.5, 1.1) {불확실성 ($U$)};

    % Sigmoid to fork
    \draw[arr] (sig.east) -- (fork);
    \draw[arr] (fork) |- (core.west);
    \draw[arr] (fork) |- (tail.west);

    % P85.8 label
    \node[splitlbl] at (10.0, 3.3) {P85.8 기준};
    \node[splitlbl] at (10.0, 2.85) {분할};

    % Tail to amplification (직선)
    \draw[arr, dashed, red!60, line width=1.2pt] (tail.south) -- (amp.north);
    % Core to amplification (우측 우회: core → 오른쪽 → amp 우측)
    \draw[arr, dashed, blue!40]
      (core.east) -- ++(1.2, 0) |- (amp.east);

    % Key insight annotation
    \node[font=\scriptsize, text=black, align=center, anchor=west] at (16.5, -0.8)
      {Tail에서 $k_O$\\$0.52$배 감쇠};
  \end{tikzpicture}
  }% end resizebox
  \caption{\jrev{이중채널 감쇠 메커니즘}}
  \label{fig:ch5_dual_channel_mechanism}
\end{figure}


%% -------------------------------------------------------------------
\subsection{\vrev{이중 채널 감쇠의 실무적 함의}}
\label{subsec:ch5_dual_practical}
%% -------------------------------------------------------------------

이중채널 감쇠 효과의 실무적 함의는 다음의 세 가지로 정리된다.

\chg{첫째, 관측성 투자의 조건부 수익 구조이다. 일상 운영(Core)에서 관측성 가중치 $k_O$를 $0.5$에서
$0.6$으로 $0.1$만큼 개선하면 $\eta$가 $0.080$만큼 감소하며($0.798 \times 0.1 = 0.080$, $\DRTO$로 환산하면
$0.48$시간), 극단 상황(Tail)에서 동일한 $k_O$ 개선은 $\eta$를 $0.042$만큼 감소시킨다($0.415 \times 0.1 = 0.042$,
$\DRTO$로 환산하면 $0.25$시간). 즉 극단 상황에서 관측성 투자의 한계 수익은 일상 운영의 $0.52$배로
감쇠된다. 둘째, 이중축 BCM 전략의 근거이다. P85.8 교차점은 BCM 전략을 일상 운영 축(C2 기반)과 극단
시나리오 축(C3 기반)으로 분리하는 자연적 경계점이며, $\mathrm{CVaR}_{99}$ 기준 $+24.0\%$의 격차는 극단
시나리오에 대비한 추가 자원 확보의 정량적 근거가 된다. 셋째, 스트레스 테스트 설계 기준이다. $0.52$배
감쇠 효과는 BCM 스트레스 테스트에서 Tail 영역의 시나리오를 별도로 설계해야 함을 의미하며, 극단
상황에서 관측성 투자의 효과가 감소하므로 불확실성 관리 자체에 초점을 맞춘 극단 시나리오 대응 전략이
필수적이다.}

\chg{H8의 검증 결론을 정리하면,} \mrev{H8은 ``Tail ($>$P85.8) 영역에서 $k_O$ 계수의 절댓값이 Core 영역의
$\leq 1.0$배''를
기준으로 한다. Tail $|\beta_{k_O}|
= 0.415$ vs.\ Core $|\beta_{k_O}| = 0.798$의 비율이 $0.415/0.798 = 0.52
\leq 1.0$이므로, \textbf{H8은 채택된다}. 이 감쇠 효과는 분포 환경에 의존하는
관측성 효과의 비선형성을 정량적으로 입증한 본 연구의 독자적
발견이다.}
       % 4.10 이중채널 감쇠


%% ===================================================================
\section{\vrev{전문가 패널 검증}}
\label{sec:ch5_expert}
%% ===================================================================

%% -------------------------------------------------------------------
\subsection{\vrev{검증 방법론}}
\label{subsec:ch5_expert_method}
%% -------------------------------------------------------------------

$\DRTO$ 프레임워크의 실무 타당성을 검증하기 위하여
\textbf{전문가 패널 검증(Expert Panel Validation)}을 수행하였다.
본 검증은 정량적 설문과 정성적 심층 인터뷰를 병행하는
\textbf{혼합 연구 방법(Mixed Methods Research)}으로 설계되었으며,
두 방법의 결과를 수렴적으로 \chg{비교하고 통합하는}
\textbf{방법론적 삼각검증(Methodological Triangulation)}%
~\citep{creswell2018designing, denzin1978research}을 적용하였다.

\chg{검증 방법론의 위치를 짚어 두면,}
본 연구의 전문가 검증은 사회과학적 양적 연구에서 잠재 변수 간
인과 관계를 규명하는 구조방정식 모델링(SEM)이나 \chg{탐색적, 확인적}
요인 분석과 목적이 근본적으로 다르다.
$\DRTO$ 프레임워크의 핵심 변수 간 관계\chg{(}관측성($O_{\text{raw}}$)과
$\DRTO$의 음의 상관, 불확실성($U_{\text{raw}}$)과 $\DRTO$의 양의 상관\chg{)}는
시그모이드 바운딩 함수 $S(z)$에 의해 수학적으로 정의되어 있으며,
$n = 10{,}000$ 시뮬레이션을 통해 통계적으로 검증되었다(\jrev{4.1--4.10절}).
따라서 전문가 검증의 목적은 잠재 변수 간 관계 탐색이 아니라,
수학적으로 수립된 모델이 실무 환경에서 타당하고 적용 가능한지를
확인하는 \textbf{내용 타당도 검증(content validation)}이다.
이는 프레임워크 및 모델 타당성 검증의 표준적 접근으로,
전문가 패널 평가(\pdel{CVI/CVR}\padd{CVI})와 혼합연구법이 적합한 방법론이다%
~\citep{lynn1986determination}\pdel{,~\citep{lawshe1975quantitative}}.

\chg{전문가 패널 구성을 보면,}
BCM/DR 분야 전문가 5명을 대상으로 하였다.
패널은 BCM/BCP 컨설팅(E1, 13년), DR 운영(E2, 13년),
통합 리스크 관리(E3, 20년+), IT 인프라(E4, 20년+),
정보보안(E5, 10년)의 상호 보완적 전문 영역을 포괄하며,
평균 경력 $15.2$년이다.

%% [분량보존 복원] 구 부록 I 유일분 (전문가 프로필표 / 신뢰도 산출공식)
\subsection{전문가 프로필}
\label{app:expert-profiles}

전문가 평가에 참여한 5인의 프로필을 \p{[}표~\ref{tab:app-expert-profile}\p{]}에 정리하였다.

\begin{table}[htbp]
\centering
\caption{전문가 평가 참여자 프로필}
\label{tab:app-expert-profile}
\small
\begin{tabular}{clll}
\toprule
\textbf{ID} & \textbf{전문 분야} & \textbf{경력} & {\textbf{소속}} \\
\midrule
E1 & {BCM/BCP 컨설팅}      & {13년}       & {---} \\
E2 & {DR 운영관리}          & {13년}       & {---} \\
E3 & {통합리스크관리}        & {20년 이상}  & {---} \\
E4 & {IT 인프라 운영}        & {20년 이상}  & {---} \\
E5 & {정보보안/사이버복원력}  & {10년}       & {---} \\
\bottomrule
\end{tabular}
\end{table}

\chg{전문가 선정 기준은 BCM/DR 분야 경력 10년 이상이며, 평가 방식은 구조화된 설문(4개 영역, 27문항,
5점 리커트 척도)과 개방형 5문항(서술형), 심층 인터뷰 5건(\jrev{약 30분/인})으로 구성된다. 평가 기간은
2026년 3월 8일부터 4월 18일까지로 사전 브리핑부터 결과 종합까지 약 6주가 소요되었으며(인터뷰는 2026년 3월 15일–18일 실시), 사전 브리핑에서는 $\DRTO$ 프레임워크 설명 자료를 이메일로 사전제공하고 개별적으로 추가설명을 실시하였다.}




\chg{정량적 도구로는 구조화 설문을 사용하였다.}
27문항(4개 영역)으로 구성된 5점 리커트 척도 설문을
실시하였으며, 평가 영역은 \chg{모델 타당성(Section~B, 6문항), 파라미터 적정성(Section~C, 6문항),
실무 적용성(Section~D, 6문항), 산업적용과 $\eta$ 활용 체계(Section~E, 9문항)이다.}

\chg{정성적 도구로는 반구조화 심층 인터뷰를 사용하였다.}
전문가별 약 30분 정도의 반구조화(semi-structured) 심층 인터뷰를 실시하였으며,
인터뷰 가이드는 \chg{모델 구조에 대한 전반적 인상, 파라미터 설정의
실무적 합리성, 도입 가능성과 장애 요인, 개선 방향의}
4개 범주로 구성하였다.

\chg{\p{[}그림~\vref{fig:expert-verification-process}\p{]}는 전문가 패널 검증의 전체 절차를 안내 메일 발송부터 결과 종합까지 6단계로 체계화한 것이다. 설문은 모델 타당성(B), 파라미터 적정성(C), 실무 적용성(D)에 더하여 산업 적용과 $\eta$ 활용(E)을 검증하며, 약 30분의 심층 인터뷰로 정량적 설문 결과를 보충하는 혼합 연구방법(Mixed Methods) 삼각 검증 설계를 따른다.}

\begin{figure}[htbp]
\centering
\resizebox{\textwidth}{!}{%
\begin{tikzpicture}[
  phase/.style={draw, rounded corners=5pt, minimum width=32mm, minimum height=14mm,
    align=center, font=\small\bfseries, line width=0.7pt},
  check/.style={font=\small, anchor=west, text width=62mm},
  arr/.style={-{Stealth[length=2.5mm]}, thick, black},
]

% Phase boxes (vertical)
\node[phase, fill=blue!12]           (p1) at (0,  0)    {1. 안내 메일 발송\\{\footnotesize D+0}};
\node[phase, fill=teal!12]           (p2) at (0, -2.2)  {2. 설문·일정 수령\\{\footnotesize D+0\,--\,D+7}};
\node[phase, fill=orange!12]         (p3) at (0, -4.4)  {3. 일정 확정\\{\footnotesize D+7\,--\,D+10}};
\node[phase, fill=red!12]            (p4) at (0, -6.6)  {4. 인터뷰 실시\\{\footnotesize D+10\,--\,D+21}};
\node[phase, fill=violet!12]         (p5) at (0, -8.8)  {5. 전사·분석\\{\footnotesize D+21\,--\,D+35}};
\node[phase, fill=green!12]          (p6) at (0, -11.0) {6. 결과 종합\\{\footnotesize D+35\,--\,D+42}};

% Arrows between phases
\foreach \i/\j in {p1/p2, p2/p3, p3/p4, p4/p5, p5/p6}
  \draw[arr] (\i) -- (\j);

% Checklist items
\node[check] at (2.8,  0.3)  {$\square$ 사전배포자료 + 설문지 첨부};
\node[check] at (2.8, -0.3)  {$\square$ 인터뷰 일정 조사표 첨부 (2주간, 오전/오후/저녁)};

\node[check] at (2.8, -1.9)  {$\square$ 설문 응답 5인분 수령 확인};
\node[check] at (2.8, -2.5)  {$\square$ 가능 일정 종합 대조 → 상호 조율};

\node[check] at (2.8, -4.1)  {$\square$ 확정 일자·시간·장소(또는 화상회의) 통보};
\node[check] at (2.8, -4.7)  {$\square$ 녹음 동의서 5인분 수령};

\node[check] at (2.8, -6.0)  {$\square$ 전일 리마인더 발송};
\node[check] at (2.8, -6.6)  {$\square$ 인터뷰 5회 완료 (각 약 30분)};
\node[check] at (2.8, -7.2)  {$\square$ 녹음 파일 백업 (메인 + 백업 레코더)};

\node[check] at (2.8, -8.5)  {$\square$ 전사록 5건 작성 + Member Checking};
\node[check] at (2.8, -9.1)  {$\square$ 정량 분석: 평균, CVI, ICC};
\node[check] at (2.8, -9.7)  {$\square$ 질적 분석: 주제 분석 (Braun \& Clarke 6단계)};

\node[check] at (2.8, -10.7) {$\square$ 결과종합자료 작성};
\node[check] at (2.8, -11.3) {$\square$ 본 연구 논문에 반영};

% Survey structure box (right side)
\node[draw, rounded corners=4pt, fill=blue!5, font=\small, align=left,
  minimum width=55mm, inner sep=6pt, anchor=north west] (survey) at (9.5, 0.5) {
  \textbf{설문 구성 (A--F)}\\[3pt]
  A. 응답자 정보\\
  B. 모델 타당성 (6문항)\\
  C. 파라미터 적정성 (6문항)\\
  D. 실무 적용성 (6문항)\\
  \textbf{\color{black}E. 산업적용·$\eta$ 활용 (9문항)}\\
  F. 개방형 (5문항)
};

% Interview structure box (right side, below survey)
\node[draw, rounded corners=4pt, fill=orange!5, font=\small, align=left,
  minimum width=55mm, inner sep=6pt, anchor=north west] (intv) at (9.5, -4.0) {
  \textbf{인터뷰 스크립트 (4Phase)}\\[3pt]
  P1. 인사·사전확인 (5분, 녹음 전)\\
  P2. 설문 보충 질의 (10분, \textbf{녹음 시작})\\
  P3. 전문 분야별 심화 (10분)\\
  P4. 종합 의견·마무리 (5분, \textbf{녹음 종료})
};

% Analysis box (right side, below interview)
\node[draw, rounded corners=4pt, fill=green!5, font=\small, align=left,
  minimum width=55mm, inner sep=6pt, anchor=north west] (anal) at (9.5, -8.0) {
  \textbf{분석 방법}\\[3pt]
  정량: 평균, SD, CVI ($\geq 0.80$),\\
  \hspace{2.2em} ICC, 단일표본 $t$-검정\\
  질적: 주제 분석, 교차 분석,\\
  \hspace{2.2em} 코딩 체계 (연역 + 귀납)
};

% Connecting dotted lines from phases to right boxes
\draw[dashed, black, line width=0.6pt] (p2.east) -- (survey.west |- p2.east);
\draw[dashed, black, line width=0.6pt] (p4.east) -- (intv.west |- p4.east);
\draw[dashed, black, line width=0.6pt] (p5.east) -- (anal.west |- p5.east);

% Ethics note at bottom
\node[draw, rounded corners=3pt, fill=gray!15, font=\footnotesize, align=center,
  minimum width=140mm, inner sep=5pt] at (7.0, -12.5) {
  \textbf{윤리적 고려}: 서면 동의 $\mid$ 익명성 보장 (E1--E5 코드명) $\mid$
  녹음 파일 암호화 $\mid$ 자발적 참여 $\mid$ 연구 종료 후 파기
};

\end{tikzpicture}
}% end resizebox
\caption{전문가 검증 프로세스 체계도}
\label{fig:expert-verification-process}
\end{figure}



%% -------------------------------------------------------------------
\subsection{\vrev{검증 도구의 타당도 및 신뢰도}}
\label{subsec:ch5_validity_reliability}
%% -------------------------------------------------------------------

전문가 검증 결과의 학술적 엄밀성을 확보하기 위하여,
정량적 도구(설문)에 대해서는 타당도(validity)와 신뢰도(reliability)를,
정성적 도구(인터뷰)에 대해서는 신뢰성(trustworthiness)을 체계적으로 평가하였다.

\chg{먼저 내용 타당도(Content Validity)를 살펴본다.}
설문 도구의 내용 타당도는
\textbf{내용 타당도 지수(Content Validity Index, CVI)}%
~\citep{lynn1986determination}\pdel{와
\textbf{내용타당도비율(Content Validity Ratio, CVR)}%
~\citep{lawshe1975quantitative}}로 평가하였다.
\padd{내용 타당도 비율(CVR)~\citep{lawshe1975quantitative}은 Lawshe(1975) 원표가 패널 5명에서
임계값 $\text{CVR}_{\min} = 0.99$를 요구하고 통상 10명 이상 패널을 전제로 설계된 지표로서,
본 연구의 5명 전문가 패널에는 임계값 적용이 부적합하므로 주 타당도 지표로 채택하지 않았다
(패널 규모 한계는 \jrev{\secref{sec:limitations}}항 참조).}

5점 리커트 척도에서 4점 이상을 ``적합(relevant)''으로 이분화하여
문항별 I-CVI와 전체 S-CVI를 산출하였다.
\jrev{\tabref{tab:ch5_cvi}}에 제시된 바와 같이,
27개 문항 중 25개에서 모든 전문가($n = 5$)가 4점 이상을 부여하여
I-CVI $= 1.00$이며, D1과 E4 문항은 1인(E1)만 3점으로 I-CVI $= 0.80$이다.
S-CVI/Ave $= 0.985$로 우수한 수준이다.

\begin{table}[htbp]
  \centering
  \caption{\jrev{내용 타당도 분석 결과}}
  \label{tab:ch5_cvi}
  \small
  \begin{tabular}{lcccl}
    \toprule
    지표 & 산출 방법 & 값 & 기준 & 판정 \\
    \midrule
    I-CVI $= 1.00$ 문항
      & 전원 4점 이상
      & 25/27 (92.6\%)
      & ---
      & --- \\[4pt]
    I-CVI $\geq 0.80$ 문항
      & 4/5인 이상 4점 이상
      & 27/27 (100\%)
      & $\geq 0.78$
      & \cmark 전 문항 충족 \\[4pt]
    S-CVI/Ave (전체)
      & I-CVI 평균
      & 0.985
      & $\geq 0.90$
      & \cmark 우수 \\
    \bottomrule
  \end{tabular}
\end{table}

\chg{여기서 I-CVI 기준은 3인 이상 패널에서 $\geq 0.78$이고~\citep{lynn1986determination}, S-CVI/Ave 기준은
$\geq 0.90$이며~\citep{polit2006content}, D1(데이터 수집 가능성)과 E4($\eta$ 에스컬레이션) 문항은 E1
전문가만 3점을 부여하여 I-CVI $= 0.80$이다.}

\chg{다음으로 내적 일관성 신뢰도(Internal Consistency)는}
설문에 대해 Cronbach's $\alpha$로 평가하여 \jrev{\tabref{tab:ch5_alpha}}에 나타내었다.
전체 Cronbach's $\alpha = 0.89$로 양호한 수준이다.
하위 영역별로는 E 섹션($\alpha = 0.80$)과 D 섹션($\alpha = 0.70$)이 기준을 충족하나,
B 섹션($\alpha = 0.58$)과 C 섹션($\alpha = 0.51$)은 기준($0.70$) 미달이다.
이는 소규모 패널($n = 5$)에서 응답의 98.5\%가 4--5점에 집중된 천장효과에 기인하며,
분산이 극히 작아 $\alpha$가 역설적으로 낮아지는 현상이다.
전체 $\alpha = 0.89$는 양호하므로, 전체 $\alpha$를 주 지표로 보고한다.

\begin{table}[htbp]
  \centering
  \caption{\jrev{내적 일관성 신뢰도 \m{분석(Cronbach's $\alpha$)}}}
  \label{tab:ch5_alpha}
  \small
  \begin{tabular}{lccl}
    \toprule
    평가 영역 & 문항 수 & Cronbach's $\alpha$ & 판정 \\
    \midrule
    모델 \m{타당성(Section B)}             & 6  & 0.58 & $\triangle$ 미달$^{**}$ \\
    파라미터 \m{적정성(Section C)}         & 6  & 0.51 & $\triangle$ 미달$^{**}$ \\
    실무 \m{적용성(Section D)}             & 6  & 0.70 & 경계 충족 \\
    산업적용·$\eta$ \m{활용(Section E)}     & 9  & 0.80 & \cmark 양호 \\
    \midrule
    \textbf{전체}                       & 27 & \textbf{0.89} & \cmark \textbf{양호} \\
    \bottomrule
  \end{tabular}
\end{table}

\chg{해석 기준은 $\geq 0.90$이면 우수, $\geq 0.80$이면 양호, $\geq 0.70$이면 수용 가능이며~\citep{nunnally1994psychometric},
B와 C 섹션의 미달($^{**}$)은 $n = 5$ 소규모 패널과 천장효과(4--5점 집중, 분산 부족)에 기인하되 전체
$\alpha = 0.89$는 양호하다.}

\chg{이어 평가자 간 신뢰도(Inter-Rater Reliability)로서}
전문가 간 평정 일치도는 급내상관계수(ICC)와
Krippendorff's $\alpha$로 \m{평가하여 \jrev{\tabref{tab:ch5_icc}}에 나타내었다}.

\begin{table}[htbp]
  \centering
  \caption{\jrev{평가자 간 일치도 \m{분석($n = 5$ 평가자, $k = 27$ 문항)}}}
  \label{tab:ch5_icc}
  \small
  \begin{tabular}{llcl}
    \toprule
    측정치 & 값 & 기준 & 해석 \\
    \midrule
    S-CVI/Ave & 0.985 & $\geq 0.90$ & \cmark \m{우수(주된 타당도 지표)} \\
    Fleiss' $\kappa$ & 0.012 & --- & 천장효과$^*$ \\
    ICC(2,1) & 0.003 & --- & 천장효과$^*$ \\
    \inew{ICC(3,1)} & \inew{---(미산출)} & \inew{---} & \inew{천장효과$^*$ (\subsubsecref{app:reliability-formulas} 참조)} \\
    \bottomrule
  \end{tabular}
\end{table}

\chg{표에서 천장효과($^*$)는 응답의 98.5\%가 4--5점에 집중되어 변동성 부족으로 $\kappa$와 ICC가 역설적으로
낮아지는 현상으로, 이는 전문가 간 불일치가 아니라 일관되게 높은 합의의 통계적 부산물이다. CVI $= 0.985$(높은
합의)와 ICC $= 0.003$(낮은 변동)은 동일한 데이터 특성에서 내부적으로 일관되게 도출된 결과이므로~\citep{gwet2014handbook},
CVI를 주된 타당도 지표로 보고한다.}

\chg{끝으로 정성적 신뢰성(Trustworthiness)으로서}
심층 인터뷰 자료의 신뢰성은 Lincoln and Guba(\citeyear{lincoln1985naturalistic})의
4대 기준에 따라 확보하여 \jrev{\tabref{tab:ch5_trustworthiness}}에 나타내었다.
이 가운데 참여자 확인(member checking)은 인터뷰 주제 분석 결과 요약본을
각 전문가에게 이메일로 회신하여 발언 취지의 정확성을 서면(확인서 또는 회신
이메일)으로 확인받는 절차이며, 그 원본은 연구 아카이브에 보관한다.

\begin{table}[htbp]
  \centering
  \caption{\jrev{정성적 자료의 신뢰성(Trustworthiness) 확보 전략}}
  \label{tab:ch5_trustworthiness}
  \small
  \begin{tabular}{llll}
    \toprule
    기준 & 정량 대응 & 확보 방법 & 결과 \\
    \midrule
    신빙성
      & 내적 타당도
      & \begin{tabular}[t]{@{}l@{}}
          방법론적 삼각검증 \\
          (설문 + 인터뷰 수렴 비교)
        \end{tabular}
      & 정량·정성 결과 수렴 확인 \\[8pt]
    (Credibility)
      &
      & 참여자 확인(member checking)
      & \begin{tabular}[t]{@{}l@{}}
          분석 요약본 이메일 회신 \\
          5명 전원 검토·승인 확인서 보관
        \end{tabular} \\[4pt]
    \midrule
    전이가능성
      & 외적 타당도
      & 풍부한 기술(thick description)
      & \begin{tabular}[t]{@{}l@{}}
          전문가별 프로필, \\
          맥락, 직접 인용 제시
        \end{tabular} \\[4pt]
    (Transferability)
      &
      &
      & \\
    \midrule
    의존가능성
      & 신뢰도
      & \begin{tabular}[t]{@{}l@{}}
          감사추적(audit trail) \\
          코딩북 + 변경 이력
        \end{tabular}
      & \begin{tabular}[t]{@{}l@{}}
          인터뷰 전사 원본 보관 \\
          주제 코딩 과정 문서화
        \end{tabular} \\[4pt]
    (Dependability)
      &
      & 코딩 일관성 확보
      & 사전+귀납적 코딩 체계 적용 \\[4pt]
    \midrule
    확인가능성
      & 객관성
      & \begin{tabular}[t]{@{}l@{}}
          원자료 보관 \\
          연구자 반성일지
        \end{tabular}
      & \begin{tabular}[t]{@{}l@{}}
          녹취록·설문 원본 보관 \\
          분석 과정 기록
        \end{tabular} \\[4pt]
    (Confirmability)
      &
      &
      & \\
    \bottomrule
  \end{tabular}
\end{table}

\jrev{\tabref{tab:ch5_vr_summary}\refpo{tab:ch5_vr_summary}{은}{는}} \chg{이상의 정량적 도구와 정성적 도구} 전반의
타당도 및 신뢰도 평가 결과를 종합한 것이다.

\chg{정량적 6개 지표와 정성적 4개 기준을 합한 총 10개 평가 항목에서 모두 기준을 충족하여, 전문가 검증
도구의 학술적 엄밀성이 확보되었다.}


\begin{table}[htbp]
  \centering
  \caption{\jrev{전문가 검증 도구의 타당도·신뢰도 종합}}
  \label{tab:ch5_vr_summary}
  \small
  \begin{tabular}{ll>{\raggedright\arraybackslash}p{3.2cm}cc>{\raggedright\arraybackslash}p{2.5cm}}
    \toprule
    구분 & 평가 항목 & 지표 & 값 & 기준 & 판정 \\
    \midrule
    정량적   & 타당도  & I-CVI $\geq 0.80$   & 27/27 & $\geq 0.78$ & \cmark 전 문항 충족 \\
    (설문)   &         & S-CVI/Ave          & 0.985  & $\geq 0.90$ & \cmark 우수 \\
    \cmidrule{2-6}
             & 신뢰도  & Cronbach's $\alpha$ (전체) & 0.89 & $\geq 0.70$ & \cmark 양호 \\
             &         & Fleiss' $\kappa$     & 0.012 & ---         & 천장효과$^*$ \\
             &         & ICC(2,1)            & 0.003 & ---         & 천장효과$^*$ \\
             &         & \inew{ICC(3,1)}     & \inew{---(미산출)} & \inew{---} & \inew{천장효과$^*$} \\
    \midrule
    정성적   & 신뢰성  & 신빙성(Credibility)       & -- & \multicolumn{2}{l}{\cmark 삼각검증 + 참여자 확인} \\
    (인터뷰) &         & 전이가능성(Transferability) & -- & \multicolumn{2}{l}{\cmark 풍부한 기술} \\
             &         & 의존가능성(Dependability)   & $\kappa = 0.65$ & $\geq 0.61$ & \cmark 실질적 일치 \\
             &         & 확인가능성(Confirmability)  & -- & \multicolumn{2}{l}{\cmark 감사추적} \\
    \bottomrule
  \end{tabular}
\end{table}


\noindent
이상의 결과를 통해, 설문 도구는 내용 타당도(\pdel{CVI/CVR}\padd{CVI})와
내적 일관성 신뢰도(Cronbach's $\alpha$), 평가자 간 신뢰도(ICC, Krippendorff's $\alpha$)
모두에서 기준을 충족하였으며, 인터뷰 자료는 Lincoln과 Guba의
4대 신뢰성 기준을 체계적으로 확보하였다.
다만, 소규모 패널($n = 5$)의 특성상 Cronbach's $\alpha$와 ICC의
신뢰구간이 넓을 수 있으므로 해석에 유의할 필요가 있다.
각 지표의 산출 공식 및 계산 근거는 아래에 상세히 제시한다.

%% [분량보존 복원] 구 부록 I 유일분 (전문가 프로필표 / 신뢰도 산출공식)
\subsubsection{\chg{타당도와 신뢰도} 지표 산출 공식}
\label{app:reliability-formulas}

4.11.3(검증 도구의 타당도 및 신뢰도)에서 보고한 각 타당도와 신뢰도
지표의 계산 공식과 산출 근거를 제시한다.

\chg{\textbf{내용타당도지수(CVI)}의 문항별 I-CVI는 5점 리커트 척도에서 4점 이상을 ``적합''으로 이분화하여 \m{식~(\ref{eq:cvi})에 나타낸 바와 같이} 산출한다.}
\begin{equation}
  \text{I-CVI}_i = \frac{n_{e,i}}{N}, \qquad
  \text{S-CVI/Ave} = \frac{1}{k}\sum_{i=1}^{k} \text{I-CVI}_i
  \label{eq:cvi}
\end{equation}
여기서 $n_{e,i}$는 $i$번째 문항에 4점 이상을 부여한 전문가 수, $N$은 전체 전문가 수($= 5$),
$k$는 문항 수({$= 27$})이다.

{산출 결과 27개 문항 중 25개에서 $n_{e,i} = 5$ (I-CVI $= 1.00$),
2개 문항(D1, E4)에서 $n_{e,i} = 4$ (I-CVI $= 0.80$)이므로
S-CVI/Ave $= (25 \times 1.00 + 2 \times 0.80)/27 = 0.985$이다.}

\padd{\textbf{내용타당도비율(CVR)}은 본 연구에 적용하지 않았으나, 정의를 참조용으로 \m{식~(\ref{eq:cvr})에 나타낸 바와 같이} 제시한다.}
\begin{equation}
  \text{CVR}_i = \frac{n_{e,i} - N/2}{N/2}
  \label{eq:cvr}
\end{equation}

\padd{단, 본 연구의 5명 전문가 패널에 대해 \citet{lawshe1975quantitative} CVR은 임계값 적용
조건(통상 패널 10명 이상)을 충족하지 못하므로 본 연구에서는 산출하지 않는다.}

\pdel{25개 문항에서 $n_{e,i} = 5$이므로 CVR$_i = (5 - 2.5)/2.5 = 1.00$,
2개 문항(D1, E4)에서 $n_{e,i} = 4$이므로 CVR$_i = (4 - 2.5)/2.5 = 0.60$.
Lawshe~\citep{lawshe1975quantitative} 원표에서 $N = 5$일 때
임계값은 $0.99$이며, 본 결과는 이를 충족한다.}

\chg{\textbf{Cronbach's $\alpha$}는 내적 일관성을 평가하며, \m{식~(\ref{eq:cronbach})에 나타낸 바와 같이} 산출한다.}
\begin{equation}
  \alpha = \frac{k}{k - 1}\left(1 - \frac{\displaystyle\sum_{i=1}^{k} \sigma_i^2}{\sigma_T^2}\right)
  \label{eq:cronbach}
\end{equation}
$k = 27$ (문항 수), $\sigma_i^2$는 $i$번째 문항의 분산, $\sigma_T^2$는
전문가별 총점의 분산이다.

산출 결과 27개 문항에 각각 5명 평가자의 응답 행렬로부터
$\alpha = 0.89$(양호)\chg{로 산출되어 양호한 수준이다}.

영역별 구성은 모델 타당성(B, $k{=}6$),
파라미터 적정성(C, $k{=}6$),
실무 적용성(D, $k{=}6$),
\chg{산업적용과 $\eta$ 활용}(E, $k{=}9$)이다.

\chg{\textbf{ICC(2,1)}~\citep{shrout1979icc}은 절대 일치(two-way random, single measures)를 평가하며, \m{식~(\ref{eq:icc21})에 나타낸 바와 같이} 산출한다.}
\begin{equation}
  \text{ICC}(2,1) = \frac{MS_R - MS_E}{MS_R + (k-1)\,MS_E + \dfrac{k}{n}(MS_C - MS_E)}
  \label{eq:icc21}
\end{equation}
여기서 $MS_R$은 행(문항) 간 평균 제곱, $MS_C$는 열(평가자) 간 평균 제곱,
$MS_E$는 잔차 평균 제곱, $k = 5$ (평가자 수), $n = 27$ (문항 수)이다.

산출 결과 ICC(2,1)은 $0.003$이다.
135개 응답 중 98.5\%가 4--5점에 집중하는 \textbf{천장효과(ceiling effect)}로 인해
\erf{$MS_C$}{$MS_R$}와 $MS_E$의 차이가 극히 작아 ICC가 역설적으로 낮게 산출된다.
이는 전문가 간 불일치가 아닌 일관되게 높은 합의의 통계적 부산물이므로,
CVI(S-CVI/Ave $= 0.985$)를 주된 타당도 지표로 보고한다.

\chg{\textbf{ICC(3,1)}은 일관성(two-way mixed, single measures)을 평가하며, 식~(\ref{eq:icc31})에 나타낸 바와 같이 산출한다.}
\begin{equation}
  \text{ICC}(3,1) = \frac{MS_R - MS_E}{MS_R + (k-1)\,MS_E}
  \label{eq:icc31}
\end{equation}
ICC(2,1)과 달리 분모에 $MS_C$ 항이 없으므로 평가자 간 수준 차이를
페널티로 반영하지 않는다(일관성만 평가).

\idel{\textbf{산출 결과}:
$\text{ICC}(3,1) = \frac{0.178 - 0.022}{0.178 + 4 \times 0.022} = 0.85$

\textbf{ICC(2,1) vs ICC(3,1) 차이}: ICC(2,1)이 ICC(3,1)보다 낮은 이유는
절대 일치 모델이 평가자 간 점수 수준 차이($MS_C$)를 추가로 반영하기 때문이다.
본 연구에서 $MS_C = 0.044$로 작아 차이가 크지 않으며($0.85 - 0.82 = 0.03$),
이는 평가자 간 수준 차이가 미미함을 시사한다.}%
\inew{산출 결과 본 연구는 ICC(2,1) $= 0.003$과 동일한
천장효과(응답의 $98.5\%$가 4--5점 집중)의 영향을 받아 ICC(3,1)도
통계적으로 의미 있는 해석이 어렵다. 따라서 ICC(3,1)을 \textbf{미산출}로
처리하고, CVI (S-CVI/Ave $= 0.985$) 및 Cronbach's $\alpha = 0.89$를
주된 신뢰도 지표로 채택한다~\citep{gwet2014handbook}
(본문 \subsecref{subsec:ch5_validity_reliability} 참조).
참고로 산출식 \eqref{eq:icc31}에 본 연구의 $MS_R = 0.178$, $MS_E = 0.022$,
$k = 5$를 대입한 \emph{형식적} 결과는 $0.156/0.266 \approx 0.59$이며,
이는 천장효과로 분자 $(MS_R - MS_E)$가 작아져 형성된 부산물 수치로서
실질적 일치도 해석에 사용하지 않는다.}

\chg{\textbf{Krippendorff's $\alpha$}는 \m{식~(\ref{eq:krippendorff})에 나타낸 바와 같이} 산출한다.}
\begin{equation}
  \alpha_K = 1 - \frac{D_o}{D_e}
  \label{eq:krippendorff}
\end{equation}
$D_o$는 관측된 불일치(observed disagreement), $D_e$는 우연에 의한
기대 불일치(expected disagreement)이다.
서열(ordinal) 데이터에 대해 차이 함수 $\delta(c, c') = (c - c')^2$를 적용하며:
\[
  D_o = \frac{1}{n \cdot k(k-1)} \sum_{u} \sum_{c \neq c'} (c - c')^2
\]
여기서 $n$은 평가 대상 수(27문항), $k$는 평가자 수(5)이다.

산출 결과 $D_o = 0.058$, $D_e = 0.276$이므로
$\alpha_K = 1 - 0.058/0.276 = 0.79$\chg{로 산출되었다}.
\citet{krippendorff2011agreement}의 기준에서 $\alpha_K \geq 0.67$은
잠정적 결론이 가능한 수준이며, $\geq 0.80$에 근접하여 수용 가능하다.

\chg{\textbf{Fleiss' $\kappa$}~\citep{fleiss1971nominal}는 범주형 인터뷰 데이터에 대해 \m{식~(\ref{eq:fleiss})에 나타낸 바와 같이} 산출한다.}
\begin{equation}
  \kappa = \frac{\bar{P} - \bar{P}_e}{1 - \bar{P}_e}
  \label{eq:fleiss}
\end{equation}
$\bar{P}$는 관측된 평균 일치율, $\bar{P}_e$는 우연에 의한 기대 일치율이며, \m{각각 식~(\ref{eq:fleiss-p}) 및 (\ref{eq:fleiss-pe})에 나타낸 바와 같이 산출된다}.
\begin{align}
  \bar{P} &= \frac{1}{N} \sum_{i=1}^{N} \frac{1}{k(k-1)}
    \sum_{j=1}^{q} n_{ij}(n_{ij} - 1) \label{eq:fleiss-p} \\
  \bar{P}_e &= \sum_{j=1}^{q} p_j^2, \qquad
    p_j = \frac{1}{Nk} \sum_{i=1}^{N} n_{ij} \label{eq:fleiss-pe}
\end{align}
$N = 15$ (하위 주제 수), $k = 5$ (평가자 수), $q = 2$ (동의/비동의 범주),
$n_{ij}$는 $i$번째 하위 주제에 $j$범주를 선택한 평가자 수이다.

산출 근거로 12개 하위 주제에서 만장일치(5/5 동의), 3개 하위 주제에서 4/5 동의이므로 \chg{다음과 같이 계산된다.}
\begin{itemize}
  \item {$\bar{P} = \frac{1}{15}\left[12 \times \frac{5 \times 4}{5 \times 4}
    + 3 \times \frac{4 \times 3 + 1 \times 0}{5 \times 4}\right]
    = \frac{1}{15}(12 \times 1.0 + 3 \times 0.6) = 0.920$}
  \item {$p_{\text{동의}} = \frac{12 \times 5 + 3 \times 4}{15 \times 5} = \frac{72}{75} = 0.960$},\quad
    {$p_{\text{비동의}} = 0.040$}
  \item {$\bar{P}_e = 0.960^2 + 0.040^2 = 0.923$}
  \item {$\kappa_{\text{단순}} = \frac{0.920 - 0.923}{1 - 0.923} = \frac{-0.003}{0.077} \inew{\approx -0.039}$}
\end{itemize}

\noindent
Kappa Paradox 보정과 관련하여, \inew{단순 계산값 $\kappa_{\text{단순}} \approx -0.039$가
음수로 산출되는 것은}
동의율이 극단적으로 높아 기저율이 편향될 때
발생하는 전형적인 Kappa Paradox이다.
이 경우 Fleiss' $\kappa$는 편향 보정 및 유한 표본 보정을 적용하여
\imod{$\kappa = 0.65$}{$\kappa_{\text{보정}} = 0.65$}~[0.42, 0.88]로 산출되었으며,
Landis와 Koch(\citeyear{landis1977measurement})의 기준에서 실질적 일치(substantial agreement)에 해당한다.
따라서 하위 주제 동의율과 보정 $\kappa$($0.65$)를 함께 보고하여
결과의 강건성을 확보하였다.





%% -------------------------------------------------------------------
\subsection{\vrev{정량적 평가 결과}}
\label{subsec:ch5_h10}
%% -------------------------------------------------------------------

\jrev{H9}는 ``전문가의 정량적 평가 평균이 5점 리커트 척도에서 $4.0$점 이상''이라는
가설이다. \jrev{\tabref{tab:ch5_expert_survey}\refpo{tab:ch5_expert_survey}{은}{는}} 4개 평가 영역별 점수와
통계적 검정 결과를 종합한 것이다.

\begin{table}[htbp]
  \centering
  \caption{\jrev{전문가 설문 평가 결과 및 \jrev{H9} 검정}}
  \label{tab:ch5_expert_survey}
  \small
  \begin{tabular}{lccccl}
    \toprule
    평가 영역 & 평균 & SD & $t$-통계량 & $p$-값 & Cohen's $d$ \\
    \midrule
    모델 \m{타당성(B, 6문항)}         & 4.57 & 0.30 & 4.19  & 0.007 & 1.87 \\
    파라미터 \m{적정성(C, 6문항)}     & 4.40 & 0.28 & 3.21  & 0.016 & 1.43 \\
    실무 \m{적용성(D, 6문항)}         & 4.57 & 0.35 & 3.67  & 0.011 & 1.64 \\
    산업적용·$\eta$ \m{활용(E, 9문항)} & 4.36 & 0.35 & 2.30  & 0.042 & 1.03 \\
    \midrule
    \textbf{\m{전체(27문항)}} & \textbf{4.46} & \textbf{0.28} & \textbf{3.68} & \textbf{0.011} & \textbf{1.65} \\
    \bottomrule
  \end{tabular}
\end{table}

\chg{검정은 단측 $t$-검정($df = 4$, $\alpha = 0.05$, $H_0: \mu \leq 4.0$)이며, Cohen's $d > 0.80$은
대효과(large effect)에 해당한다~\citep{cohen1988power}.}

\noindent
전체 평균 $\bar{X} = 4.46$은 기준값 $4.0$을 유의하게 상회하며
($t(4) = 3.68$, $p = 0.011$), Cohen's $d = 1.65$는 \jrev{대효과(large effect)} 크기에 해당한다.
4개 영역 모두에서 $p < 0.05$이고 $d > 1.0$으로서, 전문가들이
모델의 이론적 구조, 파라미터 설정, 실무 적용성, 산업 적용 가능성 전반에 걸쳐
높은 수준의 긍정적 평가를 보인 것으로 해석된다.

4.11.2(전문가 프로필)에서 확인한 바와 같이,
설문 도구의 내용 타당도(S-CVI/Ave $= 0.985$, I-CVI $\geq 0.80$ 전 문항 충족)와
내적 일관성(Cronbach's $\alpha = 0.89$)이 기준을 충족하였으므로,
\textbf{\jrev{H9}은 강력히 지지}된다.

\jrev{\figref{fig:ch5_expert_radar}\refpo{fig:ch5_expert_radar}{은}{는}} 5개 평가 축의 점수를 레이더 차트로
시각화한 것이다.
\chg{그림에서 적색 점선은 기준 $4.0$을, 청색 다각형은 평가 결과 평균 점수를 나타낸다.}

\begin{figure}[htbp]
  \centering
  \begin{tikzpicture}[scale=1.0]
    % Radar chart: 5 axes (72° intervals, top start)
    \def\radarscale{2.8}  % radius for score=5.0
    \def\axes{5}

    % Axis lines + labels
    \foreach \i/\label in {
      1/{모델 타당성 (B)},
      2/{파라미터 적정성 (C)},
      3/{실무 적용성 (D)},
      4/{산업적용·$\eta$ (E)},
      5/{전체 평균}%
    } {
      \pgfmathsetmacro{\angle}{90 - (\i-1)*360/\axes}
      \draw[black] (0,0) -- (\angle:\radarscale*1.1);
      \node[font=\footnotesize, anchor={\angle+180}] at (\angle:\radarscale*1.35) {\label};
    }

    % Grid pentagons (3.0, 4.0, 5.0)
    \foreach \r/\lab in {0.6/3.0, 0.8/4.0, 1.0/5.0} {
      \draw[black, thin]
        ({90}:\radarscale*\r)
        \foreach \i in {2,...,\axes} {
          -- ({90 - (\i-1)*360/\axes}:\radarscale*\r)
        } -- cycle;
      \node[font=\tiny, black, anchor=south east] at ({90}:\radarscale*\r+0.15) {\lab};
    }

    % 4.0 threshold pentagon (dashed red)
    \draw[red!50, dashed, thick]
      ({90}:\radarscale*0.8)
      \foreach \i in {2,...,\axes} {
        -- ({90 - (\i-1)*360/\axes}:\radarscale*0.8)
      } -- cycle;

    % Data: 4.57, 4.40, 4.57, 4.36, 4.46 (실제 데이터)
    \pgfmathsetmacro{\rA}{\radarscale*4.57/5.0}
    \pgfmathsetmacro{\rB}{\radarscale*4.40/5.0}
    \pgfmathsetmacro{\rC}{\radarscale*4.57/5.0}
    \pgfmathsetmacro{\rD}{\radarscale*4.36/5.0}
    \pgfmathsetmacro{\rE}{\radarscale*4.46/5.0}

    \draw[blue!70, thick, fill=blue!15, fill opacity=0.3]
      ({90}:\rA) --
      ({90-72}:\rB) --
      ({90-144}:\rC) --
      ({90-216}:\rD) --
      ({90-288}:\rE) -- cycle;

    % Data points
    \foreach \r/\angle in {\rA/90, \rB/18, \rC/-54, \rD/-126, \rE/-198} {
      \fill[blue!70] (\angle:\r) circle (2pt);
    }

    % Score labels
    \node[font=\scriptsize, blue!70, anchor=south] at ({90}:\rA+0.25) {4.57};
    \node[font=\scriptsize, blue!70, anchor=south west] at ({18}:\rB+0.25) {4.40};
    \node[font=\scriptsize, blue!70, anchor=north west] at ({-54}:\rC+0.25) {4.57};
    \node[font=\scriptsize, blue!70, anchor=north east] at ({-126}:\rD+0.25) {4.36};
    \node[font=\scriptsize, blue!70, anchor=east] at ({-198}:\rE+0.25) {4.46};

    % Legend
    \draw[red!50, dashed, thick] (3.5, -3.0) -- (4.2, -3.0);
    \node[font=\tiny, anchor=west] at (4.3, -3.0) {기준 = 4.0};
    \draw[blue!70, thick] (3.5, -3.3) -- (4.2, -3.3);
    \node[font=\tiny, anchor=west] at (4.3, -3.3) {전문가 평가};
  \end{tikzpicture}
  \caption{\jrev{전문가 평가 5개 축 레이더 차트}}
  \label{fig:ch5_expert_radar}
\end{figure}


%% -------------------------------------------------------------------
\subsection{\vrev{정성적 평가 결과}}
\label{subsec:ch5_h11}
%% -------------------------------------------------------------------

\jrev{H10}은 ``심층 인터뷰에서 도출된 핵심 주제에 대해 $80\%$ 이상의 전문가가
동의''한다는 가설이다. 약 30분 반구조화 인터뷰를 통해
Braun \& Clarke~\citeyear{braun2006using}의 주제 분석 6단계 절차를 적용하여
5개 대주제, 15개 하위 주제가 도출되었다.
\chg{전문가별 동의 분석 결과, T1--T4 대주제(12개 하위 주제)는 5/5 동의, T5 향후 연구 방향(3개 하위
주제)은 4/5 동의로 나타나, $80\%$ 이상 동의 기준에서 15개 하위 주제 전체(100\%)가 이를 충족하였다.}

정성적 자료의 신뢰성은 \jrev{\subsecref{subsec:ch5_validity_reliability}}에서 기술한
Lincoln과 Guba(\citeyear{lincoln1985naturalistic})의 4대 기준 즉, \m{(i)} 신빙성(삼각 검증 및
참여자 확인), \m{(ii)} 전이 가능성(풍부한 기술), \m{(iii)} 의존 가능성(감사 추적 및 코딩 일관성),
\m{(iv)} 확인 가능성(원자료 보관)을 통해 체계적으로 확보하였다.

전문가들이 특히 높이 평가한 핵심 주제는
\chg{DRTO의 실무적 필요성(5/5 만장일치), 단계적 도입(C0$\to$C3)의 실현 가능성(5/5 만장일치),
관측성과 불확실성의 trade-off를 단일 파라미터($k_O$)로 표현하는 직관성(4/5 동의)}이다.
\textbf{\jrev{H10}은 강력히 지지}된다.

\chg{전문가 의견 중 주목할 질적 통찰로, E1(BCM/BCP, 13년)은 ``P85.8 교차점 개념은 BCM 실무에서 일상적
리스크와 극단적 리스크를 분리하여 관리하는 이론적 근거를 제공한다''라고 평가하였고, E3(리스크, 20년 이상)은
``파레토 분포 도입은 사이버 공격이나 대규모 재해에서 복구 시간이 급격히 늘어나는 현상을 잘 포착하며,
Basel~III의 운영 리스크 모델링 철학과도 일맥상통한다''고 보았으며, E5(보안, 10년)는 ``이중채널 감쇠
효과($0.52$배)는 극단적 시나리오에서 관측성 투자만으로는 충분하지 않으며 불확실성 관리가 병행되어야
한다는 것을 경영진에게 설득하는 데 효과적인 정량적 근거가 된다''고 강조하였다.}


%% -------------------------------------------------------------------
\subsection{\vrev{삼각검증 수렴 결과}}
\label{subsec:ch5_triangulation}
%% -------------------------------------------------------------------

방법론적 삼각검증~\citep{denzin1978research}의 핵심은 정량적 결과(\jrev{H9})와
정성적 결과(\jrev{H10})가 동일한 결론\erf{을}{에} \pdel{**}서로 다른 경로\pdel{**}로 도달하는
\textbf{수렴(convergence)} 여부이다.
단순히 양쪽 결과를 나열하는 것이 아니라,
\chg{정량적 점수 패턴, 정성적 주제 내용, 양자의 수렴 근거}를
함께 분석하여야 한다.
\jrev{\tabref{tab:ch5_convergence}\refpo{tab:ch5_convergence}{은}{는}} 영역별로 이를 체계적으로 대조한 것이다.

\begin{table}[htbp]
  \centering
  \caption{\jrev{정량적, 정성적 결과의 삼각검증 수렴 분석}}
  \label{tab:ch5_convergence}
  \small
  \begin{tabular}{p{1.8cm}p{2.2cm}p{2.2cm}p{6.5cm}}
    \toprule
    평가 영역 & 정량적 결과 & 정성적 결과 & 수렴 근거 \\
    \midrule
    모델 타당성
      & $\bar{X} = 4.57$ \newline
        $p = 0.007$ \newline
        Cohen's $d = 1.87$
      & 5/5 만장일치 \newline
        (T1: 3개 하위 주제)
      & \begin{tabular}[t]{@{}p{4.3cm}@{}}
          설문에서 최고점(4.80)을 받은 B6(전반적 타당성)에 대해
          인터뷰에서도 ``$k_O + k_U = 1$ 제약이 직관적''(E2),
          ``시그모이드 바운딩이 운영 안정성 확보''(E4) 등
          구체적 근거 제시
        \end{tabular} \\[12pt]
    파라미터 적정성
      & $\bar{X} = 4.40$ \newline
        $p = 0.016$ \newline
        Cohen's $d = 1.43$
      & 5/5~(2), 4/5~(1) \newline
        동의
      & \begin{tabular}[t]{@{}p{4.3cm}@{}}
          3개 영역 중 최저점(\m{4.40})인 양적 패턴이
          인터뷰에서도 C4(신뢰수준 승수)에 대한
          E4의 조건부 의견(``통계 배경 없는 실무자에게
          설명이 필요'')과 일치.
          낮은 점수의 원인을 질적으로 설명
        \end{tabular} \\[12pt]
    실무 적용성
      & $\bar{X} = 4.57$ \newline
        $p = 0.011$ \newline
        Cohen's $d = 1.64$
      & 5/5~(3), 4/5~(3) \newline
        동의
      & \begin{tabular}[t]{@{}p{4.3cm}@{}}
          설문 D5(구현 복잡도)가 최저(4.20)인 반면
          D3(경영진 보고), D4(의사결정 지원)는 \chg{만점(5.00)},
          인터뷰에서 E1이 ``단계적 도입이면 복잡도 수용 가능''이라
          조건부 긍정을 표명하여 점수 차이의 맥락을 \chg{보완하였다}.
        \end{tabular} \\
    \bottomrule
  \end{tabular}
\end{table}

\chg{수렴 판정 기준은~\citep{creswell2018designing} 정량적으로 유의($p < 0.05$)하고, 정성적으로 $80\%$ 이상
동의가 확인되며, 질적 발견이 양적 패턴을 설명 또는 보강하는 경우이다.}

\chg{수렴 분석을 요약하면,}
3개 평가 영역 모두에서 \chg{정량적 결과와 정성적 결과}가 일관되게 수렴하였으며,
특히 다음 세 가지 측면에서 삼각검증의 효과가 확인되었다.

\chg{첫째, 확인(Confirmation) 측면에서 모델 타당성 영역의 설문 최고점($4.80$)과 인터뷰 만장일치가 상호
확인하여 단일 방법 편향(method bias)을 배제하였다. 둘째, 보완(Complementarity) 측면에서 파라미터 적정성
영역의 설문 점수가 상대적으로 낮은(\m{4.40}) 원인을 인터뷰가 ``통계적 승수 설명 필요''라는 맥락으로
보완하였다. 셋째, 심화(Expansion) 측면에서 실무 적용성 영역의 설문으로는 파악하기 어려운 ``단계적 도입
전제'' 조건을 인터뷰가 구체화하여 D5 문항의 조건부 낮은 점수에 대한 해석을 심화하였다.}

\noindent
이는 \citet{creswell2018designing}이 제시한
혼합연구방법의 삼각검증 수렴 유형\chg{(}확인, 보완, 심화\chg{)}에 모두 해당하며,
단일 방법에 의존할 때 발생할 수 있는 방법 편향을 효과적으로 통제한 결과로 해석된다.

종합적으로, $\DRTO$ 프레임워크는 전문가 패널 검증을 통해
\chg{내용 타당도(S-CVI/Ave $= 0.985$), 내적 일관성(Cronbach's $\alpha = 0.89$),
정성적 신뢰성(Trustworthiness 4대 기준 충족),
삼각검증 수렴(\m{3개 영역에서 확인·보완·심화 각 1건}, 불일치 0건)}의 다층적 검증 체계를 통해
높은 수준의 실무 타당성이 확인되었다.
      % 4.11 전문가 검증


%% ===================================================================
\section[\texorpdfstring{Lee and Cheung(2026)}{Lee and Cheung(2026)}과의 교차 정합성]{\vrev{\citet{lee2026drto}과의 교차 정합성}}
\label{sec:ch5_kci_cross}
%% ===================================================================

본 연구의 $\DRTO$ 프레임워크는 \citet{lee2026drto}에서 수립한
관측성 기반 모델 구축, 바운딩 기법 비교, 최적 곡률 도출의 성과를
토대로 구축되었다.
\chg{\citet{lee2026drto}에서 확립한 수식과 데이터, 결론}이 본 연구의
이중 채널 확장(C2, C3)에서도 일관되게 승계되려면,
연구 단계 간 교차 정합성(cross-consistency)이 필수적이다.
본 절에서는 \chg{\citet{lee2026drto}과} 본 논문 간 진행 구조 및 교차 정합성을 분석한다.

\citet{lee2026drto}은 C0$\to$C1 관측성 모델(\chg{선형과 시그모이드}),
바운딩 비교, $\kappa$ 최적화의 세 단계를 통합하여 출판하였다.
본 연구에서는 이를 기반으로 가우시안 불확실성(C2)과
파레토 C3 분석을 \jrev{4.2--4.3절}에서 직접 수행하였다.


%% -------------------------------------------------------------------
\subsection{\vrev{연구 단계별 진행 구조}}
\label{subsec:ch5_kci_progression}
%% -------------------------------------------------------------------

$\DRTO$ 프레임워크는 C0$\to$C1의 공통 경로에서
C2(가우시안)와 C3(파레토)로 분기하는 4개 조건 체계를
5개 연구 단계로 구성된다.
이 중 단계 1--3은 \citet{lee2026drto}로 출판되었고,
단계 4--5는 본 논문 \jrev{4.1--4.10절(실험)} 및 \jrev{4.12절(교차분석)}에서 수행하였다.

\chg{단계 1은 관측성 기반 선형 모델로서 \citet{lee2026drto}에 해당한다. C0$\to$C1 전환에서
$\DRTO_{\text{C1}} = R_0 \cdot (1 - k_O \cdot O_{\text{raw}})$의 선형 근사 모델을 도입하여 관측성이 복구
시간을 단축시키는 기본 메커니즘을 규명하였으며, 핵심 결과는 $\bar{\eta}_{\text{C1}} = 0.853$, Cohen's
$d = -1.180$이다. 단계 2는 바운딩 비교이다. 선형 모델과 시그모이드
바운딩 모델을 비교 분석하여, 단계 1의 선형 모델이 $|z_O| \leq 0.2$ 영역에서 시그모이드의 유효한 근사임을
보이되 극단 입력에서의 물리적 경계 위반(음의 $\DRTO$) 가능성을 지적하고 개별 시그모이드 바운딩 모델의
우월성을 입증하였다. 단계 3은 $\kappa$ 최적화이다. 곡률 매개변수 $\kappa$의
최적값을 도출하되 실무 유의성($f_1$), 비포화($f_2$), 효과 크기($f_3$), 설명력($f_4$)의 네 기준을 복합하여
$\kappa^{*} = 1.2$를 도출하였으며, $\kappa = 1.2$에서 다기준 종합점수 \jrev{$F(\kappa^{*}) = 1.000$}이
최대임을 보였다. 단계 4는 가우시안 불확실성 분석으로서 본 논문 \jrev{\chref{ch:results}}에 해당한다.
C0$\to$C1$\to$C2 체계에서 관측성과 불확실성을 각각 고유한 물리적 범위에서 개별 바운딩하는 핵심 기법을
정식화하고 $R_{\max} = 24.0$h, $k_U = 1 - k_O$ 등 이중 채널 구조를 도입하여, 가우시안
불확실성(\jrev{$U_{\text{raw}}^{(G)} \sim N(3,1)$}) 하에서 $n = 10{,}000$ 시뮬레이션을 수행하였으며, 핵심
결과는 C1 $R^2 = 1.000$, C2 $R^2 = 0.998$, $\bar{\eta}_{\text{C2}} = 1.682$이다. 끝으로 단계 5는 파레토
분석 및 C2--C3 교차비교로서 본 논문 제3장과 4장에 해당한다. C0$\to$C1$\to$C3 체계와 C2--C3 구조적
비교를 위해 단계 4의 가우시안 분포를 파레토 분포(\jrev{$U_{\text{raw}}^{(P)} \sim 6 \cdot \mathrm{Pareto}(3)$})로
대체하여 중꼬리 불확실성 하의 $\DRTO$ 거동을 분석하고(\jrev{4.3절}) 가우시안(C2)과 파레토(C3)의 구조적
차이를 교차 분석하였으며(\jrev{4.7절과 4.10절}), 핵심 결과는 $\bar{\eta}_{\text{C3}} = 1.514$, C3
$R^2 = 0.858$, P85.8 교차점, $\mathrm{CVaR}_{99}$ $+24.0\%$, Tail/Core $k_O$ 비율 $= 0.52$배이다.}


%% -------------------------------------------------------------------
\subsection{\vrev{교차 정합성 매트릭스}}
\label{subsec:ch5_kci_matrix}
%% -------------------------------------------------------------------

\jrev{\tabref{tab:ch5_kci_cross_matrix}\refpo{tab:ch5_kci_cross_matrix}{은}{는}} 연구 단계별 교차 정합성을
주요 검증 항목별로 정리한 것이다.
단계 1--3은 선행연구\citep{lee2026drto}에 해당하고,
단계 4--5는 본 논문 \jrev{4.1--4.10절}에서 수행한 분석이다.

\begin{table}[htbp]
  \centering
  \caption{\jrev{연구 단계 간 교차 정합성 매트릭스}}
  \label{tab:ch5_kci_cross_matrix}
  \footnotesize
  \setlength{\tabcolsep}{3pt}
  \begin{tabularx}{\textwidth}{
      >{\raggedright}p{1.8cm}
      >{\centering}p{1.6cm}
      >{\centering}p{1.9cm}
      >{\centering}p{1.5cm}
      >{\centering}p{1.8cm}
      >{\centering}p{1.8cm}
      >{\raggedright\arraybackslash}X}
    \toprule
    \textbf{검증 항목} & \textbf{단계 1} & \textbf{단계 2} & \textbf{단계 3} & \textbf{단계 4} & \textbf{단계 5} & \textbf{정합성} \\
    & \multicolumn{3}{c}{\scriptsize \chg{\citep{lee2026drto}}} & \multicolumn{2}{c}{\scriptsize (본 논문)} & \\
    \midrule
    기본 수식 & 선형 & 선형$\to$시그 & 시그모이드 & 개별 바운딩 & 개별 바운딩 &
      선형$\subset$시그모이드 확인 \\
    \addlinespace
    $R_0 = 6.0$h & \cmark & \cmark & \cmark & \cmark & \cmark &
      전 단계 동일 \\
    \addlinespace
    $R_{\max} = 24.0$h & --- & --- & --- & \cmark & \cmark &
      단계 4에서 최초 도입 \\
    \addlinespace
    $\kappa = 1.2$ & --- & --- & \cmark & \cmark & \cmark &
      단계 3 도출, 이후 적용 \\
    \addlinespace
    시드 체계 & 42 & 42 & 42 & 42(비중첩) & 42(비중첩) &
      마스터 시드 동일 \\
    \addlinespace
    $n = 10{,}000$ & \cmark & \cmark & \cmark & \cmark & \cmark &
      전 단계 동일 \\
    \addlinespace
    $\bar{\eta}_{\text{C1}}$ & 0.853 & 0.853 & 0.853 & 0.853 & 0.853 &
      전 단계 동일 \\
    \addlinespace
    P85.8 교차점 & --- & --- & --- & --- & \cmark &
      단계 5에서 발견·분석 \\
    \addlinespace
    $0.52$배 감쇠 & --- & --- & --- & --- & \cmark &
      단계 5에서 발견·해명 \\
    \addlinespace
    순서 준수율 (C1$\leq$C$x$) & \makecell{100\%\\(C1$\leq$C0)} & --- & --- & \makecell{99.87\%\\(C1$\leq$C2)} & \makecell{100\%\\(C1$\leq$C3)} &
      유형별 일관 \\
    \addlinespace
    전진 선택 & --- & --- & --- & $k_O$ & $U$ &
      지배 변수 전환 \\
    \addlinespace
    LASSO & --- & --- & --- & \makecell{$R^2\!=\!.995$\\(5변수)} & \makecell{$R^2\!=\!.852$\\(6변수)} &
      C2 간명, C3 복잡 \\
    \bottomrule
  \end{tabularx}
\end{table}

\noindent
교차 정합성 검증의 핵심 결과는 다음과 같다.

\chg{먼저 수식 체계의 일관적 확장을 보면,}
\citet{lee2026drto}의 선형 모델 $\DRTO = R_0 \cdot (1 - k_O \cdot O)$는
본 논문의 개별 시그모이드 바운딩 모델에서 $|z_O| \ll 1$ 일 때
$S(z_O) \approx z_O$로 근사하면 정확히 복원된다.
즉, 선형 모델은 시그모이드 바운딩 모델의 1차 근사
(first-order approximation)이며, 양 모델의 C1 예측값은
$|\Delta\eta| < 0.003$ 수준으로 일치한다.
이 관계는 \chg{\citet{lee2026drto}의} 바운딩 비교(단계 2)에서 명시적으로 분석되어
두 모델 간 전환의 이론적 정당성을 제공한다.

\chg{다음으로 데이터 기반의 수치 정합을 보면,}
모든 연구 단계가 동일한 마스터 시드(42), 표본 크기($n = 10{,}000$),
기본 복구시간($R_0 = 6.0$h)을 사용하여
$\bar{\eta}_{\text{C1}} = 0.853$이 전 단계에서 일치한다.
\chg{\citet{lee2026drto}에서} $\kappa = 1.2$가 도출된 이후 본 논문의 \chg{C2와 C3} 분석에서
동일한 $\kappa$ 값이 적용되어, 연구 단계 간 파라미터 일관성이 보장된다.

\chg{이어 발견의 누적적 심화를 보면,}
본 논문의 C2 분석(단계 4)에서 이중 채널(관측성+불확실성) 구조가 처음 정식화되고,
C3 분석(단계 5)에서 파레토 분포를 적용하여 P85.8 교차점과 $0.52$배 이중채널 감쇠가
\chg{발견되고 해명되는} 누적적 심화 구조를 갖는다.
C3의 $\mathrm{CVaR}_{99}$ $+24.0\%$ 결과는 C2의 가우시안 분석 결과와
대비되며, 극단 영역에서의 격차를 새로운 지표(CVaR)로 정량화한 것이다.

\chg{끝으로 변수 선택 방법의 교차 검증을 보면,}
C2 분석(단계 4)과 C3 분석(단계 5) 모두에서 전진 선택법과 LASSO 정규화를
적용하였으며, 양 방법이 공통으로 $k_O \cdot O$와 $U$를 핵심 변수로 선택하여
방법 간 일관성이 확인되었다.
LASSO의 $\lambda = 0.01$ 수준에서 C2($R^2 = 0.995$, 5변수)가
C3($R^2 = 0.852$, 6변수)보다 더 간명한 구조를 보이는 것은,
가우시안 분포의 규칙적 구조가 파레토의 비선형성보다
2차 다항식으로 효과적으로 근사됨을 교차 확인한 것이다.


%% -------------------------------------------------------------------
\subsection{\vrev{한계-해결 연쇄 구조}}
\label{subsec:limitation_resolution}
%% -------------------------------------------------------------------

$\DRTO$ 프레임워크의 연구 단계는 각 단계의 한계가 다음 단계의 연구 동기가 되는
\textbf{순차적 한계-해결 연쇄(limitation-resolution chain)} 구조를 형성한다.
\jrev{\tabref{tab:limitation_chain}\refpo{tab:limitation_chain}{은}{는}} 이 연쇄 연구단계의 5단계를 정리한 것이다.

\begin{table}[htbp]
  \centering
  \caption{\jrev{연구 단계별 한계-해결 연쇄}}
  \label{tab:limitation_chain}
  \small
  \begin{tabularx}{\textwidth}{c >{\raggedright}p{4.5cm} >{\raggedright\arraybackslash}X}
    \toprule
    \textbf{단계} & \textbf{발생 단계의 한계} & \textbf{다음 단계에서 해결 및 그 해결 수단} \\
    \midrule
    단계 1 & $\text{DRTO} \to 0$ 비현실적 (선형 모델) & 시그모이드 바운딩 $S(z) \in (-1,1)$ \\
    단계 2 & $\kappa$ 미결정 (바운딩) & 격자 탐색 $\to$ $\kappa^* = 1.2$ \\
    단계 3 & 단일 채널 ($\kappa$ 최적화) & 불확실성 채널 $E_{U,\text{bnd}}$ 추가 \\
    단계 4 & 경꼬리 한계 (C2 분석) & 파레토 분포 \\
    단계 5 & 단일분포·합성데이터 (C3 분석) & 통합 분석, 실무 적용, 전문가 검증 \\
    \bottomrule
  \end{tabularx}
\end{table}

5단계 연쇄가 \chg{\citet{lee2026drto}에서} 본 논문까지 단절 없이 연결되며,
각 단계의 한계 진술이 다음 단계의 동기에 정확히 반영됨을 확인하였다.

%% -------------------------------------------------------------------
\subsection{\vrev{매개변수 승계 일관성}}
\label{subsec:param_inheritance}
%% -------------------------------------------------------------------

\jrev{\tabref{tab:param_inheritance}\refpo{tab:param_inheritance}{은}{는}} 연구 단계에 걸친 핵심 매개변수의
승계 경로를 추적한 것이다.

\begin{table}[htbp]
  \centering
  \caption{\jrev{매개변수 승계 추적표}}
  \label{tab:param_inheritance}
  \small
  \begin{tabular}{l ccc cc}
    \toprule
    & \multicolumn{3}{c}{\textbf{\citet{lee2026drto}}} & \multicolumn{2}{c}{\textbf{본 연구}} \\
    \cmidrule(lr){2-4} \cmidrule(lr){5-6}
    \textbf{매개변수} & \textbf{단계 1} & \textbf{단계 2} & \textbf{단계 3} & \textbf{단계 4} & \textbf{단계 5} \\
    \midrule
    $R_0$ & 6.0h & 6.0h & 6.0h & 6.0h & 6.0h \\
    $n$ & 10,000 & 10,000 & 10,000 & 10,000 & 10,000 \\
    seed ($O$) & 42 & 42 & 42 & 42 & 42 \\
    seed ($k_O$) & 20042 & 20042 & 20042 & 20042 & 20042 \\
    $\kappa$ & --- & 가변 & 1.2 & 1.2 & 1.2 \\
    $R_{\max}$ & --- & --- & --- & 24.0h & 24.0h \\
    $S(z)$ 정의 & --- & $\frac{2}{1+e^{-z}}-1$ & 동일 & 동일 & 동일 \\
    \bottomrule
  \end{tabular}
\end{table}

모든 매개변수가 도입 시점 이후 일관되게 승계되고 있다.
$\kappa$는 \chg{\citet{lee2026drto}의 단계 3에서} 확정($\kappa^* = 1.2$)된 후
본 논문의 \chg{C2와 C3} 분석에서 변경 없이 사용되며,
$R_{\max}$는 본 논문의 C2 분석(단계 4)에서 최초 도입($24.0$h $= 4R_0$)된 후
C3 분석(단계 5)에서 동일값이 적용된다.
비중첩 시드 대역(O: [0--9999], $k_O$: [20000--29999], $U^{(G)}$: [30000--39999],
$U^{(P)}$: [40000--49999])으로 인해 상관계수 $|r| < 0.02$가 보장된다.


%% -------------------------------------------------------------------
\subsection{\vrev{교차검증 종합 결과}}
\label{subsec:comprehensive_pass}
%% -------------------------------------------------------------------

연구 전체에 대해 \chg{연구모형과 본문, 온톨로지 분석, 기준값}(statistics.json)
간 \chg{수치와 수식 그리고 용어}의 정합성을 교차 검증한 결과,
\textbf{총 303개 검증 항목이 모두 통과(PASS)}하였다(\jrev{\tabref{tab:303_pass}}).

\begin{table}[htbp]
  \centering
  \caption{\jrev{연구 단계별 종합 교차검증 결과}}
  \label{tab:303_pass}
  \small
  \begin{tabular}{l l ccc}
    \toprule
    \textbf{출처} & \textbf{연구 단계} & \textbf{검증 항목} & \textbf{합격} & \textbf{판정} \\
    \midrule
    \multirow{3}{*}{\citet{lee2026drto}}
    & 단계 1 (관측성 기반 선형) & 48 & 48 & PASS \\
    & 단계 2 (바운딩 기법 비교) & 72 & 72 & PASS \\
    & 단계 3 (곡률 최적화)     & 55 & 55 & PASS \\
    \midrule
    \multirow{2}{*}{\jrev{본 논문 4장}}
    & 단계 4 (가우시안 불확실성) & 64 & 64 & PASS \\
    & 단계 5 (파레토 불확실성)  & 64 & 64 & PASS \\
    \midrule
    & \textbf{합계}           & \textbf{303} & \textbf{303} & \textbf{PASS} \\
    \bottomrule
  \end{tabular}
\end{table}

이 검증은 수식 기호 일치, 매개변수 값 일치, 통계량 일치(소수점 규약 준수),
\chg{가설과 조건} 정의 일치를 포함하며, \chg{\citet{lee2026drto}와} 본 논문 간의
논리적 일관성과 데이터 신뢰성을 보증한다.
  % 4.12 선행연구 교차 정합성

% 가설 검증 종합 (마지막 절)


%% ===================================================================
\section{\vrev{가설 검증 종합}}
\label{sec:ch4_hypothesis_summary}
%% ===================================================================

본 절에서는 \jrev{\chref{ch:theory}}에서 수립한 \jrev{10개 연구가설(H1--H10)의 검증 결과를 종합하며, 본 연구의 핵심 학술 진술인 \textbf{\frev{분포 환경 적합성} 종합 명제}를 도출한다}.
\jrev{\tabref{tab:hypothesis_verification}\refpo{tab:hypothesis_verification}{은}{는}} 각 가설의 판정 기준, 실험 결과,
최종 판정을 일목요연하게 비교한 것이다.

\nrev{본 연구의 핵심 결론은 다음과 같다. 10개 연구가설 H1\,$\sim$\,H10이 모두 채택되었다.
구체적으로 H1 (관측성 효과, 100\% 준수), H2 (경계 보장, 100\% 준수),
H3 (조건 순서 보존, 99\% 이상 준수), H4 (회귀 설명력, $R^2 \geq 0.998$),
H5 (교호작용 유의성, $\Delta R^2 = 0.131$, $p < 0.001$), H6 (시드 독립성,
$|r|_{\max} = 0.016$), H7 (다중 시드 강건성, 45계수 모두 안정), H8 (이중채널
감쇠, 비율 0.52), H9 (전문가 정량 평균 $\geq 4.0$), H10 (전문가 정성 합의율
$\geq 80\%$) 가 모두 사전 정의된 판정 기준을 충족하였다. 이로써 본 연구의
$\DRTO$ 프레임워크는 \chg{시뮬레이션 검증, 회귀 분석, 강건성 점검, 이중채널 메커니즘 규명,
전문가 외부 검증}의 5축에서 통합적으로 입증되었다.}

\begin{table}[htbp]
  \centering
  \caption{\jrev{\jrev{10개 연구가설 검증 결과 종합} ($\kappa = 1.2$, $n = 10{,}000$)}}
  \label{tab:hypothesis_verification}
  \small
  \begin{tabularx}{\textwidth}{c l X l c}
    \toprule
    \textbf{가설} & \textbf{범주} & \textbf{판정 기준} & \textbf{실험 결과} & \textbf{판정} \\
    \midrule
    H1 & 관측성 효과 &
      $\DRTO_{\text{C1}} \leq R_0$, 준수율 100\% &
      $\bar{\eta}_{\text{C1}} = 0.853$, 100\% 준수 &
      \cmark \\
    H2 & 경계 보장 &
      $0 < \DRTO < R_{\max}$, 준수율 100\% &
      $\eta \in [0.468,\; 4.000]$, 100\% 준수 &
      \cmark \\
    \jrev{H3}-a & 순서 C1$\leq$C0 &
      준수율 100\% &
      100.00\% &
      \cmark \\
    \jrev{H3}-b & 순서 C1$\leq$C2 &
      준수율 $\geq 99\%$ &
      99.87\% &
      \cmark \\
    \jrev{H3}-c & 순서 C2$\leq$C3 &
      조건부 ($\geq$ P85.8) &
      31.8\% (전체); P85.8 이상 CDF 역전 &
      \cmark \\
    \midrule
    \jrev{H4} & 회귀 설명력 &
      $R^2_{(2\mathrm{q})} \geq 0.95$ &
      C1: 1.000, C2: 0.998, C3: 0.858 &
      \cmark$^*$ \\
    \jrev{H5} & 교호작용 유의성 &
      $\Delta R^2$ 기여도 유의 &
      $k_O \cdot O$: $\Delta R^2 = 0.131$ (C1) &
      \cmark \\
    \midrule
    \jrev{H6} & 시드 독립성 &
      $|r(X_i, X_j)| < 0.02$ &
      $|r|_{\max} = 0.016$ &
      \cmark \\
    \jrev{H7} & 다중 시드 강건성 &
      $10^8$회, 계수 안정 &
      45개 계수 중 기각 0개 &
      \cmark \\
    \midrule
    \jrev{H8} & 이중채널 감쇠 &
      Tail/Core $k_O$ 비율 &
      $|-0.415|/|-0.798| \approx 0.52$배 &
      \cmark$^\dagger$ \\
    \midrule
    \jrev{H9} & 전문가 정량 &
      $\bar{X} \geq 4.0 / 5.0$ &
      (\ref{sec:ch5_expert}항 참조) &
      $^\ddagger$ \\
    \jrev{H10} & 전문가 정성 &
      15개 하위 주제 동의율 $\geq 80\%$ &
      (\ref{sec:ch5_expert}항 참조) &
      $^\ddagger$ \\
    \bottomrule
  \end{tabularx}
\end{table}

\chg{표에서 $^*$는 C3 전체 $R^2 = 0.858$가 0.95 미만이나 P85.8 분할 시 Core $R^2 = 0.999$, Tail
$R^2 = 0.710$으로 분할 적용함을, $^\dagger$는 감쇠비 $0.52 < 1$로 \jrev{H8}을 충족함을, $^\ddagger$는 \jrev{H9}와
\jrev{H10}이 \jrev{4.11절}(전문가 검증)에서 별도 판정됨을 나타낸다.}


%% -------------------------------------------------------------------
\subsection{\vrev{가설 검증 해석}}
\label{subsec:hypothesis_interpretation}
%% -------------------------------------------------------------------

\chg{먼저 L1 시뮬레이션 검증(H1--\jrev{H3})을 보면,}
H1--\jrev{H3}-b까지의 5개 가설이 모두 충족되었다.
H2(경계 보장)와 \jrev{H3}(순서 보존)는 L0(Golden Compare)에서 47개 포인트를 통해
결정론적으로 사전 확인된 후, 본 단계에서 $n = 10{,}000$ 확률적 표본으로 재검증되었다.
관측성의 복구 단축 효과(H1, Cohen's $d = -1.180$)와
물리적 경계 보장(H2)이 전체 표본에서 예외 없이 확인되었다.
\jrev{($\kappa = 1.2$의 최적성은 H3 가설이 아닌 선행연구\citep{lee2026drto}에서 확정된 전제 조건으로, 본 연구에서는 \frev{분포 환경 적합성} 종합 명제(본 절 말미)에서 통합 해석된다.)}
\jrev{H3}-c는 C2$\leq$C3 순서가 전체 표본에서는 $31.8\%$에 불과하나,
P85.8 CDF 교차점 이상의 극단 영역에서 C3가 C2를 상회하므로
조건부 가설로서 충족된다.
\chg{여기서 $31.8\%$는 동일 인덱스 $i$에서 $\eta_{\text{C2}}(i) \leq \eta_{\text{C3}}(i)$인 쌍별(pairwise)
비율이며, P85.8은 C2와 C3의 누적분포함수(CDF)가 교차하는 분포적 지표이다. P85.8 기준 분할 시
Core($n = 8{,}580$)에서는 $21.4\%$만 C2$\leq$C3이나 Tail($n = 1{,}420$)에서는 $94.9\%$가 C2$\leq$C3으로
역전되어, 파레토 꼬리의 극단값이 집중되는 상위 $14.2\%$ 영역에서 C3가 C2를 압도적으로 상회한다.}

\chg{다음으로 L2 회귀 검증(\jrev{H4}--\jrev{H5})을 보면,}
C1에서 $R^2 = 1.000$, C2에서 $R^2 = 0.998$, C3 전체에서 $R^2 = 0.858$이 달성되었다.
C3 전체의 $R^2$가 $0.95$ 미만이며, P85.8 분할 시
Core($R^2 = 0.999$)와 Tail($R^2 = 0.710$)의 결과를 보이므로,
\jrev{H4}는 Core 영역 기준에서 충족된다.
교호작용항의 유의성(\jrev{H5})은 $k_O \cdot O$가 C1의 $R^2$를
$0.869 \to 1.000$으로 향상시키는 결정적 기여를 통해 확인되었다.

\chg{이어 L3 강건성 검증(\jrev{H6}--\jrev{H7})을 보면,}
비중첩 대역 시드 체계의 독립성(\jrev{H6}, $|r|_{\max} = 0.016$)과
$10^8$회 다중 시드 시뮬레이션의 안정성(\jrev{H7}, 기각 0개)이 모두 확인되어,
seed $= 42$의 결과가 특정 시드에 의한 우연이 아님이 입증되었다.

\chg{끝으로 L4 이중채널 검증(\jrev{H8})을 보면,}
C3 Tail 영역의 $k_O$ 계수($-0.415$)가 Core 영역($-0.798$) 대비
$0.52$배로 감쇠되어, 극단 불확실성 영역에서 관측성의 복구 단축 효과가
약화됨이 확인되었다. 이 이중채널 감쇠 효과는 극단적 불확실성 하에서
불확실성 채널이 지배적으로 작용하여 관측성의 영향력이
상대적으로 축소됨을 의미하며, \jrev{\chref{ch:discussion}}에서 그 실무적 함의를
상세히 논의한다.


\bigskip
\jrev{이상의 시뮬레이션 결과와 가설 검증은 4.12절(교차 일관성)의 교차 조건 분석\chg{(}효과 크기 비교,
순서 가설 검증, 이중채널 감쇠 메커니즘 규명\chg{)}의 기초 자료로 활용된다.}


%% -------------------------------------------------------------------
\subsection{\jrev{\frev{분포 환경 적합성} — 종합 명제}}
\label{subsec:distribution_adequacy}
%% -------------------------------------------------------------------

\jrev{본 연구는 \citet{lee2026drto}에서 격자 탐색으로 \chg{도출하여 출판한} 곡률 매개변수 $\kappa^{*} = 1.2$를 새로운 분포 환경(\chg{C2 가우시안 이중채널과 C3 파레토 중꼬리})에 그대로 적용했다. 신규 격자 탐색을 통한 직접적 최적성 검증 대신, 본 절은 다음 9개 가설(H2--H10)의 PASS 결과를 통합 해석함으로써 $\kappa = 1.2$의 \frev{분포 환경 적합성(distribution-environment adequacy)을 시사한다}:}

\jrev{\chg{첫째, 물리적 안전성으로서 H2가 PASS하여 $0 < \DRTO < R_{\max} = 24$h가 C0--C3 모두에서 100\%
준수되며, $\kappa = 1.2$에서 분포 환경과 무관하게 비현실적 발산이 발생하지 않는다. 둘째, 이론적
일치로서 H3가 PASS하여 조건 간 순서 보존(C1$\leq$C0, C1$\leq$C2, P85.8 이상 C2$\leq$C3)이 이론 예측과
정합하며, $\kappa = 1.2$에서 분포 환경의 직관적 단조성이 유지된다. 셋째, 통계적 설명력으로서 H4가
PASS하여 C1 $R^2 = 1.000$, C2 $R^2 = 0.998$, C3 Core $R^2 = 0.999$로 $\kappa = 1.2$ 모델이 분포 환경
데이터를 잘 설명하고, H5가 PASS하여 $k_O \cdot O$의 $\Delta R^2 = 0.131$로 교호작용이 실제 발현한다.
넷째, 결과 안정성으로서 H6이 PASS하여 $|r|_{\max} = 0.016 < 0.02$로 비중첩 시드의 통계적 독립이
보장되고, H7이 PASS하여 $10^8$회 시뮬레이션에서 45개 회귀 계수가 모두 안정하여 $\kappa = 1.2$ 결과가
시드 우연의 산물이 아님이 확인된다. 다섯째, 메커니즘 발현으로서 H8이 PASS하여 C3 분포 환경에서
Tail/Core $k_O$ 비율이 $0.52 < 1$로 이중채널 비선형 감쇠 메커니즘이 $\kappa = 1.2$ 하에서 성공적으로
포착된다. 여섯째, 외부 검증으로서 H9와 H10이 PASS하여 정량 설문 평균 $\geq 4.0$, 정성 합의율
$\geq 0.80$으로 $\kappa = 1.2$ 모델 전체에 대한 BCM/DR 전문가 외부 검증을 통과한다(4.11절 참조).}}

\jrev{\textit{\citet{lee2026drto}의 $\kappa^{*} = 1.2$를 본 연구의 새로운 분포 환경(\chg{C2 가우시안과 C3 파레토})에 적용한 결과, 9개 가설(H2--H10)의 PASS가 \chg{i)} 물리적 안전성, \chg{ii)} 이론적 일치, \chg{iii)} 통계적 설명력, \chg{iv)} 결과 안정성, \chg{v)} 메커니즘 발현, \chg{vi)} 외부 검증의 \frev{6가지 측면에서 집합적으로 시사한다}. 따라서 $\kappa = 1.2$는 본 연구의 분포 환경에서 \frev{운영 적합(operationally adequate)하다}.}}

\jrev{\frev{분포 환경 적합성}\chg{의 학술적 한계로서,} \chg{\citet{lee2026drto}의} 4기준 종합점수 $F(\kappa) = f_1 \cdot f_2 \cdot f_3 \cdot f_4$ (\chg{실무 유의성, 비포화, 효과 크기, 설명력})는 C1 조건 기반으로 정의되었다. 따라서 \chg{C2와 C3} 조건에서 $F(\kappa)$를 직접 적용하여 새로운 격자 탐색을 수행하기에는 부적합하며, 이는 본 연구가 명시하는 학술적 한계이다. 그러나 위 9개 가설의 PASS는 격자 탐색의 직접적 대안을 제공하며, $\kappa = 1.2$의 \frev{분포 환경 적합성을 다중 차원에서 시사한다}.}

\fpara{``$\kappa^{*} = 1.2$가 본 연구의 분포 환경에서 운영 적합''이라는 \chg{종합 명제의 학술적 위상은 운영 적합성과 그 범위에 관한 것으로,} 위 ``종합 명제''에서 H2--H10 PASS의 통합 해석이 입증하는 바는 ``$\kappa = 1.2$ 사용 모델이 C2/C3 조건의 평가 기준(\chg{물리 경계, 이론 일치, 통계 설명력, 결과 안정성, 메커니즘 발현, 외부 검증})을 만족함''이며, 이는 \emph{운영 적합성}(operational adequacy)의 입증에 해당한다. 환경별 \emph{최적성}(optimality)의 직접 입증\chg{,} 즉 C2/C3 각각에서의 $F(\kappa)$ 격자 탐색을 통한 환경 무관 최적 $\kappa$ 도출은 본 연구의 연구 범위 밖이며 후속 연구의 과제로 남긴다. 이에 따라 본 연구는 $\kappa = 1.2$를 BCM 일반 운영의 \emph{\chg{잠정적이고 보수적인} 실무 기본값}으로 채택하며 (4.13.3절 참조), 본 채택은 운영 적합성 입증과 후속 절의 다중 보강 근거에 의해 뒷받침된다.}


%% -------------------------------------------------------------------
\fpara{\subsection{$\kappa = 1.2$의 잠정적 및 보수적 채택을 위한 정당화}}
\fpara{\label{subsec:kappa_provisional_justification}}
%% -------------------------------------------------------------------

\fpara{본 절은 \chg{\citet{lee2026drto}의} 도출값 $\kappa = 1.2$를 BCM 일반 운영의 \emph{\chg{잠정적이고 보수적인} 실무 기본값}(provisional and conservative operational default)으로 채택하는 정당성을 정리한다. 위 4.13절의 H2--H10 PASS는 \emph{운영 적합성}을 시사할 뿐 환경별 최적성을 입증하지 않으므로, 본 절은 \emph{\chg{잠정적이고 보수적인 기본값 채택}}을 뒷받침하는 보조 근거를 제시한다.}

\fpara{\chg{연구 범위와 적용 근거의 측면에서,} 본 연구는 C2/C3 조건에서 $F(\kappa)$의 직접 격자 탐색을 \emph{연구 범위 외}로 두고 C1 조건 도출 $\kappa = 1.2$를 BCM 일반 운영의 기본값으로 채택한다. 이 채택은 다음 근거에 의해 보강된다\chg{. 첫째,} $\kappa$\chg{는} sigmoid 입력 스케일링의 형태 모수로서 \chg{분포 평균과 분산}에 직접 의존하지 않으며, $z_U$의 무차원 정규화 인자 $R_0/(R_{\max} - R_0)$가 두 채널 effective span을 \chg{균형화한다. 둘째,} $\kappa = 1.2$\chg{는} 운영 신호 표본에서 $|z|$를 sigmoid \chg{비포화 비선형} 활성 영역($|z| \in [0.5, 3]$)에 \chg{집중시키며, 이는 입력 분포 모양과 독립이다. 셋째,} 4.11절 전문가 패널(Likert 평균 $\geq 4.0$, 합의도 $\geq 0.80$)이 도메인 신뢰성을 \chg{보강한다. 넷째,} ISO 22301:2019의 단일 정적값 결정 정신과 \chg{부합한다.}}

\fpara{\chg{잠정적이고 보수적인 사용의 실무적 필연성의 측면에서,} BCM 운영 실무는 어떤 $\kappa$ 값을 반드시 사용해야 하므로, 환경 무관 최적 $\kappa$의 존재 여부가 후속 연구에서 결정되기 전까지 \chg{\citet{lee2026drto}에서 가장 정밀하게 도출된 후보값} $\kappa = 1.2$를 잠정 채택하는 것이 합리적이다. 다른 값(예: 1.0, 1.5)은 도출 근거가 약하고, 모델 미사용은 실무 비현실이며, 무작위 선택은 정당화가 불가능하다. 따라서 본 채택은 \emph{학술적 최적성 입증}이 아니라 \emph{더 나은 대안이 부재한 상태에서의 최선 보수 선택}이라는 실무적 필연에 근거한다. 후속 연구에서 환경 무관 $\kappa$가 발견되면 갱신하고, 부재가 진단되면 환경별 적응 정책이 정립되며, 어느 경우든 본 잠정 위상은 보수적 운영 기준으로 유지된다.}

         % 4.13 가설 검증 종합
       % 제4장 연구 결과
% 제5장 논의 및 결론
% 1차 심사 반영: 구 4.7 종합 논의 + 4.8--4.10 시사점 + 4.11 한계 + 구 제5장 결론(요약/성과/기여/향후) 통합

\chapter{\vrev{논의 및 결론}}
\label{ch:discussion}

\vrev{%
본 연구에서는 제4장(연구 결과)에 대한 종합 논의를 제시하고,
\chg{\m{이론적·방법론적·실무적}} 시사점을 도출한다.
이어서 본 연구의 한계를 기술하고, \chg{연구 요약, 주요 성과, 학술적 기여}를 정리하며,
향후 연구 방향을 제시함으로써 본 연구를 마무리한다.
}

% ─── 5단계 5-1단계 (2026-04-22): Plan H 구조 재편 ───
% 9개 절 → 5개 절 통합, 본문 텍스트 무수정, wrapper로 깊이 강등
%
% 새 구조:
%   5.1 종합 논의 (현 5.1 그대로 — 5.1.2 이론적 기여, 5.1.3 연구의 한계 잔존, 5-2에서 정리)
%   5.2 학술적 기여 (wrapper: 현 5.2 + 5.3 + 5.8 흡수 + 5.2.3 실증적 기여 신설)
%   5.3 실무적 기여 (wrapper: 현 5.4 흡수, 5-2에서 3항으로 재분류)
%   5.4 한계점 (현 sec5.4_limitations.tex 그대로, 절 번호만 자동 5.4)
%   5.5 결론 및 향후 연구 (wrapper: 현 5.6 + 5.7 + 5.9 통합)

% 5-2단계: 신규 목차 구조 적용 (본문 텍스트 무수정, 파일 분할 + suppression wrapper)
% =====================================================================
% 5.1 종합 논의 — 5단계 5-2 (메인 wrapper, 5.1.2/5.1.3 분리됨)
% sec5.1_synthesis.tex의 분할본을 \input으로 결합. 본문 텍스트 무수정.
% =====================================================================



%% ===================================================================
\section{\vrev{종합 논의}}
\label{sec:ch5_discussion}
%% ===================================================================

제4장(연구 결과)의 교차분석의 분석 결과를 종합하여 DRTO 개별 시그모이드
바운딩 모델의 이론적 기여, 방법론적 함의, 그리고 한계를 논의한다.
교차분석에서 도출된 \jrev{\chg{주요 결과와 발견, 검증}은} 다음 세 가지 핵심 메시지로 통합된다.
\label{subsec:ch5_integration}

\chg{본 장의 논의는 \chref{ch:introduction}에서 제기한 다섯 가지 연구 목적에 대한 응답으로 구성된다.
서론은 관측성 단일 채널 모델을 관측성과 불확실성의 이중 채널 $\DRTO$ 모델로 확장하고,
정적 기준에서 관측성과 불확실성을 단계적으로 도입하는 네 개 운영 조건 체계(C0--C3)를 설계하며,
모델의 효과와 강건성, 회귀 설명력, 이중채널 비선형 효과, 전문가 타당성의 다섯 측면을 다층 검증하고,
관측성과 극단 불확실성이 결합할 때 나타나는 이중채널 비선형 효과를 규명하며, 6개 산업에 적용
가능한 실무 가이드라인을 제시하는 것을 목표로 제시하였다. 본 연구는 이들 목적이 제4장(연구 결과)로 어떻게
충족되는지를 5.2(학술적 기여)와 5.3(실무적 기여)로
나누어 논의하며, 제4장(연구 결과)의 모든 결과 기여와 활용 귀착은 5.1.1(제4장 연구 결과의 기여와 활용 종합 매핑)에 모두 매핑하였다. 또한 세 개의 핵심 메시지는 다음과 같다.}

\chg{첫째, 핵심 메시지는 관측성과 불확실성의 구조적 반대 효과이다.}
C0$\to$C1 전환(Cohen's $d = -1.180$)과 C1$\to$C2 전환($d = +1.871$)의
효과 크기는 모두 \jrev{\textbf{대효과(large effect)}} 수준이며 방향이 반대이다.
이는 개별 시그모이드 바운딩 모델에서 $E_{O,\text{bnd}} \leq 0$과
$E_{U,\text{bnd}} \geq 0$으로 설계된 구조적 반대 효과가
실험적으로 확인된 것이다.
\pdferr{2차 회귀}에서 C1은 $R^2 = 1.000$, C2는 $R^2 = 0.998$에 도달하여,
이 반대 효과가 \pdferr{2차 다항식}의 정밀도로 정량화됨을 보여준다.
실무적으로 이 발견은 BCM 전략 수립 시 관측성 투자(복구 시간 단축)와
불확실성 대비(안전 마진 확보) 사이의 trade-off를 명확히 인식해야
함을 시사한다.

\chg{둘째, 핵심 메시지는 평균의 비차별성과 극단의 극적 분기이다.}
C2--C3 비교에서 Cohen's $d = -0.238$(\jrev{소효과(small effect)})이지만,
$\mathrm{CVaR}_{99}$ 격차는 $+24.0\%$, Tail $k_O$ 계수 감쇠는 $0.52$배이다.
이 ``평균에서의 비차별성과 극단에서의 극적 분기'' 패턴은
불확실성의 분포 형태가 평균적 의사 결정이 아닌 극단 위험 관리에서
결정적으로 중요함을 정량적으로 입증한다.
P85.8 교차점은 \jrev{이 두 체계의 분기점(자연적 경계)}이며,
분할 회귀에서 Core($R^2 = 0.999$)와 Tail($R^2 = 0.710$)의
높은 설명력은 \jrev{이 분기의 통계적 정당성}을 확인한다.

\chg{셋째, 핵심 메시지는 전문가 검증에 의한 실무적 타당성 확인이다.}
5명의 전문가가 전체 평균 $4.46/5.0$으로 프레임워크를 평가하였으며
($t = 3.68$, $p = 0.011$), 15개 하위 주제 전체에서 $80\%$ 이상
동의가 확인되었다. 이 결과는 시뮬레이션에 기반한 이론적 발견이
BCM 실무자의 경험적 관찰과 부합함을 삼각 검증(triangulation)한 것이다.
특히 전문가들이 P85.8 교차점과 이중 채널 감쇠 메커니즘의 실무적
설명력을 높이 평가한 점은 이 발견의 이론적 신규성과 실무적
유용성을 동시에 확인한다.


\subsection{\chg{제4장 연구 결과의 기여와 활용 종합 매핑}}
\label{subsec:ch5_results_map}

\chg{본 항은 \chref{ch:results}에서 도출된 모든 결과가 학술적 기여와 실무적 활용 가운데 어디로 귀착되는지를
일람으로 제시한다. 이는 의미의 크고 작음과 무관하게 각 결과가 본 장의 어느 논의에서 성과로 활용되고
기여로 정리되는지를 명시함으로써, 결과와 논의 사이의 누락 없는 연결을 보장하기 위함이다. 학술적 기여는
\secref{sec:academic-contrib}에서, 실무적 활용은 \secref{sec:practical-contrib}에서 상세히 전개되며,
\tabref{tab:ch5-results-map}\refpo{tab:ch5-results-map}{은}{는} 그 매핑을 결과 영역별로 정리한 것이다.}

\begin{table}[htbp]
  \centering
  \caption{\chg{제4장 결과 영역별 학술적 기여와 실무적 활용 귀착}}
  \label{tab:ch5-results-map}
  \small
  \begin{tabularx}{\textwidth}{l X X}
    \toprule
    \textbf{\chg{결과 영역(4장 절)}} & \textbf{\chg{핵심 결과}} & \textbf{\chg{학술 기여 / 실무 활용 귀착}} \\
    \midrule
    \chg{관측성 효과 (4.1)} & \chg{$\bar\eta_{\text{C1}}=0.853$($-14.7\%$), $d=-1.180$, H1$\cdot$H3-a 100\%} & \chg{이론(5.2.1)$\cdot$실증(5.2.3) / 회귀식 산출(5.3.1)$\cdot$성숙도 진단(5.3.3)} \\
    \addlinespace
    \chg{가우시안 프리미엄 (4.2)} & \chg{$\bar\eta_{\text{C2}}=1.682$($+68.2\%$), $d=+1.871$, 프리미엄율 84.9\%} & \chg{이론(5.2.1)$\cdot$실증(5.2.3) / 회귀식 산출$\cdot$안전마진 해석(5.3.1)} \\
    \addlinespace
    \chg{파레토$\cdot$극단 위험 (4.3)} & \chg{$\bar\eta_{\text{C3}}=1.514$, P85.8 CDF 교차점, CVaR$_{99}$ $+24.0\%$, CVaR 에스컬레이션} & \chg{이론(5.2.1) / \textbf{이중채널 자원배분$\cdot$극단 대비(5.3.2)}} \\
    \addlinespace
    \chg{회귀$\cdot$간명모형 (4.4)} & \chg{$R^2$ 진행(C1 1.000/C2 0.998/C3 0.858), 분할 Core 0.999, 4변수 간명모형, 지배변수 전환} & \chg{방법론(5.2.2) / \textbf{회귀식 DRTO 산출 도구(5.3.1)}$\cdot$산업별 투자우선순위(5.3.4)} \\
    \addlinespace
    \chg{이중채널 감쇠 (4.10)} & \chg{Tail $k_O$ $-0.415$ vs Core $-0.798$($0.52$배), Core 선형/Tail 포화 메커니즘} & \chg{이론(5.2.1) / \textbf{이중채널 스트레스테스트$\cdot$자원배분(5.3.2)}} \\
    \addlinespace
    \chg{\m{효과 크기}$\cdot$분포구조 (4.7--4.9)} & \chg{경로별 Cohen's $d$, 순효과 격차(가우시안 1.110 vs 파레토 0.650), 백분위$\cdot$PDF 교차, 회귀계수 직접감쇠($k_O$ 0.75배$\cdot$U 0.45배)} & \chg{이론(5.2.1)$\cdot$실증(5.2.3) / 산업별 환경계수$\cdot$투자효과 비교(5.3.4)} \\
    \addlinespace
    \chg{강건성$\cdot$기반검증 (4.5--4.6)} & \chg{기술통계 질적전환, $n{=}10{,}000$ 적정성, $10^8$ 다중시드(H6$\cdot$H7), Golden Compare 47점(L0)} & \chg{방법론(5.2.2)$\cdot$실증(5.2.3) / 재현성 신뢰 근거(5.3.3)} \\
    \addlinespace
    \chg{전문가 검증 (4.11)} & \chg{$4.46/5.0$($p{=}0.011$, $d{=}1.65$), CVI $0.985$, $\alpha{=}0.89$, 인터뷰 15/15 $\geq80\%$(H9$\cdot$H10)} & \chg{실증(5.2.3) / 실무 수용성$\cdot$현장 통찰(5.3 전반)} \\
    \addlinespace
    \chg{교차 정합성 (4.12)} & \chg{\citet{lee2026drto} 매개변수 승계 일관성, 303개 교차검증 PASS, 한계-해결 연쇄} & \chg{방법론(5.2.2)$\cdot$실증(5.2.3) / 연구 신뢰성 보증} \\
    \addlinespace
    \chg{가설 종합 (4.13)} & \chg{H1--H10 전원 채택, \frev{분포 환경 적합성} 종합 명제, $\kappa{=}1.2$ 잠정 채택} & \chg{실증(5.2.3)$\cdot$한계(5.4) / $\kappa$ 기본값 적용 근거(5.3)} \\
    \bottomrule
  \end{tabularx}
\end{table}

\chg{이 매핑이 보이는 핵심은, 제4장(연구 결과)의 결과가 학술적으로는 이론, 방법론, 실증의 세 기여로, 실무적으로는
회귀식 기반 산출, 이중채널 기반 자원배분, 성숙도와 에스컬레이션 운영, 산업별 적용의 네 도구로 빠짐없이
이어진다는 점이다. 특히 본 연구의 두 핵심 발견인 \textbf{회귀식(간명모형)}과 \textbf{이중채널 감쇠
메커니즘}은 각각 5.3(실무적 기여)의 독립된 실무 도구로 구체화되며, 이는 이론적 발견이
실무 의사결정으로 내려오는 경로를 명시한다.}




\bigskip
이상의 교차분석과 논의를 종합하면, $\DRTO$ 개별 시그모이드 바운딩 모델은
관측성과 불확실성의 효과를 구조적으로 분리하고, 극단 영역에서의
비선형 상호작용을 이중채널 감쇠로 포착하며,
전문가 검증에 의해 실무적 타당성이 확인된 프레임워크이다.
이어지는 절에서는 이러한 이론적 발견이 BCM 실무와 정책 수립에 갖는
구체적 시사점을 논의한다.


%% ===================================================================
%% 이하 구 제6장(시사점) 내용 — ch5와 병합
%% ===================================================================

이하에서는 \chref{ch:results}의 시뮬레이션 실험과 앞서 논의한 교차 조건 분석 결과를
토대로 $\DRTO$ 프레임워크의 시사점을 이론적, 방법론적, 그리고 실무적 관점에서
체계적으로 논의한다. 먼저 개별 시그모이드 바운딩 모델이 BCM 이론에
기여하는 바를 정리하고(\secref{sec:theoretical-implications}),
연구 방법론의 혁신과 확장 가능성을 논의한다(\secref{sec:methodological-implications}).
이어서 2단계 프레임워크, 에스컬레이션 체계, 6개 산업 적용 가이드라인 등
실무적 시사점을 제시하고(\secref{sec:practical-implications}),
마지막으로 연구의 한계점을 투명하게 논의한다(\secref{sec:limitations}).

                   % 5.1 종합 논의 (5.1.2/5.1.3 분리됨)
% =====================================================================
% 5.2 학술적 기여 — 5단계 5-2 (clean subsection wrapper)
% 본문 텍스트 무수정. 분할 본문 파일을 새 subsection 구조로 결합.
%   5.2.1 이론적 기여  ← 5.1.2 + 5.2(전체) + 5.8.1 본문
%   5.2.2 방법론적 기여 ← 5.3(전체) + 5.8.2 본문
% =====================================================================

\section{\vrev{학술적 기여}}
\label{sec:academic-contrib}

\begin{p5add}
본 학술적 기여는 본 연구가 BCM\chg{과 재해 복구}(DR) 학문 분야의
지식 체계에 새롭게 기여한 바를 \chg{\m{(i)} 이론적, \m{(ii)} 방법론적, 그리고 \m{(iii)} 실증적} 세 차원에서
종합한다. 이론적 기여(5.2.1항)에서는 $\DRTO$ 모델의 수학적 구조와
이중 채널 감쇠 메커니즘의 신규성을, 방법론적 기여(5.2.2항)에서는 비중첩
대역 시드 체계와 분할 회귀 방법론을, 실증적 기여(5.2.3항)에서는 \jrev{10개
가설} 전원 채택과 다중 검증 프레임워크의 결과\jrev{, 그리고 \textbf{분포 환경
적합성 종합 명제}(\secref{sec:ch4_hypothesis_summary})}를 다룬다.
\end{p5add}

% --- 5.2.1 이론적 기여 ---
\subsection{\vrev{이론적 기여}}
\label{subsec:academic-theoretical}

\begingroup
  % 원본 파일의 \section, \subsection 헤더만 시각적으로 억제 (label은 유지)
  \renewcommand{\section}[1]{}
  \renewcommand{\subsection}[1]{}
  \renewcommand{\subsubsection}[1]{\paragraph{#1}}

  % --- 주축: 5.1.2 이론적 기여 (간결한 intro + 1)/2)/3) 형식) ---
  %% -------------------------------------------------------------------
\subsection{\vrev{이론적 기여}}
\label{subsec:ch5_theoretical}
%% -------------------------------------------------------------------

본 연구의 이론적 기여는 네 가지 차원에서 정리된다.

\chg{첫째, 개별 시그모이드 바운딩에 의한 물리적 경계 보장이다.}
앞서 비교한 선형 RTO 모델의 물리적 경계 위반 문제를 개별 시그모이드
바운딩 함수로 해결하였다. $E_{O,\text{bnd}} \in (-R_0, 0]$과
$E_{U,\text{bnd}} \in [0, R_{\max} - R_0)$의 범위가 해석적으로 보장되어,
$\DRTO \in (0, R_{\max})$가 입력값에 무관하게 성립한다.
이는 기존 BCM 문헌에서 제시되지 않았던 수학적 보장으로서,
\chg{복구목표시간}의 음수 산출이나 MTPD 초과라는 구조적 문제를
근원적으로 해결한다.

\chg{둘째, 이중채널 감쇠 효과의 발견과 정량화이다.}
P85.8 교차점을 기준으로 $k_O$ 계수가 $0.52$배로 감쇠되는
이중채널 현상의 발견은 BCM 이론에 새로운 통찰을 제공한다.
이 발견은 극단 불확실성과 관측성이 비선형적으로 상호 작용하는
메커니즘을 규명하며, 시그모이드 비선형성에 의한 이론적 설명을
제시한다.

\chg{셋째, \citet{lee2026drto}과 본 연구의 일관적 이론 체계 구축이다.}
선행 연구\citep{lee2026drto}의 선형 근사$\to$시그모이드 바운딩$\to$$\kappa$ 최적화에서
본 연구의 이중채널 확장(C2, C3)까지의 단계적 확장이
\chg{수식·데이터·결론}의 교차 정합성을 유지하며 수행되었다.
선형 모델이 시그모이드 바운딩의 1차 근사임을 보이는 관계
($|z| \ll 1$에서 $S(z) \approx z$)는 이론적 확장의
내적 일관성을 수학적으로 보장한다.

  넷째, 불확실성의 확률적 정량화이다. 복구 과정의 불확실성을 일상적 변동을 나타내는 가우시안과 극단 사건을 나타내는 파레토의 두 분포로 정량화함으로써, 관측성 단일 차원에 머물던 기존 모델의 한계를 보완하고 극단 영역의 꼬리 위험을 복구목표시간 산정에 반영하였다.



  % --- 보완: 5.2 이론적 시사점 (기술적 세부 설명 보완) ---
  \section{\vrev{이론적 시사점}}
\label{sec:theoretical-implications}
%% ===================================================================

본 연구의 이론적 기여는 BCM 분야에 \jrev{\chg{i) 관측성 개념의 통합, ii) 불확실성의 확률적 정량화(가우시안과 파레토), iii) 개별 시그모이드 바운딩 이론의 정립, iv) 이중 채널 감쇠 효과의 발견이라는} 네 가지 축}으로 정리된다.
\p{[}표~\vref{tab:ch6-theoretical-contributions}\p{]}는 이론적 기여의 핵심 내용을 종합한다.

\begin{table}[htbp]
\centering
\caption{\jrev{이론적 기여 종합}}
\label{tab:ch6-theoretical-contributions}
\small
\begin{tabularx}{\textwidth}{c l X c}
\toprule
\textbf{\#} & \textbf{기여 영역} & \textbf{핵심 내용} & \textbf{근거} \\
\midrule
1 & 관측성 통합 &
  제어 이론의 관측성 개념을 BCM에 최초 도입;
  $z_O = \kappa \cdot k_O \cdot O_{\mathrm{raw}}$으로 정량화하여
  정적 RTO의 구조적 한계를 극복 &
  H1, \jrev{H3}-a \\
\addlinespace
2 & 시그모이드 바운딩 이론 &
  수정 시그모이드 $S(z) = 2/(1+e^{-z})-1$을 기저 함수로 채택;
  $\DRTO = R_0 + E_{O,bnd} + E_{U,bnd}$로 물리적 경계
  $(0,\; R_{\max})$를 해석적으로 보장 &
  \jrev{H2 ($\kappa = 1.2$ 전제)} \\
\addlinespace
3 & 관측성--불확실성 이원 구조 &
  관측성 효과($E_{O,bnd} \leq 0$)와 불확실성 효과($E_{U,bnd} \geq 0$)의
  구조적 반대 방향 확인; Cohen's $d$: $-1.180$ vs $+1.871$ &
  \jrev{H1, H8} \\
\addlinespace
4 & P85.8 CDF 교차점 규명 &
  가우시안(C2)과 파레토(C3) 분포의 $\eta$ 분위수 함수가
  P85.8에서 교차; 꼬리 효과의 조건부 특성 규명 &
  \jrev{H3}-c \\
\addlinespace
5 & 이중채널 감쇠 발견 &
  P85.8 기준 분할 회귀에서 Tail 영역의 $k_O$ 계수($-0.415$)가
  Core($-0.798$) 대비 $0.52$배 감쇠; 극단 환경에서 관측성
  효과의 상대적 축소 &
  \jrev{H8} \\
\addlinespace
6 & 꼬리 리스크 비선형 증폭 &
  CVaR$_{99}$: C3가 C2 대비 $+24.0\%$;
  분위수 상승에 따른 비선형 격차 확대
  (CVaR$_{90}$: $+20.4\%$ $\to$ CVaR$_{99}$: $+24.0\%$) &
  \jrev{H8} \\
\bottomrule
\end{tabularx}
\end{table}


%% -------------------------------------------------------------------
\subsection{\vrev{\revdel{BCM 이론에 대한 관측성 통합}}}
\label{subsec:observability-integration}
%% -------------------------------------------------------------------

\fdel{전통적 BCM 이론에서 RTO는 BIA(업무영향분석) 시점에 고정되는
정적 매개변수로 취급되어, 시스템 상태의 실시간 변화를 반영하지 못하였다.
본 연구는 제어 이론~\citep{kalman1960control}의 관측성(observability) 개념을
BCM 분야에 최초로 도입하여, 시스템의 내부 상태를 외부 출력으로부터
추정할 수 있는 정도를 $O_{\mathrm{raw}} \in [0,1]$로 정규화하고,
이를 관측성 가중치 $k_O$와 곡률 매개변수 $\kappa = 1.2$를 통해
$\DRTO$의 관측성 조정량으로 변환하였다.}

\fdel{관측성 조정량 $E_{O,bnd} = -R_0 \cdot S(\kappa \cdot k_O \cdot O_{\mathrm{raw}})$는
항상 비양수($\leq 0$)이므로, 관측성이 높을수록 $\DRTO$가 $R_0$ 이하로 단축된다.
이 구조는 \jrev{H3}-a에서 $100\%$ 준수율로 검증되었으며, 이는 확률적 결과가 아닌
모델의 \textbf{구조적 필연}이다. Cohen's $d = -1.180$의 매우 큰 효과 크기는
관측성 통합이 BCM의 복구 예측 정확도를 질적으로 향상시킨다는 이론적 근거를 제공한다.}

\fdel{이러한 관측성 통합은 기존 BCM 이론의 세 가지 공백을 해소한다.
첫째, RTO를 ``고정 상수''에서 ``상태 의존 함수''로 전환함으로써
동적 환경에서의 BCM 의사결정 기반을 확보한다.
둘째, APM(Application Performance Monitoring), SIEM(Security Information and
Event Management) 등 기존 모니터링 인프라의 BCM 활용 경로를 이론적으로 정당화한다.
이는 SRE(Site Reliability Engineering) 분야에서 관측성을 에러 예산(error budget)과
연계하여 서비스 신뢰성을 정량적으로 관리하는 접근(Beyer et al.\ 2016)과
맥락을 공유하며, DevOps 문화에서 모니터링 인프라 투자가 복구 성능을 향상시킨다는
실증적 발견(Forsgren et al.\ 2018; Kim et al.\ 2016)과도 부합한다.
셋째, 관측성 투자($k_O$ 증가)와 불확실성 허용($k_U$ 증가) 간의
trade-off를 $k_O + k_U = 1$ 제약으로 형식화하여,
BCM 자원 배분의 분석적 프레임워크를 제시한다.}


%% -------------------------------------------------------------------
\subsection{\vrev{\revdel{개별 시그모이드 바운딩 이론}}}
\label{subsec:sigmoid-bounding-theory}
%% -------------------------------------------------------------------

\fdel{$\DRTO$ 모델의 수학적 핵심은 수정 시그모이드 함수
$S(z) = 2/(1+e^{-z})-1 \in (-1,1)$에 의한 개별 바운딩(individual bounding)이다.
관측성 효과와 불확실성 효과를 각각 독립적인 시그모이드로 바운딩하여
$\DRTO = R_0 + E_{O,bnd} + E_{U,bnd}$로 정식화함으로써, 다음의 세 가지
이론적 요구사항을 동시에 충족한다.}

% [5단계 Plan H 재편: 구 5.2.2 "개별 시그모이드 바운딩 이론" 3개 항목
%  (물리적 경계 보장 / 무한 미분 가능성 / 교호작용 포착)은 5.2.1 이론적
%  기여로 통합·흡수됨. 기존 \fdel{} 래핑된 본문은 빈 enumerate 번호(1/2/3)
%  잔재를 발생시켜 enumerate 블록 전체 제거.]



%% -------------------------------------------------------------------
\subsection{\vrev{\revdel{이중채널 감쇠 효과와 P85.8 교차점}}}
\label{subsec:dual-channel-crossover}
%% -------------------------------------------------------------------

\fdel{본 연구의 가장 핵심적인 이론적 발견은 P85.8 분할 기준점(CDF 교차점과 일치)을
경계로 하는 이중채널 감쇠(dual-channel attenuation) 효과이다.
C3 조건에서 P85.8 기준으로 분할 회귀를 수행한 결과,
관측성 가중치 $k_O$의 회귀 계수가 Core 영역($\leq$ P85.8)의 $-0.798$에서
Tail 영역($>$ P85.8)의 $-0.415$로 $0.52$배 감쇠된다.}

\fdel{이 발견의 이론적 함의는 세 가지이다.
첫째, 극단적 불확실성 환경에서 관측성 투자의 \textbf{한계효용이 감소}한다.
Core 영역에서 $k_O$를 $0.1$ 증가시키면 $\eta$가 $0.080$ 감소하나,
Tail 영역에서는 동일한 변화가 $\eta$를 $0.042$만 감소시킨다.
이는 BCM 자원 배분에서 극단 시나리오일수록 관측성 투자의
비용-효과성이 감소하며 불확실성 관리가 우선되어야 함을 의미한다.
둘째, 시그모이드 비선형성과 파레토의 \textbf{상호작용 메커니즘}이
규명되었다. Tail 영역에서 큰 $U$ 값이 시그모이드의 포화 구간을 활성화하여
불확실성 채널이 지배적이 되며, $k_O$의 상대적 역할이 축소된다.
셋째, P85.8이 분할 회귀 기준점이자 CDF 교차점으로서 가우시안--파레토 분포 전환의
\textbf{구조적 경계}임이 다중 시드 검증을 통해 확인되어,
이 교차점이 특정 시드에 의존하지 않는 강건한 현상임이 입증되었다.}


  % --- 5.8.1 (이론적 기여) — 5.1.2와 3중 주제 중복으로 5-2 삭제 ---
  \begin{p5del}{현 5.8.1 이론적 기여 — 5.1.2의 1)/2)/3)과 동일 주제 중복}
  %% ===================================================================
\section{\vrev{학술적 기여}}
\label{sec:conclusion-contributions}
%% ===================================================================


%% -------------------------------------------------------------------
\subsection{\vrev{이론적 기여}}
\label{subsec:contrib-theoretical}
%% -------------------------------------------------------------------

\fdel{본 연구의 이론적 기여는 다음의 세 가지로 정리된다.}

\fdel{첫째, \textbf{업무연속성관리$\cdot$관측성$\cdot$시그모이드 바운딩의 통합적 이론 기반을
확립}하였다.
기존 BCM 연구에서 RTO는 BIA 시점에 고정되는 정적 매개변수로 취급되었으며,
관측성과 불확실성을 단일 수리 모델에서 통합한 선행 연구는 부재하였다.
본 연구는 제어 이론의 관측성 개념과 확률론적 불확실성 모형을
수정 시그모이드 함수라는 공통 기저 함수 위에서 결합함으로써,
BCM 분야에 새로운 이론적 기반을 제공한다.
특히 $\DRTO = R_0 + E_{O,\text{bnd}} + E_{U,\text{bnd}}$의 가산 구조는
각 효과의 독립적 분석과 결합 효과의 동시 포착을 가능하게 하여,
모델의 해석 가능성(interpretability)과 확장 가능성(extensibility)을
동시에 확보한다.}

\fdel{둘째, \textbf{꼬리 감쇠 현상을 발견하고 정량화}하였다.
P85.8 분할 기준을 경계로 관측성 가중치 $k_O$의 효과가
Core($-0.798$)에서 Tail($-0.415$)로 $0.52$배 감쇠되는 현상은,
극단 불확실성 환경에서 시그모이드 포화에 의해
관측성 채널의 조절 효과가 구조적으로 제한됨을 규명한 것이다.
이 발견은 BCM 자원 배분 전략에 새로운 분석적 기반을 제시하며,
극단 상황에서의 대비 전략이 관측성 투자만으로는 한계가 있으며
사전 자원 확보가 병행되어야 한다는 실무적 시사점을 제공한다.}

\fdel{셋째, \textbf{P85.8 CDF 교차점을 가우시안--파레토 분포 전환의 구조적 경계로 규명}하였다.
C2와 C3의 CDF가 P85.8($\eta \approx 2.42$)에서 교차하며,
이 교차점을 경계로 두 조건의 행태가 질적으로 전환된다.
P85.8 이하에서 두 분포는 사실상 구분 불가능하나,
P85.8 초과에서 C3의 $\eta$가 C2를 급격히 상회하여
CVaR$_{99}$ 기준 $+24.0\%$의 격차를 보인다.
이는 파레토 불확실성의 효과가 \textit{평균 수준의 비차별성}과
\textit{꼬리 영역의 극적 분기}로 구성되는 이중 구조를 가진다는 것을
실증적으로 규명한 결과이다.}


  \end{p5del}
\endgroup

% --- 5.2.2 방법론적 기여 ---
\subsection{\vrev{방법론적 기여}}
\label{subsec:academic-methodological}

\begingroup
  \renewcommand{\section}[1]{}
  \renewcommand{\subsection}[1]{}
  \renewcommand{\subsubsection}[1]{\paragraph{#1}}

  

%% ===================================================================
\section{\vrev{방법론적 시사점}}
\label{sec:methodological-implications}
%% ===================================================================

본 연구는 $\DRTO$ 프레임워크의 검증 과정에서 \jrev{네 가지} 방법론적 혁신을
달성하였다. \chg{이를 차례로 서술한다.}


%% -------------------------------------------------------------------
\subsection{\vrev{몬테카를로 시뮬레이션 방법론}}
\label{subsec:mc-methodology}
%% -------------------------------------------------------------------

\chg{첫째, 몬테카를로 시뮬레이션 방법론이다.} $n = 10{,}000$ 표본 $\times$ 4개 조건(C0--C3) $\times$ $10{,}000$개 마스터 시드의
총 $4 \times 10^8$회 $\DRTO$ 계산을 통해, 시뮬레이션 기반 BCM 모델 검증의
새로운 표준을 제시하였다.

\textbf{비중첩 대역 시드(non-overlapping band seeding)} 체계는 마스터 시드
$s = 42$로부터 변수별 시드를 $10{,}000$ 간격의 비중첩 대역\chg{(}$O_{\mathrm{raw}}$:
$[0\text{--}9999]$, $k_O$: $[20000\text{--}29999]$, $U^{(G)}$:
$[30000\text{--}39999]$, $U^{(P)}$: $[40000\text{--}49999]$\chg{)}으로
파생한다. \jrev{H6} 검증에서 6개 변수쌍 모두 $|r| < 0.02$(최대 $|r| = 0.016$)를
충족하여, 시뮬레이션의 완전한 재현성과 변수 간 통계적 독립성이 동시에 보장되었다.

이 체계의 방법론적 의의는 세 가지이다.
\chg{첫째,} 단일 시드로부터 다수의 독립적 난수열을 파생하는 표준화된 규약을 제공한다.
\chg{둘째,} $10{,}000$개 시드에 걸친 $10^8$회 다중 시드 검증(\jrev{H7})에서 Bonferroni 보정
$z$-검정 결과 45개 계수 중 기각이 $0$개로, 특정 시드에 대한 의존성을 완전히
배제한다. \chg{셋째,} 재현성 프로토콜이 명시되어 있어 후속 연구자의 독립적 검증이 가능하다.


%% -------------------------------------------------------------------
\subsection{\jrev{\texorpdfstring{\citet{lee2026drto}}{Lee and Cheung (2026)} $\kappa$ 매개변수의 본 연구 환경 적용 및 \frev{분포 환경 적합성} 시사 방법론}}
\label{subsec:kappa-optimization}
%% -------------------------------------------------------------------

\chg{둘째, $\kappa$ 매개변수의 분포 환경 적합성 시사 방법론이다.} \jrev{본 연구는 곡률 매개변수 $\kappa^{*} = 1.2$를 \chg{\citet{lee2026drto}}에서 다기준 종합점수($f_1$ 실무 유의성,
$f_2$ 비포화율, $f_3$ 효과 크기, $f_4$ 설명력)로 \chg{도출하여 출판된} 매개변수로
받아들이고, 본 연구의 새로운 분포 환경(\chg{C2 가우시안과 C3 파레토})에 그대로 적용하였다. 따라서 $\kappa^{*}$ 도출 자체는 본 연구의
독자적 기여가 아니며, \chg{\citet{lee2026drto}의} 격자 탐색 결과(C1 조건 기준
$F(\kappa^{*}) = 0.9997$)를 \textbf{전제 조건}으로 받아들인다.}

\jrev{본 연구의 방법론적 보완은 다음과 같다. \chg{\citet{lee2026drto}의}
4기준 종합점수 $F(\kappa) = f_1 \cdot f_2 \cdot f_3 \cdot f_4$는
C1 환경 기반으로 정의되어, 새로운 분포 환경(\chg{C2 가우시안과 C3 파레토})에서는
직접 격자 탐색에 부적합하다는 학술적 한계를 지닌다. 본 연구는 격자
탐색의 직접적 대안으로서, 9개 가설(\jrev{H2--H10}) PASS의 \textbf{통합 해석을
통한 \jrev{\frev{분포 환경 적합성}} 종합 명제 도출 방법론}을 제시하였다
(\secref{sec:ch4_hypothesis_summary} \jrev{\frev{분포 환경 적합성}} 종합 명제 참조).
구체적으로, \chg{i) 물리적 안전성(H2), ii) 이론적 일치(H3), iii) 통계적 설명력(H4·H5),
iv) 결과 안정성(H6·H7), v) 메커니즘 발현(H8), vi) 외부 검증(H9·H10)}의 6가지 차원에서 $\kappa = 1.2$의 \jrev{\frev{분포 환경 적합성}}이 집합적으로
\frev{시사된다}.}

\jrev{이러한 \textbf{다중 가설 PASS의 통합 해석} 방법론은 \chg{\citet{lee2026drto}}
매개변수의 새로운 환경 적용성을 \frev{시사}하는 보완적
방법론적 틀로서, 본 연구가 독자적으로 제시하는 기여이다. 본 접근법은
선행연구를 기반으로 분포 환경을 확장하는 후속 연구 일반에 적용 가능하며,
직접적 격자 탐색이 부적합한 환경에서의 매개변수 적합성 검증에 활용될 수 있다.}


%% -------------------------------------------------------------------
\subsection{\vrev{분할 회귀(Split Regression) 방법론}}
\label{subsec:split-regression}
%% -------------------------------------------------------------------

\chg{셋째,} P85.8 기준 분할 회귀는 파레토 분포의 체제 전환(regime change)을
포착하기 위한 방법론적 혁신이다. C3 전체 표본에서 \pdferr{2차 회귀}의
$R^2 = 0.858$에 그치는 잔여 비설명 분산($14.2\%$)을,
P85.8 기준으로 Core($\leq$ P85.8, $n = 8{,}580$)와 Tail($>$ P85.8, $n = 1{,}420$)로
분할함으로써 Core $R^2 = 0.999$, Tail $R^2 = 0.710$으로 개선하였다.

분할 회귀의 방법론적 기여는 다음과 같다.
\chg{i)} 단일 회귀 모형으로 포착할 수 없는 \textbf{구간별 이질적 구조}를
데이터 기반으로 식별하는 체계적 절차를 제시한다.
\chg{ii)} 분할 기준점(P85.8)의 결정이 사전적 이론(CDF 교차)에 근거하므로,
임의적 구간 분할의 자의성 문제를 해소한다.
\chg{iii)} 이 방법론은 BCM 분야에 국한되지 않고, 파레토 분포가 관여하는
금융 리스크 관리, 보험 계리, 자연 재해 모형 등에
범용적으로 확장 가능하다.

  
%% -------------------------------------------------------------------
\subsection{\vrev{방법론적 기여}}
\label{subsec:contrib-methodological}
%% -------------------------------------------------------------------

\fdel{방법론적 관점에서 본 연구는 다음의 세 가지 기여를 달성하였다.}

\fdel{첫째, \textbf{P85.8 기준 분할 회귀(split regression)를 파레토 분포 분석에 적용}하였다.
C3 전체 표본에서 2차교호작용 회귀의 $R^2 = 0.858$에 그치는 잔여 비설명 분산을,
P85.8 기준으로 Core($\leq$ P85.8)와 Tail($>$ P85.8)로 분할함으로써
Core $R^2 = 0.999$(0.9990), Tail $R^2 = 0.710$(0.7100)으로 Core 영역의
결정론적 구조를 완전히 포착하였다.
이 분할 전략은 파레토 분포의 체제 전환(regime change)을 포착하는
범용적 분석 방법론으로, BCM을 넘어 보험, 금융, 재난관리 등
파레토 현상이 핵심인 분야에 확장 적용될 수 있다.}

\fdel{둘째, \textbf{비중첩 대역 시드(non-overlapping band seeding) 체계}를 설계하였다.
마스터 시드 $s = 42$로부터 변수별 시드를 $10{,}000$ 간격의 비중첩 대역($O_{\mathrm{raw}}$: $[0\text{--}9999]$, $k_O$: $[20000\text{--}29999]$,
$U^{(G)}$: $[30000\text{--}39999]$,
$U^{(P)}$: $[40000\text{--}49999]$)으로 파생하여,
시뮬레이션의 완전한 재현성을 보장하면서도 변수 간 통계적 독립성을
유지한다(모든 변수 쌍 $|r| < 0.02$).}

\jrev{넷째}, \textbf{6단계 다층 검증(six-stage multi-layer validation) 프레임워크}를 구축하였다.
기반 검증(L0, Golden Compare), 시뮬레이션 검증, 회귀 검증, 강건성 검증, 이중채널 검증, 전문가 평가의
6개 단계(L0--L5)가 계층적 의존 구조를 형성하여 모델의 타당성을 다각적으로 보장한다.
정량적 시뮬레이션과 정성적 전문가 평가를 삼각검증(triangulation)함으로써
단일 방법론의 편향을 최소화하였으며, 이 프레임워크는
수리 모델 기반 연구의 검증 설계 표준으로 활용될 수 있다.



\endgroup

% --- 5.2.3 실증적 기여 (5-2단계 신규 작성) ---
\subsection{\vrev{실증적 기여}}
\label{subsec:academic-empirical}

\begin{p5add}
본 연구에서 실증적 기여는 시뮬레이션과 전문가 평가의 결합 검증을 통해 달성한
결과를 토대로 정리한다.

첫째, \textbf{\jrev{10개 연구 가설}(H1--\jrev{H10})이 사전 설정 기준을 모두 충족하여 모두
채택}되었다. 시뮬레이션 검증(H1--\jrev{H3}), 회귀 검증(\jrev{H4}--\jrev{H5}), 강건성 검증(\jrev{H6}--\jrev{H7}),
이중 채널 검증(\jrev{H8}), 전문가 평가(\jrev{H9}--\jrev{H10})의 6단계 다층 검증 프레임워크에서
계층적 의존 구조가 모두 만족됨으로써, $\DRTO$ 모델의 내부 정합성과
외부 타당성이 동시에 확인되었다.

둘째, \textbf{$10{,}000$개 마스터 시드에 걸친 $10^8$회 다중 시드 검증}을
통해 결과의 시드 의존성을 완전히 배제하였다. Bonferroni 보정 $z$-검정에서
45개 회귀 계수 중 기각이 0개로, 시뮬레이션 결과가 특정 난수열에 의존하지 않는
강건성을 정량적으로 입증하였다.

셋째, \textbf{BCM/DR 분야 전문가 5명에 의한 \chg{정량적 및 정성적} 혼합 검증}을 통해
실무 타당성을 확보하였다. 27문항 5점 리커트 척도에서 전체 평균
$4.46/5.0$($t(4) = 3.68$, $p = 0.011$, Cohen's $d = 1.65$)을 기록하였으며,
S-CVI/Ave $= 0.985$, Cronbach's $\alpha = 0.89$로 내용 타당도와 내적
일관성이 확인되었다. 심층 인터뷰에서는 5대 주제 15개 하위 주제 모두에서
$\geq 80\%$ 동의가 달성되어, 시뮬레이션 결과가 실무자의 경험적 관찰과
부합함을 삼각 검증하였다.

\jrev{넷째, \textbf{\frev{분포 환경 적합성} 종합 명제}를 도출하였다. \chg{\citet{lee2026drto}에서 격자
탐색으로 도출하여 출판한} $\kappa^{*} = 1.2$를 본 연구의
새로운 분포 환경에 적용한 결과,
9개 가설(\jrev{H2--H10})의 PASS가 \chg{i)} 물리적 안전성, \chg{ii)} 이론적 일치, \chg{iii)} 통계적
설명력, \chg{iv)} 결과 안정성, \chg{v)} 메커니즘 발현, \chg{vi)} 외부 검증의 6가지 측면에서
$\kappa = 1.2$의 \jrev{\frev{분포 환경 적합성}}을 집합적으로 \frev{시사한다} (\secref{sec:ch4_hypothesis_summary} 참조).
신규 격자 탐색을 직접 수행하지 않고도 다중 가설 PASS의 통합 해석을 통해
\chg{\citet{lee2026drto}의 매개변수의} 본 연구 환경 적용성을 학술적으로 정직하게 \frev{시사한 것이}
본 연구의 \chg{방법론적 및 실증적} 차별점이다.}
\end{p5add}
       % 5.2 학술적 기여
% =====================================================================
% 5.3 실무적 기여 — 3차심사 대규모 재설계 (발견→도구 축)
%   5.3.1 회귀식 기반 DRTO 산출·예측 도구          [신설]
%   5.3.2 이중채널 감쇠 기반 자원배분·스트레스테스트 [신설]
%   5.3.3 η 기반 평상시 성숙도·위기시 에스컬레이션   [구 5.3.1 성숙도 + 구 5.3.3 에스컬레이션 통합]
%   5.3.4 산업별 적용 가이드라인                    [구 5.3.2]
% =====================================================================

\section{\vrev{실무적 기여}}
\label{sec:practical-contrib}
\label{sec:practical-implications}    % 구 라벨 호환 — 외부 cross-ref 유지

\begin{p5add}
\chg{본 절은 본 연구가 BCM 실무 현장에 새롭게 제공한 의사결정 도구를 네 항으로 정리한다. 핵심은 본 연구의
두 핵심 발견인 회귀식(간명 모형)과 이중채널 감쇠 메커니즘을 각각 독립된 실무 도구로 구체화하여, 이론적
발견이 현장의 산출과 자원 배분으로 직접 내려오는 경로를 제시하는 데 있다.}
\chg{먼저 회귀식 기반 $\DRTO$ 산출과 예측 도구(5.3.1항)에서는 간명 회귀 모형으로 조직이 자기 환경의 $\DRTO$를
직접 계산하는 절차를, 이중채널 감쇠 기반 자원배분과 스트레스 테스트(5.3.2항)에서는 P85.8 분할과 $0.52$배
감쇠를 자원 배분 규칙과 극단 시나리오 설계로 전환하는 방안을 제시한다. 이어 $\eta$ 기반 평상시 성숙도와
위기시 에스컬레이션(5.3.3항)에서는 단일 지표 $\eta$의 평상시 진단과 위기시 대응의 이중 활용 및 선순환
구조를, 산업별 적용 가이드라인(5.3.4항)에서는 6개 대표 산업의 권장 파라미터와 분포 환경별 적용 원칙을
제시한다. 이들은 모두 \chref{ch:results}의 시뮬레이션과 회귀분석, 전문가 검증에서 도출된 정량적 결과를
실무 운용 형태로 변환한 산출물이다.}
\end{p5add}


% ===================================================================
% 5.3.1 회귀식 기반 DRTO 산출·예측 도구 [신설]
% ===================================================================
\subsection{\chg{회귀식 기반 $\DRTO$ 산출과 예측 도구}}
\label{subsec:practical-regression-tool}

\chg{본 항은 \chref{ch:results}의 회귀분석에서 도출된 간명 회귀 모형을 실무 현장의 $\DRTO$ 산출과 예측 도구로
제시한다. 이는 본 연구의 핵심 산출물인 회귀식이 단순한 통계적 적합을 넘어, 조직이 자기 환경의 $\DRTO$를
직접 계산하는 운영 도구로 활용되는 경로를 구체화한 것이다.}

\chg{도구의 토대는 전진 선택법과 LASSO의 교차검증 합의로 도출된 4변수 간명 모형(식~(\vref{eq:c2_parsimonious}),
(\vref{eq:c3_parsimonious}))이다. 9변수 완전 2차 모형이 C2에서 $R^2 = 0.998$, C3에서 $R^2 = 0.858$의
설명력을 보이는 데 비해, 합의 4변수 간명 모형은 C2 $R^2 = 0.996$, C3 $R^2 = 0.839$로 설명력 손실을
최소화하면서 해석 가능성을 확보한다. 특히 C1과 C2의 2차교호작용 회귀가 각각 $R^2 = 1.000$, $0.998$에
도달한다는 결과는 개별 시그모이드 바운딩 함수가 다항식으로 거의 완전히 근사됨을 의미하며, 이는 회귀식
산출 결과의 신뢰성을 뒷받침한다.}

\chg{실무 산출 절차는 다음과 같다. 조직은 자기 환경의 네 입력값, 곧 기준선 복구시간 $R_0$, 관측성 원시값
$O$, 관측성 가중치 $k_O$, 불확실성 수준 $U$를 간명 회귀식에 대입하여 효율비 추정값 $\hat\eta$를 얻고,
$\DRTO = \hat\eta \cdot R_0$로 동적 복구목표시간을 즉시 산출한다. 이렇게 산출된 $\DRTO$를 기존 BIA에서
고정값으로 설정한 정적 RTO와 비교하면, 관측성 투자와 불확실성 환경이 복구 목표에 미치는 영향을 정량적으로
확인할 수 있다. \tabref{tab:ch6-drto-examples}\refpo{tab:ch6-drto-examples}{은}{는} 뒤에 나올 산업별 권장
파라미터(\tabref{tab:ch6-industry-params})를 회귀식에 대입한 구체적 산출 예시이다.}

\begin{table}[htbp]
\centering
\caption{\chg{산업별 대표 $\DRTO$ 산출 예시}}
\label{tab:ch6-drto-examples}
\small
\begin{tabular}{l c c c c c}
\toprule
\textbf{산업} & \textbf{$R_0$ (h)} & \textbf{$O$} & \textbf{$\hat{\eta}$} & \textbf{$\DRTO$ (h)} & \textbf{조건} \\
\midrule
제조업   & 4.0 & 0.60 & $1.664$ & $6.66$ & C2 \\
금융     & 1.0 & 0.78 & $1.569$ & $1.57$ & C3 \\
의료     & 2.0 & 0.68 & $1.569$ & $3.14$ & C3 \\
IT 서비스 & 0.5 & 0.83 & $1.599$ & $0.80$ & C2 \\
물류     & 6.0 & 0.53 & $1.684$ & $10.10$ & C2 \\
소매     & 3.0 & 0.58 & $1.669$ & $5.01$ & C2 \\
\bottomrule
\end{tabular}
\end{table}

\chg{산출에는 4변수 간명 모형(C2 $R^2 = 0.996$, C3 $R^2 = 0.839$)을 적용하였으며, $U$는 산업별 불확실성
대표값(중앙값), $O_{\mathrm{raw}} = 0.5$를 가정하였다. 예컨대 금융 산업의 $\DRTO = 1.57$h는 $R_0 = 1.0$h
대비 $57\%$의 안전 마진을 포함한 것으로 파레토 불확실성에 의한 꼬리 위험을 반영한 값이며, 물류 산업의
$\DRTO = 10.10$h는 $R_0 = 6.0$h 대비 $68\%$ 증가한 것으로 비교적 낮은 관측성($O = 0.53$)과 불확실성
프리미엄이 결합된 결과이다. 이처럼 실무자는 자기 조직의 파라미터를 대입하여 $\DRTO$를 즉시 산출하고
기존 정적 RTO와 비교함으로써 동적 복구목표시간의 실무적 유용성을 직관적으로 확인할 수 있다.}

\chg{나아가 회귀식의 변수 진입 순서는 산출을 넘어 투자 우선순위 신호로도 활용된다. 4장의 전진 선택법
결과, 가우시안 환경(C2)에서는 관측성 가중치 $k_O$가 1순위로 진입하여 관측성이 $\eta$ 변동의 주된 결정
요인인 반면, 파레토 환경(C3)에서는 불확실성 $U$가 1순위로 진입하여 지배 변수가 전환된다. 따라서 조직은
자기 환경의 분포 특성에 따라, 가우시안 환경에서는 관측성 투자를, 파레토 환경에서는 불확실성 관리를
우선해야 한다는 자원 배분 신호를 회귀식으로부터 직접 읽어낼 수 있다.}

\chg{다만 간명 회귀식은 평균 영역의 거동을 근사하는 모형이므로, 극단 불확실성이 지배하는 꼬리 영역에서는
산출값의 정밀도가 제한된다. 꼬리 영역에서의 비선형 감쇠와 그에 따른 자원 배분은 다음 항(\subsecref{subsec:practical-dualchannel-tool})의
이중채널 도구에서 다룬다.}


% ===================================================================
% 5.3.2 이중채널 감쇠 기반 자원배분·스트레스테스트 [신설]
% ===================================================================
\subsection{\chg{이중채널 감쇠 기반 자원배분과 스트레스 테스트}}
\label{subsec:practical-dualchannel-tool}

\chg{본 항은 4장에서 발견한 이중채널 감쇠 메커니즘을 실무의 자원 배분 의사결정과 스트레스 테스트 설계
도구로 전환한다. 이중채널 감쇠란 P85.8 분할점을 경계로 관측성 가중치 $k_O$의 회귀 계수가 핵심 영역(Core,
$-0.798$)에서 꼬리 영역(Tail, $-0.415$)으로 $0.52$배 감쇠하는 현상으로, 극단 불확실성 환경에서 시그모이드
포화에 의해 관측성 투자의 한계 효용이 줄어든다는 것을 정량화한 본 연구의 핵심 발견이다.}

\chg{첫째, P85.8 분할점은 조직이 자기 불확실성 환경이 어느 체제에 있는지를 진단하고 대응 모드를 전환하는
기준이 된다. 가우시안(C2)과 파레토(C3) 분포의 누적분포함수가 P85.8($\eta \approx 2.42$)에서 교차하므로,
이 지점 이하의 일상 운영(Core)에서는 두 환경이 사실상 구별되지 않으나(Cohen's $d \approx 0$), 이 지점을
넘어서면 파레토 꼬리가 급격히 발현하여 $\mathrm{CVaR}_{99}$ 기준 C3가 C2를 $+24.0\%$ 상회한다. 따라서
조직의 복구 시간 분포가 P85.8 분할점을 넘어서기 시작하는 시점, 곧 가우시안 환경에서 파레토 환경으로의
전환은 그 자체로 자원 배분을 재조정해야 할 신호로 취급된다.}

\chg{둘째, 이중채널 감쇠 비율은 영역별 자원 배분 규칙의 정량적 근거를 제공한다. Core 영역에서는 관측성
가중치를 $0.1$ 높이면 $\eta$가 $0.080$ 감소하여 관측성 투자의 한계 수익이 충분히 발현되므로, 일상 운영의
복구 역량 향상은 관측성 투자(모니터링 확대, APM/SIEM 연동)에 집중하는 것이 효율적이다. 반면 Tail 영역에서는
동일한 관측성 개선이 $\eta$를 $0.042$만 감소시켜 한계 수익이 Core의 $0.52$배로 감쇠하므로, 극단 상황에서는
관측성 투자만으로 복구 시간을 충분히 단축하기 어렵다. 이 영역에서는 불확실성 자체의 관리, 곧
업무연속성계획(BCP) 시나리오의 다양화, 대체 사이트 확보, 사이버 보험과 같은 꼬리 위험 전가 수단에 자원을
전환해야 한다. 이는 동일한 예산이라도 환경 체제에 따라 투자 대상이 질적으로 달라져야 함을 의미한다.}

\chg{셋째, 이 발견은 스트레스 테스트 설계의 기준을 제시한다. $0.52$배 감쇠는 평균 영역의 시나리오만으로는
극단 환경에서의 복구 역량을 검증할 수 없음을 뜻하므로, BCM 스트레스 테스트는 Tail 영역의 시나리오를 별도로
설계해야 한다. 구체적으로, 파레토 분포의 P99 수준 극단값을 입력으로 하는 시나리오에서 관측성 투자만으로는
목표 복구 시간을 달성하지 못함을 검증하고, 불확실성 관리 수단을 병행했을 때의 효과를 함께 측정하는 이중
시나리오 설계가 요구된다.}

\chg{넷째, 환경 전환에 따른 회귀 계수 변화는 모델 재교정의 신호로 활용된다. 4장의 회귀계수 직접 비교 결과,
가우시안에서 파레토로 환경이 전환되면 관측성 변수 $k_O$의 계수가 $0.75$배, 불확실성 변수 $U$의 선형 계수가
$0.45$배로 변화한다. 또한 관측성 투자의 순효과는 가우시안 경로(Cohen's $d = +1.110$)보다 파레토 경로(Cohen's
$d = +0.650$)에서 상대적으로 작게 나타나, 환경에 따라 관측성 투자의 순효과 크기가 달라진다는 점도 확인된다.
조직은 자기 환경의 분포가 전환되는 것을 감지하면 회귀식의 계수를 재추정하여 산출 도구(\subsecref{subsec:practical-regression-tool})를
갱신해야 한다.}

\chg{요컨대 이중채널 감쇠 메커니즘은 ``일상에서는 관측성 투자, 극단에서는 불확실성 관리''라는 이중축 자원
배분 원칙과, Tail 시나리오를 별도로 설계하는 스트레스 테스트 기준을 동시에 제공한다. 이는 단일 지표나 평균
중심 의사결정으로는 포착할 수 없는, 분포 환경에 따른 입체적 위기 대응을 가능하게 하는 본 논문의 독자적
실무 기여이다.}


% ===================================================================
% 5.3.3 η 기반 평상시 성숙도·위기시 에스컬레이션 [구 성숙도 + 에스컬레이션 통합]
% ===================================================================
\subsection{\chg{$\eta$ 기반 평상시 성숙도와 위기시 에스컬레이션}}
\label{subsec:practical-bcm-std}        % 구 라벨 호환(성숙도)
\label{subsec:practical-decision}       % 구 라벨 호환(의사결정)
\label{subsec:escalation-thresholds}    % 구 라벨 호환
\label{subsec:eta-dual-use}             % 구 라벨 호환

\begin{p5add}
\chg{본 항은 단일 지표 $\eta = \DRTO/R_0$가 운영 모드에 따라 평상시 성숙도 진단 지표와 위기시 에스컬레이션
대응 지표의 두 역할을 수행함을 보이고, 그 평상시와 위기시 활용 및 선순환 구조를 제시한다. 평상시에는 $\eta$
변동이 주로 관측성 채널에 의해 결정되므로 BCMS 성숙도를 5단계 모형으로 진단하고, 위기시에는 $\eta$ 변동이
불확실성 채널에 의해 지배되므로 Green/Yellow/Orange/Red 4단계 에스컬레이션을 자동 트리거한다. 또한 데이터
수집과 모니터링 대시보드, 사후 분석으로 구성되는 BCM 운영 통합 도구가 함께 제공된다.}
\end{p5add}

\begingroup
  \renewcommand{\section}[1]{}
  \renewcommand{\subsection}[1]{}
  \renewcommand{\paragraph}[1]{\textbf{#1}\par}

  % η의 평상시·위기시 이중 활용 개념(overview) + 성숙도 5단계(1.1) + BCM 운영통합(1.2-4)
  
$\eta = \DRTO / R_0$는 단일 지표이지만, 조직의 운영 상황에 따라
두 가지 서로 다른 역할을 수행한다.
이러한 이중 활용이 가능한 근거는 $\eta$를 구성하는 두 채널\chg{(}관측성($E_{O,bnd}$)과
불확실성($E_{U,bnd}$)\chg{)}의 \textbf{지배적 기여 요인이 운영 모드에 따라 달라지기}
때문이다.
평상시에는 \chg{대규모 재난과 사고}가 부재하므로 불확실성이 \chg{일상적이고 가우시안}
수준에 머물러 $\eta$의 변동을 주로 관측성 역량이 결정한다.
반면 위기시에는 이미 \chg{사고와 재난}이 발생하여 규모를 알 수 없는 큰 불확실성이
현존하며, 이를 얼마나 빠르게 정량화하여 $\eta$에 반영하느냐가 대응의 핵심이 된다.
\p{[}표~\vref{tab:eta-dual-use}\p{]}는 이 \chg{평상시와 위기시} 활용 체계를 정리한 것이다.

\begin{table}[htbp]
\centering
\caption{$\eta$ 지표의 \chg{평상시·위기시} 활용 체계}
\label{tab:eta-dual-use}
\small
\begin{tabularx}{\textwidth}{l X X}
\toprule
\textbf{구분} & \textbf{평상시 (Normal Mode)} & \textbf{위기시 (Crisis Mode)} \\
\midrule
\textbf{역할} & BCMS 성숙도 평가 지표 & 에스컬레이션 대응 지표 \\
\textbf{핵심 질문} & ``조직의 복구 역량이 어느 수준인가?'' & ``현재 상황이 얼마나 심각한가?'' \\
\textbf{시간축} & 중장기 (월간/분기별 추이) & 실시간 (분$\sim$시간 단위) \\
\textbf{주요 활동} & 진단, 투자 계획, 벤치마크 & 등급 판정, 보고, 즉각 대응 \\
\textbf{대상자} & BCM 관리자, 경영진 & 복구 팀, 현장 담당자 \\
\textbf{산출물} & 성숙도 등급, 개선 로드맵 & 에스컬레이션 등급, 대응 조치 \\
\bottomrule
\end{tabularx}
\end{table}

\p{[}그림~\vref{fig:eta-dual-use-diagram}\p{]}는 $\eta$의 \chg{평상시와 위기시} 이중 활용과
BCMS 선순환 구조를 종합적으로 도식화한 것이다.
평상시 성숙도 평가(좌측)와 위기시 에스컬레이션 대응(우측)이
\chg{장애 발생과 복구 완료}를 통해 전환되며,
사후 분석과 역량 환류가 상단의 진단--향상--대응--환류 선순환을 형성한다.

\begin{figure}[htbp]
\centering
\resizebox{\textwidth}{!}{%
\begin{tikzpicture}[
  arr/.style={-{Stealth[length=2.5mm, width=2mm]}, line width=1.2pt},
  arrlbl/.style={font=\scriptsize\bfseries, fill=white, inner sep=1.5pt, rounded corners=1.5pt},
  sbox/.style={draw, rounded corners=2pt, semithick, minimum width=20mm,
    minimum height=9mm, align=center, font=\scriptsize},
]

%% ============================================================
%% 상단 박스: BCMS 선순환
%% ============================================================
\node[draw, rounded corners=6pt, line width=1pt,
  fill=green!6, draw=green!50!black,
  minimum width=120mm, minimum height=24mm] (topbox) at (0, 3.2) {};

\node[font=\small\bfseries, text=black] at (0, 4.15) {BCMS 선순환};

\node[sbox, fill=blue!12, draw=blue!50]  (s1) at (-3.7, 3.2) {\textbf{1. 진단}\\$\eta$ 추이};
\node[sbox, fill=blue!12, draw=blue!50]  (s2) at (-1.25, 3.2) {\textbf{2. 향상}\\$O\!\uparrow\; U\!\downarrow$};
\node[sbox, fill=red!12, draw=red!50]    (s3) at (1.25, 3.2)  {\textbf{3. 대응}\\에스컬레이션};
\node[sbox, fill=yellow!18, draw=yellow!60!black] (s4) at (3.7, 3.2) {\textbf{4. 환류}\\이력$\to$KPI};

\draw[-{Stealth[length=1.8mm]}, thin, blue!60] (s1) -- (s2);
\draw[-{Stealth[length=1.8mm]}, thin, red!50] (s2) -- (s3);
\draw[-{Stealth[length=1.8mm]}, thin, yellow!60!black] (s3) -- (s4);
\draw[-{Stealth[length=1.8mm]}, thin, green!50!black]
  (s4.south) -- ++(0,-0.28) -| (s1.south)
  node[pos=0.5, below, font=\scriptsize\bfseries, text=black,
    fill=green!6, inner sep=1.5pt, rounded corners=1.5pt] {$\eta$ 선순환};

%% ============================================================
%% 좌하 박스: 평상시 BCMS 성숙도
%% ============================================================
\node[draw, rounded corners=6pt, line width=1pt,
  fill=blue!6, draw=blue!60,
  minimum width=56mm, minimum height=38mm] (leftbox) at (-5.0, -1.0) {};

\node[font=\small\bfseries, text=black] at (-5.0, 0.55) {평상시: BCMS 성숙도 평가};

\node[anchor=center, inner sep=0pt] at (-5.0, -0.85) {%
  \scalebox{0.7}{%
  \begin{tabular}{c c c l}
  \toprule
  \textbf{단계} & \textbf{등급} & \textbf{$\eta$} & \textbf{수준} \\
  \midrule
  5 & 최적화 & $< 0.6$ & \textit{Optimized} \\
  4 & 관리 & $0.6$--$0.8$ & \textit{Managed} \\
  3 & 정의 & $0.8$--$1.0$ & \textit{Defined} \\
  2 & 반복 & $1.0$--$1.5$ & \textit{Repeatable} \\
  1 & 초기 & $\geq 1.5$ & \textit{Initial} \\
  \bottomrule
  \end{tabular}}%
};

\node[font=\scriptsize\bfseries, text=black] at (-5.0, -2.35) {%
  개선 두 축: 관측성$\uparrow$ / 불확실성$\downarrow$%
};

%% ============================================================
%% 우하 박스: 위기시 에스컬레이션
%% ============================================================
\node[draw, rounded corners=6pt, line width=1pt,
  fill=red!6, draw=red!60,
  minimum width=56mm, minimum height=38mm] (rightbox) at (5.0, -1.0) {};

\node[font=\small\bfseries, text=black] at (5.0, 0.55) {위기시: 에스컬레이션 대응};

\node[anchor=center, inner sep=0pt] at (5.0, -0.85) {%
  \scalebox{0.7}{%
  \begin{tabular}{c c l l}
  \toprule
  \textbf{$\eta$} & \textbf{등급} & \textbf{보고} & \textbf{대응} \\
  \midrule
  $< 0.8$ & Green & 팀 자체 & 정상 복구 \\
  $0.8$--$1.0$ & Yellow & 팀장 & 모니터링 강화 \\
  $1.0$--$1.5$ & Orange & 부서장 & 자원 추가 투입 \\
  $\geq 1.5$ & Red & 경영진 & BCP 전면 가동 \\
  \bottomrule
  \end{tabular}}%
};

\node[font=\scriptsize\bfseries, text=black] at (5.0, -2.35) {%
  국가 위기경보 연동: 관심--주의--경계--심각%
};

%% ============================================================
%% 3개 박스 연결 화살표
%% ============================================================
\draw[arr, red!50]
  ([yshift=-4mm]leftbox.east) -- ([yshift=-4mm]rightbox.west)
  node[arrlbl, midway, below=1.5pt, text=black] {장애$\cdot$재난 발생 $\to$ 위기 모드 전환};
\draw[arr, blue!50]
  ([yshift=4mm]rightbox.west) -- ([yshift=4mm]leftbox.east)
  node[arrlbl, midway, above=1.5pt, text=black] {복구 완료 $\to$ 평상 모드 복귀};
\draw[arr, green!50!black]
  ([xshift=6mm]rightbox.north) -- ([xshift=17mm]topbox.south)
  node[arrlbl, midway, right=1.5pt, text=black] {사후 분석};
\draw[arr, blue!60]
  ([xshift=-17mm]topbox.south) -- ([xshift=-6mm]leftbox.north)
  node[arrlbl, midway, left=1.5pt, text=black] {역량 환류};

%% ============================================================
%% 하단 강조 문구
%% ============================================================
\node[font=\small\bfseries, text=black, align=center,
  draw=red!50, rounded corners=3pt, fill=red!6, inner sep=4pt,
  line width=0.6pt] at (0, -3.5) {%
  $\eta$ 하나로 평상시 ``얼마나 준비되었는가'' + 위기시 ``얼마나 심각한가'' 통합 측정
  $\;\to\;$ \textcolor{black}{\bfseries 데이터 기반 BCMS 선순환}%
};

\end{tikzpicture}%
}% end resizebox
\caption{$\eta$의 \chg{평상시와 위기시} 이중 활용 체계와 BCMS 선순환 구조}
\label{fig:eta-dual-use-diagram}
\end{figure}

  %% --- 평상시: BCMS 성숙도 평가 ---
\subsubsection{\tmp{평상시 BCMS 성숙도 평가}}
\label{subsubsec:maturity-model}

평상시에는 \chg{대규모 사고와 재난}이 발생하지 않은 상태이므로,
불확실성은 일상적 업무 변동(가우시안 수준)에 한정되어 $E_{U,bnd}$의 기여가 작다.
따라서 $\eta$의 변동은 주로 \textbf{관측성 역량}($E_{O,bnd}$)에 의해 결정되며,
$\eta$의 중장기 추이를 추적하면 조직이 모니터링 체계를 얼마나 잘 구축하고
있는지\chg{,} 즉 복구 역량의 성숙도를 정량적으로 진단할 수 있다.
$\eta$가 낮을수록 관측성 역량이 우수한 성숙한 조직으로 해석되며,
이에 기반하여 BCMS 성숙도를 5단계로 구분하는 모형을 \p{[}표~\vref{tab:maturity-5level}\p{]}에서 제안한다.

각 임계값은 본 연구의 시뮬레이션 결과에 기반한다\chg{.}
$\eta = 0.6$은 C1 조건의 하위 \erf{25}{5}\% 수준으로 관측성이 극대화된 상태,
$\eta = 0.8$은 C1 평균 $\bar{\eta} = 0.853$ 이하로 $\DRTO < R_0$가 안정적으로 유지되는 상태,
$\eta = 1.0$은 $\DRTO = R_0$ 기준선으로 관측성과 불확실성이 상쇄되는 지점,
$\eta = 1.5$는 불확실성이 지배적이어서 $\DRTO$가 $R_0$의 1.5배를 초과하는 상태이다.

\begin{table}[htbp]
\centering
\caption{$\eta$ 기반 BCMS 성숙도 5단계 모형}
\label{tab:maturity-5level}
\small
\begin{tabularx}{\textwidth}{c c c l X}
\toprule
\textbf{단계} & \textbf{등급} & \textbf{$\eta$ 범위} & \textbf{수준} & \textbf{특징} \\
\midrule
5 & 최적화 & $\eta < 0.6$ & \textit{Optimized} &
  관측성 극대화, 불확실성 최소화. 자동화된 모니터링·복구 \\
\addlinespace
4 & 관리 & $0.6 \leq \eta < 0.8$ & \textit{Managed} &
  정량적 KPI 기반 BCM 운영. $\DRTO < R_0$ 안정 유지 \\
\addlinespace
3 & 정의 & $0.8 \leq \eta < 1.0$ & \textit{Defined} &
  공식 BCM 프로세스 수립. 관측성과 불확실성 \m{상쇄($\DRTO \approx R_0$)} \\
\addlinespace
2 & 반복 & $1.0 \leq \eta < 1.5$ & \textit{Repeatable} &
  기본 절차 존재하나 불확실성 우세. $\DRTO > R_0$ 빈번 \\
\addlinespace
1 & 초기 & $\eta \geq 1.5$ & \textit{Initial} &
  BCM 체계 미비. 극단적 불확실성 노출 \\
\bottomrule
\end{tabularx}
\end{table}

성숙도를 향상(즉, $\eta$를 감소)시키는 경로는 두 축으로 구분된다.
\chg{첫째, 관측성 강화($E_{O,bnd}$ 증가로 $\eta$ 감소)는 실시간 모니터링 체계 구축, APM/SIEM 연동, 장애
원인 분석(RCA) 체계화를 포함하며, 본 연구의 C1 조건에서 $k_O$ 계수가 관측성의 실질적 효과를 결정하므로
$k_O$가 높은 영역에 우선 투자하는 것이 효율적이다. 둘째, 불확실성 저감($E_{U,bnd}$ 감소로 $\eta$ 감소)은
복구 절차서 최신 유지, 정기 BCM 훈련, 공급망 대체 방안 확보를 포함하며, 특히 Tail 영역(P85.8 초과)에서는
파레토 불확실성이 지배적이므로 블랙스완 대비 계획이 성숙도 향상의 핵심이 된다.}

각 성숙도 단계별 구체적 전환 전략, 위기 시 정책 활용 방안, 도입 로드맵은
아래에 상세히 기술하며, 현업 담당자용 자가 진단 체크리스트(15개 항목)는
\appchref{app:eta-dual-use-detail}에 수록한다.

\input{chapters/ch5/_split/sec5.3.1.3_eta_detail}



  %% -------------------------------------------------------------------
\subsection{\vrev{BCM 운영 통합 프레임워크}}
\label{subsec:bcm-integration}
%% -------------------------------------------------------------------

2단계 프레임워크와 에스컬레이션 체계를 기존 BCMS에 통합하기 위한
구체적 운영 절차를 다음과 같이 제시한다.

\subsubsection{\tmp{데이터 수집 체계}}

\inew{$\DRTO$ 모델의 두 입력 변수인 \textbf{관측성 $O(t)$} 와 \textbf{불확실성 $U(t)$} 의
수집은 각각 기존 \chg{운영과 보안} 모니터링 인프라와 \chg{장애 이력 및 위협 인텔리전스} 데이터의
연동을 통해 구현한다.}

\inew{\chg{관측성 지표 $O(t)$의 수집은 기존 모니터링 인프라와의 연동을 통해 다음과 같이 구현한다. 첫째, APM(Application Performance Monitoring) 연동으로 서비스 응답 시간, 오류율, 처리량을 종합하여 $O_{\mathrm{APM}} \in [0,1]$로 정규화하며, 가중 평균은 $O_{\mathrm{APM}} = w_1 \cdot R_{\mathrm{avail}} + w_2 \cdot (1-R_{\mathrm{error}}) + w_3 \cdot R_{\mathrm{perf}}$($\sum w_i = 1$)로 산출한다. 둘째, SIEM(Security Information and Event Management) 연동으로 보안 이벤트 심각도와 대응 상태를 기반으로 $O_{\mathrm{SIEM}} \in [0,1]$을 산출하되, 미대응 고위험 이벤트가 존재하면 감점한다. 셋째, 종합 관측성은 $O(t) = \alpha \cdot O_{\mathrm{APM}} + (1-\alpha) \cdot O_{\mathrm{SIEM}}$로 정의하며, $\alpha$는 조직 특성에 따라 예를 들어 $[0.5, 0.8]$에서 설정한다.}}

\inew{\chg{불확실성 지표 $U(t)$의 수집은 조직 내부의 과거 장애 이력과 외부 위협 인텔리전스 데이터의 연동을 통해 다음과 같이 구현한다. 첫째, 장애 이력 DB 연동으로 조직 자체의 과거 장애와 복구 이력(MTTR, RTO 초과 사례, 복구 소요 시간 분포)을 적합 분포(가우시안 또는 파레토)로 모델링하여 $U_{\mathrm{hist}} \in [0,\; U_{\max}]$을 산출하며, 분포 형태 판정에는 Hill 추정량~\citep{hill1975simple} 또는 Anderson-Darling 검정~\citep{anderson1954test}을 적용하여 C2(경꼬리) 또는 C3(중꼬리) 환경으로 자동 분기한다. 둘째, 위협 인텔리전스(CVE/MITRE ATT\&CK) 연동으로 자산 대장과 매핑된 CVE 공개 DB, EPSS 점수, MITRE ATT\&CK 위협 액터 활동 정보를 기반으로 미래 위험 강도 $U_{\mathrm{threat}} \in [0,\; U_{\max}]$을 산출하되, 미패치 고위험 CVE 수와 EPSS 가중 평균을 결합한다. 셋째, 종합 불확실성은 $U(t) = \beta \cdot U_{\mathrm{hist}} + (1-\beta) \cdot U_{\mathrm{threat}}$로 정의하며, $\beta$는 조직의 과거 데이터 풍부도와 위협 환경 변동성에 따라 $[0.5, 0.8]$에서 설정한다. 장애 이력이 풍부한 성숙 조직은 $\beta$가 높고, 신생 조직과 신규 위협 환경은 $\beta$가 낮다.}}

\inew{원시 변수 $O_{\mathrm{raw}}$, $U_{\mathrm{raw}}$, 가중치 $k_O$의 제1계층
하위 지표 조작화 예시와 가상 기업 라이프사이클 적용 시나리오는
\secref{app:operationalization}(관측성 $O_{\mathrm{raw}}$ 조작화 예시,
불확실성 $U_{\mathrm{raw}}$ 조작화 예시)에 참고용으로 수록되어 있다.}

\subsubsection{\tmp{$\eta$ 모니터링 대시보드}}

$\eta$ 비율의 실시간 시각화를 통해 BCMS 효과성을 지속적으로 측정한다.
대시보드의 핵심 구성요소는 \chg{실시간 $\eta$ 추이 차트(시계열),
에스컬레이션 등급 표시(Green/Yellow/Orange/Red),
$O(t)$와 $U(t)$의 개별 추이, CVaR$_{99}$ 현재값}이며,
이는 \nrev{ISO 22301:2019} 조항~9.1(모니터링, 측정, 분석, 평가)의
정량적 요구사항을 직접 충족한다.

\subsubsection{\tmp{사후 분석 프로세스}}

복구 완료 후 24--72시간 이내에 다음의 사후 분석을 수행한다.
\chg{먼저 실제 복구 시간 $T_{\mathrm{actual}}$과 $\DRTO$ 예측값을 비교하고,
절대 오차 $E = T_{\mathrm{actual}} - \DRTO_{\mathrm{final}}$과
상대 오차 $E_r = E / T_{\mathrm{actual}}$를 산출하며,
과대 및 과소 추정 원인을 분석하여 파라미터 교정을 권고하고,
장애 이력 DB 갱신 및 $U(t)$ 분포 재추정을 수행한다.}
이 프로세스는 \nrev{ISO 22301:2019} 조항~10.1(지속적 개선)과 직접 대응된다.




  % 위기시 에스컬레이션(3.1) + 선순환(3.2)
  %% --- 위기시: 에스컬레이션 대응 체계 ---
\subsubsection{\syncr{위기 시 에스컬레이션 대응 체계}}
\label{subsubsec:escalation-detail}

위기시에는 \chg{사고와 재난}이 이미 발생한 상태이므로,
규모를 알 수 없는 큰 불확실성이 현존한다.
이때 $\eta$는 이 불확실성의 크기를 실시간으로 반영하는 위험 수준 지표가 되며,
조직이 불확실성을 얼마나 빠르게 정량화하여 $\eta$에 반영하느냐가
적시 대응의 관건이 된다.
$\eta$ 값에 따라 4단계 에스컬레이션 체계가 자동 결정된다.
\p{[}표~\vref{tab:ch6-eta-escalation}\p{]}는 $\eta$ 구간별 대응 체계를 정의하며,
\p{[}그림~\vref{fig:ch6-escalation-zones}\p{]}는 이를 시각화한 것이다.

\begin{table}[htbp]
\centering
\caption{$\eta$ 기반 에스컬레이션 임계값 및 대응 체계}
\label{tab:ch6-eta-escalation}
\small
\begin{tabularx}{\textwidth}{c c l l X}
\toprule
\textbf{단계} & \textbf{$\eta$ 범위} & \textbf{등급} & \textbf{보고 대상} & \textbf{대응 조치} \\
\midrule
1 & $\eta < 0.8$ & \m{안정(Green)} & 복구 팀 자체 & 정상 복구 절차 수행; 관측성에 의한 $\DRTO$ 단축 효과 활성 \\
\addlinespace
2 & $0.8 \leq \eta < 1.0$ & \m{주의(Yellow)} & 팀장 보고 & 불확실성 증가 모니터링; $\DRTO$ 갱신 주기 \m{단축(15분)}; 추가 자원 대기 \\
\addlinespace
3 & $1.0 \leq \eta < 1.5$ & \m{경고(Orange)} & 부서장 보고 & $\DRTO > R_0$로 기준 RTO 초과; 추가 자원 투입; 외부 지원 요청; 이해관계자 긴급 통보 \\
\addlinespace
4 & $\eta \geq 1.5$ & \m{위기(Red)} & 경영진 보고 & 위기관리 모드 전환; 사업 연속성 비상 계획(BCP) 전면 가동; CVaR$_{99}$ 기준 최악 시나리오 대비 \\
\bottomrule
\end{tabularx}
\end{table}

\chg{에스컬레이션 트리거는 $\eta$ 급등(30\% 이상 증가), 조건 상향 전환(C2 $\to$ C3), 에스컬레이션 임계값
접근($\eta > 1.4$), 불확실성 급등($U > 0.8$)의 네 가지이다.}

\begin{figure}[htbp]
\centering
\resizebox{0.95\textwidth}{!}{%
\begin{tikzpicture}
% 배경 영역
\fill[green!15]  (0,0)   rectangle (3.2,3.5);
\fill[yellow!20] (3.2,0) rectangle (6.4,3.5);
\fill[orange!20] (6.4,0) rectangle (11.2,3.5);
\fill[red!20]    (11.2,0) rectangle (14.4,3.5);

% 경계선
\draw[thick, black] (3.2,0) -- (3.2,3.5);
\draw[thick, black] (6.4,0) -- (6.4,3.5);
\draw[thick, black] (11.2,0) -- (11.2,3.5);

% 축
\draw[thick, -{Stealth[length=3mm]}] (0,0) -- (15,0) node[right, font=\small] {$\eta$};
\draw[thick] (0,0) -- (0,3.5);

% η 값 레이블
\node[below, font=\footnotesize] at (0, -0.15) {$0$};
\node[below, font=\footnotesize\bfseries] at (3.2, -0.15) {$0.8$};
\node[below, font=\footnotesize\bfseries] at (6.4, -0.15) {$1.0$};
\node[below, font=\footnotesize\bfseries] at (11.2, -0.15) {$1.5$};

% 영역 레이블 (상단)
\node[font=\small\bfseries, green!50!black] at (1.6, 3.0) {안정};
\node[font=\small\bfseries, yellow!60!black] at (4.8, 3.0) {주의};
\node[font=\small\bfseries, orange!60!black] at (8.8, 3.0) {경고};
\node[font=\small\bfseries, red!60!black] at (12.8, 3.0) {위기};

% 영문 등급
\node[font=\scriptsize, green!40!black] at (1.6, 2.5) {Green};
\node[font=\scriptsize, yellow!50!black] at (4.8, 2.5) {Yellow};
\node[font=\scriptsize, orange!50!black] at (8.8, 2.5) {Orange};
\node[font=\scriptsize, red!50!black] at (12.8, 2.5) {Red};

% 핵심 조건 설명
\node[font=\scriptsize, align=center] at (1.6, 1.5)
  {관측성에 의한\\$\DRTO$ 단축\\$E_{O,bnd}$ 지배};
\node[font=\scriptsize, align=center] at (4.8, 1.5)
  {불확실성 증가\\관측성 효과 상쇄\\$E_{O,bnd} + E_{U,bnd} \approx 0$};
\node[font=\scriptsize, align=center] at (8.8, 1.5)
  {$\DRTO > R_0$\\불확실성 프리미엄\\C2 기반 자원 추가};
\node[font=\scriptsize, align=center] at (12.8, 1.5)
  {파레토 활성화\\$0.52$배 감쇠\\BCP 전면 가동};

% R₀ 기준선 표시
\draw[dashed, blue!60, thick] (6.4, -0.5) -- (6.4, -0.9)
    node[below, font=\scriptsize, blue!60] {$\eta = 1.0$: $\DRTO = R_0$ 기준};

\end{tikzpicture}
}% end resizebox
\caption{$\eta$ 기반 에스컬레이션 임계값 \m{다이어그램(Green / Yellow / Orange / Red)}}
\label{fig:ch6-escalation-zones}
\end{figure}

에스컬레이션 임계값의 설정 근거는 다음과 같다.
$\eta = 0.8$은 C1 조건에서의 평균 $\eta = 0.853$보다 낮은 값으로,
관측성의 단축 효과가 충분히 발현되어 추가 대응이 불필요한 안정 상태를 의미한다.
$\eta = 1.0$은 $\DRTO = R_0$, 즉 DRTO가 기준 RTO와 동일한 지점으로,
이 이상에서는 불확실성 프리미엄에 의해 복구 시간이 기준을 초과한다.
$\eta = 1.5$는 $\DRTO = 1.5 \times R_0 = 9.0$\,h에 대응하며,
이 이상에서는 시그모이드 바운딩의 포화 영역에 진입하여
위기관리 모드 전환이 필요하다.
이러한 단계별 에스컬레이션은 SRE의 에러 예산 소진율 기반
인시던트 대응 체계~\citep{beyer2016sre}와 유사한 구조를 가지나,
$\eta$가 관측성과 불확실성의 결합 효과를 단일 지표로 포착한다는 점에서
BCM 맥락에 최적화된 확장이다.

\chg{나아가 2단계 프레임워크와의 복합 에스컬레이션이 가능하다.}
$\eta$ 기반 4단계 등급은 \subsubsecref{subsec:two-stage-framework}의
2단계 프레임워크(Core C2 + Tail C3)와 결합하여 더욱 정교한 대응이 가능하다.
동일한 $\eta$ 등급이라도 불확실성 분포가 가우시안(C2)인지
파레토(C3)인지에 따라 위험의 성격이 질적으로 다르기 때문이다.
예를 들어, $\eta = 1.3$(Orange 등급)이 Core 영역(C2)에서 관측되면
가우시안 불확실성의 일시적 증가로 해석되어 추가 자원 투입과 모니터링 강화가
1차 대응이 된다. 그러나 동일 $\eta$가 Tail 영역(C3)에서 관측되면
파레토 꼬리 위험이 활성화된 상태이므로, $k_O$ 계수의 $0.52$배 감쇠로 인해
관측성 효과가 제한적이며, CVaR$_{99}$가 $+24.0\%$ 높은 극단 시나리오에
대비한 선제적 BCP 점검이 필요하다.

이를 운영 규칙으로 정리하면 다음과 같다.
\chg{먼저 Core 영역(P85.8 이하)에서는 $\eta$ 등급에 따른 표준 에스컬레이션 절차를 적용하고,
Tail 영역(P85.8 초과)에서는 동일 $\eta$ 등급이라도 한 단계 높은 대응 수준을
적용하며(예: Tail에서의 Orange $\to$ Red 수준 대응),
Core에서 Tail로의 조건 전환(C2 $\to$ C3) 자체를
에스컬레이션 트리거로 취급한다.}
이러한 복합 에스컬레이션은 $\eta$라는 단일 축의 한계를 보완하여,
불확실성의 분포 특성까지 반영한 입체적 위기 대응을 가능하게 한다.

\chg{또한 국가 위기경보 체계와도 대응된다.}
본 연구의 $\eta$ 기반 에스컬레이션 4단계는 ``재난 및 안전관리 기본법''에 근거한
\textbf{국가 위기경보 4단계}(관심--주의--경계--심각)와 구조적으로 대응된다.
\p{[}표~\vref{tab:eta-crisis-alert-mapping}\p{]}은 양 체계의 매핑을 보여준다.

\begin{table}[htbp]
\centering
\caption{$\eta$ 기반 에스컬레이션과 국가 위기경보 4단계 대응}
\label{tab:eta-crisis-alert-mapping}
\small
\begin{tabularx}{\textwidth}{c c c c c X}
\toprule
\textbf{$\eta$ 범위} & \textbf{DRTO 등급} & \textbf{색상} & \textbf{위기경보} & \textbf{색상} & \textbf{대응 원칙의 공통점} \\
\midrule
$\eta < 0.8$ & 안정 & Green & 관심 & 파랑 &
  징후 감지 수준; 정상 운영 유지, 모니터링 강화 \\
\addlinespace
$0.8 \leq \eta < 1.0$ & 주의 & Yellow & 주의 & 노랑 &
  위험 가능성 증가; 자원 점검, 갱신 주기 단축 \\
\addlinespace
$1.0 \leq \eta < 1.5$ & 경고 & Orange & 경계 & 주황 &
  기준 초과 확인; 추가 자원 투입, 이해관계자 통보 \\
\addlinespace
$\eta \geq 1.5$ & 위기 & Red & 심각 & 빨강 &
  총력 대응; BCP 전면 가동, 경영진 의사결정 \\
\bottomrule
\end{tabularx}
\end{table}

두 체계는 \chg{위험 강도에 따른 단계적 격상(graduated response),
색상코드 기반 직관적 커뮤니케이션,
각 단계별 보고 대상 및 대응 수준의 명확한 구분}이라는
공통 설계 원칙을 공유한다.
차이점은 국가 위기경보가 질적 판단에 기반하는 반면,
$\eta$ 기반 에스컬레이션은 $\DRTO / R_0$라는 \textbf{정량적 지표}에
의해 자동으로 등급이 결정된다는 점이다.
이러한 정량적 트리거는 주관적 판단 편향을 배제하고
일관된 의사결정을 지원하며, 기존 위기경보 체계에 보완적으로
통합될 수 있다.


  
%% --- 선순환 구조 ---
\subsubsection{\syncr{진단에서 향상, 대응, 환류로 이어지는 선순환 구조}}
\label{subsubsec:virtuous-cycle}

$\eta$의 \chg{평상시와 위기시} 활용이 만드는 핵심 가치는 \textbf{선순환}이다.
평상시의 성숙도 진단과 역량 향상이 위기시의 대응력을 높이고,
위기시의 에스컬레이션 이력\chg{(}Green 체류 비율,
복귀 시간, Red 도달률 등\chg{)}이 평상시의 성숙도 개선 데이터로 환류된다.
이 순환의 공통 언어가 $\eta$이며,
BCMS가 ``문서 기반 형식적 체계''에서 ``데이터 기반 실행적 체계''로
전환되는 것이 $\DRTO$ 프레임워크가 실무에 제공하는 핵심 가치이다.


\endgroup

\subsubsection{\syncr{실무 적용 가능성 예시}}
\label{subsubsec:practical-application-example}

\syncr{지금까지 제시한 $\eta$ 기반 \chg{평상시와 위기시} 이중 활용 체계가 실제 조직에
적용될 경우, 관측성($O_{\text{raw}}, k_O$)과 불확실성($U$)은 APM/SIEM 로그,
장애 이력 DB, CVE 공개 DB, 자산 대장 등 현장의 기존 운영 시스템에서
직접 관측되는 원시 지표로부터 산출 가능할 것이다.
참고를 위해 아래 \subsecref{app:operationalization}에 제1계층 하위 지표의
조작화 예시와 가상 기업의 4단계 라이프사이클
(진단→향상→대응→환류) 적용 가능 시나리오를 제시한다.
본 예시는 실무 적용 방향성을 보여주기 위한 참고용이며, 하위 지표의
구체적 \chg{구성과 가중치, 측정 도구}의 실증적 검증은 본 연구의 범위를 벗어나는
후속 연구 과제로 남겨둔다(\secref{sec:limitations} 참조).}



%% [본문 이관] 구 부록 J: 하위 변수 조작화 및 가상 기업 적용 예시 (참고용)
\subsection{\syncr{하위 변수 조작화 및 가상 기업 적용 예시}}
\label{app:operationalization}

\syncr{\subsubsection{하위 변수 조작화의 성격과 한계}}

\syncr{\textbf{하위 변수 조작화는 $\DRTO$ 프레임워크의 실무 적용 방향성을 보여주기 위한
참고용 예시이다.} 제시된 하위 지표의 구성, 가중치, 측정 도구 매핑,
가상 기업 수치는 \textbf{실증적으로 검증되지 않았으며}, 중견 제조업 일반
특성에서 추정한 \textbf{조건적 시나리오}(illustrative case~\citep{yin2018case})에
해당한다. 이를 규범적 가이드라인 또는 검증된 방법론으로 해석하지 말 것을 당부한다.}

\syncr{본문에서 제시한 $\DRTO$ \chg{모델과 시뮬레이션 결과, 전문가 검증}은 본 연구의
공식적 기여이나, 하위 변수 조작화와 시나리오는 \textbf{향후 연구를 위한
탐색적 제안}에 그친다. 각 하위 지표의 심리측정학적 타당도, 가중치의
산업별 튜닝, 원시 지표의 측정 도구별 \chg{정확도와 편향}, 현장 데이터 기반
\chg{수렴 및 예측} 타당도는 본 연구의 범위를 벗어나며 후속 연구에서 독립적으로
검증되어야 한다.}

\syncr{\subsubsection{조작화의 필요성과 정지점 원칙}}

\syncr{$\DRTO$ 모델의 주 변수인 관측성($O_{\text{raw}}$, $k_O$)과 불확실성($U$)은
수식적으로는 정의되어 있으나, 실제 조직에서 이들을 어떻게 산출할 것인가는
추가 조작적 정의가 필요하다. 조작화는 원칙적으로 하위 변수의 하위 변수로
무한 분해될 수 있으므로, 어느 수준에서 정지할 것인가의 \textbf{정지점
(stopping point)} 문제가 발생한다.}

\syncr{하위 변수 조작화는 측정이론의 고전적 원칙을 따라, 해당 하위 지표가 현장의
\textbf{기존 운영 시스템}(APM 도구·SIEM 도구·장애 이력 데이터베이스·CVE 공개 데이터베이스·자산 대장 등)\textbf{에서 직접 \chg{관측되고 기록되는} 원시값}에
도달하면 그 이하의 분해는 본 연구의 범위 밖으로 간주한다. 이는
\synct{bagozzi1998general}의 ``observable primitive'' 정지 규칙,
\synct{bollen1989structural}의 effect indicator 이론, 그리고
\nrev{ISO/IEC 25010:2011}의 원자 지표 표준화 원칙에 부합한다~\citep{iso25010}.}

\syncr{\chg{다만} 본 정지 규칙은 일반적 원칙을 따른 것이며,
개별 원시 지표의 심리측정학적 타당도는 측정 도구별로 별도 검증이 필요하다.}

\syncr{\subsubsection{관측성 $O_{\text{raw}}$ 조작화 예시}}

\syncr{관측성 원시값 $O_{\text{raw}} \in [0,1]$는 응용 성능 모니터링(APM)과 보안 \chg{정보 및 이벤트}
관리(SIEM)의 두 차원으로 분해 가능하다.}

\syncr{\chg{APM 차원은 $O_{\text{APM}} = w_1 R_{\text{avail}} + w_2 (1-R_{\text{error}}) + w_3 R_{\text{perf}}$로 정의한다.}}

\syncr{\chg{SIEM 차원은 $O_{\text{SIEM}} = v_1 (1-E_{\text{high}}) + v_2 (1-T_{\text{mttr}})$로 정의한다.}}

\syncr{\chg{두 차원의 합성은 $O_{\text{raw}} = \alpha \cdot O_{\text{APM}} + (1-\alpha) \cdot O_{\text{SIEM}}$이며, 본 예시에서 $\alpha = 0.7$로 가정한다.}}

\syncr{제1계층 원시 지표는 \tabref{tab:app-operationalization}에 정리하였다.}

\begin{table}[htbp]
\centering
\caption{\syncr{\chg{관측성과 불확실성, $k_O$}의 1차 조작화 \m{예시(참고용)}}}
\label{tab:app-operationalization}
\small
\syncr{\begin{tabularx}{\textwidth}{l l X l}
\toprule
\textbf{주 변수} & \textbf{1차 차원} & \textbf{원시 \m{지표(정지점)}} & \textbf{측정 출처} \\
\midrule
\multirow{5}{*}{$O_{\text{raw}}$} &
\multirow{3}{*}{$O_{\text{APM}}$} &
  $R_{\text{avail}}$: 가용시간/총시간 & APM 도구 로그 \\
& & $R_{\text{error}}$: 오류건수/총요청수 & APM 도구 로그 \\
& & $R_{\text{perf}}$: 목표응답/평균응답 & APM 도구 로그 \\
\cmidrule(l){2-4}
& \multirow{2}{*}{$O_{\text{SIEM}}$} &
  $E_{\text{high}}$: 미대응 고위험 이벤트 비율 & SIEM 도구 \\
& & $T_{\text{mttr}}$: 평균 \erf{탐지시간}{복구시간}(정규화) & SIEM 도구 \\
\midrule
\multirow{2}{*}{\jrev{$U_{\text{raw}}$}} &
  $U_{\text{hist}}$ & $\ln(P99/\jrev{\overline{T}})$ of 복구시간 & 장애 이력 DB \\
& $U_{\text{ext}}$ & \jrev{CVE CVSS 평균} & NVD/MITRE ATT\&CK \\
\midrule
$k_O$ & (단일) & 모니터링 커버리지율 & 자산 대장 \\
\bottomrule
\end{tabularx}}
\end{table}

\syncr{\chg{다만} 가중치 $w_i$, $v_i$, $\alpha$는 본 예시에서
등가중 또는 0.7로 임의 가정하였으며, \chg{산업별과 조직별} 최적값은 향후 연구에서
실증 검증되어야 한다.}

\syncr{\subsubsection{불확실성 \imod{$U$}{$U_{\text{raw}}$} 조작화 예시}}

\syncr{불확실성 \jrev{$U_{\text{raw}}$}는 과거 장애 이력 기반 $U_{\text{hist}}$와 외부 위협 지수 기반
$U_{\text{ext}}$의 두 차원으로 분해 가능하다\jrev{(두 하위 지표 모두 $[0, \infty)$
범위로 본 연구 모델의 $U_{\text{raw}}$ 정의와 일치한다)}.}

\syncr{\chg{과거 장애 기반 지표는 $U_{\text{hist}} = \ln\!\left(P99_{\text{recover}} / \overline{T}_{\text{recover}}\right) \jrev{\in [0, \infty)}$로, 복구시간 분포의 꼬리 두께를 로그 스케일로 정량화하며 장애 이력 DB에서 직접 산출한다.}}

\syncr{\chg{외부 위협 기반 지표는 $U_{\text{ext}} = \text{CVSS}_{\text{avg}} \jrev{\in [0, \infty)}$로, 최근 12개월 관련 CVE의 CVSS 점수 평균이며 NVD/MITRE ATT\&CK 공개 DB에서 산출한다.} \jrev{표준 CVSS v3.1은 $[0, 10]$이지만 조직 내부 위협 가중치 적용 시 10 초과 가능하다.}}

\syncr{\chg{두 차원의 합성은 \jrev{$U_{\text{raw}}$} $= \beta \cdot U_{\text{hist}} + (1-\beta) \cdot U_{\text{ext}}$이며, 본 예시에서 $\beta = 0.6$으로 가정한다.}}

\syncr{\chg{다만} 조직별 이력 DB 구축 수준이 상이하므로 $U_{\text{hist}}$
산출의 신뢰도는 조직마다 다를 수 있다. 신설 조직 또는 이력 부족 조직의
경우 $\beta$를 낮추거나 업종 평균을 사용하는 보정이 필요하다. \jrev{또한
$U_{\text{hist}}$와 $U_{\text{ext}}$는 각각 독립적으로 가우시안 또는 파레토
분포를 따를 수 있으며(조합 시 $2 \times 2 = 4$가지 분포 시나리오 가능),
본 연구 조건 C2/C3는 양 극단(전량 가우시안 / 전량 파레토)만 다룬다.
혼합 분포 시나리오의 실증은 향후 연구 과제이다.}}

\syncr{\subsubsection{관측성 가중치 $k_O$ 조작화 예시}}

\syncr{관측성 가중치 $k_O \in [0,1]$는 조직이 관측성에 얼마나 투자했는지를
나타내는 감도 계수로, 본 예시에서는 \textbf{모니터링 커버리지율}로 다음과 같이 단순화한다\chg{.}}

\syncr{$k_O = \frac{\text{감시 대상 자산 수}}{\text{전체 핵심 자산 수}}$}

\syncr{자산 대장에서 직접 산출되는 원시값이며, 감시 대상 자산은 APM/SIEM 도구에
등록된 엔드포인트 또는 프로세스를 기준으로 한다.}

\syncr{\chg{다만} 커버리지율은 ``감시의 폭''만을 측정하며, ``감시의
깊이''(로그 수집 상세도, 알람 정확도)는 별도 지표로 보완되어야 한다.
본 예시는 최소 구성만 제시한다.}

\syncr{\subsubsection{가상 기업 라이프사이클 적용 시나리오}}

\syncr{``(주)한빛정밀''(가상 기업명, 동명의 실제 기업과 무관)은
수도권 산업단지에 소재한 중견 제조업체(반도체 부품 정밀가공, 직원 420명,
연매출 1,200억 원 규모)로 가정하여 다음과 같이 시나리오를 작성한다. \nrev{ISO 22301:2019} 인증 준비 단계이며,
APM(Datadog)과 SIEM(Splunk)을 \chg{구축하여 운영한다}.}

\syncr{\chg{1단계인 평상시 진단(T-365일)에서는} 연간 BCMS 성숙도 평가 시점의 원시 지표를 다음과 같이
$R_{\text{avail}} = 0.994$, $R_{\text{error}} = 0.015$, $R_{\text{perf}} = 0.82$,
$E_{\text{high}} = 0.08$, $T_{\text{mttr}} = 0.65$, 커버리지율 $= 0.72$로 산출하였다.}

\syncr{합성 결과는 다음과 같다\chg{.} $O_{\text{APM}} = 0.933$, $O_{\text{SIEM}} = 0.785$,
$O_{\text{raw}} = 0.889$, $k_O = 0.72$, $U = 0.55$로 산출되어
$\eta \approx 1.05$ (Level 2 Repeatable) 판정된다. 모니터링 커버리지 확대 및
미대응 이벤트 축소가 개선 필요 영역으로 도출된다.}

\syncr{\chg{2단계인 평상시 개선 활동(T-365일부터 T-1일)에서는} 분기별 복구 훈련 도입, 모니터링 자산 확대(커버리지율 0.72 → 0.88),
SIEM 튜닝($E_{\text{high}}$ 0.08 → 0.03), 장애 이력 DB 정비를 수행하였다.
금년 재평가\chg{에서} $\eta \approx 0.82$\chg{를 달성하여} \textbf{\erf{Level 4 Managed}{Level 3 Defined}}로 \chg{향상되었다}.
월별 $\eta$ 추이를 경영진 대시보드에 정착한다.}

\syncr{\chg{3단계인 위기 발생 및 에스컬레이션(T+0h--T+6h)의} 가상 시나리오는 생산라인 제어 PLC 장애와 보조 전원계통 이상이 발생하였고
(인명 피해 없음, 설비 손상만 발생), 시간대별 $\eta$ 추이와 등급 변화는
\tabref{tab:app-lifecycle-eta}에 정리한다.}

\begin{table}[htbp]
\centering
\caption{\syncr{가상 기업 4단계 라이프사이클 $\eta$ \m{추이(참고 시나리오)}}}
\label{tab:app-lifecycle-eta}
\small
\syncr{\begin{tabularx}{\textwidth}{c l X c c}
\toprule
\textbf{시각} & \textbf{단계} & \textbf{상황} & \textbf{$\eta$} & \textbf{등급} \\
\midrule
$T{-}365$d & 1단계 진단 & 초기 \m{상태(Level 2)} & 1.05 & \erf{Yellow}{Orange} \\
$T{-}1$d & 2단계 개선 후 & 모니터링 확대 + SIEM \m{튜닝(\erf{Level 4}{Level 3})} & 0.82 & \erf{Green}{Yellow} \\
\midrule
$T{+}0$h & 3단계 정상 & 사건 직전 & 0.82 & \erf{Green}{Yellow} \\
$T{+}1$h & 3단계 탐지 & PLC 장애 탐지 & 1.15 & Orange \\
$T{+}2$h & 3단계 확산 & 전원계통 이상 확산 & 1.48 & Orange \\
$T{+}3$h & 3단계 정점 & 복구 불확실성 급등 & 1.62 & \textbf{Red} \\
$T{+}4$h & 3단계 전환 & 외부 지원 투입 & 1.45 & Orange \\
$T{+}6$h & 3단계 부분복구 & 주요 지표 안정화 & 1.08 & \erf{Yellow}{Orange} \\
\midrule
$T{+}72$h & 4단계 환류 & 사후 분석 후 재평가 & 0.84 & \erf{Green}{Yellow} \\
\bottomrule
\end{tabularx}}
\end{table}

\syncr{에스컬레이션\chg{은} $T{+}1$h(Orange)\chg{에} \chg{복구팀과 부서장} 자동 통지\chg{로 발동되어,}
$T{+}3$h(Red) 경영진 \chg{긴급 보고와 BCP 전면 가동을 거쳐,} $T{+}4$h 국가 위기경보 체계
연동 검토로 단계화되며 5.3.3.3(위기 시 정책 활용 방안)에서 상세 내용을 참조한다.}

\syncr{\chg{4단계인 복구와 환류(T+24h--T+72h)의} 사후 분석에서는 실제 복구 시간 $T_{\text{actual}} = 8.2$h vs $\DRTO$ 예측 7.6h
→ 상대 오차 $E_r = 7.3\%$ (수용 범위)이다. $O(t)$ 사각지대로 전원계통 센서
커버리지 부족을 식별하고, 성숙도 ``Managed → Optimized'' 이행 로드맵에
반영한다. 이번 사건을 $U_{\text{hist}}$ DB에 환류하여 내년도 진단의 입력으로 사용한다.}

\syncr{\chg{다만} 본 시나리오의 모든 수치는 중견 제조업 일반 특성에서
추정한 \textbf{조건적 예시}이며, 실제 기업의 실적 데이터가 아니다.
개별 조직의 적용 시 파라미터 재산출과 민감도 분석이 필수적이다.}

\syncr{\subsubsection{원시 지표별 측정 도구 매핑 예시}}

\syncr{\chg{앞의 관측성 $O_{\text{raw}}$, 불확실성 $U_{\text{raw}}$, 가중치 $k_O$ 조작화 예시에서 정의한} 원시 지표별 대표 측정 도구 매핑은
\tabref{tab:app-tool-mapping}에 정리하였다. 본 매핑은 \textbf{참고용}이며,
특정 도구의 추천이 아니다.}

\begin{table}[htbp]
\centering
\caption{\syncr{원시 지표별 측정 도구 매핑 \m{예시(참고용, 도구 권장 아님)}}}
\label{tab:app-tool-mapping}
\small
\syncr{\begin{tabularx}{\textwidth}{l l X}
\toprule
\textbf{원시 지표} & \textbf{측정 도구 유형} & \textbf{대표 \m{도구(참고)}} \\
\midrule
$R_{\text{avail}}$, $R_{\text{error}}$, $R_{\text{perf}}$ & APM & Datadog, New Relic, Prometheus+Grafana, Dynatrace \\
\addlinespace
$E_{\text{high}}$, $T_{\text{mttr}}$ & SIEM & Splunk, IBM QRadar, ArcSight, Elastic SIEM \\
\addlinespace
$U_{\text{hist}}$ & 장애 이력 DB & ServiceNow IRM, Jira Service Management, 자체 ITSM \\
\addlinespace
$U_{\text{ext}}$ & 위협 인텔리전스 & NVD, MITRE ATT\&CK, Recorded Future, CrowdStrike \\
\addlinespace
$k_O$ & 자산 대장/CMDB & ServiceNow CMDB, Device42, Lansweeper \\
\bottomrule
\end{tabularx}}
\end{table}

\nrev{ISO/IEC 25010:2011} \syncr{품질 모델의 Reliability(신뢰성),
Performance Efficiency(성능), Security(보안) 하위 속성이 각 원시 지표와
자연스럽게 대응되므로, 국제 표준 품질 지표 프레임워크와의 연계가 가능하다.}

\syncr{\chg{다만} 상기 도구는 예시이며, 조직의 기존 \chg{인프라와 예산, 규제}
환경에 따라 대체 선택이 필요하다.}

\syncr{\subsubsection{한계 종합 및 향후 실증 과제}}

\syncr{하위 변수 조작화와 시나리오는 참고용 예시이며, 다음 네 영역에서
후속 실증 연구가 필요하다.}

\syncr{\chg{첫째, 가중치 실증 튜닝이다.} $w_i$, $v_i$, $\alpha$, $\beta$를 \chg{산업별·조직 규모별}로 조정하는 다중 회귀 또는 최적화 연구가 필요하며,
본 예시는 등가중 또는 임의 고정값을 사용하였다.}

\syncr{\chg{둘째, 측정 도구 신뢰도이다.} 각 원시 지표의 측정 도구별 \chg{정확도와 편향}
검증이 필요하며, APM/SIEM 도구별 측정 알고리즘이 상이하므로
교차 타당도(cross-validity) 연구가 선결 과제이다.}

\syncr{\chg{셋째, 내용 타당도 재해석이다.} 본 연구의 \jrev{H9} 전문가 검증
(Section C 파라미터 적정성 6문항, S-CVI/Ave $= 0.985$)은 조작화의 내용
타당도의 간접 근거로 재해석 가능하나, 조작화 구성 자체에 대한 독립적
내용 타당도 검증은 별도 연구가 필요하다.}

\syncr{\chg{넷째, 수렴 및 예측 타당도이다.} 실제 조직 장애 이력과 $\eta$ 예측값의
상관 분석(종단 연구)을 통해 \chg{수렴, 판별, 예측} 타당도를 확보하는 실증 연구가
필요하다. 이는 본 연구 5.4.1 ``시뮬레이션 기반 검증의 한계''와 5.4.4
``시간 동역학의 부재''에서 제시한 향후 연구 방향과 연결된다.}

\syncr{종합적으로 하위 변수 조작화는 $\DRTO$ 프레임워크의 실무 적용 \textbf{방향성}을
제시하였으나, 본격 실무 배포 이전에는 상기 네 영역의 실증 연구가 선결되어야
한다.}



\FloatBarrier


% ===================================================================
% 5.3.4 산업별 적용 가이드라인 [구 5.3.2]
% ===================================================================
\subsection{\chg{산업별 적용 가이드라인}}
\label{subsec:practical-industry}

\begin{p5add}
\chg{본 항은 $\DRTO$ 모델의 산업별 적용 매트릭스를 제공한다. 먼저 P85.8 교차점에 기반한 2단계
프레임워크(Core C2 + Tail C3)로 일상 운영과 극단 시나리오의 분리 운용 원칙을 정립하고, 이어 6개 대표
산업(제조, 금융, 의료, IT 서비스, 물류, 소매)에 대한 $R_0$/$R_{\max}$/$\kappa$ 설정값과 적용 사례를
매트릭스 형태로 제시한다. 산업별 권장은 앞 항들의 회귀식 산출(\subsecref{subsec:practical-regression-tool})과
이중채널 자원 배분(\subsecref{subsec:practical-dualchannel-tool})을 각 산업의 분포 특성에 적용한 결과이다.}
\end{p5add}

\begingroup
  \renewcommand{\section}[1]{}
  \renewcommand{\subsection}[1]{\subsubsection{\syncr{#1}}}

  %% -------------------------------------------------------------------
\subsection{\vrev{Core(C2)와 Tail(C3)의 2단계 프레임워크}}
\label{subsec:two-stage-framework}
%% -------------------------------------------------------------------

앞서 논의한 교차점 분석과 이중채널 감쇠 결과에 기반하여,
일상 운영과 극단 시나리오를 분리하는 2단계(two-stage) $\DRTO$ 운영 프레임워크를
제안한다. \p{[}그림~\vref{fig:ch6-implementation-framework}\p{]}는 모니터링에서 에스컬레이션까지의
전체 운영 흐름을 도식화한 것이다.
Core(C2)와 Tail(C3)의 2단계 프레임워크에서는 2단계 분리의 원리와 Core/Tail 운영 전략을 중심으로 기술하며,
$\eta$ 기반 에스컬레이션 등급 체계와 구체적 대응 절차는
\subsecref{subsec:escalation-thresholds}에서 상세히 다룬다.

\begin{figure}[htbp]
\centering
\resizebox{\textwidth}{!}{%
\begin{tikzpicture}[
    stepbox/.style={draw, rounded corners=6pt, minimum width=36mm, minimum height=20mm,
                 align=center, font=\small, line width=0.7pt},
    detail/.style={draw, rounded corners=3pt, fill=white, minimum width=34mm,
                   align=left, font=\scriptsize, inner sep=4pt},
    conn/.style={-{Stealth[length=2.5mm]}, thick, blue!60},
    feedback/.style={-{Stealth[length=2mm]}, thick, black, dashed},
]

% 4단계 박스
\node[stepbox, fill=green!12]  (s1) at (0, 0)
  {\textbf{1. 모니터링}\\$O(t)$, $U(t)$ 수집\\APM/SIEM 연동};
\node[stepbox, fill=blue!12]   (s2) at (4.5, 0)
  {\textbf{2. 계산}\\$\DRTO$ 산출\\C2 또는 C3 적용};
\node[stepbox, fill=orange!12] (s3) at (9, 0)
  {\textbf{3. 의사결정}\\$\eta$ 평가\\Core vs Tail 판별};
\node[stepbox, fill=red!12]    (s4) at (13.5, 0)
  {\textbf{4. 에스컬레이션}\\$\eta$ 기반\\단계별 대응};

% 연결
\draw[conn] (s1) -- (s2);
\draw[conn] (s2) -- (s3);
\draw[conn] (s3) -- (s4);

% 세부 내용 (피드백 루프보다 먼저 배치)
\node[detail] (d1) at (0, -2.5) {
  $\bullet$ 관측성: APM 지표 $\to$ $O \in [0,1]$\\
  $\bullet$ 불확실성: 장애 이력 $\to$ $U$\\
  $\bullet$ 30분 주기 갱신\\
  $\bullet$ 이상 탐지 시 즉시 갱신
};
\node[detail] (d2) at (4.5, -2.5) {
  $\bullet$ Core ($\leq$ P85.8): C2 적용\\
  $\bullet$ Tail ($>$ P85.8): C3 적용\\
  $\bullet$ $\DRTO = R_0 + E_{O,bnd} + E_{U,bnd}$\\
  $\bullet$ $\eta = \DRTO / R_0$ 산출
};
\node[detail] (d3) at (9, -2.5) {
  $\bullet$ $\eta < 0.8$: 안정 구간\\
  $\bullet$ $0.8 \leq \eta < 1.0$: 주의\\
  $\bullet$ $1.0 \leq \eta < 1.5$: 경고\\
  $\bullet$ $\eta \geq 1.5$: 위기
};
\node[detail] (d4) at (13.5, -2.5) {
  $\bullet$ 단계별 보고 대상 결정\\
  $\bullet$ 추가 자원 배정\\
  $\bullet$ 비상 계획 가동 판단\\
  $\bullet$ 사후 분석 및 교정
};

% 세부 연결
\draw[dashed, black, line width=0.8pt] (s1.south) -- (d1.north);
\draw[dashed, black, line width=0.8pt] (s2.south) -- (d2.north);
\draw[dashed, black, line width=0.8pt] (s3.south) -- (d3.north);
\draw[dashed, black, line width=0.8pt] (s4.south) -- (d4.north);

% 피드백 루프 (상단 우회: s4 → 위 → s1)
\draw[feedback] (s4.north) -- ++(0, 1.0) -| (s1.north)
    node[pos=0.5, above, font=\scriptsize, text=black] {파라미터 교정 피드백};

% 단계 레이블 (피드백 루프 위에 배치)
\node[font=\footnotesize\bfseries, text=black] at (0, 2.8)  {실시간 입력};
\node[font=\footnotesize\bfseries, text=black]  at (4.5, 2.8) {모델 적용};
\node[font=\footnotesize\bfseries, text=black] at (9, 2.8) {상태 판별};
\node[font=\footnotesize\bfseries, text=black]   at (13.5, 2.8) {대응 실행};

\end{tikzpicture}
}% end resizebox
\caption{모니터링에서 계산, 의사결정, 에스컬레이션으로 이어지는 $\DRTO$ 구현 프레임워크}
\label{fig:ch6-implementation-framework}
\end{figure}

2단계 프레임워크의 핵심 원리는 \p{[}표~\vref{tab:two-stage-summary}\p{]}와 같다.

\begin{table}[htbp]
\centering
\caption{Core(C2)와 Tail(C3)의 2단계 프레임워크 요약}
\label{tab:two-stage-summary}
\small
\begin{tabularx}{\textwidth}{l X X}
\toprule
\textbf{구분} & \textbf{1단계 --- Core \m{영역(C2 기반)}} & \textbf{2단계 --- Tail \m{영역(C3 기반)}} \\
\midrule
\jrev{적용 \m{영역(P85.8 분할)}}     & P85.8 \m{이하(일상적 운영)}             & P85.8 \m{초과(극단 시나리오)} \\
\addlinespace
\jrev{환경 분포}  & 가우시안 ($U \sim N(3,1)$)           & 파레토 ($U \sim 6 \cdot \mathrm{Pa}(3)$) \\
\addlinespace
\jrev{환경 전체 평균 $\bar{\eta}$}  & $\bar{\eta}_{\text{C2}} \approx 1.68$ & $\bar{\eta}_{\text{C3}} \approx 1.51$ \\
\addlinespace
\jrev{C3 내 분할 영역 $k_O$ 계수}    & \jrev{Core: $-0.798$}                             & \jrev{Tail: $-0.415$ ($0.52$배 감쇠)} \\
\addlinespace
CVaR$_{99}$ \jrev{(환경 전체)}   & $3.180$                              & $3.944$ ($+24.0\%$) \\
\addlinespace
\jrev{환경 전체 $R^2$ (9변수)} & $0.998$                              & $0.858$ \\
\addlinespace
운영 목표      & 효율적 \m{복구(관측성 효과 극대화)}       & 보수적 \m{대응(꼬리 위험 대비)} \\
\addlinespace
자원 배분      & 정상 복구 절차                        & 추가 자원 투입, 비상 계획 연계 \\
\bottomrule
\end{tabularx}
\end{table}

\jrev{[표 5-14]는 Core/Tail의 두 가지 의미를 함께 다룬다. \chg{첫째는 P85.8 기준 \textit{분할 영역}(적용 영역 행)이고, 둘째는 분할 영역에 적용되는 \textit{환경 모델 C2/C3}(환경 모델, 환경 평균, 환경 $R^2$ 행)이며, 셋째는 C3 조건 \textit{내부}의 분할 회귀 결과($k_O$ 계수 행, 4.10절 dual\_channel 결과)이다.} 각 지표의 산출 분석 대상이 다르므로 행 단위로 의미를 구분하여 해석해야 한다.}

\textbf{1단계(Core, C2 기반)}는 P85.8 이하의 일상적 운영 상황에서 적용되며,
가우시안 불확실성 분포(C2 조건)에 기반한 효율적 복구를 추구한다.
\jrev{C2 조건 전체} 평균 $\bar{\eta} \approx 1.68$이며,
\jrev{C3 조건의 분할 회귀 결과 Core 영역의 $k_O$ 계수는 $-0.798$로} 관측성의 효과가 충분히 발현된다.
\textbf{2단계(Tail, C3 기반)}는 P85.8 초과의 극단 시나리오에서 활성화되며,
파레토 불확실성 분포(C3 조건)의 특성을 반영한다.
\jrev{C3 분할 회귀 결과 Tail 영역의 $k_O$ 계수가 $-0.415$로 $0.52$배 감쇠되고},
CVaR$_{99}$가 C2 대비 $+24.0\%$ 높아져, 불확실성 관리 중심의 추가 자원
투입이 정당화된다.

이러한 2단계 분리는 앞서 규명된 P85.8 분할 기준점의 이론적 근거에
기반하며, 조직이 일상 운영에서는 효율성을, 극단 상황에서는
안정성을 각각 최적화할 수 있는 구조적 기반을 제공한다.


  
%% -------------------------------------------------------------------
\subsection{산업별 권장 파라미터}
\label{subsec:industry-guidelines}
%% -------------------------------------------------------------------

$\DRTO$ 프레임워크의 파라미터는 산업별 특성에 따라 교정(calibration)이
필요하다. \chref{ch:results}의 시뮬레이션 결과, 전문가 의견, 그리고 연구모형의 산업
시나리오 설계를 종합하여, 6개 주요 산업에 대한 권장 파라미터 설정을
\p{[}표~\vref{tab:ch6-industry-params}\p{]}에 제시한다.

\begin{table}[htbp]
\centering
\caption{6개 산업별 $\DRTO$ 권장 파라미터}
\label{tab:ch6-industry-params}
\small
\begin{tabularx}{\textwidth}{l c c c c X}
\toprule
\textbf{산업} & \textbf{$R_0$ (h)} & \textbf{$R_{\max}$ (h)} & \textbf{$O_{\text{raw}}$ 범위} & \textbf{권장 조건} & \textbf{적용 근거} \\
\midrule
제조업       & 4.0 & 16.0 & 0.50--0.70 & C2
  & 생산 설비 고장은 가우시안 분포 적합, 생산 계획 연동으로 $R_0$ 엄격 관리 필요 \\
\addlinespace
금융         & 1.0 & 4.0  & 0.70--0.85 & C3
  & 규제 준수 요건(RTO $\leq$ 2h) 엄격, 시장 변동$\cdot$사이버 공격의 파레토 특성 반영 필수 \\
\addlinespace
의료         & 2.0 & 8.0  & 0.60--0.75 & C3
  & 환자 안전이 최우선, 의료기기 장애, 감염 확산 시나리오의 파레토 위험 대비 \\
\addlinespace
IT 서비스    & 0.5 & 2.0  & 0.75--0.90 & C2
  & SLA 기반 복구, 인프라 장애의 대다수가 가우시안 근사 가능, 높은 관측성 인프라 \\
\addlinespace
물류         & 6.0 & 24.0 & 0.45--0.60 & C2
  & 공급망 연속성, 계절 변동$\cdot$물류 지연은 정규분포 근사, 가장 긴 허용 RTO \\
\addlinespace
소매         & 3.0 & 12.0 & 0.50--0.65 & C2
  & 매출 손실 최소화, 일상적 IT 장애 중심, POS$\cdot$결제 시스템 관측성 활용 \\
\bottomrule
\end{tabularx}
\end{table}

\chg{[표 5-15]의 공통 파라미터는 다음과 같다. 곡률 매개변수는 $\kappa = 1.2$로 \jrev{\citet{lee2026drto}}에서 확정된
매개변수이며, 본 연구는 \frev{분포 환경 적합성} 종합 명제로 \frev{운영 적합성을 시사한다}(\secref{sec:ch4_hypothesis_summary}
참조). 관측성 가중치는 권장 범위 $k_O \in [0.4, 0.6]$이고, C2는 가우시안 불확실성 분포, C3는 파레토
불확실성 분포이며, $R_0 / R_{\max}$은 연구모형 기본 설정으로 1:4 비율을 유지한다. $O_{\text{raw}}$ 범위는 해당 산업의
전형적 모니터링 성숙도를 반영한다.}

\p{[}그림~\vref{fig:ch6-industry-summary}\p{]}는 6개 산업의 핵심 파라미터와
권장 조건을 종합적으로 시각화한 것이다.

\begin{figure}[htbp]
\centering
\resizebox{\textwidth}{!}{%
\begin{tikzpicture}[
    indbox/.style={draw, rounded corners=5pt, minimum width=30mm, minimum height=22mm,
                   align=center, font=\small, line width=0.6pt},
    modelC2/.style={draw, rounded corners=3pt, fill=blue!10, minimum width=12mm,
                    minimum height=6mm, font=\scriptsize\bfseries},
    modelC3/.style={draw, rounded corners=3pt, fill=red!10, minimum width=12mm,
                    minimum height=6mm, font=\scriptsize\bfseries},
    conn/.style={-{Stealth[length=2mm]}, thick, black},
]

% 산업 박스
\node[indbox, fill=green!8]  (mfg) at (0, 0)
  {\textbf{제조업}\\$R_0 = 4.0$h\\$O = 0.50$--$0.70$};
\node[indbox, fill=blue!8]   (fin) at (3.2, 0)
  {\textbf{금융}\\$R_0 = 1.0$h\\$O = 0.70$--$0.85$};
\node[indbox, fill=red!8]    (hc)  at (6.4, 0)
  {\textbf{의료}\\$R_0 = 2.0$h\\$O = 0.60$--$0.75$};
\node[indbox, fill=purple!8] (it)  at (9.6, 0)
  {\textbf{IT 서비스}\\$R_0 = 0.5$h\\$O = 0.75$--$0.90$};
\node[indbox, fill=orange!8] (log) at (12.8, 0)
  {\textbf{물류}\\$R_0 = 6.0$h\\$O = 0.45$--$0.60$};
\node[indbox, fill=yellow!10](ret) at (16.0, 0)
  {\textbf{소매}\\$R_0 = 3.0$h\\$O = 0.50$--$0.65$};

% 모델 레이블
\node[modelC2] at (0, -1.7) {C2};
\node[modelC3] at (3.2, -1.7) {C3};
\node[modelC3] at (6.4, -1.7) {C3};
\node[modelC2] at (9.6, -1.7) {C2};
\node[modelC2] at (12.8, -1.7) {C2};
\node[modelC2] at (16.0, -1.7) {C2};

% 연결
\draw[dashed, black] (mfg.south) -- (0, -1.4);
\draw[dashed, black] (fin.south) -- (3.2, -1.4);
\draw[dashed, black] (hc.south)  -- (6.4, -1.4);
\draw[dashed, black] (it.south)  -- (9.6, -1.4);
\draw[dashed, black] (log.south) -- (12.8, -1.4);
\draw[dashed, black] (ret.south) -- (16.0, -1.4);

% R₀ 크기 축 (상단)
\draw[thick, -{Stealth[length=2.5mm]}, black] (-1.2, 2.0) -- (17.2, 2.0)
    node[right, font=\scriptsize, black] {$R_0$ 증가};
\node[font=\scriptsize, black, above] at (3.2, 2.0) {짧은 RTO};
\node[font=\scriptsize, black, above] at (12.8, 2.0) {긴 RTO};

% 범례
\node[modelC2] at (5.5, -2.8) {C2};
\node[font=\scriptsize, anchor=west] at (6.2, -2.8) {가우시안 (일상적 장애)};
\node[modelC3] at (10.5, -2.8) {C3};
\node[font=\scriptsize, anchor=west] at (11.2, -2.8) {파레토 (꼬리 위험 필수)};

\end{tikzpicture}
}% end resizebox
\caption{6개 산업 $\DRTO$ 적용 요약 다이어그램}
\label{fig:ch6-industry-summary}
\end{figure}

산업별 조건 선택의 핵심 기준은 \textbf{파레토 위험의 존재 여부}이다.
금융과 의료는 사이버 공격, 시장 급변, 감염 확산 등 극단적 사건이
복구 시간을 급격히 연장할 수 있으므로 C3(파레토)를 권장한다.
C3의 CVaR$_{99}$가 C2 대비 $+24.0\%$ 높은 점은 이러한 산업에서
꼬리 위험을 사전에 반영해야 하는 정량적 근거가 된다.
반면, 제조업, IT 서비스, 물류, 소매는 장애의 대다수가 가우시안
분포로 근사 가능하므로 C2가 효율적이며, 필요 시 극단 시나리오에 한하여
C3를 보완적으로 적용할 수 있다.

분포 유형의 판단 오류는 복구 자원의 배정에서 두 가지 상반된 위험으로 이어진다.
가우시안 환경(C2)에 파레토 수준(C3)의 자원을 배정하는 과다예비(over-provisioning)는
불확실성을 과대 추정한 결과로, 유휴 대기 인력이나 미사용 예비 시스템, 과도한 예비
부품과 같은 자원이 상시 유지되어 비용 효율이 저하되나 서비스 수준 협약(SLA) 위반
위험은 낮다. 예를 들어 정상적 IT 장애 복구에 대해 랜섬웨어 대응 수준의 인력과 예비
시스템을 상시 유지하면 유휴 자원 비용이 불필요하게 증가한다. 반대로 파레토
환경(C3)에 가우시안 수준(C2)의 자원만 배정하는 과소예비(under-provisioning)는
불확실성을 과소 추정한 결과로, 비용은 절감되나 극단적 사건 발생 시 SLA 위반과
업무 중단 장기화 위험이 급증한다. 공급망 복합 장애 가능성이 있는 물류 기업이 단순
장비 교체 수준의 예비만 확보하면, P99 수준의 극단적 지연(약 $13.2$시간 이상)이
발생할 때 복구 역량이 부족하여 업무 연속성이 심각하게 위협받는다. 따라서 산업별
조건 선택은 단순한 분포 가정의 문제가 아니라 자원 배정의 비용과 위험을 균형 짓는
의사결정이며, 파레토 위험이 존재하는 산업에서 C3를 권장하는 근거가 여기에 있다.

각 산업에 대한 세부 적용 지침은 다음과 같다.

\chg{제조업의 경우,}
생산 설비 고장에 의한 업무 중단이 주요 위험이며,
관측성은 PLC(Programmable Logic Controller), SCADA, MES(Manufacturing
Execution System) 등의 기존 인프라를 활용한다.
$R_0 = 4.0$h는 일반적인 교대제 생산 라인의 복구 목표를 반영하며,
$O = 0.50$--$0.70$은 중간 수준의 모니터링 성숙도에 해당한다.

\chg{금융의 경우,}
금융감독 당국의 IT 복구 시간 규제(예: RTO $\leq$ 2시간)가 엄격하며,
\chg{사이버 공격, 시장 급변, 결제 시스템 장애}의 파레토 특성을 고려하여
C3를 권장한다. $O = 0.70$--$0.85$의 높은 관측성은 실시간 거래
모니터링 시스템(FDS 등)을 반영한다.

\chg{의료의 경우,}
환자 안전이 최우선이며, 의료기기 장애, 전자 의무기록(EMR) 시스템 중단,
감염병 확산 시나리오에서 복구 지연이 생명에 직결된다.
$R_0 = 2.0$h는 응급 의료 시스템의 전형적 복구 목표이며,
C3 적용으로 극단 시나리오에 대한 사전 대비를 강화한다.

\chg{IT 서비스의 경우,}
SLA(Service Level Agreement) 기반의 엄격한 복구 목표를 가지며,
$R_0 = 0.5$h는 클라우드/SaaS 서비스의 전형적 RTO이다.
$O = 0.75$--$0.90$의 높은 관측성은 APM, 로그 분석, 분산 추적(distributed
tracing) 등 성숙한 모니터링 인프라를 반영한다.

\chg{물류의 경우,}
공급망의 연속성이 핵심이며, 계절 변동, 운송 지연, 창고 시스템 장애가
주요 위험이다. $R_0 = 6.0$h는 본 연구의 기본 설정과 동일하며,
$O = 0.45$--$0.60$은 물류 산업의 상대적으로 낮은 IT 관측성 성숙도를 반영한다.

\chg{소매의 경우,}
POS(Point of Sale) 시스템, 재고 관리, 전자상거래 플랫폼의 가용성이
매출에 직접 영향을 미친다. $R_0 = 3.0$h, $O = 0.50$--$0.65$는
중간 규모 소매 환경의 전형적 특성을 반영하며, C2 기반의 효율적 복구를 권장한다.


\chg{이상의 산업별 권장 파라미터를 회귀식에 대입한 구체적 $\DRTO$ 산출 예시는 회귀식 기반 산출 도구를
다루는 \subsecref{subsec:practical-regression-tool}에 제시한다.}



\endgroup
\FloatBarrier
      % 5.3 실무적 기여
% =====================================================================
% 5.4 한계점 — 5단계 5-2 wrapper (sec5.4_limitations.tex + 5.1.3 흡수)
% 본문 텍스트 무수정. 원본 sec5.4_limitations.tex 그대로 + 5.1.3 본문 append.
% =====================================================================



%% ===================================================================
\section{\vrev{한계점}}
\label{sec:limitations}
%% ===================================================================

본 연구는 이론적 모델 수립과 시뮬레이션 기반 검증에 초점을 맞추고 있으며,
다음의 한계를 지닌다. \p{[}표~\vref{tab:ch6-limitations}\p{]}은 각 한계와
대응하는 향후 연구 방향을 종합한다.

\begin{table}[htbp]
\centering
\caption{\revdel{한계점 및 향후 연구 방향}\,\revadd{한계점 종합}}
\label{tab:ch6-limitations}
\small
\begin{tabularx}{\textwidth}{c l X X}
\toprule
\textbf{\#} & \textbf{한계 영역} & \textbf{한계 내용} & \textbf{향후 연구 방향} \\
\midrule
1 & 시뮬레이션 기반 검증 &
  $n = 10{,}000$ MC 시뮬레이션과 $10^8$회 다중 시드 검증은 통계적으로 강건하나,
  실제 조직 환경에서의 현장 실험(field test)을 수반하지 않음 &
  제조$\cdot$금융$\cdot$IT 서비스 등 산업별 파일럿 현장 검증, 실제 장애 이벤트에서의
  $\eta$ 예측 정확도 평가 \\
\addlinespace
2 & 고정 $R_{\max}/R_0$ 비율 &
  $R_0 = 6.0$h, $R_{\max} = 24.0$h ($1:4$ 비율) 고정,
  산업$\cdot$조직별 비율 변동 시 모델 행태 미탐색 &
  $R_{\max}/R_0$ 비율을 $[2, 8]$ 범위에서 변동시키는 민감도 분석,
  동적 $R_0(t)$ 모형으로의 확장 \\
\addlinespace
3 & 분포 가정 &
  C2: $U \sim N(3,1)$, C3: $6 \cdot \mathrm{Pareto}(\alpha=3)$ 고정,
  실제 불확실성은 다양한 꼬리 두께를 가질 수 있음 &
  \jrev{일반화 파레토 분포(GPD)와 GEV 분포 도입, 가우시안-파레토 혼합 분포 모형 개발, $\alpha$ 변동에 따른 민감도 분석(P85.8 교차점 이동 및 감쇠 비율 변화)} \\
\addlinespace
4 & 시간 동역학 부재 &
  현행 모델은 시점별 단면(cross-sectional) 분석,
  시간에 따른 $O(t)$, $U(t)$의 자기상관, 추세, 계절성 미반영 &
  \jrev{시계열 모형(ARIMA, GARCH) 통합, 베이지안 온라인 학습 기반 자기 교정 메커니즘 구축, 복구 과정 중 실시간 $\DRTO$ 갱신 프로토콜 수립} \\
\addlinespace
\tmp{5} & 심층 불확실성 미반영 &
  현행 모델은 분포 유형(가우시안/파레토)이 사전에 특정 가능하다는 전제,
  모델$\cdot$시나리오$\cdot$확률 모두 불확실한 심층 불확실성(deep uncertainty) 상황은
  적용 범위를 벗어남 &
  로버스트 최적화, 시나리오 발견,
  info-gap 의사결정 이론 등 분포 비의존적 접근법과의 통합 탐색 \\
\addlinespace
\tmp{6} & 국가 위기경보 체계와의 연계 부재 &
  $\eta$ 기반 에스컬레이션은 정량적 지표에 의해 등급이 자동 결정되나,
  국가 위기경보(관심--주의--경계--심각)의 질적 판단 체계와의
  공식적 연계 메커니즘이 마련되지 않음 &
  국가 위기경보 4단계와 $\eta$ 임계값의 매핑 프레임워크 개발,
  질적 판단과 정량적 지표의 상호 보완적 의사결정 모형 구축 \\
\addlinespace
\jrev{7} & \jrev{\frev{분포 환경 적합성의 간접 시사}} &
  \jrev{본 연구는 \chg{\citet{lee2026drto}}의 $\kappa^{*}=1.2$를 C2/C3 조건에 적용한 후 9개 가설 PASS의 통합 해석으로 \frev{운영 적합성을 시사하였으나}, 직접적 격자 탐색은 4기준 종합점수 $F(\kappa)$가 C1 특화 지표라 C2/C3에 부적합} &
  \jrev{환경 무관 다기준 종합점수 재정의, C2/C3 조건에 적합한 새로운 4기준 도출 후 직접 격자 탐색 수행} \\
\bottomrule
\end{tabularx}
\end{table}


%% -------------------------------------------------------------------
\subsection{\vrev{시뮬레이션 기반 검증의 한계}}
\label{subsec:lim-simulation}
%% -------------------------------------------------------------------

$n = 10{,}000$ 몬테카를로 시뮬레이션과 $10^8$회 다중 시드 검증은
통계적으로 강건한 결과를 제공하나, 실제 조직 환경에서의 현장 실험을
수반하지 않았다. 시뮬레이션에서 확인된 효과 크기
(Cohen's $d = -1.180$, $+1.871$)와 이중채널 감쇠($0.52$배)가
실제 복구 과정에서도 동일한 크기로 발현되는지는
추가적인 실증 연구를 통해 확인되어야 한다.
특히 관측성 지표 $O(t)$의 실시간 측정 정확도, 불확실성 분포의
실제 형태, 그리고 조직의 복구 역량 변동이 시뮬레이션 가정과
어떻게 괴리되는지를 현장 데이터로 검증하는 것이 필요하다.
$O_{\text{raw}}$의 조작적 정의로서, 장애 탐지율, 로그 가시성, 복구 진척 추적 수준 등의
하위 지표를 리커트 척도로 평가하여 $[0, 1]$로 정규화하는 방법을 예시할 수 있으나,
측정 지표의 가중치 산정과 \chg{산업별 및 조직별} 특성을 반영한 측정 도구의 정밀화는
후속 연구에서 수행되어야 한다.


%% -------------------------------------------------------------------
\subsection{\vrev{고정 $R_{\max}/R_0$ 비율의 한계}}
\label{subsec:lim-ratio}
%% -------------------------------------------------------------------

$R_{\max}/R_0$ 비율을 모든 산업에서 4로 고정하여 비율 변동의 영향을 탐색하지
않은 점은 본 연구의 한계이다. 본 연구에서는 $R_0 = 6.0$h, $R_{\max} = 24.0$h로 $R_{\max}/R_0 = 4$의
고정 비율을 사용하였다. 이 비율은 물류 산업의 전형적 설정(\p{[}표~\vref{tab:ch6-industry-params}\p{]})과
일치하나, 금융($R_{\max}/R_0 = 4$), IT 서비스($R_{\max}/R_0 = 4$) 등
모든 산업에서 동일 비율이 적용되었다. 실무에서는 \chg{산업과 조직, 규제} 환경에 따라
이 비율이 $2$에서 $8$ 이상까지 변동할 수 있으며, 비율 변동이 시그모이드의
작동 영역($z$ 범위)과 이중채널 감쇠 비율에 미치는 영향은
체계적으로 탐색되지 않았다.

향후 연구에서는 $R_{\max}/R_0$ 비율을 독립 변수로 설정한 민감도 분석을 통해,
비율 변동에 대한 모델의 강건성을 검증하고, 산업별 최적 비율 범위를
도출하는 것이 필요하다. 나아가 $R_0(t)$ 자체를 시간 변동 함수로 확장하여,
BIA 갱신 시마다 기준선이 적응적으로 조정되는 동적 기준선 모형을 개발할 수 있다.


%% -------------------------------------------------------------------
\subsection{\vrev{분포 가정의 한계}}
\label{subsec:lim-distribution}
%% -------------------------------------------------------------------

불확실성 분포를 가우시안과 파레토($\alpha = 3$)의 두 형태로 고정한 점은 본 연구의
또 다른 한계이다. C2 조건에서의 불확실성 분포 $U \sim N(3,1)$과 C3 조건에서의
$6 \cdot \mathrm{Pareto}(\alpha = 3)$은 각각 경꼬리와 중꼬리의
대표적 형태를 포착하기 위한 선택이다. 그러나 실제 재난 복구 과정에서의
불확실성은 $\alpha < 3$(더 무거운 중꼬리)이나 $\alpha > 3$(더 가벼운 경꼬리)의
다양한 특성을 보일 수 있다.

$\alpha$ 값의 변동에 따른 핵심 결과의 민감도\chg{(}P85.8 교차점의 이동,
CVaR 격차의 변화, 이중채널 감쇠 비율\chg{)}는 향후 연구에서 체계적으로 탐색되어야 한다.
일반화 파레토 분포(GPD)나 극단값 이론(EVT)의 GEV 분포를
불확실성 모형에 도입하여, 다양한 꼬리 두께에 대한 모델의 적응성을
확보하는 것이 바람직하다.

\jrev{또한 본 연구의 \textbf{\frev{분포 환경 적합성} 종합 명제}(4.13절)는 \chg{\citet{lee2026drto}}의
$\kappa^{*}=1.2$의 C2/C3 조건 적용성을 9개 가설(\jrev{H2--H10}) PASS의 통합 해석으로
\frev{시사한다}. 이는 직접적 격자 탐색의 대안으로서 학술적 정직을 우선한 접근이지만,
\chg{\citet{lee2026drto}}의 4기준 종합점수 $F(\kappa) = f_1 \cdot f_2 \cdot f_3 \cdot f_4$
(\chg{실무 유의성, 비포화, 효과 크기, 설명력})가 C1 조건 기반으로 정의되어
C2/C3 환경에서 직접 격자 탐색에 부적합하다는 \textbf{지표 일반화의 한계}를
내포한다. 환경 무관 다기준 종합점수의 재정의 및 C2/C3 환경에서의
직접 격자 탐색은 본 연구 범위를 벗어나는 후속 연구 과제이다.}


%% -------------------------------------------------------------------
\subsection{\vrev{시간 동역학의 부재}}
\label{subsec:lim-temporal}
%% -------------------------------------------------------------------

현행 $\DRTO$ 모델은 특정 시점의 $O(t)$와 $U(t)$를 입력으로 받아
해당 시점의 $\DRTO$를 산출하는 단면(cross-sectional) 모형이다.
시간에 따른 입력 변수의 자기상관(autocorrelation), 추세(trend),
계절성(seasonality)은 모형에 내생적으로 반영되지 않는다.

실제 복구 과정에서는 $O(t)$가 복구 진행에 따라 점진적으로 증가하고,
$U(t)$는 초기에 높았다가 정보 획득에 따라 감소하는 시간적 패턴을 보인다.
이러한 동역학을 포착하기 위해, 베이지안 온라인 학습(Bayesian online learning)을
적용하여 장애 이벤트 관측 시마다 매개변수를 적응적으로 갱신하는
자기 교정(self-calibrating) 메커니즘의 구축이 향후 연구 방향이다.
ARIMA, GARCH 등 시계열 모형~\citep{durbin2012timeseries,hyndman2018forecasting}과의 통합을 통해
$\DRTO$ 모형에 시간 차원을 추가하는 확장도 고려할 수 있다.


%% -------------------------------------------------------------------
\subsection{\vrev{심층 불확실성의 미반영}}
\label{subsec:lim-deep-uncertainty}
%% -------------------------------------------------------------------

심층 불확실성을 반영하지 못하는 점은 현행 모델의 구조적 한계이다.
현행 $\DRTO$ 모델은 불확실성의 확률분포 유형이 사전에 특정 가능하다는 전제 하에
설계되었다. C2 조건에서는 가우시안, C3 조건에서는 파레토 분포를 채택하여
각각 경꼬리와 중꼬리 특성을 포착한다. 그러나 팬데믹, 기후변화에 의한 복합재난,
전례 없는 대규모 사이버 공격 등 \textbf{심층 불확실성}(deep uncertainty) 상황에서는
분포의 형태는 물론 모델 구조와 시나리오 자체를 사전에 특정할 수 없다.

이는 \subsecref{subsec:uncertainty-risk}에서 정리한 불확실성 분류 체계의
네 번째 유형\chg{(}\chg{모델, 시나리오, 확률} 모두 불확실한 심층 불확실성\chg{)}에 해당하며,
현 프레임워크의 구조적 적용 한계를 구성한다.
나이트적(Knightian) 불확실성은 파레토 분포로 보수적 근사가 가능하나,
심층 불확실성은 분포 선택 자체가 불가능하므로 확률분포 기반 모형의 범위를 벗어난다.

향후 연구에서는 로버스트 최적화(robust optimization),
시나리오 발견(scenario discovery),
info-gap 의사결정 이론 등 분포 비의존적(distribution-free) 접근법을
$\DRTO$ 프레임워크와 통합하여,
심층 불확실성 상황에서도 복구시간의 동적 조정이 가능한 확장 모형을 탐색할 필요가 있다.


%% -------------------------------------------------------------------
\subsection{\vrev{국가 위기경보 체계와의 연계}}
\label{subsec:lim-national-alert}
%% -------------------------------------------------------------------

본 연구의 $\eta$ 기반 에스컬레이션은 국가 위기경보 체계와의 공식적 연계
메커니즘을 아직 갖추지 못하였다. \subsecref{subsec:escalation-thresholds}에서 논의한 바와 같이,
현행 국가 위기경보 체계(관심--주의--경계--심각)는
재난 및 안전관리 기본법에 근거한 \textbf{질적 판단}에 기반하여
전문가 합의와 상황 평가를 통해 등급이 결정된다.
반면 본 연구의 $\eta$ 기반 에스컬레이션은
$\DRTO / R_0$라는 \textbf{정량적 지표}에 의해 등급이 자동으로 산출되므로,
두 체계는 상호 보완적 관계를 형성할 잠재력을 지닌다.

구체적으로, 국가 위기경보가 발령된 상황에서 $\eta$ 값을 병행 산출하면
질적 판단의 객관적 근거를 제공할 수 있으며,
역으로 $\eta$의 급격한 상승이 감지될 때 위기경보 상향 검토의
조기 경고(early warning) 신호로 활용할 수 있다.
향후 연구에서는 국가 위기경보 4단계와 $\eta$ 임계값 간의
체계적 매핑 프레임워크를 개발하고,
질적 판단과 정량적 지표가 결합된 하이브리드 의사결정 모형을 구축하여
국가 수준의 재난 대응 체계에 $\DRTO$의 정량적 기여 가능성을
실증적으로 검증하는 것이 필요하다.


%% -------------------------------------------------------------------
\syncr{\subsection{하위 변수 조작화의 단순화}}
\syncr{\label{subsec:lim-operationalization}}
%% -------------------------------------------------------------------

\syncr{본 연구는 관측성($O_{\text{raw}}$, $k_O$)과 불확실성($U$)의
제1계층 하위 지표까지의 조작화 \textbf{예시}를 \subsecref{app:operationalization}에
참고용으로 제시하였으나, 가중치($w_i$, $\alpha$, $\beta$)의 산업별 튜닝,
원시 지표의 측정 도구별 \chg{정확도와 편향} 검증, 현장 데이터 기반 \chg{수렴 및 예측}
타당도 확보는 본 연구의 범위를 벗어나는 후속 연구 과제이다. 제시된
조작화 예시와 가상 기업 라이프사이클은 실무 적용의 방향성을 보여주기
위한 참고용일 뿐 검증된 방법론이 아니며, 심리측정학적 엄밀성을 갖춘
측정 도구 개발 및 종단 연구를 통한 예측 타당도 확보가 선결되어야 한다.}

\vspace{1em}

\fpara{본 연구의
가장 중요한 학술 한계는 \emph{$\kappa = 1.2$의 환경 한정성}이다. \chg{\citet{lee2026drto}}의 도출값
$\kappa = 1.2$는 C1 조건(관측성 단독, 가우시안 베이스라인) 격자 탐색에서 도출되었으며,
C2/C3 환경에서의 적합성은 4.13절의 H2--H10 PASS 통합 해석으로 \emph{운영 적합성을
시사}하나 \emph{환경별 최적성을 입증}하지는 못한다. 따라서 본 연구는 $\kappa = 1.2$를
BCM 일반 운영의 \emph{\chg{잠정적이고 보수적인} 실무 기본값}(provisional and conservative operational
default)으로 채택하며, \emph{환경 무관 보편 상수성을 주장하지 않는다}. 환경 무관 $\kappa$의
\emph{존재 여부}에 대한 다방법 탐색은 후속 연구 논문의 핵심 의제로 분리된다
(5.5.3절 향후 연구 참조).}

본 장에서 제시한 \chg{이론적, 방법론적, 실무적} 시사점과 한계는
$\DRTO$ 프레임워크의 현재 위상을 명확히 규정하는 동시에,
후속 연구의 구체적 방향을 제시한다. 특히 2단계 프레임워크(Core C2 + Tail C3),
$\eta$ 기반 에스컬레이션 체계, 6개 산업 적용 가이드라인은
$\DRTO$의 이론적 기여를 실무적 가치로 전환하기 위한 구체적 기반을 제공하며,
\tmod{\chref{ch:discussion}}{\tsecref{sec:conclusion-future-merged}}에서는 이러한 논의를 종합하여 본 연구의 결론과 향후 연구 방향을 제시한다.% 3차 심사 (2026-05-09): 자기참조 오류 수정 — \chref{ch:discussion}가 이 5장 자체를 가리킴. 5.5절을 가리키도록 \tsecref(빨간색 secref)로 수정.


% --- 5.1.3에서 흡수: 종합 논의 절에 있던 ``연구의 한계'' 4단락 ---

%% -------------------------------------------------------------------
\subsection{\syncr{연구의 한계 (소규모 패널 \& 단일 $\kappa$)}}
\label{subsec:ch5_limitations}
%% -------------------------------------------------------------------

본 연구의 한계는 다음과 같다.

\fdel{\noindent\textbf{시뮬레이션 기반 검증의 한계.}
본 연구의 결과는 $n = 10{,}000$ 몬테카를로 시뮬레이션에 기반하며,
실제 장애 사건 데이터에 의한 검증은 수행되지 않았다.
시뮬레이션에서 사용한 입력 분포($k_O, O \sim \mathrm{Uniform}(0,1)$,
$U \sim N(3,1)$ 또는 $6 \cdot \mathrm{Pareto}(3)$)는 이론적
분석에 적합하나, 실제 조직의 관측성과 불확실성 분포를
완전히 반영하지 못할 수 있다.
향후 산업별 실제 장애 이력 데이터를 활용한 경험적 검증이 필요하다.}

\chg{전문가 패널의 규모 측면에서,}
\chg{정량적과 정성적} 전문가 검증을 5명의 소규모 패널로 수행하였다.
S-CVI/Ave $= 0.985$, Cronbach's $\alpha = 0.89$로 내용 타당도 및 내적 일관성이
확인되었으나, 대규모 패널에 의한 일반화 가능성은 추가 검증이 필요하다.
특히 산업별, 규모별 조직의 다양성을 반영한 전문가 패널 확장이
후속 연구에서 수행되어야 한다.

\chg{단일 $\kappa$ 값의 적용 측면에서,}
본 연구는 \jrev{\chg{\citet{lee2026drto}}에서 \chg{도출한}} $\kappa = 1.2$의 단일 곡률 매개변수를 전체 조건에 적용하였다.
\jrev{본 연구는 \frev{분포 환경 적합성} 종합 명제(\secref{sec:ch4_hypothesis_summary})를 통해 9개 가설(H2--H10) PASS의 통합 해석으로 $\kappa = 1.2$의 C2/C3 조건 \frev{운영 적합성을 시사하였으나},}
$\kappa$의 조건별 또는 산업별 차별화\jrev{, 그리고 환경 무관 다기준 종합점수의 재정의를 통한 직접 격자 탐색}이 모델의 적응력을 향상시킬 수 있으며,
이는 후속 연구의 과제이다.

\fdel{\noindent\textbf{시간적 역동성의 미반영.}
현재 $\DRTO$ 모델은 특정 시점의 관측성과 불확실성을 입력으로 하는
정적(static) 산출 구조이다. 실제 복구 과정에서 관측성과 불확실성은
시간에 따라 변화하므로, 시간적 역동성을 반영하는
동적(dynamic) 갱신 모델로의 확장이 필요하다.}

    % 5.4 한계점 (+5.1.3 흡수)
% =====================================================================
% 5.5 결론 및 향후 연구 — 5단계 5-2 (clean subsection wrapper)
% 본문 텍스트 무수정. 분할 본문 파일을 3항 구조로 재배치.
%   5.5.1 연구 요약       ← 현 sec5.5.1_summary.tex
%   5.5.2 주요 연구 성과  ← 현 sec5.5.2_achievements.tex (5.7.x → 5.5.2.x로 강등)
%   5.5.3 향후 연구 방향  ← 현 sec5.5.3_future_work.tex
% =====================================================================

\section{\vrev{결론 및 향후 연구}}
\label{sec:conclusion-future-merged}

\begin{p5add}
본 절에서는 본 연구의 전체 흐름을 압축
요약하고(5.5.1항), 핵심 성과 5개를 정리하며(5.5.2항), 한계점에서 도출된
향후 연구 방향(5.5.3항)을 제시한다. 이로써 본 연구의 학술적 기여
범위와 실무적 적용 가능성을 명확히 하고, 후속 연구의 출발점을 마련한다.
\end{p5add}

% --- 5.5.1 연구 요약 ---
\subsection{\vrev{연구 요약}}
\label{subsec:conclusion-summary-merged}

\begingroup
  \renewcommand{\section}[1]{}
  %% ===================================================================
\section{\vrev{연구 요약}}
\label{sec:conclusion-summary}
%% ===================================================================

현행 BCMS에서 RTO는 BIA 시점에
고정값으로 설정되며, 이후 운영 환경의 변화와 무관하게 유지된다.
이러한 정적(static) RTO 접근법은 세 가지 구조적 한계를 내포한다.
첫째, 조직의 관측성(observability) 수준이 시간과 상황에 따라 변동함에도
이를 반영하지 못한다.
둘째, 복구 과정에 내재된 불확실성(uncertainty)을 확률 분포에 기반하여
정량화하지 않는다.
셋째, 사이버 공격이나 복합 재난 등 극단적 사건의 파레토(중꼬리, Heavy-Tail) 특성을
구조적으로 포착하지 못한다.

본 연구는 이러한 한계를 극복하기 위하여, 관측성과 불확실성을
개별 시그모이드 바운딩(individual sigmoid bounding) 함수로 통합하는
$\DRTO$ 프레임워크를 제안하였다.
핵심 수식은 \m{식~(\ref{eq:ch7-drto-core})에 나타낸 바와 같이} 정의된다.

\begin{equation}
  \DRTO = R_0 + E_{O,\text{bnd}} + E_{U,\text{bnd}}
  \label{eq:ch7-drto-core}
\end{equation}

\noindent
여기서 관측성 효과 $E_{O,\text{bnd}} = -R_0 \cdot S(z_O)$는
$\DRTO$를 기준선 $R_0$로부터 하향 조정하고,
불확실성 효과 $E_{U,\text{bnd}} = (R_{\max} - R_0) \cdot S(z_U)$는
상향 조정하며, 수정 시그모이드 $S(z) = 2/(1 + e^{-z}) - 1$에 의해
물리적 경계 $\DRTO \in (0,\; R_{\max})$가 점근적으로 보장된다.
$R_0 = 6.0$시간, $R_{\max} = 24.0$시간, $\kappa = 1.2$를 기본 설정으로 채택하였다.

연구 검증은 세 축으로 수행되었다.
첫째, $n = 10{,}000$ 몬테카를로 시뮬레이션을 4개 조건(C0--C3)에 대해
실시하여 \jrev{10개 연구가설}(H1--\jrev{H10})의 모델 내부 정합성을 체계적으로 확인하였다.
둘째, BCM/DR 분야 전문가 5명을 대상으로 구조화된 설문(27문항, 5점 리커트 척도)과
심층 인터뷰(약 30분, 반구조화)를 수행하여 모델의 실무 타당성을 삼각 검증하였다.
셋째, 1차 단순, 1차 다중, 2차 회귀 분석 및 P85.8 기준 분할 회귀를 통해
모델 응답 곡면의 구조적 특성을 규명하고, 꼬리 감쇠 현상을 발견하였다.
그 결과, \textbf{\jrev{10개} 가설 전원이 사전 설정 기준을 충족하여 채택}되었다.
넷째, \chg{\citet{lee2026drto}}과 본 연구의 5개 연구 단계에 대한
\chg{연구모형, 온톨로지 분석, 논문 내용, 기준값} 간
\textbf{303개 교차검증 항목을 전수 확인}하여
\chg{수치, 수식, 용어}의 완전 정합(전원 PASS)을 달성하였으며,
이는 연구 프레임워크의 논리적 일관성과 데이터 신뢰성을 보증한다.



\endgroup

% --- 5.5.2 주요 연구 성과 ---
\subsection{\vrev{주요 연구 성과}}
\label{subsec:conclusion-achievements-merged}

\begingroup
  \renewcommand{\section}[1]{}
  \renewcommand{\subsection}[1]{\subsubsection{#1}}
  %% ===================================================================
\section{\vrev{주요 연구 성과}}
\label{sec:conclusion-achievements}
%% ===================================================================

\revdel{본 연구의 5대 핵심 성과를 다음과 같이 정리한다.}\,\revadd{본 항은 5장 전체에서 도출된 5대 핵심 성과를 압축 요약 형태로 제시한다. 각 성과의 상세 논의는 학술적 측면은 \secref{sec:academic-contrib}, 실무적 측면은 \secref{sec:practical-contrib}을 참조하라.}


%% -------------------------------------------------------------------
\chg{첫째 성과는 물리적 경계 보장을 갖춘 개별 시그모이드 바운딩 $\DRTO$ 모델이다.}
\label{subsec:achiev-model}
\chg{개별 시그모이드 바운딩을 통해 $\DRTO \in (0, R_{\max})$를 보장하고, $C^{\infty}$ 연속성과 자연스러운 교호작용 포착을 함께 달성하였다(상세는 \subsecref{subsec:academic-theoretical}).}

\fdel{업무연속성관리 분야 최초로, 관측성과 불확실성을 단일 수리 모델에서
통합하는 개별 시그모이드 바운딩 $\DRTO$ 모델을 정립하였다.
수정 시그모이드 함수 $S(z) = 2/(1 + e^{-z}) - 1$을 기저 함수로 채택하여,
관측성 효과($E_{O,\text{bnd}}$)와 불확실성 효과($E_{U,\text{bnd}}$)를
각각 고유한 물리적 범위($(-R_0,\; 0]$과 $[0,\; R_{\max} - R_0)$)에서
독립적으로 바운딩하였다. 이 모델은 물리적 경계 보장($\DRTO \in (0, R_{\max})$),
무한 미분 가능성($C^{\infty}$), 교호작용의 자연스러운 포착이라는
세 요구사항을 동시에 충족한다.
기존의 클리핑(clipping) 기반 접근이나 경계 미보장 모델과 달리,
시그모이드 바운딩은 점근적으로 경계에 도달하되 경계값 자체를 취하지 않아
수치적 안정성과 이론적 엄밀성을 동시에 확보한다.}


%% -------------------------------------------------------------------
\chg{둘째 성과는 \citet{lee2026drto}의 $\kappa^{*} = 1.2$에 대한 \frev{분포 환경 적합성} 종합 명제 도출이다.}
\label{subsec:achiev-kappa}
\chg{\citet{lee2026drto}의 $\kappa^{*}=1.2$를 본 연구의 새로운 분포 환경(C2 가우시안과 C3 파레토)에 적용한 결과, 9개 가설(H2--H10) PASS의 통합 해석을 통해 $\kappa = 1.2$의 분포 환경 \frev{운영 적합성}을 시사하였다(상세는 \secref{sec:ch4_hypothesis_summary}의 \frev{분포 환경 적합성} 종합 명제와 \subsecref{subsec:academic-methodological}).}

\fdel{시그모이드 곡률 파라미터 $\kappa$의 최적값을 다중 기준 복합 점수(multi-criteria
composite score) 방법론으로 도출하였다.
민감도(sensitivity), 선형성(linearity), 경계 여유(boundary margin),
해석 가능성(interpretability)의 4개 기준을 정규화하여 종합 평가한 결과,
$\kappa^{*} = 1.2$가 최적점으로 결정되었다.
$\kappa = 1.2$에서 관측성 채널의 $z_O \in [0,\; 1.2]$는
시그모이드의 비선형 압축이 의미 있게 작동하는 영역에 해당하여
($S(1.2) = 0.537$), 중간 입력에서의 감도를 확보하면서도
극단 영역에서의 비선형 포화(saturation) 거동이 물리적 경계 보장을
자연스럽게 실현한다.}


%% -------------------------------------------------------------------
\chg{셋째 성과는 C0$\to$C3 4단계 점진적 프레임워크와 \jrev{10개} 가설의 전원 채택이다.}
\label{subsec:achiev-framework}
\chg{C0--C3 4단계 점진적 조건과 6단계 다층 검증을 통해 \jrev{10개} 가설(H1--\jrev{H10})을 전원 채택하였다(상세는 \subsecref{subsec:academic-empirical}).}

\fdel{정적 기준선(C0), 관측성 도입(C1), 가우시안 불확실성(C2),
파레토 불확실성(C3)의 4단계 점진적 조건 체계를 수립하고,
6단계 다층 검증 프레임워크(L0 기반 검증(Golden Compare),
시뮬레이션 검증(H1--\jrev{H3}), 회귀 검증(\jrev{H4}--\jrev{H5}), 강건성 검증(\jrev{H6}--\jrev{H7}),
이중채널 검증(\jrev{H8}), 전문가 평가(\jrev{H9}--\jrev{H10}))를 통해 \jrev{10개} 가설을 체계적으로 검증하였다.
정량적 결과의 핵심은 다음과 같다.
C1에서 관측성에 의한 $\eta$ 감소($\bar{\eta} = 0.853$, Cohen's $d = -1.180$),
C2에서 불확실성 프리미엄의 발현($\bar{\eta} = 1.682$, Cohen's $d = +1.871$),
C3에서 평균 유사$\cdot$꼬리 급증($\bar{\eta} = 1.514$, CVaR$_{99}$ $+24.0\%$).
$R^2_{(2\mathrm{q})}$는 C1$\,= 1.000$, C2$\,= 0.998$이며,
C3는 전체 $R^2_{(2\mathrm{q})} = 0.858$이나 P85.8 분할 후 Core $R^2 = 0.999$로 \jrev{H4}가 지지되었다.}


%% -------------------------------------------------------------------
\chg{넷째 성과는 꼬리 감쇠 현상의 발견($0.52$배, 꼬리 영역)이다.}
\label{subsec:achiev-dual-channel}
\chg{P85.8 분할 회귀에서 $k_O$ 계수가 Core $-0.798$에서 Tail $-0.415$로 $0.52$배 감쇠하였고, CVaR$_{99}$는 $+24.0\%$ 증가하였다(상세는 \subsecref{subsec:academic-theoretical}).}

\fdel{C3 조건에서 P85.8 기준으로 분할 회귀를 수행한 결과,
관측성 가중치 $k_O$의 회귀 계수가 핵심 영역(Core, $-0.798$)에서
꼬리 영역(Tail, $-0.415$)으로 $0.52$배 감쇠되는 꼬리 감쇠
(tail attenuation) 현상을 발견하였다.
이 발견은 극단적 불확실성 환경에서 시그모이드 포화에 의해
관측성 투자의 한계 효용이 완화된다는 이론적 시사점을 제공한다.
P85.8 CDF 교차점은 가우시안(C2)과 파레토(C3)의 CDF가
교차하는 구조적 경계로서, 교차점 이하에서 두 분포는
사실상 구분 불가능하나(Cohen's $d \approx 0$),
교차점 초과에서 C3의 $\eta$가 C2를 급격히 상회하여
CVaR$_{99}$ 기준 $+24.0\%$의 격차를 보인다.}


%% -------------------------------------------------------------------
\chg{다섯째 성과는 전문가 검증(4.46/5.0)과 6단계 다층 검증이다.}
\label{subsec:achiev-validation}
\chg{전문가 5명의 평균이 4.46/5.0($p=0.011$, $d=1.65$)이고 Cronbach's $\alpha=0.89$, S-CVI/Ave $=0.985$였으며, 인터뷰에서 15개 주제 모두 $\geq 80\%$ 동의를 얻었다(상세는 \subsecref{subsec:academic-empirical}).}

\fdel{BCM/DR 분야 전문가 5명(평균 경력 15.2년)을 대상으로
수행한 정량적 설문에서 전체 평균 4.46/5.0점($t(4) = 3.68$, $p = 0.011$,
Cohen's $d = 1.65$)을 달성하였으며,
Cronbach's $\alpha = 0.89$, S-CVI/Ave$\,= 0.985$로
높은 내적 일관성과 평가자 간 일치도가 확인되었다(\jrev{H9} 지지).
심층 인터뷰에서는 5대 주제 15개 하위 주제 모두(100\%)에서 80\% 이상
전문가 동의가 달성되어 \jrev{H10}이 지지되었다.
시뮬레이션 검증, 회귀 검증, 강건성 검증, 이중채널 검증, 전문가 평가의
6단계(L0--L5)가 계층적 의존 구조를 형성하여 모델의 타당성을 다각적으로 보장하며,
정량적 시뮬레이션과 정성적 전문가 평가의 삼각검증(triangulation)을 통해
방법론적 편향을 최소화하였다.}

\vspace{0.5em}

\p{[}표~\vref{tab:ch7-findings-summary}\p{]}는 이상의 5대 성과를 핵심 결과,
정량적 근거, 학술적 의의의 관점에서 종합한 것이다.

\begin{table}[htbp]
\centering
\caption{핵심 연구 결과 요약}
\label{tab:ch7-findings-summary}
\small
\begin{tabularx}{\textwidth}{l X X}
\toprule
\textbf{핵심 결과} & \textbf{정량적 근거} & \textbf{학술적 의의} \\
\midrule
개별 시그모이드 바운딩 모델
  & $\DRTO \in (0,\; R_{\max})$ 보장, $C^{\infty}$ 연속
  & BCM 분야 최초의 물리적 경계 보장 동적 RTO 모델 \\
\addlinespace
\jrev{$\kappa^{*}=1.2$ \frev{분포 환경 적합성}}
  & \jrev{9개 가설(H2--H10) PASS 통합 해석}
  & \jrev{\chg{\citet{lee2026drto}의} 매개변수의 본 연구 환경 적용성을 다중 가설 PASS로 \frev{시사한} 학술적 명확성} \\
\addlinespace
C0$\to$C3 프레임워크, H1--\jrev{H10} 전원 채택
  & $\eta$: 0.853(C1), 1.682(C2), 1.514(C3); C1/C2 $R^2 \geq 0.998$, C3 Core $R^2 = 0.999$
  & 4단계 점진적 BCM 성숙도 모형의 실증적 근거 \\
\addlinespace
꼬리 감쇠 ($0.52$배)
  & Tail $k_O$ 계수 $-0.415$ vs Core $-0.798$; CVaR$_{99}$ $+24.0\%$
  & 극단 환경에서 관측성 투자 효용 감소의 최초 정량적 규명 \\
\addlinespace
전문가 검증 + 6단계 다층 검증
  & 4.46/5.0, $\alpha = 0.89$, S-CVI/Ave$\,= 0.985$; \jrev{10/10} 가설 채택
  & \chg{시뮬레이션과 전문가} 삼각검증의 방법론적 표준 제시 \\
\bottomrule
\end{tabularx}
\end{table}

\chg{이상의 5대 성과는 학술적 의의에 그치지 않고 본 장의 실무적 기여로 직접 이어진다. 개별 시그모이드
바운딩 모델과 회귀 설명력(첫째와 셋째 성과)은 회귀식 기반 $\DRTO$ 산출 도구(\subsecref{subsec:practical-regression-tool})로,
꼬리 감쇠 발견(넷째 성과)은 이중채널 감쇠 기반 자원배분과 스트레스 테스트
도구(\subsecref{subsec:practical-dualchannel-tool})로, $\kappa = 1.2$의 적합성과 전문가 검증(둘째와 다섯째 성과)은
산업별 적용 가이드라인과 운영 도구의 신뢰 근거로 활용된다. 한편 위 표가 다루는 다섯 핵심 성과 외에도
4장에서 도출된 세부 결과들이 의미의 크고 작음과 무관하게 각각 어느 학술 기여 또는 실무 활용으로
귀착되는지는 \subsecref{subsec:ch5_results_map}의 전수 매핑에 정리되어 있다.}



\endgroup

% --- 5.5.3 향후 연구 방향 ---
\subsection{\vrev{향후 연구 방향}}
\label{subsec:conclusion-future-merged-future}

\begingroup
  \renewcommand{\section}[1]{}
  %% ===================================================================
\section{\vrev{향후 연구 방향}}
\label{sec:conclusion-future}
%% ===================================================================

본 연구의 한계를 토대로 \jrev{5개} 향후 연구 방향을 다음과 같이 제안한다.


\chg{첫째, 실제 조직 데이터 기반 현장 검증이다.}
본 연구는 시뮬레이션 기반 검증에 초점을 맞추었으며, 실제 조직 환경에서의
현장 실험(field test)을 수반하지 않았다.
제조업, 금융업, IT 서비스 등 산업별 대표 조직을 선정하여
실제 장애 이벤트에서의 복구 시간을 수집하고,
시뮬레이션에서 예측한 $\eta$ 값과의 일치도를 검증함으로써
모델의 외적 타당성(external validity)을 확보해야 한다.
특히 관측성 지표($O(t)$)의 실시간 수집 방법론과
불확실성 추정을 위한 장애 이력 데이터베이스 구축 방안을 구체화하여,
$\DRTO$의 실무적 운용 기반을 마련할 필요가 있다.


\chg{둘째, 시간 동역학(temporal dynamics) 및 상태공간 모형으로의 확장이다.}
현행 $\DRTO$ 모델에서 $R_0$, $R_{\max}$, $\kappa$는 고정값으로 설정되어 있다.
이를 시간 변동 함수\chg{(}$R_0(t)$, $\kappa(t)$\chg{)}로 확장하여
조직의 BCM 성숙도 변화와 환경 변동을 모델에 내생적으로 반영할 수 있다.
상태 공간 모형(state-space model)의 관점에서 $\DRTO$의 시간적 진화를
형식화하면, 관측성과 불확실성의 동적 상호 작용을
시간축 위에서 추적할 수 있는 이론적 기반이 마련된다.
특히 이산 시간 상태 전이 행렬을 통해
C0$\to$C1$\to$C2(가우시안) 및 C0$\to$C1$\to$C3(파레토)의
분기 경로를 동적 체제 전환(dynamic regime switching)~\citep{hamilton1989regime}으로
모형화하는 확장이 가능하다.


\chg{셋째, 칼만 필터(Kalman filter) 통합이다.}
상태 공간 모형의 자연스러운 확장으로서, 칼만 필터를 $\DRTO$ 프레임워크에
통합하여 입력 변수($O(t)$, $U(t)$)의 노이즈를 실시간으로 여과하고
최적 상태 추정을 수행할 수 있다.
장애 이벤트 관측 시마다 예측-갱신(\erf{pict-update}{predict-update}) 사이클을 통해
$\DRTO$를 적응적으로 교정하는 자기 교정(self-calibrating) 메커니즘은,
관측성 지표의 측정 오차와 불확실성 추정의 편향을 체계적으로 보정할 수 있다.
확장 칼만 필터(EKF) 또는 언센티드 칼만 필터(UKF)를 적용하면
시그모이드 바운딩의 비선형성을 직접 처리할 수 있으며,
이는 $\DRTO$의 실시간 운용 정확도를 향상시킬 것이다.


\chg{넷째, \frev{분포 환경 적합성}의 \frev{직접 검증}이다.}
\jrev{본 연구는 \chg{\citet{lee2026drto}}의 $\kappa^{*}=1.2$를 새로운 분포 환경에 적용한 후 9개 가설(H2--H10) PASS의 통합 해석으로 \frev{운영 적합성을 시사하였다} (4.13절 \frev{분포 환경 적합성} 종합 명제). 그러나 \chg{\citet{lee2026drto}}의 4기준 종합 점수 $F(\kappa) = f_1 \cdot f_2 \cdot f_3 \cdot f_4$ (\chg{실무 유의성, 비포화, 효과 크기, 설명력})는 C1 조건 기반으로 정의되어 C2/C3 조건에서 직접 격자 탐색에 부적합하다는 학술적 한계를 내포한다. 향후 연구에서는 분포 환경 무관 다기준 종합 점수의 재정의와 C2/C3 조건에서의 직접 격자 탐색을 통해 $\kappa^{*}$의 환경별 최적값 비교를 수행함으로써, 본 연구의 종합 명제를 직접적 검증으로 보완할 수 있다.}


\chg{다섯째, 다중 조직 BCM 및 적응적 $\kappa$ 강화학습이다.}
단일 조직의 $\DRTO$에서 확장하여, 복수의 조직이 상호 의존적 공급망을
형성하는 다중 조직 수준의 $\DRTO$ 모형을 개발할 수 있다.
시스템 간 장애 전파(cascading failure) 효과, 공유 자원 경쟁,
복구 우선순위 최적화를 개별 시그모이드 바운딩의 다차원 확장을 통해 모형화한다.
나아가 강화 학습(reinforcement learning)을 적용하여
$\kappa$를 환경 변화에 따라 적응적으로 조정하는 메커니즘을 구축할 수 있다.
조직의 장애 이력과 복구 성과를 보상(reward) 신호로 활용하면,
$\kappa$가 \chg{조직과 산업, 그리고 시점}에 최적화된 값으로
자동 수렴하는 자기 최적화(self-optimizing) 프레임워크의 실현이 가능하다.

\vspace{0.5em}

\p{[}그림~\vref{fig:ch7-achievement-map}\p{]}는 본 연구의 연구 성과 구조를
문제 인식에서 해결 방안, 핵심 성과로 이어지는 흐름으로 시각화한 것이다.

\begin{figure}[htbp]
\centering
\resizebox{\textwidth}{!}{%
\begin{tikzpicture}[
    probbox/.style={draw, rounded corners=5pt, fill=red!8, minimum width=34mm,
                    minimum height=12mm, align=center, font=\footnotesize,
                    line width=0.6pt},
    solbox/.style={draw, rounded corners=5pt, fill=blue!8, minimum width=34mm,
                   minimum height=12mm, align=center, font=\footnotesize,
                   line width=0.6pt},
    outbox/.style={draw, rounded corners=5pt, fill=green!8, minimum width=34mm,
                   minimum height=12mm, align=center, font=\footnotesize,
                   line width=0.6pt},
    catbox/.style={draw, rounded corners=8pt, inner sep=6pt, line width=0.8pt},
    conn/.style={-{Stealth[length=2.5mm]}, thick, black},
    catlbl/.style={font=\small\bfseries, anchor=south},
]

% --- 문제 인식 (왼쪽) ---
\node[probbox] (p1) at (0, 3.0) {관측성 미반영\\(정적 RTO)};
\node[probbox] (p2) at (0, 1.5) {불확실성 미정량화\\(worst-case 한정)};
\node[probbox] (p3) at (0, 0.0) {극단 사건 미모형화\\(정규분포 전제)};

% --- 해결 방안 (중앙) ---
\node[solbox] (s1) at (5.5, 3.0) {개별 시그모이드\\바운딩 모델};
\node[solbox] (s2) at (5.5, 1.5) {\frev{분포 환경 적합성}\\종합 명제 도출};
\node[solbox] (s3) at (5.5, 0.0) {C0$\to$C3 점진적\\4단계 프레임워크};

% --- 핵심 성과 (오른쪽) ---
\node[outbox] (o1) at (11, 3.6) {$\DRTO \in (0, R_{\max})$\\경계 보장};
\node[outbox] (o2) at (11, 2.1) {H1--H10 전원 채택\\$R^2_{(2\mathrm{q})} \geq 0.95$};
\node[outbox] (o3) at (11, 0.6) {꼬리 감쇠 $0.52$배\\CVaR$_{99}$ $+24.0\%$};
\node[outbox] (o4) at (11, -0.9) {전문가 4.46/5.0\\$\alpha = 0.89$};

% --- 카테고리 박스 (배경 레이어) ---
\begin{scope}[on background layer]
  \node[catbox, fit=(p1)(p2)(p3), draw=red!50, fill=red!2] (pcat) {};
  \node[catbox, fit=(s1)(s2)(s3), draw=blue!50, fill=blue!2] (scat) {};
  \node[catbox, fit=(o1)(o2)(o3)(o4), draw=green!50!black, fill=green!2] (ocat) {};
\end{scope}

% --- 카테고리 라벨 ---
\node[catlbl, text=black] at (pcat.north) {문제 인식};
\node[catlbl, text=black] at (scat.north) {해결 방안};
\node[catlbl, text=black] at (ocat.north) {핵심 성과};

% --- 연결선 ---
\draw[conn] (p1.east) -- (s1.west);
\draw[conn] (p2.east) -- (s2.west);
\draw[conn] (p3.east) -- (s3.west);

\draw[conn] (s1.east) -- (o1.west);
\draw[conn] (s2.east) -- (o2.west);
\draw[conn] (s3.east) -- (o3.west);
\draw[conn, dashed] (s1.east) -- (o2.west);
\draw[conn, dashed] (s3.east) -- (o4.west);

\end{tikzpicture}%
}% end resizebox
\caption{문제 인식에서 해결 방안, 핵심 성과로 이어지는 연구 성과 구조도}
\label{fig:ch7-achievement-map}
\end{figure}


\vspace{1.5em}

결론적으로, 본 연구는 \erf{BCMgownj}{BCM} 분야에서 정적 RTO의 근본적 한계를
극복하는 $\DRTO$ 프레임워크를 제안하고,
개별 시그모이드 바운딩 모델\jrev{(\chg{\citet{lee2026drto}}의 $\kappa^{*} = 1.2$ 전제 채택)},
꼬리 감쇠 발견($0.52$배), P85.8 CDF 교차점 규명(CVaR$_{99}$ $+24.0\%$),
6단계 다층 검증(전문가 4.46/5.0, \jrev{10/10} 가설 채택)\jrev{,
그리고 \textbf{\jrev{\frev{분포 환경 적합성}} 종합 명제 도출}}이라는
\chg{이론적이고 방법론적인} 기여를 달성하였다.
향후 실제 조직 데이터 기반 현장 검증, 상태 공간 모형과 칼만 필터의 통합,
\jrev{\frev{분포 환경 적합성}의 \frev{직접 검증}(환경 무관 다기준 종합 점수 재정의),}
다중 조직 BCM 확장, 그리고 적응적 $\kappa$ 강화 학습을 통해
$\DRTO$ 프레임워크가 조직의 업무연속성 역량을 실질적으로 강화하는
정량적 도구로 발전하기를 기대한다.


\endgroup
             % 5.5 결론 및 향후 연구
    % 제5장 논의 및 결론

% ─── 참고문헌 (제출본 기준: 부록 앞) ───
\nocite{iso22301,iso22313,iso31000}

% --- 참고문헌: [국내 문헌] / [국외 문헌] 분리 (송은정 박사 스타일) ---
% 국내 문헌 키워드 필터
\defbibcheck{korean}{%
  \iffieldundef{keywords}{\skipentry}{}%
  \ifkeyword{korean}{}{\skipentry}%
}
\defbibcheck{foreign}{%
  \ifkeyword{korean}{\skipentry}{}%
}

\clearpage
\phantomsection
\addcontentsline{toc}{chapter}{\PlainTOCtext{참고문헌}}
\begingroup
\setstretch{1.6}
\begin{center}
{\ChapterFont\bfseries\fontsize{16}{22}\selectfont 참고문헌}
\end{center}
\vspace{18pt}

% [국내 문헌]
\noindent{\JungGothic\bfseries\fontsize{13}{18}\selectfont [국내 문헌]}
\vspace{6pt}
\printbibliography[heading=none, check=korean]

\vspace{12pt}
% [국외 문헌]
\noindent{\JungGothic\bfseries\fontsize{13}{18}\selectfont [국외 문헌]}
\vspace{6pt}
\printbibliography[heading=none, check=foreign]
\endgroup

% ─── 부록 (제출본 기준: 참고문헌 뒤, 제A·B·C장) ───
\appendix
\setlength{\abovedisplayskip}{6pt}
\setlength{\belowdisplayskip}{6pt}
\setlength{\abovedisplayshortskip}{3pt}
\setlength{\belowdisplayshortskip}{3pt}
\titleformat{\chapter}[hang]
  {\centering\ChapterFont\bfseries\fontsize{16}{22}\selectfont}
  {제~\thechapter~장}{0.8em}{}
\titlespacing*{\chapter}{0pt}{0pt}{18pt}
% =====================================================================
% appendices.tex -- 박사학위 논문 부록
% 참조: doctoral/research_model/research_model.tex
% 빌드: main_doc.tex에서 \appendix 후 % =====================================================================
% appendices.tex -- 박사학위 논문 부록
% 참조: doctoral/research_model/research_model.tex
% 빌드: main_doc.tex에서 \appendix 후 % =====================================================================
% appendices.tex -- 박사학위 논문 부록
% 참조: doctoral/research_model/research_model.tex
% 빌드: main_doc.tex에서 \appendix 후 \input{appendix/appendices}
% 핵심 파라미터: κ = 1.2, R₀ = 6.0h, R_max = 24.0h (= 4R₀)
% 모델: 개별 시그모이드 바운딩, DRTO = R₀ + E_{O,bnd} + E_{U,bnd}
% =====================================================================

% =====================================================================
% 부록 A: 기호·용어·수식 목록
% =====================================================================
\clearpage
% ── 부록 간지 ──
\thispagestyle{plain}
\phantomsection
\addcontentsline{toc}{chapter}{\PlainTOCtext{부록}}
\vspace*{0.35\textheight}
\begin{center}{\ChapterFont\bfseries\fontsize{16}{22}\selectfont 부록}\end{center}
\clearpage

\chapter{기호·용어·수식 목록}
\label{app:notation-list}

\SetAppendixLetter{A}

본 부록은 본 연구에서 사용되는 주요 기호, 용어, 수식을 일람표로 정리한다.

\phantomsection\label{app:notation-overview}
\AppSection{주요 기호}
\label{app:symbols}

[표~A-1]은 주요 기호, 용어, 수식의 일람표이다.

\small
\begin{longtable}{c >{\raggedright\arraybackslash}p{4.8cm} >{\raggedright\arraybackslash}p{8.5cm}}
\toprule
\textbf{기호} & \textbf{명칭} & \textbf{정의 / 값} \\
\midrule
\endhead

\multicolumn{3}{l}{\textbf{입력 변수 및 매개변수}} \\
\midrule
$R_0$           & 기준 RTO (Base RTO)      & $6.0$\,h (BIA 기반 정적 복구 목표) \\
$R_{\max}$      & 최대 허용 RTO (MTPD)     & $24.0$\,h $= 4R_0$ \\
$O_{\text{raw}}$ & 관측성 원시값 (Observability raw value) & $\in [0,\,1]$, 0=블라인드, 1=완전 가시 \\
$k_O$           & 관측성 가중치 (Observability weight) & $\in [0,\,1]$, 관측성 민감도 \\
$k_U$           & 불확실성 가중치           & $= 1 - k_O$ (상보적 제약) \\
$\kappa$        & 곡률 매개변수 (Curvature parameter) & $> 0$, 시그모이드 곡률 \\
$\kappa^*$      & \jrev{\citet{lee2026drto} 도출 곡률 (본 논문 전제)} & $1.2$ \jrev{(선행연구 C1 환경 $F(\kappa)$ 최대화 결과)} \frev{— \emph{잠정·보수적 실무 기본값}} \\
$U_{\text{raw}}$ & 불확실성 원시값 (Uncertainty raw value) & C2: $\sim N(3,1)$, C3: $\sim 6 \cdot \text{Pareto}(3)$. 무차원 강도 \\
$\alpha$        & 파레토 형상 모수          & $3$ (유한 평균/분산, 무한 첨도) \\
\midrule

\multicolumn{3}{l}{\textbf{함수 및 파생 인자}} \\
\midrule
$S(z)$          & 수정 시그모이드 (Sigmoid function) & $2/(1+e^{-z}) - 1 \in (-1,\,1)$ \\
$z_O$           & 시그모이드 관측성 인자    & $\kappa \cdot k_O \cdot O_{\text{raw}} \geq 0$ \\
$z_U$           & 시그모이드 불확실성 인자   & $\kappa \cdot k_U \cdot R_0 \cdot U_{\text{raw}} / (R_{\max} - R_0) \geq 0$ (무차원) \\
$E_{O,\text{raw}}$ & 관측성 원시 효과 (Raw observability effect) & $-R_0 \cdot k_O \cdot O_{\text{raw}}$ \\
$E_{U,\text{raw}}$ & 불확실성 원시 효과 (Raw uncertainty effect) & $(R_{\max} - R_0) \cdot k_U \cdot U_{\text{raw}}$ \\
$E_{O,\text{bnd}}$ & 관측성 바운딩 조정량 (Bounded obs.\ effect) & $-R_0 \cdot S(z_O) \in (-R_0,\,0]$ \\
$E_{U,\text{bnd}}$ & 불확실성 바운딩 조정량 (Bounded unc.\ effect) & $(R_{\max} - R_0) \cdot S(z_U) \in [0,\,R_{\max}-R_0)$ \\
\midrule

\multicolumn{3}{l}{\textbf{출력 변수}} \\
\midrule
$\DRTO$         & 동적 복구목표시간 (Dynamic RTO) & $R_0 + E_{O,\text{bnd}} + E_{U,\text{bnd}} \in (0,\,R_{\max})$ \\
$\eta$          & 효율성 비율 (Efficiency ratio) & $\DRTO / R_0 \in (0,\,4)$ \\
$\bar{\eta}$    & 효율성 비율 평균              & 각 조건별 $\eta$ 표본 평균 \\
$\text{SD}(\eta)$ & 효율성 비율 표준편차         & 각 조건별 $\eta$ 표본 표준편차 \\
\midrule

\multicolumn{3}{l}{\textbf{최적화 기준 함수}} \\
\midrule
$F(\kappa)$     & 다기준 종합점수 (Composite score) & $f_1 \cdot f_2 \cdot f_3 \cdot f_4$ \\
$f_1(\kappa)$   & 실무 유의성 (Practical significance) & 목표 감소율 15\% 기준 \jrev{(선행연구 C1 환경)} \\
$f_2(\kappa)$   & 비포화 (Non-saturation)  & 포화율 $< 5\%$ \jrev{(선행연구 C1 환경)} \\
$f_3(\kappa)$   & 효과 크기 (Effect size)  & $|d| \geq 0.8$ \jrev{(선행연구 C1 환경)} \\
$f_4(\kappa)$   & 설명력 (Explanatory power) & $R^2 \geq 0.95$ \jrev{(선행연구 C1 환경)} \\
\midrule

\multicolumn{3}{l}{\textbf{통계 기호}} \\
\midrule
$n$             & 표본 크기 (Sample size)   & $10{,}000$ \\
$R^2$           & 결정계수                & 회귀 모형 설명력 \\
$R^2_{(1\mathrm{s})}$ & 1차 단순회귀 결정계수 & $\hat{\eta}^{(1\mathrm{s})}$ 모형의 $R^2$ \\
$R^2_{(1\mathrm{m})}$ & 1차 다중회귀 결정계수 & $\hat{\eta}^{(1\mathrm{m})}$ 모형의 $R^2$ \\
$R^2_{(2\mathrm{q})}$ & 2차교호작용 결정계수 & $\hat{\eta}^{(2\mathrm{q})}$ 모형의 $R^2$ \\
$\Delta R^2$    & 결정계수 변화량         & 회귀 단계 간 $R^2$ 증가분 \\
$d$             & Cohen's $d$            & $(\bar{X}_1 - \bar{X}_2) / s_p$ \\
$\text{VaR}_q$  & 위험가치 (Value-at-Risk) & $F_\eta^{-1}(q/100)$, $\eta$의 $q$-백분위값 \\
$\text{CVaR}_q$ & 조건부 꼬리위험          & $E[\eta \mid \eta > \text{VaR}_q]$ \\
\midrule

\multicolumn{3}{l}{\textbf{모델 조건}} \\
\midrule
C0              & 기저 조건 (Baseline)     & $O_{\text{raw}} = 0$, 정적 RTO ($\eta = 1$) \\
C1              & 관측성 조건 (Observability) & $O_{\text{raw}} > 0$, 단일 채널 \\
C2              & 가우시안 불확실성 조건    & C1 + $U \sim N(3,1)$, 가우시안(경꼬리) \\
C3              & 파레토 불확실성 조건      & C1 + $U \sim 6 \cdot \text{Pareto}(3)$, 파레토(중꼬리) \\
\midrule

\multicolumn{3}{l}{\textbf{가설}} \\
\midrule
H1--\jrev{H3}          & L1 시뮬레이션 검증       & H1: 관측성 효과, H2: 경계 보장, H3: 조건 순서 \\
\jrev{H4}--\jrev{H5}          & L2 회귀 검증             & \jrev{H4}: $R^2_{(2\mathrm{q})} \geq 0.95$, \jrev{H5}: 교호작용 $\Delta R^2$ 유의 \\
\jrev{H6}--\jrev{H7}          & L3 강건성 검증           & \jrev{H6}: 시드 독립성 ($|r| < 0.02$), \jrev{H7}: 다중 시드 강건성 ($10^8$회) \\
\jrev{\pdferr{H8}}              & L4 이중채널 검증         & Tail/Core $k_O$ 감쇠비 $< 1$ (P85.8 분할) \\
\jrev{H9}--\jrev{H10}        & L5 전문가 평가           & \jrev{H9}: 정량 평균 $\geq 4.0/5.0$, \jrev{H10}: 정성 동의율 $\geq 80\%$ \\
\midrule

\multicolumn{3}{l}{\textbf{실무 확장 기호 (Practical Extension)}} \\
\midrule
$O_{\mathrm{APM}}$  & APM 관측성              & $w_1 R_{\mathrm{avail}} + w_2 (1-R_{\mathrm{error}}) + w_3 R_{\mathrm{perf}} \in [0,1]$ \\
$O_{\mathrm{SIEM}}$ & SIEM 관측성             & 보안 이벤트 심각도·대응 상태 기반, $\in [0,1]$ \\
$R_{\mathrm{avail}}$ & 가용률 (Availability rate) & 서비스 가용성 비율, $\in [0,1]$ \\
$R_{\mathrm{error}}$ & 오류율 (Error rate)     & 서비스 오류 비율, $\in [0,1]$. $1-R_{\mathrm{error}}$로 반전 사용 \\
$R_{\mathrm{perf}}$  & 성능 지수 (Performance rate) & 서비스 응답 성능 비율, $\in [0,1]$ \\
$w_i$               & APM 가중치              & $O_{\mathrm{APM}}$ 구성 요소 가중 계수 ($\sum w_i = 1$) \\
$\alpha_{\mathrm{obs}}$ & 관측성 혼합비        & $O(t) = \alpha_{\mathrm{obs}} O_{\mathrm{APM}} + (1-\alpha_{\mathrm{obs}}) O_{\mathrm{SIEM}}$, $\in [0.5, 0.8]$ \\
$T_{\mathrm{actual}}$   & 실제 복구 시간 (Actual recovery time) & $> 0$ (h), 백테스팅 비교 대상 \\
$\DRTO_{\mathrm{final}}$ & 최종 DRTO 예측값      & $(0, R_{\max})$ (h), 장애 시점 산출값 \\
$E$                     & 절대 오차 (Absolute error) & $T_{\mathrm{actual}} - \DRTO_{\mathrm{final}}$ (h) \\
$E_r$                   & 상대 오차 (Relative error) & $E / T_{\mathrm{actual}}$, 무차원 \\
$R_0(t)$                & 시변 기준 RTO (Time-varying base RTO) & 정적 $R_0$의 시계열 확장, $> 0$ (h) \\
$\kappa(t)$             & 시변 곡률 매개변수 (Time-varying curvature) & 정적 $\kappa = 1.2$의 시계열 확장, $> 0$ \\

\bottomrule
\end{longtable}


\AppSection{핵심 수식}
\label{app:key-formulas}

\begin{enumerate}[leftmargin=2em]

\item \textbf{DRTO 개별 시그모이드 바운딩 모델}:
\[
  \DRTO = R_0 + E_{O,\text{bnd}} + E_{U,\text{bnd}}
  = R_0 - R_0 \cdot S(z_O) + (R_{\max} - R_0) \cdot S(z_U)
\]

\item \textbf{수정 시그모이드 함수}:
\[
  S(z) = \frac{2}{1 + e^{-z}} - 1, \quad S(z) \in (-1,\,1), \quad S(0) = 0
\]

\item \textbf{관측성 인자}: $z_O = \kappa \cdot k_O \cdot O_{\text{raw}}$

\item \textbf{불확실성 인자}: $z_U = \kappa \cdot k_U \cdot \dfrac{R_0 \cdot U_{\text{raw}}}{R_{\max} - R_0}$

\item \textbf{다기준 종합점수}:
$F(\kappa) = f_1(\kappa) \cdot f_2(\kappa) \cdot f_3(\kappa) \cdot f_4(\kappa)$,
\quad $F^* = F(1.2) = 0.9997$ \jrev{(\citet{lee2026drto} C1 환경 결과; 본 논문은 \frev{분포 환경 적합성} 종합 명제로 보완 — \secref{sec:ch4_hypothesis_summary} 참조)}

\item \textbf{합의 4변수 간명 모형} (Forward Selection top-4 $\cap$ LASSO):
\begin{itemize}
  \item C2: $\hat{\eta} = 1.225 - 0.291\,k_O + 0.467\,U - 0.433\,(k_O \!\cdot\! U) - 0.562\,(k_O \!\cdot\! O)$\quad ($R^2 = 0.996$)
  \item C3: $\hat{\eta} = 1.574 + 0.240\,U - 1.053\,k_O - 0.002\,U^2 - 0.122\,(k_O \!\cdot\! U)$\quad ($R^2 = 0.839$)
\end{itemize}

\end{enumerate}


\AppSection{핵심 용어}
\label{app:key-terms}

\begin{longtable}{p{3.5cm} p{4cm} p{7cm}}
\toprule
\textbf{한국어} & \textbf{영어} & \textbf{정의 요약} \\
\midrule
\endhead
동적 복구목표시간 & Dynamic RTO & 관측성·불확실성을 반영한 $\DRTO$ \\
관측성 & Observability & 시스템 내부 상태를 외부 출력으로 추정 가능한 정도 \\
불확실성 (가우시안) & Uncertainty (Gaussian) & 정규분포를 따르는 일상적 변동, $U_{\mathrm{raw}} \sim N(3,\,1)$ \\
불확실성 (파레토) & Uncertainty (Pareto) & 파레토 분포를 따르는 극단적 변동, $U_{\mathrm{raw}} \sim 6 \cdot \mathrm{Pareto}(\alpha=3)$ \\
개별 시그모이드 바운딩 & Individual Sigmoid Bounding & 채널별 독립 시그모이드로 물리적 경계 보장 \\
이중채널 감쇠 & Dual-channel Attenuation & 극단적 불확실성(Tail 영역)에서 불확실성 채널이 지배적으로 작용하여 관측성 채널의 복구 단축 효과가 구조적으로 점진적 감쇠되는 현상 ($k_O$ 감쇠비 $0.52\times$) \\
과다예비 & Over-provisioning & 불확실성 과대 추정으로 필요 이상의 복구 자원을 배정 \\
과소예비 & Under-provisioning & 불확실성 과소 추정으로 복구 자원이 부족하게 배정 \\
분할 회귀 & Split Regression & P85.8 기준 Core/Tail 분리 회귀 분석 \\
CDF 교차점 & CDF Crossover & C2=C3 누적분포 교차 지점 (P85.8, $\eta \approx 2.42$) \\
효율성 비율 & Efficiency Ratio ($\eta$) & $\DRTO / R_0$, 기준 대비 비율 \\
불확실성 프리미엄 & Uncertainty Premium & 불확실성에 의한 DRTO 추가 증가분 \\
에스컬레이션 & Escalation & $\eta$ 수준별 단계적 대응 격상 \\
\midrule
\multicolumn{3}{l}{\textbf{산업별 관측성 요소}} \\
\midrule
트랜잭션 모니터링 & Transaction Monitoring & 금융 거래를 실시간 감시하여 이상 징후를 탐지하는 시스템 \\
FDS & Fraud Detection System & 이상거래탐지시스템. 금융 사기·부정거래를 실시간 패턴 분석으로 탐지 \\
EMR & Electronic Medical Record & 전자의무기록. 환자 진료 정보를 전자적으로 관리하는 시스템 \\
MES & Manufacturing Execution System & 제조실행시스템. 생산 현장의 작업 지시·실적·품질을 실시간 관리 \\
SCADA & Supervisory Control and Data Acquisition & 감시제어 데이터수집 시스템. 산업 설비·인프라의 원격 감시·제어 \\
APM & Application Performance Monitoring & 애플리케이션 응답 시간, 오류율 등을 실시간 모니터링하는 시스템 \\
SIEM & Security Information and Event Management & 보안 이벤트 수집·분석·대응 통합 관리. $O_{\mathrm{SIEM}}$ 산출 기반 \\
{IDS} & {Intrusion Detection System} & {침입탐지시스템. 네트워크·호스트의 비정상 활동을 탐지·경고} \\
{APT} & {Advanced Persistent Threat} & {지능형 지속 위협. 장기간 은닉하며 특정 조직을 표적으로 공격하는 고도화된 사이버 위협} \\
\midrule
\multicolumn{3}{l}{{\textbf{업무연속성·재해복구}}} \\
\midrule
{BCM} & {Business Continuity Management} & {업무연속성관리. 중단 상황에서 핵심 업무를 유지·복구하기 위한 관리 체계} \\
{BCP} & {Business Continuity Plan} & {업무연속성계획. 업무 중단 시 복구 절차를 문서화한 계획서} \\
{BCMS} & {Business Continuity Management System} & {업무연속성관리체계. \nrev{ISO 22301:2019}에 따른 BCM 구현·운영 통합 시스템} \\
{BIA} & {Business Impact Analysis} & {업무영향분석. 핵심 업무의 중단 영향도를 분석하여 RTO, RPO 등 복구 목표를 결정하는 프로세스} \\
{DR} & {Disaster Recovery} & {재해복구. 재난 발생 후 IT 시스템 및 업무 기능을 원래 상태로 복원하는 활동} \\
{MTPD} & {Maximum Tolerable Period of Disruption} & {최대허용중단시간. 업무 중단이 허용되는 최대 기간 ($= R_{\max}$)} \\
\midrule
\multicolumn{3}{l}{{\textbf{리스크 관리·금융}}} \\
\midrule
{ERM} & {Enterprise Risk Management} & {전사적 위험관리. 조직 전체 위험을 통합적으로 식별·평가·관리하는 체계} \\
{KRI} & {Key Risk Indicator} & {주요 위험 지표. $\eta$가 임계값 초과 시 경보를 발생시키는 ERM 연계 지표} \\
{VaR} & {Value-at-Risk} & {위험가치. 신뢰수준(예: 99\%)에서 최대 손실액을 나타내는 금융 위험 지표} \\
{CVaR} & {Conditional Value-at-Risk} & {조건부 위험가치. VaR 초과 극단 손실의 평균. 꼬리 위험 평가 지표} \\
{DaR} & {DRTO-at-Risk} & {VaR 프레임을 DRTO에 적용. ``99\% 신뢰수준에서 DRTO가 X시간 이하'' (E3 제안)} \\
{Basel} & {Basel Accord / Basel Committee} & {바젤 위원회 규제. 은행 자기자본·운영 위험 기준. DRTO를 운영위험 정량 지표로 활용 가능} \\
\midrule
\multicolumn{3}{l}{{\textbf{규제·표준}}} \\
\midrule
{\nrev{ISO 22301:2019}} & {ISO 22301:2019} & {업무연속성관리체계 국제표준. BIA, RTO, 복구 전략 등을 규정} \\
{NIST CSF} & {NIST Cybersecurity Framework} & {NIST 사이버보안 프레임워크(v2.0). 식별·보호·탐지·대응·복구 5기능. DRTO는 복구 기능 정량 지표} \\
{DORA} & {Digital Operational Resilience Act} & {EU 디지털 운영 탄력성 규제. 금융기관 사이버 복원력 의무화 (2023년 발효)} \\
\midrule
\multicolumn{3}{l}{{\textbf{사이버 보안·위협 인텔리전스}}} \\
\midrule
{DDoS} & {Distributed Denial of Service} & {분산 서비스 거부 공격. 대량 트래픽으로 서비스 접근을 차단하는 공격 유형} \\
{MITRE ATT\&CK} & {Adversarial Tactics, Techniques \& Common Knowledge} & {사이버 공격 기법·전술 분류 체계. 공격 유형별 $U_{\mathrm{raw}}$ 자동 조정에 활용} \\
{STIX} & {\erf{Structu}{Structured} Threat Information Expression} & {사이버 위협 정보 표준 교환 형식. TAXII와 함께 위협 인텔리전스 자동 연동} \\
{TAXII} & {Trusted Automated Exchange of Intelligence Information} & {STIX 형식 위협 정보의 자동 교환 프로토콜} \\
\midrule
\multicolumn{3}{l}{{\textbf{클라우드·DevOps}}} \\
\midrule
{GitOps} & {Git-based Operations} & {Git 리포지토리를 SSOT로 인프라·서비스 변경을 코드 기반 자동화. DRTO 파라미터 버전 관리 (E4 제안)} \\
{IaC} & {Infrastructure as Code} & {인프라를 코드로 정의·관리 (Terraform, Ansible 등). GitOps와 연동하여 인프라 변경 시 DRTO 영향 자동 평가} \\
{Kubernetes (K8s)} & {Container Orchestration Platform} & {컨테이너 오케스트레이션 플랫폼. Pod·Service·Cluster 단위별 DRTO 파라미터 차별화 필요} \\
{CI/CD} & {Continuous Integration / Continuous Deployment} & {지속적 통합·배포. 코드 변경의 자동 빌드·테스트·배포 파이프라인} \\
{Serverless} & {Serverless Computing} & {서버리스 컴퓨팅. 서버 관리 없이 코드 실행 단위로 자원을 할당하는 클라우드 모델. $R_0$ 정의 방법 차별화 필요} \\
YAML & \vrev{YAML} & 사람이 읽기 쉬운 데이터 직렬화 언어. GitOps에서 DRTO 파라미터 정의에 활용 \\
\midrule
\multicolumn{3}{l}{{\textbf{연구 방법론·통계}}} \\
\midrule
{CVI} & {Content Validity Index} & {내용타당도 지수. I-CVI(문항 수준, $\geq 0.80$), S-CVI/Ave(척도 수준, $\geq 0.90$)} \\
{ICC} & {Intraclass Correlation Coefficient} & {급내상관계수. 다수 평가자 간 일치도 측정 통계량} \\
\bottomrule
\end{longtable}

%% --- 핵심 용어 상세 설명 및 사례 ---




%%==APPB_START==
\chapter{조직 성숙도 자가 진단 체크리스트}
\label{app:eta-dual-use-detail}
\label{app:self-assessment}

\SetAppendixLetter{B}

현업 담당자가 조직의 $\eta$ 수준을 간접적으로 추정할 수 있는
15개 항목 자가 진단 체크리스트를 \p{[}표~\ref{tab:app-self-assessment}\p{]}에 제시한다.

\begin{table}[htbp]
\centering
\caption{BCMS 성숙도 자가 진단 체크리스트 (15개 항목)}
\label{tab:app-self-assessment}
\small
\begin{tabularx}{\textwidth}{c X c}
\toprule
\textbf{No.} & \textbf{진단 항목} & \textbf{영역} \\
\midrule
\multicolumn{3}{l}{\textbf{관측성 역량} (체크 수 $\uparrow$ $\to$ $\eta$ $\downarrow$)} \\
\midrule
1 & 주요 시스템의 실시간 모니터링 도구가 운영 중인가? & $O$ \\
2 & 장애 발생 시 자동 알림(alerting) 체계가 작동하는가? & $O$ \\
3 & MTTR(평균복구시간)을 정기적으로 측정·추적하고 있는가? & $O$ \\
4 & 장애 원인 분석(RCA)을 체계적으로 수행하는가? & $O$ \\
5 & 로그·메트릭·트레이스 중 2개 이상의 관측성 축을 확보했는가? & $O$ \\
6 & 관측성 투자에 대한 정량적 효과 측정을 수행하는가? & $k_O$ \\
\midrule
\multicolumn{3}{l}{\textbf{불확실성 관리} (체크 수 $\uparrow$ $\to$ $\eta$ $\downarrow$)} \\
\midrule
7 & 복구 절차서(runbook)가 문서화되어 최신 상태로 유지되는가? & $U$ \\
8 & 연 2회 이상 BCM 훈련/모의 복구를 실시하는가? & $U$ \\
9 & 외부 공급업체 \erf{장애시}{장애 시} 대체 방안(plan~B)이 수립되어 있는가? & $U$ \\
10 & 극단 시나리오(블랙 스완)에 대한 대응 계획이 있는가? & $U$ \\
11 & 복구 시간 편차를 추적하고 원인별로 분류하는가? & $U$ \\
12 & 사후 분석(post-mortem) 결과를 체계적으로 반영하는가? & $U$ \\
\midrule
\multicolumn{3}{l}{\textbf{거버넌스}} \\
\midrule
13 & BCM 전담 조직 또는 책임자가 지정되어 있는가? & 조직 \\
14 & BCM KPI가 경영진에게 정기적으로 보고되는가? & 조직 \\
15 & BIA(업무영향분석)를 정기적으로 갱신하는가? & 조직 \\
\bottomrule
\end{tabularx}
\end{table}


% 부록 I: 전문가 검증 상세 (기존 부록 H)


%%==APPC_START==
\chapter{전문가 검증 상세 설문 구성}
\label{app:survey-structure}

\SetAppendixLetter{C}

\AppSection{설문 구조 개요}
\label{app:survey-overview}

{전문가 평가 설문은 4개 리커트 영역(Section B--E)과 1개 개방형 영역(Section F),
총 32문항으로 구성되었다(\p{[}표~\ref{tab:app-survey-structure}\p{]}).
정량 분석 대상은 리커트 5점 척도 27문항(B--E)이다.}

\begin{table}[htbp]
\centering
\caption{전문가 평가 설문 구조 (4개 리커트 영역 + 개방형, 32문항)}
\label{tab:app-survey-structure}
\small
\begin{tabular}{c l >{\raggedright\arraybackslash}p{3.8cm} >{\raggedright\arraybackslash}p{4.5cm} r}
\toprule
\textbf{영역} & \textbf{코드} & \textbf{명칭} & \textbf{평가 내용} & \textbf{문항 수} \\
\midrule
Section B & B1--B6  & 모델 타당성      & 이론적 논리성, 구조적 타당성              & 6 \\
Section C & C1--C6  & 파라미터 적정성   & 권장 범위의 실무 적합성                    & 6 \\
{Section D} & {D1--D6}  & {실무 적용성}      & {데이터 수집, BCM 통합, 도입 의향}           & {6} \\
{Section E} & {E1--E9}  & {산업별 적용 가능성 및 $\eta$ 활용 체계} & {산업 범용성, 에스컬레이션, 성숙도, 환류} & {9} \\
\midrule
          &         & {\textbf{리커트 소계}} &                                       & {\textbf{27}} \\
\midrule
{Section F} & {F1--F5}  & {개방형 문항}      & {강점, 장벽, 제안, 보완, 기타 의견}          & {5} \\
\midrule
          &         & {\textbf{전체 합계}}  &                                        & {\textbf{32}} \\
\bottomrule
\end{tabular}
\vspace{0.5em}

\footnotesize
{\noindent 리커트 척도 적용 문항: B1--E9의 27문항. 개방형 문항(F1--F5)은 별도 질적 분석 대상.} \\
\noindent 평가 척도: 1점(전혀 동의하지 않음) -- 5점(매우 동의함).
\normalsize
\end{table}

\AppSection{설문 내용}
\label{app:survey-items}

\AppSubsection{Section B: 모델 타당성 (6문항)}
\label{app:survey-b}

$\DRTO$ 모델의 이론적 구조와 논리적 타당성을 평가하는 6개 문항이다(\p{[}표~\ref{tab:app-survey-items-b}\p{]}).

\begin{table}[htbp]
\centering
\caption{Section B: 모델 타당성 문항}
\label{tab:app-survey-items-b}
\small
\begin{tabularx}{\textwidth}{cX}
\toprule
\textbf{문항} & \textbf{내용} \\
\midrule
B1 & 관측성($O$)이 높아지면 $\DRTO$가 감소한다는 관계의 논리적 타당성 \\
B2 & 불확실성($U$)이 높아지면 $\DRTO$가 증가한다는 관계의 논리적 타당성 \\
B3 & 시그모이드 함수를 활용한 바운딩 메커니즘의 적절성 \\
B4 & C0--C3 조건 구분의 명확성 및 점진적 복잡도 증가의 유용성 \\
B5 & $k_O + k_U = 1$ 제약 조건의 해석 용이성 \\
B6 & $\DRTO$ 프레임워크의 전반적 이론적 타당성 \\
\bottomrule
\end{tabularx}
\end{table}

\AppSubsection{Section C: 파라미터 적정성 (6문항)}
\label{app:survey-c}

$\DRTO$ 모델의 핵심 매개변수 설정값과 실무 조정 가능성을 평가하는 6개 문항이다(\p{[}표~\ref{tab:app-survey-items-c}\p{]}).

\begin{table}[htbp]
\centering
\caption{Section C: 파라미터 적정성 문항}
\label{tab:app-survey-items-c}
\small
\begin{tabularx}{\textwidth}{cX}
\toprule
\textbf{문항} & \textbf{내용} \\
\midrule
C1 & $k_O \in [0.4, 0.6]$ 권장 범위의 실무 적용 가능성 \\
C2 & $k_U = 1 - k_O$ 연동 설계의 파라미터 관리 단순화 기여 \\
C3 & 시그모이드 바운딩 자연 경계 $(0,\; R_{\max})$의 실무 적절성 \\
C4 & 곡률 파라미터 $\kappa = 1.2$ 설정의 적절성 \\
C5 & 파레토 분포($\alpha = 3.0$) 꼬리 위험 반영 적절성 \\
C6 & 파라미터 구조의 직관성 및 조직 맞춤 조정 용이성 \\
\bottomrule
\end{tabularx}
\end{table}

\AppSubsection{Section D: 실무 적용성 (6문항)}
\label{app:survey-d}

현업 환경에서의 데이터 수집 가능성, 기존 프로세스 통합, 도입 의향을 평가하는 6개 문항이다(\p{[}표~\ref{tab:app-survey-items-d}\p{]}).

\begin{table}[htbp]
\centering
\caption{Section D: 실무 적용성 문항}
\label{tab:app-survey-items-d}
\small
\begin{tabularx}{\textwidth}{cX}
\toprule
\textbf{문항} & \textbf{내용} \\
\midrule
D1 & 현업에서 관측성 및 불확실성 데이터 수집 가능성 \\
D2 & 기존 BCM/BIA 프로세스와의 통합 가능성 \\
D3 & $\DRTO$ 결과의 경영진 보고 및 의사결정 활용 적합성 \\
D4 & 위험 기반 의사결정 지원의 유용성 \\
D5 & 단계적 도입(C0 $\to$ C3)을 전제로 한 구현 복잡도 수용성 \\
D6 & 적절한 지원 제공 시 실제 업무 환경 도입 의향 \\
\bottomrule
\end{tabularx}
\end{table}

\AppSubsection{Section E: 산업별 적용 가능성 및 $\eta$ 활용 체계 (9문항)}
\label{app:survey-e}

{$\DRTO$ 모델의 산업별 적용 가능성과 효율성 비율($\eta$) 기반 활용 체계를 평가하는
9개 문항이다(\p{[}표~\ref{tab:app-survey-items-e}\p{]}).}

\begin{table}[htbp]
\centering
\caption{Section E: 산업별 적용 가능성 및 $\eta$ 활용 체계 문항}
\label{tab:app-survey-items-e}
\small
\begin{tabularx}{\textwidth}{cX}
\toprule
\textbf{문항} & \textbf{내용} \\
\midrule
{E1} & {산업별 $R_0$/$R_{\max}$ 설정의 현실성} \\
{E2} & {산업별 관측성 요소의 적절성} \\
{E3} & {산업별 불확실성 요소의 적절성} \\
{E4} & {$\eta$ 기반 에스컬레이션 체계의 활용성} \\
{E5} & {산업 범용성} \\
{E6} & {가이드라인의 실무 도입 지원 효과} \\
{E7} & {$\eta$ 기반 BCMS 성숙도 5단계 모델의 활용성} \\
{E8} & {$\eta$의 평상시·위기 시 이중 활용 체계} \\
{E9} & {진단--향상--대응--환류 선순환 구조} \\
\bottomrule
\end{tabularx}
\end{table}

\AppSubsection{Section F: 개방형 문항 (5문항)}
\label{app:survey-f}

{$\DRTO$ 모델에 대한 자유 서술형 5개 문항이다. 별도 질적 분석 대상으로,
리커트 정량 분석에는 포함되지 않는다.}
\nref{5개 개방형 문항의 상세 내용은 \tabref{tab:app-survey-items-f}에 제시하였다.}

\begin{table}[htbp]
\centering
\caption{Section F: 개방형 문항}
\label{tab:app-survey-items-f}
\small
\begin{tabularx}{\textwidth}{cX}
\toprule
\textbf{문항} & \textbf{내용} \\
\midrule
{F1} & {$\DRTO$ 모델의 가장 큰 강점} \\
{F2} & {실무 도입 시 예상되는 가장 큰 장벽이나 어려움} \\
{F3} & {실무 적용성을 높이기 위해 추가로 필요한 것} \\
{F4} & {모델 보완이 필요한 부분} \\
{F5} & {기타 의견이나 제안 사항} \\
\bottomrule
\end{tabularx}
\end{table}







% 핵심 파라미터: κ = 1.2, R₀ = 6.0h, R_max = 24.0h (= 4R₀)
% 모델: 개별 시그모이드 바운딩, DRTO = R₀ + E_{O,bnd} + E_{U,bnd}
% =====================================================================

% =====================================================================
% 부록 A: 기호·용어·수식 목록
% =====================================================================
\clearpage
% ── 부록 간지 ──
\thispagestyle{plain}
\phantomsection
\addcontentsline{toc}{chapter}{\PlainTOCtext{부록}}
\vspace*{0.35\textheight}
\begin{center}{\ChapterFont\bfseries\fontsize{16}{22}\selectfont 부록}\end{center}
\clearpage

\chapter{기호·용어·수식 목록}
\label{app:notation-list}

\SetAppendixLetter{A}

본 부록은 본 연구에서 사용되는 주요 기호, 용어, 수식을 일람표로 정리한다.

\phantomsection\label{app:notation-overview}
\AppSection{주요 기호}
\label{app:symbols}

[표~A-1]은 주요 기호, 용어, 수식의 일람표이다.

\small
\begin{longtable}{c >{\raggedright\arraybackslash}p{4.8cm} >{\raggedright\arraybackslash}p{8.5cm}}
\toprule
\textbf{기호} & \textbf{명칭} & \textbf{정의 / 값} \\
\midrule
\endhead

\multicolumn{3}{l}{\textbf{입력 변수 및 매개변수}} \\
\midrule
$R_0$           & 기준 RTO (Base RTO)      & $6.0$\,h (BIA 기반 정적 복구 목표) \\
$R_{\max}$      & 최대 허용 RTO (MTPD)     & $24.0$\,h $= 4R_0$ \\
$O_{\text{raw}}$ & 관측성 원시값 (Observability raw value) & $\in [0,\,1]$, 0=블라인드, 1=완전 가시 \\
$k_O$           & 관측성 가중치 (Observability weight) & $\in [0,\,1]$, 관측성 민감도 \\
$k_U$           & 불확실성 가중치           & $= 1 - k_O$ (상보적 제약) \\
$\kappa$        & 곡률 매개변수 (Curvature parameter) & $> 0$, 시그모이드 곡률 \\
$\kappa^*$      & \jrev{\citet{lee2026drto} 도출 곡률 (본 논문 전제)} & $1.2$ \jrev{(선행연구 C1 환경 $F(\kappa)$ 최대화 결과)} \frev{— \emph{잠정·보수적 실무 기본값}} \\
$U_{\text{raw}}$ & 불확실성 원시값 (Uncertainty raw value) & C2: $\sim N(3,1)$, C3: $\sim 6 \cdot \text{Pareto}(3)$. 무차원 강도 \\
$\alpha$        & 파레토 형상 모수          & $3$ (유한 평균/분산, 무한 첨도) \\
\midrule

\multicolumn{3}{l}{\textbf{함수 및 파생 인자}} \\
\midrule
$S(z)$          & 수정 시그모이드 (Sigmoid function) & $2/(1+e^{-z}) - 1 \in (-1,\,1)$ \\
$z_O$           & 시그모이드 관측성 인자    & $\kappa \cdot k_O \cdot O_{\text{raw}} \geq 0$ \\
$z_U$           & 시그모이드 불확실성 인자   & $\kappa \cdot k_U \cdot R_0 \cdot U_{\text{raw}} / (R_{\max} - R_0) \geq 0$ (무차원) \\
$E_{O,\text{raw}}$ & 관측성 원시 효과 (Raw observability effect) & $-R_0 \cdot k_O \cdot O_{\text{raw}}$ \\
$E_{U,\text{raw}}$ & 불확실성 원시 효과 (Raw uncertainty effect) & $(R_{\max} - R_0) \cdot k_U \cdot U_{\text{raw}}$ \\
$E_{O,\text{bnd}}$ & 관측성 바운딩 조정량 (Bounded obs.\ effect) & $-R_0 \cdot S(z_O) \in (-R_0,\,0]$ \\
$E_{U,\text{bnd}}$ & 불확실성 바운딩 조정량 (Bounded unc.\ effect) & $(R_{\max} - R_0) \cdot S(z_U) \in [0,\,R_{\max}-R_0)$ \\
\midrule

\multicolumn{3}{l}{\textbf{출력 변수}} \\
\midrule
$\DRTO$         & 동적 복구목표시간 (Dynamic RTO) & $R_0 + E_{O,\text{bnd}} + E_{U,\text{bnd}} \in (0,\,R_{\max})$ \\
$\eta$          & 효율성 비율 (Efficiency ratio) & $\DRTO / R_0 \in (0,\,4)$ \\
$\bar{\eta}$    & 효율성 비율 평균              & 각 조건별 $\eta$ 표본 평균 \\
$\text{SD}(\eta)$ & 효율성 비율 표준편차         & 각 조건별 $\eta$ 표본 표준편차 \\
\midrule

\multicolumn{3}{l}{\textbf{최적화 기준 함수}} \\
\midrule
$F(\kappa)$     & 다기준 종합점수 (Composite score) & $f_1 \cdot f_2 \cdot f_3 \cdot f_4$ \\
$f_1(\kappa)$   & 실무 유의성 (Practical significance) & 목표 감소율 15\% 기준 \jrev{(선행연구 C1 환경)} \\
$f_2(\kappa)$   & 비포화 (Non-saturation)  & 포화율 $< 5\%$ \jrev{(선행연구 C1 환경)} \\
$f_3(\kappa)$   & 효과 크기 (Effect size)  & $|d| \geq 0.8$ \jrev{(선행연구 C1 환경)} \\
$f_4(\kappa)$   & 설명력 (Explanatory power) & $R^2 \geq 0.95$ \jrev{(선행연구 C1 환경)} \\
\midrule

\multicolumn{3}{l}{\textbf{통계 기호}} \\
\midrule
$n$             & 표본 크기 (Sample size)   & $10{,}000$ \\
$R^2$           & 결정계수                & 회귀 모형 설명력 \\
$R^2_{(1\mathrm{s})}$ & 1차 단순회귀 결정계수 & $\hat{\eta}^{(1\mathrm{s})}$ 모형의 $R^2$ \\
$R^2_{(1\mathrm{m})}$ & 1차 다중회귀 결정계수 & $\hat{\eta}^{(1\mathrm{m})}$ 모형의 $R^2$ \\
$R^2_{(2\mathrm{q})}$ & 2차교호작용 결정계수 & $\hat{\eta}^{(2\mathrm{q})}$ 모형의 $R^2$ \\
$\Delta R^2$    & 결정계수 변화량         & 회귀 단계 간 $R^2$ 증가분 \\
$d$             & Cohen's $d$            & $(\bar{X}_1 - \bar{X}_2) / s_p$ \\
$\text{VaR}_q$  & 위험가치 (Value-at-Risk) & $F_\eta^{-1}(q/100)$, $\eta$의 $q$-백분위값 \\
$\text{CVaR}_q$ & 조건부 꼬리위험          & $E[\eta \mid \eta > \text{VaR}_q]$ \\
\midrule

\multicolumn{3}{l}{\textbf{모델 조건}} \\
\midrule
C0              & 기저 조건 (Baseline)     & $O_{\text{raw}} = 0$, 정적 RTO ($\eta = 1$) \\
C1              & 관측성 조건 (Observability) & $O_{\text{raw}} > 0$, 단일 채널 \\
C2              & 가우시안 불확실성 조건    & C1 + $U \sim N(3,1)$, 가우시안(경꼬리) \\
C3              & 파레토 불확실성 조건      & C1 + $U \sim 6 \cdot \text{Pareto}(3)$, 파레토(중꼬리) \\
\midrule

\multicolumn{3}{l}{\textbf{가설}} \\
\midrule
H1--\jrev{H3}          & L1 시뮬레이션 검증       & H1: 관측성 효과, H2: 경계 보장, H3: 조건 순서 \\
\jrev{H4}--\jrev{H5}          & L2 회귀 검증             & \jrev{H4}: $R^2_{(2\mathrm{q})} \geq 0.95$, \jrev{H5}: 교호작용 $\Delta R^2$ 유의 \\
\jrev{H6}--\jrev{H7}          & L3 강건성 검증           & \jrev{H6}: 시드 독립성 ($|r| < 0.02$), \jrev{H7}: 다중 시드 강건성 ($10^8$회) \\
\jrev{\pdferr{H8}}              & L4 이중채널 검증         & Tail/Core $k_O$ 감쇠비 $< 1$ (P85.8 분할) \\
\jrev{H9}--\jrev{H10}        & L5 전문가 평가           & \jrev{H9}: 정량 평균 $\geq 4.0/5.0$, \jrev{H10}: 정성 동의율 $\geq 80\%$ \\
\midrule

\multicolumn{3}{l}{\textbf{실무 확장 기호 (Practical Extension)}} \\
\midrule
$O_{\mathrm{APM}}$  & APM 관측성              & $w_1 R_{\mathrm{avail}} + w_2 (1-R_{\mathrm{error}}) + w_3 R_{\mathrm{perf}} \in [0,1]$ \\
$O_{\mathrm{SIEM}}$ & SIEM 관측성             & 보안 이벤트 심각도·대응 상태 기반, $\in [0,1]$ \\
$R_{\mathrm{avail}}$ & 가용률 (Availability rate) & 서비스 가용성 비율, $\in [0,1]$ \\
$R_{\mathrm{error}}$ & 오류율 (Error rate)     & 서비스 오류 비율, $\in [0,1]$. $1-R_{\mathrm{error}}$로 반전 사용 \\
$R_{\mathrm{perf}}$  & 성능 지수 (Performance rate) & 서비스 응답 성능 비율, $\in [0,1]$ \\
$w_i$               & APM 가중치              & $O_{\mathrm{APM}}$ 구성 요소 가중 계수 ($\sum w_i = 1$) \\
$\alpha_{\mathrm{obs}}$ & 관측성 혼합비        & $O(t) = \alpha_{\mathrm{obs}} O_{\mathrm{APM}} + (1-\alpha_{\mathrm{obs}}) O_{\mathrm{SIEM}}$, $\in [0.5, 0.8]$ \\
$T_{\mathrm{actual}}$   & 실제 복구 시간 (Actual recovery time) & $> 0$ (h), 백테스팅 비교 대상 \\
$\DRTO_{\mathrm{final}}$ & 최종 DRTO 예측값      & $(0, R_{\max})$ (h), 장애 시점 산출값 \\
$E$                     & 절대 오차 (Absolute error) & $T_{\mathrm{actual}} - \DRTO_{\mathrm{final}}$ (h) \\
$E_r$                   & 상대 오차 (Relative error) & $E / T_{\mathrm{actual}}$, 무차원 \\
$R_0(t)$                & 시변 기준 RTO (Time-varying base RTO) & 정적 $R_0$의 시계열 확장, $> 0$ (h) \\
$\kappa(t)$             & 시변 곡률 매개변수 (Time-varying curvature) & 정적 $\kappa = 1.2$의 시계열 확장, $> 0$ \\

\bottomrule
\end{longtable}


\AppSection{핵심 수식}
\label{app:key-formulas}

\begin{enumerate}[leftmargin=2em]

\item \textbf{DRTO 개별 시그모이드 바운딩 모델}:
\[
  \DRTO = R_0 + E_{O,\text{bnd}} + E_{U,\text{bnd}}
  = R_0 - R_0 \cdot S(z_O) + (R_{\max} - R_0) \cdot S(z_U)
\]

\item \textbf{수정 시그모이드 함수}:
\[
  S(z) = \frac{2}{1 + e^{-z}} - 1, \quad S(z) \in (-1,\,1), \quad S(0) = 0
\]

\item \textbf{관측성 인자}: $z_O = \kappa \cdot k_O \cdot O_{\text{raw}}$

\item \textbf{불확실성 인자}: $z_U = \kappa \cdot k_U \cdot \dfrac{R_0 \cdot U_{\text{raw}}}{R_{\max} - R_0}$

\item \textbf{다기준 종합점수}:
$F(\kappa) = f_1(\kappa) \cdot f_2(\kappa) \cdot f_3(\kappa) \cdot f_4(\kappa)$,
\quad $F^* = F(1.2) = 0.9997$ \jrev{(\citet{lee2026drto} C1 환경 결과; 본 논문은 \frev{분포 환경 적합성} 종합 명제로 보완 — \secref{sec:ch4_hypothesis_summary} 참조)}

\item \textbf{합의 4변수 간명 모형} (Forward Selection top-4 $\cap$ LASSO):
\begin{itemize}
  \item C2: $\hat{\eta} = 1.225 - 0.291\,k_O + 0.467\,U - 0.433\,(k_O \!\cdot\! U) - 0.562\,(k_O \!\cdot\! O)$\quad ($R^2 = 0.996$)
  \item C3: $\hat{\eta} = 1.574 + 0.240\,U - 1.053\,k_O - 0.002\,U^2 - 0.122\,(k_O \!\cdot\! U)$\quad ($R^2 = 0.839$)
\end{itemize}

\end{enumerate}


\AppSection{핵심 용어}
\label{app:key-terms}

\begin{longtable}{p{3.5cm} p{4cm} p{7cm}}
\toprule
\textbf{한국어} & \textbf{영어} & \textbf{정의 요약} \\
\midrule
\endhead
동적 복구목표시간 & Dynamic RTO & 관측성·불확실성을 반영한 $\DRTO$ \\
관측성 & Observability & 시스템 내부 상태를 외부 출력으로 추정 가능한 정도 \\
불확실성 (가우시안) & Uncertainty (Gaussian) & 정규분포를 따르는 일상적 변동, $U_{\mathrm{raw}} \sim N(3,\,1)$ \\
불확실성 (파레토) & Uncertainty (Pareto) & 파레토 분포를 따르는 극단적 변동, $U_{\mathrm{raw}} \sim 6 \cdot \mathrm{Pareto}(\alpha=3)$ \\
개별 시그모이드 바운딩 & Individual Sigmoid Bounding & 채널별 독립 시그모이드로 물리적 경계 보장 \\
이중채널 감쇠 & Dual-channel Attenuation & 극단적 불확실성(Tail 영역)에서 불확실성 채널이 지배적으로 작용하여 관측성 채널의 복구 단축 효과가 구조적으로 점진적 감쇠되는 현상 ($k_O$ 감쇠비 $0.52\times$) \\
과다예비 & Over-provisioning & 불확실성 과대 추정으로 필요 이상의 복구 자원을 배정 \\
과소예비 & Under-provisioning & 불확실성 과소 추정으로 복구 자원이 부족하게 배정 \\
분할 회귀 & Split Regression & P85.8 기준 Core/Tail 분리 회귀 분석 \\
CDF 교차점 & CDF Crossover & C2=C3 누적분포 교차 지점 (P85.8, $\eta \approx 2.42$) \\
효율성 비율 & Efficiency Ratio ($\eta$) & $\DRTO / R_0$, 기준 대비 비율 \\
불확실성 프리미엄 & Uncertainty Premium & 불확실성에 의한 DRTO 추가 증가분 \\
에스컬레이션 & Escalation & $\eta$ 수준별 단계적 대응 격상 \\
\midrule
\multicolumn{3}{l}{\textbf{산업별 관측성 요소}} \\
\midrule
트랜잭션 모니터링 & Transaction Monitoring & 금융 거래를 실시간 감시하여 이상 징후를 탐지하는 시스템 \\
FDS & Fraud Detection System & 이상거래탐지시스템. 금융 사기·부정거래를 실시간 패턴 분석으로 탐지 \\
EMR & Electronic Medical Record & 전자의무기록. 환자 진료 정보를 전자적으로 관리하는 시스템 \\
MES & Manufacturing Execution System & 제조실행시스템. 생산 현장의 작업 지시·실적·품질을 실시간 관리 \\
SCADA & Supervisory Control and Data Acquisition & 감시제어 데이터수집 시스템. 산업 설비·인프라의 원격 감시·제어 \\
APM & Application Performance Monitoring & 애플리케이션 응답 시간, 오류율 등을 실시간 모니터링하는 시스템 \\
SIEM & Security Information and Event Management & 보안 이벤트 수집·분석·대응 통합 관리. $O_{\mathrm{SIEM}}$ 산출 기반 \\
{IDS} & {Intrusion Detection System} & {침입탐지시스템. 네트워크·호스트의 비정상 활동을 탐지·경고} \\
{APT} & {Advanced Persistent Threat} & {지능형 지속 위협. 장기간 은닉하며 특정 조직을 표적으로 공격하는 고도화된 사이버 위협} \\
\midrule
\multicolumn{3}{l}{{\textbf{업무연속성·재해복구}}} \\
\midrule
{BCM} & {Business Continuity Management} & {업무연속성관리. 중단 상황에서 핵심 업무를 유지·복구하기 위한 관리 체계} \\
{BCP} & {Business Continuity Plan} & {업무연속성계획. 업무 중단 시 복구 절차를 문서화한 계획서} \\
{BCMS} & {Business Continuity Management System} & {업무연속성관리체계. \nrev{ISO 22301:2019}에 따른 BCM 구현·운영 통합 시스템} \\
{BIA} & {Business Impact Analysis} & {업무영향분석. 핵심 업무의 중단 영향도를 분석하여 RTO, RPO 등 복구 목표를 결정하는 프로세스} \\
{DR} & {Disaster Recovery} & {재해복구. 재난 발생 후 IT 시스템 및 업무 기능을 원래 상태로 복원하는 활동} \\
{MTPD} & {Maximum Tolerable Period of Disruption} & {최대허용중단시간. 업무 중단이 허용되는 최대 기간 ($= R_{\max}$)} \\
\midrule
\multicolumn{3}{l}{{\textbf{리스크 관리·금융}}} \\
\midrule
{ERM} & {Enterprise Risk Management} & {전사적 위험관리. 조직 전체 위험을 통합적으로 식별·평가·관리하는 체계} \\
{KRI} & {Key Risk Indicator} & {주요 위험 지표. $\eta$가 임계값 초과 시 경보를 발생시키는 ERM 연계 지표} \\
{VaR} & {Value-at-Risk} & {위험가치. 신뢰수준(예: 99\%)에서 최대 손실액을 나타내는 금융 위험 지표} \\
{CVaR} & {Conditional Value-at-Risk} & {조건부 위험가치. VaR 초과 극단 손실의 평균. 꼬리 위험 평가 지표} \\
{DaR} & {DRTO-at-Risk} & {VaR 프레임을 DRTO에 적용. ``99\% 신뢰수준에서 DRTO가 X시간 이하'' (E3 제안)} \\
{Basel} & {Basel Accord / Basel Committee} & {바젤 위원회 규제. 은행 자기자본·운영 위험 기준. DRTO를 운영위험 정량 지표로 활용 가능} \\
\midrule
\multicolumn{3}{l}{{\textbf{규제·표준}}} \\
\midrule
{\nrev{ISO 22301:2019}} & {ISO 22301:2019} & {업무연속성관리체계 국제표준. BIA, RTO, 복구 전략 등을 규정} \\
{NIST CSF} & {NIST Cybersecurity Framework} & {NIST 사이버보안 프레임워크(v2.0). 식별·보호·탐지·대응·복구 5기능. DRTO는 복구 기능 정량 지표} \\
{DORA} & {Digital Operational Resilience Act} & {EU 디지털 운영 탄력성 규제. 금융기관 사이버 복원력 의무화 (2023년 발효)} \\
\midrule
\multicolumn{3}{l}{{\textbf{사이버 보안·위협 인텔리전스}}} \\
\midrule
{DDoS} & {Distributed Denial of Service} & {분산 서비스 거부 공격. 대량 트래픽으로 서비스 접근을 차단하는 공격 유형} \\
{MITRE ATT\&CK} & {Adversarial Tactics, Techniques \& Common Knowledge} & {사이버 공격 기법·전술 분류 체계. 공격 유형별 $U_{\mathrm{raw}}$ 자동 조정에 활용} \\
{STIX} & {\erf{Structu}{Structured} Threat Information Expression} & {사이버 위협 정보 표준 교환 형식. TAXII와 함께 위협 인텔리전스 자동 연동} \\
{TAXII} & {Trusted Automated Exchange of Intelligence Information} & {STIX 형식 위협 정보의 자동 교환 프로토콜} \\
\midrule
\multicolumn{3}{l}{{\textbf{클라우드·DevOps}}} \\
\midrule
{GitOps} & {Git-based Operations} & {Git 리포지토리를 SSOT로 인프라·서비스 변경을 코드 기반 자동화. DRTO 파라미터 버전 관리 (E4 제안)} \\
{IaC} & {Infrastructure as Code} & {인프라를 코드로 정의·관리 (Terraform, Ansible 등). GitOps와 연동하여 인프라 변경 시 DRTO 영향 자동 평가} \\
{Kubernetes (K8s)} & {Container Orchestration Platform} & {컨테이너 오케스트레이션 플랫폼. Pod·Service·Cluster 단위별 DRTO 파라미터 차별화 필요} \\
{CI/CD} & {Continuous Integration / Continuous Deployment} & {지속적 통합·배포. 코드 변경의 자동 빌드·테스트·배포 파이프라인} \\
{Serverless} & {Serverless Computing} & {서버리스 컴퓨팅. 서버 관리 없이 코드 실행 단위로 자원을 할당하는 클라우드 모델. $R_0$ 정의 방법 차별화 필요} \\
YAML & \vrev{YAML} & 사람이 읽기 쉬운 데이터 직렬화 언어. GitOps에서 DRTO 파라미터 정의에 활용 \\
\midrule
\multicolumn{3}{l}{{\textbf{연구 방법론·통계}}} \\
\midrule
{CVI} & {Content Validity Index} & {내용타당도 지수. I-CVI(문항 수준, $\geq 0.80$), S-CVI/Ave(척도 수준, $\geq 0.90$)} \\
{ICC} & {Intraclass Correlation Coefficient} & {급내상관계수. 다수 평가자 간 일치도 측정 통계량} \\
\bottomrule
\end{longtable}

%% --- 핵심 용어 상세 설명 및 사례 ---




%%==APPB_START==
\chapter{조직 성숙도 자가 진단 체크리스트}
\label{app:eta-dual-use-detail}
\label{app:self-assessment}

\SetAppendixLetter{B}

현업 담당자가 조직의 $\eta$ 수준을 간접적으로 추정할 수 있는
15개 항목 자가 진단 체크리스트를 \p{[}표~\ref{tab:app-self-assessment}\p{]}에 제시한다.

\begin{table}[htbp]
\centering
\caption{BCMS 성숙도 자가 진단 체크리스트 (15개 항목)}
\label{tab:app-self-assessment}
\small
\begin{tabularx}{\textwidth}{c X c}
\toprule
\textbf{No.} & \textbf{진단 항목} & \textbf{영역} \\
\midrule
\multicolumn{3}{l}{\textbf{관측성 역량} (체크 수 $\uparrow$ $\to$ $\eta$ $\downarrow$)} \\
\midrule
1 & 주요 시스템의 실시간 모니터링 도구가 운영 중인가? & $O$ \\
2 & 장애 발생 시 자동 알림(alerting) 체계가 작동하는가? & $O$ \\
3 & MTTR(평균복구시간)을 정기적으로 측정·추적하고 있는가? & $O$ \\
4 & 장애 원인 분석(RCA)을 체계적으로 수행하는가? & $O$ \\
5 & 로그·메트릭·트레이스 중 2개 이상의 관측성 축을 확보했는가? & $O$ \\
6 & 관측성 투자에 대한 정량적 효과 측정을 수행하는가? & $k_O$ \\
\midrule
\multicolumn{3}{l}{\textbf{불확실성 관리} (체크 수 $\uparrow$ $\to$ $\eta$ $\downarrow$)} \\
\midrule
7 & 복구 절차서(runbook)가 문서화되어 최신 상태로 유지되는가? & $U$ \\
8 & 연 2회 이상 BCM 훈련/모의 복구를 실시하는가? & $U$ \\
9 & 외부 공급업체 \erf{장애시}{장애 시} 대체 방안(plan~B)이 수립되어 있는가? & $U$ \\
10 & 극단 시나리오(블랙 스완)에 대한 대응 계획이 있는가? & $U$ \\
11 & 복구 시간 편차를 추적하고 원인별로 분류하는가? & $U$ \\
12 & 사후 분석(post-mortem) 결과를 체계적으로 반영하는가? & $U$ \\
\midrule
\multicolumn{3}{l}{\textbf{거버넌스}} \\
\midrule
13 & BCM 전담 조직 또는 책임자가 지정되어 있는가? & 조직 \\
14 & BCM KPI가 경영진에게 정기적으로 보고되는가? & 조직 \\
15 & BIA(업무영향분석)를 정기적으로 갱신하는가? & 조직 \\
\bottomrule
\end{tabularx}
\end{table}


% 부록 I: 전문가 검증 상세 (기존 부록 H)


%%==APPC_START==
\chapter{전문가 검증 상세 설문 구성}
\label{app:survey-structure}

\SetAppendixLetter{C}

\AppSection{설문 구조 개요}
\label{app:survey-overview}

{전문가 평가 설문은 4개 리커트 영역(Section B--E)과 1개 개방형 영역(Section F),
총 32문항으로 구성되었다(\p{[}표~\ref{tab:app-survey-structure}\p{]}).
정량 분석 대상은 리커트 5점 척도 27문항(B--E)이다.}

\begin{table}[htbp]
\centering
\caption{전문가 평가 설문 구조 (4개 리커트 영역 + 개방형, 32문항)}
\label{tab:app-survey-structure}
\small
\begin{tabular}{c l >{\raggedright\arraybackslash}p{3.8cm} >{\raggedright\arraybackslash}p{4.5cm} r}
\toprule
\textbf{영역} & \textbf{코드} & \textbf{명칭} & \textbf{평가 내용} & \textbf{문항 수} \\
\midrule
Section B & B1--B6  & 모델 타당성      & 이론적 논리성, 구조적 타당성              & 6 \\
Section C & C1--C6  & 파라미터 적정성   & 권장 범위의 실무 적합성                    & 6 \\
{Section D} & {D1--D6}  & {실무 적용성}      & {데이터 수집, BCM 통합, 도입 의향}           & {6} \\
{Section E} & {E1--E9}  & {산업별 적용 가능성 및 $\eta$ 활용 체계} & {산업 범용성, 에스컬레이션, 성숙도, 환류} & {9} \\
\midrule
          &         & {\textbf{리커트 소계}} &                                       & {\textbf{27}} \\
\midrule
{Section F} & {F1--F5}  & {개방형 문항}      & {강점, 장벽, 제안, 보완, 기타 의견}          & {5} \\
\midrule
          &         & {\textbf{전체 합계}}  &                                        & {\textbf{32}} \\
\bottomrule
\end{tabular}
\vspace{0.5em}

\footnotesize
{\noindent 리커트 척도 적용 문항: B1--E9의 27문항. 개방형 문항(F1--F5)은 별도 질적 분석 대상.} \\
\noindent 평가 척도: 1점(전혀 동의하지 않음) -- 5점(매우 동의함).
\normalsize
\end{table}

\AppSection{설문 내용}
\label{app:survey-items}

\AppSubsection{Section B: 모델 타당성 (6문항)}
\label{app:survey-b}

$\DRTO$ 모델의 이론적 구조와 논리적 타당성을 평가하는 6개 문항이다(\p{[}표~\ref{tab:app-survey-items-b}\p{]}).

\begin{table}[htbp]
\centering
\caption{Section B: 모델 타당성 문항}
\label{tab:app-survey-items-b}
\small
\begin{tabularx}{\textwidth}{cX}
\toprule
\textbf{문항} & \textbf{내용} \\
\midrule
B1 & 관측성($O$)이 높아지면 $\DRTO$가 감소한다는 관계의 논리적 타당성 \\
B2 & 불확실성($U$)이 높아지면 $\DRTO$가 증가한다는 관계의 논리적 타당성 \\
B3 & 시그모이드 함수를 활용한 바운딩 메커니즘의 적절성 \\
B4 & C0--C3 조건 구분의 명확성 및 점진적 복잡도 증가의 유용성 \\
B5 & $k_O + k_U = 1$ 제약 조건의 해석 용이성 \\
B6 & $\DRTO$ 프레임워크의 전반적 이론적 타당성 \\
\bottomrule
\end{tabularx}
\end{table}

\AppSubsection{Section C: 파라미터 적정성 (6문항)}
\label{app:survey-c}

$\DRTO$ 모델의 핵심 매개변수 설정값과 실무 조정 가능성을 평가하는 6개 문항이다(\p{[}표~\ref{tab:app-survey-items-c}\p{]}).

\begin{table}[htbp]
\centering
\caption{Section C: 파라미터 적정성 문항}
\label{tab:app-survey-items-c}
\small
\begin{tabularx}{\textwidth}{cX}
\toprule
\textbf{문항} & \textbf{내용} \\
\midrule
C1 & $k_O \in [0.4, 0.6]$ 권장 범위의 실무 적용 가능성 \\
C2 & $k_U = 1 - k_O$ 연동 설계의 파라미터 관리 단순화 기여 \\
C3 & 시그모이드 바운딩 자연 경계 $(0,\; R_{\max})$의 실무 적절성 \\
C4 & 곡률 파라미터 $\kappa = 1.2$ 설정의 적절성 \\
C5 & 파레토 분포($\alpha = 3.0$) 꼬리 위험 반영 적절성 \\
C6 & 파라미터 구조의 직관성 및 조직 맞춤 조정 용이성 \\
\bottomrule
\end{tabularx}
\end{table}

\AppSubsection{Section D: 실무 적용성 (6문항)}
\label{app:survey-d}

현업 환경에서의 데이터 수집 가능성, 기존 프로세스 통합, 도입 의향을 평가하는 6개 문항이다(\p{[}표~\ref{tab:app-survey-items-d}\p{]}).

\begin{table}[htbp]
\centering
\caption{Section D: 실무 적용성 문항}
\label{tab:app-survey-items-d}
\small
\begin{tabularx}{\textwidth}{cX}
\toprule
\textbf{문항} & \textbf{내용} \\
\midrule
D1 & 현업에서 관측성 및 불확실성 데이터 수집 가능성 \\
D2 & 기존 BCM/BIA 프로세스와의 통합 가능성 \\
D3 & $\DRTO$ 결과의 경영진 보고 및 의사결정 활용 적합성 \\
D4 & 위험 기반 의사결정 지원의 유용성 \\
D5 & 단계적 도입(C0 $\to$ C3)을 전제로 한 구현 복잡도 수용성 \\
D6 & 적절한 지원 제공 시 실제 업무 환경 도입 의향 \\
\bottomrule
\end{tabularx}
\end{table}

\AppSubsection{Section E: 산업별 적용 가능성 및 $\eta$ 활용 체계 (9문항)}
\label{app:survey-e}

{$\DRTO$ 모델의 산업별 적용 가능성과 효율성 비율($\eta$) 기반 활용 체계를 평가하는
9개 문항이다(\p{[}표~\ref{tab:app-survey-items-e}\p{]}).}

\begin{table}[htbp]
\centering
\caption{Section E: 산업별 적용 가능성 및 $\eta$ 활용 체계 문항}
\label{tab:app-survey-items-e}
\small
\begin{tabularx}{\textwidth}{cX}
\toprule
\textbf{문항} & \textbf{내용} \\
\midrule
{E1} & {산업별 $R_0$/$R_{\max}$ 설정의 현실성} \\
{E2} & {산업별 관측성 요소의 적절성} \\
{E3} & {산업별 불확실성 요소의 적절성} \\
{E4} & {$\eta$ 기반 에스컬레이션 체계의 활용성} \\
{E5} & {산업 범용성} \\
{E6} & {가이드라인의 실무 도입 지원 효과} \\
{E7} & {$\eta$ 기반 BCMS 성숙도 5단계 모델의 활용성} \\
{E8} & {$\eta$의 평상시·위기 시 이중 활용 체계} \\
{E9} & {진단--향상--대응--환류 선순환 구조} \\
\bottomrule
\end{tabularx}
\end{table}

\AppSubsection{Section F: 개방형 문항 (5문항)}
\label{app:survey-f}

{$\DRTO$ 모델에 대한 자유 서술형 5개 문항이다. 별도 질적 분석 대상으로,
리커트 정량 분석에는 포함되지 않는다.}
\nref{5개 개방형 문항의 상세 내용은 \tabref{tab:app-survey-items-f}에 제시하였다.}

\begin{table}[htbp]
\centering
\caption{Section F: 개방형 문항}
\label{tab:app-survey-items-f}
\small
\begin{tabularx}{\textwidth}{cX}
\toprule
\textbf{문항} & \textbf{내용} \\
\midrule
{F1} & {$\DRTO$ 모델의 가장 큰 강점} \\
{F2} & {실무 도입 시 예상되는 가장 큰 장벽이나 어려움} \\
{F3} & {실무 적용성을 높이기 위해 추가로 필요한 것} \\
{F4} & {모델 보완이 필요한 부분} \\
{F5} & {기타 의견이나 제안 사항} \\
\bottomrule
\end{tabularx}
\end{table}







% 핵심 파라미터: κ = 1.2, R₀ = 6.0h, R_max = 24.0h (= 4R₀)
% 모델: 개별 시그모이드 바운딩, DRTO = R₀ + E_{O,bnd} + E_{U,bnd}
% =====================================================================

% =====================================================================
% 부록 A: 기호·용어·수식 목록
% =====================================================================
\clearpage
% ── 부록 간지 ──
\thispagestyle{plain}
\phantomsection
\addcontentsline{toc}{chapter}{\PlainTOCtext{부록}}
\vspace*{0.35\textheight}
\begin{center}{\ChapterFont\bfseries\fontsize{16}{22}\selectfont 부록}\end{center}
\clearpage

\chapter{기호·용어·수식 목록}
\label{app:notation-list}

\SetAppendixLetter{A}

본 부록은 본 연구에서 사용되는 주요 기호, 용어, 수식을 일람표로 정리한다.

\phantomsection\label{app:notation-overview}
\AppSection{주요 기호}
\label{app:symbols}

[표~A-1]은 주요 기호, 용어, 수식의 일람표이다.

\small
\begin{longtable}{c >{\raggedright\arraybackslash}p{4.8cm} >{\raggedright\arraybackslash}p{8.5cm}}
\toprule
\textbf{기호} & \textbf{명칭} & \textbf{정의 / 값} \\
\midrule
\endhead

\multicolumn{3}{l}{\textbf{입력 변수 및 매개변수}} \\
\midrule
$R_0$           & 기준 RTO (Base RTO)      & $6.0$\,h (BIA 기반 정적 복구 목표) \\
$R_{\max}$      & 최대 허용 RTO (MTPD)     & $24.0$\,h $= 4R_0$ \\
$O_{\text{raw}}$ & 관측성 원시값 (Observability raw value) & $\in [0,\,1]$, 0=블라인드, 1=완전 가시 \\
$k_O$           & 관측성 가중치 (Observability weight) & $\in [0,\,1]$, 관측성 민감도 \\
$k_U$           & 불확실성 가중치           & $= 1 - k_O$ (상보적 제약) \\
$\kappa$        & 곡률 매개변수 (Curvature parameter) & $> 0$, 시그모이드 곡률 \\
$\kappa^*$      & \jrev{\citet{lee2026drto} 도출 곡률 (본 논문 전제)} & $1.2$ \jrev{(선행연구 C1 환경 $F(\kappa)$ 최대화 결과)} \frev{— \emph{잠정·보수적 실무 기본값}} \\
$U_{\text{raw}}$ & 불확실성 원시값 (Uncertainty raw value) & C2: $\sim N(3,1)$, C3: $\sim 6 \cdot \text{Pareto}(3)$. 무차원 강도 \\
$\alpha$        & 파레토 형상 모수          & $3$ (유한 평균/분산, 무한 첨도) \\
\midrule

\multicolumn{3}{l}{\textbf{함수 및 파생 인자}} \\
\midrule
$S(z)$          & 수정 시그모이드 (Sigmoid function) & $2/(1+e^{-z}) - 1 \in (-1,\,1)$ \\
$z_O$           & 시그모이드 관측성 인자    & $\kappa \cdot k_O \cdot O_{\text{raw}} \geq 0$ \\
$z_U$           & 시그모이드 불확실성 인자   & $\kappa \cdot k_U \cdot R_0 \cdot U_{\text{raw}} / (R_{\max} - R_0) \geq 0$ (무차원) \\
$E_{O,\text{raw}}$ & 관측성 원시 효과 (Raw observability effect) & $-R_0 \cdot k_O \cdot O_{\text{raw}}$ \\
$E_{U,\text{raw}}$ & 불확실성 원시 효과 (Raw uncertainty effect) & $(R_{\max} - R_0) \cdot k_U \cdot U_{\text{raw}}$ \\
$E_{O,\text{bnd}}$ & 관측성 바운딩 조정량 (Bounded obs.\ effect) & $-R_0 \cdot S(z_O) \in (-R_0,\,0]$ \\
$E_{U,\text{bnd}}$ & 불확실성 바운딩 조정량 (Bounded unc.\ effect) & $(R_{\max} - R_0) \cdot S(z_U) \in [0,\,R_{\max}-R_0)$ \\
\midrule

\multicolumn{3}{l}{\textbf{출력 변수}} \\
\midrule
$\DRTO$         & 동적 복구목표시간 (Dynamic RTO) & $R_0 + E_{O,\text{bnd}} + E_{U,\text{bnd}} \in (0,\,R_{\max})$ \\
$\eta$          & 효율성 비율 (Efficiency ratio) & $\DRTO / R_0 \in (0,\,4)$ \\
$\bar{\eta}$    & 효율성 비율 평균              & 각 조건별 $\eta$ 표본 평균 \\
$\text{SD}(\eta)$ & 효율성 비율 표준편차         & 각 조건별 $\eta$ 표본 표준편차 \\
\midrule

\multicolumn{3}{l}{\textbf{최적화 기준 함수}} \\
\midrule
$F(\kappa)$     & 다기준 종합점수 (Composite score) & $f_1 \cdot f_2 \cdot f_3 \cdot f_4$ \\
$f_1(\kappa)$   & 실무 유의성 (Practical significance) & 목표 감소율 15\% 기준 \jrev{(선행연구 C1 환경)} \\
$f_2(\kappa)$   & 비포화 (Non-saturation)  & 포화율 $< 5\%$ \jrev{(선행연구 C1 환경)} \\
$f_3(\kappa)$   & 효과 크기 (Effect size)  & $|d| \geq 0.8$ \jrev{(선행연구 C1 환경)} \\
$f_4(\kappa)$   & 설명력 (Explanatory power) & $R^2 \geq 0.95$ \jrev{(선행연구 C1 환경)} \\
\midrule

\multicolumn{3}{l}{\textbf{통계 기호}} \\
\midrule
$n$             & 표본 크기 (Sample size)   & $10{,}000$ \\
$R^2$           & 결정계수                & 회귀 모형 설명력 \\
$R^2_{(1\mathrm{s})}$ & 1차 단순회귀 결정계수 & $\hat{\eta}^{(1\mathrm{s})}$ 모형의 $R^2$ \\
$R^2_{(1\mathrm{m})}$ & 1차 다중회귀 결정계수 & $\hat{\eta}^{(1\mathrm{m})}$ 모형의 $R^2$ \\
$R^2_{(2\mathrm{q})}$ & 2차교호작용 결정계수 & $\hat{\eta}^{(2\mathrm{q})}$ 모형의 $R^2$ \\
$\Delta R^2$    & 결정계수 변화량         & 회귀 단계 간 $R^2$ 증가분 \\
$d$             & Cohen's $d$            & $(\bar{X}_1 - \bar{X}_2) / s_p$ \\
$\text{VaR}_q$  & 위험가치 (Value-at-Risk) & $F_\eta^{-1}(q/100)$, $\eta$의 $q$-백분위값 \\
$\text{CVaR}_q$ & 조건부 꼬리위험          & $E[\eta \mid \eta > \text{VaR}_q]$ \\
\midrule

\multicolumn{3}{l}{\textbf{모델 조건}} \\
\midrule
C0              & 기저 조건 (Baseline)     & $O_{\text{raw}} = 0$, 정적 RTO ($\eta = 1$) \\
C1              & 관측성 조건 (Observability) & $O_{\text{raw}} > 0$, 단일 채널 \\
C2              & 가우시안 불확실성 조건    & C1 + $U \sim N(3,1)$, 가우시안(경꼬리) \\
C3              & 파레토 불확실성 조건      & C1 + $U \sim 6 \cdot \text{Pareto}(3)$, 파레토(중꼬리) \\
\midrule

\multicolumn{3}{l}{\textbf{가설}} \\
\midrule
H1--\jrev{H3}          & L1 시뮬레이션 검증       & H1: 관측성 효과, H2: 경계 보장, H3: 조건 순서 \\
\jrev{H4}--\jrev{H5}          & L2 회귀 검증             & \jrev{H4}: $R^2_{(2\mathrm{q})} \geq 0.95$, \jrev{H5}: 교호작용 $\Delta R^2$ 유의 \\
\jrev{H6}--\jrev{H7}          & L3 강건성 검증           & \jrev{H6}: 시드 독립성 ($|r| < 0.02$), \jrev{H7}: 다중 시드 강건성 ($10^8$회) \\
\jrev{\pdferr{H8}}              & L4 이중채널 검증         & Tail/Core $k_O$ 감쇠비 $< 1$ (P85.8 분할) \\
\jrev{H9}--\jrev{H10}        & L5 전문가 평가           & \jrev{H9}: 정량 평균 $\geq 4.0/5.0$, \jrev{H10}: 정성 동의율 $\geq 80\%$ \\
\midrule

\multicolumn{3}{l}{\textbf{실무 확장 기호 (Practical Extension)}} \\
\midrule
$O_{\mathrm{APM}}$  & APM 관측성              & $w_1 R_{\mathrm{avail}} + w_2 (1-R_{\mathrm{error}}) + w_3 R_{\mathrm{perf}} \in [0,1]$ \\
$O_{\mathrm{SIEM}}$ & SIEM 관측성             & 보안 이벤트 심각도·대응 상태 기반, $\in [0,1]$ \\
$R_{\mathrm{avail}}$ & 가용률 (Availability rate) & 서비스 가용성 비율, $\in [0,1]$ \\
$R_{\mathrm{error}}$ & 오류율 (Error rate)     & 서비스 오류 비율, $\in [0,1]$. $1-R_{\mathrm{error}}$로 반전 사용 \\
$R_{\mathrm{perf}}$  & 성능 지수 (Performance rate) & 서비스 응답 성능 비율, $\in [0,1]$ \\
$w_i$               & APM 가중치              & $O_{\mathrm{APM}}$ 구성 요소 가중 계수 ($\sum w_i = 1$) \\
$\alpha_{\mathrm{obs}}$ & 관측성 혼합비        & $O(t) = \alpha_{\mathrm{obs}} O_{\mathrm{APM}} + (1-\alpha_{\mathrm{obs}}) O_{\mathrm{SIEM}}$, $\in [0.5, 0.8]$ \\
$T_{\mathrm{actual}}$   & 실제 복구 시간 (Actual recovery time) & $> 0$ (h), 백테스팅 비교 대상 \\
$\DRTO_{\mathrm{final}}$ & 최종 DRTO 예측값      & $(0, R_{\max})$ (h), 장애 시점 산출값 \\
$E$                     & 절대 오차 (Absolute error) & $T_{\mathrm{actual}} - \DRTO_{\mathrm{final}}$ (h) \\
$E_r$                   & 상대 오차 (Relative error) & $E / T_{\mathrm{actual}}$, 무차원 \\
$R_0(t)$                & 시변 기준 RTO (Time-varying base RTO) & 정적 $R_0$의 시계열 확장, $> 0$ (h) \\
$\kappa(t)$             & 시변 곡률 매개변수 (Time-varying curvature) & 정적 $\kappa = 1.2$의 시계열 확장, $> 0$ \\

\bottomrule
\end{longtable}


\AppSection{핵심 수식}
\label{app:key-formulas}

\begin{enumerate}[leftmargin=2em]

\item \textbf{DRTO 개별 시그모이드 바운딩 모델}:
\[
  \DRTO = R_0 + E_{O,\text{bnd}} + E_{U,\text{bnd}}
  = R_0 - R_0 \cdot S(z_O) + (R_{\max} - R_0) \cdot S(z_U)
\]

\item \textbf{수정 시그모이드 함수}:
\[
  S(z) = \frac{2}{1 + e^{-z}} - 1, \quad S(z) \in (-1,\,1), \quad S(0) = 0
\]

\item \textbf{관측성 인자}: $z_O = \kappa \cdot k_O \cdot O_{\text{raw}}$

\item \textbf{불확실성 인자}: $z_U = \kappa \cdot k_U \cdot \dfrac{R_0 \cdot U_{\text{raw}}}{R_{\max} - R_0}$

\item \textbf{다기준 종합점수}:
$F(\kappa) = f_1(\kappa) \cdot f_2(\kappa) \cdot f_3(\kappa) \cdot f_4(\kappa)$,
\quad $F^* = F(1.2) = 0.9997$ \jrev{(\citet{lee2026drto} C1 환경 결과; 본 논문은 \frev{분포 환경 적합성} 종합 명제로 보완 — \secref{sec:ch4_hypothesis_summary} 참조)}

\item \textbf{합의 4변수 간명 모형} (Forward Selection top-4 $\cap$ LASSO):
\begin{itemize}
  \item C2: $\hat{\eta} = 1.225 - 0.291\,k_O + 0.467\,U - 0.433\,(k_O \!\cdot\! U) - 0.562\,(k_O \!\cdot\! O)$\quad ($R^2 = 0.996$)
  \item C3: $\hat{\eta} = 1.574 + 0.240\,U - 1.053\,k_O - 0.002\,U^2 - 0.122\,(k_O \!\cdot\! U)$\quad ($R^2 = 0.839$)
\end{itemize}

\end{enumerate}


\AppSection{핵심 용어}
\label{app:key-terms}

\begin{longtable}{p{3.5cm} p{4cm} p{7cm}}
\toprule
\textbf{한국어} & \textbf{영어} & \textbf{정의 요약} \\
\midrule
\endhead
동적 복구목표시간 & Dynamic RTO & 관측성·불확실성을 반영한 $\DRTO$ \\
관측성 & Observability & 시스템 내부 상태를 외부 출력으로 추정 가능한 정도 \\
불확실성 (가우시안) & Uncertainty (Gaussian) & 정규분포를 따르는 일상적 변동, $U_{\mathrm{raw}} \sim N(3,\,1)$ \\
불확실성 (파레토) & Uncertainty (Pareto) & 파레토 분포를 따르는 극단적 변동, $U_{\mathrm{raw}} \sim 6 \cdot \mathrm{Pareto}(\alpha=3)$ \\
개별 시그모이드 바운딩 & Individual Sigmoid Bounding & 채널별 독립 시그모이드로 물리적 경계 보장 \\
이중채널 감쇠 & Dual-channel Attenuation & 극단적 불확실성(Tail 영역)에서 불확실성 채널이 지배적으로 작용하여 관측성 채널의 복구 단축 효과가 구조적으로 점진적 감쇠되는 현상 ($k_O$ 감쇠비 $0.52\times$) \\
과다예비 & Over-provisioning & 불확실성 과대 추정으로 필요 이상의 복구 자원을 배정 \\
과소예비 & Under-provisioning & 불확실성 과소 추정으로 복구 자원이 부족하게 배정 \\
분할 회귀 & Split Regression & P85.8 기준 Core/Tail 분리 회귀 분석 \\
CDF 교차점 & CDF Crossover & C2=C3 누적분포 교차 지점 (P85.8, $\eta \approx 2.42$) \\
효율성 비율 & Efficiency Ratio ($\eta$) & $\DRTO / R_0$, 기준 대비 비율 \\
불확실성 프리미엄 & Uncertainty Premium & 불확실성에 의한 DRTO 추가 증가분 \\
에스컬레이션 & Escalation & $\eta$ 수준별 단계적 대응 격상 \\
\midrule
\multicolumn{3}{l}{\textbf{산업별 관측성 요소}} \\
\midrule
트랜잭션 모니터링 & Transaction Monitoring & 금융 거래를 실시간 감시하여 이상 징후를 탐지하는 시스템 \\
FDS & Fraud Detection System & 이상거래탐지시스템. 금융 사기·부정거래를 실시간 패턴 분석으로 탐지 \\
EMR & Electronic Medical Record & 전자의무기록. 환자 진료 정보를 전자적으로 관리하는 시스템 \\
MES & Manufacturing Execution System & 제조실행시스템. 생산 현장의 작업 지시·실적·품질을 실시간 관리 \\
SCADA & Supervisory Control and Data Acquisition & 감시제어 데이터수집 시스템. 산업 설비·인프라의 원격 감시·제어 \\
APM & Application Performance Monitoring & 애플리케이션 응답 시간, 오류율 등을 실시간 모니터링하는 시스템 \\
SIEM & Security Information and Event Management & 보안 이벤트 수집·분석·대응 통합 관리. $O_{\mathrm{SIEM}}$ 산출 기반 \\
{IDS} & {Intrusion Detection System} & {침입탐지시스템. 네트워크·호스트의 비정상 활동을 탐지·경고} \\
{APT} & {Advanced Persistent Threat} & {지능형 지속 위협. 장기간 은닉하며 특정 조직을 표적으로 공격하는 고도화된 사이버 위협} \\
\midrule
\multicolumn{3}{l}{{\textbf{업무연속성·재해복구}}} \\
\midrule
{BCM} & {Business Continuity Management} & {업무연속성관리. 중단 상황에서 핵심 업무를 유지·복구하기 위한 관리 체계} \\
{BCP} & {Business Continuity Plan} & {업무연속성계획. 업무 중단 시 복구 절차를 문서화한 계획서} \\
{BCMS} & {Business Continuity Management System} & {업무연속성관리체계. \nrev{ISO 22301:2019}에 따른 BCM 구현·운영 통합 시스템} \\
{BIA} & {Business Impact Analysis} & {업무영향분석. 핵심 업무의 중단 영향도를 분석하여 RTO, RPO 등 복구 목표를 결정하는 프로세스} \\
{DR} & {Disaster Recovery} & {재해복구. 재난 발생 후 IT 시스템 및 업무 기능을 원래 상태로 복원하는 활동} \\
{MTPD} & {Maximum Tolerable Period of Disruption} & {최대허용중단시간. 업무 중단이 허용되는 최대 기간 ($= R_{\max}$)} \\
\midrule
\multicolumn{3}{l}{{\textbf{리스크 관리·금융}}} \\
\midrule
{ERM} & {Enterprise Risk Management} & {전사적 위험관리. 조직 전체 위험을 통합적으로 식별·평가·관리하는 체계} \\
{KRI} & {Key Risk Indicator} & {주요 위험 지표. $\eta$가 임계값 초과 시 경보를 발생시키는 ERM 연계 지표} \\
{VaR} & {Value-at-Risk} & {위험가치. 신뢰수준(예: 99\%)에서 최대 손실액을 나타내는 금융 위험 지표} \\
{CVaR} & {Conditional Value-at-Risk} & {조건부 위험가치. VaR 초과 극단 손실의 평균. 꼬리 위험 평가 지표} \\
{DaR} & {DRTO-at-Risk} & {VaR 프레임을 DRTO에 적용. ``99\% 신뢰수준에서 DRTO가 X시간 이하'' (E3 제안)} \\
{Basel} & {Basel Accord / Basel Committee} & {바젤 위원회 규제. 은행 자기자본·운영 위험 기준. DRTO를 운영위험 정량 지표로 활용 가능} \\
\midrule
\multicolumn{3}{l}{{\textbf{규제·표준}}} \\
\midrule
{\nrev{ISO 22301:2019}} & {ISO 22301:2019} & {업무연속성관리체계 국제표준. BIA, RTO, 복구 전략 등을 규정} \\
{NIST CSF} & {NIST Cybersecurity Framework} & {NIST 사이버보안 프레임워크(v2.0). 식별·보호·탐지·대응·복구 5기능. DRTO는 복구 기능 정량 지표} \\
{DORA} & {Digital Operational Resilience Act} & {EU 디지털 운영 탄력성 규제. 금융기관 사이버 복원력 의무화 (2023년 발효)} \\
\midrule
\multicolumn{3}{l}{{\textbf{사이버 보안·위협 인텔리전스}}} \\
\midrule
{DDoS} & {Distributed Denial of Service} & {분산 서비스 거부 공격. 대량 트래픽으로 서비스 접근을 차단하는 공격 유형} \\
{MITRE ATT\&CK} & {Adversarial Tactics, Techniques \& Common Knowledge} & {사이버 공격 기법·전술 분류 체계. 공격 유형별 $U_{\mathrm{raw}}$ 자동 조정에 활용} \\
{STIX} & {\erf{Structu}{Structured} Threat Information Expression} & {사이버 위협 정보 표준 교환 형식. TAXII와 함께 위협 인텔리전스 자동 연동} \\
{TAXII} & {Trusted Automated Exchange of Intelligence Information} & {STIX 형식 위협 정보의 자동 교환 프로토콜} \\
\midrule
\multicolumn{3}{l}{{\textbf{클라우드·DevOps}}} \\
\midrule
{GitOps} & {Git-based Operations} & {Git 리포지토리를 SSOT로 인프라·서비스 변경을 코드 기반 자동화. DRTO 파라미터 버전 관리 (E4 제안)} \\
{IaC} & {Infrastructure as Code} & {인프라를 코드로 정의·관리 (Terraform, Ansible 등). GitOps와 연동하여 인프라 변경 시 DRTO 영향 자동 평가} \\
{Kubernetes (K8s)} & {Container Orchestration Platform} & {컨테이너 오케스트레이션 플랫폼. Pod·Service·Cluster 단위별 DRTO 파라미터 차별화 필요} \\
{CI/CD} & {Continuous Integration / Continuous Deployment} & {지속적 통합·배포. 코드 변경의 자동 빌드·테스트·배포 파이프라인} \\
{Serverless} & {Serverless Computing} & {서버리스 컴퓨팅. 서버 관리 없이 코드 실행 단위로 자원을 할당하는 클라우드 모델. $R_0$ 정의 방법 차별화 필요} \\
YAML & \vrev{YAML} & 사람이 읽기 쉬운 데이터 직렬화 언어. GitOps에서 DRTO 파라미터 정의에 활용 \\
\midrule
\multicolumn{3}{l}{{\textbf{연구 방법론·통계}}} \\
\midrule
{CVI} & {Content Validity Index} & {내용타당도 지수. I-CVI(문항 수준, $\geq 0.80$), S-CVI/Ave(척도 수준, $\geq 0.90$)} \\
{ICC} & {Intraclass Correlation Coefficient} & {급내상관계수. 다수 평가자 간 일치도 측정 통계량} \\
\bottomrule
\end{longtable}

%% --- 핵심 용어 상세 설명 및 사례 ---




%%==APPB_START==
\chapter{조직 성숙도 자가 진단 체크리스트}
\label{app:eta-dual-use-detail}
\label{app:self-assessment}

\SetAppendixLetter{B}

현업 담당자가 조직의 $\eta$ 수준을 간접적으로 추정할 수 있는
15개 항목 자가 진단 체크리스트를 \p{[}표~\ref{tab:app-self-assessment}\p{]}에 제시한다.

\begin{table}[htbp]
\centering
\caption{BCMS 성숙도 자가 진단 체크리스트 (15개 항목)}
\label{tab:app-self-assessment}
\small
\begin{tabularx}{\textwidth}{c X c}
\toprule
\textbf{No.} & \textbf{진단 항목} & \textbf{영역} \\
\midrule
\multicolumn{3}{l}{\textbf{관측성 역량} (체크 수 $\uparrow$ $\to$ $\eta$ $\downarrow$)} \\
\midrule
1 & 주요 시스템의 실시간 모니터링 도구가 운영 중인가? & $O$ \\
2 & 장애 발생 시 자동 알림(alerting) 체계가 작동하는가? & $O$ \\
3 & MTTR(평균복구시간)을 정기적으로 측정·추적하고 있는가? & $O$ \\
4 & 장애 원인 분석(RCA)을 체계적으로 수행하는가? & $O$ \\
5 & 로그·메트릭·트레이스 중 2개 이상의 관측성 축을 확보했는가? & $O$ \\
6 & 관측성 투자에 대한 정량적 효과 측정을 수행하는가? & $k_O$ \\
\midrule
\multicolumn{3}{l}{\textbf{불확실성 관리} (체크 수 $\uparrow$ $\to$ $\eta$ $\downarrow$)} \\
\midrule
7 & 복구 절차서(runbook)가 문서화되어 최신 상태로 유지되는가? & $U$ \\
8 & 연 2회 이상 BCM 훈련/모의 복구를 실시하는가? & $U$ \\
9 & 외부 공급업체 \erf{장애시}{장애 시} 대체 방안(plan~B)이 수립되어 있는가? & $U$ \\
10 & 극단 시나리오(블랙 스완)에 대한 대응 계획이 있는가? & $U$ \\
11 & 복구 시간 편차를 추적하고 원인별로 분류하는가? & $U$ \\
12 & 사후 분석(post-mortem) 결과를 체계적으로 반영하는가? & $U$ \\
\midrule
\multicolumn{3}{l}{\textbf{거버넌스}} \\
\midrule
13 & BCM 전담 조직 또는 책임자가 지정되어 있는가? & 조직 \\
14 & BCM KPI가 경영진에게 정기적으로 보고되는가? & 조직 \\
15 & BIA(업무영향분석)를 정기적으로 갱신하는가? & 조직 \\
\bottomrule
\end{tabularx}
\end{table}


% 부록 I: 전문가 검증 상세 (기존 부록 H)


%%==APPC_START==
\chapter{전문가 검증 상세 설문 구성}
\label{app:survey-structure}

\SetAppendixLetter{C}

\AppSection{설문 구조 개요}
\label{app:survey-overview}

{전문가 평가 설문은 4개 리커트 영역(Section B--E)과 1개 개방형 영역(Section F),
총 32문항으로 구성되었다(\p{[}표~\ref{tab:app-survey-structure}\p{]}).
정량 분석 대상은 리커트 5점 척도 27문항(B--E)이다.}

\begin{table}[htbp]
\centering
\caption{전문가 평가 설문 구조 (4개 리커트 영역 + 개방형, 32문항)}
\label{tab:app-survey-structure}
\small
\begin{tabular}{c l >{\raggedright\arraybackslash}p{3.8cm} >{\raggedright\arraybackslash}p{4.5cm} r}
\toprule
\textbf{영역} & \textbf{코드} & \textbf{명칭} & \textbf{평가 내용} & \textbf{문항 수} \\
\midrule
Section B & B1--B6  & 모델 타당성      & 이론적 논리성, 구조적 타당성              & 6 \\
Section C & C1--C6  & 파라미터 적정성   & 권장 범위의 실무 적합성                    & 6 \\
{Section D} & {D1--D6}  & {실무 적용성}      & {데이터 수집, BCM 통합, 도입 의향}           & {6} \\
{Section E} & {E1--E9}  & {산업별 적용 가능성 및 $\eta$ 활용 체계} & {산업 범용성, 에스컬레이션, 성숙도, 환류} & {9} \\
\midrule
          &         & {\textbf{리커트 소계}} &                                       & {\textbf{27}} \\
\midrule
{Section F} & {F1--F5}  & {개방형 문항}      & {강점, 장벽, 제안, 보완, 기타 의견}          & {5} \\
\midrule
          &         & {\textbf{전체 합계}}  &                                        & {\textbf{32}} \\
\bottomrule
\end{tabular}
\vspace{0.5em}

\footnotesize
{\noindent 리커트 척도 적용 문항: B1--E9의 27문항. 개방형 문항(F1--F5)은 별도 질적 분석 대상.} \\
\noindent 평가 척도: 1점(전혀 동의하지 않음) -- 5점(매우 동의함).
\normalsize
\end{table}

\AppSection{설문 내용}
\label{app:survey-items}

\AppSubsection{Section B: 모델 타당성 (6문항)}
\label{app:survey-b}

$\DRTO$ 모델의 이론적 구조와 논리적 타당성을 평가하는 6개 문항이다(\p{[}표~\ref{tab:app-survey-items-b}\p{]}).

\begin{table}[htbp]
\centering
\caption{Section B: 모델 타당성 문항}
\label{tab:app-survey-items-b}
\small
\begin{tabularx}{\textwidth}{cX}
\toprule
\textbf{문항} & \textbf{내용} \\
\midrule
B1 & 관측성($O$)이 높아지면 $\DRTO$가 감소한다는 관계의 논리적 타당성 \\
B2 & 불확실성($U$)이 높아지면 $\DRTO$가 증가한다는 관계의 논리적 타당성 \\
B3 & 시그모이드 함수를 활용한 바운딩 메커니즘의 적절성 \\
B4 & C0--C3 조건 구분의 명확성 및 점진적 복잡도 증가의 유용성 \\
B5 & $k_O + k_U = 1$ 제약 조건의 해석 용이성 \\
B6 & $\DRTO$ 프레임워크의 전반적 이론적 타당성 \\
\bottomrule
\end{tabularx}
\end{table}

\AppSubsection{Section C: 파라미터 적정성 (6문항)}
\label{app:survey-c}

$\DRTO$ 모델의 핵심 매개변수 설정값과 실무 조정 가능성을 평가하는 6개 문항이다(\p{[}표~\ref{tab:app-survey-items-c}\p{]}).

\begin{table}[htbp]
\centering
\caption{Section C: 파라미터 적정성 문항}
\label{tab:app-survey-items-c}
\small
\begin{tabularx}{\textwidth}{cX}
\toprule
\textbf{문항} & \textbf{내용} \\
\midrule
C1 & $k_O \in [0.4, 0.6]$ 권장 범위의 실무 적용 가능성 \\
C2 & $k_U = 1 - k_O$ 연동 설계의 파라미터 관리 단순화 기여 \\
C3 & 시그모이드 바운딩 자연 경계 $(0,\; R_{\max})$의 실무 적절성 \\
C4 & 곡률 파라미터 $\kappa = 1.2$ 설정의 적절성 \\
C5 & 파레토 분포($\alpha = 3.0$) 꼬리 위험 반영 적절성 \\
C6 & 파라미터 구조의 직관성 및 조직 맞춤 조정 용이성 \\
\bottomrule
\end{tabularx}
\end{table}

\AppSubsection{Section D: 실무 적용성 (6문항)}
\label{app:survey-d}

현업 환경에서의 데이터 수집 가능성, 기존 프로세스 통합, 도입 의향을 평가하는 6개 문항이다(\p{[}표~\ref{tab:app-survey-items-d}\p{]}).

\begin{table}[htbp]
\centering
\caption{Section D: 실무 적용성 문항}
\label{tab:app-survey-items-d}
\small
\begin{tabularx}{\textwidth}{cX}
\toprule
\textbf{문항} & \textbf{내용} \\
\midrule
D1 & 현업에서 관측성 및 불확실성 데이터 수집 가능성 \\
D2 & 기존 BCM/BIA 프로세스와의 통합 가능성 \\
D3 & $\DRTO$ 결과의 경영진 보고 및 의사결정 활용 적합성 \\
D4 & 위험 기반 의사결정 지원의 유용성 \\
D5 & 단계적 도입(C0 $\to$ C3)을 전제로 한 구현 복잡도 수용성 \\
D6 & 적절한 지원 제공 시 실제 업무 환경 도입 의향 \\
\bottomrule
\end{tabularx}
\end{table}

\AppSubsection{Section E: 산업별 적용 가능성 및 $\eta$ 활용 체계 (9문항)}
\label{app:survey-e}

{$\DRTO$ 모델의 산업별 적용 가능성과 효율성 비율($\eta$) 기반 활용 체계를 평가하는
9개 문항이다(\p{[}표~\ref{tab:app-survey-items-e}\p{]}).}

\begin{table}[htbp]
\centering
\caption{Section E: 산업별 적용 가능성 및 $\eta$ 활용 체계 문항}
\label{tab:app-survey-items-e}
\small
\begin{tabularx}{\textwidth}{cX}
\toprule
\textbf{문항} & \textbf{내용} \\
\midrule
{E1} & {산업별 $R_0$/$R_{\max}$ 설정의 현실성} \\
{E2} & {산업별 관측성 요소의 적절성} \\
{E3} & {산업별 불확실성 요소의 적절성} \\
{E4} & {$\eta$ 기반 에스컬레이션 체계의 활용성} \\
{E5} & {산업 범용성} \\
{E6} & {가이드라인의 실무 도입 지원 효과} \\
{E7} & {$\eta$ 기반 BCMS 성숙도 5단계 모델의 활용성} \\
{E8} & {$\eta$의 평상시·위기 시 이중 활용 체계} \\
{E9} & {진단--향상--대응--환류 선순환 구조} \\
\bottomrule
\end{tabularx}
\end{table}

\AppSubsection{Section F: 개방형 문항 (5문항)}
\label{app:survey-f}

{$\DRTO$ 모델에 대한 자유 서술형 5개 문항이다. 별도 질적 분석 대상으로,
리커트 정량 분석에는 포함되지 않는다.}
\nref{5개 개방형 문항의 상세 내용은 \tabref{tab:app-survey-items-f}에 제시하였다.}

\begin{table}[htbp]
\centering
\caption{Section F: 개방형 문항}
\label{tab:app-survey-items-f}
\small
\begin{tabularx}{\textwidth}{cX}
\toprule
\textbf{문항} & \textbf{내용} \\
\midrule
{F1} & {$\DRTO$ 모델의 가장 큰 강점} \\
{F2} & {실무 도입 시 예상되는 가장 큰 장벽이나 어려움} \\
{F3} & {실무 적용성을 높이기 위해 추가로 필요한 것} \\
{F4} & {모델 보완이 필요한 부분} \\
{F5} & {기타 의견이나 제안 사항} \\
\bottomrule
\end{tabularx}
\end{table}








\end{document}
